\documentclass[12pt, ukrainian]{article}
\setlength\textwidth{190mm} \setlength\hoffset{-25mm} 
\setlength\textheight{277mm} \setlength\voffset{-28mm} 
%\setlength\textwidth{190mm}%\setlength\voffset{-38mm} \setlength\hoffset{-38mm}
%\setlength\parindent{0mm}
\pagestyle{empty}
\usepackage[T2A]{fontenc}
\usepackage[utf8]{inputenc}
\usepackage[ukrainian]{babel}
\usepackage{graphicx}
\usepackage{caption}
\usepackage{verbatim, listings, clrscode3e}
\usepackage{fancyvrb}
\DeclareCaptionFormat{nocaption}{}
\captionsetup{format=nocaption,aboveskip=0pt,belowskip=0pt}
\usepackage[Export]{adjustbox} 
\adjustboxset{max size={0.9\linewidth}{0.9\paperheight}}
\def\code#1{\Verb!#1!}
\begin{document}
%begin

{{
\begin {large}

\newpage
1. Що буде надруковано за виконання?

\begin{verbatim}
print('R', end=' ')
c=('a','b','c','d','e',0,1,2,3,4)
print(c[-1])
\end{verbatim}

%enditem


\newpage
2. Що буде надруковано за виконання?

\begin{verbatim}
print('R', end=' ')
c='0123456789'
print(c[14:13:3])
\end{verbatim}

%enditem



\newpage
3. Що буде надруковано за виконання?

\begin{verbatim}
print('R', end=' ')
f=[10,11,12,13,14]
print(f.pop(0), end=' ')
print(f)
\end{verbatim}

%enditem



\newpage
4. Що буде надруковано за виконання?

\begin{verbatim}
print('R', end=' ')
d=[10,11,12,13,14]
d[7:]=[2,3]
print(d)
\end{verbatim}

%enditem



\newpage
5. Укажіть ТИП значення виразу.\par
\begin{verbatim}
{1:3, 5:8}
\end{verbatim}
%enditem


\newpage
6. Що буде надруковано за виконання?\par
\begin{verbatim}
try:
    d = {51:5, 44:8, 35:8, 44:0}
    d[44] = 5
    for x in d :
        print(x, sep=' ', end=' ')
except:
    print('ERR')
\end{verbatim}
%enditem


\newpage
7. Нехай змінна \kw{s} позначає дуже довгу послідовність цілих, по якій можна ітерувати. 

Написати вираз-генератор, що задає підпослідовність її елементів, що є точними квадратами. 

Обмеження: усі елементи шуканої підпослідовності одночасно в пам'яті розміщені бути не можуть. 
%enditem


\newpage
8. Що буде надруковано за виконання?
\begin{verbatim}
d = 0
k = 1

class C:
    d = 2
    
    def __init__(self):
        self.d = 3
        d = 4
        C.d = 5

obj = C()
try:
    print(d, end=' ')
    print(k, end=' ')
    print(C.d, end=' ')
    print(obj.d)
except:
    print('ERR')
\end{verbatim}
%enditem


\newpage
9. Що буде надруковано за виконання?
\begin{verbatim}
class A0:
    def a(self): print(1, end=' ')
    def g(self): print(2, end=' ')

class A1(A0):
    def g(self): print(3, end=' ')

try:
    A0().a()
    A1().g()
except:
    print('ERR')
\end{verbatim}
%enditem


\newpage
10. Що буде надруковано за виконання?
\begin{verbatim}
class A:
    c = 1

    def _g1(self): print(2, end=' ')
    def g2(self): print(3, end=' ')

    def main(self):
        try: print(self.c, end=' ')
        except: print('E1', end=' ')

        try: self._g1()
        except: print('E2', end=' ')

        try: self.g2()
        except: print('E3', end=' ')


class B(A):
    c = 4

    def _g1(self): print(5, end=' ')
    def g2(self): print(6, end=' ')

try: B().main()
except: print('E4', end=' ')
\end{verbatim}
%enditem


\newpage
11. Змінні \kw{next} та \kw{prev} вузла списку позначають наступний та попередній за ним вузол відповідно. Починаючи з голови, у список записано послідовні цілі числа від~$-70$ до~$46$. Нехай змінна \kw{p} позначає вузол, що зберігає значення~$9$.
Яке значення зберігається у вузлі \kw{p.prev.next.prev.prev.next} ?
%enditem


\newpage
12. 
\begin{center}
	\adjustimage{max size={0.8\linewidth}{0.8\paperheight}}
	{BinTree/trees_exp/150.png}
\end{center}
Указати послідовність міток вершин дерева, зображеного на рис., що утвориться в результаті обходу дерева в оберненому порядку. Мітки записувати через пробіл.\par Алгоритм починає роботу з кореня (на рис. це верхня вершина).
%enditem


\newpage
13. 
\begin{center}
	\adjustimage{max size={0.9\linewidth}{0.9\paperheight}}
	{Graphs/AdjList/324.png}
\end{center}
Для графа, зображеного на рис., вказати список суміжності вершини 0.\par
Формат оформлення відповіді:\par
список суміжності виписувати через пробіл на одному рядку, якщо пробілів десь буде більше одного - не важливо, упорядкування вершин у списку також значення не має.
%enditem


\newpage
14. Для графа на рисунку вказати послідовність, в якій алгоритм пошуку в глибину буде відмічати відвідані вершини, якщо списки суміжності вершин переглядаються алгоритмом за зростанням міток вершин та вершини відмічаються на першому вход\-женні у вершину.\par
Пошук розпочинається з вершини 0.\par

\begin{center}
	\adjustimage{max size={0.9\linewidth}{0.9\paperheight}}
	{Graphs/Traversals/219.png}
\end{center}
%enditem


\newpage
15. Для графа на рисунку з вершини 2 запущено алгоритм пошуку в ширину, що будує кістякове дерево. Списки суміжності вершин переглядаються алгоритмом за зростанням міток вершин.\par
\begin{center}
	\adjustimage{max size={0.9\linewidth}{0.9\paperheight}}
	{Graphs/Traversals/161.png}
\end{center}
\par Які з наведених ребер входять до складу отриманого кістякового дерева?\par
1) (2, 5)\par
2) (1, 5)\par
3) (3, 6)\par
4) (0, 4)\par
5) (1, 6)\par
6) (4, 5)\par
{\vskip 15pt} Номери відповідних ребер записати за зростанням через пробіл.


%enditem


\newpage
16. 

Рядки журналу через двокрапку містять ім'я користувача (ідентифікатор), час події,тип події, код події, наприклад
\begin{verbatim}
s="username : 2022-05-05 13:26 : ERROR : 404"
\end{verbatim}
у коді 
\begin{verbatim}
res = re.fullmatch('???', s) 
s = ''
code_str = ??? 
\end{verbatim}
чим треба замінити кожен з ???, щоб значенням code\_str був код події? 


\vspace{10mm}
\textbf{Завдання 17} 

Задано програму, що містить код 1-2 класів, можливе успадкування, та клієнтський код.
У наведеному коді невеликі частини закриті. 
Також наведено вихід програми.
Необхідно відновити закриті частини, щоб отриманий код відповідав наведеному виходу програми.


\textbf{Завдання 18-19} 

Кілька запитань за завданнями для самостійної роботи.
\vspace{10mm}
Скоріш за все, що завдань буде трохи менше, а деякі попередні будуть переформульовані в стилі двох останніх.

%enditem

\end {large}}

%end

\newpage
%begin

{{
\begin {large}

\newpage
1. Що буде надруковано за виконання?

\begin{verbatim}
print('R', end=' ')
e=('a','b','c','d','e',0,1)
print(e[-4])
\end{verbatim}

%enditem


\newpage
2. Що буде надруковано за виконання?

\begin{verbatim}
print('R', end=' ')
h=['0','1','2','3','4','5','6','7','8','9']
print(h[::-3])
\end{verbatim}

%enditem



\newpage
3. Що буде надруковано за виконання?

\begin{verbatim}
print('R', end=' ')
d=['0','1','2','3','4']
d.append('11')
print(d)
\end{verbatim}

%enditem



\newpage
4. Що буде надруковано за виконання?

\begin{verbatim}
print('R', end=' ')
h=['0','1','2','3','4']
h[:3]='13'
print(h)
\end{verbatim}

%enditem



\newpage
5. Укажіть ТИП значення виразу.\par
\begin{verbatim}
{y:13 for y in range(-7,7)}
\end{verbatim}
%enditem


\newpage
6. Що буде надруковано за виконання?\par
\begin{verbatim}
try:
    t = {81:9, 42:5, 32:0, 81:1}
    t[67] = 1
    for x in t.items():
        print(x, sep=' ', end=' ')
except:
    print('ERR')
\end{verbatim}
%enditem


\newpage
7. Нехай змінна \kw{s} позначає послідовність цілих, по якій можна ітерувати. 

Написати вираз, що перевіряє, чи правда, що послідовність  містить хоча б одне додатнє число.
%enditem


\newpage
8. Що буде надруковано за виконання?
%Оригинал был утерян. Требует отладки!
\begin{verbatim}
__c = 0

class D:
    __c = 1
    
    def __init__(self):
        self.__c = 2
        __c = 3
        D.__c = 4

obj = D()
D.__c = 5
try:
    print(__c, end=' ')
    print(obj.__c, end=' ')
    print(D.__c, end=' ')
    print(obj.__c, end=' ')
    print(obj._D__c)
except:
    print('ERR')
\end{verbatim}
%enditem


\newpage
9. Що буде надруковано за виконання?
\begin{verbatim}
class D0:
    def b(self): print(1, end=' ')
    def h(self): print(2, end=' ')

class D1(D0):
    def h(self): print(3, end=' ')

try:
    D1().b()
    D0().h()
except:
    print('ERR')
\end{verbatim}
%enditem


\newpage
10. Що буде надруковано за виконання?
\begin{verbatim}
class A:
    a = 1

    def __init__(self):
        c = 2
        a = 3

class B(A):
    c = 4
    b = 5

    def __init__(self):
        super().__init__()
        self.b = 6

try:
    obj = B()
    obj.a = obj.b + 10
    B.a = 13
except: print('EE', end=' ')

try: print(obj.a, end= ' ')
except: print('E1', end=' ')

try: print(A.a, end= ' ')
except: print('E2', end=' ')

try: print(B.a, end= ' ')
except: print('E3', end=' ')
\end{verbatim}
%enditem


\newpage
11. Змінні \kw{next} та \kw{prev} вузла списку позначають наступний та попередній за ним вузол відповідно. Починаючи з голови, у список записано послідовні цілі числа від~$-60$ до~$47$. Нехай змінна \kw{p} позначає вузол, що зберігає значення~$-5$.
Яке значення зберігається у вузлі \kw{p.next.prev.next.next.prev} ?
%enditem


\newpage
12. 
\begin{center}
	\adjustimage{max size={0.8\linewidth}{0.8\paperheight}}
	{BinTree/trees_exp/42.png}
\end{center}
Указати послідовність міток вершин дерева, зображеного на рис., що утвориться в результаті обходу дерева в прямому порядку. Мітки записувати через пробіл.\par Алгоритм починає роботу з кореня (на рис. це верхня вершина).
%enditem


\newpage
13. 
\begin{center}
	\adjustimage{max size={0.9\linewidth}{0.9\paperheight}}
	{Graphs/AdjList/340.png}
\end{center}
Для графа, зображеного на рис., вказати список суміжності вершини 1.\par
Формат оформлення відповіді:\par
список суміжності виписувати через пробіл на одному рядку, якщо пробілів десь буде більше одного - не важливо, упорядкування вершин у списку також значення не має.
%enditem


\newpage
14. Для графа на рисунку з вершини 6 запущено алгоритм пошуку в глибину, що будує кістякове дерево. Списки суміжності вершин переглядаються алгоритмом за зростанням міток вершин.\par
\begin{center}
	\adjustimage{max size={0.9\linewidth}{0.9\paperheight}}
	{Graphs/Traversals/312.png}
\end{center}
\par Які з наведених ребер входять до складу отриманого кістякового дерева?\par
1) (0, 4)\par
2) (1, 3)\par
3) (0, 1)\par
4) (4, 5)\par
5) (4, 6)\par
6) (0, 5)\par
{\vskip 15pt} Номери відповідних ребер записати за зростанням через пробіл.

%enditem


\newpage
15. Для графа на рисунку з вершини 3 запущено алгоритм пошуку в ширину, що будує кістякове дерево. Списки суміжності вершин переглядаються алгоритмом за зростанням міток вершин.\par
\begin{center}
	\adjustimage{max size={0.9\linewidth}{0.9\paperheight}}
	{Graphs/Traversals/101.png}
\end{center}
\par Які з наведених ребер входять до складу отриманого кістякового дерева?\par
1) (0, 1)\par
2) (1, 5)\par
3) (0, 5)\par
4) (1, 3)\par
5) (1, 4)\par
6) (3, 5)\par
{\vskip 15pt} Номери відповідних ребер записати за зростанням через пробіл.


%enditem


\newpage
16. 

Рядки журналу через двокрапку містять ім'я користувача (ідентифікатор), час події,тип події, код події, наприклад
\begin{verbatim}
s="username : 2022-05-05 13:26 : ERROR : 404"
\end{verbatim}
у коді 
\begin{verbatim}
res = re.fullmatch('???', s) 
s = ''
code_str = ??? 
\end{verbatim}
чим треба замінити кожен з ???, щоб значенням code\_str був код події? 


\vspace{10mm}
\textbf{Завдання 17} 

Задано програму, що містить код 1-2 класів, можливе успадкування, та клієнтський код.
У наведеному коді невеликі частини закриті. 
Також наведено вихід програми.
Необхідно відновити закриті частини, щоб отриманий код відповідав наведеному виходу програми.


\textbf{Завдання 18-19} 

Кілька запитань за завданнями для самостійної роботи.
\vspace{10mm}
Скоріш за все, що завдань буде трохи менше, а деякі попередні будуть переформульовані в стилі двох останніх.

%enditem

\end {large}}

%end

\newpage
%begin

{{
\begin {large}

\newpage
1. Що буде надруковано за виконання?

\begin{verbatim}
print('R', end=' ')
h=('a','b','c','d','e',0,1,2,4)
print(h[0])
\end{verbatim}

%enditem


\newpage
2. Що буде надруковано за виконання?

\begin{verbatim}
print('R', end=' ')
g=(0,1,2,3,4,5,6,7,8,9)
print(g[12:14:-2])
\end{verbatim}

%enditem



\newpage
3. Що буде надруковано за виконання?

\begin{verbatim}
print('R', end=' ')
g=['0','1','2','3','4']
g+=[4,5]
print(g)
\end{verbatim}

%enditem



\newpage
4. Що буде надруковано за виконання?

\begin{verbatim}
print('R', end=' ')
b=[10,11,12,13,14]
del b[-9:-8]
print(b)
\end{verbatim}

%enditem



\newpage
5. Укажіть ТИП значення виразу.\par
\begin{verbatim}
{y:None for y in range(-8,3)}
\end{verbatim}
%enditem


\newpage
6. Що буде надруковано за виконання?\par
\begin{verbatim}
try:
    s = {50:3, 70:4, 87:3, 36:2, 87:6}
    s[87] = 2
    for x in s.values():
        print(x, sep=' ', end=' ')
except:
    print('ERR')
\end{verbatim}
%enditem


\newpage
7. Нехай змінна \kw{s} позначає послідовність цілих, по якій можна ітерувати.

Написати вираз, що перевіряє, чи правда, що всі елементи послідовності є точними квадратами.
%enditem


\newpage
8. Що буде надруковано за виконання?
%Оригинал был утерян. Требует отладки!
\begin{verbatim}
a = 0

class C:
    a = 1
    
    def __init__(self):
        self.a = 2
        a = 3
        C.a = 4

obj = C()
obj.a = 5
try:
    print(a, end=' ')
    print(obj.a, end=' ')
    print(C.a, end=' ')
    print(obj.a, end=' ')
    print(obj._C__a)
except:
    print('ERR')
\end{verbatim}
%enditem


\newpage
9. Що буде надруковано за виконання?
\begin{verbatim}
class C0:
    def d(self): print(1, end=' ')
    def f(self): print(2, end=' ')

class C1(C0):
    def d(self): print(3, end=' ')

try:
    C1().f()
    C1().d()
except:
    print('ERR')
\end{verbatim}
%enditem


\newpage
10. Що буде надруковано за виконання?
\begin{verbatim}
class A:
    a = 1

    def _h1(self): print(2, end=' ')
    def __h2(self): print(3, end=' ')

    def main(self):
        try: print(self.a, end=' ')
        except: print('E1', end=' ')

        try: self._h1()
        except: print('E2', end=' ')

        try: self.__h2()
        except: print('E3', end=' ')


class B(A):
    a = 4

    def _h1(self): print(5, end=' ')
    def __h2(self): print(6, end=' ')

try: B().main()
except: print('E4', end=' ')
\end{verbatim}
%enditem


\newpage
11. Змінні \kw{next} та \kw{prev} вузла списку позначають наступний та попередній за ним вузол відповідно. Починаючи з голови, у список записано послідовні цілі числа від~$-58$ до~$48$. Нехай змінна \kw{p} позначає вузол, що зберігає значення~$3$.
Яке значення зберігається у вузлі \kw{p.next.prev.prev.next.prev} ?
%enditem


\newpage
12. 
\begin{center}
	\adjustimage{max size={0.8\linewidth}{0.8\paperheight}}
	{BinTree/trees_exp/18.png}
\end{center}
Указати послідовність міток вершин дерева, зображеного на рис., що утвориться в результаті обходу дерева в оберненому порядку. Мітки записувати через пробіл.\par Алгоритм починає роботу з кореня (на рис. це верхня вершина).
%enditem


\newpage
13. 
\begin{center}
	\adjustimage{max size={0.9\linewidth}{0.9\paperheight}}
	{Graphs/AdjList/108.png}
\end{center}
Для графа, зображеного на рис., вказати список суміжності вершини 3.\par
Формат оформлення відповіді:\par
список суміжності виписувати через пробіл на одному рядку, якщо пробілів десь буде більше одного - не важливо, упорядкування вершин у списку також значення не має.
%enditem


\newpage
14. Для графа на рисунку вказати послідовність, в якій алгоритм пошуку в глибину буде відмічати відвідані вершини, якщо списки суміжності вершин переглядаються алгоритмом за зростанням міток вершин та вершини відмічаються на першому вход\-женні у вершину.\par
Пошук розпочинається з вершини 3.\par

\begin{center}
	\adjustimage{max size={0.9\linewidth}{0.9\paperheight}}
	{Graphs/Traversals/222.png}
\end{center}
%enditem


\newpage
15. Для графа на рисунку з вершини 4 запущено алгоритм пошуку в ширину, що будує кістякове дерево. Списки суміжності вершин переглядаються алгоритмом за зростанням міток вершин.\par
\begin{center}
	\adjustimage{max size={0.9\linewidth}{0.9\paperheight}}
	{Graphs/Traversals/206.png}
\end{center}
\par Які з наведених ребер входять до складу отриманого кістякового дерева?\par
1) (1, 6)\par
2) (2, 6)\par
3) (3, 6)\par
4) (0, 5)\par
5) (3, 5)\par
6) (0, 1)\par
{\vskip 15pt} Номери відповідних ребер записати за зростанням через пробіл.


%enditem


\newpage
16. 

Рядки журналу через двокрапку містять ім'я користувача (ідентифікатор), час події,тип події, код події, наприклад
\begin{verbatim}
s="username : 2022-05-05 13:26 : ERROR : 404"
\end{verbatim}
у коді 
\begin{verbatim}
res = re.fullmatch('???', s) 
s = ''
code_str = ??? 
\end{verbatim}
чим треба замінити кожен з ???, щоб значенням code\_str був код події? 


\vspace{10mm}
\textbf{Завдання 17} 

Задано програму, що містить код 1-2 класів, можливе успадкування, та клієнтський код.
У наведеному коді невеликі частини закриті. 
Також наведено вихід програми.
Необхідно відновити закриті частини, щоб отриманий код відповідав наведеному виходу програми.


\textbf{Завдання 18-19} 

Кілька запитань за завданнями для самостійної роботи.
\vspace{10mm}
Скоріш за все, що завдань буде трохи менше, а деякі попередні будуть переформульовані в стилі двох останніх.

%enditem

\end {large}}

%end

\newpage
%begin

{{
\begin {large}

\newpage
1. Що буде надруковано за виконання?

\begin{verbatim}
print('R', end=' ')
d=('a','b','c','d','e',0,1,2)
print(d[-12])
\end{verbatim}

%enditem


\newpage
2. Що буде надруковано за виконання?

\begin{verbatim}
print('R', end=' ')
e=('0','1','2','3','4','5','6','7','8','9')
print(e[-1::2])
\end{verbatim}

%enditem



\newpage
3. Що буде надруковано за виконання?

\begin{verbatim}
print('R', end=' ')
e=[10,11,12,13,14]
print(e.index(12), end=' ')
\end{verbatim}

%enditem



\newpage
4. Що буде надруковано за виконання?

\begin{verbatim}
print('R', end=' ')
c=['0','1','2','3','4']
del c[8:7]
print(c)
\end{verbatim}

%enditem



\newpage
5. Укажіть ТИП значення виразу.\par
\begin{verbatim}
[13,2,4]+[3,5]
\end{verbatim}
%enditem


\newpage
6. Що буде надруковано за виконання?\par
\begin{verbatim}
try:
    s = {88:5, 81:6, 84:2, 42:3}
    s[39] = 3
    for x, y in s.items():
        print(x, y, sep=' ', end=' ')
except:
    print('ERR')
\end{verbatim}
%enditem


\newpage
7. Нехай змінна \kw{s} позначає дуже довгу послідовність цілих, по якій можна ітерувати. 

Написати вираз-генератор, що задає підпослідовність її елементів, що більші за задане число num. 

Обмеження: усі елементи шуканої підпослідовності одночасно в пам'яті розміщені бути не можуть. 
%enditem


\newpage
8. Що буде надруковано за виконання?
%Оригинал был утерян. Требует отладки!
\begin{verbatim}
_b = 0

class A:
    _b = 1
    
    def __init__(self):
        self._b = 2
        _b = 3
        A._b = 4

obj = A()
obj._b = 5
try:
    print(_b, end=' ')
    print(obj._b, end=' ')
    print(A._b, end=' ')
    print(obj._b, end=' ')
    print(obj._A__b)
except:
    print('ERR')
\end{verbatim}
%enditem


\newpage
9. Що буде надруковано за виконання?
\begin{verbatim}
class B0:
    def c(self): print(1, end=' ')
    def h(self): print(2, end=' ')

class B1(B0):
    def h(self): print(3, end=' ')

try:
    B1().c()
    B1().h()
except:
    print('ERR')
\end{verbatim}
%enditem


\newpage
10. Що буде надруковано за виконання?
\begin{verbatim}
class A:
    b = 1

    def __init__(self):
        c = 2
        c = 3

class B(A):
    b = 4
    a = 5

    def __init__(self):
        super().__init__()
        self.a = 6

try:
    obj = B()
    obj.c = obj.b + 10
    B.c = 13
except: print('EE', end=' ')

try: print(obj.c, end= ' ')
except: print('E1', end=' ')

try: print(A.c, end= ' ')
except: print('E2', end=' ')

try: print(B.c, end= ' ')
except: print('E3', end=' ')
\end{verbatim}
%enditem


\newpage
11. Змінні \kw{next} та \kw{prev} вузла списку позначають наступний та попередній за ним вузол відповідно. Починаючи з голови, у список записано послідовні цілі числа від~$-56$ до~$30$. Нехай змінна \kw{p} позначає вузол, що зберігає значення~$-9$.
Яке значення зберігається у вузлі \kw{p.prev.next.next.prev.next} ?
%enditem


\newpage
12. 
\begin{center}
	\adjustimage{max size={0.8\linewidth}{0.8\paperheight}}
	{BinTree/trees_exp/12.png}
\end{center}
Указати послідовність міток вершин дерева, зображеного на рис., що утвориться в результаті обходу дерева в прямому порядку. Мітки записувати через пробіл.\par Алгоритм починає роботу з кореня (на рис. це верхня вершина).
%enditem


\newpage
13. 
\begin{center}
	\adjustimage{max size={0.9\linewidth}{0.9\paperheight}}
	{Graphs/AdjList/292.png}
\end{center}
Для графа, зображеного на рис., вказати список суміжності вершини 3.\par
Формат оформлення відповіді:\par
список суміжності виписувати через пробіл на одному рядку, якщо пробілів десь буде більше одного - не важливо, упорядкування вершин у списку також значення не має.
%enditem


\newpage
14. Для графа на рисунку вказати послідовність, в якій алгоритм пошуку в глибину буде відмічати відвідані вершини, якщо списки суміжності вершин переглядаються алгоритмом за зростанням міток вершин та вершини відмічаються на першому вход\-женні у вершину.\par
Пошук розпочинається з вершини 6.\par

\begin{center}
	\adjustimage{max size={0.9\linewidth}{0.9\paperheight}}
	{Graphs/Traversals/57.png}
\end{center}
%enditem


\newpage
15. Для графа на рисунку з вершини 3 запущено алгоритм пошуку в ширину, що будує кістякове дерево. Списки суміжності вершин переглядаються алгоритмом за зростанням міток вершин.\par
\begin{center}
	\adjustimage{max size={0.9\linewidth}{0.9\paperheight}}
	{Graphs/Traversals/274.png}
\end{center}
\par Які з наведених ребер входять до складу отриманого кістякового дерева?\par
1) (4, 6)\par
2) (0, 4)\par
3) (2, 4)\par
4) (3, 4)\par
5) (4, 5)\par
6) (1, 4)\par
{\vskip 15pt} Номери відповідних ребер записати за зростанням через пробіл.


%enditem


\newpage
16. 

Рядки журналу через двокрапку містять ім'я користувача (ідентифікатор), час події,тип події, код події, наприклад
\begin{verbatim}
s="username : 2022-05-05 13:26 : ERROR : 404"
\end{verbatim}
у коді 
\begin{verbatim}
res = re.fullmatch('???', s) 
s = ''
code_str = ??? 
\end{verbatim}
чим треба замінити кожен з ???, щоб значенням code\_str був код події? 


\vspace{10mm}
\textbf{Завдання 17} 

Задано програму, що містить код 1-2 класів, можливе успадкування, та клієнтський код.
У наведеному коді невеликі частини закриті. 
Також наведено вихід програми.
Необхідно відновити закриті частини, щоб отриманий код відповідав наведеному виходу програми.


\textbf{Завдання 18-19} 

Кілька запитань за завданнями для самостійної роботи.
\vspace{10mm}
Скоріш за все, що завдань буде трохи менше, а деякі попередні будуть переформульовані в стилі двох останніх.

%enditem

\end {large}}

%end

\newpage
%begin

{{
\begin {large}

\newpage
1. Що буде надруковано за виконання?

\begin{verbatim}
print('R', end=' ')
b=('a','b','c','d','e',0,1,2,3)
print(b[-9])
\end{verbatim}

%enditem


\newpage
2. Що буде надруковано за виконання?

\begin{verbatim}
print('R', end=' ')
a=[0,1,2,3,4,5,6,7,8,9]
print(a[-7::-1])
\end{verbatim}

%enditem



\newpage
3. Що буде надруковано за виконання?

\begin{verbatim}
print('R', end=' ')
a=[10,11,12,13,14]
a.extend(15)
print(a)
\end{verbatim}

%enditem



\newpage
4. Що буде надруковано за виконання?

\begin{verbatim}
print('R', end=' ')
a=['0','1','2','3','4']
a[-7:-9]=()
print(a)
\end{verbatim}

%enditem



\newpage
5. Укажіть ТИП значення виразу.\par
\begin{verbatim}
{(1,4), 8}
\end{verbatim}
%enditem


\newpage
6. Що буде надруковано за виконання?\par
\begin{verbatim}
try:
    t = {25:9, 37:4, 69:3, 25:8}
    t[25] = 0
    for x in t.items():
        print(*x, sep=' ', end=' ')
except:
    print('ERR')
\end{verbatim}
%enditem


\newpage
7. Нехай змінна \kw{s} позначає дуже довгу послідовність цілих, по якій можна ітерувати. 

Написати вираз-генератор, що задає підпослідовність її елементів, що діляться на 3. 

Обмеження: усі елементи шуканої підпослідовності одночасно в пам'яті розміщені бути не можуть. 
%enditem


\newpage
8. Що буде надруковано за виконання?
%Оригинал был утерян. Требует отладки!
\begin{verbatim}
__d = 0

class B:
    __d = 1
    
    def __init__(self):
        self.__d = 2
        __d = 3
        B.__d = 4

obj = B()
B.__d = 5
try:
    print(__d, end=' ')
    print(obj.__d, end=' ')
    print(B.__d, end=' ')
    print(obj.__d, end=' ')
    print(obj._B__d)
except:
    print('ERR')
\end{verbatim}
%enditem


\newpage
9. Що буде надруковано за виконання?
\begin{verbatim}
class D0:
    def c(self): print(1, end=' ')
    def f(self): print(2, end=' ')

class D1(D0):
    def f(self): print(3, end=' ')

try:
    D0().c()
    D1().f()
except:
    print('ERR')
\end{verbatim}
%enditem


\newpage
10. Що буде надруковано за виконання?
\begin{verbatim}
class A:
    b = 1

    def __f1(self): print(2, end=' ')
    def _f2(self): print(3, end=' ')

    def main(self):
        try: print(self.b, end=' ')
        except: print('E1', end=' ')

        try: self.__f1()
        except: print('E2', end=' ')

        try: self._f2()
        except: print('E3', end=' ')


class B(A):
    b = 4

    def __f1(self): print(5, end=' ')
    def _f2(self): print(6, end=' ')

try: B().main()
except: print('E4', end=' ')
\end{verbatim}
%enditem


\newpage
11. Змінні \kw{next} та \kw{prev} вузла списку позначають наступний та попередній за ним вузол відповідно. Починаючи з голови, у список записано послідовні цілі числа від~$-33$ до~$50$. Нехай змінна \kw{p} позначає вузол, що зберігає значення~$-1$.
Яке значення зберігається у вузлі \kw{p.next.prev.next.next.next} ?
%enditem


\newpage
12. 
\begin{center}
	\adjustimage{max size={0.8\linewidth}{0.8\paperheight}}
	{BinTree/trees_exp/61.png}
\end{center}
Указати послідовність міток вершин дерева, зображеного на рис., що утвориться в результаті обходу дерева в прямому порядку. Мітки записувати через пробіл.\par Алгоритм починає роботу з кореня (на рис. це верхня вершина).
%enditem


\newpage
13. 
\begin{center}
	\adjustimage{max size={0.9\linewidth}{0.9\paperheight}}
	{Graphs/AdjList/285.png}
\end{center}
Для графа, зображеного на рис., вказати список суміжності вершини 1.\par
Формат оформлення відповіді:\par
список суміжності виписувати через пробіл на одному рядку, якщо пробілів десь буде більше одного - не важливо, упорядкування вершин у списку також значення не має.
%enditem


\newpage
14. Для графа на рисунку вказати послідовність, в якій алгоритм пошуку в глибину буде відмічати відвідані вершини, якщо списки суміжності вершин переглядаються алгоритмом за зростанням міток вершин та вершини відмічаються на першому вход\-женні у вершину.\par
Пошук розпочинається з вершини 3.\par

\begin{center}
	\adjustimage{max size={0.9\linewidth}{0.9\paperheight}}
	{Graphs/Traversals/319.png}
\end{center}
%enditem


\newpage
15. Для графа на рисунку з вершини 0 запущено алгоритм пошуку в ширину, що будує кістякове дерево. Списки суміжності вершин переглядаються алгоритмом за зростанням міток вершин.\par
\begin{center}
	\adjustimage{max size={0.9\linewidth}{0.9\paperheight}}
	{Graphs/Traversals/336.png}
\end{center}
\par Які з наведених ребер входять до складу отриманого кістякового дерева?\par
1) (1, 6)\par
2) (0, 5)\par
3) (1, 2)\par
4) (5, 6)\par
5) (0, 4)\par
6) (2, 4)\par
{\vskip 15pt} Номери відповідних ребер записати за зростанням через пробіл.


%enditem


\newpage
16. 

Рядки журналу через двокрапку містять ім'я користувача (ідентифікатор), час події,тип події, код події, наприклад
\begin{verbatim}
s="username : 2022-05-05 13:26 : ERROR : 404"
\end{verbatim}
у коді 
\begin{verbatim}
res = re.fullmatch('???', s) 
s = ''
code_str = ??? 
\end{verbatim}
чим треба замінити кожен з ???, щоб значенням code\_str був код події? 


\vspace{10mm}
\textbf{Завдання 17} 

Задано програму, що містить код 1-2 класів, можливе успадкування, та клієнтський код.
У наведеному коді невеликі частини закриті. 
Також наведено вихід програми.
Необхідно відновити закриті частини, щоб отриманий код відповідав наведеному виходу програми.


\textbf{Завдання 18-19} 

Кілька запитань за завданнями для самостійної роботи.
\vspace{10mm}
Скоріш за все, що завдань буде трохи менше, а деякі попередні будуть переформульовані в стилі двох останніх.

%enditem

\end {large}}

%end

\newpage
%begin

{{
\begin {large}

\newpage
1. Що буде надруковано за виконання?

\begin{verbatim}
print('R', end=' ')
a=('a','b','c','d','e',0,1,2)
print(a[7])
\end{verbatim}

%enditem


\newpage
2. Що буде надруковано за виконання?

\begin{verbatim}
print('R', end=' ')
d=['0','1','2','3','4','5','6','7','8','9']
print(d[:15:3])
\end{verbatim}

%enditem



\newpage
3. Що буде надруковано за виконання?

\begin{verbatim}
print('R', end=' ')
h=['0','1','2','3','4']
h.remove('3')
print(h)
\end{verbatim}

%enditem



\newpage
4. Що буде надруковано за виконання?

\begin{verbatim}
print('R', end=' ')
e=[10,11,12,13,14]
e[:-4]='12'
print(e)
\end{verbatim}

%enditem



\newpage
5. Укажіть ТИП значення виразу.\par
\begin{verbatim}
{21,2,4}
\end{verbatim}
%enditem


\newpage
6. Що буде надруковано за виконання?\par
\begin{verbatim}
try:
    d = {53:8, 88:9, 90:8, 90:9}
    d[88] = 3
    for x in d.keys():
        print(x, sep=' ', end=' ')
except:
    print('ERR')
\end{verbatim}
%enditem


\newpage
7. Записати ВИРАЗ, значенням якого є множина, що складається з квадратних корен\-ів цілих чисел від 18 до 2003 (включно), що діляться на 7.
%enditem


\newpage
8. Що буде надруковано за виконання?
\begin{verbatim}
c = 0
m = 1

class D:
    m = 2
    
    def __init__(self):
        global m
        self.m = 3
        m = 4
        D.m = 5

obj = D()
try:
    print(c, end=' ')
    print(m, end=' ')
    print(D.m, end=' ')
    print(obj.m)
except:
    print('ERR')
\end{verbatim}
%enditem


\newpage
9. Що буде надруковано за виконання?
\begin{verbatim}
class C0:
    def a(self): print(1, end=' ')
    def g(self): print(2, end=' ')

class C1(C0):
    def g(self): print(3, end=' ')

try:
    C1().a()
    C1().g()
except:
    print('ERR')
\end{verbatim}
%enditem


\newpage
10. Що буде надруковано за виконання?
\begin{verbatim}
class A:
    d = 1

    def g1(self): print(2, end=' ')
    def _g2(self): print(3, end=' ')

    def main(self):
        try: print(self.d, end=' ')
        except: print('E1', end=' ')

        try: self.g1()
        except: print('E2', end=' ')

        try: self._g2()
        except: print('E3', end=' ')


class B(A):
    d = 4

    def g1(self): print(5, end=' ')
    def _g2(self): print(6, end=' ')

try: B().main()
except: print('E4', end=' ')
\end{verbatim}
%enditem


\newpage
11. Змінні \kw{next} та \kw{prev} вузла списку позначають наступний та попередній за ним вузол відповідно. Починаючи з голови, у список записано послідовні цілі числа від~$-25$ до~$29$. Нехай змінна \kw{p} позначає вузол, що зберігає значення~$0$.
Яке значення зберігається у вузлі \kw{p.prev.next.prev.prev.prev} ?
%enditem


\newpage
12. 
\begin{center}
	\adjustimage{max size={0.8\linewidth}{0.8\paperheight}}
	{BinTree/trees_exp/60.png}
\end{center}
Указати послідовність міток вершин дерева, зображеного на рис., що утвориться в результаті обходу дерева в оберненому порядку. Мітки записувати через пробіл.\par Алгоритм починає роботу з кореня (на рис. це верхня вершина).
%enditem


\newpage
13. 
\begin{center}
	\adjustimage{max size={0.9\linewidth}{0.9\paperheight}}
	{Graphs/AdjList/346.png}
\end{center}
Для графа, зображеного на рис., вказати список суміжності вершини 5.\par
Формат оформлення відповіді:\par
список суміжності виписувати через пробіл на одному рядку, якщо пробілів десь буде більше одного - не важливо, упорядкування вершин у списку також значення не має.
%enditem


\newpage
14. Для графа на рисунку вказати послідовність, в якій алгоритм пошуку в глибину буде відмічати відвідані вершини, якщо списки суміжності вершин переглядаються алгоритмом за зростанням міток вершин та вершини відмічаються на першому вход\-женні у вершину.\par
Пошук розпочинається з вершини 3.\par

\begin{center}
	\adjustimage{max size={0.9\linewidth}{0.9\paperheight}}
	{Graphs/Traversals/166.png}
\end{center}
%enditem


\newpage
15. Для графа на рисунку з вершини 0 запущено алгоритм пошуку в ширину, що будує кістякове дерево. Списки суміжності вершин переглядаються алгоритмом за зростанням міток вершин.\par
\begin{center}
	\adjustimage{max size={0.9\linewidth}{0.9\paperheight}}
	{Graphs/Traversals/75.png}
\end{center}
\par Які з наведених ребер входять до складу отриманого кістякового дерева?\par
1) (2, 6)\par
2) (2, 3)\par
3) (0, 4)\par
4) (1, 6)\par
5) (2, 4)\par
6) (3, 4)\par
{\vskip 15pt} Номери відповідних ребер записати за зростанням через пробіл.


%enditem


\newpage
16. 

Рядки журналу через двокрапку містять ім'я користувача (ідентифікатор), час події,тип події, код події, наприклад
\begin{verbatim}
s="username : 2022-05-05 13:26 : ERROR : 404"
\end{verbatim}
у коді 
\begin{verbatim}
res = re.fullmatch('???', s) 
s = ''
code_str = ??? 
\end{verbatim}
чим треба замінити кожен з ???, щоб значенням code\_str був код події? 


\vspace{10mm}
\textbf{Завдання 17} 

Задано програму, що містить код 1-2 класів, можливе успадкування, та клієнтський код.
У наведеному коді невеликі частини закриті. 
Також наведено вихід програми.
Необхідно відновити закриті частини, щоб отриманий код відповідав наведеному виходу програми.


\textbf{Завдання 18-19} 

Кілька запитань за завданнями для самостійної роботи.
\vspace{10mm}
Скоріш за все, що завдань буде трохи менше, а деякі попередні будуть переформульовані в стилі двох останніх.

%enditem

\end {large}}

%end

\newpage
%begin

{{
\begin {large}

\newpage
1. Що буде надруковано за виконання?

\begin{verbatim}
print('R', end=' ')
f=('a','b','c','d','e',0,1)
print(f[-13])
\end{verbatim}

%enditem


\newpage
2. Що буде надруковано за виконання?

\begin{verbatim}
print('R', end=' ')
f='0123456789'
print(f[:-11:-2])
\end{verbatim}

%enditem



\newpage
3. Що буде надруковано за виконання?

\begin{verbatim}
print('R', end=' ')
b=[10,11,12,13,14]
b.insert(1, '10')
print(b)
\end{verbatim}

%enditem



\newpage
4. Що буде надруковано за виконання?

\begin{verbatim}
print('R', end=' ')
g=['0','1','2','3','4']
g[-3:-2]='10'
print(g)
\end{verbatim}

%enditem



\newpage
5. Укажіть ТИП значення виразу.\par
\begin{verbatim}
[{1,2}, 5]
\end{verbatim}
%enditem


\newpage
6. Що буде надруковано за виконання?\par
\begin{verbatim}
try:
    t = {69:2, 22:3, 52:1, 22:1}
    t[65] = 9
    for x in t.keys():
        print(x, sep=' ', end=' ')
except:
    print('ERR')
\end{verbatim}
%enditem


\newpage
7. 
Записати ВИРАЗ, значенням якого є список, що складається з цілих чисел від 15 до 2013 (включно), причому спочатку за спаданням записано всі, що не діляться на 7, а потім решту за зростанням.

%enditem


\newpage
8. Що буде надруковано за виконання?
\begin{verbatim}
b = 0
j = 1

class A:
    b = 2
    
    def __init__(self):
        global b
        self.b = 3
        b = 4
        A.b = 5

obj = A()
try:
    print(b, end=' ')
    print(j, end=' ')
    print(A.b, end=' ')
    print(obj.b)
except:
    print('ERR')
\end{verbatim}
%enditem


\newpage
9. Що буде надруковано за виконання?
\begin{verbatim}
class A0:
    def b(self): print(1, end=' ')
    def h(self): print(2, end=' ')

class A1(A0):
    def b(self): print(3, end=' ')

try:
    A1().h()
    A0().b()
except:
    print('ERR')
\end{verbatim}
%enditem


\newpage
10. Що буде надруковано за виконання?
\begin{verbatim}
class A:
    c = 1

    def __init__(self):
        b = 2
        b = 3

class B(A):
    c = 4
    a = 5

    def __init__(self):
        super().__init__()
        self.a = 6

try:
    obj = B()
    obj.a = obj.c + 10
    B.a = 13
except: print('EE', end=' ')

try: print(obj.a, end= ' ')
except: print('E1', end=' ')

try: print(A.a, end= ' ')
except: print('E2', end=' ')

try: print(B.a, end= ' ')
except: print('E3', end=' ')
\end{verbatim}
%enditem


\newpage
11. Змінні \kw{next} та \kw{prev} вузла списку позначають наступний та попередній за ним вузол відповідно. Починаючи з голови, у список записано послідовні цілі числа від~$-54$ до~$32$. Нехай змінна \kw{p} позначає вузол, що зберігає значення~$-7$.
Яке значення зберігається у вузлі \kw{p.next.prev.next.prev.next} ?
%enditem


\newpage
12. 
\begin{center}
	\adjustimage{max size={0.8\linewidth}{0.8\paperheight}}
	{BinTree/trees_exp/76.png}
\end{center}
Указати послідовність міток вершин дерева, зображеного на рис., що утвориться в результаті обходу дерева в прямому порядку. Мітки записувати через пробіл.\par Алгоритм починає роботу з кореня (на рис. це верхня вершина).
%enditem


\newpage
13. 
\begin{center}
	\adjustimage{max size={0.9\linewidth}{0.9\paperheight}}
	{Graphs/AdjList/5.png}
\end{center}
Для графа, зображеного на рис., вказати список суміжності вершини 1.\par
Формат оформлення відповіді:\par
список суміжності виписувати через пробіл на одному рядку, якщо пробілів десь буде більше одного - не важливо, упорядкування вершин у списку також значення не має.
%enditem


\newpage
14. Для графа на рисунку з вершини 5 запущено алгоритм пошуку в глибину, що будує кістякове дерево. Списки суміжності вершин переглядаються алгоритмом за зростанням міток вершин.\par
\begin{center}
	\adjustimage{max size={0.9\linewidth}{0.9\paperheight}}
	{Graphs/Traversals/10.png}
\end{center}
\par Які з наведених ребер входять до складу отриманого кістякового дерева?\par
1) (3, 5)\par
2) (0, 2)\par
3) (3, 4)\par
4) (0, 6)\par
5) (0, 1)\par
6) (1, 6)\par
{\vskip 15pt} Номери відповідних ребер записати за зростанням через пробіл.

%enditem


\newpage
15. Для графа на рисунку з вершини 4 запущено алгоритм пошуку в ширину, що будує кістякове дерево. Списки суміжності вершин переглядаються алгоритмом за зростанням міток вершин.\par
\begin{center}
	\adjustimage{max size={0.9\linewidth}{0.9\paperheight}}
	{Graphs/Traversals/239.png}
\end{center}
\par Які з наведених ребер входять до складу отриманого кістякового дерева?\par
1) (1, 4)\par
2) (2, 4)\par
3) (3, 5)\par
4) (0, 2)\par
5) (0, 5)\par
6) (0, 4)\par
{\vskip 15pt} Номери відповідних ребер записати за зростанням через пробіл.


%enditem


\newpage
16. 

Рядки журналу через двокрапку містять ім'я користувача (ідентифікатор), час події,тип події, код події, наприклад
\begin{verbatim}
s="username : 2022-05-05 13:26 : ERROR : 404"
\end{verbatim}
у коді 
\begin{verbatim}
res = re.fullmatch('???', s) 
s = ''
code_str = ??? 
\end{verbatim}
чим треба замінити кожен з ???, щоб значенням code\_str був код події? 


\vspace{10mm}
\textbf{Завдання 17} 

Задано програму, що містить код 1-2 класів, можливе успадкування, та клієнтський код.
У наведеному коді невеликі частини закриті. 
Також наведено вихід програми.
Необхідно відновити закриті частини, щоб отриманий код відповідав наведеному виходу програми.


\textbf{Завдання 18-19} 

Кілька запитань за завданнями для самостійної роботи.
\vspace{10mm}
Скоріш за все, що завдань буде трохи менше, а деякі попередні будуть переформульовані в стилі двох останніх.

%enditem

\end {large}}

%end

\newpage
%begin

{{
\begin {large}

\newpage
1. Що буде надруковано за виконання?

\begin{verbatim}
print('R', end=' ')
g=('a','b','c','d','e',0,1,2,3)
print(g[1])
\end{verbatim}

%enditem


\newpage
2. Що буде надруковано за виконання?

\begin{verbatim}
print('R', end=' ')
b=('0','1','2','3','4','5','6','7','8','9')
print(b[-5:-12:-1])
\end{verbatim}

%enditem



\newpage
3. Що буде надруковано за виконання?

\begin{verbatim}
print('R', end=' ')
c=['0','1','2','3','4']
print(c.index(2), end=' ')
\end{verbatim}

%enditem



\newpage
4. Що буде надруковано за виконання?

\begin{verbatim}
print('R', end=' ')
f=[10,11,12,13,14]
del f[-8:]
print(f)
\end{verbatim}

%enditem



\newpage
5. Укажіть ТИП значення виразу.\par
\begin{verbatim}
[7,2,3,4]
\end{verbatim}
%enditem


\newpage
6. Що буде надруковано за виконання?\par
\begin{verbatim}
try:
    s = {85:5, 73:2, 33:2, 73:1}
    s[12] = 5
    for x in s.items():
        print(x, sep=' ', end=' ')
except:
    print('ERR')
\end{verbatim}
%enditem


\newpage
7. Нехай змінна \kw{s} позначає дуже довгу послідовність цілих, по якій можна ітерувати. 

Написати вираз-генератор, що задає підпослідовність її елементів, що не діляться на жодне з чисел 2, 7, 5. 

Обмеження: усі елементи шуканої підпослідовності одночасно в пам'яті розміщені бути не можуть. 
%enditem


\newpage
8. Що буде надруковано за виконання?
%Оригинал был утерян. Требует отладки!
\begin{verbatim}
_d = 0

class B:
    _d = 1
    
    def __init__(self):
        self._d = 2
        _d = 3
        B._d = 4

obj = B()
B._d = 5
try:
    print(_d, end=' ')
    print(obj._d, end=' ')
    print(B._d, end=' ')
    print(obj._d, end=' ')
    print(obj._B__d)
except:
    print('ERR')
\end{verbatim}
%enditem


\newpage
9. Що буде надруковано за виконання?
\begin{verbatim}
class B0:
    def d(self): print(1, end=' ')
    def f(self): print(2, end=' ')

class B1(B0):
    def d(self): print(3, end=' ')

try:
    B1().f()
    B1().d()
except:
    print('ERR')
\end{verbatim}
%enditem


\newpage
10. Що буде надруковано за виконання?
\begin{verbatim}
class A:
    b = 1

    def __init__(self):
        a = 2
        b = 3

class B(A):
    a = 4
    c = 5

    def __init__(self):
        super().__init__()
        self.c = 6

try:
    obj = B()
    obj.a = obj.b + 10
    B.a = 13
except: print('EE', end=' ')

try: print(obj.a, end= ' ')
except: print('E1', end=' ')

try: print(A.a, end= ' ')
except: print('E2', end=' ')

try: print(B.a, end= ' ')
except: print('E3', end=' ')
\end{verbatim}
%enditem


\newpage
11. Змінні \kw{next} та \kw{prev} вузла списку позначають наступний та попередній за ним вузол відповідно. Починаючи з голови, у список записано послідовні цілі числа від~$-65$ до~$31$. Нехай змінна \kw{p} позначає вузол, що зберігає значення~$6$.
Яке значення зберігається у вузлі \kw{p.prev.next.prev.next.prev} ?
%enditem


\newpage
12. 
\begin{center}
	\adjustimage{max size={0.8\linewidth}{0.8\paperheight}}
	{BinTree/trees_exp/134.png}
\end{center}
Указати послідовність міток вершин дерева, зображеного на рис., що утвориться в результаті обходу дерева в оберненому порядку. Мітки записувати через пробіл.\par Алгоритм починає роботу з кореня (на рис. це верхня вершина).
%enditem


\newpage
13. 
\begin{center}
	\adjustimage{max size={0.9\linewidth}{0.9\paperheight}}
	{Graphs/AdjList/229.png}
\end{center}
Для графа, зображеного на рис., вказати список суміжності вершини 2.\par
Формат оформлення відповіді:\par
список суміжності виписувати через пробіл на одному рядку, якщо пробілів десь буде більше одного - не важливо, упорядкування вершин у списку також значення не має.
%enditem


\newpage
14. Для графа на рисунку вказати послідовність, в якій алгоритм пошуку в глибину буде відмічати відвідані вершини, якщо списки суміжності вершин переглядаються алгоритмом за зростанням міток вершин та вершини відмічаються на першому вход\-женні у вершину.\par
Пошук розпочинається з вершини 1.\par

\begin{center}
	\adjustimage{max size={0.9\linewidth}{0.9\paperheight}}
	{Graphs/Traversals/314.png}
\end{center}
%enditem


\newpage
15. Для графа на рисунку з вершини 3 запущено алгоритм пошуку в ширину, що будує кістякове дерево. Списки суміжності вершин переглядаються алгоритмом за зростанням міток вершин.\par
\begin{center}
	\adjustimage{max size={0.9\linewidth}{0.9\paperheight}}
	{Graphs/Traversals/45.png}
\end{center}
\par Які з наведених ребер входять до складу отриманого кістякового дерева?\par
1) (0, 4)\par
2) (1, 3)\par
3) (4, 6)\par
4) (2, 3)\par
5) (2, 4)\par
6) (0, 3)\par
{\vskip 15pt} Номери відповідних ребер записати за зростанням через пробіл.


%enditem


\newpage
16. 

Рядки журналу через двокрапку містять ім'я користувача (ідентифікатор), час події,тип події, код події, наприклад
\begin{verbatim}
s="username : 2022-05-05 13:26 : ERROR : 404"
\end{verbatim}
у коді 
\begin{verbatim}
res = re.fullmatch('???', s) 
s = ''
code_str = ??? 
\end{verbatim}
чим треба замінити кожен з ???, щоб значенням code\_str був код події? 


\vspace{10mm}
\textbf{Завдання 17} 

Задано програму, що містить код 1-2 класів, можливе успадкування, та клієнтський код.
У наведеному коді невеликі частини закриті. 
Також наведено вихід програми.
Необхідно відновити закриті частини, щоб отриманий код відповідав наведеному виходу програми.


\textbf{Завдання 18-19} 

Кілька запитань за завданнями для самостійної роботи.
\vspace{10mm}
Скоріш за все, що завдань буде трохи менше, а деякі попередні будуть переформульовані в стилі двох останніх.

%enditem

\end {large}}

%end

\newpage
%begin

{{
\begin {large}

\newpage
1. Що буде надруковано за виконання?

\begin{verbatim}
print('R', end=' ')
c=('a','b','c','d','e',0,1,2,3,4)
print(c[-3])
\end{verbatim}

%enditem


\newpage
2. Що буде надруковано за виконання?

\begin{verbatim}
print('R', end=' ')
e=(0,1,2,3,4,5,6,7,8,9)
print(e[2:-13:-3])
\end{verbatim}

%enditem



\newpage
3. Що буде надруковано за виконання?

\begin{verbatim}
print('R', end=' ')
f=['0','1','2','3','4']
f.remove('6')
print(f)
\end{verbatim}

%enditem



\newpage
4. Що буде надруковано за виконання?

\begin{verbatim}
print('R', end=' ')
f=['0','1','2','3','4']
f[-1:0]='11'
print(f)
\end{verbatim}

%enditem



\newpage
5. Укажіть ТИП значення виразу.\par
\begin{verbatim}
[x for x in range(7)]
\end{verbatim}
%enditem


\newpage
6. Що буде надруковано за виконання?\par
\begin{verbatim}
try:
    d = {87:0, 35:7, 54:9, 87:9}
    d[39] = 7
    for x in d.values():
        print(x, sep=' ', end=' ')
except:
    print('ERR')
\end{verbatim}
%enditem


\newpage
7. Нехай змінна \kw{s} позначає послідовність цілих, по якій можна ітерувати.

Написати вираз, що перевіряє, чи правда, що всі елементи послідовності є цілими.
%enditem


\newpage
8. Що буде надруковано за виконання?
\begin{verbatim}
a = 0
n = 1

class B:
    n = 2
    
    def __init__(self):
        self.n = 3
        n = 4
        B.n = 5

obj = B()
try:
    print(a, end=' ')
    print(n, end=' ')
    print(B.n, end=' ')
    print(obj.n)
except:
    print('ERR')
\end{verbatim}
%enditem


\newpage
9. Що буде надруковано за виконання?
\begin{verbatim}
class D0:
    def b(self): print(1, end=' ')
    def g(self): print(2, end=' ')

class D1(D0):
    def g(self): print(3, end=' ')

try:
    D1().b()
    D0().g()
except:
    print('ERR')
\end{verbatim}
%enditem


\newpage
10. Що буде надруковано за виконання?
\begin{verbatim}
class A:
    c = 1

    def __init__(self):
        a = 2
        a = 3

class B(A):
    a = 4
    b = 5

    def __init__(self):
        super().__init__()
        self.b = 6

try:
    obj = B()
    obj.c = obj.c + 10
    B.c = 13
except: print('EE', end=' ')

try: print(obj.c, end= ' ')
except: print('E1', end=' ')

try: print(A.c, end= ' ')
except: print('E2', end=' ')

try: print(B.c, end= ' ')
except: print('E3', end=' ')
\end{verbatim}
%enditem


\newpage
11. Змінні \kw{next} та \kw{prev} вузла списку позначають наступний та попередній за ним вузол відповідно. Починаючи з голови, у список записано послідовні цілі числа від~$-28$ до~$40$. Нехай змінна \kw{p} позначає вузол, що зберігає значення~$4$.
Яке значення зберігається у вузлі \kw{p.next.prev.prev.prev.prev} ?
%enditem


\newpage
12. 
\begin{center}
	\adjustimage{max size={0.8\linewidth}{0.8\paperheight}}
	{BinTree/trees_exp/171.png}
\end{center}
Указати послідовність міток вершин дерева, зображеного на рис., що утвориться в результаті обходу дерева в прямому порядку. Мітки записувати через пробіл.\par Алгоритм починає роботу з кореня (на рис. це верхня вершина).
%enditem


\newpage
13. 
\begin{center}
	\adjustimage{max size={0.9\linewidth}{0.9\paperheight}}
	{Graphs/AdjList/278.png}
\end{center}
Для графа, зображеного на рис., вказати список суміжності вершини 0.\par
Формат оформлення відповіді:\par
список суміжності виписувати через пробіл на одному рядку, якщо пробілів десь буде більше одного - не важливо, упорядкування вершин у списку також значення не має.
%enditem


\newpage
14. Для графа на рисунку вказати послідовність, в якій алгоритм пошуку в глибину буде відмічати відвідані вершини, якщо списки суміжності вершин переглядаються алгоритмом за зростанням міток вершин та вершини відмічаються на першому вход\-женні у вершину.\par
Пошук розпочинається з вершини 1.\par

\begin{center}
	\adjustimage{max size={0.9\linewidth}{0.9\paperheight}}
	{Graphs/Traversals/216.png}
\end{center}
%enditem


\newpage
15. Для графа на рисунку з вершини 5 запущено алгоритм пошуку в ширину, що будує кістякове дерево. Списки суміжності вершин переглядаються алгоритмом за зростанням міток вершин.\par
\begin{center}
	\adjustimage{max size={0.9\linewidth}{0.9\paperheight}}
	{Graphs/Traversals/307.png}
\end{center}
\par Які з наведених ребер входять до складу отриманого кістякового дерева?\par
1) (2, 3)\par
2) (1, 6)\par
3) (5, 6)\par
4) (4, 5)\par
5) (1, 3)\par
6) (1, 2)\par
{\vskip 15pt} Номери відповідних ребер записати за зростанням через пробіл.


%enditem


\newpage
16. 

Рядки журналу через двокрапку містять ім'я користувача (ідентифікатор), час події,тип події, код події, наприклад
\begin{verbatim}
s="username : 2022-05-05 13:26 : ERROR : 404"
\end{verbatim}
у коді 
\begin{verbatim}
res = re.fullmatch('???', s) 
s = ''
code_str = ??? 
\end{verbatim}
чим треба замінити кожен з ???, щоб значенням code\_str був код події? 


\vspace{10mm}
\textbf{Завдання 17} 

Задано програму, що містить код 1-2 класів, можливе успадкування, та клієнтський код.
У наведеному коді невеликі частини закриті. 
Також наведено вихід програми.
Необхідно відновити закриті частини, щоб отриманий код відповідав наведеному виходу програми.


\textbf{Завдання 18-19} 

Кілька запитань за завданнями для самостійної роботи.
\vspace{10mm}
Скоріш за все, що завдань буде трохи менше, а деякі попередні будуть переформульовані в стилі двох останніх.

%enditem

\end {large}}

%end

\newpage
%begin

{{
\begin {large}

\newpage
1. Що буде надруковано за виконання?

\begin{verbatim}
print('R', end=' ')
b=('a','b','c','d','e',0,1,2,4)
print(b[8])
\end{verbatim}

%enditem


\newpage
2. Що буде надруковано за виконання?

\begin{verbatim}
print('R', end=' ')
a=[0,1,2,3,4,5,6,7,8,9]
print(a[1:5:2])
\end{verbatim}

%enditem



\newpage
3. Що буде надруковано за виконання?

\begin{verbatim}
print('R', end=' ')
a=[10,11,12,13,14]
a+='13'
print(a)
\end{verbatim}

%enditem



\newpage
4. Що буде надруковано за виконання?

\begin{verbatim}
print('R', end=' ')
e=[10,11,12,13,14]
e[:-5]=[1,0]
print(e)
\end{verbatim}

%enditem



\newpage
5. Укажіть ТИП значення виразу.\par
\begin{verbatim}
{x for x in range(7) if el%3==1}
\end{verbatim}
%enditem


\newpage
6. Що буде надруковано за виконання?\par
\begin{verbatim}
try:
    t = {23:3, 22:1, 17:5, 17:5}
    t[87] = 4
    for x in t :
        print(x, sep=' ', end=' ')
except:
    print('ERR')
\end{verbatim}
%enditem


\newpage
7. Записати ВИРАЗ, значенням якого є словник, ключами якого є цілі числа від 11 до 2006 (включно), що не діляться на жодне з чисел 2, 3, а ключам відповідають їх квадратні корені  
%enditem


\newpage
8. Що буде надруковано за виконання?
\begin{verbatim}
a = 0
h = 1

class D:
    a = 2
    
    def __init__(self):
        global a
        self.h = 3
        a = 4
        D.a = 5

obj = D()
try:
    print(a, end=' ')
    print(h, end=' ')
    print(D.a, end=' ')
    print(obj.h)
except:
    print('ERR')
\end{verbatim}
%enditem


\newpage
9. Що буде надруковано за виконання?
\begin{verbatim}
class C0:
    def a(self): print(1, end=' ')
    def h(self): print(2, end=' ')

class C1(C0):
    def h(self): print(3, end=' ')

try:
    C1().a()
    C1().h()
except:
    print('ERR')
\end{verbatim}
%enditem


\newpage
10. Що буде надруковано за виконання?
\begin{verbatim}
class A:
    a = 1

    def _f1(self): print(2, end=' ')
    def _f2(self): print(3, end=' ')

    def main(self):
        try: print(self.a, end=' ')
        except: print('E1', end=' ')

        try: self._f1()
        except: print('E2', end=' ')

        try: self._f2()
        except: print('E3', end=' ')


class B(A):
    a = 4

    def _f1(self): print(5, end=' ')
    def _f2(self): print(6, end=' ')

try: B().main()
except: print('E4', end=' ')
\end{verbatim}
%enditem


\newpage
11. Змінні \kw{next} та \kw{prev} вузла списку позначають наступний та попередній за ним вузол відповідно. Починаючи з голови, у список записано послідовні цілі числа від~$-35$ до~$24$. Нехай змінна \kw{p} позначає вузол, що зберігає значення~$7$.
Яке значення зберігається у вузлі \kw{p.prev.next.next.next.next} ?
%enditem


\newpage
12. 
\begin{center}
	\adjustimage{max size={0.8\linewidth}{0.8\paperheight}}
	{BinTree/trees_exp/100.png}
\end{center}
Указати послідовність міток вершин дерева, зображеного на рис., що утвориться в результаті обходу дерева в оберненому порядку. Мітки записувати через пробіл.\par Алгоритм починає роботу з кореня (на рис. це верхня вершина).
%enditem


\newpage
13. 
\begin{center}
	\adjustimage{max size={0.9\linewidth}{0.9\paperheight}}
	{Graphs/AdjList/59.png}
\end{center}
Для графа, зображеного на рис., вказати список суміжності вершини 3.\par
Формат оформлення відповіді:\par
список суміжності виписувати через пробіл на одному рядку, якщо пробілів десь буде більше одного - не важливо, упорядкування вершин у списку також значення не має.
%enditem


\newpage
14. Для графа на рисунку вказати послідовність, в якій алгоритм пошуку в глибину буде відмічати відвідані вершини, якщо списки суміжності вершин переглядаються алгоритмом за зростанням міток вершин та вершини відмічаються на першому вход\-женні у вершину.\par
Пошук розпочинається з вершини 2.\par

\begin{center}
	\adjustimage{max size={0.9\linewidth}{0.9\paperheight}}
	{Graphs/Traversals/192.png}
\end{center}
%enditem


\newpage
15. Для графа на рисунку з вершини 0 запущено алгоритм пошуку в ширину, що будує кістякове дерево. Списки суміжності вершин переглядаються алгоритмом за зростанням міток вершин.\par
\begin{center}
	\adjustimage{max size={0.9\linewidth}{0.9\paperheight}}
	{Graphs/Traversals/148.png}
\end{center}
\par Які з наведених ребер входять до складу отриманого кістякового дерева?\par
1) (1, 5)\par
2) (0, 5)\par
3) (4, 6)\par
4) (0, 1)\par
5) (3, 4)\par
6) (2, 6)\par
{\vskip 15pt} Номери відповідних ребер записати за зростанням через пробіл.


%enditem


\newpage
16. 

Рядки журналу через двокрапку містять ім'я користувача (ідентифікатор), час події,тип події, код події, наприклад
\begin{verbatim}
s="username : 2022-05-05 13:26 : ERROR : 404"
\end{verbatim}
у коді 
\begin{verbatim}
res = re.fullmatch('???', s) 
s = ''
code_str = ??? 
\end{verbatim}
чим треба замінити кожен з ???, щоб значенням code\_str був код події? 


\vspace{10mm}
\textbf{Завдання 17} 

Задано програму, що містить код 1-2 класів, можливе успадкування, та клієнтський код.
У наведеному коді невеликі частини закриті. 
Також наведено вихід програми.
Необхідно відновити закриті частини, щоб отриманий код відповідав наведеному виходу програми.


\textbf{Завдання 18-19} 

Кілька запитань за завданнями для самостійної роботи.
\vspace{10mm}
Скоріш за все, що завдань буде трохи менше, а деякі попередні будуть переформульовані в стилі двох останніх.

%enditem

\end {large}}

%end

\newpage
%begin

{{
\begin {large}

\newpage
1. Що буде надруковано за виконання?

\begin{verbatim}
print('R', end=' ')
a=('a','b','c','d','e',0,1,2,3)
print(a[4])
\end{verbatim}

%enditem


\newpage
2. Що буде надруковано за виконання?

\begin{verbatim}
print('R', end=' ')
f=(0,1,2,3,4,5,6,7,8,9)
print(f[:-14:-2])
\end{verbatim}

%enditem



\newpage
3. Що буде надруковано за виконання?

\begin{verbatim}
print('R', end=' ')
b=[10,11,12,13,14]
b.append([1,0])
print(b)
\end{verbatim}

%enditem



\newpage
4. Що буде надруковано за виконання?

\begin{verbatim}
print('R', end=' ')
g=['0','1','2','3','4']
del g[4:]
print(g)
\end{verbatim}

%enditem



\newpage
5. Укажіть ТИП значення виразу.\par
\begin{verbatim}
(2, 7, -4)
\end{verbatim}
%enditem


\newpage
6. Що буде надруковано за виконання?\par
\begin{verbatim}
try:
    d = {52:2, 42:4, 21:9, 42:8}
    d[69] = 1
    for x, y in d.items():
        print(x, y, sep=' ', end=' ')
except:
    print('ERR')
\end{verbatim}
%enditem


\newpage
7. Задано словник \kw{d} з цілими ключами та значеннями. Записати ВИРАЗ, значенням якого є множина, що складається  з тих значень, чиї ключі невід'єм\-ні.
%enditem


\newpage
8. Що буде надруковано за виконання?
\begin{verbatim}
b = 0
n = 1

class B:
    n = 2
    
    def __init__(self):
        self.b = 3
        n = 4
        B.n = 5

obj = B()
try:
    print(b, end=' ')
    print(n, end=' ')
    print(B.n, end=' ')
    print(obj.b)
except:
    print('ERR')
\end{verbatim}
%enditem


\newpage
9. Що буде надруковано за виконання?
\begin{verbatim}
class B0:
    def c(self): print(1, end=' ')
    def f(self): print(2, end=' ')

class B1(B0):
    def c(self): print(3, end=' ')

try:
    B0().f()
    B1().c()
except:
    print('ERR')
\end{verbatim}
%enditem


\newpage
10. Що буде надруковано за виконання?
\begin{verbatim}
class A:
    b = 1

    def _h1(self): print(2, end=' ')
    def __h2(self): print(3, end=' ')

    def main(self):
        try: print(self.b, end=' ')
        except: print('E1', end=' ')

        try: self._h1()
        except: print('E2', end=' ')

        try: self.__h2()
        except: print('E3', end=' ')


class B(A):
    b = 4

    def _h1(self): print(5, end=' ')
    def __h2(self): print(6, end=' ')

try: B().main()
except: print('E4', end=' ')
\end{verbatim}
%enditem


\newpage
11. Змінні \kw{next} та \kw{prev} вузла списку позначають наступний та попередній за ним вузол відповідно. Починаючи з голови, у список записано послідовні цілі числа від~$-48$ до~$21$. Нехай змінна \kw{p} позначає вузол, що зберігає значення~$-6$.
Яке значення зберігається у вузлі \kw{p.next.prev.next.prev.next} ?
%enditem


\newpage
12. 
\begin{center}
	\adjustimage{max size={0.8\linewidth}{0.8\paperheight}}
	{BinTree/trees_exp/103.png}
\end{center}
Указати послідовність міток вершин дерева, зображеного на рис., що утвориться в результаті обходу дерева в прямому порядку. Мітки записувати через пробіл.\par Алгоритм починає роботу з кореня (на рис. це верхня вершина).
%enditem


\newpage
13. 
\begin{center}
	\adjustimage{max size={0.9\linewidth}{0.9\paperheight}}
	{Graphs/AdjList/287.png}
\end{center}
Для графа, зображеного на рис., вказати список суміжності вершини 3.\par
Формат оформлення відповіді:\par
список суміжності виписувати через пробіл на одному рядку, якщо пробілів десь буде більше одного - не важливо, упорядкування вершин у списку також значення не має.
%enditem


\newpage
14. Для графа на рисунку вказати послідовність, в якій алгоритм пошуку в глибину буде відмічати відвідані вершини, якщо списки суміжності вершин переглядаються алгоритмом за зростанням міток вершин та вершини відмічаються на першому вход\-женні у вершину.\par
Пошук розпочинається з вершини 1.\par

\begin{center}
	\adjustimage{max size={0.9\linewidth}{0.9\paperheight}}
	{Graphs/Traversals/347.png}
\end{center}
%enditem


\newpage
15. Для графа на рисунку з вершини 1 запущено алгоритм пошуку в ширину, що будує кістякове дерево. Списки суміжності вершин переглядаються алгоритмом за зростанням міток вершин.\par
\begin{center}
	\adjustimage{max size={0.9\linewidth}{0.9\paperheight}}
	{Graphs/Traversals/26.png}
\end{center}
\par Які з наведених ребер входять до складу отриманого кістякового дерева?\par
1) (2, 3)\par
2) (2, 5)\par
3) (1, 6)\par
4) (0, 4)\par
5) (2, 4)\par
6) (0, 3)\par
{\vskip 15pt} Номери відповідних ребер записати за зростанням через пробіл.


%enditem


\newpage
16. 

Рядки журналу через двокрапку містять ім'я користувача (ідентифікатор), час події,тип події, код події, наприклад
\begin{verbatim}
s="username : 2022-05-05 13:26 : ERROR : 404"
\end{verbatim}
у коді 
\begin{verbatim}
res = re.fullmatch('???', s) 
s = ''
code_str = ??? 
\end{verbatim}
чим треба замінити кожен з ???, щоб значенням code\_str був код події? 


\vspace{10mm}
\textbf{Завдання 17} 

Задано програму, що містить код 1-2 класів, можливе успадкування, та клієнтський код.
У наведеному коді невеликі частини закриті. 
Також наведено вихід програми.
Необхідно відновити закриті частини, щоб отриманий код відповідав наведеному виходу програми.


\textbf{Завдання 18-19} 

Кілька запитань за завданнями для самостійної роботи.
\vspace{10mm}
Скоріш за все, що завдань буде трохи менше, а деякі попередні будуть переформульовані в стилі двох останніх.

%enditem

\end {large}}

%end

\newpage
%begin

{{
\begin {large}

\newpage
1. Що буде надруковано за виконання?

\begin{verbatim}
print('R', end=' ')
g=('a','b','c','d','e',0,1,2,3,4)
print(g[-8])
\end{verbatim}

%enditem


\newpage
2. Що буде надруковано за виконання?

\begin{verbatim}
print('R', end=' ')
d=('0','1','2','3','4','5','6','7','8','9')
print(d[5:2:2])
\end{verbatim}

%enditem



\newpage
3. Що буде надруковано за виконання?

\begin{verbatim}
print('R', end=' ')
c=['0','1','2','3','4']
print(c.pop(4), end=' ')
print(c)
\end{verbatim}

%enditem



\newpage
4. Що буде надруковано за виконання?

\begin{verbatim}
print('R', end=' ')
c=[10,11,12,13,14]
c[-6:1]=(14)
print(c)
\end{verbatim}

%enditem



\newpage
5. Укажіть ТИП значення виразу.\par
\begin{verbatim}
{x for x in range(7)}
\end{verbatim}
%enditem


\newpage
6. Що буде надруковано за виконання?\par
\begin{verbatim}
try:
    s = {42:8, 89:4, 64:2, 84:3, 42:6}
    s[84] = 4
    for x in s.items():
        print(*x, sep=' ', end=' ')
except:
    print('ERR')
\end{verbatim}
%enditem


\newpage
7. 
Записати ВИРАЗ, значенням якого є список, що складається з кубів цілих чисел від 12 до 2005 (включно), що не діляться на жодне з чисел 2, 3.
%enditem


\newpage
8. Що буде надруковано за виконання?
%Оригинал был утерян. Требует отладки!
\begin{verbatim}
b = 0

class C:
    b = 1
    
    def __init__(self):
        self.b = 2
        b = 3
        C.b = 4

obj = C()
obj.b = 5
try:
    print(b, end=' ')
    print(obj.b, end=' ')
    print(C.b, end=' ')
    print(obj.b, end=' ')
    print(obj._C__b)
except:
    print('ERR')
\end{verbatim}
%enditem


\newpage
9. Що буде надруковано за виконання?
\begin{verbatim}
class A0:
    def d(self): print(1, end=' ')
    def g(self): print(2, end=' ')

class A1(A0):
    def g(self): print(3, end=' ')

try:
    A1().d()
    A1().g()
except:
    print('ERR')
\end{verbatim}
%enditem


\newpage
10. Що буде надруковано за виконання?
\begin{verbatim}
class A:
    a = 1

    def __init__(self):
        b = 2
        a = 3

class B(A):
    a = 4
    c = 5

    def __init__(self):
        super().__init__()
        self.c = 6

try:
    obj = B()
    obj.b = obj.a + 10
    B.b = 13
except: print('EE', end=' ')

try: print(obj.b, end= ' ')
except: print('E1', end=' ')

try: print(A.b, end= ' ')
except: print('E2', end=' ')

try: print(B.b, end= ' ')
except: print('E3', end=' ')
\end{verbatim}
%enditem


\newpage
11. Змінні \kw{next} та \kw{prev} вузла списку позначають наступний та попередній за ним вузол відповідно. Починаючи з голови, у список записано послідовні цілі числа від~$-23$ до~$25$. Нехай змінна \kw{p} позначає вузол, що зберігає значення~$2$.
Яке значення зберігається у вузлі \kw{p.prev.next.prev.next.prev} ?
%enditem


\newpage
12. 
\begin{center}
	\adjustimage{max size={0.8\linewidth}{0.8\paperheight}}
	{BinTree/trees_exp/178.png}
\end{center}
Указати послідовність міток вершин дерева, зображеного на рис., що утвориться в результаті обходу дерева в оберненому порядку. Мітки записувати через пробіл.\par Алгоритм починає роботу з кореня (на рис. це верхня вершина).
%enditem


\newpage
13. 
\begin{center}
	\adjustimage{max size={0.9\linewidth}{0.9\paperheight}}
	{Graphs/AdjList/326.png}
\end{center}
Для графа, зображеного на рис., вказати список суміжності вершини 1.\par
Формат оформлення відповіді:\par
список суміжності виписувати через пробіл на одному рядку, якщо пробілів десь буде більше одного - не важливо, упорядкування вершин у списку також значення не має.
%enditem


\newpage
14. Для графа на рисунку вказати послідовність, в якій алгоритм пошуку в глибину буде відмічати відвідані вершини, якщо списки суміжності вершин переглядаються алгоритмом за зростанням міток вершин та вершини відмічаються на першому вход\-женні у вершину.\par
Пошук розпочинається з вершини 4.\par

\begin{center}
	\adjustimage{max size={0.9\linewidth}{0.9\paperheight}}
	{Graphs/Traversals/206.png}
\end{center}
%enditem


\newpage
15. Для графа на рисунку з вершини 6 запущено алгоритм пошуку в ширину, що будує кістякове дерево. Списки суміжності вершин переглядаються алгоритмом за зростанням міток вершин.\par
\begin{center}
	\adjustimage{max size={0.9\linewidth}{0.9\paperheight}}
	{Graphs/Traversals/267.png}
\end{center}
\par Які з наведених ребер входять до складу отриманого кістякового дерева?\par
1) (1, 6)\par
2) (4, 5)\par
3) (3, 4)\par
4) (2, 5)\par
5) (0, 6)\par
6) (3, 5)\par
{\vskip 15pt} Номери відповідних ребер записати за зростанням через пробіл.


%enditem


\newpage
16. 

Рядки журналу через двокрапку містять ім'я користувача (ідентифікатор), час події,тип події, код події, наприклад
\begin{verbatim}
s="username : 2022-05-05 13:26 : ERROR : 404"
\end{verbatim}
у коді 
\begin{verbatim}
res = re.fullmatch('???', s) 
s = ''
code_str = ??? 
\end{verbatim}
чим треба замінити кожен з ???, щоб значенням code\_str був код події? 


\vspace{10mm}
\textbf{Завдання 17} 

Задано програму, що містить код 1-2 класів, можливе успадкування, та клієнтський код.
У наведеному коді невеликі частини закриті. 
Також наведено вихід програми.
Необхідно відновити закриті частини, щоб отриманий код відповідав наведеному виходу програми.


\textbf{Завдання 18-19} 

Кілька запитань за завданнями для самостійної роботи.
\vspace{10mm}
Скоріш за все, що завдань буде трохи менше, а деякі попередні будуть переформульовані в стилі двох останніх.

%enditem

\end {large}}

%end

\newpage
%begin

{{
\begin {large}

\newpage
1. Що буде надруковано за виконання?

\begin{verbatim}
print('R', end=' ')
h=('a','b','c','d','e',0,1)
print(h[-14])
\end{verbatim}

%enditem


\newpage
2. Що буде надруковано за виконання?

\begin{verbatim}
print('R', end=' ')
c='0123456789'
print(c[0:9:-1])
\end{verbatim}

%enditem



\newpage
3. Що буде надруковано за виконання?

\begin{verbatim}
print('R', end=' ')
h=[10,11,12,13,14]
h.insert(2, [1,2])
print(h)
\end{verbatim}

%enditem



\newpage
4. Що буде надруковано за виконання?

\begin{verbatim}
print('R', end=' ')
b=['0','1','2','3','4']
b[:2]=[12]
print(b)
\end{verbatim}

%enditem



\newpage
5. Укажіть ТИП значення виразу.\par
\begin{verbatim}
{y:y*y for y in range(-7,7)}
\end{verbatim}
%enditem


\newpage
6. Що буде надруковано за виконання?\par
\begin{verbatim}
try:
    d = {55:7, 84:2, 38:7, 84:6}
    d[58] = 8
    for x in d.items():
        print(*x, sep=' ', end=' ')
except:
    print('ERR')
\end{verbatim}
%enditem


\newpage
7. Нехай змінна \kw{s} позначає послідовність цілих, по якій можна ітерувати.

Написати вираз, що перевіряє, чи правда, що всі елементи послідовності є від’ємни\-ми.
%enditem


\newpage
8. Що буде надруковано за виконання?
%Оригинал был утерян. Требует отладки!
\begin{verbatim}
_c = 0

class A:
    _c = 1
    
    def __init__(self):
        self._c = 2
        _c = 3
        A._c = 4

obj = A()
obj._c = 5
try:
    print(_c, end=' ')
    print(obj._c, end=' ')
    print(A._c, end=' ')
    print(obj._c, end=' ')
    print(obj._A__c)
except:
    print('ERR')
\end{verbatim}
%enditem


\newpage
9. Що буде надруковано за виконання?
\begin{verbatim}
class B0:
    def b(self): print(1, end=' ')
    def g(self): print(2, end=' ')

class B1(B0):
    def g(self): print(3, end=' ')

try:
    B1().b()
    B1().g()
except:
    print('ERR')
\end{verbatim}
%enditem


\newpage
10. Що буде надруковано за виконання?
\begin{verbatim}
class A:
    c = 1

    def __init__(self):
        b = 2
        c = 3

class B(A):
    c = 4
    a = 5

    def __init__(self):
        super().__init__()
        self.a = 6

try:
    obj = B()
    obj.b = obj.a + 10
    B.b = 13
except: print('EE', end=' ')

try: print(obj.b, end= ' ')
except: print('E1', end=' ')

try: print(A.b, end= ' ')
except: print('E2', end=' ')

try: print(B.b, end= ' ')
except: print('E3', end=' ')
\end{verbatim}
%enditem


\newpage
11. Змінні \kw{next} та \kw{prev} вузла списку позначають наступний та попередній за ним вузол відповідно. Починаючи з голови, у список записано послідовні цілі числа від~$-57$ до~$37$. Нехай змінна \kw{p} позначає вузол, що зберігає значення~$1$.
Яке значення зберігається у вузлі \kw{p.next.prev.prev.prev.prev} ?
%enditem


\newpage
12. 
\begin{center}
	\adjustimage{max size={0.8\linewidth}{0.8\paperheight}}
	{BinTree/trees_exp/33.png}
\end{center}
Указати послідовність міток вершин дерева, зображеного на рис., що утвориться в результаті обходу дерева в оберненому порядку. Мітки записувати через пробіл.\par Алгоритм починає роботу з кореня (на рис. це верхня вершина).
%enditem


\newpage
13. 
\begin{center}
	\adjustimage{max size={0.9\linewidth}{0.9\paperheight}}
	{Graphs/AdjList/197.png}
\end{center}
Для графа, зображеного на рис., вказати список суміжності вершини 0.\par
Формат оформлення відповіді:\par
список суміжності виписувати через пробіл на одному рядку, якщо пробілів десь буде більше одного - не важливо, упорядкування вершин у списку також значення не має.
%enditem


\newpage
14. Для графа на рисунку з вершини 0 запущено алгоритм пошуку в глибину, що будує кістякове дерево. Списки суміжності вершин переглядаються алгоритмом за зростанням міток вершин.\par
\begin{center}
	\adjustimage{max size={0.9\linewidth}{0.9\paperheight}}
	{Graphs/Traversals/59.png}
\end{center}
\par Які з наведених ребер входять до складу отриманого кістякового дерева?\par
1) (3, 4)\par
2) (0, 4)\par
3) (4, 6)\par
4) (3, 6)\par
5) (3, 5)\par
6) (1, 3)\par
{\vskip 15pt} Номери відповідних ребер записати за зростанням через пробіл.

%enditem


\newpage
15. Для графа на рисунку з вершини 3 запущено алгоритм пошуку в ширину, що будує кістякове дерево. Списки суміжності вершин переглядаються алгоритмом за зростанням міток вершин.\par
\begin{center}
	\adjustimage{max size={0.9\linewidth}{0.9\paperheight}}
	{Graphs/Traversals/94.png}
\end{center}
\par Які з наведених ребер входять до складу отриманого кістякового дерева?\par
1) (0, 4)\par
2) (1, 4)\par
3) (0, 1)\par
4) (2, 3)\par
5) (1, 2)\par
6) (3, 5)\par
{\vskip 15pt} Номери відповідних ребер записати за зростанням через пробіл.


%enditem


\newpage
16. 

Рядки журналу через двокрапку містять ім'я користувача (ідентифікатор), час події,тип події, код події, наприклад
\begin{verbatim}
s="username : 2022-05-05 13:26 : ERROR : 404"
\end{verbatim}
у коді 
\begin{verbatim}
res = re.fullmatch('???', s) 
s = ''
code_str = ??? 
\end{verbatim}
чим треба замінити кожен з ???, щоб значенням code\_str був код події? 


\vspace{10mm}
\textbf{Завдання 17} 

Задано програму, що містить код 1-2 класів, можливе успадкування, та клієнтський код.
У наведеному коді невеликі частини закриті. 
Також наведено вихід програми.
Необхідно відновити закриті частини, щоб отриманий код відповідав наведеному виходу програми.


\textbf{Завдання 18-19} 

Кілька запитань за завданнями для самостійної роботи.
\vspace{10mm}
Скоріш за все, що завдань буде трохи менше, а деякі попередні будуть переформульовані в стилі двох останніх.

%enditem

\end {large}}

%end

\newpage
%begin

{{
\begin {large}

\newpage
1. Що буде надруковано за виконання?

\begin{verbatim}
print('R', end=' ')
e=('a','b','c','d','e',0,1,2,4)
print(e[-7])
\end{verbatim}

%enditem


\newpage
2. Що буде надруковано за виконання?

\begin{verbatim}
print('R', end=' ')
h=['0','1','2','3','4','5','6','7','8','9']
print(h[11:-1:3])
\end{verbatim}

%enditem



\newpage
3. Що буде надруковано за виконання?

\begin{verbatim}
print('R', end=' ')
d=['0','1','2','3','4']
d.extend('12')
print(d)
\end{verbatim}

%enditem



\newpage
4. Що буде надруковано за виконання?

\begin{verbatim}
print('R', end=' ')
h=[10,11,12,13,14]
del h[2:-1]
print(h)
\end{verbatim}

%enditem



\newpage
5. Укажіть ТИП значення виразу.\par
\begin{verbatim}
(11,2,3)
\end{verbatim}
%enditem


\newpage
6. Що буде надруковано за виконання?\par
\begin{verbatim}
try:
    s = {43:9, 88:6, 83:6, 80:1, 80:1}
    s[46] = 8
    for x, y in s.items():
        print(x, y, sep=' ', end=' ')
except:
    print('ERR')
\end{verbatim}
%enditem


\newpage
7. Нехай змінна \kw{s} позначає послідовність цілих, по якій можна ітерувати.

Написати вираз, що перевіряє, чи правда, що всі елементи послідовності діляться на 11.
%enditem


\newpage
8. Що буде надруковано за виконання?
\begin{verbatim}
d = 0
k = 1

class C:
    d = 2
    
    def __init__(self):
        self.k = 3
        d = 4
        C.d = 5

obj = C()
try:
    print(d, end=' ')
    print(k, end=' ')
    print(C.d, end=' ')
    print(obj.k)
except:
    print('ERR')
\end{verbatim}
%enditem


\newpage
9. Що буде надруковано за виконання?
\begin{verbatim}
class D0:
    def c(self): print(1, end=' ')
    def h(self): print(2, end=' ')

class D1(D0):
    def c(self): print(3, end=' ')

try:
    D1().h()
    D1().c()
except:
    print('ERR')
\end{verbatim}
%enditem


\newpage
10. Що буде надруковано за виконання?
\begin{verbatim}
class A:
    c = 1

    def g1(self): print(2, end=' ')
    def _g2(self): print(3, end=' ')

    def main(self):
        try: print(self.c, end=' ')
        except: print('E1', end=' ')

        try: self.g1()
        except: print('E2', end=' ')

        try: self._g2()
        except: print('E3', end=' ')


class B(A):
    c = 4

    def g1(self): print(5, end=' ')
    def _g2(self): print(6, end=' ')

try: B().main()
except: print('E4', end=' ')
\end{verbatim}
%enditem


\newpage
11. Змінні \kw{next} та \kw{prev} вузла списку позначають наступний та попередній за ним вузол відповідно. Починаючи з голови, у список записано послідовні цілі числа від~$-37$ до~$49$. Нехай змінна \kw{p} позначає вузол, що зберігає значення~$-8$.
Яке значення зберігається у вузлі \kw{p.prev.next.next.next.next} ?
%enditem


\newpage
12. 
\begin{center}
	\adjustimage{max size={0.8\linewidth}{0.8\paperheight}}
	{BinTree/trees_exp/101.png}
\end{center}
Указати послідовність міток вершин дерева, зображеного на рис., що утвориться в результаті обходу дерева в прямому порядку. Мітки записувати через пробіл.\par Алгоритм починає роботу з кореня (на рис. це верхня вершина).
%enditem


\newpage
13. 
\begin{center}
	\adjustimage{max size={0.9\linewidth}{0.9\paperheight}}
	{Graphs/AdjList/257.png}
\end{center}
Для графа, зображеного на рис., вказати список суміжності вершини 7.\par
Формат оформлення відповіді:\par
список суміжності виписувати через пробіл на одному рядку, якщо пробілів десь буде більше одного - не важливо, упорядкування вершин у списку також значення не має.
%enditem


\newpage
14. Для графа на рисунку вказати послідовність, в якій алгоритм пошуку в глибину буде відмічати відвідані вершини, якщо списки суміжності вершин переглядаються алгоритмом за зростанням міток вершин та вершини відмічаються на першому вход\-женні у вершину.\par
Пошук розпочинається з вершини 4.\par

\begin{center}
	\adjustimage{max size={0.9\linewidth}{0.9\paperheight}}
	{Graphs/Traversals/342.png}
\end{center}
%enditem


\newpage
15. Для графа на рисунку з вершини 1 запущено алгоритм пошуку в ширину, що будує кістякове дерево. Списки суміжності вершин переглядаються алгоритмом за зростанням міток вершин.\par
\begin{center}
	\adjustimage{max size={0.9\linewidth}{0.9\paperheight}}
	{Graphs/Traversals/203.png}
\end{center}
\par Які з наведених ребер входять до складу отриманого кістякового дерева?\par
1) (1, 2)\par
2) (0, 5)\par
3) (0, 3)\par
4) (0, 4)\par
5) (0, 2)\par
6) (4, 5)\par
{\vskip 15pt} Номери відповідних ребер записати за зростанням через пробіл.


%enditem


\newpage
16. 

Рядки журналу через двокрапку містять ім'я користувача (ідентифікатор), час події,тип події, код події, наприклад
\begin{verbatim}
s="username : 2022-05-05 13:26 : ERROR : 404"
\end{verbatim}
у коді 
\begin{verbatim}
res = re.fullmatch('???', s) 
s = ''
code_str = ??? 
\end{verbatim}
чим треба замінити кожен з ???, щоб значенням code\_str був код події? 


\vspace{10mm}
\textbf{Завдання 17} 

Задано програму, що містить код 1-2 класів, можливе успадкування, та клієнтський код.
У наведеному коді невеликі частини закриті. 
Також наведено вихід програми.
Необхідно відновити закриті частини, щоб отриманий код відповідав наведеному виходу програми.


\textbf{Завдання 18-19} 

Кілька запитань за завданнями для самостійної роботи.
\vspace{10mm}
Скоріш за все, що завдань буде трохи менше, а деякі попередні будуть переформульовані в стилі двох останніх.

%enditem

\end {large}}

%end

\newpage
%begin

{{
\begin {large}

\newpage
1. Що буде надруковано за виконання?

\begin{verbatim}
print('R', end=' ')
f=('a','b','c','d','e',0,1,2)
print(f[2])
\end{verbatim}

%enditem


\newpage
2. Що буде надруковано за виконання?

\begin{verbatim}
print('R', end=' ')
b=[0,1,2,3,4,5,6,7,8,9]
print(b[::-3])
\end{verbatim}

%enditem



\newpage
3. Що буде надруковано за виконання?

\begin{verbatim}
print('R', end=' ')
g=['0','1','2','3','4']
g+=[2,3]
print(g)
\end{verbatim}

%enditem



\newpage
4. Що буде надруковано за виконання?

\begin{verbatim}
print('R', end=' ')
a=[10,11,12,13,14]
a[5:8]=[1,2]
print(a)
\end{verbatim}

%enditem



\newpage
5. Укажіть ТИП значення виразу.\par
\begin{verbatim}
{x:0 for x in range(7)}
\end{verbatim}
%enditem


\newpage
6. Що буде надруковано за виконання?\par
\begin{verbatim}
try:
    t = {32:5, 59:8, 68:9, 38:8, 68:9}
    t[84] = 0
    for x in t :
        print(x, sep=' ', end=' ')
except:
    print('ERR')
\end{verbatim}
%enditem


\newpage
7. Нехай змінна \kw{s} позначає дуже довгу послідовність, по якій можна ітерувати. 

Написати вираз-генератор, що задає підпослідовність її елементів, що є числами типу \kw{int}.

Обмеження: усі елементи шуканої підпослідовності одночасно в пам'яті розміщені бути не можуть. 
%enditem


\newpage
8. Що буде надруковано за виконання?
%Оригинал был утерян. Требует отладки!
\begin{verbatim}
a = 0

class D:
    a = 1
    
    def __init__(self):
        self.a = 2
        a = 3
        D.a = 4

obj = D()
D.a = 5
try:
    print(a, end=' ')
    print(obj.a, end=' ')
    print(D.a, end=' ')
    print(obj.a, end=' ')
    print(obj._D__a)
except:
    print('ERR')
\end{verbatim}
%enditem


\newpage
9. Що буде надруковано за виконання?
\begin{verbatim}
class A0:
    def a(self): print(1, end=' ')
    def f(self): print(2, end=' ')

class A1(A0):
    def a(self): print(3, end=' ')

try:
    A1().f()
    A0().a()
except:
    print('ERR')
\end{verbatim}
%enditem


\newpage
10. Що буде надруковано за виконання?
\begin{verbatim}
class A:
    c = 1

    def __init__(self):
        a = 2
        a = 3

class B(A):
    a = 4
    b = 5

    def __init__(self):
        super().__init__()
        self.b = 6

try:
    obj = B()
    obj.c = obj.c + 10
    B.c = 13
except: print('EE', end=' ')

try: print(obj.c, end= ' ')
except: print('E1', end=' ')

try: print(A.c, end= ' ')
except: print('E2', end=' ')

try: print(B.c, end= ' ')
except: print('E3', end=' ')
\end{verbatim}
%enditem


\newpage
11. Змінні \kw{next} та \kw{prev} вузла списку позначають наступний та попередній за ним вузол відповідно. Починаючи з голови, у список записано послідовні цілі числа від~$-30$ до~$39$. Нехай змінна \kw{p} позначає вузол, що зберігає значення~$8$.
Яке значення зберігається у вузлі \kw{p.next.next.prev.prev.next} ?
%enditem


\newpage
12. 
\begin{center}
	\adjustimage{max size={0.8\linewidth}{0.8\paperheight}}
	{BinTree/trees_exp/27.png}
\end{center}
Указати послідовність міток вершин дерева, зображеного на рис., що утвориться в результаті обходу дерева в оберненому порядку. Мітки записувати через пробіл.\par Алгоритм починає роботу з кореня (на рис. це верхня вершина).
%enditem


\newpage
13. 
\begin{center}
	\adjustimage{max size={0.9\linewidth}{0.9\paperheight}}
	{Graphs/AdjList/337.png}
\end{center}
Для графа, зображеного на рис., вказати список суміжності вершини 0.\par
Формат оформлення відповіді:\par
список суміжності виписувати через пробіл на одному рядку, якщо пробілів десь буде більше одного - не важливо, упорядкування вершин у списку також значення не має.
%enditem


\newpage
14. Для графа на рисунку вказати послідовність, в якій алгоритм пошуку в глибину буде відмічати відвідані вершини, якщо списки суміжності вершин переглядаються алгоритмом за зростанням міток вершин та вершини відмічаються на першому вход\-женні у вершину.\par
Пошук розпочинається з вершини 5.\par

\begin{center}
	\adjustimage{max size={0.9\linewidth}{0.9\paperheight}}
	{Graphs/Traversals/132.png}
\end{center}
%enditem


\newpage
15. Для графа на рисунку з вершини 4 запущено алгоритм пошуку в ширину, що будує кістякове дерево. Списки суміжності вершин переглядаються алгоритмом за зростанням міток вершин.\par
\begin{center}
	\adjustimage{max size={0.9\linewidth}{0.9\paperheight}}
	{Graphs/Traversals/302.png}
\end{center}
\par Які з наведених ребер входять до складу отриманого кістякового дерева?\par
1) (1, 3)\par
2) (4, 5)\par
3) (1, 6)\par
4) (2, 6)\par
5) (5, 6)\par
6) (0, 2)\par
{\vskip 15pt} Номери відповідних ребер записати за зростанням через пробіл.


%enditem


\newpage
16. 

Рядки журналу через двокрапку містять ім'я користувача (ідентифікатор), час події,тип події, код події, наприклад
\begin{verbatim}
s="username : 2022-05-05 13:26 : ERROR : 404"
\end{verbatim}
у коді 
\begin{verbatim}
res = re.fullmatch('???', s) 
s = ''
code_str = ??? 
\end{verbatim}
чим треба замінити кожен з ???, щоб значенням code\_str був код події? 


\vspace{10mm}
\textbf{Завдання 17} 

Задано програму, що містить код 1-2 класів, можливе успадкування, та клієнтський код.
У наведеному коді невеликі частини закриті. 
Також наведено вихід програми.
Необхідно відновити закриті частини, щоб отриманий код відповідав наведеному виходу програми.


\textbf{Завдання 18-19} 

Кілька запитань за завданнями для самостійної роботи.
\vspace{10mm}
Скоріш за все, що завдань буде трохи менше, а деякі попередні будуть переформульовані в стилі двох останніх.

%enditem

\end {large}}

%end

\newpage
%begin

{{
\begin {large}

\newpage
1. Що буде надруковано за виконання?

\begin{verbatim}
print('R', end=' ')
d=('a','b','c','d','e',0,1,2,3,4)
print(d[-5])
\end{verbatim}

%enditem


\newpage
2. Що буде надруковано за виконання?

\begin{verbatim}
print('R', end=' ')
g=(0,1,2,3,4,5,6,7,8,9)
print(g[-14::-3])
\end{verbatim}

%enditem



\newpage
3. Що буде надруковано за виконання?

\begin{verbatim}
print('R', end=' ')
e=[10,11,12,13,14]
print(e.pop(-2), end=' ')
print(e)
\end{verbatim}

%enditem



\newpage
4. Що буде надруковано за виконання?

\begin{verbatim}
print('R', end=' ')
d=['0','1','2','3','4']
del d[-5:]
print(d)
\end{verbatim}

%enditem



\newpage
5. Укажіть ТИП значення виразу.\par
\begin{verbatim}
{x:x for x in range(7)}
\end{verbatim}
%enditem


\newpage
6. Що буде надруковано за виконання?\par
\begin{verbatim}
try:
    d = {23:6, 89:8, 55:8, 21:9}
    d[54] = 7
    for x in d.values():
        print(x, sep=' ', end=' ')
except:
    print('ERR')
\end{verbatim}
%enditem


\newpage
7. Нехай змінна \kw{s} позначає дуже довгу послідовність цілих, по якій можна ітерувати. 

Написати вираз-генератор, що задає підпослідовність її елементів, що є додатни\-ми. 

Обмеження: усі елементи шуканої підпослідовності одночасно в пам'яті розміщені бути не можуть. 
%enditem


\newpage
8. Що буде надруковано за виконання?
%Оригинал был утерян. Требует отладки!
\begin{verbatim}
__a = 0

class B:
    __a = 1
    
    def __init__(self):
        self.__a = 2
        __a = 3
        B.__a = 4

obj = B()
obj.__a = 5
try:
    print(__a, end=' ')
    print(obj.__a, end=' ')
    print(B.__a, end=' ')
    print(obj.__a, end=' ')
    print(obj._B__a)
except:
    print('ERR')
\end{verbatim}
%enditem


\newpage
9. Що буде надруковано за виконання?
\begin{verbatim}
class C0:
    def d(self): print(1, end=' ')
    def g(self): print(2, end=' ')

class C1(C0):
    def g(self): print(3, end=' ')

try:
    C0().d()
    C1().g()
except:
    print('ERR')
\end{verbatim}
%enditem


\newpage
10. Що буде надруковано за виконання?
\begin{verbatim}
class A:
    d = 1

    def _h1(self): print(2, end=' ')
    def _h2(self): print(3, end=' ')

    def main(self):
        try: print(self.d, end=' ')
        except: print('E1', end=' ')

        try: self._h1()
        except: print('E2', end=' ')

        try: self._h2()
        except: print('E3', end=' ')


class B(A):
    d = 4

    def _h1(self): print(5, end=' ')
    def _h2(self): print(6, end=' ')

try: B().main()
except: print('E4', end=' ')
\end{verbatim}
%enditem


\newpage
11. Змінні \kw{next} та \kw{prev} вузла списку позначають наступний та попередній за ним вузол відповідно. Починаючи з голови, у список записано послідовні цілі числа від~$-69$ до~$44$. Нехай змінна \kw{p} позначає вузол, що зберігає значення~$-2$.
Яке значення зберігається у вузлі \kw{p.prev.prev.next.next.prev} ?
%enditem


\newpage
12. 
\begin{center}
	\adjustimage{max size={0.8\linewidth}{0.8\paperheight}}
	{BinTree/trees_exp/38.png}
\end{center}
Указати послідовність міток вершин дерева, зображеного на рис., що утвориться в результаті обходу дерева в прямому порядку. Мітки записувати через пробіл.\par Алгоритм починає роботу з кореня (на рис. це верхня вершина).
%enditem


\newpage
13. 
\begin{center}
	\adjustimage{max size={0.9\linewidth}{0.9\paperheight}}
	{Graphs/AdjList/97.png}
\end{center}
Для графа, зображеного на рис., вказати список суміжності вершини 1.\par
Формат оформлення відповіді:\par
список суміжності виписувати через пробіл на одному рядку, якщо пробілів десь буде більше одного - не важливо, упорядкування вершин у списку також значення не має.
%enditem


\newpage
14. Для графа на рисунку з вершини 2 запущено алгоритм пошуку в глибину, що будує кістякове дерево. Списки суміжності вершин переглядаються алгоритмом за зростанням міток вершин.\par
\begin{center}
	\adjustimage{max size={0.9\linewidth}{0.9\paperheight}}
	{Graphs/Traversals/308.png}
\end{center}
\par Які з наведених ребер входять до складу отриманого кістякового дерева?\par
1) (0, 2)\par
2) (3, 6)\par
3) (2, 3)\par
4) (0, 3)\par
5) (0, 6)\par
6) (2, 6)\par
{\vskip 15pt} Номери відповідних ребер записати за зростанням через пробіл.

%enditem


\newpage
15. Для графа на рисунку з вершини 4 запущено алгоритм пошуку в ширину, що будує кістякове дерево. Списки суміжності вершин переглядаються алгоритмом за зростанням міток вершин.\par
\begin{center}
	\adjustimage{max size={0.9\linewidth}{0.9\paperheight}}
	{Graphs/Traversals/90.png}
\end{center}
\par Які з наведених ребер входять до складу отриманого кістякового дерева?\par
1) (0, 6)\par
2) (3, 6)\par
3) (0, 3)\par
4) (4, 5)\par
5) (1, 5)\par
6) (2, 5)\par
{\vskip 15pt} Номери відповідних ребер записати за зростанням через пробіл.


%enditem


\newpage
16. 

Рядки журналу через двокрапку містять ім'я користувача (ідентифікатор), час події,тип події, код події, наприклад
\begin{verbatim}
s="username : 2022-05-05 13:26 : ERROR : 404"
\end{verbatim}
у коді 
\begin{verbatim}
res = re.fullmatch('???', s) 
s = ''
code_str = ??? 
\end{verbatim}
чим треба замінити кожен з ???, щоб значенням code\_str був код події? 


\vspace{10mm}
\textbf{Завдання 17} 

Задано програму, що містить код 1-2 класів, можливе успадкування, та клієнтський код.
У наведеному коді невеликі частини закриті. 
Також наведено вихід програми.
Необхідно відновити закриті частини, щоб отриманий код відповідав наведеному виходу програми.


\textbf{Завдання 18-19} 

Кілька запитань за завданнями для самостійної роботи.
\vspace{10mm}
Скоріш за все, що завдань буде трохи менше, а деякі попередні будуть переформульовані в стилі двох останніх.

%enditem

\end {large}}

%end

\newpage
%begin

{{
\begin {large}

\newpage
1. Що буде надруковано за виконання?

\begin{verbatim}
print('R', end=' ')
h=('a','b','c','d','e',0,1,2)
print(h[-11])
\end{verbatim}

%enditem


\newpage
2. Що буде надруковано за виконання?

\begin{verbatim}
print('R', end=' ')
e=['0','1','2','3','4','5','6','7','8','9']
print(e[7::-2])
\end{verbatim}

%enditem



\newpage
3. Що буде надруковано за виконання?

\begin{verbatim}
print('R', end=' ')
a=[10,11,12,13,14]
a.insert(3, [4,5])
print(a)
\end{verbatim}

%enditem



\newpage
4. Що буде надруковано за виконання?

\begin{verbatim}
print('R', end=' ')
h=[10,11,12,13,14]
h[6:-7]=[4,5]
print(h)
\end{verbatim}

%enditem



\newpage
5. Укажіть ТИП значення виразу.\par
\begin{verbatim}
{x*x:x for x in range(-7,7)}
\end{verbatim}
%enditem


\newpage
6. Що буде надруковано за виконання?\par
\begin{verbatim}
try:
    t = {15:3, 69:9, 30:4, 30:8}
    t[69] = 9
    for x in t.keys():
        print(x, sep=' ', end=' ')
except:
    print('ERR')
\end{verbatim}
%enditem


\newpage
7. Нехай змінна \kw{s} позначає послідовність цілих, по якій можна ітерувати.

Написати вираз, що перевіряє, чи правда, що всі елементи послідовності є додатни\-ми.
%enditem


\newpage
8. Що буде надруковано за виконання?
\begin{verbatim}
c = 0
h = 1

class A:
    h = 2
    
    def __init__(self):
        global h
        self.c = 3
        h = 4
        A.h = 5

obj = A()
try:
    print(c, end=' ')
    print(h, end=' ')
    print(A.h, end=' ')
    print(obj.c)
except:
    print('ERR')
\end{verbatim}
%enditem


\newpage
9. Що буде надруковано за виконання?
\begin{verbatim}
class D0:
    def d(self): print(1, end=' ')
    def f(self): print(2, end=' ')

class D1(D0):
    def d(self): print(3, end=' ')

try:
    D1().f()
    D0().d()
except:
    print('ERR')
\end{verbatim}
%enditem


\newpage
10. Що буде надруковано за виконання?
\begin{verbatim}
class A:
    a = 1

    def __init__(self):
        c = 2
        c = 3

class B(A):
    c = 4
    b = 5

    def __init__(self):
        super().__init__()
        self.b = 6

try:
    obj = B()
    obj.b = obj.a + 10
    B.b = 13
except: print('EE', end=' ')

try: print(obj.b, end= ' ')
except: print('E1', end=' ')

try: print(A.b, end= ' ')
except: print('E2', end=' ')

try: print(B.b, end= ' ')
except: print('E3', end=' ')
\end{verbatim}
%enditem


\newpage
11. Змінні \kw{next} та \kw{prev} вузла списку позначають наступний та попередній за ним вузол відповідно. Починаючи з голови, у список записано послідовні цілі числа від~$-24$ до~$41$. Нехай змінна \kw{p} позначає вузол, що зберігає значення~$5$.
Яке значення зберігається у вузлі \kw{p.prev.prev.next.next.next} ?
%enditem


\newpage
12. 
\begin{center}
	\adjustimage{max size={0.8\linewidth}{0.8\paperheight}}
	{BinTree/trees_exp/125.png}
\end{center}
Указати послідовність міток вершин дерева, зображеного на рис., що утвориться в результаті обходу дерева в оберненому порядку. Мітки записувати через пробіл.\par Алгоритм починає роботу з кореня (на рис. це верхня вершина).
%enditem


\newpage
13. 
\begin{center}
	\adjustimage{max size={0.9\linewidth}{0.9\paperheight}}
	{Graphs/AdjList/80.png}
\end{center}
Для графа, зображеного на рис., вказати список суміжності вершини 0.\par
Формат оформлення відповіді:\par
список суміжності виписувати через пробіл на одному рядку, якщо пробілів десь буде більше одного - не важливо, упорядкування вершин у списку також значення не має.
%enditem


\newpage
14. Для графа на рисунку з вершини 2 запущено алгоритм пошуку в глибину, що будує кістякове дерево. Списки суміжності вершин переглядаються алгоритмом за зростанням міток вершин.\par
\begin{center}
	\adjustimage{max size={0.9\linewidth}{0.9\paperheight}}
	{Graphs/Traversals/173.png}
\end{center}
\par Які з наведених ребер входять до складу отриманого кістякового дерева?\par
1) (2, 3)\par
2) (4, 5)\par
3) (2, 6)\par
4) (0, 4)\par
5) (4, 6)\par
6) (1, 4)\par
{\vskip 15pt} Номери відповідних ребер записати за зростанням через пробіл.

%enditem


\newpage
15. Для графа на рисунку з вершини 2 запущено алгоритм пошуку в ширину, що будує кістякове дерево. Списки суміжності вершин переглядаються алгоритмом за зростанням міток вершин.\par
\begin{center}
	\adjustimage{max size={0.9\linewidth}{0.9\paperheight}}
	{Graphs/Traversals/4.png}
\end{center}
\par Які з наведених ребер входять до складу отриманого кістякового дерева?\par
1) (4, 5)\par
2) (1, 2)\par
3) (1, 3)\par
4) (1, 4)\par
5) (0, 4)\par
6) (3, 4)\par
{\vskip 15pt} Номери відповідних ребер записати за зростанням через пробіл.


%enditem


\newpage
16. 

Рядки журналу через двокрапку містять ім'я користувача (ідентифікатор), час події,тип події, код події, наприклад
\begin{verbatim}
s="username : 2022-05-05 13:26 : ERROR : 404"
\end{verbatim}
у коді 
\begin{verbatim}
res = re.fullmatch('???', s) 
s = ''
code_str = ??? 
\end{verbatim}
чим треба замінити кожен з ???, щоб значенням code\_str був код події? 


\vspace{10mm}
\textbf{Завдання 17} 

Задано програму, що містить код 1-2 класів, можливе успадкування, та клієнтський код.
У наведеному коді невеликі частини закриті. 
Також наведено вихід програми.
Необхідно відновити закриті частини, щоб отриманий код відповідав наведеному виходу програми.


\textbf{Завдання 18-19} 

Кілька запитань за завданнями для самостійної роботи.
\vspace{10mm}
Скоріш за все, що завдань буде трохи менше, а деякі попередні будуть переформульовані в стилі двох останніх.

%enditem

\end {large}}

%end

\newpage
%begin

{{
\begin {large}

\newpage
1. Що буде надруковано за виконання?

\begin{verbatim}
print('R', end=' ')
d=('a','b','c','d','e',0,1,2,4)
print(d[-2])
\end{verbatim}

%enditem


\newpage
2. Що буде надруковано за виконання?

\begin{verbatim}
print('R', end=' ')
f=[0,1,2,3,4,5,6,7,8,9]
print(f[-8:-7:-1])
\end{verbatim}

%enditem



\newpage
3. Що буде надруковано за виконання?

\begin{verbatim}
print('R', end=' ')
d=['0','1','2','3','4']
d.append([1,2])
print(d)
\end{verbatim}

%enditem



\newpage
4. Що буде надруковано за виконання?

\begin{verbatim}
print('R', end=' ')
f=['0','1','2','3','4']
f[0:-3]=[2,3]
print(f)
\end{verbatim}

%enditem



\newpage
5. Укажіть ТИП значення виразу.\par
\begin{verbatim}
{(1,4), 8}
\end{verbatim}
%enditem


\newpage
6. Що буде надруковано за виконання?\par
\begin{verbatim}
try:
    s = {21:7, 59:5, 31:7, 70:9, 21:6}
    s[67] = 2
    for x in s.items():
        print(x, sep=' ', end=' ')
except:
    print('ERR')
\end{verbatim}
%enditem


\newpage
7. Нехай змінна \kw{s} позначає дуже довгу послідовність, по якій можна ітерувати. 

Написати вираз-генератор, що задає підпослідовність її елементів, що не є числами типу \kw{float}.

Обмеження: усі елементи шуканої підпослідовності одночасно в пам'яті розміщені бути не можуть. 
%enditem


\newpage
8. Що буде надруковано за виконання?
%Оригинал был утерян. Требует отладки!
\begin{verbatim}
__c = 0

class C:
    __c = 1
    
    def __init__(self):
        self.__c = 2
        __c = 3
        C.__c = 4

obj = C()
C.__c = 5
try:
    print(__c, end=' ')
    print(obj.__c, end=' ')
    print(C.__c, end=' ')
    print(obj.__c, end=' ')
    print(obj._C__c)
except:
    print('ERR')
\end{verbatim}
%enditem


\newpage
9. Що буде надруковано за виконання?
\begin{verbatim}
class C0:
    def b(self): print(1, end=' ')
    def h(self): print(2, end=' ')

class C1(C0):
    def h(self): print(3, end=' ')

try:
    C1().b()
    C1().h()
except:
    print('ERR')
\end{verbatim}
%enditem


\newpage
10. Що буде надруковано за виконання?
\begin{verbatim}
class A:
    b = 1

    def __init__(self):
        a = 2
        b = 3

class B(A):
    b = 4
    c = 5

    def __init__(self):
        super().__init__()
        self.c = 6

try:
    obj = B()
    obj.c = obj.a + 10
    B.c = 13
except: print('EE', end=' ')

try: print(obj.c, end= ' ')
except: print('E1', end=' ')

try: print(A.c, end= ' ')
except: print('E2', end=' ')

try: print(B.c, end= ' ')
except: print('E3', end=' ')
\end{verbatim}
%enditem


\newpage
11. Змінні \kw{next} та \kw{prev} вузла списку позначають наступний та попередній за ним вузол відповідно. Починаючи з голови, у список записано послідовні цілі числа від~$-64$ до~$16$. Нехай змінна \kw{p} позначає вузол, що зберігає значення~$-4$.
Яке значення зберігається у вузлі \kw{p.next.next.prev.prev.prev} ?
%enditem


\newpage
12. 
\begin{center}
	\adjustimage{max size={0.8\linewidth}{0.8\paperheight}}
	{BinTree/trees_exp/59.png}
\end{center}
Указати послідовність міток вершин дерева, зображеного на рис., що утвориться в результаті обходу дерева в прямому порядку. Мітки записувати через пробіл.\par Алгоритм починає роботу з кореня (на рис. це верхня вершина).
%enditem


\newpage
13. 
\begin{center}
	\adjustimage{max size={0.9\linewidth}{0.9\paperheight}}
	{Graphs/AdjList/323.png}
\end{center}
Для графа, зображеного на рис., вказати список суміжності вершини 6.\par
Формат оформлення відповіді:\par
список суміжності виписувати через пробіл на одному рядку, якщо пробілів десь буде більше одного - не важливо, упорядкування вершин у списку також значення не має.
%enditem


\newpage
14. Для графа на рисунку вказати послідовність, в якій алгоритм пошуку в глибину буде відмічати відвідані вершини, якщо списки суміжності вершин переглядаються алгоритмом за зростанням міток вершин та вершини відмічаються на першому вход\-женні у вершину.\par
Пошук розпочинається з вершини 2.\par

\begin{center}
	\adjustimage{max size={0.9\linewidth}{0.9\paperheight}}
	{Graphs/Traversals/161.png}
\end{center}
%enditem


\newpage
15. Для графа на рисунку з вершини 0 запущено алгоритм пошуку в ширину, що будує кістякове дерево. Списки суміжності вершин переглядаються алгоритмом за зростанням міток вершин.\par
\begin{center}
	\adjustimage{max size={0.9\linewidth}{0.9\paperheight}}
	{Graphs/Traversals/232.png}
\end{center}
\par Які з наведених ребер входять до складу отриманого кістякового дерева?\par
1) (1, 4)\par
2) (1, 2)\par
3) (0, 2)\par
4) (3, 5)\par
5) (3, 4)\par
6) (1, 3)\par
{\vskip 15pt} Номери відповідних ребер записати за зростанням через пробіл.


%enditem


\newpage
16. 

Рядки журналу через двокрапку містять ім'я користувача (ідентифікатор), час події,тип події, код події, наприклад
\begin{verbatim}
s="username : 2022-05-05 13:26 : ERROR : 404"
\end{verbatim}
у коді 
\begin{verbatim}
res = re.fullmatch('???', s) 
s = ''
code_str = ??? 
\end{verbatim}
чим треба замінити кожен з ???, щоб значенням code\_str був код події? 


\vspace{10mm}
\textbf{Завдання 17} 

Задано програму, що містить код 1-2 класів, можливе успадкування, та клієнтський код.
У наведеному коді невеликі частини закриті. 
Також наведено вихід програми.
Необхідно відновити закриті частини, щоб отриманий код відповідав наведеному виходу програми.


\textbf{Завдання 18-19} 

Кілька запитань за завданнями для самостійної роботи.
\vspace{10mm}
Скоріш за все, що завдань буде трохи менше, а деякі попередні будуть переформульовані в стилі двох останніх.

%enditem

\end {large}}

%end

\newpage
%begin

{{
\begin {large}

\newpage
1. Що буде надруковано за виконання?

\begin{verbatim}
print('R', end=' ')
e=('a','b','c','d','e',0,1,2,3)
print(e[-11])
\end{verbatim}

%enditem


\newpage
2. Що буде надруковано за виконання?

\begin{verbatim}
print('R', end=' ')
h='0123456789'
print(h[-13:1:3])
\end{verbatim}

%enditem



\newpage
3. Що буде надруковано за виконання?

\begin{verbatim}
print('R', end=' ')
f=[10,11,12,13,14]
f.extend([1,0])
print(f)
\end{verbatim}

%enditem



\newpage
4. Що буде надруковано за виконання?

\begin{verbatim}
print('R', end=' ')
e=[10,11,12,13,14]
e[-4:5]='10'
print(e)
\end{verbatim}

%enditem



\newpage
5. Укажіть ТИП значення виразу.\par
\begin{verbatim}
{y:None for y in range(-8,3)}
\end{verbatim}
%enditem


\newpage
6. Що буде надруковано за виконання?\par
\begin{verbatim}
try:
    t = {26:3, 11:4, 17:6, 17:2}
    t[63] = 1
    for x in t.items():
        print(x, sep=' ', end=' ')
except:
    print('ERR')
\end{verbatim}
%enditem


\newpage
7. Нехай змінна \kw{s} позначає послідовність цілих, по якій можна ітерувати. 

Написати вираз, що перевіряє, чи правда, що послідовність  містить числа, більші за 5.
%enditem


\newpage
8. Що буде надруковано за виконання?
%Оригинал был утерян. Требует отладки!
\begin{verbatim}
d = 0

class A:
    d = 1
    
    def __init__(self):
        self.d = 2
        d = 3
        A.d = 4

obj = A()
A.d = 5
try:
    print(d, end=' ')
    print(obj.d, end=' ')
    print(A.d, end=' ')
    print(obj.d, end=' ')
    print(obj._A__d)
except:
    print('ERR')
\end{verbatim}
%enditem


\newpage
9. Що буде надруковано за виконання?
\begin{verbatim}
class A0:
    def c(self): print(1, end=' ')
    def g(self): print(2, end=' ')

class A1(A0):
    def c(self): print(3, end=' ')

try:
    A0().g()
    A1().c()
except:
    print('ERR')
\end{verbatim}
%enditem


\newpage
10. Що буде надруковано за виконання?
\begin{verbatim}
class A:
    b = 1

    def __init__(self):
        c = 2
        b = 3

class B(A):
    c = 4
    a = 5

    def __init__(self):
        super().__init__()
        self.a = 6

try:
    obj = B()
    obj.b = obj.c + 10
    B.b = 13
except: print('EE', end=' ')

try: print(obj.b, end= ' ')
except: print('E1', end=' ')

try: print(A.b, end= ' ')
except: print('E2', end=' ')

try: print(B.b, end= ' ')
except: print('E3', end=' ')
\end{verbatim}
%enditem


\newpage
11. Змінні \kw{next} та \kw{prev} вузла списку позначають наступний та попередній за ним вузол відповідно. Починаючи з голови, у список записано послідовні цілі числа від~$-32$ до~$33$. Нехай змінна \kw{p} позначає вузол, що зберігає значення~$-3$.
Яке значення зберігається у вузлі \kw{p.next.prev.next.next.prev} ?
%enditem


\newpage
12. 
\begin{center}
	\adjustimage{max size={0.8\linewidth}{0.8\paperheight}}
	{BinTree/trees_exp/107.png}
\end{center}
Указати послідовність міток вершин дерева, зображеного на рис., що утвориться в результаті обходу дерева в оберненому порядку. Мітки записувати через пробіл.\par Алгоритм починає роботу з кореня (на рис. це верхня вершина).
%enditem


\newpage
13. 
\begin{center}
	\adjustimage{max size={0.9\linewidth}{0.9\paperheight}}
	{Graphs/AdjList/157.png}
\end{center}
Для графа, зображеного на рис., вказати список суміжності вершини 6.\par
Формат оформлення відповіді:\par
список суміжності виписувати через пробіл на одному рядку, якщо пробілів десь буде більше одного - не важливо, упорядкування вершин у списку також значення не має.
%enditem


\newpage
14. Для графа на рисунку вказати послідовність, в якій алгоритм пошуку в глибину буде відмічати відвідані вершини, якщо списки суміжності вершин переглядаються алгоритмом за зростанням міток вершин та вершини відмічаються на першому вход\-женні у вершину.\par
Пошук розпочинається з вершини 1.\par

\begin{center}
	\adjustimage{max size={0.9\linewidth}{0.9\paperheight}}
	{Graphs/Traversals/242.png}
\end{center}
%enditem


\newpage
15. Для графа на рисунку з вершини 4 запущено алгоритм пошуку в ширину, що будує кістякове дерево. Списки суміжності вершин переглядаються алгоритмом за зростанням міток вершин.\par
\begin{center}
	\adjustimage{max size={0.9\linewidth}{0.9\paperheight}}
	{Graphs/Traversals/253.png}
\end{center}
\par Які з наведених ребер входять до складу отриманого кістякового дерева?\par
1) (3, 5)\par
2) (2, 5)\par
3) (3, 4)\par
4) (2, 3)\par
5) (1, 6)\par
6) (0, 6)\par
{\vskip 15pt} Номери відповідних ребер записати за зростанням через пробіл.


%enditem


\newpage
16. 

Рядки журналу через двокрапку містять ім'я користувача (ідентифікатор), час події,тип події, код події, наприклад
\begin{verbatim}
s="username : 2022-05-05 13:26 : ERROR : 404"
\end{verbatim}
у коді 
\begin{verbatim}
res = re.fullmatch('???', s) 
s = ''
code_str = ??? 
\end{verbatim}
чим треба замінити кожен з ???, щоб значенням code\_str був код події? 


\vspace{10mm}
\textbf{Завдання 17} 

Задано програму, що містить код 1-2 класів, можливе успадкування, та клієнтський код.
У наведеному коді невеликі частини закриті. 
Також наведено вихід програми.
Необхідно відновити закриті частини, щоб отриманий код відповідав наведеному виходу програми.


\textbf{Завдання 18-19} 

Кілька запитань за завданнями для самостійної роботи.
\vspace{10mm}
Скоріш за все, що завдань буде трохи менше, а деякі попередні будуть переформульовані в стилі двох останніх.

%enditem

\end {large}}

%end

\newpage
%begin

{{
\begin {large}

\newpage
1. Що буде надруковано за виконання?

\begin{verbatim}
print('R', end=' ')
c=('a','b','c','d','e',0,1)
print(c[5])
\end{verbatim}

%enditem


\newpage
2. Що буде надруковано за виконання?

\begin{verbatim}
print('R', end=' ')
c=('0','1','2','3','4','5','6','7','8','9')
print(c[:4:2])
\end{verbatim}

%enditem



\newpage
3. Що буде надруковано за виконання?

\begin{verbatim}
print('R', end=' ')
c=['0','1','2','3','4']
print(c.index(1), end=' ')
\end{verbatim}

%enditem



\newpage
4. Що буде надруковано за виконання?

\begin{verbatim}
print('R', end=' ')
b=['0','1','2','3','4']
b[1:]='11'
print(b)
\end{verbatim}

%enditem



\newpage
5. Укажіть ТИП значення виразу.\par
\begin{verbatim}
{x:0 for x in range(7)}
\end{verbatim}
%enditem


\newpage
6. Що буде надруковано за виконання?\par
\begin{verbatim}
try:
    d = {79:6, 68:8, 41:3, 68:4}
    d[79] = 5
    for x in d.values():
        print(x, sep=' ', end=' ')
except:
    print('ERR')
\end{verbatim}
%enditem


\newpage
7. Нехай змінна \kw{s} позначає послідовність цілих, по якій можна ітерувати. 

Написати вираз, що перевіряє, чи правда, що послідовність  містить число 35.
%enditem


\newpage
8. Що буде надруковано за виконання?
%Оригинал был утерян. Требует отладки!
\begin{verbatim}
_b = 0

class D:
    _b = 1
    
    def __init__(self):
        self._b = 2
        _b = 3
        D._b = 4

obj = D()
obj._b = 5
try:
    print(_b, end=' ')
    print(obj._b, end=' ')
    print(D._b, end=' ')
    print(obj._b, end=' ')
    print(obj._D__b)
except:
    print('ERR')
\end{verbatim}
%enditem


\newpage
9. Що буде надруковано за виконання?
\begin{verbatim}
class B0:
    def a(self): print(1, end=' ')
    def f(self): print(2, end=' ')

class B1(B0):
    def f(self): print(3, end=' ')

try:
    B1().a()
    B1().f()
except:
    print('ERR')
\end{verbatim}
%enditem


\newpage
10. Що буде надруковано за виконання?
\begin{verbatim}
class A:
    a = 1

    def __init__(self):
        b = 2
        b = 3

class B(A):
    a = 4
    c = 5

    def __init__(self):
        super().__init__()
        self.c = 6

try:
    obj = B()
    obj.b = obj.a + 10
    B.b = 13
except: print('EE', end=' ')

try: print(obj.b, end= ' ')
except: print('E1', end=' ')

try: print(A.b, end= ' ')
except: print('E2', end=' ')

try: print(B.b, end= ' ')
except: print('E3', end=' ')
\end{verbatim}
%enditem


\newpage
11. Змінні \kw{next} та \kw{prev} вузла списку позначають наступний та попередній за ним вузол відповідно. Починаючи з голови, у список записано послідовні цілі числа від~$-21$ до~$35$. Нехай змінна \kw{p} позначає вузол, що зберігає значення~$5$.
Яке значення зберігається у вузлі \kw{p.prev.next.prev.prev.next} ?
%enditem


\newpage
12. 
\begin{center}
	\adjustimage{max size={0.8\linewidth}{0.8\paperheight}}
	{BinTree/trees_exp/44.png}
\end{center}
Указати послідовність міток вершин дерева, зображеного на рис., що утвориться в результаті обходу дерева в прямому порядку. Мітки записувати через пробіл.\par Алгоритм починає роботу з кореня (на рис. це верхня вершина).
%enditem


\newpage
13. 
\begin{center}
	\adjustimage{max size={0.9\linewidth}{0.9\paperheight}}
	{Graphs/AdjList/305.png}
\end{center}
Для графа, зображеного на рис., вказати список суміжності вершини 7.\par
Формат оформлення відповіді:\par
список суміжності виписувати через пробіл на одному рядку, якщо пробілів десь буде більше одного - не важливо, упорядкування вершин у списку також значення не має.
%enditem


\newpage
14. Для графа на рисунку вказати послідовність, в якій алгоритм пошуку в глибину буде відмічати відвідані вершини, якщо списки суміжності вершин переглядаються алгоритмом за зростанням міток вершин та вершини відмічаються на першому вход\-женні у вершину.\par
Пошук розпочинається з вершини 6.\par

\begin{center}
	\adjustimage{max size={0.9\linewidth}{0.9\paperheight}}
	{Graphs/Traversals/124.png}
\end{center}
%enditem


\newpage
15. Для графа на рисунку з вершини 5 запущено алгоритм пошуку в ширину, що будує кістякове дерево. Списки суміжності вершин переглядаються алгоритмом за зростанням міток вершин.\par
\begin{center}
	\adjustimage{max size={0.9\linewidth}{0.9\paperheight}}
	{Graphs/Traversals/31.png}
\end{center}
\par Які з наведених ребер входять до складу отриманого кістякового дерева?\par
1) (0, 5)\par
2) (2, 4)\par
3) (0, 6)\par
4) (0, 4)\par
5) (3, 4)\par
6) (1, 5)\par
{\vskip 15pt} Номери відповідних ребер записати за зростанням через пробіл.


%enditem


\newpage
16. 

Рядки журналу через двокрапку містять ім'я користувача (ідентифікатор), час події,тип події, код події, наприклад
\begin{verbatim}
s="username : 2022-05-05 13:26 : ERROR : 404"
\end{verbatim}
у коді 
\begin{verbatim}
res = re.fullmatch('???', s) 
s = ''
code_str = ??? 
\end{verbatim}
чим треба замінити кожен з ???, щоб значенням code\_str був код події? 


\vspace{10mm}
\textbf{Завдання 17} 

Задано програму, що містить код 1-2 класів, можливе успадкування, та клієнтський код.
У наведеному коді невеликі частини закриті. 
Також наведено вихід програми.
Необхідно відновити закриті частини, щоб отриманий код відповідав наведеному виходу програми.


\textbf{Завдання 18-19} 

Кілька запитань за завданнями для самостійної роботи.
\vspace{10mm}
Скоріш за все, що завдань буде трохи менше, а деякі попередні будуть переформульовані в стилі двох останніх.

%enditem

\end {large}}

%end

\newpage
%begin

{{
\begin {large}

\newpage
1. Що буде надруковано за виконання?

\begin{verbatim}
print('R', end=' ')
a=('a','b','c','d','e',0,1,2)
print(a[-12])
\end{verbatim}

%enditem


\newpage
2. Що буде надруковано за виконання?

\begin{verbatim}
print('R', end=' ')
b=(0,1,2,3,4,5,6,7,8,9)
print(b[:-4:3])
\end{verbatim}

%enditem



\newpage
3. Що буде надруковано за виконання?

\begin{verbatim}
print('R', end=' ')
g=[10,11,12,13,14]
g.remove(14)
print(g)
\end{verbatim}

%enditem



\newpage
4. Що буде надруковано за виконання?

\begin{verbatim}
print('R', end=' ')
c=['0','1','2','3','4']
del c[-2:]
print(c)
\end{verbatim}

%enditem



\newpage
5. Укажіть ТИП значення виразу.\par
\begin{verbatim}
[13,2,4]+[3,5]
\end{verbatim}
%enditem


\newpage
6. Що буде надруковано за виконання?\par
\begin{verbatim}
try:
    s = {65:0, 27:4, 88:6, 85:8}
    s[31] = 8
    for x in s :
        print(x, sep=' ', end=' ')
except:
    print('ERR')
\end{verbatim}
%enditem


\newpage
7. Нехай змінна \kw{s} позначає дуже довгу послідовність, по якій можна ітерувати. 

Написати вираз-генератор, що задає підпослідовність її елементів, що є рядками типу \kw{str}.

Обмеження: усі елементи шуканої підпослідовності одночасно в пам'яті розміщені бути не можуть. 
%enditem


\newpage
8. Що буде надруковано за виконання?
%Оригинал был утерян. Требует отладки!
\begin{verbatim}
_c = 0

class B:
    _c = 1
    
    def __init__(self):
        self._c = 2
        _c = 3
        B._c = 4

obj = B()
B._c = 5
try:
    print(_c, end=' ')
    print(obj._c, end=' ')
    print(B._c, end=' ')
    print(obj._c, end=' ')
    print(obj._B__c)
except:
    print('ERR')
\end{verbatim}
%enditem


\newpage
9. Що буде надруковано за виконання?
\begin{verbatim}
class A0:
    def d(self): print(1, end=' ')
    def h(self): print(2, end=' ')

class A1(A0):
    def h(self): print(3, end=' ')

try:
    A1().d()
    A0().h()
except:
    print('ERR')
\end{verbatim}
%enditem


\newpage
10. Що буде надруковано за виконання?
\begin{verbatim}
class A:
    a = 1

    def _f1(self): print(2, end=' ')
    def f2(self): print(3, end=' ')

    def main(self):
        try: print(self.a, end=' ')
        except: print('E1', end=' ')

        try: self._f1()
        except: print('E2', end=' ')

        try: self.f2()
        except: print('E3', end=' ')


class B(A):
    a = 4

    def _f1(self): print(5, end=' ')
    def f2(self): print(6, end=' ')

try: B().main()
except: print('E4', end=' ')
\end{verbatim}
%enditem


\newpage
11. Змінні \kw{next} та \kw{prev} вузла списку позначають наступний та попередній за ним вузол відповідно. Починаючи з голови, у список записано послідовні цілі числа від~$-31$ до~$15$. Нехай змінна \kw{p} позначає вузол, що зберігає значення~$-1$.
Яке значення зберігається у вузлі \kw{p.prev.prev.next.next.prev} ?
%enditem


\newpage
12. 
\begin{center}
	\adjustimage{max size={0.8\linewidth}{0.8\paperheight}}
	{BinTree/trees_exp/123.png}
\end{center}
Указати послідовність міток вершин дерева, зображеного на рис., що утвориться в результаті обходу дерева в оберненому порядку. Мітки записувати через пробіл.\par Алгоритм починає роботу з кореня (на рис. це верхня вершина).
%enditem


\newpage
13. 
\begin{center}
	\adjustimage{max size={0.9\linewidth}{0.9\paperheight}}
	{Graphs/AdjList/102.png}
\end{center}
Для графа, зображеного на рис., вказати список суміжності вершини 0.\par
Формат оформлення відповіді:\par
список суміжності виписувати через пробіл на одному рядку, якщо пробілів десь буде більше одного - не важливо, упорядкування вершин у списку також значення не має.
%enditem


\newpage
14. Для графа на рисунку вказати послідовність, в якій алгоритм пошуку в глибину буде відмічати відвідані вершини, якщо списки суміжності вершин переглядаються алгоритмом за зростанням міток вершин та вершини відмічаються на першому вход\-женні у вершину.\par
Пошук розпочинається з вершини 3.\par

\begin{center}
	\adjustimage{max size={0.9\linewidth}{0.9\paperheight}}
	{Graphs/Traversals/136.png}
\end{center}
%enditem


\newpage
15. Для графа на рисунку з вершини 2 запущено алгоритм пошуку в ширину, що будує кістякове дерево. Списки суміжності вершин переглядаються алгоритмом за зростанням міток вершин.\par
\begin{center}
	\adjustimage{max size={0.9\linewidth}{0.9\paperheight}}
	{Graphs/Traversals/134.png}
\end{center}
\par Які з наведених ребер входять до складу отриманого кістякового дерева?\par
1) (2, 5)\par
2) (1, 2)\par
3) (0, 1)\par
4) (0, 4)\par
5) (3, 6)\par
6) (2, 6)\par
{\vskip 15pt} Номери відповідних ребер записати за зростанням через пробіл.


%enditem


\newpage
16. 

Рядки журналу через двокрапку містять ім'я користувача (ідентифікатор), час події,тип події, код події, наприклад
\begin{verbatim}
s="username : 2022-05-05 13:26 : ERROR : 404"
\end{verbatim}
у коді 
\begin{verbatim}
res = re.fullmatch('???', s) 
s = ''
code_str = ??? 
\end{verbatim}
чим треба замінити кожен з ???, щоб значенням code\_str був код події? 


\vspace{10mm}
\textbf{Завдання 17} 

Задано програму, що містить код 1-2 класів, можливе успадкування, та клієнтський код.
У наведеному коді невеликі частини закриті. 
Також наведено вихід програми.
Необхідно відновити закриті частини, щоб отриманий код відповідав наведеному виходу програми.


\textbf{Завдання 18-19} 

Кілька запитань за завданнями для самостійної роботи.
\vspace{10mm}
Скоріш за все, що завдань буде трохи менше, а деякі попередні будуть переформульовані в стилі двох останніх.

%enditem

\end {large}}

%end

\newpage
%begin

{{
\begin {large}

\newpage
1. Що буде надруковано за виконання?

\begin{verbatim}
print('R', end=' ')
g=('a','b','c','d','e',0,1,2,3)
print(g[6])
\end{verbatim}

%enditem


\newpage
2. Що буде надруковано за виконання?

\begin{verbatim}
print('R', end=' ')
g=['0','1','2','3','4','5','6','7','8','9']
print(g[::-1])
\end{verbatim}

%enditem



\newpage
3. Що буде надруковано за виконання?

\begin{verbatim}
print('R', end=' ')
b=['0','1','2','3','4']
b+='10'
print(b)
\end{verbatim}

%enditem



\newpage
4. Що буде надруковано за виконання?

\begin{verbatim}
print('R', end=' ')
a=[10,11,12,13,14]
del a[3:-6]
print(a)
\end{verbatim}

%enditem



\newpage
5. Укажіть ТИП значення виразу.\par
\begin{verbatim}
[x for x in range(7)]
\end{verbatim}
%enditem


\newpage
6. Що буде надруковано за виконання?\par
\begin{verbatim}
try:
    s = {65:6, 81:1, 72:0, 72:4}
    s[81] = 1
    for x in s.keys():
        print(x, sep=' ', end=' ')
except:
    print('ERR')
\end{verbatim}
%enditem


\newpage
7. Нехай змінна \kw{s} позначає дуже довгу послідовність, по якій можна ітерувати. 

Написати вираз-генератор, що задає підпослідовність її елементів, що не є числами типу \kw{int}.

Обмеження: усі елементи шуканої підпослідовності одночасно в пам'яті розміщені бути не можуть. 
%enditem


\newpage
8. Що буде надруковано за виконання?
%Оригинал был утерян. Требует отладки!
\begin{verbatim}
b = 0

class A:
    b = 1
    
    def __init__(self):
        self.b = 2
        b = 3
        A.b = 4

obj = A()
obj.b = 5
try:
    print(b, end=' ')
    print(obj.b, end=' ')
    print(A.b, end=' ')
    print(obj.b, end=' ')
    print(obj._A__b)
except:
    print('ERR')
\end{verbatim}
%enditem


\newpage
9. Що буде надруковано за виконання?
\begin{verbatim}
class C0:
    def c(self): print(1, end=' ')
    def f(self): print(2, end=' ')

class C1(C0):
    def f(self): print(3, end=' ')

try:
    C0().c()
    C1().f()
except:
    print('ERR')
\end{verbatim}
%enditem


\newpage
10. Що буде надруковано за виконання?
\begin{verbatim}
class A:
    a = 1

    def __init__(self):
        c = 2
        a = 3

class B(A):
    c = 4
    b = 5

    def __init__(self):
        super().__init__()
        self.b = 6

try:
    obj = B()
    obj.c = obj.b + 10
    B.c = 13
except: print('EE', end=' ')

try: print(obj.c, end= ' ')
except: print('E1', end=' ')

try: print(A.c, end= ' ')
except: print('E2', end=' ')

try: print(B.c, end= ' ')
except: print('E3', end=' ')
\end{verbatim}
%enditem


\newpage
11. Змінні \kw{next} та \kw{prev} вузла списку позначають наступний та попередній за ним вузол відповідно. Починаючи з голови, у список записано послідовні цілі числа від~$-40$ до~$20$. Нехай змінна \kw{p} позначає вузол, що зберігає значення~$-2$.
Яке значення зберігається у вузлі \kw{p.next.next.prev.prev.next} ?
%enditem


\newpage
12. 
\begin{center}
	\adjustimage{max size={0.8\linewidth}{0.8\paperheight}}
	{BinTree/trees_exp/75.png}
\end{center}
Указати послідовність міток вершин дерева, зображеного на рис., що утвориться в результаті обходу дерева в прямому порядку. Мітки записувати через пробіл.\par Алгоритм починає роботу з кореня (на рис. це верхня вершина).
%enditem


\newpage
13. 
\begin{center}
	\adjustimage{max size={0.9\linewidth}{0.9\paperheight}}
	{Graphs/AdjList/54.png}
\end{center}
Для графа, зображеного на рис., вказати список суміжності вершини 9.\par
Формат оформлення відповіді:\par
список суміжності виписувати через пробіл на одному рядку, якщо пробілів десь буде більше одного - не важливо, упорядкування вершин у списку також значення не має.
%enditem


\newpage
14. Для графа на рисунку вказати послідовність, в якій алгоритм пошуку в глибину буде відмічати відвідані вершини, якщо списки суміжності вершин переглядаються алгоритмом за зростанням міток вершин та вершини відмічаються на першому вход\-женні у вершину.\par
Пошук розпочинається з вершини 5.\par

\begin{center}
	\adjustimage{max size={0.9\linewidth}{0.9\paperheight}}
	{Graphs/Traversals/275.png}
\end{center}
%enditem


\newpage
15. Для графа на рисунку з вершини 4 запущено алгоритм пошуку в ширину, що будує кістякове дерево. Списки суміжності вершин переглядаються алгоритмом за зростанням міток вершин.\par
\begin{center}
	\adjustimage{max size={0.9\linewidth}{0.9\paperheight}}
	{Graphs/Traversals/149.png}
\end{center}
\par Які з наведених ребер входять до складу отриманого кістякового дерева?\par
1) (0, 5)\par
2) (5, 6)\par
3) (2, 4)\par
4) (0, 1)\par
5) (0, 4)\par
6) (0, 2)\par
{\vskip 15pt} Номери відповідних ребер записати за зростанням через пробіл.


%enditem


\newpage
16. 

Рядки журналу через двокрапку містять ім'я користувача (ідентифікатор), час події,тип події, код події, наприклад
\begin{verbatim}
s="username : 2022-05-05 13:26 : ERROR : 404"
\end{verbatim}
у коді 
\begin{verbatim}
res = re.fullmatch('???', s) 
s = ''
code_str = ??? 
\end{verbatim}
чим треба замінити кожен з ???, щоб значенням code\_str був код події? 


\vspace{10mm}
\textbf{Завдання 17} 

Задано програму, що містить код 1-2 класів, можливе успадкування, та клієнтський код.
У наведеному коді невеликі частини закриті. 
Також наведено вихід програми.
Необхідно відновити закриті частини, щоб отриманий код відповідав наведеному виходу програми.


\textbf{Завдання 18-19} 

Кілька запитань за завданнями для самостійної роботи.
\vspace{10mm}
Скоріш за все, що завдань буде трохи менше, а деякі попередні будуть переформульовані в стилі двох останніх.

%enditem

\end {large}}

%end

\newpage
%begin

{{
\begin {large}

\newpage
1. Що буде надруковано за виконання?

\begin{verbatim}
print('R', end=' ')
b=('a','b','c','d','e',0,1,2,4)
print(b[-10])
\end{verbatim}

%enditem


\newpage
2. Що буде надруковано за виконання?

\begin{verbatim}
print('R', end=' ')
d='0123456789'
print(d[3:-3:2])
\end{verbatim}

%enditem



\newpage
3. Що буде надруковано за виконання?

\begin{verbatim}
print('R', end=' ')
h=['0','1','2','3','4']
h.insert(-4, 19)
print(h)
\end{verbatim}

%enditem



\newpage
4. Що буде надруковано за виконання?

\begin{verbatim}
print('R', end=' ')
d=['0','1','2','3','4']
d[-6:4]='13'
print(d)
\end{verbatim}

%enditem



\newpage
5. Укажіть ТИП значення виразу.\par
\begin{verbatim}
{x:x for x in range(7)}
\end{verbatim}
%enditem


\newpage
6. Що буде надруковано за виконання?\par
\begin{verbatim}
try:
    t = {77:2, 81:3, 83:1, 70:1, 70:1}
    t[77] = 4
    for x, y in t.items():
        print(x, y, sep=' ', end=' ')
except:
    print('ERR')
\end{verbatim}
%enditem


\newpage
7. Нехай змінна \kw{s} позначає послідовність цілих, по якій можна ітерувати.

Написати вираз, що перевіряє, чи правда, що всі елементи послідовності не діляться на 8.
%enditem


\newpage
8. Що буде надруковано за виконання?
%Оригинал был утерян. Требует отладки!
\begin{verbatim}
__a = 0

class C:
    __a = 1
    
    def __init__(self):
        self.__a = 2
        __a = 3
        C.__a = 4

obj = C()
C.__a = 5
try:
    print(__a, end=' ')
    print(obj.__a, end=' ')
    print(C.__a, end=' ')
    print(obj.__a, end=' ')
    print(obj._C__a)
except:
    print('ERR')
\end{verbatim}
%enditem


\newpage
9. Що буде надруковано за виконання?
\begin{verbatim}
class B0:
    def a(self): print(1, end=' ')
    def h(self): print(2, end=' ')

class B1(B0):
    def h(self): print(3, end=' ')

try:
    B1().a()
    B1().h()
except:
    print('ERR')
\end{verbatim}
%enditem


\newpage
10. Що буде надруковано за виконання?
\begin{verbatim}
class A:
    b = 1

    def __g1(self): print(2, end=' ')
    def _g2(self): print(3, end=' ')

    def main(self):
        try: print(self.b, end=' ')
        except: print('E1', end=' ')

        try: self.__g1()
        except: print('E2', end=' ')

        try: self._g2()
        except: print('E3', end=' ')


class B(A):
    b = 4

    def __g1(self): print(5, end=' ')
    def _g2(self): print(6, end=' ')

try: B().main()
except: print('E4', end=' ')
\end{verbatim}
%enditem


\newpage
11. Змінні \kw{next} та \kw{prev} вузла списку позначають наступний та попередній за ним вузол відповідно. Починаючи з голови, у список записано послідовні цілі числа від~$-18$ до~$19$. Нехай змінна \kw{p} позначає вузол, що зберігає значення~$-4$.
Яке значення зберігається у вузлі \kw{p.prev.next.next.prev.prev} ?
%enditem


\newpage
12. 
\begin{center}
	\adjustimage{max size={0.8\linewidth}{0.8\paperheight}}
	{BinTree/trees_exp/102.png}
\end{center}
Указати послідовність міток вершин дерева, зображеного на рис., що утвориться в результаті обходу дерева в прямому порядку. Мітки записувати через пробіл.\par Алгоритм починає роботу з кореня (на рис. це верхня вершина).
%enditem


\newpage
13. 
\begin{center}
	\adjustimage{max size={0.9\linewidth}{0.9\paperheight}}
	{Graphs/AdjList/320.png}
\end{center}
Для графа, зображеного на рис., вказати список суміжності вершини 1.\par
Формат оформлення відповіді:\par
список суміжності виписувати через пробіл на одному рядку, якщо пробілів десь буде більше одного - не важливо, упорядкування вершин у списку також значення не має.
%enditem


\newpage
14. Для графа на рисунку вказати послідовність, в якій алгоритм пошуку в глибину буде відмічати відвідані вершини, якщо списки суміжності вершин переглядаються алгоритмом за зростанням міток вершин та вершини відмічаються на першому вход\-женні у вершину.\par
Пошук розпочинається з вершини 2.\par

\begin{center}
	\adjustimage{max size={0.9\linewidth}{0.9\paperheight}}
	{Graphs/Traversals/234.png}
\end{center}
%enditem


\newpage
15. Для графа на рисунку з вершини 5 запущено алгоритм пошуку в ширину, що будує кістякове дерево. Списки суміжності вершин переглядаються алгоритмом за зростанням міток вершин.\par
\begin{center}
	\adjustimage{max size={0.9\linewidth}{0.9\paperheight}}
	{Graphs/Traversals/86.png}
\end{center}
\par Які з наведених ребер входять до складу отриманого кістякового дерева?\par
1) (0, 2)\par
2) (3, 5)\par
3) (0, 1)\par
4) (2, 6)\par
5) (0, 4)\par
6) (2, 5)\par
{\vskip 15pt} Номери відповідних ребер записати за зростанням через пробіл.


%enditem


\newpage
16. 

Рядки журналу через двокрапку містять ім'я користувача (ідентифікатор), час події,тип події, код події, наприклад
\begin{verbatim}
s="username : 2022-05-05 13:26 : ERROR : 404"
\end{verbatim}
у коді 
\begin{verbatim}
res = re.fullmatch('???', s) 
s = ''
code_str = ??? 
\end{verbatim}
чим треба замінити кожен з ???, щоб значенням code\_str був код події? 


\vspace{10mm}
\textbf{Завдання 17} 

Задано програму, що містить код 1-2 класів, можливе успадкування, та клієнтський код.
У наведеному коді невеликі частини закриті. 
Також наведено вихід програми.
Необхідно відновити закриті частини, щоб отриманий код відповідав наведеному виходу програми.


\textbf{Завдання 18-19} 

Кілька запитань за завданнями для самостійної роботи.
\vspace{10mm}
Скоріш за все, що завдань буде трохи менше, а деякі попередні будуть переформульовані в стилі двох останніх.

%enditem

\end {large}}

%end

\newpage
%begin

{{
\begin {large}

\newpage
1. Що буде надруковано за виконання?

\begin{verbatim}
print('R', end=' ')
f=('a','b','c','d','e',0,1,2,3,4)
print(f[13])
\end{verbatim}

%enditem


\newpage
2. Що буде надруковано за виконання?

\begin{verbatim}
print('R', end=' ')
a=[0,1,2,3,4,5,6,7,8,9]
print(a[-9:10:-3])
\end{verbatim}

%enditem



\newpage
3. Що буде надруковано за виконання?

\begin{verbatim}
print('R', end=' ')
e=[10,11,12,13,14]
e.append('13')
print(e)
\end{verbatim}

%enditem



\newpage
4. Що буде надруковано за виконання?

\begin{verbatim}
print('R', end=' ')
g=[10,11,12,13,14]
g[4:6]=()
print(g)
\end{verbatim}

%enditem



\newpage
5. Укажіть ТИП значення виразу.\par
\begin{verbatim}
{x for x in range(7) if el%3==1}
\end{verbatim}
%enditem


\newpage
6. Що буде надруковано за виконання?\par
\begin{verbatim}
try:
    d = {14:1, 63:9, 17:9, 17:5}
    d[87] = 3
    for x in d.items():
        print(*x, sep=' ', end=' ')
except:
    print('ERR')
\end{verbatim}
%enditem


\newpage
7. 
 Задано числову послідовність \kw{s} довжини 6000. Записати ВИРАЗ, значенням якого є список, що складається з останньої чверті послідовності, записаної через один елемент.
%enditem


\newpage
8. Що буде надруковано за виконання?
\begin{verbatim}
d = 0
j = 1

class A:
    j = 2
    
    def __init__(self):
        self.j = 3
        j = 4
        A.j = 5

obj = A()
try:
    print(d, end=' ')
    print(j, end=' ')
    print(A.j, end=' ')
    print(obj.j)
except:
    print('ERR')
\end{verbatim}
%enditem


\newpage
9. Що буде надруковано за виконання?
\begin{verbatim}
class D0:
    def b(self): print(1, end=' ')
    def g(self): print(2, end=' ')

class D1(D0):
    def g(self): print(3, end=' ')

try:
    D1().b()
    D1().g()
except:
    print('ERR')
\end{verbatim}
%enditem


\newpage
10. Що буде надруковано за виконання?
\begin{verbatim}
class A:
    d = 1

    def __f1(self): print(2, end=' ')
    def _f2(self): print(3, end=' ')

    def main(self):
        try: print(self.d, end=' ')
        except: print('E1', end=' ')

        try: self.__f1()
        except: print('E2', end=' ')

        try: self._f2()
        except: print('E3', end=' ')


class B(A):
    d = 4

    def __f1(self): print(5, end=' ')
    def _f2(self): print(6, end=' ')

try: B().main()
except: print('E4', end=' ')
\end{verbatim}
%enditem


\newpage
11. Змінні \kw{next} та \kw{prev} вузла списку позначають наступний та попередній за ним вузол відповідно. Починаючи з голови, у список записано послідовні цілі числа від~$-26$ до~$34$. Нехай змінна \kw{p} позначає вузол, що зберігає значення~$-8$.
Яке значення зберігається у вузлі \kw{p.next.prev.prev.next.next} ?
%enditem


\newpage
12. 
\begin{center}
	\adjustimage{max size={0.8\linewidth}{0.8\paperheight}}
	{BinTree/trees_exp/20.png}
\end{center}
Указати послідовність міток вершин дерева, зображеного на рис., що утвориться в результаті обходу дерева в оберненому порядку. Мітки записувати через пробіл.\par Алгоритм починає роботу з кореня (на рис. це верхня вершина).
%enditem


\newpage
13. 
\begin{center}
	\adjustimage{max size={0.9\linewidth}{0.9\paperheight}}
	{Graphs/AdjList/261.png}
\end{center}
Для графа, зображеного на рис., вказати список суміжності вершини 1.\par
Формат оформлення відповіді:\par
список суміжності виписувати через пробіл на одному рядку, якщо пробілів десь буде більше одного - не важливо, упорядкування вершин у списку також значення не має.
%enditem


\newpage
14. Для графа на рисунку вказати послідовність, в якій алгоритм пошуку в глибину буде відмічати відвідані вершини, якщо списки суміжності вершин переглядаються алгоритмом за зростанням міток вершин та вершини відмічаються на першому вход\-женні у вершину.\par
Пошук розпочинається з вершини 1.\par

\begin{center}
	\adjustimage{max size={0.9\linewidth}{0.9\paperheight}}
	{Graphs/Traversals/324.png}
\end{center}
%enditem


\newpage
15. Для графа на рисунку з вершини 5 запущено алгоритм пошуку в ширину, що будує кістякове дерево. Списки суміжності вершин переглядаються алгоритмом за зростанням міток вершин.\par
\begin{center}
	\adjustimage{max size={0.9\linewidth}{0.9\paperheight}}
	{Graphs/Traversals/81.png}
\end{center}
\par Які з наведених ребер входять до складу отриманого кістякового дерева?\par
1) (2, 4)\par
2) (0, 5)\par
3) (0, 4)\par
4) (1, 5)\par
5) (1, 6)\par
6) (3, 4)\par
{\vskip 15pt} Номери відповідних ребер записати за зростанням через пробіл.


%enditem


\newpage
16. 

Рядки журналу через двокрапку містять ім'я користувача (ідентифікатор), час події,тип події, код події, наприклад
\begin{verbatim}
s="username : 2022-05-05 13:26 : ERROR : 404"
\end{verbatim}
у коді 
\begin{verbatim}
res = re.fullmatch('???', s) 
s = ''
code_str = ??? 
\end{verbatim}
чим треба замінити кожен з ???, щоб значенням code\_str був код події? 


\vspace{10mm}
\textbf{Завдання 17} 

Задано програму, що містить код 1-2 класів, можливе успадкування, та клієнтський код.
У наведеному коді невеликі частини закриті. 
Також наведено вихід програми.
Необхідно відновити закриті частини, щоб отриманий код відповідав наведеному виходу програми.


\textbf{Завдання 18-19} 

Кілька запитань за завданнями для самостійної роботи.
\vspace{10mm}
Скоріш за все, що завдань буде трохи менше, а деякі попередні будуть переформульовані в стилі двох останніх.

%enditem

\end {large}}

%end

\newpage
%begin

{{
\begin {large}

\newpage
1. Що буде надруковано за виконання?

\begin{verbatim}
print('R', end=' ')
g=('a','b','c','d','e',0,1)
print(g[-9])
\end{verbatim}

%enditem


\newpage
2. Що буде надруковано за виконання?

\begin{verbatim}
print('R', end=' ')
e=('0','1','2','3','4','5','6','7','8','9')
print(e[:-5:-2])
\end{verbatim}

%enditem



\newpage
3. Що буде надруковано за виконання?

\begin{verbatim}
print('R', end=' ')
d=['0','1','2','3','4']
d.remove('7')
print(d)
\end{verbatim}

%enditem



\newpage
4. Що буде надруковано за виконання?

\begin{verbatim}
print('R', end=' ')
h=[10,11,12,13,14]
h[:6]=[1,0]
print(h)
\end{verbatim}

%enditem



\newpage
5. Укажіть ТИП значення виразу.\par
\begin{verbatim}
{y:13 for y in range(-7,7)}
\end{verbatim}
%enditem


\newpage
6. Що буде надруковано за виконання?\par
\begin{verbatim}
try:
    d = {28:2, 34:9, 44:2, 34:7}
    d[34] = 3
    for x in d.values():
        print(x, sep=' ', end=' ')
except:
    print('ERR')
\end{verbatim}
%enditem


\newpage
7. Нехай змінна \kw{s} позначає дуже довгу послідовність, по якій можна ітерувати. 

Написати вираз-генератор, що задає підпослідовність її елементів, що є числами типу \kw{float}.

Обмеження: усі елементи шуканої підпослідовності одночасно в пам'яті розміщені бути не можуть. 
%enditem


\newpage
8. Що буде надруковано за виконання?
\begin{verbatim}
a = 0
m = 1

class B:
    a = 2
    
    def __init__(self):
        global a
        self.a = 3
        a = 4
        B.a = 5

obj = B()
try:
    print(a, end=' ')
    print(m, end=' ')
    print(B.a, end=' ')
    print(obj.a)
except:
    print('ERR')
\end{verbatim}
%enditem


\newpage
9. Що буде надруковано за виконання?
\begin{verbatim}
class C0:
    def d(self): print(1, end=' ')
    def g(self): print(2, end=' ')

class C1(C0):
    def d(self): print(3, end=' ')

try:
    C0().g()
    C1().d()
except:
    print('ERR')
\end{verbatim}
%enditem


\newpage
10. Що буде надруковано за виконання?
\begin{verbatim}
class A:
    c = 1

    def __init__(self):
        a = 2
        a = 3

class B(A):
    c = 4
    b = 5

    def __init__(self):
        super().__init__()
        self.b = 6

try:
    obj = B()
    obj.a = obj.c + 10
    B.a = 13
except: print('EE', end=' ')

try: print(obj.a, end= ' ')
except: print('E1', end=' ')

try: print(A.a, end= ' ')
except: print('E2', end=' ')

try: print(B.a, end= ' ')
except: print('E3', end=' ')
\end{verbatim}
%enditem


\newpage
11. Змінні \kw{next} та \kw{prev} вузла списку позначають наступний та попередній за ним вузол відповідно. Починаючи з голови, у список записано послідовні цілі числа від~$-63$ до~$45$. Нехай змінна \kw{p} позначає вузол, що зберігає значення~$1$.
Яке значення зберігається у вузлі \kw{p.prev.prev.next.prev.prev} ?
%enditem


\newpage
12. 
\begin{center}
	\adjustimage{max size={0.8\linewidth}{0.8\paperheight}}
	{BinTree/trees_exp/45.png}
\end{center}
Указати послідовність міток вершин дерева, зображеного на рис., що утвориться в результаті обходу дерева в оберненому порядку. Мітки записувати через пробіл.\par Алгоритм починає роботу з кореня (на рис. це верхня вершина).
%enditem


\newpage
13. 
\begin{center}
	\adjustimage{max size={0.9\linewidth}{0.9\paperheight}}
	{Graphs/AdjList/111.png}
\end{center}
Для графа, зображеного на рис., вказати список суміжності вершини 8.\par
Формат оформлення відповіді:\par
список суміжності виписувати через пробіл на одному рядку, якщо пробілів десь буде більше одного - не важливо, упорядкування вершин у списку також значення не має.
%enditem


\newpage
14. Для графа на рисунку з вершини 3 запущено алгоритм пошуку в глибину, що будує кістякове дерево. Списки суміжності вершин переглядаються алгоритмом за зростанням міток вершин.\par
\begin{center}
	\adjustimage{max size={0.9\linewidth}{0.9\paperheight}}
	{Graphs/Traversals/82.png}
\end{center}
\par Які з наведених ребер входять до складу отриманого кістякового дерева?\par
1) (0, 2)\par
2) (0, 3)\par
3) (5, 6)\par
4) (1, 5)\par
5) (4, 6)\par
6) (1, 6)\par
{\vskip 15pt} Номери відповідних ребер записати за зростанням через пробіл.

%enditem


\newpage
15. Для графа на рисунку з вершини 6 запущено алгоритм пошуку в ширину, що будує кістякове дерево. Списки суміжності вершин переглядаються алгоритмом за зростанням міток вершин.\par
\begin{center}
	\adjustimage{max size={0.9\linewidth}{0.9\paperheight}}
	{Graphs/Traversals/84.png}
\end{center}
\par Які з наведених ребер входять до складу отриманого кістякового дерева?\par
1) (3, 4)\par
2) (3, 6)\par
3) (1, 4)\par
4) (2, 4)\par
5) (3, 5)\par
6) (5, 6)\par
{\vskip 15pt} Номери відповідних ребер записати за зростанням через пробіл.


%enditem


\newpage
16. 

Рядки журналу через двокрапку містять ім'я користувача (ідентифікатор), час події,тип події, код події, наприклад
\begin{verbatim}
s="username : 2022-05-05 13:26 : ERROR : 404"
\end{verbatim}
у коді 
\begin{verbatim}
res = re.fullmatch('???', s) 
s = ''
code_str = ??? 
\end{verbatim}
чим треба замінити кожен з ???, щоб значенням code\_str був код події? 


\vspace{10mm}
\textbf{Завдання 17} 

Задано програму, що містить код 1-2 класів, можливе успадкування, та клієнтський код.
У наведеному коді невеликі частини закриті. 
Також наведено вихід програми.
Необхідно відновити закриті частини, щоб отриманий код відповідав наведеному виходу програми.


\textbf{Завдання 18-19} 

Кілька запитань за завданнями для самостійної роботи.
\vspace{10mm}
Скоріш за все, що завдань буде трохи менше, а деякі попередні будуть переформульовані в стилі двох останніх.

%enditem

\end {large}}

%end

\newpage
%begin

{{
\begin {large}

\newpage
1. Що буде надруковано за виконання?

\begin{verbatim}
print('R', end=' ')
a=('a','b','c','d','e',0,1,2,3)
print(a[-6])
\end{verbatim}

%enditem


\newpage
2. Що буде надруковано за виконання?

\begin{verbatim}
print('R', end=' ')
h=(0,1,2,3,4,5,6,7,8,9)
print(h[-12::2])
\end{verbatim}

%enditem



\newpage
3. Що буде надруковано за виконання?

\begin{verbatim}
print('R', end=' ')
e=[10,11,12,13,14]
print(e.pop(), end=' ')
print(e)
\end{verbatim}

%enditem



\newpage
4. Що буде надруковано за виконання?

\begin{verbatim}
print('R', end=' ')
d=['0','1','2','3','4']
d[-5:-8]=[4,5]
print(d)
\end{verbatim}

%enditem



\newpage
5. Укажіть ТИП значення виразу.\par
\begin{verbatim}
{21,2,4}
\end{verbatim}
%enditem


\newpage
6. Що буде надруковано за виконання?\par
\begin{verbatim}
try:
    s = {85:4, 79:9, 29:7, 85:3}
    s[18] = 0
    for x in s.keys():
        print(x, sep=' ', end=' ')
except:
    print('ERR')
\end{verbatim}
%enditem


\newpage
7. Нехай змінна \kw{s} позначає послідовність цілих, по якій можна ітерувати. 

Написати вираз, що перевіряє, чи правда, що послідовність  містить число 13.
%enditem


\newpage
8. Що буде надруковано за виконання?
\begin{verbatim}
c = 0
m = 1

class D:
    c = 2
    
    def __init__(self):
        self.c = 3
        c = 4
        D.c = 5

obj = D()
try:
    print(c, end=' ')
    print(m, end=' ')
    print(D.c, end=' ')
    print(obj.c)
except:
    print('ERR')
\end{verbatim}
%enditem


\newpage
9. Що буде надруковано за виконання?
\begin{verbatim}
class A0:
    def b(self): print(1, end=' ')
    def h(self): print(2, end=' ')

class A1(A0):
    def b(self): print(3, end=' ')

try:
    A1().h()
    A1().b()
except:
    print('ERR')
\end{verbatim}
%enditem


\newpage
10. Що буде надруковано за виконання?
\begin{verbatim}
class A:
    c = 1

    def _h1(self): print(2, end=' ')
    def _h2(self): print(3, end=' ')

    def main(self):
        try: print(self.c, end=' ')
        except: print('E1', end=' ')

        try: self._h1()
        except: print('E2', end=' ')

        try: self._h2()
        except: print('E3', end=' ')


class B(A):
    c = 4

    def _h1(self): print(5, end=' ')
    def _h2(self): print(6, end=' ')

try: B().main()
except: print('E4', end=' ')
\end{verbatim}
%enditem


\newpage
11. Змінні \kw{next} та \kw{prev} вузла списку позначають наступний та попередній за ним вузол відповідно. Починаючи з голови, у список записано послідовні цілі числа від~$-55$ до~$36$. Нехай змінна \kw{p} позначає вузол, що зберігає значення~$3$.
Яке значення зберігається у вузлі \kw{p.next.next.prev.next.next} ?
%enditem


\newpage
12. 
\begin{center}
	\adjustimage{max size={0.8\linewidth}{0.8\paperheight}}
	{BinTree/trees_exp/159.png}
\end{center}
Указати послідовність міток вершин дерева, зображеного на рис., що утвориться в результаті обходу дерева в прямому порядку. Мітки записувати через пробіл.\par Алгоритм починає роботу з кореня (на рис. це верхня вершина).
%enditem


\newpage
13. 
\begin{center}
	\adjustimage{max size={0.9\linewidth}{0.9\paperheight}}
	{Graphs/AdjList/191.png}
\end{center}
Для графа, зображеного на рис., вказати список суміжності вершини 3.\par
Формат оформлення відповіді:\par
список суміжності виписувати через пробіл на одному рядку, якщо пробілів десь буде більше одного - не важливо, упорядкування вершин у списку також значення не має.
%enditem


\newpage
14. Для графа на рисунку з вершини 4 запущено алгоритм пошуку в глибину, що будує кістякове дерево. Списки суміжності вершин переглядаються алгоритмом за зростанням міток вершин.\par
\begin{center}
	\adjustimage{max size={0.9\linewidth}{0.9\paperheight}}
	{Graphs/Traversals/90.png}
\end{center}
\par Які з наведених ребер входять до складу отриманого кістякового дерева?\par
1) (0, 1)\par
2) (2, 5)\par
3) (5, 6)\par
4) (0, 3)\par
5) (2, 6)\par
6) (4, 6)\par
{\vskip 15pt} Номери відповідних ребер записати за зростанням через пробіл.

%enditem


\newpage
15. Для графа на рисунку з вершини 6 запущено алгоритм пошуку в ширину, що будує кістякове дерево. Списки суміжності вершин переглядаються алгоритмом за зростанням міток вершин.\par
\begin{center}
	\adjustimage{max size={0.9\linewidth}{0.9\paperheight}}
	{Graphs/Traversals/181.png}
\end{center}
\par Які з наведених ребер входять до складу отриманого кістякового дерева?\par
1) (0, 1)\par
2) (5, 6)\par
3) (1, 6)\par
4) (0, 5)\par
5) (3, 6)\par
6) (1, 4)\par
{\vskip 15pt} Номери відповідних ребер записати за зростанням через пробіл.


%enditem


\newpage
16. 

Рядки журналу через двокрапку містять ім'я користувача (ідентифікатор), час події,тип події, код події, наприклад
\begin{verbatim}
s="username : 2022-05-05 13:26 : ERROR : 404"
\end{verbatim}
у коді 
\begin{verbatim}
res = re.fullmatch('???', s) 
s = ''
code_str = ??? 
\end{verbatim}
чим треба замінити кожен з ???, щоб значенням code\_str був код події? 


\vspace{10mm}
\textbf{Завдання 17} 

Задано програму, що містить код 1-2 класів, можливе успадкування, та клієнтський код.
У наведеному коді невеликі частини закриті. 
Також наведено вихід програми.
Необхідно відновити закриті частини, щоб отриманий код відповідав наведеному виходу програми.


\textbf{Завдання 18-19} 

Кілька запитань за завданнями для самостійної роботи.
\vspace{10mm}
Скоріш за все, що завдань буде трохи менше, а деякі попередні будуть переформульовані в стилі двох останніх.

%enditem

\end {large}}

%end

\newpage
%begin

{{
\begin {large}

\newpage
1. Що буде надруковано за виконання?

\begin{verbatim}
print('R', end=' ')
b=('a','b','c','d','e',0,1,2,3,4)
print(b[-2])
\end{verbatim}

%enditem


\newpage
2. Що буде надруковано за виконання?

\begin{verbatim}
print('R', end=' ')
g='0123456789'
print(g[-15:-15:-2])
\end{verbatim}

%enditem



\newpage
3. Що буде надруковано за виконання?

\begin{verbatim}
print('R', end=' ')
b=[10,11,12,13,14]
print(b.index(10), end=' ')
\end{verbatim}

%enditem



\newpage
4. Що буде надруковано за виконання?

\begin{verbatim}
print('R', end=' ')
c=[10,11,12,13,14]
del c[:-1]
print(c)
\end{verbatim}

%enditem



\newpage
5. Укажіть ТИП значення виразу.\par
\begin{verbatim}
(2, 7, -4)
\end{verbatim}
%enditem


\newpage
6. Що буде надруковано за виконання?\par
\begin{verbatim}
try:
    t = {23:7, 82:3, 77:3, 43:7, 43:0}
    t[77] = 7
    for x in t :
        print(x, sep=' ', end=' ')
except:
    print('ERR')
\end{verbatim}
%enditem


\newpage
7. 
 Задано числову послідовність \kw{s} довжини 6000. Записати ВИРАЗ, значенням якого є список, що складається з кожного третього елемента послідовності, записаних в обернено\-му порядку.
%enditem


\newpage
8. Що буде надруковано за виконання?
%Оригинал был утерян. Требует отладки!
\begin{verbatim}
_d = 0

class D:
    _d = 1
    
    def __init__(self):
        self._d = 2
        _d = 3
        D._d = 4

obj = D()
obj._d = 5
try:
    print(_d, end=' ')
    print(obj._d, end=' ')
    print(D._d, end=' ')
    print(obj._d, end=' ')
    print(obj._D__d)
except:
    print('ERR')
\end{verbatim}
%enditem


\newpage
9. Що буде надруковано за виконання?
\begin{verbatim}
class B0:
    def a(self): print(1, end=' ')
    def f(self): print(2, end=' ')

class B1(B0):
    def f(self): print(3, end=' ')

try:
    B1().a()
    B0().f()
except:
    print('ERR')
\end{verbatim}
%enditem


\newpage
10. Що буде надруковано за виконання?
\begin{verbatim}
class A:
    c = 1

    def _g1(self): print(2, end=' ')
    def __g2(self): print(3, end=' ')

    def main(self):
        try: print(self.c, end=' ')
        except: print('E1', end=' ')

        try: self._g1()
        except: print('E2', end=' ')

        try: self.__g2()
        except: print('E3', end=' ')


class B(A):
    c = 4

    def _g1(self): print(5, end=' ')
    def __g2(self): print(6, end=' ')

try: B().main()
except: print('E4', end=' ')
\end{verbatim}
%enditem


\newpage
11. Змінні \kw{next} та \kw{prev} вузла списку позначають наступний та попередній за ним вузол відповідно. Починаючи з голови, у список записано послідовні цілі числа від~$-41$ до~$27$. Нехай змінна \kw{p} позначає вузол, що зберігає значення~$8$.
Яке значення зберігається у вузлі \kw{p.prev.next.prev.next.next} ?
%enditem


\newpage
12. 
\begin{center}
	\adjustimage{max size={0.8\linewidth}{0.8\paperheight}}
	{BinTree/trees_exp/17.png}
\end{center}
Указати послідовність міток вершин дерева, зображеного на рис., що утвориться в результаті обходу дерева в прямому порядку. Мітки записувати через пробіл.\par Алгоритм починає роботу з кореня (на рис. це верхня вершина).
%enditem


\newpage
13. 
\begin{center}
	\adjustimage{max size={0.9\linewidth}{0.9\paperheight}}
	{Graphs/AdjList/62.png}
\end{center}
Для графа, зображеного на рис., вказати список суміжності вершини 7.\par
Формат оформлення відповіді:\par
список суміжності виписувати через пробіл на одному рядку, якщо пробілів десь буде більше одного - не важливо, упорядкування вершин у списку також значення не має.
%enditem


\newpage
14. Для графа на рисунку з вершини 3 запущено алгоритм пошуку в глибину, що будує кістякове дерево. Списки суміжності вершин переглядаються алгоритмом за зростанням міток вершин.\par
\begin{center}
	\adjustimage{max size={0.9\linewidth}{0.9\paperheight}}
	{Graphs/Traversals/100.png}
\end{center}
\par Які з наведених ребер входять до складу отриманого кістякового дерева?\par
1) (2, 5)\par
2) (1, 2)\par
3) (2, 4)\par
4) (2, 3)\par
5) (0, 2)\par
6) (0, 6)\par
{\vskip 15pt} Номери відповідних ребер записати за зростанням через пробіл.

%enditem


\newpage
15. Для графа на рисунку з вершини 0 запущено алгоритм пошуку в ширину, що будує кістякове дерево. Списки суміжності вершин переглядаються алгоритмом за зростанням міток вершин.\par
\begin{center}
	\adjustimage{max size={0.9\linewidth}{0.9\paperheight}}
	{Graphs/Traversals/59.png}
\end{center}
\par Які з наведених ребер входять до складу отриманого кістякового дерева?\par
1) (3, 4)\par
2) (3, 6)\par
3) (3, 5)\par
4) (4, 5)\par
5) (0, 4)\par
6) (1, 4)\par
{\vskip 15pt} Номери відповідних ребер записати за зростанням через пробіл.


%enditem


\newpage
16. 

Рядки журналу через двокрапку містять ім'я користувача (ідентифікатор), час події,тип події, код події, наприклад
\begin{verbatim}
s="username : 2022-05-05 13:26 : ERROR : 404"
\end{verbatim}
у коді 
\begin{verbatim}
res = re.fullmatch('???', s) 
s = ''
code_str = ??? 
\end{verbatim}
чим треба замінити кожен з ???, щоб значенням code\_str був код події? 


\vspace{10mm}
\textbf{Завдання 17} 

Задано програму, що містить код 1-2 класів, можливе успадкування, та клієнтський код.
У наведеному коді невеликі частини закриті. 
Також наведено вихід програми.
Необхідно відновити закриті частини, щоб отриманий код відповідав наведеному виходу програми.


\textbf{Завдання 18-19} 

Кілька запитань за завданнями для самостійної роботи.
\vspace{10mm}
Скоріш за все, що завдань буде трохи менше, а деякі попередні будуть переформульовані в стилі двох останніх.

%enditem

\end {large}}

%end

\newpage
%begin

{{
\begin {large}

\newpage
1. Що буде надруковано за виконання?

\begin{verbatim}
print('R', end=' ')
e=('a','b','c','d','e',0,1)
print(e[3])
\end{verbatim}

%enditem


\newpage
2. Що буде надруковано за виконання?

\begin{verbatim}
print('R', end=' ')
b=('0','1','2','3','4','5','6','7','8','9')
print(b[9:3:-3])
\end{verbatim}

%enditem



\newpage
3. Що буде надруковано за виконання?

\begin{verbatim}
print('R', end=' ')
g=['0','1','2','3','4']
g.extend('11')
print(g)
\end{verbatim}

%enditem



\newpage
4. Що буде надруковано за виконання?

\begin{verbatim}
print('R', end=' ')
b=['0','1','2','3','4']
b[-7:]=(13)
print(b)
\end{verbatim}

%enditem



\newpage
5. Укажіть ТИП значення виразу.\par
\begin{verbatim}
[{1,2}, 5]
\end{verbatim}
%enditem


\newpage
6. Що буде надруковано за виконання?\par
\begin{verbatim}
try:
    d = {87:9, 26:1, 34:9, 59:5, 87:9}
    d[46] = 8
    for x, y in d.items():
        print(x, y, sep=' ', end=' ')
except:
    print('ERR')
\end{verbatim}
%enditem


\newpage
7. Нехай змінна \kw{s} позначає дуже довгу послідовність цілих, по якій можна ітерувати. 

Написати вираз-генератор, що задає підпослідовність її елементів, що менші за задане число num. 

Обмеження: усі елементи шуканої підпослідовності одночасно в пам'яті розміщені бути не можуть. 
%enditem


\newpage
8. Що буде надруковано за виконання?
\begin{verbatim}
b = 0
k = 1

class C:
    k = 2
    
    def __init__(self):
        global k
        self.k = 3
        k = 4
        C.k = 5

obj = C()
try:
    print(b, end=' ')
    print(k, end=' ')
    print(C.k, end=' ')
    print(obj.k)
except:
    print('ERR')
\end{verbatim}
%enditem


\newpage
9. Що буде надруковано за виконання?
\begin{verbatim}
class D0:
    def c(self): print(1, end=' ')
    def g(self): print(2, end=' ')

class D1(D0):
    def c(self): print(3, end=' ')

try:
    D1().g()
    D1().c()
except:
    print('ERR')
\end{verbatim}
%enditem


\newpage
10. Що буде надруковано за виконання?
\begin{verbatim}
class A:
    d = 1

    def _f1(self): print(2, end=' ')
    def f2(self): print(3, end=' ')

    def main(self):
        try: print(self.d, end=' ')
        except: print('E1', end=' ')

        try: self._f1()
        except: print('E2', end=' ')

        try: self.f2()
        except: print('E3', end=' ')


class B(A):
    d = 4

    def _f1(self): print(5, end=' ')
    def f2(self): print(6, end=' ')

try: B().main()
except: print('E4', end=' ')
\end{verbatim}
%enditem


\newpage
11. Змінні \kw{next} та \kw{prev} вузла списку позначають наступний та попередній за ним вузол відповідно. Починаючи з голови, у список записано послідовні цілі числа від~$-16$ до~$26$. Нехай змінна \kw{p} позначає вузол, що зберігає значення~$-7$.
Яке значення зберігається у вузлі \kw{p.next.prev.next.prev.prev} ?
%enditem


\newpage
12. 
\begin{center}
	\adjustimage{max size={0.8\linewidth}{0.8\paperheight}}
	{BinTree/trees_exp/161.png}
\end{center}
Указати послідовність міток вершин дерева, зображеного на рис., що утвориться в результаті обходу дерева в оберненому порядку. Мітки записувати через пробіл.\par Алгоритм починає роботу з кореня (на рис. це верхня вершина).
%enditem


\newpage
13. 
\begin{center}
	\adjustimage{max size={0.9\linewidth}{0.9\paperheight}}
	{Graphs/AdjList/226.png}
\end{center}
Для графа, зображеного на рис., вказати список суміжності вершини 6.\par
Формат оформлення відповіді:\par
список суміжності виписувати через пробіл на одному рядку, якщо пробілів десь буде більше одного - не важливо, упорядкування вершин у списку також значення не має.
%enditem


\newpage
14. Для графа на рисунку з вершини 6 запущено алгоритм пошуку в глибину, що будує кістякове дерево. Списки суміжності вершин переглядаються алгоритмом за зростанням міток вершин.\par
\begin{center}
	\adjustimage{max size={0.9\linewidth}{0.9\paperheight}}
	{Graphs/Traversals/67.png}
\end{center}
\par Які з наведених ребер входять до складу отриманого кістякового дерева?\par
1) (0, 5)\par
2) (1, 5)\par
3) (0, 6)\par
4) (0, 2)\par
5) (1, 4)\par
6) (0, 4)\par
{\vskip 15pt} Номери відповідних ребер записати за зростанням через пробіл.

%enditem


\newpage
15. Для графа на рисунку з вершини 1 запущено алгоритм пошуку в ширину, що будує кістякове дерево. Списки суміжності вершин переглядаються алгоритмом за зростанням міток вершин.\par
\begin{center}
	\adjustimage{max size={0.9\linewidth}{0.9\paperheight}}
	{Graphs/Traversals/213.png}
\end{center}
\par Які з наведених ребер входять до складу отриманого кістякового дерева?\par
1) (1, 5)\par
2) (3, 6)\par
3) (1, 2)\par
4) (4, 6)\par
5) (1, 4)\par
6) (3, 4)\par
{\vskip 15pt} Номери відповідних ребер записати за зростанням через пробіл.


%enditem


\newpage
16. 

Рядки журналу через двокрапку містять ім'я користувача (ідентифікатор), час події,тип події, код події, наприклад
\begin{verbatim}
s="username : 2022-05-05 13:26 : ERROR : 404"
\end{verbatim}
у коді 
\begin{verbatim}
res = re.fullmatch('???', s) 
s = ''
code_str = ??? 
\end{verbatim}
чим треба замінити кожен з ???, щоб значенням code\_str був код події? 


\vspace{10mm}
\textbf{Завдання 17} 

Задано програму, що містить код 1-2 класів, можливе успадкування, та клієнтський код.
У наведеному коді невеликі частини закриті. 
Також наведено вихід програми.
Необхідно відновити закриті частини, щоб отриманий код відповідав наведеному виходу програми.


\textbf{Завдання 18-19} 

Кілька запитань за завданнями для самостійної роботи.
\vspace{10mm}
Скоріш за все, що завдань буде трохи менше, а деякі попередні будуть переформульовані в стилі двох останніх.

%enditem

\end {large}}

%end

\newpage
%begin

{{
\begin {large}

\newpage
1. Що буде надруковано за виконання?

\begin{verbatim}
print('R', end=' ')
d=('a','b','c','d','e',0,1,2,4)
print(d[-5])
\end{verbatim}

%enditem


\newpage
2. Що буде надруковано за виконання?

\begin{verbatim}
print('R', end=' ')
d=['0','1','2','3','4','5','6','7','8','9']
print(d[-6::3])
\end{verbatim}

%enditem



\newpage
3. Що буде надруковано за виконання?

\begin{verbatim}
print('R', end=' ')
a=['0','1','2','3','4']
print(a.pop(1), end=' ')
print(a)
\end{verbatim}

%enditem



\newpage
4. Що буде надруковано за виконання?

\begin{verbatim}
print('R', end=' ')
a=[10,11,12,13,14]
a[-9:1]=[15]
print(a)
\end{verbatim}

%enditem



\newpage
5. Укажіть ТИП значення виразу.\par
\begin{verbatim}
{1:3, 5:8}
\end{verbatim}
%enditem


\newpage
6. Що буде надруковано за виконання?\par
\begin{verbatim}
try:
    t = {18:9, 12:7, 35:5, 35:4}
    t[39] = 2
    for x in t.items():
        print(x, sep=' ', end=' ')
except:
    print('ERR')
\end{verbatim}
%enditem


\newpage
7. Нехай змінна \kw{s} позначає дуже довгу послідовність, по якій можна ітерувати. 

Написати вираз-генератор, що задає підпослідовність її елементів, що не є рядками типу \kw{str}.

Обмеження: усі елементи шуканої підпослідовності одночасно в пам'яті розміщені бути не можуть. 
%enditem


\newpage
8. Що буде надруковано за виконання?
\begin{verbatim}
a = 0
j = 1

class A:
    j = 2
    
    def __init__(self):
        global j
        self.a = 3
        j = 4
        A.j = 5

obj = A()
try:
    print(a, end=' ')
    print(j, end=' ')
    print(A.j, end=' ')
    print(obj.a)
except:
    print('ERR')
\end{verbatim}
%enditem


\newpage
9. Що буде надруковано за виконання?
\begin{verbatim}
class A0:
    def a(self): print(1, end=' ')
    def h(self): print(2, end=' ')

class A1(A0):
    def a(self): print(3, end=' ')

try:
    A1().h()
    A1().a()
except:
    print('ERR')
\end{verbatim}
%enditem


\newpage
10. Що буде надруковано за виконання?
\begin{verbatim}
class A:
    a = 1

    def __init__(self):
        b = 2
        a = 3

class B(A):
    b = 4
    c = 5

    def __init__(self):
        super().__init__()
        self.c = 6

try:
    obj = B()
    obj.b = obj.a + 10
    B.b = 13
except: print('EE', end=' ')

try: print(obj.b, end= ' ')
except: print('E1', end=' ')

try: print(A.b, end= ' ')
except: print('E2', end=' ')

try: print(B.b, end= ' ')
except: print('E3', end=' ')
\end{verbatim}
%enditem


\newpage
11. Змінні \kw{next} та \kw{prev} вузла списку позначають наступний та попередній за ним вузол відповідно. Починаючи з голови, у список записано послідовні цілі числа від~$-17$ до~$17$. Нехай змінна \kw{p} позначає вузол, що зберігає значення~$6$.
Яке значення зберігається у вузлі \kw{p.prev.prev.prev.next.next} ?
%enditem


\newpage
12. 
\begin{center}
	\adjustimage{max size={0.8\linewidth}{0.8\paperheight}}
	{BinTree/trees_exp/94.png}
\end{center}
Указати послідовність міток вершин дерева, зображеного на рис., що утвориться в результаті обходу дерева в прямому порядку. Мітки записувати через пробіл.\par Алгоритм починає роботу з кореня (на рис. це верхня вершина).
%enditem


\newpage
13. 
\begin{center}
	\adjustimage{max size={0.9\linewidth}{0.9\paperheight}}
	{Graphs/AdjList/259.png}
\end{center}
Для графа, зображеного на рис., вказати список суміжності вершини 1.\par
Формат оформлення відповіді:\par
список суміжності виписувати через пробіл на одному рядку, якщо пробілів десь буде більше одного - не важливо, упорядкування вершин у списку також значення не має.
%enditem


\newpage
14. Для графа на рисунку вказати послідовність, в якій алгоритм пошуку в глибину буде відмічати відвідані вершини, якщо списки суміжності вершин переглядаються алгоритмом за зростанням міток вершин та вершини відмічаються на першому вход\-женні у вершину.\par
Пошук розпочинається з вершини 4.\par

\begin{center}
	\adjustimage{max size={0.9\linewidth}{0.9\paperheight}}
	{Graphs/Traversals/38.png}
\end{center}
%enditem


\newpage
15. Для графа на рисунку з вершини 2 запущено алгоритм пошуку в ширину, що будує кістякове дерево. Списки суміжності вершин переглядаються алгоритмом за зростанням міток вершин.\par
\begin{center}
	\adjustimage{max size={0.9\linewidth}{0.9\paperheight}}
	{Graphs/Traversals/299.png}
\end{center}
\par Які з наведених ребер входять до складу отриманого кістякового дерева?\par
1) (4, 5)\par
2) (3, 6)\par
3) (0, 6)\par
4) (1, 3)\par
5) (0, 1)\par
6) (1, 5)\par
{\vskip 15pt} Номери відповідних ребер записати за зростанням через пробіл.


%enditem


\newpage
16. 

Рядки журналу через двокрапку містять ім'я користувача (ідентифікатор), час події,тип події, код події, наприклад
\begin{verbatim}
s="username : 2022-05-05 13:26 : ERROR : 404"
\end{verbatim}
у коді 
\begin{verbatim}
res = re.fullmatch('???', s) 
s = ''
code_str = ??? 
\end{verbatim}
чим треба замінити кожен з ???, щоб значенням code\_str був код події? 


\vspace{10mm}
\textbf{Завдання 17} 

Задано програму, що містить код 1-2 класів, можливе успадкування, та клієнтський код.
У наведеному коді невеликі частини закриті. 
Також наведено вихід програми.
Необхідно відновити закриті частини, щоб отриманий код відповідав наведеному виходу програми.


\textbf{Завдання 18-19} 

Кілька запитань за завданнями для самостійної роботи.
\vspace{10mm}
Скоріш за все, що завдань буде трохи менше, а деякі попередні будуть переформульовані в стилі двох останніх.

%enditem

\end {large}}

%end

\newpage
%begin

{{
\begin {large}

\newpage
1. Що буде надруковано за виконання?

\begin{verbatim}
print('R', end=' ')
c=('a','b','c','d','e',0,1,2)
print(c[4])
\end{verbatim}

%enditem


\newpage
2. Що буде надруковано за виконання?

\begin{verbatim}
print('R', end=' ')
c=[0,1,2,3,4,5,6,7,8,9]
print(c[-10:-8:-1])
\end{verbatim}

%enditem



\newpage
3. Що буде надруковано за виконання?

\begin{verbatim}
print('R', end=' ')
c=[10,11,12,13,14]
c.remove(13)
print(c)
\end{verbatim}

%enditem



\newpage
4. Що буде надруковано за виконання?

\begin{verbatim}
print('R', end=' ')
f=['0','1','2','3','4']
del f[-8:7]
print(f)
\end{verbatim}

%enditem



\newpage
5. Укажіть ТИП значення виразу.\par
\begin{verbatim}
{x*x:x for x in range(-7,7)}
\end{verbatim}
%enditem


\newpage
6. Що буде надруковано за виконання?\par
\begin{verbatim}
try:
    s = {65:4, 74:6, 10:5, 46:9, 74:8}
    s[51] = 1
    for x in s.items():
        print(*x, sep=' ', end=' ')
except:
    print('ERR')
\end{verbatim}
%enditem


\newpage
7. Нехай змінна \kw{s} позначає послідовність цілих, по якій можна ітерувати. 

Написати вираз, що перевіряє, чи правда, що послідовність  містить хоча б одне число, що ділиться на 7.
%enditem


\newpage
8. Що буде надруковано за виконання?
%Оригинал был утерян. Требует отладки!
\begin{verbatim}
__a = 0

class B:
    __a = 1
    
    def __init__(self):
        self.__a = 2
        __a = 3
        B.__a = 4

obj = B()
B.__a = 5
try:
    print(__a, end=' ')
    print(obj.__a, end=' ')
    print(B.__a, end=' ')
    print(obj.__a, end=' ')
    print(obj._B__a)
except:
    print('ERR')
\end{verbatim}
%enditem


\newpage
9. Що буде надруковано за виконання?
\begin{verbatim}
class B0:
    def d(self): print(1, end=' ')
    def f(self): print(2, end=' ')

class B1(B0):
    def f(self): print(3, end=' ')

try:
    B1().d()
    B0().f()
except:
    print('ERR')
\end{verbatim}
%enditem


\newpage
10. Що буде надруковано за виконання?
\begin{verbatim}
class A:
    a = 1

    def h1(self): print(2, end=' ')
    def _h2(self): print(3, end=' ')

    def main(self):
        try: print(self.a, end=' ')
        except: print('E1', end=' ')

        try: self.h1()
        except: print('E2', end=' ')

        try: self._h2()
        except: print('E3', end=' ')


class B(A):
    a = 4

    def h1(self): print(5, end=' ')
    def _h2(self): print(6, end=' ')

try: B().main()
except: print('E4', end=' ')
\end{verbatim}
%enditem


\newpage
11. Змінні \kw{next} та \kw{prev} вузла списку позначають наступний та попередній за ним вузол відповідно. Починаючи з голови, у список записано послідовні цілі числа від~$-53$ до~$23$. Нехай змінна \kw{p} позначає вузол, що зберігає значення~$2$.
Яке значення зберігається у вузлі \kw{p.next.next.next.prev.prev} ?
%enditem


\newpage
12. 
\begin{center}
	\adjustimage{max size={0.8\linewidth}{0.8\paperheight}}
	{BinTree/trees_exp/174.png}
\end{center}
Указати послідовність міток вершин дерева, зображеного на рис., що утвориться в результаті обходу дерева в оберненому порядку. Мітки записувати через пробіл.\par Алгоритм починає роботу з кореня (на рис. це верхня вершина).
%enditem


\newpage
13. 
\begin{center}
	\adjustimage{max size={0.9\linewidth}{0.9\paperheight}}
	{Graphs/AdjList/49.png}
\end{center}
Для графа, зображеного на рис., вказати список суміжності вершини 3.\par
Формат оформлення відповіді:\par
список суміжності виписувати через пробіл на одному рядку, якщо пробілів десь буде більше одного - не важливо, упорядкування вершин у списку також значення не має.
%enditem


\newpage
14. Для графа на рисунку з вершини 1 запущено алгоритм пошуку в глибину, що будує кістякове дерево. Списки суміжності вершин переглядаються алгоритмом за зростанням міток вершин.\par
\begin{center}
	\adjustimage{max size={0.9\linewidth}{0.9\paperheight}}
	{Graphs/Traversals/111.png}
\end{center}
\par Які з наведених ребер входять до складу отриманого кістякового дерева?\par
1) (0, 2)\par
2) (0, 1)\par
3) (5, 6)\par
4) (1, 6)\par
5) (2, 6)\par
6) (2, 4)\par
{\vskip 15pt} Номери відповідних ребер записати за зростанням через пробіл.

%enditem


\newpage
15. Для графа на рисунку з вершини 3 запущено алгоритм пошуку в ширину, що будує кістякове дерево. Списки суміжності вершин переглядаються алгоритмом за зростанням міток вершин.\par
\begin{center}
	\adjustimage{max size={0.9\linewidth}{0.9\paperheight}}
	{Graphs/Traversals/102.png}
\end{center}
\par Які з наведених ребер входять до складу отриманого кістякового дерева?\par
1) (4, 6)\par
2) (1, 6)\par
3) (1, 4)\par
4) (3, 5)\par
5) (0, 3)\par
6) (0, 6)\par
{\vskip 15pt} Номери відповідних ребер записати за зростанням через пробіл.


%enditem


\newpage
16. 

Рядки журналу через двокрапку містять ім'я користувача (ідентифікатор), час події,тип події, код події, наприклад
\begin{verbatim}
s="username : 2022-05-05 13:26 : ERROR : 404"
\end{verbatim}
у коді 
\begin{verbatim}
res = re.fullmatch('???', s) 
s = ''
code_str = ??? 
\end{verbatim}
чим треба замінити кожен з ???, щоб значенням code\_str був код події? 


\vspace{10mm}
\textbf{Завдання 17} 

Задано програму, що містить код 1-2 класів, можливе успадкування, та клієнтський код.
У наведеному коді невеликі частини закриті. 
Також наведено вихід програми.
Необхідно відновити закриті частини, щоб отриманий код відповідав наведеному виходу програми.


\textbf{Завдання 18-19} 

Кілька запитань за завданнями для самостійної роботи.
\vspace{10mm}
Скоріш за все, що завдань буде трохи менше, а деякі попередні будуть переформульовані в стилі двох останніх.

%enditem

\end {large}}

%end

\newpage
%begin

{{
\begin {large}

\newpage
1. Що буде надруковано за виконання?

\begin{verbatim}
print('R', end=' ')
f=('a','b','c','d','e',0,1,2,4)
print(f[-11])
\end{verbatim}

%enditem


\newpage
2. Що буде надруковано за виконання?

\begin{verbatim}
print('R', end=' ')
a=[0,1,2,3,4,5,6,7,8,9]
print(a[:-10:-3])
\end{verbatim}

%enditem



\newpage
3. Що буде надруковано за виконання?

\begin{verbatim}
print('R', end=' ')
h=['0','1','2','3','4']
h.append([2,3])
print(h)
\end{verbatim}

%enditem



\newpage
4. Що буде надруковано за виконання?

\begin{verbatim}
print('R', end=' ')
g=['0','1','2','3','4']
g[:-7]=[1,2]
print(g)
\end{verbatim}

%enditem



\newpage
5. Укажіть ТИП значення виразу.\par
\begin{verbatim}
[7,2,3,4]
\end{verbatim}
%enditem


\newpage
6. Що буде надруковано за виконання?\par
\begin{verbatim}
try:
    s = {21:9, 57:5, 89:5, 89:9}
    s[32] = 8
    for x, y in s.items():
        print(x, y, sep=' ', end=' ')
except:
    print('ERR')
\end{verbatim}
%enditem


\newpage
7. Нехай змінна \kw{s} позначає дуже довгу послідовність цілих, по якій можна ітерувати. 

Написати вираз-генератор, що задає підпослідовність її елементів, що є від'єм\-ни\-ми. 

Обмеження: усі елементи шуканої підпослідовності одночасно в пам'яті розміщені бути не можуть. 
%enditem


\newpage
8. Що буде надруковано за виконання?
\begin{verbatim}
c = 0
h = 1

class D:
    c = 2
    
    def __init__(self):
        self.h = 3
        c = 4
        D.c = 5

obj = D()
try:
    print(c, end=' ')
    print(h, end=' ')
    print(D.c, end=' ')
    print(obj.h)
except:
    print('ERR')
\end{verbatim}
%enditem


\newpage
9. Що буде надруковано за виконання?
\begin{verbatim}
class C0:
    def b(self): print(1, end=' ')
    def f(self): print(2, end=' ')

class C1(C0):
    def b(self): print(3, end=' ')

try:
    C1().f()
    C1().b()
except:
    print('ERR')
\end{verbatim}
%enditem


\newpage
10. Що буде надруковано за виконання?
\begin{verbatim}
class A:
    b = 1

    def _g1(self): print(2, end=' ')
    def _g2(self): print(3, end=' ')

    def main(self):
        try: print(self.b, end=' ')
        except: print('E1', end=' ')

        try: self._g1()
        except: print('E2', end=' ')

        try: self._g2()
        except: print('E3', end=' ')


class B(A):
    b = 4

    def _g1(self): print(5, end=' ')
    def _g2(self): print(6, end=' ')

try: B().main()
except: print('E4', end=' ')
\end{verbatim}
%enditem


\newpage
11. Змінні \kw{next} та \kw{prev} вузла списку позначають наступний та попередній за ним вузол відповідно. Починаючи з голови, у список записано послідовні цілі числа від~$-59$ до~$28$. Нехай змінна \kw{p} позначає вузол, що зберігає значення~$-5$.
Яке значення зберігається у вузлі \kw{p.prev.next.prev.next.next} ?
%enditem


\newpage
12. 
\begin{center}
	\adjustimage{max size={0.8\linewidth}{0.8\paperheight}}
	{BinTree/trees_exp/148.png}
\end{center}
Указати послідовність міток вершин дерева, зображеного на рис., що утвориться в результаті обходу дерева в оберненому порядку. Мітки записувати через пробіл.\par Алгоритм починає роботу з кореня (на рис. це верхня вершина).
%enditem


\newpage
13. 
\begin{center}
	\adjustimage{max size={0.9\linewidth}{0.9\paperheight}}
	{Graphs/AdjList/325.png}
\end{center}
Для графа, зображеного на рис., вказати список суміжності вершини 5.\par
Формат оформлення відповіді:\par
список суміжності виписувати через пробіл на одному рядку, якщо пробілів десь буде більше одного - не важливо, упорядкування вершин у списку також значення не має.
%enditem


\newpage
14. Для графа на рисунку з вершини 5 запущено алгоритм пошуку в глибину, що будує кістякове дерево. Списки суміжності вершин переглядаються алгоритмом за зростанням міток вершин.\par
\begin{center}
	\adjustimage{max size={0.9\linewidth}{0.9\paperheight}}
	{Graphs/Traversals/315.png}
\end{center}
\par Які з наведених ребер входять до складу отриманого кістякового дерева?\par
1) (0, 3)\par
2) (2, 3)\par
3) (1, 5)\par
4) (3, 6)\par
5) (1, 2)\par
6) (0, 6)\par
{\vskip 15pt} Номери відповідних ребер записати за зростанням через пробіл.

%enditem


\newpage
15. Для графа на рисунку з вершини 5 запущено алгоритм пошуку в ширину, що будує кістякове дерево. Списки суміжності вершин переглядаються алгоритмом за зростанням міток вершин.\par
\begin{center}
	\adjustimage{max size={0.9\linewidth}{0.9\paperheight}}
	{Graphs/Traversals/178.png}
\end{center}
\par Які з наведених ребер входять до складу отриманого кістякового дерева?\par
1) (2, 5)\par
2) (1, 5)\par
3) (0, 3)\par
4) (0, 4)\par
5) (0, 2)\par
6) (0, 5)\par
{\vskip 15pt} Номери відповідних ребер записати за зростанням через пробіл.


%enditem


\newpage
16. 

Рядки журналу через двокрапку містять ім'я користувача (ідентифікатор), час події,тип події, код події, наприклад
\begin{verbatim}
s="username : 2022-05-05 13:26 : ERROR : 404"
\end{verbatim}
у коді 
\begin{verbatim}
res = re.fullmatch('???', s) 
s = ''
code_str = ??? 
\end{verbatim}
чим треба замінити кожен з ???, щоб значенням code\_str був код події? 


\vspace{10mm}
\textbf{Завдання 17} 

Задано програму, що містить код 1-2 класів, можливе успадкування, та клієнтський код.
У наведеному коді невеликі частини закриті. 
Також наведено вихід програми.
Необхідно відновити закриті частини, щоб отриманий код відповідав наведеному виходу програми.


\textbf{Завдання 18-19} 

Кілька запитань за завданнями для самостійної роботи.
\vspace{10mm}
Скоріш за все, що завдань буде трохи менше, а деякі попередні будуть переформульовані в стилі двох останніх.

%enditem

\end {large}}

%end

\newpage
%begin

{{
\begin {large}

\newpage
1. Що буде надруковано за виконання?

\begin{verbatim}
print('R', end=' ')
h=('a','b','c','d','e',0,1,2,3,4)
print(h[10])
\end{verbatim}

%enditem


\newpage
2. Що буде надруковано за виконання?

\begin{verbatim}
print('R', end=' ')
f=(0,1,2,3,4,5,6,7,8,9)
print(f[-4:7:3])
\end{verbatim}

%enditem



\newpage
3. Що буде надруковано за виконання?

\begin{verbatim}
print('R', end=' ')
f=[10,11,12,13,14]
f.extend('12')
print(f)
\end{verbatim}

%enditem



\newpage
4. Що буде надруковано за виконання?

\begin{verbatim}
print('R', end=' ')
e=[10,11,12,13,14]
e[8:5]='12'
print(e)
\end{verbatim}

%enditem



\newpage
5. Укажіть ТИП значення виразу.\par
\begin{verbatim}
{x for x in range(7)}
\end{verbatim}
%enditem


\newpage
6. Що буде надруковано за виконання?\par
\begin{verbatim}
try:
    t = {72:9, 84:9, 73:7, 28:4, 84:1}
    t[72] = 6
    for x in t.keys():
        print(x, sep=' ', end=' ')
except:
    print('ERR')
\end{verbatim}
%enditem


\newpage
7. Нехай змінна \kw{s} позначає послідовність цілих, по якій можна ітерувати. 

Написати вираз, що перевіряє, чи правда, що послідовність  містить від’ємні числа.
%enditem


\newpage
8. Що буде надруковано за виконання?
\begin{verbatim}
b = 0
n = 1

class B:
    b = 2
    
    def __init__(self):
        self.b = 3
        b = 4
        B.b = 5

obj = B()
try:
    print(b, end=' ')
    print(n, end=' ')
    print(B.b, end=' ')
    print(obj.b)
except:
    print('ERR')
\end{verbatim}
%enditem


\newpage
9. Що буде надруковано за виконання?
\begin{verbatim}
class D0:
    def c(self): print(1, end=' ')
    def h(self): print(2, end=' ')

class D1(D0):
    def c(self): print(3, end=' ')

try:
    D0().h()
    D1().c()
except:
    print('ERR')
\end{verbatim}
%enditem


\newpage
10. Що буде надруковано за виконання?
\begin{verbatim}
class A:
    b = 1

    def __init__(self):
        a = 2
        a = 3

class B(A):
    b = 4
    c = 5

    def __init__(self):
        super().__init__()
        self.c = 6

try:
    obj = B()
    obj.b = obj.c + 10
    B.b = 13
except: print('EE', end=' ')

try: print(obj.b, end= ' ')
except: print('E1', end=' ')

try: print(A.b, end= ' ')
except: print('E2', end=' ')

try: print(B.b, end= ' ')
except: print('E3', end=' ')
\end{verbatim}
%enditem


\newpage
11. Змінні \kw{next} та \kw{prev} вузла списку позначають наступний та попередній за ним вузол відповідно. Починаючи з голови, у список записано послідовні цілі числа від~$-50$ до~$42$. Нехай змінна \kw{p} позначає вузол, що зберігає значення~$7$.
Яке значення зберігається у вузлі \kw{p.next.prev.next.prev.prev} ?
%enditem


\newpage
12. 
\begin{center}
	\adjustimage{max size={0.8\linewidth}{0.8\paperheight}}
	{BinTree/trees_exp/111.png}
\end{center}
Указати послідовність міток вершин дерева, зображеного на рис., що утвориться в результаті обходу дерева в прямому порядку. Мітки записувати через пробіл.\par Алгоритм починає роботу з кореня (на рис. це верхня вершина).
%enditem


\newpage
13. 
\begin{center}
	\adjustimage{max size={0.9\linewidth}{0.9\paperheight}}
	{Graphs/AdjList/333.png}
\end{center}
Для графа, зображеного на рис., вказати список суміжності вершини 2.\par
Формат оформлення відповіді:\par
список суміжності виписувати через пробіл на одному рядку, якщо пробілів десь буде більше одного - не важливо, упорядкування вершин у списку також значення не має.
%enditem


\newpage
14. Для графа на рисунку вказати послідовність, в якій алгоритм пошуку в глибину буде відмічати відвідані вершини, якщо списки суміжності вершин переглядаються алгоритмом за зростанням міток вершин та вершини відмічаються на першому вход\-женні у вершину.\par
Пошук розпочинається з вершини 2.\par

\begin{center}
	\adjustimage{max size={0.9\linewidth}{0.9\paperheight}}
	{Graphs/Traversals/271.png}
\end{center}
%enditem


\newpage
15. Для графа на рисунку з вершини 1 запущено алгоритм пошуку в ширину, що будує кістякове дерево. Списки суміжності вершин переглядаються алгоритмом за зростанням міток вершин.\par
\begin{center}
	\adjustimage{max size={0.9\linewidth}{0.9\paperheight}}
	{Graphs/Traversals/7.png}
\end{center}
\par Які з наведених ребер входять до складу отриманого кістякового дерева?\par
1) (2, 4)\par
2) (0, 2)\par
3) (1, 3)\par
4) (0, 4)\par
5) (1, 2)\par
6) (0, 6)\par
{\vskip 15pt} Номери відповідних ребер записати за зростанням через пробіл.


%enditem


\newpage
16. 

Рядки журналу через двокрапку містять ім'я користувача (ідентифікатор), час події,тип події, код події, наприклад
\begin{verbatim}
s="username : 2022-05-05 13:26 : ERROR : 404"
\end{verbatim}
у коді 
\begin{verbatim}
res = re.fullmatch('???', s) 
s = ''
code_str = ??? 
\end{verbatim}
чим треба замінити кожен з ???, щоб значенням code\_str був код події? 


\vspace{10mm}
\textbf{Завдання 17} 

Задано програму, що містить код 1-2 класів, можливе успадкування, та клієнтський код.
У наведеному коді невеликі частини закриті. 
Також наведено вихід програми.
Необхідно відновити закриті частини, щоб отриманий код відповідав наведеному виходу програми.


\textbf{Завдання 18-19} 

Кілька запитань за завданнями для самостійної роботи.
\vspace{10mm}
Скоріш за все, що завдань буде трохи менше, а деякі попередні будуть переформульовані в стилі двох останніх.

%enditem

\end {large}}

%end

\newpage
%begin

{{
\begin {large}

\newpage
1. Що буде надруковано за виконання?

\begin{verbatim}
print('R', end=' ')
b=('a','b','c','d','e',0,1,2,3)
print(b[1])
\end{verbatim}

%enditem


\newpage
2. Що буде надруковано за виконання?

\begin{verbatim}
print('R', end=' ')
h=['0','1','2','3','4','5','6','7','8','9']
print(h[10:12:2])
\end{verbatim}

%enditem



\newpage
3. Що буде надруковано за виконання?

\begin{verbatim}
print('R', end=' ')
c=[10,11,12,13,14]
print(c.index(16), end=' ')
\end{verbatim}

%enditem



\newpage
4. Що буде надруковано за виконання?

\begin{verbatim}
print('R', end=' ')
f=['0','1','2','3','4']
del f[-1:]
print(f)
\end{verbatim}

%enditem



\newpage
5. Укажіть ТИП значення виразу.\par
\begin{verbatim}
{y:y*y for y in range(-7,7)}
\end{verbatim}
%enditem


\newpage
6. Що буде надруковано за виконання?\par
\begin{verbatim}
try:
    d = {47:5, 82:2, 75:0, 47:3}
    d[50] = 4
    for x in d.items():
        print(x, sep=' ', end=' ')
except:
    print('ERR')
\end{verbatim}
%enditem


\newpage
7. Нехай змінна \kw{s} позначає дуже довгу послідовність, по якій можна ітерувати. 

Написати вираз-генератор, що задає підпослідовність її елементів, що є числами типу \kw{float}.

Обмеження: усі елементи шуканої підпослідовності одночасно в пам'яті розміщені бути не можуть. 
%enditem


\newpage
8. Що буде надруковано за виконання?
%Оригинал был утерян. Требует отладки!
\begin{verbatim}
b = 0

class C:
    b = 1
    
    def __init__(self):
        self.b = 2
        b = 3
        C.b = 4

obj = C()
obj.b = 5
try:
    print(b, end=' ')
    print(obj.b, end=' ')
    print(C.b, end=' ')
    print(obj.b, end=' ')
    print(obj._C__b)
except:
    print('ERR')
\end{verbatim}
%enditem


\newpage
9. Що буде надруковано за виконання?
\begin{verbatim}
class D0:
    def a(self): print(1, end=' ')
    def g(self): print(2, end=' ')

class D1(D0):
    def g(self): print(3, end=' ')

try:
    D1().a()
    D1().g()
except:
    print('ERR')
\end{verbatim}
%enditem


\newpage
10. Що буде надруковано за виконання?
\begin{verbatim}
class A:
    a = 1

    def f1(self): print(2, end=' ')
    def _f2(self): print(3, end=' ')

    def main(self):
        try: print(self.a, end=' ')
        except: print('E1', end=' ')

        try: self.f1()
        except: print('E2', end=' ')

        try: self._f2()
        except: print('E3', end=' ')


class B(A):
    a = 4

    def f1(self): print(5, end=' ')
    def _f2(self): print(6, end=' ')

try: B().main()
except: print('E4', end=' ')
\end{verbatim}
%enditem


\newpage
11. Змінні \kw{next} та \kw{prev} вузла списку позначають наступний та попередній за ним вузол відповідно. Починаючи з голови, у список записано послідовні цілі числа від~$-47$ до~$38$. Нехай змінна \kw{p} позначає вузол, що зберігає значення~$9$.
Яке значення зберігається у вузлі \kw{p.next.next.next.next.prev} ?
%enditem


\newpage
12. 
\begin{center}
	\adjustimage{max size={0.8\linewidth}{0.8\paperheight}}
	{BinTree/trees_exp/50.png}
\end{center}
Указати послідовність міток вершин дерева, зображеного на рис., що утвориться в результаті обходу дерева в оберненому порядку. Мітки записувати через пробіл.\par Алгоритм починає роботу з кореня (на рис. це верхня вершина).
%enditem


\newpage
13. 
\begin{center}
	\adjustimage{max size={0.9\linewidth}{0.9\paperheight}}
	{Graphs/AdjList/271.png}
\end{center}
Для графа, зображеного на рис., вказати список суміжності вершини 0.\par
Формат оформлення відповіді:\par
список суміжності виписувати через пробіл на одному рядку, якщо пробілів десь буде більше одного - не важливо, упорядкування вершин у списку також значення не має.
%enditem


\newpage
14. Для графа на рисунку з вершини 4 запущено алгоритм пошуку в глибину, що будує кістякове дерево. Списки суміжності вершин переглядаються алгоритмом за зростанням міток вершин.\par
\begin{center}
	\adjustimage{max size={0.9\linewidth}{0.9\paperheight}}
	{Graphs/Traversals/139.png}
\end{center}
\par Які з наведених ребер входять до складу отриманого кістякового дерева?\par
1) (4, 5)\par
2) (0, 1)\par
3) (1, 6)\par
4) (1, 4)\par
5) (1, 3)\par
6) (0, 5)\par
{\vskip 15pt} Номери відповідних ребер записати за зростанням через пробіл.

%enditem


\newpage
15. Для графа на рисунку з вершини 2 запущено алгоритм пошуку в ширину, що будує кістякове дерево. Списки суміжності вершин переглядаються алгоритмом за зростанням міток вершин.\par
\begin{center}
	\adjustimage{max size={0.9\linewidth}{0.9\paperheight}}
	{Graphs/Traversals/182.png}
\end{center}
\par Які з наведених ребер входять до складу отриманого кістякового дерева?\par
1) (0, 5)\par
2) (2, 3)\par
3) (4, 5)\par
4) (5, 6)\par
5) (3, 6)\par
6) (1, 4)\par
{\vskip 15pt} Номери відповідних ребер записати за зростанням через пробіл.


%enditem


\newpage
16. 

Рядки журналу через двокрапку містять ім'я користувача (ідентифікатор), час події,тип події, код події, наприклад
\begin{verbatim}
s="username : 2022-05-05 13:26 : ERROR : 404"
\end{verbatim}
у коді 
\begin{verbatim}
res = re.fullmatch('???', s) 
s = ''
code_str = ??? 
\end{verbatim}
чим треба замінити кожен з ???, щоб значенням code\_str був код події? 


\vspace{10mm}
\textbf{Завдання 17} 

Задано програму, що містить код 1-2 класів, можливе успадкування, та клієнтський код.
У наведеному коді невеликі частини закриті. 
Також наведено вихід програми.
Необхідно відновити закриті частини, щоб отриманий код відповідав наведеному виходу програми.


\textbf{Завдання 18-19} 

Кілька запитань за завданнями для самостійної роботи.
\vspace{10mm}
Скоріш за все, що завдань буде трохи менше, а деякі попередні будуть переформульовані в стилі двох останніх.

%enditem

\end {large}}

%end

\newpage
%begin

{{
\begin {large}

\newpage
1. Що буде надруковано за виконання?

\begin{verbatim}
print('R', end=' ')
f=('a','b','c','d','e',0,1)
print(f[-4])
\end{verbatim}

%enditem


\newpage
2. Що буде надруковано за виконання?

\begin{verbatim}
print('R', end=' ')
a=('0','1','2','3','4','5','6','7','8','9')
print(a[-11:0:-1])
\end{verbatim}

%enditem



\newpage
3. Що буде надруковано за виконання?

\begin{verbatim}
print('R', end=' ')
d=['0','1','2','3','4']
d.insert(4, [1,0])
print(d)
\end{verbatim}

%enditem



\newpage
4. Що буде надруковано за виконання?

\begin{verbatim}
print('R', end=' ')
h=[10,11,12,13,14]
h[-4:0]=()
print(h)
\end{verbatim}

%enditem



\newpage
5. Укажіть ТИП значення виразу.\par
\begin{verbatim}
(11,2,3)
\end{verbatim}
%enditem


\newpage
6. Що буде надруковано за виконання?\par
\begin{verbatim}
try:
    d = {30:7, 55:0, 79:3, 88:3, 55:6}
    d[55] = 8
    for x in d :
        print(x, sep=' ', end=' ')
except:
    print('ERR')
\end{verbatim}
%enditem


\newpage
7. Нехай змінна \kw{s} позначає дуже довгу послідовність, по якій можна ітерувати. 

Написати вираз-генератор, що задає підпослідовність її елементів, що не є числами типу \kw{int}.

Обмеження: усі елементи шуканої підпослідовності одночасно в пам'яті розміщені бути не можуть. 
%enditem


\newpage
8. Що буде надруковано за виконання?
%Оригинал был утерян. Требует отладки!
\begin{verbatim}
__d = 0

class D:
    __d = 1
    
    def __init__(self):
        self.__d = 2
        __d = 3
        D.__d = 4

obj = D()
obj.__d = 5
try:
    print(__d, end=' ')
    print(obj.__d, end=' ')
    print(D.__d, end=' ')
    print(obj.__d, end=' ')
    print(obj._D__d)
except:
    print('ERR')
\end{verbatim}
%enditem


\newpage
9. Що буде надруковано за виконання?
\begin{verbatim}
class A0:
    def b(self): print(1, end=' ')
    def g(self): print(2, end=' ')

class A1(A0):
    def g(self): print(3, end=' ')

try:
    A1().b()
    A1().g()
except:
    print('ERR')
\end{verbatim}
%enditem


\newpage
10. Що буде надруковано за виконання?
\begin{verbatim}
class A:
    c = 1

    def _h1(self): print(2, end=' ')
    def h2(self): print(3, end=' ')

    def main(self):
        try: print(self.c, end=' ')
        except: print('E1', end=' ')

        try: self._h1()
        except: print('E2', end=' ')

        try: self.h2()
        except: print('E3', end=' ')


class B(A):
    c = 4

    def _h1(self): print(5, end=' ')
    def h2(self): print(6, end=' ')

try: B().main()
except: print('E4', end=' ')
\end{verbatim}
%enditem


\newpage
11. Змінні \kw{next} та \kw{prev} вузла списку позначають наступний та попередній за ним вузол відповідно. Починаючи з голови, у список записано послідовні цілі числа від~$-62$ до~$43$. Нехай змінна \kw{p} позначає вузол, що зберігає значення~$0$.
Яке значення зберігається у вузлі \kw{p.prev.prev.prev.prev.next} ?
%enditem


\newpage
12. 
\begin{center}
	\adjustimage{max size={0.8\linewidth}{0.8\paperheight}}
	{BinTree/trees_exp/112.png}
\end{center}
Указати послідовність міток вершин дерева, зображеного на рис., що утвориться в результаті обходу дерева в прямому порядку. Мітки записувати через пробіл.\par Алгоритм починає роботу з кореня (на рис. це верхня вершина).
%enditem


\newpage
13. 
\begin{center}
	\adjustimage{max size={0.9\linewidth}{0.9\paperheight}}
	{Graphs/AdjList/138.png}
\end{center}
Для графа, зображеного на рис., вказати список суміжності вершини 8.\par
Формат оформлення відповіді:\par
список суміжності виписувати через пробіл на одному рядку, якщо пробілів десь буде більше одного - не важливо, упорядкування вершин у списку також значення не має.
%enditem


\newpage
14. Для графа на рисунку з вершини 6 запущено алгоритм пошуку в глибину, що будує кістякове дерево. Списки суміжності вершин переглядаються алгоритмом за зростанням міток вершин.\par
\begin{center}
	\adjustimage{max size={0.9\linewidth}{0.9\paperheight}}
	{Graphs/Traversals/184.png}
\end{center}
\par Які з наведених ребер входять до складу отриманого кістякового дерева?\par
1) (0, 3)\par
2) (0, 6)\par
3) (1, 3)\par
4) (1, 5)\par
5) (3, 5)\par
6) (4, 6)\par
{\vskip 15pt} Номери відповідних ребер записати за зростанням через пробіл.

%enditem


\newpage
15. Для графа на рисунку з вершини 0 запущено алгоритм пошуку в ширину, що будує кістякове дерево. Списки суміжності вершин переглядаються алгоритмом за зростанням міток вершин.\par
\begin{center}
	\adjustimage{max size={0.9\linewidth}{0.9\paperheight}}
	{Graphs/Traversals/233.png}
\end{center}
\par Які з наведених ребер входять до складу отриманого кістякового дерева?\par
1) (0, 6)\par
2) (2, 4)\par
3) (0, 3)\par
4) (0, 1)\par
5) (2, 3)\par
6) (0, 2)\par
{\vskip 15pt} Номери відповідних ребер записати за зростанням через пробіл.


%enditem


\newpage
16. 

Рядки журналу через двокрапку містять ім'я користувача (ідентифікатор), час події,тип події, код події, наприклад
\begin{verbatim}
s="username : 2022-05-05 13:26 : ERROR : 404"
\end{verbatim}
у коді 
\begin{verbatim}
res = re.fullmatch('???', s) 
s = ''
code_str = ??? 
\end{verbatim}
чим треба замінити кожен з ???, щоб значенням code\_str був код події? 


\vspace{10mm}
\textbf{Завдання 17} 

Задано програму, що містить код 1-2 класів, можливе успадкування, та клієнтський код.
У наведеному коді невеликі частини закриті. 
Також наведено вихід програми.
Необхідно відновити закриті частини, щоб отриманий код відповідав наведеному виходу програми.


\textbf{Завдання 18-19} 

Кілька запитань за завданнями для самостійної роботи.
\vspace{10mm}
Скоріш за все, що завдань буде трохи менше, а деякі попередні будуть переформульовані в стилі двох останніх.

%enditem

\end {large}}

%end

\newpage
%begin

{{
\begin {large}

\newpage
1. Що буде надруковано за виконання?

\begin{verbatim}
print('R', end=' ')
h=('a','b','c','d','e',0,1,2)
print(h[-13])
\end{verbatim}

%enditem


\newpage
2. Що буде надруковано за виконання?

\begin{verbatim}
print('R', end=' ')
e='0123456789'
print(e[13:8:-2])
\end{verbatim}

%enditem



\newpage
3. Що буде надруковано за виконання?

\begin{verbatim}
print('R', end=' ')
b=['0','1','2','3','4']
b+=[2,3]
print(b)
\end{verbatim}

%enditem



\newpage
4. Що буде надруковано за виконання?

\begin{verbatim}
print('R', end=' ')
b=[10,11,12,13,14]
del b[:-2]
print(b)
\end{verbatim}

%enditem



\newpage
5. Укажіть ТИП значення виразу.\par
\begin{verbatim}
{x for x in range(7)}
\end{verbatim}
%enditem


\newpage
6. Що буде надруковано за виконання?\par
\begin{verbatim}
try:
    t = {28:0, 17:0, 58:4, 58:3}
    t[38] = 2
    for x in t.values():
        print(x, sep=' ', end=' ')
except:
    print('ERR')
\end{verbatim}
%enditem


\newpage
7. Нехай змінна \kw{s} позначає дуже довгу послідовність, по якій можна ітерувати. 

Написати вираз-генератор, що задає підпослідовність її елементів, що є рядками типу \kw{str}.

Обмеження: усі елементи шуканої підпослідовності одночасно в пам'яті розміщені бути не можуть. 
%enditem


\newpage
8. Що буде надруковано за виконання?
\begin{verbatim}
d = 0
k = 1

class C:
    k = 2
    
    def __init__(self):
        global k
        self.k = 3
        k = 4
        C.k = 5

obj = C()
try:
    print(d, end=' ')
    print(k, end=' ')
    print(C.k, end=' ')
    print(obj.k)
except:
    print('ERR')
\end{verbatim}
%enditem


\newpage
9. Що буде надруковано за виконання?
\begin{verbatim}
class B0:
    def d(self): print(1, end=' ')
    def h(self): print(2, end=' ')

class B1(B0):
    def h(self): print(3, end=' ')

try:
    B0().d()
    B1().h()
except:
    print('ERR')
\end{verbatim}
%enditem


\newpage
10. Що буде надруковано за виконання?
\begin{verbatim}
class A:
    d = 1

    def _g1(self): print(2, end=' ')
    def __g2(self): print(3, end=' ')

    def main(self):
        try: print(self.d, end=' ')
        except: print('E1', end=' ')

        try: self._g1()
        except: print('E2', end=' ')

        try: self.__g2()
        except: print('E3', end=' ')


class B(A):
    d = 4

    def _g1(self): print(5, end=' ')
    def __g2(self): print(6, end=' ')

try: B().main()
except: print('E4', end=' ')
\end{verbatim}
%enditem


\newpage
11. Змінні \kw{next} та \kw{prev} вузла списку позначають наступний та попередній за ним вузол відповідно. Починаючи з голови, у список записано послідовні цілі числа від~$-20$ до~$18$. Нехай змінна \kw{p} позначає вузол, що зберігає значення~$-9$.
Яке значення зберігається у вузлі \kw{p.next.prev.next.next.next} ?
%enditem


\newpage
12. 
\begin{center}
	\adjustimage{max size={0.8\linewidth}{0.8\paperheight}}
	{BinTree/trees_exp/51.png}
\end{center}
Указати послідовність міток вершин дерева, зображеного на рис., що утвориться в результаті обходу дерева в оберненому порядку. Мітки записувати через пробіл.\par Алгоритм починає роботу з кореня (на рис. це верхня вершина).
%enditem


\newpage
13. 
\begin{center}
	\adjustimage{max size={0.9\linewidth}{0.9\paperheight}}
	{Graphs/AdjList/79.png}
\end{center}
Для графа, зображеного на рис., вказати список суміжності вершини 2.\par
Формат оформлення відповіді:\par
список суміжності виписувати через пробіл на одному рядку, якщо пробілів десь буде більше одного - не важливо, упорядкування вершин у списку також значення не має.
%enditem


\newpage
14. Для графа на рисунку вказати послідовність, в якій алгоритм пошуку в глибину буде відмічати відвідані вершини, якщо списки суміжності вершин переглядаються алгоритмом за зростанням міток вершин та вершини відмічаються на першому вход\-женні у вершину.\par
Пошук розпочинається з вершини 5.\par

\begin{center}
	\adjustimage{max size={0.9\linewidth}{0.9\paperheight}}
	{Graphs/Traversals/193.png}
\end{center}
%enditem


\newpage
15. Для графа на рисунку з вершини 2 запущено алгоритм пошуку в ширину, що будує кістякове дерево. Списки суміжності вершин переглядаються алгоритмом за зростанням міток вершин.\par
\begin{center}
	\adjustimage{max size={0.9\linewidth}{0.9\paperheight}}
	{Graphs/Traversals/316.png}
\end{center}
\par Які з наведених ребер входять до складу отриманого кістякового дерева?\par
1) (2, 3)\par
2) (5, 6)\par
3) (3, 5)\par
4) (3, 4)\par
5) (0, 4)\par
6) (4, 6)\par
{\vskip 15pt} Номери відповідних ребер записати за зростанням через пробіл.


%enditem


\newpage
16. 

Рядки журналу через двокрапку містять ім'я користувача (ідентифікатор), час події,тип події, код події, наприклад
\begin{verbatim}
s="username : 2022-05-05 13:26 : ERROR : 404"
\end{verbatim}
у коді 
\begin{verbatim}
res = re.fullmatch('???', s) 
s = ''
code_str = ??? 
\end{verbatim}
чим треба замінити кожен з ???, щоб значенням code\_str був код події? 


\vspace{10mm}
\textbf{Завдання 17} 

Задано програму, що містить код 1-2 класів, можливе успадкування, та клієнтський код.
У наведеному коді невеликі частини закриті. 
Також наведено вихід програми.
Необхідно відновити закриті частини, щоб отриманий код відповідав наведеному виходу програми.


\textbf{Завдання 18-19} 

Кілька запитань за завданнями для самостійної роботи.
\vspace{10mm}
Скоріш за все, що завдань буде трохи менше, а деякі попередні будуть переформульовані в стилі двох останніх.

%enditem

\end {large}}

%end

\newpage
%begin

{{
\begin {large}

\newpage
1. Що буде надруковано за виконання?

\begin{verbatim}
print('R', end=' ')
d=('a','b','c','d','e',0,1,2,3)
print(d[-6])
\end{verbatim}

%enditem


\newpage
2. Що буде надруковано за виконання?

\begin{verbatim}
print('R', end=' ')
g=('0','1','2','3','4','5','6','7','8','9')
print(g[-3::-3])
\end{verbatim}

%enditem



\newpage
3. Що буде надруковано за виконання?

\begin{verbatim}
print('R', end=' ')
g=[10,11,12,13,14]
g.append('13')
print(g)
\end{verbatim}

%enditem



\newpage
4. Що буде надруковано за виконання?

\begin{verbatim}
print('R', end=' ')
d=['0','1','2','3','4']
d[-2:8]='11'
print(d)
\end{verbatim}

%enditem



\newpage
5. Укажіть ТИП значення виразу.\par
\begin{verbatim}
{x:0 for x in range(7)}
\end{verbatim}
%enditem


\newpage
6. Що буде надруковано за виконання?\par
\begin{verbatim}
try:
    s = {42:0, 55:3, 54:8, 70:7, 42:3}
    s[39] = 5
    for x in s.items():
        print(*x, sep=' ', end=' ')
except:
    print('ERR')
\end{verbatim}
%enditem


\newpage
7. Нехай змінна \kw{s} позначає дуже довгу послідовність, по якій можна ітерувати. 

Написати вираз-генератор, що задає підпослідовність її елементів, що не є рядками типу \kw{str}.

Обмеження: усі елементи шуканої підпослідовності одночасно в пам'яті розміщені бути не можуть. 
%enditem


\newpage
8. Що буде надруковано за виконання?
\begin{verbatim}
d = 0
j = 1

class A:
    d = 2
    
    def __init__(self):
        global d
        self.j = 3
        d = 4
        A.d = 5

obj = A()
try:
    print(d, end=' ')
    print(j, end=' ')
    print(A.d, end=' ')
    print(obj.j)
except:
    print('ERR')
\end{verbatim}
%enditem


\newpage
9. Що буде надруковано за виконання?
\begin{verbatim}
class C0:
    def c(self): print(1, end=' ')
    def f(self): print(2, end=' ')

class C1(C0):
    def f(self): print(3, end=' ')

try:
    C1().c()
    C0().f()
except:
    print('ERR')
\end{verbatim}
%enditem


\newpage
10. Що буде надруковано за виконання?
\begin{verbatim}
class A:
    b = 1

    def __init__(self):
        c = 2
        c = 3

class B(A):
    c = 4
    a = 5

    def __init__(self):
        super().__init__()
        self.a = 6

try:
    obj = B()
    obj.a = obj.a + 10
    B.a = 13
except: print('EE', end=' ')

try: print(obj.a, end= ' ')
except: print('E1', end=' ')

try: print(A.a, end= ' ')
except: print('E2', end=' ')

try: print(B.a, end= ' ')
except: print('E3', end=' ')
\end{verbatim}
%enditem


\newpage
11. Змінні \kw{next} та \kw{prev} вузла списку позначають наступний та попередній за ним вузол відповідно. Починаючи з голови, у список записано послідовні цілі числа від~$-43$ до~$22$. Нехай змінна \kw{p} позначає вузол, що зберігає значення~$4$.
Яке значення зберігається у вузлі \kw{p.prev.next.prev.prev.prev} ?
%enditem


\newpage
12. 
\begin{center}
	\adjustimage{max size={0.8\linewidth}{0.8\paperheight}}
	{BinTree/trees_exp/139.png}
\end{center}
Указати послідовність міток вершин дерева, зображеного на рис., що утвориться в результаті обходу дерева в прямому порядку. Мітки записувати через пробіл.\par Алгоритм починає роботу з кореня (на рис. це верхня вершина).
%enditem


\newpage
13. 
\begin{center}
	\adjustimage{max size={0.9\linewidth}{0.9\paperheight}}
	{Graphs/AdjList/222.png}
\end{center}
Для графа, зображеного на рис., вказати список суміжності вершини 1.\par
Формат оформлення відповіді:\par
список суміжності виписувати через пробіл на одному рядку, якщо пробілів десь буде більше одного - не важливо, упорядкування вершин у списку також значення не має.
%enditem


\newpage
14. Для графа на рисунку вказати послідовність, в якій алгоритм пошуку в глибину буде відмічати відвідані вершини, якщо списки суміжності вершин переглядаються алгоритмом за зростанням міток вершин та вершини відмічаються на першому вход\-женні у вершину.\par
Пошук розпочинається з вершини 1.\par

\begin{center}
	\adjustimage{max size={0.9\linewidth}{0.9\paperheight}}
	{Graphs/Traversals/258.png}
\end{center}
%enditem


\newpage
15. Для графа на рисунку з вершини 5 запущено алгоритм пошуку в ширину, що будує кістякове дерево. Списки суміжності вершин переглядаються алгоритмом за зростанням міток вершин.\par
\begin{center}
	\adjustimage{max size={0.9\linewidth}{0.9\paperheight}}
	{Graphs/Traversals/346.png}
\end{center}
\par Які з наведених ребер входять до складу отриманого кістякового дерева?\par
1) (0, 1)\par
2) (2, 4)\par
3) (0, 5)\par
4) (1, 2)\par
5) (1, 5)\par
6) (1, 6)\par
{\vskip 15pt} Номери відповідних ребер записати за зростанням через пробіл.


%enditem


\newpage
16. 

Рядки журналу через двокрапку містять ім'я користувача (ідентифікатор), час події,тип події, код події, наприклад
\begin{verbatim}
s="username : 2022-05-05 13:26 : ERROR : 404"
\end{verbatim}
у коді 
\begin{verbatim}
res = re.fullmatch('???', s) 
s = ''
code_str = ??? 
\end{verbatim}
чим треба замінити кожен з ???, щоб значенням code\_str був код події? 


\vspace{10mm}
\textbf{Завдання 17} 

Задано програму, що містить код 1-2 класів, можливе успадкування, та клієнтський код.
У наведеному коді невеликі частини закриті. 
Також наведено вихід програми.
Необхідно відновити закриті частини, щоб отриманий код відповідав наведеному виходу програми.


\textbf{Завдання 18-19} 

Кілька запитань за завданнями для самостійної роботи.
\vspace{10mm}
Скоріш за все, що завдань буде трохи менше, а деякі попередні будуть переформульовані в стилі двох останніх.

%enditem

\end {large}}

%end

\newpage
%begin

{{
\begin {large}

\newpage
1. Що буде надруковано за виконання?

\begin{verbatim}
print('R', end=' ')
a=('a','b','c','d','e',0,1,2,3,4)
print(a[-3])
\end{verbatim}

%enditem


\newpage
2. Що буде надруковано за виконання?

\begin{verbatim}
print('R', end=' ')
c=['0','1','2','3','4','5','6','7','8','9']
print(c[6:-2:3])
\end{verbatim}

%enditem



\newpage
3. Що буде надруковано за виконання?

\begin{verbatim}
print('R', end=' ')
h=[10,11,12,13,14]
print(h.pop(2), end=' ')
print(h)
\end{verbatim}

%enditem



\newpage
4. Що буде надруковано за виконання?

\begin{verbatim}
print('R', end=' ')
a=['0','1','2','3','4']
a[:-5]=(15)
print(a)
\end{verbatim}

%enditem



\newpage
5. Укажіть ТИП значення виразу.\par
\begin{verbatim}
{x for x in range(7) if el%3==1}
\end{verbatim}
%enditem


\newpage
6. Що буде надруковано за виконання?\par
\begin{verbatim}
try:
    d = {83:7, 34:6, 35:3, 35:9}
    d[34] = 6
    for x in d.items():
        print(*x, sep=' ', end=' ')
except:
    print('ERR')
\end{verbatim}
%enditem


\newpage
7. Нехай змінна \kw{s} позначає дуже довгу послідовність цілих, по якій можна ітерувати. 

Написати вираз-генератор, що задає підпослідовність її елементів, що є додатни\-ми. 

Обмеження: усі елементи шуканої підпослідовності одночасно в пам'яті розміщені бути не можуть. 
%enditem


\newpage
8. Що буде надруковано за виконання?
\begin{verbatim}
a = 0
m = 1

class B:
    m = 2
    
    def __init__(self):
        self.a = 3
        m = 4
        B.m = 5

obj = B()
try:
    print(a, end=' ')
    print(m, end=' ')
    print(B.m, end=' ')
    print(obj.a)
except:
    print('ERR')
\end{verbatim}
%enditem


\newpage
9. Що буде надруковано за виконання?
\begin{verbatim}
class D0:
    def a(self): print(1, end=' ')
    def f(self): print(2, end=' ')

class D1(D0):
    def a(self): print(3, end=' ')

try:
    D1().f()
    D1().a()
except:
    print('ERR')
\end{verbatim}
%enditem


\newpage
10. Що буде надруковано за виконання?
\begin{verbatim}
class A:
    c = 1

    def __init__(self):
        b = 2
        c = 3

class B(A):
    c = 4
    a = 5

    def __init__(self):
        super().__init__()
        self.a = 6

try:
    obj = B()
    obj.c = obj.b + 10
    B.c = 13
except: print('EE', end=' ')

try: print(obj.c, end= ' ')
except: print('E1', end=' ')

try: print(A.c, end= ' ')
except: print('E2', end=' ')

try: print(B.c, end= ' ')
except: print('E3', end=' ')
\end{verbatim}
%enditem


\newpage
11. Змінні \kw{next} та \kw{prev} вузла списку позначають наступний та попередній за ним вузол відповідно. Починаючи з голови, у список записано послідовні цілі числа від~$-67$ до~$36$. Нехай змінна \kw{p} позначає вузол, що зберігає значення~$-3$.
Яке значення зберігається у вузлі \kw{p.next.next.next.prev.next} ?
%enditem


\newpage
12. 
\begin{center}
	\adjustimage{max size={0.8\linewidth}{0.8\paperheight}}
	{BinTree/trees_exp/29.png}
\end{center}
Указати послідовність міток вершин дерева, зображеного на рис., що утвориться в результаті обходу дерева в прямому порядку. Мітки записувати через пробіл.\par Алгоритм починає роботу з кореня (на рис. це верхня вершина).
%enditem


\newpage
13. 
\begin{center}
	\adjustimage{max size={0.9\linewidth}{0.9\paperheight}}
	{Graphs/AdjList/145.png}
\end{center}
Для графа, зображеного на рис., вказати список суміжності вершини 7.\par
Формат оформлення відповіді:\par
список суміжності виписувати через пробіл на одному рядку, якщо пробілів десь буде більше одного - не важливо, упорядкування вершин у списку також значення не має.
%enditem


\newpage
14. Для графа на рисунку вказати послідовність, в якій алгоритм пошуку в глибину буде відмічати відвідані вершини, якщо списки суміжності вершин переглядаються алгоритмом за зростанням міток вершин та вершини відмічаються на першому вход\-женні у вершину.\par
Пошук розпочинається з вершини 5.\par

\begin{center}
	\adjustimage{max size={0.9\linewidth}{0.9\paperheight}}
	{Graphs/Traversals/81.png}
\end{center}
%enditem


\newpage
15. Для графа на рисунку з вершини 3 запущено алгоритм пошуку в ширину, що будує кістякове дерево. Списки суміжності вершин переглядаються алгоритмом за зростанням міток вершин.\par
\begin{center}
	\adjustimage{max size={0.9\linewidth}{0.9\paperheight}}
	{Graphs/Traversals/255.png}
\end{center}
\par Які з наведених ребер входять до складу отриманого кістякового дерева?\par
1) (5, 6)\par
2) (0, 4)\par
3) (4, 6)\par
4) (2, 3)\par
5) (2, 4)\par
6) (1, 3)\par
{\vskip 15pt} Номери відповідних ребер записати за зростанням через пробіл.


%enditem


\newpage
16. 

Рядки журналу через двокрапку містять ім'я користувача (ідентифікатор), час події,тип події, код події, наприклад
\begin{verbatim}
s="username : 2022-05-05 13:26 : ERROR : 404"
\end{verbatim}
у коді 
\begin{verbatim}
res = re.fullmatch('???', s) 
s = ''
code_str = ??? 
\end{verbatim}
чим треба замінити кожен з ???, щоб значенням code\_str був код події? 


\vspace{10mm}
\textbf{Завдання 17} 

Задано програму, що містить код 1-2 класів, можливе успадкування, та клієнтський код.
У наведеному коді невеликі частини закриті. 
Також наведено вихід програми.
Необхідно відновити закриті частини, щоб отриманий код відповідав наведеному виходу програми.


\textbf{Завдання 18-19} 

Кілька запитань за завданнями для самостійної роботи.
\vspace{10mm}
Скоріш за все, що завдань буде трохи менше, а деякі попередні будуть переформульовані в стилі двох останніх.

%enditem

\end {large}}

%end

\newpage
%begin

{{
\begin {large}

\newpage
1. Що буде надруковано за виконання?

\begin{verbatim}
print('R', end=' ')
g=('a','b','c','d','e',0,1,2)
print(g[0])
\end{verbatim}

%enditem


\newpage
2. Що буде надруковано за виконання?

\begin{verbatim}
print('R', end=' ')
f='0123456789'
print(f[15::-1])
\end{verbatim}

%enditem



\newpage
3. Що буде надруковано за виконання?

\begin{verbatim}
print('R', end=' ')
a=['0','1','2','3','4']
print(a.index('8'), end=' ')
\end{verbatim}

%enditem



\newpage
4. Що буде надруковано за виконання?

\begin{verbatim}
print('R', end=' ')
e=[10,11,12,13,14]
e[7:]=[2,3]
print(e)
\end{verbatim}

%enditem



\newpage
5. Укажіть ТИП значення виразу.\par
\begin{verbatim}
{x:x for x in range(7)}
\end{verbatim}
%enditem


\newpage
6. Що буде надруковано за виконання?\par
\begin{verbatim}
try:
    t = {78:9, 45:4, 36:6, 88:0}
    t[45] = 1
    for x in t :
        print(x, sep=' ', end=' ')
except:
    print('ERR')
\end{verbatim}
%enditem


\newpage
7. Задано словник \kw{d} з цілими ключами та значеннями. Записати ВИРАЗ, значенням якого є множина, що складається  з тих значень, чиї ключі не діляться на 11.
%enditem


\newpage
8. Що буде надруковано за виконання?
%Оригинал был утерян. Требует отладки!
\begin{verbatim}
c = 0

class A:
    c = 1
    
    def __init__(self):
        self.c = 2
        c = 3
        A.c = 4

obj = A()
A.c = 5
try:
    print(c, end=' ')
    print(obj.c, end=' ')
    print(A.c, end=' ')
    print(obj.c, end=' ')
    print(obj._A__c)
except:
    print('ERR')
\end{verbatim}
%enditem


\newpage
9. Що буде надруковано за виконання?
\begin{verbatim}
class B0:
    def d(self): print(1, end=' ')
    def g(self): print(2, end=' ')

class B1(B0):
    def d(self): print(3, end=' ')

try:
    B0().g()
    B1().d()
except:
    print('ERR')
\end{verbatim}
%enditem


\newpage
10. Що буде надруковано за виконання?
\begin{verbatim}
class A:
    a = 1

    def __init__(self):
        b = 2
        b = 3

class B(A):
    a = 4
    c = 5

    def __init__(self):
        super().__init__()
        self.c = 6

try:
    obj = B()
    obj.a = obj.b + 10
    B.a = 13
except: print('EE', end=' ')

try: print(obj.a, end= ' ')
except: print('E1', end=' ')

try: print(A.a, end= ' ')
except: print('E2', end=' ')

try: print(B.a, end= ' ')
except: print('E3', end=' ')
\end{verbatim}
%enditem


\newpage
11. Змінні \kw{next} та \kw{prev} вузла списку позначають наступний та попередній за ним вузол відповідно. Починаючи з голови, у список записано послідовні цілі числа від~$-15$ до~$46$. Нехай змінна \kw{p} позначає вузол, що зберігає значення~$-6$.
Яке значення зберігається у вузлі \kw{p.prev.prev.prev.next.prev} ?
%enditem


\newpage
12. 
\begin{center}
	\adjustimage{max size={0.8\linewidth}{0.8\paperheight}}
	{BinTree/trees_exp/32.png}
\end{center}
Указати послідовність міток вершин дерева, зображеного на рис., що утвориться в результаті обходу дерева в оберненому порядку. Мітки записувати через пробіл.\par Алгоритм починає роботу з кореня (на рис. це верхня вершина).
%enditem


\newpage
13. 
\begin{center}
	\adjustimage{max size={0.9\linewidth}{0.9\paperheight}}
	{Graphs/AdjList/316.png}
\end{center}
Для графа, зображеного на рис., вказати список суміжності вершини 3.\par
Формат оформлення відповіді:\par
список суміжності виписувати через пробіл на одному рядку, якщо пробілів десь буде більше одного - не важливо, упорядкування вершин у списку також значення не має.
%enditem


\newpage
14. Для графа на рисунку вказати послідовність, в якій алгоритм пошуку в глибину буде відмічати відвідані вершини, якщо списки суміжності вершин переглядаються алгоритмом за зростанням міток вершин та вершини відмічаються на першому вход\-женні у вершину.\par
Пошук розпочинається з вершини 3.\par

\begin{center}
	\adjustimage{max size={0.9\linewidth}{0.9\paperheight}}
	{Graphs/Traversals/162.png}
\end{center}
%enditem


\newpage
15. Для графа на рисунку з вершини 4 запущено алгоритм пошуку в ширину, що будує кістякове дерево. Списки суміжності вершин переглядаються алгоритмом за зростанням міток вершин.\par
\begin{center}
	\adjustimage{max size={0.9\linewidth}{0.9\paperheight}}
	{Graphs/Traversals/342.png}
\end{center}
\par Які з наведених ребер входять до складу отриманого кістякового дерева?\par
1) (0, 1)\par
2) (3, 4)\par
3) (1, 2)\par
4) (5, 6)\par
5) (0, 3)\par
6) (2, 3)\par
{\vskip 15pt} Номери відповідних ребер записати за зростанням через пробіл.


%enditem


\newpage
16. 

Рядки журналу через двокрапку містять ім'я користувача (ідентифікатор), час події,тип події, код події, наприклад
\begin{verbatim}
s="username : 2022-05-05 13:26 : ERROR : 404"
\end{verbatim}
у коді 
\begin{verbatim}
res = re.fullmatch('???', s) 
s = ''
code_str = ??? 
\end{verbatim}
чим треба замінити кожен з ???, щоб значенням code\_str був код події? 


\vspace{10mm}
\textbf{Завдання 17} 

Задано програму, що містить код 1-2 класів, можливе успадкування, та клієнтський код.
У наведеному коді невеликі частини закриті. 
Також наведено вихід програми.
Необхідно відновити закриті частини, щоб отриманий код відповідав наведеному виходу програми.


\textbf{Завдання 18-19} 

Кілька запитань за завданнями для самостійної роботи.
\vspace{10mm}
Скоріш за все, що завдань буде трохи менше, а деякі попередні будуть переформульовані в стилі двох останніх.

%enditem

\end {large}}

%end

\newpage
%begin

{{
\begin {large}

\newpage
1. Що буде надруковано за виконання?

\begin{verbatim}
print('R', end=' ')
e=('a','b','c','d','e',0,1)
print(e[9])
\end{verbatim}

%enditem


\newpage
2. Що буде надруковано за виконання?

\begin{verbatim}
print('R', end=' ')
d=(0,1,2,3,4,5,6,7,8,9)
print(d[:11:2])
\end{verbatim}

%enditem



\newpage
3. Що буде надруковано за виконання?

\begin{verbatim}
print('R', end=' ')
e=[10,11,12,13,14]
e+='10'
print(e)
\end{verbatim}

%enditem



\newpage
4. Що буде надруковано за виконання?

\begin{verbatim}
print('R', end=' ')
c=['0','1','2','3','4']
del c[1:-6]
print(c)
\end{verbatim}

%enditem



\newpage
5. Укажіть ТИП значення виразу.\par
\begin{verbatim}
{y:13 for y in range(-7,7)}
\end{verbatim}
%enditem


\newpage
6. Що буде надруковано за виконання?\par
\begin{verbatim}
try:
    s = {54:0, 47:6, 51:5, 51:8}
    s[49] = 8
    for x in s.items():
        print(x, sep=' ', end=' ')
except:
    print('ERR')
\end{verbatim}
%enditem


\newpage
7. Задано множину \kw{x}, що складається з цілих чисел. Записати ВИРАЗ, значенням якого є множина, що складається  з парних елементів множини \kw{x}.
%enditem


\newpage
8. Що буде надруковано за виконання?
%Оригинал был утерян. Требует отладки!
\begin{verbatim}
_c = 0

class A:
    _c = 1
    
    def __init__(self):
        self._c = 2
        _c = 3
        A._c = 4

obj = A()
obj._c = 5
try:
    print(_c, end=' ')
    print(obj._c, end=' ')
    print(A._c, end=' ')
    print(obj._c, end=' ')
    print(obj._A__c)
except:
    print('ERR')
\end{verbatim}
%enditem


\newpage
9. Що буде надруковано за виконання?
\begin{verbatim}
class A0:
    def b(self): print(1, end=' ')
    def h(self): print(2, end=' ')

class A1(A0):
    def b(self): print(3, end=' ')

try:
    A1().h()
    A1().b()
except:
    print('ERR')
\end{verbatim}
%enditem


\newpage
10. Що буде надруковано за виконання?
\begin{verbatim}
class A:
    b = 1

    def __f1(self): print(2, end=' ')
    def _f2(self): print(3, end=' ')

    def main(self):
        try: print(self.b, end=' ')
        except: print('E1', end=' ')

        try: self.__f1()
        except: print('E2', end=' ')

        try: self._f2()
        except: print('E3', end=' ')


class B(A):
    b = 4

    def __f1(self): print(5, end=' ')
    def _f2(self): print(6, end=' ')

try: B().main()
except: print('E4', end=' ')
\end{verbatim}
%enditem


\newpage
11. Змінні \kw{next} та \kw{prev} вузла списку позначають наступний та попередній за ним вузол відповідно. Починаючи з голови, у список записано послідовні цілі числа від~$-44$ до~$47$. Нехай змінна \kw{p} позначає вузол, що зберігає значення~$7$.
Яке значення зберігається у вузлі \kw{p.prev.prev.next.prev.prev} ?
%enditem


\newpage
12. 
\begin{center}
	\adjustimage{max size={0.8\linewidth}{0.8\paperheight}}
	{BinTree/trees_exp/70.png}
\end{center}
Указати послідовність міток вершин дерева, зображеного на рис., що утвориться в результаті обходу дерева в прямому порядку. Мітки записувати через пробіл.\par Алгоритм починає роботу з кореня (на рис. це верхня вершина).
%enditem


\newpage
13. 
\begin{center}
	\adjustimage{max size={0.9\linewidth}{0.9\paperheight}}
	{Graphs/AdjList/203.png}
\end{center}
Для графа, зображеного на рис., вказати список суміжності вершини 2.\par
Формат оформлення відповіді:\par
список суміжності виписувати через пробіл на одному рядку, якщо пробілів десь буде більше одного - не важливо, упорядкування вершин у списку також значення не має.
%enditem


\newpage
14. Для графа на рисунку вказати послідовність, в якій алгоритм пошуку в глибину буде відмічати відвідані вершини, якщо списки суміжності вершин переглядаються алгоритмом за зростанням міток вершин та вершини відмічаються на першому вход\-женні у вершину.\par
Пошук розпочинається з вершини 1.\par

\begin{center}
	\adjustimage{max size={0.9\linewidth}{0.9\paperheight}}
	{Graphs/Traversals/318.png}
\end{center}
%enditem


\newpage
15. Для графа на рисунку з вершини 6 запущено алгоритм пошуку в ширину, що будує кістякове дерево. Списки суміжності вершин переглядаються алгоритмом за зростанням міток вершин.\par
\begin{center}
	\adjustimage{max size={0.9\linewidth}{0.9\paperheight}}
	{Graphs/Traversals/285.png}
\end{center}
\par Які з наведених ребер входять до складу отриманого кістякового дерева?\par
1) (3, 6)\par
2) (5, 6)\par
3) (4, 6)\par
4) (3, 4)\par
5) (2, 4)\par
6) (0, 1)\par
{\vskip 15pt} Номери відповідних ребер записати за зростанням через пробіл.


%enditem


\newpage
16. 

Рядки журналу через двокрапку містять ім'я користувача (ідентифікатор), час події,тип події, код події, наприклад
\begin{verbatim}
s="username : 2022-05-05 13:26 : ERROR : 404"
\end{verbatim}
у коді 
\begin{verbatim}
res = re.fullmatch('???', s) 
s = ''
code_str = ??? 
\end{verbatim}
чим треба замінити кожен з ???, щоб значенням code\_str був код події? 


\vspace{10mm}
\textbf{Завдання 17} 

Задано програму, що містить код 1-2 класів, можливе успадкування, та клієнтський код.
У наведеному коді невеликі частини закриті. 
Також наведено вихід програми.
Необхідно відновити закриті частини, щоб отриманий код відповідав наведеному виходу програми.


\textbf{Завдання 18-19} 

Кілька запитань за завданнями для самостійної роботи.
\vspace{10mm}
Скоріш за все, що завдань буде трохи менше, а деякі попередні будуть переформульовані в стилі двох останніх.

%enditem

\end {large}}

%end

\newpage
%begin

{{
\begin {large}

\newpage
1. Що буде надруковано за виконання?

\begin{verbatim}
print('R', end=' ')
c=('a','b','c','d','e',0,1,2,4)
print(c[-13])
\end{verbatim}

%enditem


\newpage
2. Що буде надруковано за виконання?

\begin{verbatim}
print('R', end=' ')
b=[0,1,2,3,4,5,6,7,8,9]
print(b[::-2])
\end{verbatim}

%enditem



\newpage
3. Що буде надруковано за виконання?

\begin{verbatim}
print('R', end=' ')
f=['0','1','2','3','4']
f.extend([4,5])
print(f)
\end{verbatim}

%enditem



\newpage
4. Що буде надруковано за виконання?

\begin{verbatim}
print('R', end=' ')
g=[10,11,12,13,14]
g[:-4]='13'
print(g)
\end{verbatim}

%enditem



\newpage
5. Укажіть ТИП значення виразу.\par
\begin{verbatim}
[7,2,3,4]
\end{verbatim}
%enditem


\newpage
6. Що буде надруковано за виконання?\par
\begin{verbatim}
try:
    s = {82:7, 56:6, 84:8, 84:9}
    s[82] = 8
    for x, y in s.items():
        print(x, y, sep=' ', end=' ')
except:
    print('ERR')
\end{verbatim}
%enditem


\newpage
7. 
Записати ВИРАЗ, значенням якого є список, що складається з цілих чисел від 19 до 2011 (включно), що не діляться на 2.
%enditem


\newpage
8. Що буде надруковано за виконання?
\begin{verbatim}
c = 0
n = 1

class C:
    n = 2
    
    def __init__(self):
        self.c = 3
        n = 4
        C.n = 5

obj = C()
try:
    print(c, end=' ')
    print(n, end=' ')
    print(C.n, end=' ')
    print(obj.c)
except:
    print('ERR')
\end{verbatim}
%enditem


\newpage
9. Що буде надруковано за виконання?
\begin{verbatim}
class C0:
    def c(self): print(1, end=' ')
    def g(self): print(2, end=' ')

class C1(C0):
    def c(self): print(3, end=' ')

try:
    C1().g()
    C0().c()
except:
    print('ERR')
\end{verbatim}
%enditem


\newpage
10. Що буде надруковано за виконання?
\begin{verbatim}
class A:
    c = 1

    def __init__(self):
        b = 2
        c = 3

class B(A):
    b = 4
    a = 5

    def __init__(self):
        super().__init__()
        self.a = 6

try:
    obj = B()
    obj.c = obj.c + 10
    B.c = 13
except: print('EE', end=' ')

try: print(obj.c, end= ' ')
except: print('E1', end=' ')

try: print(A.c, end= ' ')
except: print('E2', end=' ')

try: print(B.c, end= ' ')
except: print('E3', end=' ')
\end{verbatim}
%enditem


\newpage
11. Змінні \kw{next} та \kw{prev} вузла списку позначають наступний та попередній за ним вузол відповідно. Починаючи з голови, у список записано послідовні цілі числа від~$-39$ до~$45$. Нехай змінна \kw{p} позначає вузол, що зберігає значення~$5$.
Яке значення зберігається у вузлі \kw{p.next.next.prev.next.next} ?
%enditem


\newpage
12. 
\begin{center}
	\adjustimage{max size={0.8\linewidth}{0.8\paperheight}}
	{BinTree/trees_exp/117.png}
\end{center}
Указати послідовність міток вершин дерева, зображеного на рис., що утвориться в результаті обходу дерева в оберненому порядку. Мітки записувати через пробіл.\par Алгоритм починає роботу з кореня (на рис. це верхня вершина).
%enditem


\newpage
13. 
\begin{center}
	\adjustimage{max size={0.9\linewidth}{0.9\paperheight}}
	{Graphs/AdjList/133.png}
\end{center}
Для графа, зображеного на рис., вказати список суміжності вершини 3.\par
Формат оформлення відповіді:\par
список суміжності виписувати через пробіл на одному рядку, якщо пробілів десь буде більше одного - не важливо, упорядкування вершин у списку також значення не має.
%enditem


\newpage
14. Для графа на рисунку з вершини 6 запущено алгоритм пошуку в глибину, що будує кістякове дерево. Списки суміжності вершин переглядаються алгоритмом за зростанням міток вершин.\par
\begin{center}
	\adjustimage{max size={0.9\linewidth}{0.9\paperheight}}
	{Graphs/Traversals/150.png}
\end{center}
\par Які з наведених ребер входять до складу отриманого кістякового дерева?\par
1) (2, 6)\par
2) (3, 6)\par
3) (0, 1)\par
4) (4, 6)\par
5) (1, 6)\par
6) (0, 5)\par
{\vskip 15pt} Номери відповідних ребер записати за зростанням через пробіл.

%enditem


\newpage
15. Для графа на рисунку з вершини 5 запущено алгоритм пошуку в ширину, що будує кістякове дерево. Списки суміжності вершин переглядаються алгоритмом за зростанням міток вершин.\par
\begin{center}
	\adjustimage{max size={0.9\linewidth}{0.9\paperheight}}
	{Graphs/Traversals/146.png}
\end{center}
\par Які з наведених ребер входять до складу отриманого кістякового дерева?\par
1) (0, 5)\par
2) (0, 3)\par
3) (1, 3)\par
4) (0, 1)\par
5) (1, 6)\par
6) (1, 4)\par
{\vskip 15pt} Номери відповідних ребер записати за зростанням через пробіл.


%enditem


\newpage
16. 

Рядки журналу через двокрапку містять ім'я користувача (ідентифікатор), час події,тип події, код події, наприклад
\begin{verbatim}
s="username : 2022-05-05 13:26 : ERROR : 404"
\end{verbatim}
у коді 
\begin{verbatim}
res = re.fullmatch('???', s) 
s = ''
code_str = ??? 
\end{verbatim}
чим треба замінити кожен з ???, щоб значенням code\_str був код події? 


\vspace{10mm}
\textbf{Завдання 17} 

Задано програму, що містить код 1-2 класів, можливе успадкування, та клієнтський код.
У наведеному коді невеликі частини закриті. 
Також наведено вихід програми.
Необхідно відновити закриті частини, щоб отриманий код відповідав наведеному виходу програми.


\textbf{Завдання 18-19} 

Кілька запитань за завданнями для самостійної роботи.
\vspace{10mm}
Скоріш за все, що завдань буде трохи менше, а деякі попередні будуть переформульовані в стилі двох останніх.

%enditem

\end {large}}

%end

\newpage
%begin

{{
\begin {large}

\newpage
1. Що буде надруковано за виконання?

\begin{verbatim}
print('R', end=' ')
f=('a','b','c','d','e',0,1,2,3)
print(f[2])
\end{verbatim}

%enditem


\newpage
2. Що буде надруковано за виконання?

\begin{verbatim}
print('R', end=' ')
f=('0','1','2','3','4','5','6','7','8','9')
print(f[-2:-6:3])
\end{verbatim}

%enditem



\newpage
3. Що буде надруковано за виконання?

\begin{verbatim}
print('R', end=' ')
g=[10,11,12,13,14]
g.insert(6, 17)
print(g)
\end{verbatim}

%enditem



\newpage
4. Що буде надруковано за виконання?

\begin{verbatim}
print('R', end=' ')
h=[10,11,12,13,14]
h[6:-9]='12'
print(h)
\end{verbatim}

%enditem



\newpage
5. Укажіть ТИП значення виразу.\par
\begin{verbatim}
{(1,4), 8}
\end{verbatim}
%enditem


\newpage
6. Що буде надруковано за виконання?\par
\begin{verbatim}
try:
    t = {87:5, 26:8, 63:7, 43:7, 87:3}
    t[43] = 4
    for x in t.keys():
        print(x, sep=' ', end=' ')
except:
    print('ERR')
\end{verbatim}
%enditem


\newpage
7. Нехай змінна \kw{s} позначає дуже довгу послідовність цілих, по якій можна ітерувати. 

Написати вираз-генератор, що задає підпослідовність її елементів, що менші за задане число num. 

Обмеження: усі елементи шуканої підпослідовності одночасно в пам'яті розміщені бути не можуть. 
%enditem


\newpage
8. Що буде надруковано за виконання?
\begin{verbatim}
b = 0
h = 1

class D:
    b = 2
    
    def __init__(self):
        global b
        self.h = 3
        b = 4
        D.b = 5

obj = D()
try:
    print(b, end=' ')
    print(h, end=' ')
    print(D.b, end=' ')
    print(obj.h)
except:
    print('ERR')
\end{verbatim}
%enditem


\newpage
9. Що буде надруковано за виконання?
\begin{verbatim}
class C0:
    def a(self): print(1, end=' ')
    def h(self): print(2, end=' ')

class C1(C0):
    def a(self): print(3, end=' ')

try:
    C0().h()
    C1().a()
except:
    print('ERR')
\end{verbatim}
%enditem


\newpage
10. Що буде надруковано за виконання?
\begin{verbatim}
class A:
    a = 1

    def __init__(self):
        c = 2
        a = 3

class B(A):
    a = 4
    b = 5

    def __init__(self):
        super().__init__()
        self.b = 6

try:
    obj = B()
    obj.b = obj.a + 10
    B.b = 13
except: print('EE', end=' ')

try: print(obj.b, end= ' ')
except: print('E1', end=' ')

try: print(A.b, end= ' ')
except: print('E2', end=' ')

try: print(B.b, end= ' ')
except: print('E3', end=' ')
\end{verbatim}
%enditem


\newpage
11. Змінні \kw{next} та \kw{prev} вузла списку позначають наступний та попередній за ним вузол відповідно. Починаючи з голови, у список записано послідовні цілі числа від~$-61$ до~$29$. Нехай змінна \kw{p} позначає вузол, що зберігає значення~$-4$.
Яке значення зберігається у вузлі \kw{p.next.next.next.prev.prev} ?
%enditem


\newpage
12. 
\begin{center}
	\adjustimage{max size={0.8\linewidth}{0.8\paperheight}}
	{BinTree/trees_exp/72.png}
\end{center}
Указати послідовність міток вершин дерева, зображеного на рис., що утвориться в результаті обходу дерева в оберненому порядку. Мітки записувати через пробіл.\par Алгоритм починає роботу з кореня (на рис. це верхня вершина).
%enditem


\newpage
13. 
\begin{center}
	\adjustimage{max size={0.9\linewidth}{0.9\paperheight}}
	{Graphs/AdjList/264.png}
\end{center}
Для графа, зображеного на рис., вказати список суміжності вершини 6.\par
Формат оформлення відповіді:\par
список суміжності виписувати через пробіл на одному рядку, якщо пробілів десь буде більше одного - не важливо, упорядкування вершин у списку також значення не має.
%enditem


\newpage
14. Для графа на рисунку вказати послідовність, в якій алгоритм пошуку в глибину буде відмічати відвідані вершини, якщо списки суміжності вершин переглядаються алгоритмом за зростанням міток вершин та вершини відмічаються на першому вход\-женні у вершину.\par
Пошук розпочинається з вершини 5.\par

\begin{center}
	\adjustimage{max size={0.9\linewidth}{0.9\paperheight}}
	{Graphs/Traversals/346.png}
\end{center}
%enditem


\newpage
15. Для графа на рисунку з вершини 1 запущено алгоритм пошуку в ширину, що будує кістякове дерево. Списки суміжності вершин переглядаються алгоритмом за зростанням міток вершин.\par
\begin{center}
	\adjustimage{max size={0.9\linewidth}{0.9\paperheight}}
	{Graphs/Traversals/287.png}
\end{center}
\par Які з наведених ребер входять до складу отриманого кістякового дерева?\par
1) (0, 4)\par
2) (3, 5)\par
3) (3, 4)\par
4) (5, 6)\par
5) (0, 3)\par
6) (1, 2)\par
{\vskip 15pt} Номери відповідних ребер записати за зростанням через пробіл.


%enditem


\newpage
16. 

Рядки журналу через двокрапку містять ім'я користувача (ідентифікатор), час події,тип події, код події, наприклад
\begin{verbatim}
s="username : 2022-05-05 13:26 : ERROR : 404"
\end{verbatim}
у коді 
\begin{verbatim}
res = re.fullmatch('???', s) 
s = ''
code_str = ??? 
\end{verbatim}
чим треба замінити кожен з ???, щоб значенням code\_str був код події? 


\vspace{10mm}
\textbf{Завдання 17} 

Задано програму, що містить код 1-2 класів, можливе успадкування, та клієнтський код.
У наведеному коді невеликі частини закриті. 
Також наведено вихід програми.
Необхідно відновити закриті частини, щоб отриманий код відповідав наведеному виходу програми.


\textbf{Завдання 18-19} 

Кілька запитань за завданнями для самостійної роботи.
\vspace{10mm}
Скоріш за все, що завдань буде трохи менше, а деякі попередні будуть переформульовані в стилі двох останніх.

%enditem

\end {large}}

%end

\newpage
%begin

{{
\begin {large}

\newpage
1. Що буде надруковано за виконання?

\begin{verbatim}
print('R', end=' ')
d=('a','b','c','d','e',0,1)
print(d[-7])
\end{verbatim}

%enditem


\newpage
2. Що буде надруковано за виконання?

\begin{verbatim}
print('R', end=' ')
e='0123456789'
print(e[8:-9:2])
\end{verbatim}

%enditem



\newpage
3. Що буде надруковано за виконання?

\begin{verbatim}
print('R', end=' ')
f=['0','1','2','3','4']
f.remove('5')
print(f)
\end{verbatim}

%enditem



\newpage
4. Що буде надруковано за виконання?

\begin{verbatim}
print('R', end=' ')
f=['0','1','2','3','4']
del f[0:4]
print(f)
\end{verbatim}

%enditem



\newpage
5. Укажіть ТИП значення виразу.\par
\begin{verbatim}
(2, 7, -4)
\end{verbatim}
%enditem


\newpage
6. Що буде надруковано за виконання?\par
\begin{verbatim}
try:
    d = {56:0, 20:2, 42:0, 20:6}
    d[53] = 4
    for x in d.values():
        print(x, sep=' ', end=' ')
except:
    print('ERR')
\end{verbatim}
%enditem


\newpage
7. Нехай змінна \kw{s} позначає дуже довгу послідовність, по якій можна ітерувати. 

Написати вираз-генератор, що задає підпослідовність її елементів, що не є числами типу \kw{float}.

Обмеження: усі елементи шуканої підпослідовності одночасно в пам'яті розміщені бути не можуть. 
%enditem


\newpage
8. Що буде надруковано за виконання?
%Оригинал был утерян. Требует отладки!
\begin{verbatim}
a = 0

class B:
    a = 1
    
    def __init__(self):
        self.a = 2
        a = 3
        B.a = 4

obj = B()
B.a = 5
try:
    print(a, end=' ')
    print(obj.a, end=' ')
    print(B.a, end=' ')
    print(obj.a, end=' ')
    print(obj._B__a)
except:
    print('ERR')
\end{verbatim}
%enditem


\newpage
9. Що буде надруковано за виконання?
\begin{verbatim}
class A0:
    def c(self): print(1, end=' ')
    def f(self): print(2, end=' ')

class A1(A0):
    def c(self): print(3, end=' ')

try:
    A1().f()
    A0().c()
except:
    print('ERR')
\end{verbatim}
%enditem


\newpage
10. Що буде надруковано за виконання?
\begin{verbatim}
class A:
    d = 1

    def _h1(self): print(2, end=' ')
    def __h2(self): print(3, end=' ')

    def main(self):
        try: print(self.d, end=' ')
        except: print('E1', end=' ')

        try: self._h1()
        except: print('E2', end=' ')

        try: self.__h2()
        except: print('E3', end=' ')


class B(A):
    d = 4

    def _h1(self): print(5, end=' ')
    def __h2(self): print(6, end=' ')

try: B().main()
except: print('E4', end=' ')
\end{verbatim}
%enditem


\newpage
11. Змінні \kw{next} та \kw{prev} вузла списку позначають наступний та попередній за ним вузол відповідно. Починаючи з голови, у список записано послідовні цілі числа від~$-29$ до~$32$. Нехай змінна \kw{p} позначає вузол, що зберігає значення~$-8$.
Яке значення зберігається у вузлі \kw{p.prev.prev.prev.next.next} ?
%enditem


\newpage
12. 
\begin{center}
	\adjustimage{max size={0.8\linewidth}{0.8\paperheight}}
	{BinTree/trees_exp/129.png}
\end{center}
Указати послідовність міток вершин дерева, зображеного на рис., що утвориться в результаті обходу дерева в прямому порядку. Мітки записувати через пробіл.\par Алгоритм починає роботу з кореня (на рис. це верхня вершина).
%enditem


\newpage
13. 
\begin{center}
	\adjustimage{max size={0.9\linewidth}{0.9\paperheight}}
	{Graphs/AdjList/16.png}
\end{center}
Для графа, зображеного на рис., вказати список суміжності вершини 5.\par
Формат оформлення відповіді:\par
список суміжності виписувати через пробіл на одному рядку, якщо пробілів десь буде більше одного - не важливо, упорядкування вершин у списку також значення не має.
%enditem


\newpage
14. Для графа на рисунку вказати послідовність, в якій алгоритм пошуку в глибину буде відмічати відвідані вершини, якщо списки суміжності вершин переглядаються алгоритмом за зростанням міток вершин та вершини відмічаються на першому вход\-женні у вершину.\par
Пошук розпочинається з вершини 3.\par

\begin{center}
	\adjustimage{max size={0.9\linewidth}{0.9\paperheight}}
	{Graphs/Traversals/58.png}
\end{center}
%enditem


\newpage
15. Для графа на рисунку з вершини 3 запущено алгоритм пошуку в ширину, що будує кістякове дерево. Списки суміжності вершин переглядаються алгоритмом за зростанням міток вершин.\par
\begin{center}
	\adjustimage{max size={0.9\linewidth}{0.9\paperheight}}
	{Graphs/Traversals/254.png}
\end{center}
\par Які з наведених ребер входять до складу отриманого кістякового дерева?\par
1) (0, 4)\par
2) (1, 4)\par
3) (4, 5)\par
4) (2, 4)\par
5) (0, 6)\par
6) (1, 2)\par
{\vskip 15pt} Номери відповідних ребер записати за зростанням через пробіл.


%enditem


\newpage
16. 

Рядки журналу через двокрапку містять ім'я користувача (ідентифікатор), час події,тип події, код події, наприклад
\begin{verbatim}
s="username : 2022-05-05 13:26 : ERROR : 404"
\end{verbatim}
у коді 
\begin{verbatim}
res = re.fullmatch('???', s) 
s = ''
code_str = ??? 
\end{verbatim}
чим треба замінити кожен з ???, щоб значенням code\_str був код події? 


\vspace{10mm}
\textbf{Завдання 17} 

Задано програму, що містить код 1-2 класів, можливе успадкування, та клієнтський код.
У наведеному коді невеликі частини закриті. 
Також наведено вихід програми.
Необхідно відновити закриті частини, щоб отриманий код відповідав наведеному виходу програми.


\textbf{Завдання 18-19} 

Кілька запитань за завданнями для самостійної роботи.
\vspace{10mm}
Скоріш за все, що завдань буде трохи менше, а деякі попередні будуть переформульовані в стилі двох останніх.

%enditem

\end {large}}

%end

\newpage
%begin

{{
\begin {large}

\newpage
1. Що буде надруковано за виконання?

\begin{verbatim}
print('R', end=' ')
b=('a','b','c','d','e',0,1,2,3,4)
print(b[-10])
\end{verbatim}

%enditem


\newpage
2. Що буде надруковано за виконання?

\begin{verbatim}
print('R', end=' ')
g=[0,1,2,3,4,5,6,7,8,9]
print(g[4:6:-2])
\end{verbatim}

%enditem



\newpage
3. Що буде надруковано за виконання?

\begin{verbatim}
print('R', end=' ')
e=[10,11,12,13,14]
print(e.pop(-4), end=' ')
print(e)
\end{verbatim}

%enditem



\newpage
4. Що буде надруковано за виконання?

\begin{verbatim}
print('R', end=' ')
b=[10,11,12,13,14]
b[-3:2]=[1,0]
print(b)
\end{verbatim}

%enditem



\newpage
5. Укажіть ТИП значення виразу.\par
\begin{verbatim}
[{1,2}, 5]
\end{verbatim}
%enditem


\newpage
6. Що буде надруковано за виконання?\par
\begin{verbatim}
try:
    d = {68:6, 67:5, 30:4, 13:6, 46:1}
    d[68] = 8
    for x in d.values():
        print(x, sep=' ', end=' ')
except:
    print('ERR')
\end{verbatim}
%enditem


\newpage
7. Нехай змінна \kw{s} позначає дуже довгу послідовність цілих, по якій можна ітерувати. 

Написати вираз-генератор, що задає підпослідовність її елементів, що більші за задане число num. 

Обмеження: усі елементи шуканої підпослідовності одночасно в пам'яті розміщені бути не можуть. 
%enditem


\newpage
8. Що буде надруковано за виконання?
\begin{verbatim}
b = 0
n = 1

class A:
    n = 2
    
    def __init__(self):
        self.n = 3
        n = 4
        A.n = 5

obj = A()
try:
    print(b, end=' ')
    print(n, end=' ')
    print(A.n, end=' ')
    print(obj.n)
except:
    print('ERR')
\end{verbatim}
%enditem


\newpage
9. Що буде надруковано за виконання?
\begin{verbatim}
class B0:
    def d(self): print(1, end=' ')
    def f(self): print(2, end=' ')

class B1(B0):
    def f(self): print(3, end=' ')

try:
    B1().d()
    B1().f()
except:
    print('ERR')
\end{verbatim}
%enditem


\newpage
10. Що буде надруковано за виконання?
\begin{verbatim}
class A:
    b = 1

    def __init__(self):
        c = 2
        c = 3

class B(A):
    c = 4
    a = 5

    def __init__(self):
        super().__init__()
        self.a = 6

try:
    obj = B()
    obj.a = obj.b + 10
    B.a = 13
except: print('EE', end=' ')

try: print(obj.a, end= ' ')
except: print('E1', end=' ')

try: print(A.a, end= ' ')
except: print('E2', end=' ')

try: print(B.a, end= ' ')
except: print('E3', end=' ')
\end{verbatim}
%enditem


\newpage
11. Змінні \kw{next} та \kw{prev} вузла списку позначають наступний та попередній за ним вузол відповідно. Починаючи з голови, у список записано послідовні цілі числа від~$-66$ до~$44$. Нехай змінна \kw{p} позначає вузол, що зберігає значення~$1$.
Яке значення зберігається у вузлі \kw{p.prev.next.prev.next.next} ?
%enditem


\newpage
12. 
\begin{center}
	\adjustimage{max size={0.8\linewidth}{0.8\paperheight}}
	{BinTree/trees_exp/64.png}
\end{center}
Указати послідовність міток вершин дерева, зображеного на рис., що утвориться в результаті обходу дерева в прямому порядку. Мітки записувати через пробіл.\par Алгоритм починає роботу з кореня (на рис. це верхня вершина).
%enditem


\newpage
13. 
\begin{center}
	\adjustimage{max size={0.9\linewidth}{0.9\paperheight}}
	{Graphs/AdjList/64.png}
\end{center}
Для графа, зображеного на рис., вказати список суміжності вершини 1.\par
Формат оформлення відповіді:\par
список суміжності виписувати через пробіл на одному рядку, якщо пробілів десь буде більше одного - не важливо, упорядкування вершин у списку також значення не має.
%enditem


\newpage
14. Для графа на рисунку з вершини 3 запущено алгоритм пошуку в глибину, що будує кістякове дерево. Списки суміжності вершин переглядаються алгоритмом за зростанням міток вершин.\par
\begin{center}
	\adjustimage{max size={0.9\linewidth}{0.9\paperheight}}
	{Graphs/Traversals/332.png}
\end{center}
\par Які з наведених ребер входять до складу отриманого кістякового дерева?\par
1) (0, 1)\par
2) (2, 3)\par
3) (0, 4)\par
4) (5, 6)\par
5) (2, 5)\par
6) (0, 6)\par
{\vskip 15pt} Номери відповідних ребер записати за зростанням через пробіл.

%enditem


\newpage
15. Для графа на рисунку з вершини 1 запущено алгоритм пошуку в ширину, що будує кістякове дерево. Списки суміжності вершин переглядаються алгоритмом за зростанням міток вершин.\par
\begin{center}
	\adjustimage{max size={0.9\linewidth}{0.9\paperheight}}
	{Graphs/Traversals/19.png}
\end{center}
\par Які з наведених ребер входять до складу отриманого кістякового дерева?\par
1) (1, 5)\par
2) (2, 5)\par
3) (3, 4)\par
4) (1, 6)\par
5) (0, 6)\par
6) (2, 6)\par
{\vskip 15pt} Номери відповідних ребер записати за зростанням через пробіл.


%enditem


\newpage
16. 

Рядки журналу через двокрапку містять ім'я користувача (ідентифікатор), час події,тип події, код події, наприклад
\begin{verbatim}
s="username : 2022-05-05 13:26 : ERROR : 404"
\end{verbatim}
у коді 
\begin{verbatim}
res = re.fullmatch('???', s) 
s = ''
code_str = ??? 
\end{verbatim}
чим треба замінити кожен з ???, щоб значенням code\_str був код події? 


\vspace{10mm}
\textbf{Завдання 17} 

Задано програму, що містить код 1-2 класів, можливе успадкування, та клієнтський код.
У наведеному коді невеликі частини закриті. 
Також наведено вихід програми.
Необхідно відновити закриті частини, щоб отриманий код відповідав наведеному виходу програми.


\textbf{Завдання 18-19} 

Кілька запитань за завданнями для самостійної роботи.
\vspace{10mm}
Скоріш за все, що завдань буде трохи менше, а деякі попередні будуть переформульовані в стилі двох останніх.

%enditem

\end {large}}

%end

\newpage
%begin

{{
\begin {large}

\newpage
1. Що буде надруковано за виконання?

\begin{verbatim}
print('R', end=' ')
e=('a','b','c','d','e',0,1,2,4)
print(e[-7])
\end{verbatim}

%enditem


\newpage
2. Що буде надруковано за виконання?

\begin{verbatim}
print('R', end=' ')
b=(0,1,2,3,4,5,6,7,8,9)
print(b[::-3])
\end{verbatim}

%enditem



\newpage
3. Що буде надруковано за виконання?

\begin{verbatim}
print('R', end=' ')
b=['0','1','2','3','4']
b+='12'
print(b)
\end{verbatim}

%enditem



\newpage
4. Що буде надруковано за виконання?

\begin{verbatim}
print('R', end=' ')
g=['0','1','2','3','4']
g[2:]='10'
print(g)
\end{verbatim}

%enditem



\newpage
5. Укажіть ТИП значення виразу.\par
\begin{verbatim}
{1:3, 5:8}
\end{verbatim}
%enditem


\newpage
6. Що буде надруковано за виконання?\par
\begin{verbatim}
try:
    s = {58:5, 60:3, 72:5, 72:4}
    s[11] = 7
    for x in s.items():
        print(x, sep=' ', end=' ')
except:
    print('ERR')
\end{verbatim}
%enditem


\newpage
7. Нехай змінна \kw{s} позначає послідовність цілих, по якій можна ітерувати. 

Написати вираз, що перевіряє, чи правда, що послідовність  містить число 13.
%enditem


\newpage
8. Що буде надруковано за виконання?
\begin{verbatim}
d = 0
m = 1

class D:
    d = 2
    
    def __init__(self):
        global d
        self.d = 3
        d = 4
        D.d = 5

obj = D()
try:
    print(d, end=' ')
    print(m, end=' ')
    print(D.d, end=' ')
    print(obj.d)
except:
    print('ERR')
\end{verbatim}
%enditem


\newpage
9. Що буде надруковано за виконання?
\begin{verbatim}
class D0:
    def b(self): print(1, end=' ')
    def g(self): print(2, end=' ')

class D1(D0):
    def b(self): print(3, end=' ')

try:
    D1().g()
    D1().b()
except:
    print('ERR')
\end{verbatim}
%enditem


\newpage
10. Що буде надруковано за виконання?
\begin{verbatim}
class A:
    a = 1

    def _f1(self): print(2, end=' ')
    def f2(self): print(3, end=' ')

    def main(self):
        try: print(self.a, end=' ')
        except: print('E1', end=' ')

        try: self._f1()
        except: print('E2', end=' ')

        try: self.f2()
        except: print('E3', end=' ')


class B(A):
    a = 4

    def _f1(self): print(5, end=' ')
    def f2(self): print(6, end=' ')

try: B().main()
except: print('E4', end=' ')
\end{verbatim}
%enditem


\newpage
11. Змінні \kw{next} та \kw{prev} вузла списку позначають наступний та попередній за ним вузол відповідно. Починаючи з голови, у список записано послідовні цілі числа від~$-36$ до~$17$. Нехай змінна \kw{p} позначає вузол, що зберігає значення~$8$.
Яке значення зберігається у вузлі \kw{p.next.prev.next.prev.prev} ?
%enditem


\newpage
12. 
\begin{center}
	\adjustimage{max size={0.8\linewidth}{0.8\paperheight}}
	{BinTree/trees_exp/105.png}
\end{center}
Указати послідовність міток вершин дерева, зображеного на рис., що утвориться в результаті обходу дерева в оберненому порядку. Мітки записувати через пробіл.\par Алгоритм починає роботу з кореня (на рис. це верхня вершина).
%enditem


\newpage
13. 
\begin{center}
	\adjustimage{max size={0.9\linewidth}{0.9\paperheight}}
	{Graphs/AdjList/311.png}
\end{center}
Для графа, зображеного на рис., вказати список суміжності вершини 1.\par
Формат оформлення відповіді:\par
список суміжності виписувати через пробіл на одному рядку, якщо пробілів десь буде більше одного - не важливо, упорядкування вершин у списку також значення не має.
%enditem


\newpage
14. Для графа на рисунку вказати послідовність, в якій алгоритм пошуку в глибину буде відмічати відвідані вершини, якщо списки суміжності вершин переглядаються алгоритмом за зростанням міток вершин та вершини відмічаються на першому вход\-женні у вершину.\par
Пошук розпочинається з вершини 4.\par

\begin{center}
	\adjustimage{max size={0.9\linewidth}{0.9\paperheight}}
	{Graphs/Traversals/245.png}
\end{center}
%enditem


\newpage
15. Для графа на рисунку з вершини 1 запущено алгоритм пошуку в ширину, що будує кістякове дерево. Списки суміжності вершин переглядаються алгоритмом за зростанням міток вершин.\par
\begin{center}
	\adjustimage{max size={0.9\linewidth}{0.9\paperheight}}
	{Graphs/Traversals/318.png}
\end{center}
\par Які з наведених ребер входять до складу отриманого кістякового дерева?\par
1) (3, 4)\par
2) (2, 6)\par
3) (1, 6)\par
4) (0, 2)\par
5) (1, 3)\par
6) (2, 4)\par
{\vskip 15pt} Номери відповідних ребер записати за зростанням через пробіл.


%enditem


\newpage
16. 

Рядки журналу через двокрапку містять ім'я користувача (ідентифікатор), час події,тип події, код події, наприклад
\begin{verbatim}
s="username : 2022-05-05 13:26 : ERROR : 404"
\end{verbatim}
у коді 
\begin{verbatim}
res = re.fullmatch('???', s) 
s = ''
code_str = ??? 
\end{verbatim}
чим треба замінити кожен з ???, щоб значенням code\_str був код події? 


\vspace{10mm}
\textbf{Завдання 17} 

Задано програму, що містить код 1-2 класів, можливе успадкування, та клієнтський код.
У наведеному коді невеликі частини закриті. 
Також наведено вихід програми.
Необхідно відновити закриті частини, щоб отриманий код відповідав наведеному виходу програми.


\textbf{Завдання 18-19} 

Кілька запитань за завданнями для самостійної роботи.
\vspace{10mm}
Скоріш за все, що завдань буде трохи менше, а деякі попередні будуть переформульовані в стилі двох останніх.

%enditem

\end {large}}

%end

\newpage
%begin

{{
\begin {large}

\newpage
1. Що буде надруковано за виконання?

\begin{verbatim}
print('R', end=' ')
a=('a','b','c','d','e',0,1,2)
print(a[-8])
\end{verbatim}

%enditem


\newpage
2. Що буде надруковано за виконання?

\begin{verbatim}
print('R', end=' ')
h=['0','1','2','3','4','5','6','7','8','9']
print(h[:-1:-1])
\end{verbatim}

%enditem



\newpage
3. Що буде надруковано за виконання?

\begin{verbatim}
print('R', end=' ')
d=[10,11,12,13,14]
print(d.index(8), end=' ')
\end{verbatim}

%enditem



\newpage
4. Що буде надруковано за виконання?

\begin{verbatim}
print('R', end=' ')
c=['0','1','2','3','4']
del c[3:3]
print(c)
\end{verbatim}

%enditem



\newpage
5. Укажіть ТИП значення виразу.\par
\begin{verbatim}
(11,2,3)
\end{verbatim}
%enditem


\newpage
6. Що буде надруковано за виконання?\par
\begin{verbatim}
try:
    t = {38:6, 26:2, 74:4, 13:3, 13:7}
    t[26] = 4
    for x in t.keys():
        print(x, sep=' ', end=' ')
except:
    print('ERR')
\end{verbatim}
%enditem


\newpage
7. Нехай змінна \kw{s} позначає послідовність цілих, по якій можна ітерувати. 

Написати вираз, що перевіряє, чи правда, що послідовність  містить число 35.
%enditem


\newpage
8. Що буде надруковано за виконання?
%Оригинал был утерян. Требует отладки!
\begin{verbatim}
__b = 0

class D:
    __b = 1
    
    def __init__(self):
        self.__b = 2
        __b = 3
        D.__b = 4

obj = D()
D.__b = 5
try:
    print(__b, end=' ')
    print(obj.__b, end=' ')
    print(D.__b, end=' ')
    print(obj.__b, end=' ')
    print(obj._D__b)
except:
    print('ERR')
\end{verbatim}
%enditem


\newpage
9. Що буде надруковано за виконання?
\begin{verbatim}
class C0:
    def b(self): print(1, end=' ')
    def h(self): print(2, end=' ')

class C1(C0):
    def b(self): print(3, end=' ')

try:
    C1().h()
    C1().b()
except:
    print('ERR')
\end{verbatim}
%enditem


\newpage
10. Що буде надруковано за виконання?
\begin{verbatim}
class A:
    c = 1

    def g1(self): print(2, end=' ')
    def _g2(self): print(3, end=' ')

    def main(self):
        try: print(self.c, end=' ')
        except: print('E1', end=' ')

        try: self.g1()
        except: print('E2', end=' ')

        try: self._g2()
        except: print('E3', end=' ')


class B(A):
    c = 4

    def g1(self): print(5, end=' ')
    def _g2(self): print(6, end=' ')

try: B().main()
except: print('E4', end=' ')
\end{verbatim}
%enditem


\newpage
11. Змінні \kw{next} та \kw{prev} вузла списку позначають наступний та попередній за ним вузол відповідно. Починаючи з голови, у список записано послідовні цілі числа від~$-34$ до~$48$. Нехай змінна \kw{p} позначає вузол, що зберігає значення~$-5$.
Яке значення зберігається у вузлі \kw{p.next.next.prev.prev.prev} ?
%enditem


\newpage
12. 
\begin{center}
	\adjustimage{max size={0.8\linewidth}{0.8\paperheight}}
	{BinTree/trees_exp/176.png}
\end{center}
Указати послідовність міток вершин дерева, зображеного на рис., що утвориться в результаті обходу дерева в прямому порядку. Мітки записувати через пробіл.\par Алгоритм починає роботу з кореня (на рис. це верхня вершина).
%enditem


\newpage
13. 
\begin{center}
	\adjustimage{max size={0.9\linewidth}{0.9\paperheight}}
	{Graphs/AdjList/123.png}
\end{center}
Для графа, зображеного на рис., вказати список суміжності вершини 9.\par
Формат оформлення відповіді:\par
список суміжності виписувати через пробіл на одному рядку, якщо пробілів десь буде більше одного - не важливо, упорядкування вершин у списку також значення не має.
%enditem


\newpage
14. Для графа на рисунку вказати послідовність, в якій алгоритм пошуку в глибину буде відмічати відвідані вершини, якщо списки суміжності вершин переглядаються алгоритмом за зростанням міток вершин та вершини відмічаються на першому вход\-женні у вершину.\par
Пошук розпочинається з вершини 4.\par

\begin{center}
	\adjustimage{max size={0.9\linewidth}{0.9\paperheight}}
	{Graphs/Traversals/8.png}
\end{center}
%enditem


\newpage
15. Для графа на рисунку з вершини 5 запущено алгоритм пошуку в ширину, що будує кістякове дерево. Списки суміжності вершин переглядаються алгоритмом за зростанням міток вершин.\par
\begin{center}
	\adjustimage{max size={0.9\linewidth}{0.9\paperheight}}
	{Graphs/Traversals/224.png}
\end{center}
\par Які з наведених ребер входять до складу отриманого кістякового дерева?\par
1) (5, 6)\par
2) (2, 6)\par
3) (0, 5)\par
4) (3, 5)\par
5) (2, 4)\par
6) (0, 2)\par
{\vskip 15pt} Номери відповідних ребер записати за зростанням через пробіл.


%enditem


\newpage
16. 

Рядки журналу через двокрапку містять ім'я користувача (ідентифікатор), час події,тип події, код події, наприклад
\begin{verbatim}
s="username : 2022-05-05 13:26 : ERROR : 404"
\end{verbatim}
у коді 
\begin{verbatim}
res = re.fullmatch('???', s) 
s = ''
code_str = ??? 
\end{verbatim}
чим треба замінити кожен з ???, щоб значенням code\_str був код події? 


\vspace{10mm}
\textbf{Завдання 17} 

Задано програму, що містить код 1-2 класів, можливе успадкування, та клієнтський код.
У наведеному коді невеликі частини закриті. 
Також наведено вихід програми.
Необхідно відновити закриті частини, щоб отриманий код відповідав наведеному виходу програми.


\textbf{Завдання 18-19} 

Кілька запитань за завданнями для самостійної роботи.
\vspace{10mm}
Скоріш за все, що завдань буде трохи менше, а деякі попередні будуть переформульовані в стилі двох останніх.

%enditem

\end {large}}

%end

\newpage
%begin

{{
\begin {large}

\newpage
1. Що буде надруковано за виконання?

\begin{verbatim}
print('R', end=' ')
c=('a','b','c','d','e',0,1,2)
print(c[6])
\end{verbatim}

%enditem


\newpage
2. Що буде надруковано за виконання?

\begin{verbatim}
print('R', end=' ')
d=[0,1,2,3,4,5,6,7,8,9]
print(d[6::3])
\end{verbatim}

%enditem



\newpage
3. Що буде надруковано за виконання?

\begin{verbatim}
print('R', end=' ')
h=['0','1','2','3','4']
h.remove(0)
print(h)
\end{verbatim}

%enditem



\newpage
4. Що буде надруковано за виконання?

\begin{verbatim}
print('R', end=' ')
d=[10,11,12,13,14]
del d[5:]
print(d)
\end{verbatim}

%enditem



\newpage
5. Укажіть ТИП значення виразу.\par
\begin{verbatim}
{y:None for y in range(-8,3)}
\end{verbatim}
%enditem


\newpage
6. Що буде надруковано за виконання?\par
\begin{verbatim}
try:
    s = {32:5, 70:2, 67:7, 32:3}
    s[83] = 7
    for x in s :
        print(x, sep=' ', end=' ')
except:
    print('ERR')
\end{verbatim}
%enditem


\newpage
7. Нехай змінна \kw{s} позначає дуже довгу послідовність, по якій можна ітерувати. 

Написати вираз-генератор, що задає підпослідовність її елементів, що є числами типу \kw{int}.

Обмеження: усі елементи шуканої підпослідовності одночасно в пам'яті розміщені бути не можуть. 
%enditem


\newpage
8. Що буде надруковано за виконання?
%Оригинал был утерян. Требует отладки!
\begin{verbatim}
_d = 0

class C:
    _d = 1
    
    def __init__(self):
        self._d = 2
        _d = 3
        C._d = 4

obj = C()
obj._d = 5
try:
    print(_d, end=' ')
    print(obj._d, end=' ')
    print(C._d, end=' ')
    print(obj._d, end=' ')
    print(obj._C__d)
except:
    print('ERR')
\end{verbatim}
%enditem


\newpage
9. Що буде надруковано за виконання?
\begin{verbatim}
class D0:
    def a(self): print(1, end=' ')
    def f(self): print(2, end=' ')

class D1(D0):
    def f(self): print(3, end=' ')

try:
    D0().a()
    D1().f()
except:
    print('ERR')
\end{verbatim}
%enditem


\newpage
10. Що буде надруковано за виконання?
\begin{verbatim}
class A:
    b = 1

    def _h1(self): print(2, end=' ')
    def _h2(self): print(3, end=' ')

    def main(self):
        try: print(self.b, end=' ')
        except: print('E1', end=' ')

        try: self._h1()
        except: print('E2', end=' ')

        try: self._h2()
        except: print('E3', end=' ')


class B(A):
    b = 4

    def _h1(self): print(5, end=' ')
    def _h2(self): print(6, end=' ')

try: B().main()
except: print('E4', end=' ')
\end{verbatim}
%enditem


\newpage
11. Змінні \kw{next} та \kw{prev} вузла списку позначають наступний та попередній за ним вузол відповідно. Починаючи з голови, у список записано послідовні цілі числа від~$-45$ до~$21$. Нехай змінна \kw{p} позначає вузол, що зберігає значення~$6$.
Яке значення зберігається у вузлі \kw{p.prev.prev.next.next.next} ?
%enditem


\newpage
12. 
\begin{center}
	\adjustimage{max size={0.8\linewidth}{0.8\paperheight}}
	{BinTree/trees_exp/166.png}
\end{center}
Указати послідовність міток вершин дерева, зображеного на рис., що утвориться в результаті обходу дерева в оберненому порядку. Мітки записувати через пробіл.\par Алгоритм починає роботу з кореня (на рис. це верхня вершина).
%enditem


\newpage
13. 
\begin{center}
	\adjustimage{max size={0.9\linewidth}{0.9\paperheight}}
	{Graphs/AdjList/170.png}
\end{center}
Для графа, зображеного на рис., вказати список суміжності вершини 0.\par
Формат оформлення відповіді:\par
список суміжності виписувати через пробіл на одному рядку, якщо пробілів десь буде більше одного - не важливо, упорядкування вершин у списку також значення не має.
%enditem


\newpage
14. Для графа на рисунку вказати послідовність, в якій алгоритм пошуку в глибину буде відмічати відвідані вершини, якщо списки суміжності вершин переглядаються алгоритмом за зростанням міток вершин та вершини відмічаються на першому вход\-женні у вершину.\par
Пошук розпочинається з вершини 1.\par

\begin{center}
	\adjustimage{max size={0.9\linewidth}{0.9\paperheight}}
	{Graphs/Traversals/221.png}
\end{center}
%enditem


\newpage
15. Для графа на рисунку з вершини 0 запущено алгоритм пошуку в ширину, що будує кістякове дерево. Списки суміжності вершин переглядаються алгоритмом за зростанням міток вершин.\par
\begin{center}
	\adjustimage{max size={0.9\linewidth}{0.9\paperheight}}
	{Graphs/Traversals/49.png}
\end{center}
\par Які з наведених ребер входять до складу отриманого кістякового дерева?\par
1) (0, 5)\par
2) (3, 5)\par
3) (1, 2)\par
4) (1, 3)\par
5) (1, 6)\par
6) (2, 6)\par
{\vskip 15pt} Номери відповідних ребер записати за зростанням через пробіл.


%enditem


\newpage
16. 

Рядки журналу через двокрапку містять ім'я користувача (ідентифікатор), час події,тип події, код події, наприклад
\begin{verbatim}
s="username : 2022-05-05 13:26 : ERROR : 404"
\end{verbatim}
у коді 
\begin{verbatim}
res = re.fullmatch('???', s) 
s = ''
code_str = ??? 
\end{verbatim}
чим треба замінити кожен з ???, щоб значенням code\_str був код події? 


\vspace{10mm}
\textbf{Завдання 17} 

Задано програму, що містить код 1-2 класів, можливе успадкування, та клієнтський код.
У наведеному коді невеликі частини закриті. 
Також наведено вихід програми.
Необхідно відновити закриті частини, щоб отриманий код відповідав наведеному виходу програми.


\textbf{Завдання 18-19} 

Кілька запитань за завданнями для самостійної роботи.
\vspace{10mm}
Скоріш за все, що завдань буде трохи менше, а деякі попередні будуть переформульовані в стилі двох останніх.

%enditem

\end {large}}

%end

\newpage
%begin

{{
\begin {large}

\newpage
1. Що буде надруковано за виконання?

\begin{verbatim}
print('R', end=' ')
g=('a','b','c','d','e',0,1,2,3,4)
print(g[-9])
\end{verbatim}

%enditem


\newpage
2. Що буде надруковано за виконання?

\begin{verbatim}
print('R', end=' ')
a='0123456789'
print(a[-3::2])
\end{verbatim}

%enditem



\newpage
3. Що буде надруковано за виконання?

\begin{verbatim}
print('R', end=' ')
a=[10,11,12,13,14]
a.append('11')
print(a)
\end{verbatim}

%enditem



\newpage
4. Що буде надруковано за виконання?

\begin{verbatim}
print('R', end=' ')
a=['0','1','2','3','4']
a[5:-3]=[1,2]
print(a)
\end{verbatim}

%enditem



\newpage
5. Укажіть ТИП значення виразу.\par
\begin{verbatim}
{x*x:x for x in range(-7,7)}
\end{verbatim}
%enditem


\newpage
6. Що буде надруковано за виконання?\par
\begin{verbatim}
try:
    d = {10:9, 41:2, 58:8, 41:8}
    d[58] = 1
    for x, y in d.items():
        print(x, y, sep=' ', end=' ')
except:
    print('ERR')
\end{verbatim}
%enditem


\newpage
7. Нехай змінна \kw{s} позначає дуже довгу послідовність цілих, по якій можна ітерувати. 

Написати вираз-генератор, що задає підпослідовність її елементів, що діляться на 3. 

Обмеження: усі елементи шуканої підпослідовності одночасно в пам'яті розміщені бути не можуть. 
%enditem


\newpage
8. Що буде надруковано за виконання?
\begin{verbatim}
a = 0
k = 1

class B:
    a = 2
    
    def __init__(self):
        global a
        self.k = 3
        a = 4
        B.a = 5

obj = B()
try:
    print(a, end=' ')
    print(k, end=' ')
    print(B.a, end=' ')
    print(obj.k)
except:
    print('ERR')
\end{verbatim}
%enditem


\newpage
9. Що буде надруковано за виконання?
\begin{verbatim}
class A0:
    def c(self): print(1, end=' ')
    def h(self): print(2, end=' ')

class A1(A0):
    def h(self): print(3, end=' ')

try:
    A1().c()
    A1().h()
except:
    print('ERR')
\end{verbatim}
%enditem


\newpage
10. Що буде надруковано за виконання?
\begin{verbatim}
class A:
    d = 1

    def __f1(self): print(2, end=' ')
    def _f2(self): print(3, end=' ')

    def main(self):
        try: print(self.d, end=' ')
        except: print('E1', end=' ')

        try: self.__f1()
        except: print('E2', end=' ')

        try: self._f2()
        except: print('E3', end=' ')


class B(A):
    d = 4

    def __f1(self): print(5, end=' ')
    def _f2(self): print(6, end=' ')

try: B().main()
except: print('E4', end=' ')
\end{verbatim}
%enditem


\newpage
11. Змінні \kw{next} та \kw{prev} вузла списку позначають наступний та попередній за ним вузол відповідно. Починаючи з голови, у список записано послідовні цілі числа від~$-68$ до~$20$. Нехай змінна \kw{p} позначає вузол, що зберігає значення~$-6$.
Яке значення зберігається у вузлі \kw{p.next.next.next.next.prev} ?
%enditem


\newpage
12. 
\begin{center}
	\adjustimage{max size={0.8\linewidth}{0.8\paperheight}}
	{BinTree/trees_exp/85.png}
\end{center}
Указати послідовність міток вершин дерева, зображеного на рис., що утвориться в результаті обходу дерева в оберненому порядку. Мітки записувати через пробіл.\par Алгоритм починає роботу з кореня (на рис. це верхня вершина).
%enditem


\newpage
13. 
\begin{center}
	\adjustimage{max size={0.9\linewidth}{0.9\paperheight}}
	{Graphs/AdjList/239.png}
\end{center}
Для графа, зображеного на рис., вказати список суміжності вершини 1.\par
Формат оформлення відповіді:\par
список суміжності виписувати через пробіл на одному рядку, якщо пробілів десь буде більше одного - не важливо, упорядкування вершин у списку також значення не має.
%enditem


\newpage
14. Для графа на рисунку з вершини 3 запущено алгоритм пошуку в глибину, що будує кістякове дерево. Списки суміжності вершин переглядаються алгоритмом за зростанням міток вершин.\par
\begin{center}
	\adjustimage{max size={0.9\linewidth}{0.9\paperheight}}
	{Graphs/Traversals/270.png}
\end{center}
\par Які з наведених ребер входять до складу отриманого кістякового дерева?\par
1) (3, 5)\par
2) (1, 4)\par
3) (1, 5)\par
4) (2, 3)\par
5) (2, 4)\par
6) (0, 1)\par
{\vskip 15pt} Номери відповідних ребер записати за зростанням через пробіл.

%enditem


\newpage
15. Для графа на рисунку з вершини 3 запущено алгоритм пошуку в ширину, що будує кістякове дерево. Списки суміжності вершин переглядаються алгоритмом за зростанням міток вершин.\par
\begin{center}
	\adjustimage{max size={0.9\linewidth}{0.9\paperheight}}
	{Graphs/Traversals/3.png}
\end{center}
\par Які з наведених ребер входять до складу отриманого кістякового дерева?\par
1) (0, 6)\par
2) (3, 6)\par
3) (5, 6)\par
4) (3, 4)\par
5) (1, 2)\par
6) (4, 5)\par
{\vskip 15pt} Номери відповідних ребер записати за зростанням через пробіл.


%enditem


\newpage
16. 

Рядки журналу через двокрапку містять ім'я користувача (ідентифікатор), час події,тип події, код події, наприклад
\begin{verbatim}
s="username : 2022-05-05 13:26 : ERROR : 404"
\end{verbatim}
у коді 
\begin{verbatim}
res = re.fullmatch('???', s) 
s = ''
code_str = ??? 
\end{verbatim}
чим треба замінити кожен з ???, щоб значенням code\_str був код події? 


\vspace{10mm}
\textbf{Завдання 17} 

Задано програму, що містить код 1-2 класів, можливе успадкування, та клієнтський код.
У наведеному коді невеликі частини закриті. 
Також наведено вихід програми.
Необхідно відновити закриті частини, щоб отриманий код відповідав наведеному виходу програми.


\textbf{Завдання 18-19} 

Кілька запитань за завданнями для самостійної роботи.
\vspace{10mm}
Скоріш за все, що завдань буде трохи менше, а деякі попередні будуть переформульовані в стилі двох останніх.

%enditem

\end {large}}

%end

\newpage
%begin

{{
\begin {large}

\newpage
1. Що буде надруковано за виконання?

\begin{verbatim}
print('R', end=' ')
h=('a','b','c','d','e',0,1)
print(h[11])
\end{verbatim}

%enditem


\newpage
2. Що буде надруковано за виконання?

\begin{verbatim}
print('R', end=' ')
c=['0','1','2','3','4','5','6','7','8','9']
print(c[-4:14:-1])
\end{verbatim}

%enditem



\newpage
3. Що буде надруковано за виконання?

\begin{verbatim}
print('R', end=' ')
c=['0','1','2','3','4']
c.extend([1,2])
print(c)
\end{verbatim}

%enditem



\newpage
4. Що буде надруковано за виконання?

\begin{verbatim}
print('R', end=' ')
e=[10,11,12,13,14]
e[-1:-3]=[4,5]
print(e)
\end{verbatim}

%enditem



\newpage
5. Укажіть ТИП значення виразу.\par
\begin{verbatim}
[13,2,4]+[3,5]
\end{verbatim}
%enditem


\newpage
6. Що буде надруковано за виконання?\par
\begin{verbatim}
try:
    t = {58:9, 12:8, 57:5, 12:7}
    t[62] = 7
    for x in t.items():
        print(*x, sep=' ', end=' ')
except:
    print('ERR')
\end{verbatim}
%enditem


\newpage
7. Нехай змінна \kw{s} позначає послідовність цілих, по якій можна ітерувати. 

Написати вираз, що перевіряє, чи правда, що послідовність  містить від’ємні числа.
%enditem


\newpage
8. Що буде надруковано за виконання?
\begin{verbatim}
c = 0
j = 1

class C:
    j = 2
    
    def __init__(self):
        self.c = 3
        j = 4
        C.j = 5

obj = C()
try:
    print(c, end=' ')
    print(j, end=' ')
    print(C.j, end=' ')
    print(obj.c)
except:
    print('ERR')
\end{verbatim}
%enditem


\newpage
9. Що буде надруковано за виконання?
\begin{verbatim}
class B0:
    def d(self): print(1, end=' ')
    def g(self): print(2, end=' ')

class B1(B0):
    def d(self): print(3, end=' ')

try:
    B1().g()
    B0().d()
except:
    print('ERR')
\end{verbatim}
%enditem


\newpage
10. Що буде надруковано за виконання?
\begin{verbatim}
class A:
    b = 1

    def __init__(self):
        a = 2
        a = 3

class B(A):
    a = 4
    c = 5

    def __init__(self):
        super().__init__()
        self.c = 6

try:
    obj = B()
    obj.c = obj.c + 10
    B.c = 13
except: print('EE', end=' ')

try: print(obj.c, end= ' ')
except: print('E1', end=' ')

try: print(A.c, end= ' ')
except: print('E2', end=' ')

try: print(B.c, end= ' ')
except: print('E3', end=' ')
\end{verbatim}
%enditem


\newpage
11. Змінні \kw{next} та \kw{prev} вузла списку позначають наступний та попередній за ним вузол відповідно. Починаючи з голови, у список записано послідовні цілі числа від~$-51$ до~$19$. Нехай змінна \kw{p} позначає вузол, що зберігає значення~$3$.
Яке значення зберігається у вузлі \kw{p.prev.prev.prev.prev.next} ?
%enditem


\newpage
12. 
\begin{center}
	\adjustimage{max size={0.8\linewidth}{0.8\paperheight}}
	{BinTree/trees_exp/8.png}
\end{center}
Указати послідовність міток вершин дерева, зображеного на рис., що утвориться в результаті обходу дерева в прямому порядку. Мітки записувати через пробіл.\par Алгоритм починає роботу з кореня (на рис. це верхня вершина).
%enditem


\newpage
13. 
\begin{center}
	\adjustimage{max size={0.9\linewidth}{0.9\paperheight}}
	{Graphs/AdjList/56.png}
\end{center}
Для графа, зображеного на рис., вказати список суміжності вершини 3.\par
Формат оформлення відповіді:\par
список суміжності виписувати через пробіл на одному рядку, якщо пробілів десь буде більше одного - не важливо, упорядкування вершин у списку також значення не має.
%enditem


\newpage
14. Для графа на рисунку з вершини 2 запущено алгоритм пошуку в глибину, що будує кістякове дерево. Списки суміжності вершин переглядаються алгоритмом за зростанням міток вершин.\par
\begin{center}
	\adjustimage{max size={0.9\linewidth}{0.9\paperheight}}
	{Graphs/Traversals/66.png}
\end{center}
\par Які з наведених ребер входять до складу отриманого кістякового дерева?\par
1) (2, 4)\par
2) (4, 5)\par
3) (3, 5)\par
4) (1, 2)\par
5) (5, 6)\par
6) (3, 4)\par
{\vskip 15pt} Номери відповідних ребер записати за зростанням через пробіл.

%enditem


\newpage
15. Для графа на рисунку з вершини 0 запущено алгоритм пошуку в ширину, що будує кістякове дерево. Списки суміжності вершин переглядаються алгоритмом за зростанням міток вершин.\par
\begin{center}
	\adjustimage{max size={0.9\linewidth}{0.9\paperheight}}
	{Graphs/Traversals/185.png}
\end{center}
\par Які з наведених ребер входять до складу отриманого кістякового дерева?\par
1) (5, 6)\par
2) (1, 4)\par
3) (3, 5)\par
4) (2, 5)\par
5) (2, 6)\par
6) (3, 6)\par
{\vskip 15pt} Номери відповідних ребер записати за зростанням через пробіл.


%enditem


\newpage
16. 

Рядки журналу через двокрапку містять ім'я користувача (ідентифікатор), час події,тип події, код події, наприклад
\begin{verbatim}
s="username : 2022-05-05 13:26 : ERROR : 404"
\end{verbatim}
у коді 
\begin{verbatim}
res = re.fullmatch('???', s) 
s = ''
code_str = ??? 
\end{verbatim}
чим треба замінити кожен з ???, щоб значенням code\_str був код події? 


\vspace{10mm}
\textbf{Завдання 17} 

Задано програму, що містить код 1-2 класів, можливе успадкування, та клієнтський код.
У наведеному коді невеликі частини закриті. 
Також наведено вихід програми.
Необхідно відновити закриті частини, щоб отриманий код відповідав наведеному виходу програми.


\textbf{Завдання 18-19} 

Кілька запитань за завданнями для самостійної роботи.
\vspace{10mm}
Скоріш за все, що завдань буде трохи менше, а деякі попередні будуть переформульовані в стилі двох останніх.

%enditem

\end {large}}

%end

\newpage
%begin

{{
\begin {large}

\newpage
1. Що буде надруковано за виконання?

\begin{verbatim}
print('R', end=' ')
d=('a','b','c','d','e',0,1,2,3)
print(d[-12])
\end{verbatim}

%enditem


\newpage
2. Що буде надруковано за виконання?

\begin{verbatim}
print('R', end=' ')
g=(0,1,2,3,4,5,6,7,8,9)
print(g[:-5:-2])
\end{verbatim}

%enditem



\newpage
3. Що буде надруковано за виконання?

\begin{verbatim}
print('R', end=' ')
f=['0','1','2','3','4']
f.insert(5, [1,0])
print(f)
\end{verbatim}

%enditem



\newpage
4. Що буде надруковано за виконання?

\begin{verbatim}
print('R', end=' ')
b=[10,11,12,13,14]
del b[-5:0]
print(b)
\end{verbatim}

%enditem



\newpage
5. Укажіть ТИП значення виразу.\par
\begin{verbatim}
{21,2,4}
\end{verbatim}
%enditem


\newpage
6. Що буде надруковано за виконання?\par
\begin{verbatim}
try:
    t = {46:3, 48:3, 74:4, 79:8, 48:6}
    t[79] = 3
    for x in t.items():
        print(x, sep=' ', end=' ')
except:
    print('ERR')
\end{verbatim}
%enditem


\newpage
7. Задано словник \kw{d} з цілими ключами та значеннями. Записати ВИРАЗ, значенням якого є множина, що складається  з тих ключів словника, яким відповідають парні значення.
%enditem


\newpage
8. Що буде надруковано за виконання?
%Оригинал был утерян. Требует отладки!
\begin{verbatim}
_b = 0

class D:
    _b = 1
    
    def __init__(self):
        self._b = 2
        _b = 3
        D._b = 4

obj = D()
D._b = 5
try:
    print(_b, end=' ')
    print(obj._b, end=' ')
    print(D._b, end=' ')
    print(obj._b, end=' ')
    print(obj._D__b)
except:
    print('ERR')
\end{verbatim}
%enditem


\newpage
9. Що буде надруковано за виконання?
\begin{verbatim}
class C0:
    def b(self): print(1, end=' ')
    def h(self): print(2, end=' ')

class C1(C0):
    def b(self): print(3, end=' ')

try:
    C1().h()
    C1().b()
except:
    print('ERR')
\end{verbatim}
%enditem


\newpage
10. Що буде надруковано за виконання?
\begin{verbatim}
class A:
    c = 1

    def __init__(self):
        a = 2
        c = 3

class B(A):
    c = 4
    b = 5

    def __init__(self):
        super().__init__()
        self.b = 6

try:
    obj = B()
    obj.a = obj.b + 10
    B.a = 13
except: print('EE', end=' ')

try: print(obj.a, end= ' ')
except: print('E1', end=' ')

try: print(A.a, end= ' ')
except: print('E2', end=' ')

try: print(B.a, end= ' ')
except: print('E3', end=' ')
\end{verbatim}
%enditem


\newpage
11. Змінні \kw{next} та \kw{prev} вузла списку позначають наступний та попередній за ним вузол відповідно. Починаючи з голови, у список записано послідовні цілі числа від~$-38$ до~$22$. Нехай змінна \kw{p} позначає вузол, що зберігає значення~$4$.
Яке значення зберігається у вузлі \kw{p.next.next.next.prev.next} ?
%enditem


\newpage
12. 
\begin{center}
	\adjustimage{max size={0.8\linewidth}{0.8\paperheight}}
	{BinTree/trees_exp/58.png}
\end{center}
Указати послідовність міток вершин дерева, зображеного на рис., що утвориться в результаті обходу дерева в прямому порядку. Мітки записувати через пробіл.\par Алгоритм починає роботу з кореня (на рис. це верхня вершина).
%enditem


\newpage
13. 
\begin{center}
	\adjustimage{max size={0.9\linewidth}{0.9\paperheight}}
	{Graphs/AdjList/194.png}
\end{center}
Для графа, зображеного на рис., вказати список суміжності вершини 9.\par
Формат оформлення відповіді:\par
список суміжності виписувати через пробіл на одному рядку, якщо пробілів десь буде більше одного - не важливо, упорядкування вершин у списку також значення не має.
%enditem


\newpage
14. Для графа на рисунку з вершини 0 запущено алгоритм пошуку в глибину, що будує кістякове дерево. Списки суміжності вершин переглядаються алгоритмом за зростанням міток вершин.\par
\begin{center}
	\adjustimage{max size={0.9\linewidth}{0.9\paperheight}}
	{Graphs/Traversals/330.png}
\end{center}
\par Які з наведених ребер входять до складу отриманого кістякового дерева?\par
1) (0, 3)\par
2) (2, 4)\par
3) (0, 1)\par
4) (4, 5)\par
5) (0, 6)\par
6) (2, 3)\par
{\vskip 15pt} Номери відповідних ребер записати за зростанням через пробіл.

%enditem


\newpage
15. Для графа на рисунку з вершини 5 запущено алгоритм пошуку в ширину, що будує кістякове дерево. Списки суміжності вершин переглядаються алгоритмом за зростанням міток вершин.\par
\begin{center}
	\adjustimage{max size={0.9\linewidth}{0.9\paperheight}}
	{Graphs/Traversals/36.png}
\end{center}
\par Які з наведених ребер входять до складу отриманого кістякового дерева?\par
1) (1, 3)\par
2) (0, 6)\par
3) (5, 6)\par
4) (1, 5)\par
5) (2, 6)\par
6) (0, 5)\par
{\vskip 15pt} Номери відповідних ребер записати за зростанням через пробіл.


%enditem


\newpage
16. 

Рядки журналу через двокрапку містять ім'я користувача (ідентифікатор), час події,тип події, код події, наприклад
\begin{verbatim}
s="username : 2022-05-05 13:26 : ERROR : 404"
\end{verbatim}
у коді 
\begin{verbatim}
res = re.fullmatch('???', s) 
s = ''
code_str = ??? 
\end{verbatim}
чим треба замінити кожен з ???, щоб значенням code\_str був код події? 


\vspace{10mm}
\textbf{Завдання 17} 

Задано програму, що містить код 1-2 класів, можливе успадкування, та клієнтський код.
У наведеному коді невеликі частини закриті. 
Також наведено вихід програми.
Необхідно відновити закриті частини, щоб отриманий код відповідав наведеному виходу програми.


\textbf{Завдання 18-19} 

Кілька запитань за завданнями для самостійної роботи.
\vspace{10mm}
Скоріш за все, що завдань буде трохи менше, а деякі попередні будуть переформульовані в стилі двох останніх.

%enditem

\end {large}}

%end

\newpage
%begin

{{
\begin {large}

\newpage
1. Що буде надруковано за виконання?

\begin{verbatim}
print('R', end=' ')
f=('a','b','c','d','e',0,1,2,4)
print(f[-1])
\end{verbatim}

%enditem


\newpage
2. Що буде надруковано за виконання?

\begin{verbatim}
print('R', end=' ')
a=('0','1','2','3','4','5','6','7','8','9')
print(a[11:-2:-3])
\end{verbatim}

%enditem



\newpage
3. Що буде надруковано за виконання?

\begin{verbatim}
print('R', end=' ')
e=[10,11,12,13,14]
print(e.pop(-6), end=' ')
print(e)
\end{verbatim}

%enditem



\newpage
4. Що буде надруковано за виконання?

\begin{verbatim}
print('R', end=' ')
c=['0','1','2','3','4']
c[:-5]=[14]
print(c)
\end{verbatim}

%enditem



\newpage
5. Укажіть ТИП значення виразу.\par
\begin{verbatim}
[x for x in range(7)]
\end{verbatim}
%enditem


\newpage
6. Що буде надруковано за виконання?\par
\begin{verbatim}
try:
    s = {58:3, 61:1, 79:7, 58:6}
    s[58] = 1
    for x in s.values():
        print(x, sep=' ', end=' ')
except:
    print('ERR')
\end{verbatim}
%enditem


\newpage
7. Нехай змінна \kw{s} позначає послідовність цілих, по якій можна ітерувати.

Написати вираз, що перевіряє, чи правда, що всі елементи послідовності більші за задане число num.
%enditem


\newpage
8. Що буде надруковано за виконання?
\begin{verbatim}
d = 0
h = 1

class C:
    d = 2
    
    def __init__(self):
        global d
        self.d = 3
        d = 4
        C.d = 5

obj = C()
try:
    print(d, end=' ')
    print(h, end=' ')
    print(C.d, end=' ')
    print(obj.d)
except:
    print('ERR')
\end{verbatim}
%enditem


\newpage
9. Що буде надруковано за виконання?
\begin{verbatim}
class A0:
    def d(self): print(1, end=' ')
    def g(self): print(2, end=' ')

class A1(A0):
    def d(self): print(3, end=' ')

try:
    A0().g()
    A1().d()
except:
    print('ERR')
\end{verbatim}
%enditem


\newpage
10. Що буде надруковано за виконання?
\begin{verbatim}
class A:
    b = 1

    def __init__(self):
        c = 2
        b = 3

class B(A):
    b = 4
    a = 5

    def __init__(self):
        super().__init__()
        self.a = 6

try:
    obj = B()
    obj.a = obj.b + 10
    B.a = 13
except: print('EE', end=' ')

try: print(obj.a, end= ' ')
except: print('E1', end=' ')

try: print(A.a, end= ' ')
except: print('E2', end=' ')

try: print(B.a, end= ' ')
except: print('E3', end=' ')
\end{verbatim}
%enditem


\newpage
11. Змінні \kw{next} та \kw{prev} вузла списку позначають наступний та попередній за ним вузол відповідно. Починаючи з голови, у список записано послідовні цілі числа від~$-19$ до~$40$. Нехай змінна \kw{p} позначає вузол, що зберігає значення~$0$.
Яке значення зберігається у вузлі \kw{p.prev.prev.prev.next.prev} ?
%enditem


\newpage
12. 
\begin{center}
	\adjustimage{max size={0.8\linewidth}{0.8\paperheight}}
	{BinTree/trees_exp/53.png}
\end{center}
Указати послідовність міток вершин дерева, зображеного на рис., що утвориться в результаті обходу дерева в оберненому порядку. Мітки записувати через пробіл.\par Алгоритм починає роботу з кореня (на рис. це верхня вершина).
%enditem


\newpage
13. 
\begin{center}
	\adjustimage{max size={0.9\linewidth}{0.9\paperheight}}
	{Graphs/AdjList/297.png}
\end{center}
Для графа, зображеного на рис., вказати список суміжності вершини 0.\par
Формат оформлення відповіді:\par
список суміжності виписувати через пробіл на одному рядку, якщо пробілів десь буде більше одного - не важливо, упорядкування вершин у списку також значення не має.
%enditem


\newpage
14. Для графа на рисунку з вершини 2 запущено алгоритм пошуку в глибину, що будує кістякове дерево. Списки суміжності вершин переглядаються алгоритмом за зростанням міток вершин.\par
\begin{center}
	\adjustimage{max size={0.9\linewidth}{0.9\paperheight}}
	{Graphs/Traversals/113.png}
\end{center}
\par Які з наведених ребер входять до складу отриманого кістякового дерева?\par
1) (1, 3)\par
2) (5, 6)\par
3) (1, 4)\par
4) (2, 4)\par
5) (3, 6)\par
6) (0, 3)\par
{\vskip 15pt} Номери відповідних ребер записати за зростанням через пробіл.

%enditem


\newpage
15. Для графа на рисунку з вершини 1 запущено алгоритм пошуку в ширину, що будує кістякове дерево. Списки суміжності вершин переглядаються алгоритмом за зростанням міток вершин.\par
\begin{center}
	\adjustimage{max size={0.9\linewidth}{0.9\paperheight}}
	{Graphs/Traversals/195.png}
\end{center}
\par Які з наведених ребер входять до складу отриманого кістякового дерева?\par
1) (1, 3)\par
2) (5, 6)\par
3) (3, 6)\par
4) (1, 5)\par
5) (0, 3)\par
6) (3, 4)\par
{\vskip 15pt} Номери відповідних ребер записати за зростанням через пробіл.


%enditem


\newpage
16. 

Рядки журналу через двокрапку містять ім'я користувача (ідентифікатор), час події,тип події, код події, наприклад
\begin{verbatim}
s="username : 2022-05-05 13:26 : ERROR : 404"
\end{verbatim}
у коді 
\begin{verbatim}
res = re.fullmatch('???', s) 
s = ''
code_str = ??? 
\end{verbatim}
чим треба замінити кожен з ???, щоб значенням code\_str був код події? 


\vspace{10mm}
\textbf{Завдання 17} 

Задано програму, що містить код 1-2 класів, можливе успадкування, та клієнтський код.
У наведеному коді невеликі частини закриті. 
Також наведено вихід програми.
Необхідно відновити закриті частини, щоб отриманий код відповідав наведеному виходу програми.


\textbf{Завдання 18-19} 

Кілька запитань за завданнями для самостійної роботи.
\vspace{10mm}
Скоріш за все, що завдань буде трохи менше, а деякі попередні будуть переформульовані в стилі двох останніх.

%enditem

\end {large}}

%end

\newpage
%begin

{{
\begin {large}

\newpage
1. Що буде надруковано за виконання?

\begin{verbatim}
print('R', end=' ')
c=('a','b','c','d','e',0,1)
print(c[3])
\end{verbatim}

%enditem


\newpage
2. Що буде надруковано за виконання?

\begin{verbatim}
print('R', end=' ')
h=('0','1','2','3','4','5','6','7','8','9')
print(h[::2])
\end{verbatim}

%enditem



\newpage
3. Що буде надруковано за виконання?

\begin{verbatim}
print('R', end=' ')
c=[10,11,12,13,14]
c.insert(0, '11')
print(c)
\end{verbatim}

%enditem



\newpage
4. Що буде надруковано за виконання?

\begin{verbatim}
print('R', end=' ')
h=['0','1','2','3','4']
h[:-1]=[12]
print(h)
\end{verbatim}

%enditem



\newpage
5. Укажіть ТИП значення виразу.\par
\begin{verbatim}
{y:y*y for y in range(-7,7)}
\end{verbatim}
%enditem


\newpage
6. Що буде надруковано за виконання?\par
\begin{verbatim}
try:
    d = {10:0, 48:0, 64:4, 35:3, 10:4}
    d[76] = 9
    for x in d.items():
        print(*x, sep=' ', end=' ')
except:
    print('ERR')
\end{verbatim}
%enditem


\newpage
7. Нехай змінна \kw{s} позначає послідовність цілих, по якій можна ітерувати. 

Написати вираз, що перевіряє, чи правда, що послідовність  містить числа, більші за 5.
%enditem


\newpage
8. Що буде надруковано за виконання?
%Оригинал был утерян. Требует отладки!
\begin{verbatim}
__d = 0

class A:
    __d = 1
    
    def __init__(self):
        self.__d = 2
        __d = 3
        A.__d = 4

obj = A()
obj.__d = 5
try:
    print(__d, end=' ')
    print(obj.__d, end=' ')
    print(A.__d, end=' ')
    print(obj.__d, end=' ')
    print(obj._A__d)
except:
    print('ERR')
\end{verbatim}
%enditem


\newpage
9. Що буде надруковано за виконання?
\begin{verbatim}
class B0:
    def a(self): print(1, end=' ')
    def f(self): print(2, end=' ')

class B1(B0):
    def a(self): print(3, end=' ')

try:
    B1().f()
    B1().a()
except:
    print('ERR')
\end{verbatim}
%enditem


\newpage
10. Що буде надруковано за виконання?
\begin{verbatim}
class A:
    a = 1

    def __init__(self):
        c = 2
        c = 3

class B(A):
    c = 4
    b = 5

    def __init__(self):
        super().__init__()
        self.b = 6

try:
    obj = B()
    obj.c = obj.a + 10
    B.c = 13
except: print('EE', end=' ')

try: print(obj.c, end= ' ')
except: print('E1', end=' ')

try: print(A.c, end= ' ')
except: print('E2', end=' ')

try: print(B.c, end= ' ')
except: print('E3', end=' ')
\end{verbatim}
%enditem


\newpage
11. Змінні \kw{next} та \kw{prev} вузла списку позначають наступний та попередній за ним вузол відповідно. Починаючи з голови, у список записано послідовні цілі числа від~$-49$ до~$34$. Нехай змінна \kw{p} позначає вузол, що зберігає значення~$9$.
Яке значення зберігається у вузлі \kw{p.prev.next.prev.prev.next} ?
%enditem


\newpage
12. 
\begin{center}
	\adjustimage{max size={0.8\linewidth}{0.8\paperheight}}
	{BinTree/trees_exp/109.png}
\end{center}
Указати послідовність міток вершин дерева, зображеного на рис., що утвориться в результаті обходу дерева в прямому порядку. Мітки записувати через пробіл.\par Алгоритм починає роботу з кореня (на рис. це верхня вершина).
%enditem


\newpage
13. 
\begin{center}
	\adjustimage{max size={0.9\linewidth}{0.9\paperheight}}
	{Graphs/AdjList/177.png}
\end{center}
Для графа, зображеного на рис., вказати список суміжності вершини 3.\par
Формат оформлення відповіді:\par
список суміжності виписувати через пробіл на одному рядку, якщо пробілів десь буде більше одного - не важливо, упорядкування вершин у списку також значення не має.
%enditem


\newpage
14. Для графа на рисунку вказати послідовність, в якій алгоритм пошуку в глибину буде відмічати відвідані вершини, якщо списки суміжності вершин переглядаються алгоритмом за зростанням міток вершин та вершини відмічаються на першому вход\-женні у вершину.\par
Пошук розпочинається з вершини 2.\par

\begin{center}
	\adjustimage{max size={0.9\linewidth}{0.9\paperheight}}
	{Graphs/Traversals/41.png}
\end{center}
%enditem


\newpage
15. Для графа на рисунку з вершини 6 запущено алгоритм пошуку в ширину, що будує кістякове дерево. Списки суміжності вершин переглядаються алгоритмом за зростанням міток вершин.\par
\begin{center}
	\adjustimage{max size={0.9\linewidth}{0.9\paperheight}}
	{Graphs/Traversals/167.png}
\end{center}
\par Які з наведених ребер входять до складу отриманого кістякового дерева?\par
1) (0, 1)\par
2) (2, 6)\par
3) (4, 5)\par
4) (2, 3)\par
5) (2, 4)\par
6) (1, 6)\par
{\vskip 15pt} Номери відповідних ребер записати за зростанням через пробіл.


%enditem


\newpage
16. 

Рядки журналу через двокрапку містять ім'я користувача (ідентифікатор), час події,тип події, код події, наприклад
\begin{verbatim}
s="username : 2022-05-05 13:26 : ERROR : 404"
\end{verbatim}
у коді 
\begin{verbatim}
res = re.fullmatch('???', s) 
s = ''
code_str = ??? 
\end{verbatim}
чим треба замінити кожен з ???, щоб значенням code\_str був код події? 


\vspace{10mm}
\textbf{Завдання 17} 

Задано програму, що містить код 1-2 класів, можливе успадкування, та клієнтський код.
У наведеному коді невеликі частини закриті. 
Також наведено вихід програми.
Необхідно відновити закриті частини, щоб отриманий код відповідав наведеному виходу програми.


\textbf{Завдання 18-19} 

Кілька запитань за завданнями для самостійної роботи.
\vspace{10mm}
Скоріш за все, що завдань буде трохи менше, а деякі попередні будуть переформульовані в стилі двох останніх.

%enditem

\end {large}}

%end

\newpage
%begin

{{
\begin {large}

\newpage
1. Що буде надруковано за виконання?

\begin{verbatim}
print('R', end=' ')
g=('a','b','c','d','e',0,1,2,4)
print(g[-10])
\end{verbatim}

%enditem


\newpage
2. Що буде надруковано за виконання?

\begin{verbatim}
print('R', end=' ')
e=['0','1','2','3','4','5','6','7','8','9']
print(e[:9:-2])
\end{verbatim}

%enditem



\newpage
3. Що буде надруковано за виконання?

\begin{verbatim}
print('R', end=' ')
g=['0','1','2','3','4']
g+=[1,2]
print(g)
\end{verbatim}

%enditem



\newpage
4. Що буде надруковано за виконання?

\begin{verbatim}
print('R', end=' ')
d=[10,11,12,13,14]
del d[-3:-6]
print(d)
\end{verbatim}

%enditem



\newpage
5. Укажіть ТИП значення виразу.\par
\begin{verbatim}
(11,2,3)
\end{verbatim}
%enditem


\newpage
6. Що буде надруковано за виконання?\par
\begin{verbatim}
try:
    s = {82:5, 31:3, 13:1, 13:7}
    s[31] = 9
    for x in s.keys():
        print(x, sep=' ', end=' ')
except:
    print('ERR')
\end{verbatim}
%enditem


\newpage
7. Нехай змінна \kw{s} позначає дуже довгу послідовність цілих, по якій можна ітерувати. 

Написати вираз-генератор, що задає підпослідовність її елементів, що є від'єм\-ни\-ми. 

Обмеження: усі елементи шуканої підпослідовності одночасно в пам'яті розміщені бути не можуть. 
%enditem


\newpage
8. Що буде надруковано за виконання?
%Оригинал был утерян. Требует отладки!
\begin{verbatim}
c = 0

class B:
    c = 1
    
    def __init__(self):
        self.c = 2
        c = 3
        B.c = 4

obj = B()
obj.c = 5
try:
    print(c, end=' ')
    print(obj.c, end=' ')
    print(B.c, end=' ')
    print(obj.c, end=' ')
    print(obj._B__c)
except:
    print('ERR')
\end{verbatim}
%enditem


\newpage
9. Що буде надруковано за виконання?
\begin{verbatim}
class D0:
    def c(self): print(1, end=' ')
    def g(self): print(2, end=' ')

class D1(D0):
    def c(self): print(3, end=' ')

try:
    D1().g()
    D0().c()
except:
    print('ERR')
\end{verbatim}
%enditem


\newpage
10. Що буде надруковано за виконання?
\begin{verbatim}
class A:
    a = 1

    def __g1(self): print(2, end=' ')
    def _g2(self): print(3, end=' ')

    def main(self):
        try: print(self.a, end=' ')
        except: print('E1', end=' ')

        try: self.__g1()
        except: print('E2', end=' ')

        try: self._g2()
        except: print('E3', end=' ')


class B(A):
    a = 4

    def __g1(self): print(5, end=' ')
    def _g2(self): print(6, end=' ')

try: B().main()
except: print('E4', end=' ')
\end{verbatim}
%enditem


\newpage
11. Змінні \kw{next} та \kw{prev} вузла списку позначають наступний та попередній за ним вузол відповідно. Починаючи з голови, у список записано послідовні цілі числа від~$-46$ до~$43$. Нехай змінна \kw{p} позначає вузол, що зберігає значення~$-3$.
Яке значення зберігається у вузлі \kw{p.next.prev.next.next.prev} ?
%enditem


\newpage
12. 
\begin{center}
	\adjustimage{max size={0.8\linewidth}{0.8\paperheight}}
	{BinTree/trees_exp/93.png}
\end{center}
Указати послідовність міток вершин дерева, зображеного на рис., що утвориться в результаті обходу дерева в оберненому порядку. Мітки записувати через пробіл.\par Алгоритм починає роботу з кореня (на рис. це верхня вершина).
%enditem


\newpage
13. 
\begin{center}
	\adjustimage{max size={0.9\linewidth}{0.9\paperheight}}
	{Graphs/AdjList/68.png}
\end{center}
Для графа, зображеного на рис., вказати список суміжності вершини 7.\par
Формат оформлення відповіді:\par
список суміжності виписувати через пробіл на одному рядку, якщо пробілів десь буде більше одного - не важливо, упорядкування вершин у списку також значення не має.
%enditem


\newpage
14. Для графа на рисунку вказати послідовність, в якій алгоритм пошуку в глибину буде відмічати відвідані вершини, якщо списки суміжності вершин переглядаються алгоритмом за зростанням міток вершин та вершини відмічаються на першому вход\-женні у вершину.\par
Пошук розпочинається з вершини 2.\par

\begin{center}
	\adjustimage{max size={0.9\linewidth}{0.9\paperheight}}
	{Graphs/Traversals/303.png}
\end{center}
%enditem


\newpage
15. Для графа на рисунку з вершини 2 запущено алгоритм пошуку в ширину, що будує кістякове дерево. Списки суміжності вершин переглядаються алгоритмом за зростанням міток вершин.\par
\begin{center}
	\adjustimage{max size={0.9\linewidth}{0.9\paperheight}}
	{Graphs/Traversals/191.png}
\end{center}
\par Які з наведених ребер входять до складу отриманого кістякового дерева?\par
1) (4, 6)\par
2) (0, 4)\par
3) (1, 4)\par
4) (1, 2)\par
5) (0, 2)\par
6) (3, 4)\par
{\vskip 15pt} Номери відповідних ребер записати за зростанням через пробіл.


%enditem


\newpage
16. 

Рядки журналу через двокрапку містять ім'я користувача (ідентифікатор), час події,тип події, код події, наприклад
\begin{verbatim}
s="username : 2022-05-05 13:26 : ERROR : 404"
\end{verbatim}
у коді 
\begin{verbatim}
res = re.fullmatch('???', s) 
s = ''
code_str = ??? 
\end{verbatim}
чим треба замінити кожен з ???, щоб значенням code\_str був код події? 


\vspace{10mm}
\textbf{Завдання 17} 

Задано програму, що містить код 1-2 класів, можливе успадкування, та клієнтський код.
У наведеному коді невеликі частини закриті. 
Також наведено вихід програми.
Необхідно відновити закриті частини, щоб отриманий код відповідав наведеному виходу програми.


\textbf{Завдання 18-19} 

Кілька запитань за завданнями для самостійної роботи.
\vspace{10mm}
Скоріш за все, що завдань буде трохи менше, а деякі попередні будуть переформульовані в стилі двох останніх.

%enditem

\end {large}}

%end

\newpage
%begin

{{
\begin {large}

\newpage
1. Що буде надруковано за виконання?

\begin{verbatim}
print('R', end=' ')
b=('a','b','c','d','e',0,1,2)
print(b[-13])
\end{verbatim}

%enditem


\newpage
2. Що буде надруковано за виконання?

\begin{verbatim}
print('R', end=' ')
b=[0,1,2,3,4,5,6,7,8,9]
print(b[2:12:-3])
\end{verbatim}

%enditem



\newpage
3. Що буде надруковано за виконання?

\begin{verbatim}
print('R', end=' ')
a=['0','1','2','3','4']
a.append([2,3])
print(a)
\end{verbatim}

%enditem



\newpage
4. Що буде надруковано за виконання?

\begin{verbatim}
print('R', end=' ')
g=['0','1','2','3','4']
g[0:]=(13)
print(g)
\end{verbatim}

%enditem



\newpage
5. Укажіть ТИП значення виразу.\par
\begin{verbatim}
{1:3, 5:8}
\end{verbatim}
%enditem


\newpage
6. Що буде надруковано за виконання?\par
\begin{verbatim}
try:
    t = {47:1, 20:1, 87:1, 87:3}
    t[50] = 8
    for x, y in t.items():
        print(x, y, sep=' ', end=' ')
except:
    print('ERR')
\end{verbatim}
%enditem


\newpage
7. Записати ВИРАЗ, значенням якого є множина, що складається з цілих чисел від 19 до 2010 (включно), що не діляться на 2.
%enditem


\newpage
8. Що буде надруковано за виконання?
%Оригинал был утерян. Требует отладки!
\begin{verbatim}
_a = 0

class C:
    _a = 1
    
    def __init__(self):
        self._a = 2
        _a = 3
        C._a = 4

obj = C()
C._a = 5
try:
    print(_a, end=' ')
    print(obj._a, end=' ')
    print(C._a, end=' ')
    print(obj._a, end=' ')
    print(obj._C__a)
except:
    print('ERR')
\end{verbatim}
%enditem


\newpage
9. Що буде надруковано за виконання?
\begin{verbatim}
class B0:
    def b(self): print(1, end=' ')
    def f(self): print(2, end=' ')

class B1(B0):
    def b(self): print(3, end=' ')

try:
    B1().f()
    B1().b()
except:
    print('ERR')
\end{verbatim}
%enditem


\newpage
10. Що буде надруковано за виконання?
\begin{verbatim}
class A:
    b = 1

    def _h1(self): print(2, end=' ')
    def _h2(self): print(3, end=' ')

    def main(self):
        try: print(self.b, end=' ')
        except: print('E1', end=' ')

        try: self._h1()
        except: print('E2', end=' ')

        try: self._h2()
        except: print('E3', end=' ')


class B(A):
    b = 4

    def _h1(self): print(5, end=' ')
    def _h2(self): print(6, end=' ')

try: B().main()
except: print('E4', end=' ')
\end{verbatim}
%enditem


\newpage
11. Змінні \kw{next} та \kw{prev} вузла списку позначають наступний та попередній за ним вузол відповідно. Починаючи з голови, у список записано послідовні цілі числа від~$-22$ до~$50$. Нехай змінна \kw{p} позначає вузол, що зберігає значення~$-2$.
Яке значення зберігається у вузлі \kw{p.prev.prev.next.prev.next} ?
%enditem


\newpage
12. 
\begin{center}
	\adjustimage{max size={0.8\linewidth}{0.8\paperheight}}
	{BinTree/trees_exp/65.png}
\end{center}
Указати послідовність міток вершин дерева, зображеного на рис., що утвориться в результаті обходу дерева в прямому порядку. Мітки записувати через пробіл.\par Алгоритм починає роботу з кореня (на рис. це верхня вершина).
%enditem


\newpage
13. 
\begin{center}
	\adjustimage{max size={0.9\linewidth}{0.9\paperheight}}
	{Graphs/AdjList/181.png}
\end{center}
Для графа, зображеного на рис., вказати список суміжності вершини 7.\par
Формат оформлення відповіді:\par
список суміжності виписувати через пробіл на одному рядку, якщо пробілів десь буде більше одного - не важливо, упорядкування вершин у списку також значення не має.
%enditem


\newpage
14. Для графа на рисунку з вершини 5 запущено алгоритм пошуку в глибину, що будує кістякове дерево. Списки суміжності вершин переглядаються алгоритмом за зростанням міток вершин.\par
\begin{center}
	\adjustimage{max size={0.9\linewidth}{0.9\paperheight}}
	{Graphs/Traversals/257.png}
\end{center}
\par Які з наведених ребер входять до складу отриманого кістякового дерева?\par
1) (5, 6)\par
2) (2, 3)\par
3) (0, 5)\par
4) (0, 1)\par
5) (4, 5)\par
6) (1, 3)\par
{\vskip 15pt} Номери відповідних ребер записати за зростанням через пробіл.

%enditem


\newpage
15. Для графа на рисунку з вершини 4 запущено алгоритм пошуку в ширину, що будує кістякове дерево. Списки суміжності вершин переглядаються алгоритмом за зростанням міток вершин.\par
\begin{center}
	\adjustimage{max size={0.9\linewidth}{0.9\paperheight}}
	{Graphs/Traversals/39.png}
\end{center}
\par Які з наведених ребер входять до складу отриманого кістякового дерева?\par
1) (2, 5)\par
2) (0, 6)\par
3) (2, 6)\par
4) (2, 4)\par
5) (0, 4)\par
6) (2, 3)\par
{\vskip 15pt} Номери відповідних ребер записати за зростанням через пробіл.


%enditem


\newpage
16. 

Рядки журналу через двокрапку містять ім'я користувача (ідентифікатор), час події,тип події, код події, наприклад
\begin{verbatim}
s="username : 2022-05-05 13:26 : ERROR : 404"
\end{verbatim}
у коді 
\begin{verbatim}
res = re.fullmatch('???', s) 
s = ''
code_str = ??? 
\end{verbatim}
чим треба замінити кожен з ???, щоб значенням code\_str був код події? 


\vspace{10mm}
\textbf{Завдання 17} 

Задано програму, що містить код 1-2 класів, можливе успадкування, та клієнтський код.
У наведеному коді невеликі частини закриті. 
Також наведено вихід програми.
Необхідно відновити закриті частини, щоб отриманий код відповідав наведеному виходу програми.


\textbf{Завдання 18-19} 

Кілька запитань за завданнями для самостійної роботи.
\vspace{10mm}
Скоріш за все, що завдань буде трохи менше, а деякі попередні будуть переформульовані в стилі двох останніх.

%enditem

\end {large}}

%end

\newpage
%begin

{{
\begin {large}

\newpage
1. Що буде надруковано за виконання?

\begin{verbatim}
print('R', end=' ')
a=('a','b','c','d','e',0,1,2,3,4)
print(a[5])
\end{verbatim}

%enditem


\newpage
2. Що буде надруковано за виконання?

\begin{verbatim}
print('R', end=' ')
d=(0,1,2,3,4,5,6,7,8,9)
print(d[-7:-12:-1])
\end{verbatim}

%enditem



\newpage
3. Що буде надруковано за виконання?

\begin{verbatim}
print('R', end=' ')
b=[10,11,12,13,14]
b.extend('12')
print(b)
\end{verbatim}

%enditem



\newpage
4. Що буде надруковано за виконання?

\begin{verbatim}
print('R', end=' ')
f=[10,11,12,13,14]
f[-4:1]=[1,2]
print(f)
\end{verbatim}

%enditem



\newpage
5. Укажіть ТИП значення виразу.\par
\begin{verbatim}
{x for x in range(7)}
\end{verbatim}
%enditem


\newpage
6. Що буде надруковано за виконання?\par
\begin{verbatim}
try:
    d = {15:5, 53:7, 29:2, 29:3}
    d[34] = 1
    for x in d :
        print(x, sep=' ', end=' ')
except:
    print('ERR')
\end{verbatim}
%enditem


\newpage
7. Нехай змінна \kw{s} позначає послідовність цілих, по якій можна ітерувати. 

Написати вираз, що перевіряє, чи правда, що послідовність  містить хоча б одне число, що ділиться на 7.
%enditem


\newpage
8. Що буде надруковано за виконання?
\begin{verbatim}
c = 0
k = 1

class B:
    k = 2
    
    def __init__(self):
        self.k = 3
        k = 4
        B.k = 5

obj = B()
try:
    print(c, end=' ')
    print(k, end=' ')
    print(B.k, end=' ')
    print(obj.k)
except:
    print('ERR')
\end{verbatim}
%enditem


\newpage
9. Що буде надруковано за виконання?
\begin{verbatim}
class D0:
    def c(self): print(1, end=' ')
    def h(self): print(2, end=' ')

class D1(D0):
    def c(self): print(3, end=' ')

try:
    D1().h()
    D1().c()
except:
    print('ERR')
\end{verbatim}
%enditem


\newpage
10. Що буде надруковано за виконання?
\begin{verbatim}
class A:
    c = 1

    def _g1(self): print(2, end=' ')
    def __g2(self): print(3, end=' ')

    def main(self):
        try: print(self.c, end=' ')
        except: print('E1', end=' ')

        try: self._g1()
        except: print('E2', end=' ')

        try: self.__g2()
        except: print('E3', end=' ')


class B(A):
    c = 4

    def _g1(self): print(5, end=' ')
    def __g2(self): print(6, end=' ')

try: B().main()
except: print('E4', end=' ')
\end{verbatim}
%enditem


\newpage
11. Змінні \kw{next} та \kw{prev} вузла списку позначають наступний та попередній за ним вузол відповідно. Починаючи з голови, у список записано послідовні цілі числа від~$-42$ до~$24$. Нехай змінна \kw{p} позначає вузол, що зберігає значення~$-1$.
Яке значення зберігається у вузлі \kw{p.next.next.prev.next.prev} ?
%enditem


\newpage
12. 
\begin{center}
	\adjustimage{max size={0.8\linewidth}{0.8\paperheight}}
	{BinTree/trees_exp/145.png}
\end{center}
Указати послідовність міток вершин дерева, зображеного на рис., що утвориться в результаті обходу дерева в оберненому порядку. Мітки записувати через пробіл.\par Алгоритм починає роботу з кореня (на рис. це верхня вершина).
%enditem


\newpage
13. 
\begin{center}
	\adjustimage{max size={0.9\linewidth}{0.9\paperheight}}
	{Graphs/AdjList/283.png}
\end{center}
Для графа, зображеного на рис., вказати список суміжності вершини 7.\par
Формат оформлення відповіді:\par
список суміжності виписувати через пробіл на одному рядку, якщо пробілів десь буде більше одного - не важливо, упорядкування вершин у списку також значення не має.
%enditem


\newpage
14. Для графа на рисунку вказати послідовність, в якій алгоритм пошуку в глибину буде відмічати відвідані вершини, якщо списки суміжності вершин переглядаються алгоритмом за зростанням міток вершин та вершини відмічаються на першому вход\-женні у вершину.\par
Пошук розпочинається з вершини 1.\par

\begin{center}
	\adjustimage{max size={0.9\linewidth}{0.9\paperheight}}
	{Graphs/Traversals/18.png}
\end{center}
%enditem


\newpage
15. Для графа на рисунку з вершини 2 запущено алгоритм пошуку в ширину, що будує кістякове дерево. Списки суміжності вершин переглядаються алгоритмом за зростанням міток вершин.\par
\begin{center}
	\adjustimage{max size={0.9\linewidth}{0.9\paperheight}}
	{Graphs/Traversals/196.png}
\end{center}
\par Які з наведених ребер входять до складу отриманого кістякового дерева?\par
1) (0, 1)\par
2) (3, 4)\par
3) (1, 2)\par
4) (0, 2)\par
5) (2, 5)\par
6) (0, 6)\par
{\vskip 15pt} Номери відповідних ребер записати за зростанням через пробіл.


%enditem


\newpage
16. 

Рядки журналу через двокрапку містять ім'я користувача (ідентифікатор), час події,тип події, код події, наприклад
\begin{verbatim}
s="username : 2022-05-05 13:26 : ERROR : 404"
\end{verbatim}
у коді 
\begin{verbatim}
res = re.fullmatch('???', s) 
s = ''
code_str = ??? 
\end{verbatim}
чим треба замінити кожен з ???, щоб значенням code\_str був код події? 


\vspace{10mm}
\textbf{Завдання 17} 

Задано програму, що містить код 1-2 класів, можливе успадкування, та клієнтський код.
У наведеному коді невеликі частини закриті. 
Також наведено вихід програми.
Необхідно відновити закриті частини, щоб отриманий код відповідав наведеному виходу програми.


\textbf{Завдання 18-19} 

Кілька запитань за завданнями для самостійної роботи.
\vspace{10mm}
Скоріш за все, що завдань буде трохи менше, а деякі попередні будуть переформульовані в стилі двох останніх.

%enditem

\end {large}}

%end

\newpage
%begin

{{
\begin {large}

\newpage
1. Що буде надруковано за виконання?

\begin{verbatim}
print('R', end=' ')
e=('a','b','c','d','e',0,1,2,3)
print(e[12])
\end{verbatim}

%enditem


\newpage
2. Що буде надруковано за виконання?

\begin{verbatim}
print('R', end=' ')
c='0123456789'
print(c[-8:-11:3])
\end{verbatim}

%enditem



\newpage
3. Що буде надруковано за виконання?

\begin{verbatim}
print('R', end=' ')
h=['0','1','2','3','4']
h.remove('4')
print(h)
\end{verbatim}

%enditem



\newpage
4. Що буде надруковано за виконання?

\begin{verbatim}
print('R', end=' ')
a=['0','1','2','3','4']
a[:8]='12'
print(a)
\end{verbatim}

%enditem



\newpage
5. Укажіть ТИП значення виразу.\par
\begin{verbatim}
{x for x in range(7) if el%3==1}
\end{verbatim}
%enditem


\newpage
6. Що буде надруковано за виконання?\par
\begin{verbatim}
try:
    t = {41:2, 63:9, 40:9, 65:1, 40:0}
    t[60] = 6
    for x in t.items():
        print(*x, sep=' ', end=' ')
except:
    print('ERR')
\end{verbatim}
%enditem


\newpage
7. 
Записати ВИРАЗ, значенням якого є список, що складається з квадратних коренів цілих чисел від 11 до 2019 (включно), що діляться на 7, записаних в обернено\-му порядку.
%enditem


\newpage
8. Що буде надруковано за виконання?
%Оригинал был утерян. Требует отладки!
\begin{verbatim}
b = 0

class C:
    b = 1
    
    def __init__(self):
        self.b = 2
        b = 3
        C.b = 4

obj = C()
obj.b = 5
try:
    print(b, end=' ')
    print(obj.b, end=' ')
    print(C.b, end=' ')
    print(obj.b, end=' ')
    print(obj._C__b)
except:
    print('ERR')
\end{verbatim}
%enditem


\newpage
9. Що буде надруковано за виконання?
\begin{verbatim}
class C0:
    def d(self): print(1, end=' ')
    def g(self): print(2, end=' ')

class C1(C0):
    def d(self): print(3, end=' ')

try:
    C0().g()
    C1().d()
except:
    print('ERR')
\end{verbatim}
%enditem


\newpage
10. Що буде надруковано за виконання?
\begin{verbatim}
class A:
    d = 1

    def _h1(self): print(2, end=' ')
    def h2(self): print(3, end=' ')

    def main(self):
        try: print(self.d, end=' ')
        except: print('E1', end=' ')

        try: self._h1()
        except: print('E2', end=' ')

        try: self.h2()
        except: print('E3', end=' ')


class B(A):
    d = 4

    def _h1(self): print(5, end=' ')
    def h2(self): print(6, end=' ')

try: B().main()
except: print('E4', end=' ')
\end{verbatim}
%enditem


\newpage
11. Змінні \kw{next} та \kw{prev} вузла списку позначають наступний та попередній за ним вузол відповідно. Починаючи з голови, у список записано послідовні цілі числа від~$-52$ до~$18$. Нехай змінна \kw{p} позначає вузол, що зберігає значення~$-7$.
Яке значення зберігається у вузлі \kw{p.prev.next.next.prev.next} ?
%enditem


\newpage
12. 
\begin{center}
	\adjustimage{max size={0.8\linewidth}{0.8\paperheight}}
	{BinTree/trees_exp/9.png}
\end{center}
Указати послідовність міток вершин дерева, зображеного на рис., що утвориться в результаті обходу дерева в прямому порядку. Мітки записувати через пробіл.\par Алгоритм починає роботу з кореня (на рис. це верхня вершина).
%enditem


\newpage
13. 
\begin{center}
	\adjustimage{max size={0.9\linewidth}{0.9\paperheight}}
	{Graphs/AdjList/161.png}
\end{center}
Для графа, зображеного на рис., вказати список суміжності вершини 2.\par
Формат оформлення відповіді:\par
список суміжності виписувати через пробіл на одному рядку, якщо пробілів десь буде більше одного - не важливо, упорядкування вершин у списку також значення не має.
%enditem


\newpage
14. Для графа на рисунку з вершини 3 запущено алгоритм пошуку в глибину, що будує кістякове дерево. Списки суміжності вершин переглядаються алгоритмом за зростанням міток вершин.\par
\begin{center}
	\adjustimage{max size={0.9\linewidth}{0.9\paperheight}}
	{Graphs/Traversals/309.png}
\end{center}
\par Які з наведених ребер входять до складу отриманого кістякового дерева?\par
1) (2, 6)\par
2) (2, 5)\par
3) (2, 3)\par
4) (0, 4)\par
5) (2, 4)\par
6) (1, 4)\par
{\vskip 15pt} Номери відповідних ребер записати за зростанням через пробіл.

%enditem


\newpage
15. Для графа на рисунку з вершини 0 запущено алгоритм пошуку в ширину, що будує кістякове дерево. Списки суміжності вершин переглядаються алгоритмом за зростанням міток вершин.\par
\begin{center}
	\adjustimage{max size={0.9\linewidth}{0.9\paperheight}}
	{Graphs/Traversals/235.png}
\end{center}
\par Які з наведених ребер входять до складу отриманого кістякового дерева?\par
1) (0, 2)\par
2) (3, 5)\par
3) (0, 3)\par
4) (2, 4)\par
5) (0, 5)\par
6) (1, 2)\par
{\vskip 15pt} Номери відповідних ребер записати за зростанням через пробіл.


%enditem


\newpage
16. 

Рядки журналу через двокрапку містять ім'я користувача (ідентифікатор), час події,тип події, код події, наприклад
\begin{verbatim}
s="username : 2022-05-05 13:26 : ERROR : 404"
\end{verbatim}
у коді 
\begin{verbatim}
res = re.fullmatch('???', s) 
s = ''
code_str = ??? 
\end{verbatim}
чим треба замінити кожен з ???, щоб значенням code\_str був код події? 


\vspace{10mm}
\textbf{Завдання 17} 

Задано програму, що містить код 1-2 класів, можливе успадкування, та клієнтський код.
У наведеному коді невеликі частини закриті. 
Також наведено вихід програми.
Необхідно відновити закриті частини, щоб отриманий код відповідав наведеному виходу програми.


\textbf{Завдання 18-19} 

Кілька запитань за завданнями для самостійної роботи.
\vspace{10mm}
Скоріш за все, що завдань буде трохи менше, а деякі попередні будуть переформульовані в стилі двох останніх.

%enditem

\end {large}}

%end

\newpage
%begin

{{
\begin {large}

\newpage
1. Що буде надруковано за виконання?

\begin{verbatim}
print('R', end=' ')
h=('a','b','c','d','e',0,1)
print(h[-9])
\end{verbatim}

%enditem


\newpage
2. Що буде надруковано за виконання?

\begin{verbatim}
print('R', end=' ')
f=['0','1','2','3','4','5','6','7','8','9']
print(f[:5:-2])
\end{verbatim}

%enditem



\newpage
3. Що буде надруковано за виконання?

\begin{verbatim}
print('R', end=' ')
d=[10,11,12,13,14]
print(d.index(17), end=' ')
\end{verbatim}

%enditem



\newpage
4. Що буде надруковано за виконання?

\begin{verbatim}
print('R', end=' ')
e=[10,11,12,13,14]
e[1:-2]='10'
print(e)
\end{verbatim}

%enditem



\newpage
5. Укажіть ТИП значення виразу.\par
\begin{verbatim}
{21,2,4}
\end{verbatim}
%enditem


\newpage
6. Що буде надруковано за виконання?\par
\begin{verbatim}
try:
    d = {52:3, 37:3, 70:6, 70:5}
    d[31] = 5
    for x in d :
        print(x, sep=' ', end=' ')
except:
    print('ERR')
\end{verbatim}
%enditem


\newpage
7. Записати ВИРАЗ, значенням якого є словник, ключами якого є цілі числа від 20 до 2004 (включно), що діляться на 3, а ключам відповідають їх куби 
%enditem


\newpage
8. Що буде надруковано за виконання?
\begin{verbatim}
b = 0
m = 1

class A:
    m = 2
    
    def __init__(self):
        self.m = 3
        m = 4
        A.m = 5

obj = A()
try:
    print(b, end=' ')
    print(m, end=' ')
    print(A.m, end=' ')
    print(obj.m)
except:
    print('ERR')
\end{verbatim}
%enditem


\newpage
9. Що буде надруковано за виконання?
\begin{verbatim}
class A0:
    def a(self): print(1, end=' ')
    def f(self): print(2, end=' ')

class A1(A0):
    def a(self): print(3, end=' ')

try:
    A1().f()
    A0().a()
except:
    print('ERR')
\end{verbatim}
%enditem


\newpage
10. Що буде надруковано за виконання?
\begin{verbatim}
class A:
    a = 1

    def f1(self): print(2, end=' ')
    def _f2(self): print(3, end=' ')

    def main(self):
        try: print(self.a, end=' ')
        except: print('E1', end=' ')

        try: self.f1()
        except: print('E2', end=' ')

        try: self._f2()
        except: print('E3', end=' ')


class B(A):
    a = 4

    def f1(self): print(5, end=' ')
    def _f2(self): print(6, end=' ')

try: B().main()
except: print('E4', end=' ')
\end{verbatim}
%enditem


\newpage
11. Змінні \kw{next} та \kw{prev} вузла списку позначають наступний та попередній за ним вузол відповідно. Починаючи з голови, у список записано послідовні цілі числа від~$-27$ до~$27$. Нехай змінна \kw{p} позначає вузол, що зберігає значення~$2$.
Яке значення зберігається у вузлі \kw{p.next.prev.prev.next.prev} ?
%enditem


\newpage
12. 
\begin{center}
	\adjustimage{max size={0.8\linewidth}{0.8\paperheight}}
	{BinTree/trees_exp/2.png}
\end{center}
Указати послідовність міток вершин дерева, зображеного на рис., що утвориться в результаті обходу дерева в оберненому порядку. Мітки записувати через пробіл.\par Алгоритм починає роботу з кореня (на рис. це верхня вершина).
%enditem


\newpage
13. 
\begin{center}
	\adjustimage{max size={0.9\linewidth}{0.9\paperheight}}
	{Graphs/AdjList/53.png}
\end{center}
Для графа, зображеного на рис., вказати список суміжності вершини 1.\par
Формат оформлення відповіді:\par
список суміжності виписувати через пробіл на одному рядку, якщо пробілів десь буде більше одного - не важливо, упорядкування вершин у списку також значення не має.
%enditem


\newpage
14. Для графа на рисунку вказати послідовність, в якій алгоритм пошуку в глибину буде відмічати відвідані вершини, якщо списки суміжності вершин переглядаються алгоритмом за зростанням міток вершин та вершини відмічаються на першому вход\-женні у вершину.\par
Пошук розпочинається з вершини 0.\par

\begin{center}
	\adjustimage{max size={0.9\linewidth}{0.9\paperheight}}
	{Graphs/Traversals/130.png}
\end{center}
%enditem


\newpage
15. Для графа на рисунку з вершини 0 запущено алгоритм пошуку в ширину, що будує кістякове дерево. Списки суміжності вершин переглядаються алгоритмом за зростанням міток вершин.\par
\begin{center}
	\adjustimage{max size={0.9\linewidth}{0.9\paperheight}}
	{Graphs/Traversals/5.png}
\end{center}
\par Які з наведених ребер входять до складу отриманого кістякового дерева?\par
1) (3, 5)\par
2) (1, 2)\par
3) (4, 6)\par
4) (2, 6)\par
5) (4, 5)\par
6) (1, 5)\par
{\vskip 15pt} Номери відповідних ребер записати за зростанням через пробіл.


%enditem


\newpage
16. 

Рядки журналу через двокрапку містять ім'я користувача (ідентифікатор), час події,тип події, код події, наприклад
\begin{verbatim}
s="username : 2022-05-05 13:26 : ERROR : 404"
\end{verbatim}
у коді 
\begin{verbatim}
res = re.fullmatch('???', s) 
s = ''
code_str = ??? 
\end{verbatim}
чим треба замінити кожен з ???, щоб значенням code\_str був код події? 


\vspace{10mm}
\textbf{Завдання 17} 

Задано програму, що містить код 1-2 класів, можливе успадкування, та клієнтський код.
У наведеному коді невеликі частини закриті. 
Також наведено вихід програми.
Необхідно відновити закриті частини, щоб отриманий код відповідав наведеному виходу програми.


\textbf{Завдання 18-19} 

Кілька запитань за завданнями для самостійної роботи.
\vspace{10mm}
Скоріш за все, що завдань буде трохи менше, а деякі попередні будуть переформульовані в стилі двох останніх.

%enditem

\end {large}}

%end

\newpage
%begin

{{
\begin {large}

\newpage
1. Що буде надруковано за виконання?

\begin{verbatim}
print('R', end=' ')
g=('a','b','c','d','e',0,1,2)
print(g[-7])
\end{verbatim}

%enditem


\newpage
2. Що буде надруковано за виконання?

\begin{verbatim}
print('R', end=' ')
h=(0,1,2,3,4,5,6,7,8,9)
print(h[7:3:-1])
\end{verbatim}

%enditem



\newpage
3. Що буде надруковано за виконання?

\begin{verbatim}
print('R', end=' ')
d=['0','1','2','3','4']
d.extend('10')
print(d)
\end{verbatim}

%enditem



\newpage
4. Що буде надруковано за виконання?

\begin{verbatim}
print('R', end=' ')
a=['0','1','2','3','4']
del a[4:-4]
print(a)
\end{verbatim}

%enditem



\newpage
5. Укажіть ТИП значення виразу.\par
\begin{verbatim}
{y:y*y for y in range(-7,7)}
\end{verbatim}
%enditem


\newpage
6. Що буде надруковано за виконання?\par
\begin{verbatim}
try:
    s = {48:5, 21:5, 39:8, 21:8}
    s[14] = 8
    for x in s.items():
        print(x, sep=' ', end=' ')
except:
    print('ERR')
\end{verbatim}
%enditem


\newpage
7. Нехай змінна \kw{s} позначає дуже довгу послідовність цілих, по якій можна ітерувати. 

Написати вираз-генератор, що задає підпослідовність її елементів, що не діляться на жодне з чисел 2, 7, 5. 

Обмеження: усі елементи шуканої підпослідовності одночасно в пам'яті розміщені бути не можуть. 
%enditem


\newpage
8. Що буде надруковано за виконання?
%Оригинал был утерян. Требует отладки!
\begin{verbatim}
__d = 0

class D:
    __d = 1
    
    def __init__(self):
        self.__d = 2
        __d = 3
        D.__d = 4

obj = D()
D.__d = 5
try:
    print(__d, end=' ')
    print(obj.__d, end=' ')
    print(D.__d, end=' ')
    print(obj.__d, end=' ')
    print(obj._D__d)
except:
    print('ERR')
\end{verbatim}
%enditem


\newpage
9. Що буде надруковано за виконання?
\begin{verbatim}
class B0:
    def c(self): print(1, end=' ')
    def h(self): print(2, end=' ')

class B1(B0):
    def h(self): print(3, end=' ')

try:
    B1().c()
    B1().h()
except:
    print('ERR')
\end{verbatim}
%enditem


\newpage
10. Що буде надруковано за виконання?
\begin{verbatim}
class A:
    b = 1

    def __init__(self):
        a = 2
        a = 3

class B(A):
    b = 4
    c = 5

    def __init__(self):
        super().__init__()
        self.c = 6

try:
    obj = B()
    obj.c = obj.b + 10
    B.c = 13
except: print('EE', end=' ')

try: print(obj.c, end= ' ')
except: print('E1', end=' ')

try: print(A.c, end= ' ')
except: print('E2', end=' ')

try: print(B.c, end= ' ')
except: print('E3', end=' ')
\end{verbatim}
%enditem


\newpage
11. Змінні \kw{next} та \kw{prev} вузла списку позначають наступний та попередній за ним вузол відповідно. Починаючи з голови, у список записано послідовні цілі числа від~$-21$ до~$31$. Нехай змінна \kw{p} позначає вузол, що зберігає значення~$-9$.
Яке значення зберігається у вузлі \kw{p.next.prev.next.prev.prev} ?
%enditem


\newpage
12. 
\begin{center}
	\adjustimage{max size={0.8\linewidth}{0.8\paperheight}}
	{BinTree/trees_exp/25.png}
\end{center}
Указати послідовність міток вершин дерева, зображеного на рис., що утвориться в результаті обходу дерева в прямому порядку. Мітки записувати через пробіл.\par Алгоритм починає роботу з кореня (на рис. це верхня вершина).
%enditem


\newpage
13. 
\begin{center}
	\adjustimage{max size={0.9\linewidth}{0.9\paperheight}}
	{Graphs/AdjList/301.png}
\end{center}
Для графа, зображеного на рис., вказати список суміжності вершини 1.\par
Формат оформлення відповіді:\par
список суміжності виписувати через пробіл на одному рядку, якщо пробілів десь буде більше одного - не важливо, упорядкування вершин у списку також значення не має.
%enditem


\newpage
14. Для графа на рисунку вказати послідовність, в якій алгоритм пошуку в глибину буде відмічати відвідані вершини, якщо списки суміжності вершин переглядаються алгоритмом за зростанням міток вершин та вершини відмічаються на першому вход\-женні у вершину.\par
Пошук розпочинається з вершини 4.\par

\begin{center}
	\adjustimage{max size={0.9\linewidth}{0.9\paperheight}}
	{Graphs/Traversals/250.png}
\end{center}
%enditem


\newpage
15. Для графа на рисунку з вершини 6 запущено алгоритм пошуку в ширину, що будує кістякове дерево. Списки суміжності вершин переглядаються алгоритмом за зростанням міток вершин.\par
\begin{center}
	\adjustimage{max size={0.9\linewidth}{0.9\paperheight}}
	{Graphs/Traversals/140.png}
\end{center}
\par Які з наведених ребер входять до складу отриманого кістякового дерева?\par
1) (1, 2)\par
2) (0, 2)\par
3) (1, 5)\par
4) (1, 4)\par
5) (0, 3)\par
6) (0, 4)\par
{\vskip 15pt} Номери відповідних ребер записати за зростанням через пробіл.


%enditem


\newpage
16. 

Рядки журналу через двокрапку містять ім'я користувача (ідентифікатор), час події,тип події, код події, наприклад
\begin{verbatim}
s="username : 2022-05-05 13:26 : ERROR : 404"
\end{verbatim}
у коді 
\begin{verbatim}
res = re.fullmatch('???', s) 
s = ''
code_str = ??? 
\end{verbatim}
чим треба замінити кожен з ???, щоб значенням code\_str був код події? 


\vspace{10mm}
\textbf{Завдання 17} 

Задано програму, що містить код 1-2 класів, можливе успадкування, та клієнтський код.
У наведеному коді невеликі частини закриті. 
Також наведено вихід програми.
Необхідно відновити закриті частини, щоб отриманий код відповідав наведеному виходу програми.


\textbf{Завдання 18-19} 

Кілька запитань за завданнями для самостійної роботи.
\vspace{10mm}
Скоріш за все, що завдань буде трохи менше, а деякі попередні будуть переформульовані в стилі двох останніх.

%enditem

\end {large}}

%end

\newpage
%begin

{{
\begin {large}

\newpage
1. Що буде надруковано за виконання?

\begin{verbatim}
print('R', end=' ')
h=('a','b','c','d','e',0,1,2,3,4)
print(h[-2])
\end{verbatim}

%enditem


\newpage
2. Що буде надруковано за виконання?

\begin{verbatim}
print('R', end=' ')
c='0123456789'
print(c[-6:10:3])
\end{verbatim}

%enditem



\newpage
3. Що буде надруковано за виконання?

\begin{verbatim}
print('R', end=' ')
b=[10,11,12,13,14]
b.append('13')
print(b)
\end{verbatim}

%enditem



\newpage
4. Що буде надруковано за виконання?

\begin{verbatim}
print('R', end=' ')
g=[10,11,12,13,14]
g[2:5]=[2,3]
print(g)
\end{verbatim}

%enditem



\newpage
5. Укажіть ТИП значення виразу.\par
\begin{verbatim}
{x*x:x for x in range(-7,7)}
\end{verbatim}
%enditem


\newpage
6. Що буде надруковано за виконання?\par
\begin{verbatim}
try:
    s = {31:5, 20:6, 10:8, 29:6, 10:8}
    s[72] = 0
    for x, y in s.items():
        print(x, y, sep=' ', end=' ')
except:
    print('ERR')
\end{verbatim}
%enditem


\newpage
7. Нехай змінна \kw{s} позначає послідовність цілих, по якій можна ітерувати.

Написати вираз, що перевіряє, чи правда, що всі елементи послідовності менші за задане число num.
%enditem


\newpage
8. Що буде надруковано за виконання?
\begin{verbatim}
a = 0
j = 1

class D:
    a = 2
    
    def __init__(self):
        global a
        self.a = 3
        a = 4
        D.a = 5

obj = D()
try:
    print(a, end=' ')
    print(j, end=' ')
    print(D.a, end=' ')
    print(obj.a)
except:
    print('ERR')
\end{verbatim}
%enditem


\newpage
9. Що буде надруковано за виконання?
\begin{verbatim}
class A0:
    def a(self): print(1, end=' ')
    def f(self): print(2, end=' ')

class A1(A0):
    def a(self): print(3, end=' ')

try:
    A1().f()
    A1().a()
except:
    print('ERR')
\end{verbatim}
%enditem


\newpage
10. Що буде надруковано за виконання?
\begin{verbatim}
class A:
    c = 1

    def __init__(self):
        a = 2
        c = 3

class B(A):
    a = 4
    b = 5

    def __init__(self):
        super().__init__()
        self.b = 6

try:
    obj = B()
    obj.a = obj.b + 10
    B.a = 13
except: print('EE', end=' ')

try: print(obj.a, end= ' ')
except: print('E1', end=' ')

try: print(A.a, end= ' ')
except: print('E2', end=' ')

try: print(B.a, end= ' ')
except: print('E3', end=' ')
\end{verbatim}
%enditem


\newpage
11. Змінні \kw{next} та \kw{prev} вузла списку позначають наступний та попередній за ним вузол відповідно. Починаючи з голови, у список записано послідовні цілі числа від~$-18$ до~$35$. Нехай змінна \kw{p} позначає вузол, що зберігає значення~$8$.
Яке значення зберігається у вузлі \kw{p.prev.next.prev.next.next} ?
%enditem


\newpage
12. 
\begin{center}
	\adjustimage{max size={0.8\linewidth}{0.8\paperheight}}
	{BinTree/trees_exp/131.png}
\end{center}
Указати послідовність міток вершин дерева, зображеного на рис., що утвориться в результаті обходу дерева в оберненому порядку. Мітки записувати через пробіл.\par Алгоритм починає роботу з кореня (на рис. це верхня вершина).
%enditem


\newpage
13. 
\begin{center}
	\adjustimage{max size={0.9\linewidth}{0.9\paperheight}}
	{Graphs/AdjList/70.png}
\end{center}
Для графа, зображеного на рис., вказати список суміжності вершини 6.\par
Формат оформлення відповіді:\par
список суміжності виписувати через пробіл на одному рядку, якщо пробілів десь буде більше одного - не важливо, упорядкування вершин у списку також значення не має.
%enditem


\newpage
14. Для графа на рисунку з вершини 6 запущено алгоритм пошуку в глибину, що будує кістякове дерево. Списки суміжності вершин переглядаються алгоритмом за зростанням міток вершин.\par
\begin{center}
	\adjustimage{max size={0.9\linewidth}{0.9\paperheight}}
	{Graphs/Traversals/127.png}
\end{center}
\par Які з наведених ребер входять до складу отриманого кістякового дерева?\par
1) (0, 1)\par
2) (0, 3)\par
3) (1, 5)\par
4) (0, 6)\par
5) (2, 5)\par
6) (1, 3)\par
{\vskip 15pt} Номери відповідних ребер записати за зростанням через пробіл.

%enditem


\newpage
15. Для графа на рисунку з вершини 4 запущено алгоритм пошуку в ширину, що будує кістякове дерево. Списки суміжності вершин переглядаються алгоритмом за зростанням міток вершин.\par
\begin{center}
	\adjustimage{max size={0.9\linewidth}{0.9\paperheight}}
	{Graphs/Traversals/326.png}
\end{center}
\par Які з наведених ребер входять до складу отриманого кістякового дерева?\par
1) (1, 6)\par
2) (2, 3)\par
3) (0, 1)\par
4) (3, 5)\par
5) (4, 5)\par
6) (0, 6)\par
{\vskip 15pt} Номери відповідних ребер записати за зростанням через пробіл.


%enditem


\newpage
16. 

Рядки журналу через двокрапку містять ім'я користувача (ідентифікатор), час події,тип події, код події, наприклад
\begin{verbatim}
s="username : 2022-05-05 13:26 : ERROR : 404"
\end{verbatim}
у коді 
\begin{verbatim}
res = re.fullmatch('???', s) 
s = ''
code_str = ??? 
\end{verbatim}
чим треба замінити кожен з ???, щоб значенням code\_str був код події? 


\vspace{10mm}
\textbf{Завдання 17} 

Задано програму, що містить код 1-2 класів, можливе успадкування, та клієнтський код.
У наведеному коді невеликі частини закриті. 
Також наведено вихід програми.
Необхідно відновити закриті частини, щоб отриманий код відповідав наведеному виходу програми.


\textbf{Завдання 18-19} 

Кілька запитань за завданнями для самостійної роботи.
\vspace{10mm}
Скоріш за все, що завдань буде трохи менше, а деякі попередні будуть переформульовані в стилі двох останніх.

%enditem

\end {large}}

%end

\newpage
%begin

{{
\begin {large}

\newpage
1. Що буде надруковано за виконання?

\begin{verbatim}
print('R', end=' ')
c=('a','b','c','d','e',0,1,2,4)
print(c[5])
\end{verbatim}

%enditem


\newpage
2. Що буде надруковано за виконання?

\begin{verbatim}
print('R', end=' ')
g=('0','1','2','3','4','5','6','7','8','9')
print(g[1::-3])
\end{verbatim}

%enditem



\newpage
3. Що буде надруковано за виконання?

\begin{verbatim}
print('R', end=' ')
a=['0','1','2','3','4']
a+=16
print(a)
\end{verbatim}

%enditem



\newpage
4. Що буде надруковано за виконання?

\begin{verbatim}
print('R', end=' ')
b=[10,11,12,13,14]
b[:-9]=[1,0]
print(b)
\end{verbatim}

%enditem



\newpage
5. Укажіть ТИП значення виразу.\par
\begin{verbatim}
[7,2,3,4]
\end{verbatim}
%enditem


\newpage
6. Що буде надруковано за виконання?\par
\begin{verbatim}
try:
    d = {23:3, 71:2, 17:3, 23:0}
    d[18] = 0
    for x in d.values():
        print(x, sep=' ', end=' ')
except:
    print('ERR')
\end{verbatim}
%enditem


\newpage
7. Нехай змінна \kw{s} позначає послідовність цілих, по якій можна ітерувати. 

Написати вираз, що перевіряє, чи правда, що послідовність  містить хоча б одне додатнє число.
%enditem


\newpage
8. Що буде надруковано за виконання?
\begin{verbatim}
d = 0
h = 1

class B:
    h = 2
    
    def __init__(self):
        self.h = 3
        h = 4
        B.h = 5

obj = B()
try:
    print(d, end=' ')
    print(h, end=' ')
    print(B.h, end=' ')
    print(obj.h)
except:
    print('ERR')
\end{verbatim}
%enditem


\newpage
9. Що буде надруковано за виконання?
\begin{verbatim}
class C0:
    def d(self): print(1, end=' ')
    def h(self): print(2, end=' ')

class C1(C0):
    def h(self): print(3, end=' ')

try:
    C1().d()
    C0().h()
except:
    print('ERR')
\end{verbatim}
%enditem


\newpage
10. Що буде надруковано за виконання?
\begin{verbatim}
class A:
    a = 1

    def __init__(self):
        b = 2
        b = 3

class B(A):
    b = 4
    c = 5

    def __init__(self):
        super().__init__()
        self.c = 6

try:
    obj = B()
    obj.c = obj.b + 10
    B.c = 13
except: print('EE', end=' ')

try: print(obj.c, end= ' ')
except: print('E1', end=' ')

try: print(A.c, end= ' ')
except: print('E2', end=' ')

try: print(B.c, end= ' ')
except: print('E3', end=' ')
\end{verbatim}
%enditem


\newpage
11. Змінні \kw{next} та \kw{prev} вузла списку позначають наступний та попередній за ним вузол відповідно. Починаючи з голови, у список записано послідовні цілі числа від~$-36$ до~$41$. Нехай змінна \kw{p} позначає вузол, що зберігає значення~$7$.
Яке значення зберігається у вузлі \kw{p.next.prev.next.prev.prev} ?
%enditem


\newpage
12. 
\begin{center}
	\adjustimage{max size={0.8\linewidth}{0.8\paperheight}}
	{BinTree/trees_exp/156.png}
\end{center}
Указати послідовність міток вершин дерева, зображеного на рис., що утвориться в результаті обходу дерева в прямому порядку. Мітки записувати через пробіл.\par Алгоритм починає роботу з кореня (на рис. це верхня вершина).
%enditem


\newpage
13. 
\begin{center}
	\adjustimage{max size={0.9\linewidth}{0.9\paperheight}}
	{Graphs/AdjList/321.png}
\end{center}
Для графа, зображеного на рис., вказати список суміжності вершини 7.\par
Формат оформлення відповіді:\par
список суміжності виписувати через пробіл на одному рядку, якщо пробілів десь буде більше одного - не важливо, упорядкування вершин у списку також значення не має.
%enditem


\newpage
14. Для графа на рисунку вказати послідовність, в якій алгоритм пошуку в глибину буде відмічати відвідані вершини, якщо списки суміжності вершин переглядаються алгоритмом за зростанням міток вершин та вершини відмічаються на першому вход\-женні у вершину.\par
Пошук розпочинається з вершини 3.\par

\begin{center}
	\adjustimage{max size={0.9\linewidth}{0.9\paperheight}}
	{Graphs/Traversals/15.png}
\end{center}
%enditem


\newpage
15. Для графа на рисунку з вершини 0 запущено алгоритм пошуку в ширину, що будує кістякове дерево. Списки суміжності вершин переглядаються алгоритмом за зростанням міток вершин.\par
\begin{center}
	\adjustimage{max size={0.9\linewidth}{0.9\paperheight}}
	{Graphs/Traversals/78.png}
\end{center}
\par Які з наведених ребер входять до складу отриманого кістякового дерева?\par
1) (1, 2)\par
2) (3, 4)\par
3) (0, 5)\par
4) (4, 6)\par
5) (1, 4)\par
6) (2, 5)\par
{\vskip 15pt} Номери відповідних ребер записати за зростанням через пробіл.


%enditem


\newpage
16. 

Рядки журналу через двокрапку містять ім'я користувача (ідентифікатор), час події,тип події, код події, наприклад
\begin{verbatim}
s="username : 2022-05-05 13:26 : ERROR : 404"
\end{verbatim}
у коді 
\begin{verbatim}
res = re.fullmatch('???', s) 
s = ''
code_str = ??? 
\end{verbatim}
чим треба замінити кожен з ???, щоб значенням code\_str був код події? 


\vspace{10mm}
\textbf{Завдання 17} 

Задано програму, що містить код 1-2 класів, можливе успадкування, та клієнтський код.
У наведеному коді невеликі частини закриті. 
Також наведено вихід програми.
Необхідно відновити закриті частини, щоб отриманий код відповідав наведеному виходу програми.


\textbf{Завдання 18-19} 

Кілька запитань за завданнями для самостійної роботи.
\vspace{10mm}
Скоріш за все, що завдань буде трохи менше, а деякі попередні будуть переформульовані в стилі двох останніх.

%enditem

\end {large}}

%end

\newpage
%begin

{{
\begin {large}

\newpage
1. Що буде надруковано за виконання?

\begin{verbatim}
print('R', end=' ')
b=('a','b','c','d','e',0,1,2,3)
print(b[-8])
\end{verbatim}

%enditem


\newpage
2. Що буде надруковано за виконання?

\begin{verbatim}
print('R', end=' ')
f=[0,1,2,3,4,5,6,7,8,9]
print(f[-14::2])
\end{verbatim}

%enditem



\newpage
3. Що буде надруковано за виконання?

\begin{verbatim}
print('R', end=' ')
h=[10,11,12,13,14]
print(h.index(11), end=' ')
\end{verbatim}

%enditem



\newpage
4. Що буде надруковано за виконання?

\begin{verbatim}
print('R', end=' ')
f=['0','1','2','3','4']
del f[-9:-8]
print(f)
\end{verbatim}

%enditem



\newpage
5. Укажіть ТИП значення виразу.\par
\begin{verbatim}
{y:13 for y in range(-7,7)}
\end{verbatim}
%enditem


\newpage
6. Що буде надруковано за виконання?\par
\begin{verbatim}
try:
    t = {77:0, 65:5, 30:1, 50:8}
    t[77] = 8
    for x in t.keys():
        print(x, sep=' ', end=' ')
except:
    print('ERR')
\end{verbatim}
%enditem


\newpage
7. 
Записати ВИРАЗ, значенням якого є список, що складається з цілих чисел від 14 до 2010 (включно), причому спочатку за зростанням записано всі, що діляться на 5, а потім решту за спаданням.

%enditem


\newpage
8. Що буде надруковано за виконання?
\begin{verbatim}
c = 0
n = 1

class C:
    c = 2
    
    def __init__(self):
        global c
        self.c = 3
        c = 4
        C.c = 5

obj = C()
try:
    print(c, end=' ')
    print(n, end=' ')
    print(C.c, end=' ')
    print(obj.c)
except:
    print('ERR')
\end{verbatim}
%enditem


\newpage
9. Що буде надруковано за виконання?
\begin{verbatim}
class D0:
    def b(self): print(1, end=' ')
    def g(self): print(2, end=' ')

class D1(D0):
    def g(self): print(3, end=' ')

try:
    D0().b()
    D1().g()
except:
    print('ERR')
\end{verbatim}
%enditem


\newpage
10. Що буде надруковано за виконання?
\begin{verbatim}
class A:
    c = 1

    def __init__(self):
        b = 2
        c = 3

class B(A):
    c = 4
    a = 5

    def __init__(self):
        super().__init__()
        self.a = 6

try:
    obj = B()
    obj.c = obj.a + 10
    B.c = 13
except: print('EE', end=' ')

try: print(obj.c, end= ' ')
except: print('E1', end=' ')

try: print(A.c, end= ' ')
except: print('E2', end=' ')

try: print(B.c, end= ' ')
except: print('E3', end=' ')
\end{verbatim}
%enditem


\newpage
11. Змінні \kw{next} та \kw{prev} вузла списку позначають наступний та попередній за ним вузол відповідно. Починаючи з голови, у список записано послідовні цілі числа від~$-54$ до~$23$. Нехай змінна \kw{p} позначає вузол, що зберігає значення~$1$.
Яке значення зберігається у вузлі \kw{p.prev.next.prev.next.next} ?
%enditem


\newpage
12. 
\begin{center}
	\adjustimage{max size={0.8\linewidth}{0.8\paperheight}}
	{BinTree/trees_exp/143.png}
\end{center}
Указати послідовність міток вершин дерева, зображеного на рис., що утвориться в результаті обходу дерева в оберненому порядку. Мітки записувати через пробіл.\par Алгоритм починає роботу з кореня (на рис. це верхня вершина).
%enditem


\newpage
13. 
\begin{center}
	\adjustimage{max size={0.9\linewidth}{0.9\paperheight}}
	{Graphs/AdjList/48.png}
\end{center}
Для графа, зображеного на рис., вказати список суміжності вершини 5.\par
Формат оформлення відповіді:\par
список суміжності виписувати через пробіл на одному рядку, якщо пробілів десь буде більше одного - не важливо, упорядкування вершин у списку також значення не має.
%enditem


\newpage
14. Для графа на рисунку вказати послідовність, в якій алгоритм пошуку в глибину буде відмічати відвідані вершини, якщо списки суміжності вершин переглядаються алгоритмом за зростанням міток вершин та вершини відмічаються на першому вход\-женні у вершину.\par
Пошук розпочинається з вершини 2.\par

\begin{center}
	\adjustimage{max size={0.9\linewidth}{0.9\paperheight}}
	{Graphs/Traversals/103.png}
\end{center}
%enditem


\newpage
15. Для графа на рисунку з вершини 3 запущено алгоритм пошуку в ширину, що будує кістякове дерево. Списки суміжності вершин переглядаються алгоритмом за зростанням міток вершин.\par
\begin{center}
	\adjustimage{max size={0.9\linewidth}{0.9\paperheight}}
	{Graphs/Traversals/145.png}
\end{center}
\par Які з наведених ребер входять до складу отриманого кістякового дерева?\par
1) (0, 2)\par
2) (1, 3)\par
3) (1, 6)\par
4) (4, 5)\par
5) (0, 5)\par
6) (0, 3)\par
{\vskip 15pt} Номери відповідних ребер записати за зростанням через пробіл.


%enditem


\newpage
16. 

Рядки журналу через двокрапку містять ім'я користувача (ідентифікатор), час події,тип події, код події, наприклад
\begin{verbatim}
s="username : 2022-05-05 13:26 : ERROR : 404"
\end{verbatim}
у коді 
\begin{verbatim}
res = re.fullmatch('???', s) 
s = ''
code_str = ??? 
\end{verbatim}
чим треба замінити кожен з ???, щоб значенням code\_str був код події? 


\vspace{10mm}
\textbf{Завдання 17} 

Задано програму, що містить код 1-2 класів, можливе успадкування, та клієнтський код.
У наведеному коді невеликі частини закриті. 
Також наведено вихід програми.
Необхідно відновити закриті частини, щоб отриманий код відповідав наведеному виходу програми.


\textbf{Завдання 18-19} 

Кілька запитань за завданнями для самостійної роботи.
\vspace{10mm}
Скоріш за все, що завдань буде трохи менше, а деякі попередні будуть переформульовані в стилі двох останніх.

%enditem

\end {large}}

%end

\newpage
%begin

{{
\begin {large}

\newpage
1. Що буде надруковано за виконання?

\begin{verbatim}
print('R', end=' ')
f=('a','b','c','d','e',0,1,2,3,4)
print(f[-4])
\end{verbatim}

%enditem


\newpage
2. Що буде надруковано за виконання?

\begin{verbatim}
print('R', end=' ')
d='0123456789'
print(d[9:11:-2])
\end{verbatim}

%enditem



\newpage
3. Що буде надруковано за виконання?

\begin{verbatim}
print('R', end=' ')
f=['0','1','2','3','4']
f.remove('2')
print(f)
\end{verbatim}

%enditem



\newpage
4. Що буде надруковано за виконання?

\begin{verbatim}
print('R', end=' ')
c=[10,11,12,13,14]
c[-6:6]='13'
print(c)
\end{verbatim}

%enditem



\newpage
5. Укажіть ТИП значення виразу.\par
\begin{verbatim}
{x:x for x in range(7)}
\end{verbatim}
%enditem


\newpage
6. Що буде надруковано за виконання?\par
\begin{verbatim}
try:
    t = {58:0, 35:2, 76:9, 24:6, 24:5}
    t[55] = 2
    for x in t :
        print(x, sep=' ', end=' ')
except:
    print('ERR')
\end{verbatim}
%enditem


\newpage
7. Нехай змінна \kw{s} позначає послідовність цілих, по якій можна ітерувати. 

Написати вираз, що перевіряє, чи правда, що послідовність  містить число 35.
%enditem


\newpage
8. Що буде надруковано за виконання?
\begin{verbatim}
b = 0
j = 1

class A:
    j = 2
    
    def __init__(self):
        global j
        self.j = 3
        j = 4
        A.j = 5

obj = A()
try:
    print(b, end=' ')
    print(j, end=' ')
    print(A.j, end=' ')
    print(obj.j)
except:
    print('ERR')
\end{verbatim}
%enditem


\newpage
9. Що буде надруковано за виконання?
\begin{verbatim}
class A0:
    def d(self): print(1, end=' ')
    def f(self): print(2, end=' ')

class A1(A0):
    def f(self): print(3, end=' ')

try:
    A1().d()
    A1().f()
except:
    print('ERR')
\end{verbatim}
%enditem


\newpage
10. Що буде надруковано за виконання?
\begin{verbatim}
class A:
    c = 1

    def _f1(self): print(2, end=' ')
    def _f2(self): print(3, end=' ')

    def main(self):
        try: print(self.c, end=' ')
        except: print('E1', end=' ')

        try: self._f1()
        except: print('E2', end=' ')

        try: self._f2()
        except: print('E3', end=' ')


class B(A):
    c = 4

    def _f1(self): print(5, end=' ')
    def _f2(self): print(6, end=' ')

try: B().main()
except: print('E4', end=' ')
\end{verbatim}
%enditem


\newpage
11. Змінні \kw{next} та \kw{prev} вузла списку позначають наступний та попередній за ним вузол відповідно. Починаючи з голови, у список записано послідовні цілі числа від~$-53$ до~$30$. Нехай змінна \kw{p} позначає вузол, що зберігає значення~$-7$.
Яке значення зберігається у вузлі \kw{p.prev.prev.prev.next.next} ?
%enditem


\newpage
12. 
\begin{center}
	\adjustimage{max size={0.8\linewidth}{0.8\paperheight}}
	{BinTree/trees_exp/77.png}
\end{center}
Указати послідовність міток вершин дерева, зображеного на рис., що утвориться в результаті обходу дерева в оберненому порядку. Мітки записувати через пробіл.\par Алгоритм починає роботу з кореня (на рис. це верхня вершина).
%enditem


\newpage
13. 
\begin{center}
	\adjustimage{max size={0.9\linewidth}{0.9\paperheight}}
	{Graphs/AdjList/204.png}
\end{center}
Для графа, зображеного на рис., вказати список суміжності вершини 3.\par
Формат оформлення відповіді:\par
список суміжності виписувати через пробіл на одному рядку, якщо пробілів десь буде більше одного - не важливо, упорядкування вершин у списку також значення не має.
%enditem


\newpage
14. Для графа на рисунку вказати послідовність, в якій алгоритм пошуку в глибину буде відмічати відвідані вершини, якщо списки суміжності вершин переглядаються алгоритмом за зростанням міток вершин та вершини відмічаються на першому вход\-женні у вершину.\par
Пошук розпочинається з вершини 3.\par

\begin{center}
	\adjustimage{max size={0.9\linewidth}{0.9\paperheight}}
	{Graphs/Traversals/52.png}
\end{center}
%enditem


\newpage
15. Для графа на рисунку з вершини 6 запущено алгоритм пошуку в ширину, що будує кістякове дерево. Списки суміжності вершин переглядаються алгоритмом за зростанням міток вершин.\par
\begin{center}
	\adjustimage{max size={0.9\linewidth}{0.9\paperheight}}
	{Graphs/Traversals/143.png}
\end{center}
\par Які з наведених ребер входять до складу отриманого кістякового дерева?\par
1) (4, 6)\par
2) (5, 6)\par
3) (2, 4)\par
4) (3, 4)\par
5) (0, 2)\par
6) (0, 1)\par
{\vskip 15pt} Номери відповідних ребер записати за зростанням через пробіл.


%enditem


\newpage
16. 

Рядки журналу через двокрапку містять ім'я користувача (ідентифікатор), час події,тип події, код події, наприклад
\begin{verbatim}
s="username : 2022-05-05 13:26 : ERROR : 404"
\end{verbatim}
у коді 
\begin{verbatim}
res = re.fullmatch('???', s) 
s = ''
code_str = ??? 
\end{verbatim}
чим треба замінити кожен з ???, щоб значенням code\_str був код події? 


\vspace{10mm}
\textbf{Завдання 17} 

Задано програму, що містить код 1-2 класів, можливе успадкування, та клієнтський код.
У наведеному коді невеликі частини закриті. 
Також наведено вихід програми.
Необхідно відновити закриті частини, щоб отриманий код відповідав наведеному виходу програми.


\textbf{Завдання 18-19} 

Кілька запитань за завданнями для самостійної роботи.
\vspace{10mm}
Скоріш за все, що завдань буде трохи менше, а деякі попередні будуть переформульовані в стилі двох останніх.

%enditem

\end {large}}

%end

\newpage
%begin

{{
\begin {large}

\newpage
1. Що буде надруковано за виконання?

\begin{verbatim}
print('R', end=' ')
d=('a','b','c','d','e',0,1,2,3)
print(d[14])
\end{verbatim}

%enditem


\newpage
2. Що буде надруковано за виконання?

\begin{verbatim}
print('R', end=' ')
e=('0','1','2','3','4','5','6','7','8','9')
print(e[12:6:2])
\end{verbatim}

%enditem



\newpage
3. Що буде надруковано за виконання?

\begin{verbatim}
print('R', end=' ')
g=[10,11,12,13,14]
print(g.pop(-3), end=' ')
print(g)
\end{verbatim}

%enditem



\newpage
4. Що буде надруковано за виконання?

\begin{verbatim}
print('R', end=' ')
e=['0','1','2','3','4']
del e[-8:]
print(e)
\end{verbatim}

%enditem



\newpage
5. Укажіть ТИП значення виразу.\par
\begin{verbatim}
{(1,4), 8}
\end{verbatim}
%enditem


\newpage
6. Що буде надруковано за виконання?\par
\begin{verbatim}
try:
    d = {63:3, 57:0, 65:0, 46:3}
    d[18] = 4
    for x in d.items():
        print(x, sep=' ', end=' ')
except:
    print('ERR')
\end{verbatim}
%enditem


\newpage
7. Нехай змінна \kw{s} позначає послідовність цілих, по якій можна ітерувати. 

Написати вираз, що перевіряє, чи правда, що послідовність  містить хоча б одне додатнє число.
%enditem


\newpage
8. Що буде надруковано за виконання?
\begin{verbatim}
a = 0
k = 1

class D:
    a = 2
    
    def __init__(self):
        self.a = 3
        a = 4
        D.a = 5

obj = D()
try:
    print(a, end=' ')
    print(k, end=' ')
    print(D.a, end=' ')
    print(obj.a)
except:
    print('ERR')
\end{verbatim}
%enditem


\newpage
9. Що буде надруковано за виконання?
\begin{verbatim}
class D0:
    def a(self): print(1, end=' ')
    def h(self): print(2, end=' ')

class D1(D0):
    def h(self): print(3, end=' ')

try:
    D1().a()
    D1().h()
except:
    print('ERR')
\end{verbatim}
%enditem


\newpage
10. Що буде надруковано за виконання?
\begin{verbatim}
class A:
    a = 1

    def __init__(self):
        c = 2
        a = 3

class B(A):
    c = 4
    b = 5

    def __init__(self):
        super().__init__()
        self.b = 6

try:
    obj = B()
    obj.b = obj.c + 10
    B.b = 13
except: print('EE', end=' ')

try: print(obj.b, end= ' ')
except: print('E1', end=' ')

try: print(A.b, end= ' ')
except: print('E2', end=' ')

try: print(B.b, end= ' ')
except: print('E3', end=' ')
\end{verbatim}
%enditem


\newpage
11. Змінні \kw{next} та \kw{prev} вузла списку позначають наступний та попередній за ним вузол відповідно. Починаючи з голови, у список записано послідовні цілі числа від~$-48$ до~$37$. Нехай змінна \kw{p} позначає вузол, що зберігає значення~$5$.
Яке значення зберігається у вузлі \kw{p.next.next.next.prev.prev} ?
%enditem


\newpage
12. 
\begin{center}
	\adjustimage{max size={0.8\linewidth}{0.8\paperheight}}
	{BinTree/trees_exp/80.png}
\end{center}
Указати послідовність міток вершин дерева, зображеного на рис., що утвориться в результаті обходу дерева в прямому порядку. Мітки записувати через пробіл.\par Алгоритм починає роботу з кореня (на рис. це верхня вершина).
%enditem


\newpage
13. 
\begin{center}
	\adjustimage{max size={0.9\linewidth}{0.9\paperheight}}
	{Graphs/AdjList/347.png}
\end{center}
Для графа, зображеного на рис., вказати список суміжності вершини 5.\par
Формат оформлення відповіді:\par
список суміжності виписувати через пробіл на одному рядку, якщо пробілів десь буде більше одного - не важливо, упорядкування вершин у списку також значення не має.
%enditem


\newpage
14. Для графа на рисунку з вершини 2 запущено алгоритм пошуку в глибину, що будує кістякове дерево. Списки суміжності вершин переглядаються алгоритмом за зростанням міток вершин.\par
\begin{center}
	\adjustimage{max size={0.9\linewidth}{0.9\paperheight}}
	{Graphs/Traversals/4.png}
\end{center}
\par Які з наведених ребер входять до складу отриманого кістякового дерева?\par
1) (0, 4)\par
2) (4, 5)\par
3) (0, 2)\par
4) (1, 2)\par
5) (1, 3)\par
6) (4, 6)\par
{\vskip 15pt} Номери відповідних ребер записати за зростанням через пробіл.

%enditem


\newpage
15. Для графа на рисунку з вершини 2 запущено алгоритм пошуку в ширину, що будує кістякове дерево. Списки суміжності вершин переглядаються алгоритмом за зростанням міток вершин.\par
\begin{center}
	\adjustimage{max size={0.9\linewidth}{0.9\paperheight}}
	{Graphs/Traversals/25.png}
\end{center}
\par Які з наведених ребер входять до складу отриманого кістякового дерева?\par
1) (2, 3)\par
2) (2, 4)\par
3) (5, 6)\par
4) (2, 5)\par
5) (0, 4)\par
6) (3, 4)\par
{\vskip 15pt} Номери відповідних ребер записати за зростанням через пробіл.


%enditem


\newpage
16. 

Рядки журналу через двокрапку містять ім'я користувача (ідентифікатор), час події,тип події, код події, наприклад
\begin{verbatim}
s="username : 2022-05-05 13:26 : ERROR : 404"
\end{verbatim}
у коді 
\begin{verbatim}
res = re.fullmatch('???', s) 
s = ''
code_str = ??? 
\end{verbatim}
чим треба замінити кожен з ???, щоб значенням code\_str був код події? 


\vspace{10mm}
\textbf{Завдання 17} 

Задано програму, що містить код 1-2 класів, можливе успадкування, та клієнтський код.
У наведеному коді невеликі частини закриті. 
Також наведено вихід програми.
Необхідно відновити закриті частини, щоб отриманий код відповідав наведеному виходу програми.


\textbf{Завдання 18-19} 

Кілька запитань за завданнями для самостійної роботи.
\vspace{10mm}
Скоріш за все, що завдань буде трохи менше, а деякі попередні будуть переформульовані в стилі двох останніх.

%enditem

\end {large}}

%end

\newpage
%begin

{{
\begin {large}

\newpage
1. Що буде надруковано за виконання?

\begin{verbatim}
print('R', end=' ')
a=('a','b','c','d','e',0,1)
print(a[-11])
\end{verbatim}

%enditem


\newpage
2. Що буде надруковано за виконання?

\begin{verbatim}
print('R', end=' ')
b=['0','1','2','3','4','5','6','7','8','9']
print(b[:-7:3])
\end{verbatim}

%enditem



\newpage
3. Що буде надруковано за виконання?

\begin{verbatim}
print('R', end=' ')
e=['0','1','2','3','4']
e.insert(7, [4,5])
print(e)
\end{verbatim}

%enditem



\newpage
4. Що буде надруковано за виконання?

\begin{verbatim}
print('R', end=' ')
d=[10,11,12,13,14]
d[6:]=[4,5]
print(d)
\end{verbatim}

%enditem



\newpage
5. Укажіть ТИП значення виразу.\par
\begin{verbatim}
{y:None for y in range(-8,3)}
\end{verbatim}
%enditem


\newpage
6. Що буде надруковано за виконання?\par
\begin{verbatim}
try:
    s = {39:9, 55:3, 30:4, 70:0, 55:9}
    s[58] = 3
    for x in s.items():
        print(*x, sep=' ', end=' ')
except:
    print('ERR')
\end{verbatim}
%enditem


\newpage
7. Нехай змінна \kw{s} позначає послідовність цілих, по якій можна ітерувати.

Написати вираз, що перевіряє, чи правда, що всі елементи послідовності не діляться на 8.
%enditem


\newpage
8. Що буде надруковано за виконання?
\begin{verbatim}
a = 0
h = 1

class D:
    h = 2
    
    def __init__(self):
        global h
        self.h = 3
        h = 4
        D.h = 5

obj = D()
try:
    print(a, end=' ')
    print(h, end=' ')
    print(D.h, end=' ')
    print(obj.h)
except:
    print('ERR')
\end{verbatim}
%enditem


\newpage
9. Що буде надруковано за виконання?
\begin{verbatim}
class C0:
    def b(self): print(1, end=' ')
    def g(self): print(2, end=' ')

class C1(C0):
    def g(self): print(3, end=' ')

try:
    C0().b()
    C1().g()
except:
    print('ERR')
\end{verbatim}
%enditem


\newpage
10. Що буде надруковано за виконання?
\begin{verbatim}
class A:
    b = 1

    def g1(self): print(2, end=' ')
    def _g2(self): print(3, end=' ')

    def main(self):
        try: print(self.b, end=' ')
        except: print('E1', end=' ')

        try: self.g1()
        except: print('E2', end=' ')

        try: self._g2()
        except: print('E3', end=' ')


class B(A):
    b = 4

    def g1(self): print(5, end=' ')
    def _g2(self): print(6, end=' ')

try: B().main()
except: print('E4', end=' ')
\end{verbatim}
%enditem


\newpage
11. Змінні \kw{next} та \kw{prev} вузла списку позначають наступний та попередній за ним вузол відповідно. Починаючи з голови, у список записано послідовні цілі числа від~$-65$ до~$15$. Нехай змінна \kw{p} позначає вузол, що зберігає значення~$-3$.
Яке значення зберігається у вузлі \kw{p.next.prev.next.next.next} ?
%enditem


\newpage
12. 
\begin{center}
	\adjustimage{max size={0.8\linewidth}{0.8\paperheight}}
	{BinTree/trees_exp/126.png}
\end{center}
Указати послідовність міток вершин дерева, зображеного на рис., що утвориться в результаті обходу дерева в прямому порядку. Мітки записувати через пробіл.\par Алгоритм починає роботу з кореня (на рис. це верхня вершина).
%enditem


\newpage
13. 
\begin{center}
	\adjustimage{max size={0.9\linewidth}{0.9\paperheight}}
	{Graphs/AdjList/119.png}
\end{center}
Для графа, зображеного на рис., вказати список суміжності вершини 7.\par
Формат оформлення відповіді:\par
список суміжності виписувати через пробіл на одному рядку, якщо пробілів десь буде більше одного - не важливо, упорядкування вершин у списку також значення не має.
%enditem


\newpage
14. Для графа на рисунку вказати послідовність, в якій алгоритм пошуку в глибину буде відмічати відвідані вершини, якщо списки суміжності вершин переглядаються алгоритмом за зростанням міток вершин та вершини відмічаються на першому вход\-женні у вершину.\par
Пошук розпочинається з вершини 2.\par

\begin{center}
	\adjustimage{max size={0.9\linewidth}{0.9\paperheight}}
	{Graphs/Traversals/208.png}
\end{center}
%enditem


\newpage
15. Для графа на рисунку з вершини 2 запущено алгоритм пошуку в ширину, що будує кістякове дерево. Списки суміжності вершин переглядаються алгоритмом за зростанням міток вершин.\par
\begin{center}
	\adjustimage{max size={0.9\linewidth}{0.9\paperheight}}
	{Graphs/Traversals/308.png}
\end{center}
\par Які з наведених ребер входять до складу отриманого кістякового дерева?\par
1) (2, 6)\par
2) (3, 5)\par
3) (3, 6)\par
4) (2, 5)\par
5) (2, 4)\par
6) (1, 2)\par
{\vskip 15pt} Номери відповідних ребер записати за зростанням через пробіл.


%enditem


\newpage
16. 

Рядки журналу через двокрапку містять ім'я користувача (ідентифікатор), час події,тип події, код події, наприклад
\begin{verbatim}
s="username : 2022-05-05 13:26 : ERROR : 404"
\end{verbatim}
у коді 
\begin{verbatim}
res = re.fullmatch('???', s) 
s = ''
code_str = ??? 
\end{verbatim}
чим треба замінити кожен з ???, щоб значенням code\_str був код події? 


\vspace{10mm}
\textbf{Завдання 17} 

Задано програму, що містить код 1-2 класів, можливе успадкування, та клієнтський код.
У наведеному коді невеликі частини закриті. 
Також наведено вихід програми.
Необхідно відновити закриті частини, щоб отриманий код відповідав наведеному виходу програми.


\textbf{Завдання 18-19} 

Кілька запитань за завданнями для самостійної роботи.
\vspace{10mm}
Скоріш за все, що завдань буде трохи менше, а деякі попередні будуть переформульовані в стилі двох останніх.

%enditem

\end {large}}

%end

\newpage
%begin

{{
\begin {large}

\newpage
1. Що буде надруковано за виконання?

\begin{verbatim}
print('R', end=' ')
e=('a','b','c','d','e',0,1,2,4)
print(e[0])
\end{verbatim}

%enditem


\newpage
2. Що буде надруковано за виконання?

\begin{verbatim}
print('R', end=' ')
a=(0,1,2,3,4,5,6,7,8,9)
print(a[-11:13:-3])
\end{verbatim}

%enditem



\newpage
3. Що буде надруковано за виконання?

\begin{verbatim}
print('R', end=' ')
c=[10,11,12,13,14]
print(c.pop(3), end=' ')
print(c)
\end{verbatim}

%enditem



\newpage
4. Що буде надруковано за виконання?

\begin{verbatim}
print('R', end=' ')
h=['0','1','2','3','4']
h[-2:3]=()
print(h)
\end{verbatim}

%enditem



\newpage
5. Укажіть ТИП значення виразу.\par
\begin{verbatim}
[13,2,4]+[3,5]
\end{verbatim}
%enditem


\newpage
6. Що буде надруковано за виконання?\par
\begin{verbatim}
try:
    s = {33:2, 32:7, 78:9, 33:2}
    s[32] = 9
    for x in s.keys():
        print(x, sep=' ', end=' ')
except:
    print('ERR')
\end{verbatim}
%enditem


\newpage
7. 
 Задано числову послідовність \kw{s} довжини 6000. Записати ВИРАЗ, значенням якого є список, що складається з першої третини послідовності, записаної в оберненому порядку.
%enditem


\newpage
8. Що буде надруковано за виконання?
\begin{verbatim}
d = 0
m = 1

class B:
    d = 2
    
    def __init__(self):
        self.d = 3
        d = 4
        B.d = 5

obj = B()
try:
    print(d, end=' ')
    print(m, end=' ')
    print(B.d, end=' ')
    print(obj.d)
except:
    print('ERR')
\end{verbatim}
%enditem


\newpage
9. Що буде надруковано за виконання?
\begin{verbatim}
class B0:
    def c(self): print(1, end=' ')
    def f(self): print(2, end=' ')

class B1(B0):
    def c(self): print(3, end=' ')

try:
    B1().f()
    B0().c()
except:
    print('ERR')
\end{verbatim}
%enditem


\newpage
10. Що буде надруковано за виконання?
\begin{verbatim}
class A:
    c = 1

    def __init__(self):
        a = 2
        a = 3

class B(A):
    c = 4
    b = 5

    def __init__(self):
        super().__init__()
        self.b = 6

try:
    obj = B()
    obj.a = obj.a + 10
    B.a = 13
except: print('EE', end=' ')

try: print(obj.a, end= ' ')
except: print('E1', end=' ')

try: print(A.a, end= ' ')
except: print('E2', end=' ')

try: print(B.a, end= ' ')
except: print('E3', end=' ')
\end{verbatim}
%enditem


\newpage
11. Змінні \kw{next} та \kw{prev} вузла списку позначають наступний та попередній за ним вузол відповідно. Починаючи з голови, у список записано послідовні цілі числа від~$-40$ до~$28$. Нехай змінна \kw{p} позначає вузол, що зберігає значення~$2$.
Яке значення зберігається у вузлі \kw{p.prev.next.prev.prev.prev} ?
%enditem


\newpage
12. 
\begin{center}
	\adjustimage{max size={0.8\linewidth}{0.8\paperheight}}
	{BinTree/trees_exp/1.png}
\end{center}
Указати послідовність міток вершин дерева, зображеного на рис., що утвориться в результаті обходу дерева в оберненому порядку. Мітки записувати через пробіл.\par Алгоритм починає роботу з кореня (на рис. це верхня вершина).
%enditem


\newpage
13. 
\begin{center}
	\adjustimage{max size={0.9\linewidth}{0.9\paperheight}}
	{Graphs/AdjList/273.png}
\end{center}
Для графа, зображеного на рис., вказати список суміжності вершини 1.\par
Формат оформлення відповіді:\par
список суміжності виписувати через пробіл на одному рядку, якщо пробілів десь буде більше одного - не важливо, упорядкування вершин у списку також значення не має.
%enditem


\newpage
14. Для графа на рисунку з вершини 0 запущено алгоритм пошуку в глибину, що будує кістякове дерево. Списки суміжності вершин переглядаються алгоритмом за зростанням міток вершин.\par
\begin{center}
	\adjustimage{max size={0.9\linewidth}{0.9\paperheight}}
	{Graphs/Traversals/185.png}
\end{center}
\par Які з наведених ребер входять до складу отриманого кістякового дерева?\par
1) (2, 6)\par
2) (4, 6)\par
3) (5, 6)\par
4) (3, 5)\par
5) (2, 5)\par
6) (1, 5)\par
{\vskip 15pt} Номери відповідних ребер записати за зростанням через пробіл.

%enditem


\newpage
15. Для графа на рисунку з вершини 4 запущено алгоритм пошуку в ширину, що будує кістякове дерево. Списки суміжності вершин переглядаються алгоритмом за зростанням міток вершин.\par
\begin{center}
	\adjustimage{max size={0.9\linewidth}{0.9\paperheight}}
	{Graphs/Traversals/260.png}
\end{center}
\par Які з наведених ребер входять до складу отриманого кістякового дерева?\par
1) (5, 6)\par
2) (2, 6)\par
3) (2, 4)\par
4) (3, 5)\par
5) (0, 1)\par
6) (0, 3)\par
{\vskip 15pt} Номери відповідних ребер записати за зростанням через пробіл.


%enditem


\newpage
16. 

Рядки журналу через двокрапку містять ім'я користувача (ідентифікатор), час події,тип події, код події, наприклад
\begin{verbatim}
s="username : 2022-05-05 13:26 : ERROR : 404"
\end{verbatim}
у коді 
\begin{verbatim}
res = re.fullmatch('???', s) 
s = ''
code_str = ??? 
\end{verbatim}
чим треба замінити кожен з ???, щоб значенням code\_str був код події? 


\vspace{10mm}
\textbf{Завдання 17} 

Задано програму, що містить код 1-2 класів, можливе успадкування, та клієнтський код.
У наведеному коді невеликі частини закриті. 
Також наведено вихід програми.
Необхідно відновити закриті частини, щоб отриманий код відповідав наведеному виходу програми.


\textbf{Завдання 18-19} 

Кілька запитань за завданнями для самостійної роботи.
\vspace{10mm}
Скоріш за все, що завдань буде трохи менше, а деякі попередні будуть переформульовані в стилі двох останніх.

%enditem

\end {large}}

%end

\newpage
%begin

{{
\begin {large}

\newpage
1. Що буде надруковано за виконання?

\begin{verbatim}
print('R', end=' ')
f=('a','b','c','d','e',0,1,2)
print(f[-6])
\end{verbatim}

%enditem


\newpage
2. Що буде надруковано за виконання?

\begin{verbatim}
print('R', end=' ')
c=[0,1,2,3,4,5,6,7,8,9]
print(c[-1::-1])
\end{verbatim}

%enditem



\newpage
3. Що буде надруковано за виконання?

\begin{verbatim}
print('R', end=' ')
g=[10,11,12,13,14]
g+=[4,5]
print(g)
\end{verbatim}

%enditem



\newpage
4. Що буде надруковано за виконання?

\begin{verbatim}
print('R', end=' ')
f=['0','1','2','3','4']
f[:7]='11'
print(f)
\end{verbatim}

%enditem



\newpage
5. Укажіть ТИП значення виразу.\par
\begin{verbatim}
[x for x in range(7)]
\end{verbatim}
%enditem


\newpage
6. Що буде надруковано за виконання?\par
\begin{verbatim}
try:
    d = {83:9, 53:8, 85:7, 53:9}
    d[83] = 5
    for x in d.values():
        print(x, sep=' ', end=' ')
except:
    print('ERR')
\end{verbatim}
%enditem


\newpage
7. Нехай змінна \kw{s} позначає дуже довгу послідовність, по якій можна ітерувати. 

Написати вираз-генератор, що задає підпослідовність її елементів, що є рядками типу \kw{str}.

Обмеження: усі елементи шуканої підпослідовності одночасно в пам'яті розміщені бути не можуть. 
%enditem


\newpage
8. Що буде надруковано за виконання?
%Оригинал был утерян. Требует отладки!
\begin{verbatim}
__c = 0

class A:
    __c = 1
    
    def __init__(self):
        self.__c = 2
        __c = 3
        A.__c = 4

obj = A()
obj.__c = 5
try:
    print(__c, end=' ')
    print(obj.__c, end=' ')
    print(A.__c, end=' ')
    print(obj.__c, end=' ')
    print(obj._A__c)
except:
    print('ERR')
\end{verbatim}
%enditem


\newpage
9. Що буде надруковано за виконання?
\begin{verbatim}
class A0:
    def c(self): print(1, end=' ')
    def h(self): print(2, end=' ')

class A1(A0):
    def c(self): print(3, end=' ')

try:
    A1().h()
    A1().c()
except:
    print('ERR')
\end{verbatim}
%enditem


\newpage
10. Що буде надруковано за виконання?
\begin{verbatim}
class A:
    d = 1

    def _h1(self): print(2, end=' ')
    def h2(self): print(3, end=' ')

    def main(self):
        try: print(self.d, end=' ')
        except: print('E1', end=' ')

        try: self._h1()
        except: print('E2', end=' ')

        try: self.h2()
        except: print('E3', end=' ')


class B(A):
    d = 4

    def _h1(self): print(5, end=' ')
    def h2(self): print(6, end=' ')

try: B().main()
except: print('E4', end=' ')
\end{verbatim}
%enditem


\newpage
11. Змінні \kw{next} та \kw{prev} вузла списку позначають наступний та попередній за ним вузол відповідно. Починаючи з голови, у список записано послідовні цілі числа від~$-37$ до~$49$. Нехай змінна \kw{p} позначає вузол, що зберігає значення~$-2$.
Яке значення зберігається у вузлі \kw{p.prev.next.next.next.prev} ?
%enditem


\newpage
12. 
\begin{center}
	\adjustimage{max size={0.8\linewidth}{0.8\paperheight}}
	{BinTree/trees_exp/30.png}
\end{center}
Указати послідовність міток вершин дерева, зображеного на рис., що утвориться в результаті обходу дерева в прямому порядку. Мітки записувати через пробіл.\par Алгоритм починає роботу з кореня (на рис. це верхня вершина).
%enditem


\newpage
13. 
\begin{center}
	\adjustimage{max size={0.9\linewidth}{0.9\paperheight}}
	{Graphs/AdjList/81.png}
\end{center}
Для графа, зображеного на рис., вказати список суміжності вершини 1.\par
Формат оформлення відповіді:\par
список суміжності виписувати через пробіл на одному рядку, якщо пробілів десь буде більше одного - не важливо, упорядкування вершин у списку також значення не має.
%enditem


\newpage
14. Для графа на рисунку з вершини 1 запущено алгоритм пошуку в глибину, що будує кістякове дерево. Списки суміжності вершин переглядаються алгоритмом за зростанням міток вершин.\par
\begin{center}
	\adjustimage{max size={0.9\linewidth}{0.9\paperheight}}
	{Graphs/Traversals/2.png}
\end{center}
\par Які з наведених ребер входять до складу отриманого кістякового дерева?\par
1) (2, 3)\par
2) (1, 5)\par
3) (1, 6)\par
4) (1, 2)\par
5) (1, 3)\par
6) (5, 6)\par
{\vskip 15pt} Номери відповідних ребер записати за зростанням через пробіл.

%enditem


\newpage
15. Для графа на рисунку з вершини 3 запущено алгоритм пошуку в ширину, що будує кістякове дерево. Списки суміжності вершин переглядаються алгоритмом за зростанням міток вершин.\par
\begin{center}
	\adjustimage{max size={0.9\linewidth}{0.9\paperheight}}
	{Graphs/Traversals/282.png}
\end{center}
\par Які з наведених ребер входять до складу отриманого кістякового дерева?\par
1) (1, 5)\par
2) (2, 6)\par
3) (3, 6)\par
4) (0, 2)\par
5) (1, 3)\par
6) (0, 1)\par
{\vskip 15pt} Номери відповідних ребер записати за зростанням через пробіл.


%enditem


\newpage
16. 

Рядки журналу через двокрапку містять ім'я користувача (ідентифікатор), час події,тип події, код події, наприклад
\begin{verbatim}
s="username : 2022-05-05 13:26 : ERROR : 404"
\end{verbatim}
у коді 
\begin{verbatim}
res = re.fullmatch('???', s) 
s = ''
code_str = ??? 
\end{verbatim}
чим треба замінити кожен з ???, щоб значенням code\_str був код події? 


\vspace{10mm}
\textbf{Завдання 17} 

Задано програму, що містить код 1-2 класів, можливе успадкування, та клієнтський код.
У наведеному коді невеликі частини закриті. 
Також наведено вихід програми.
Необхідно відновити закриті частини, щоб отриманий код відповідав наведеному виходу програми.


\textbf{Завдання 18-19} 

Кілька запитань за завданнями для самостійної роботи.
\vspace{10mm}
Скоріш за все, що завдань буде трохи менше, а деякі попередні будуть переформульовані в стилі двох останніх.

%enditem

\end {large}}

%end

\newpage
%begin

{{
\begin {large}

\newpage
1. Що буде надруковано за виконання?

\begin{verbatim}
print('R', end=' ')
b=('a','b','c','d','e',0,1,2,3,4)
print(b[14])
\end{verbatim}

%enditem


\newpage
2. Що буде надруковано за виконання?

\begin{verbatim}
print('R', end=' ')
h=('0','1','2','3','4','5','6','7','8','9')
print(h[15::-1])
\end{verbatim}

%enditem



\newpage
3. Що буде надруковано за виконання?

\begin{verbatim}
print('R', end=' ')
d=['0','1','2','3','4']
d.insert(-2, [1,2])
print(d)
\end{verbatim}

%enditem



\newpage
4. Що буде надруковано за виконання?

\begin{verbatim}
print('R', end=' ')
g=[10,11,12,13,14]
del g[3:4]
print(g)
\end{verbatim}

%enditem



\newpage
5. Укажіть ТИП значення виразу.\par
\begin{verbatim}
[{1,2}, 5]
\end{verbatim}
%enditem


\newpage
6. Що буде надруковано за виконання?\par
\begin{verbatim}
try:
    t = {41:0, 73:5, 32:4, 11:4, 73:6}
    t[51] = 5
    for x, y in t.items():
        print(x, y, sep=' ', end=' ')
except:
    print('ERR')
\end{verbatim}
%enditem


\newpage
7. Нехай змінна \kw{s} позначає дуже довгу послідовність цілих, по якій можна ітерувати. 

Написати вираз-генератор, що задає підпослідовність її елементів, що є точними квадратами. 

Обмеження: усі елементи шуканої підпослідовності одночасно в пам'яті розміщені бути не можуть. 
%enditem


\newpage
8. Що буде надруковано за виконання?
%Оригинал был утерян. Требует отладки!
\begin{verbatim}
_a = 0

class B:
    _a = 1
    
    def __init__(self):
        self._a = 2
        _a = 3
        B._a = 4

obj = B()
B._a = 5
try:
    print(_a, end=' ')
    print(obj._a, end=' ')
    print(B._a, end=' ')
    print(obj._a, end=' ')
    print(obj._B__a)
except:
    print('ERR')
\end{verbatim}
%enditem


\newpage
9. Що буде надруковано за виконання?
\begin{verbatim}
class D0:
    def d(self): print(1, end=' ')
    def g(self): print(2, end=' ')

class D1(D0):
    def d(self): print(3, end=' ')

try:
    D1().g()
    D0().d()
except:
    print('ERR')
\end{verbatim}
%enditem


\newpage
10. Що буде надруковано за виконання?
\begin{verbatim}
class A:
    a = 1

    def __g1(self): print(2, end=' ')
    def _g2(self): print(3, end=' ')

    def main(self):
        try: print(self.a, end=' ')
        except: print('E1', end=' ')

        try: self.__g1()
        except: print('E2', end=' ')

        try: self._g2()
        except: print('E3', end=' ')


class B(A):
    a = 4

    def __g1(self): print(5, end=' ')
    def _g2(self): print(6, end=' ')

try: B().main()
except: print('E4', end=' ')
\end{verbatim}
%enditem


\newpage
11. Змінні \kw{next} та \kw{prev} вузла списку позначають наступний та попередній за ним вузол відповідно. Починаючи з голови, у список записано послідовні цілі числа від~$-63$ до~$39$. Нехай змінна \kw{p} позначає вузол, що зберігає значення~$-9$.
Яке значення зберігається у вузлі \kw{p.next.prev.prev.prev.next} ?
%enditem


\newpage
12. 
\begin{center}
	\adjustimage{max size={0.8\linewidth}{0.8\paperheight}}
	{BinTree/trees_exp/90.png}
\end{center}
Указати послідовність міток вершин дерева, зображеного на рис., що утвориться в результаті обходу дерева в оберненому порядку. Мітки записувати через пробіл.\par Алгоритм починає роботу з кореня (на рис. це верхня вершина).
%enditem


\newpage
13. 
\begin{center}
	\adjustimage{max size={0.9\linewidth}{0.9\paperheight}}
	{Graphs/AdjList/140.png}
\end{center}
Для графа, зображеного на рис., вказати список суміжності вершини 1.\par
Формат оформлення відповіді:\par
список суміжності виписувати через пробіл на одному рядку, якщо пробілів десь буде більше одного - не важливо, упорядкування вершин у списку також значення не має.
%enditem


\newpage
14. Для графа на рисунку вказати послідовність, в якій алгоритм пошуку в глибину буде відмічати відвідані вершини, якщо списки суміжності вершин переглядаються алгоритмом за зростанням міток вершин та вершини відмічаються на першому вход\-женні у вершину.\par
Пошук розпочинається з вершини 5.\par

\begin{center}
	\adjustimage{max size={0.9\linewidth}{0.9\paperheight}}
	{Graphs/Traversals/207.png}
\end{center}
%enditem


\newpage
15. Для графа на рисунку з вершини 6 запущено алгоритм пошуку в ширину, що будує кістякове дерево. Списки суміжності вершин переглядаються алгоритмом за зростанням міток вершин.\par
\begin{center}
	\adjustimage{max size={0.9\linewidth}{0.9\paperheight}}
	{Graphs/Traversals/306.png}
\end{center}
\par Які з наведених ребер входять до складу отриманого кістякового дерева?\par
1) (4, 6)\par
2) (3, 5)\par
3) (1, 6)\par
4) (0, 3)\par
5) (0, 1)\par
6) (0, 5)\par
{\vskip 15pt} Номери відповідних ребер записати за зростанням через пробіл.


%enditem


\newpage
16. 

Рядки журналу через двокрапку містять ім'я користувача (ідентифікатор), час події,тип події, код події, наприклад
\begin{verbatim}
s="username : 2022-05-05 13:26 : ERROR : 404"
\end{verbatim}
у коді 
\begin{verbatim}
res = re.fullmatch('???', s) 
s = ''
code_str = ??? 
\end{verbatim}
чим треба замінити кожен з ???, щоб значенням code\_str був код події? 


\vspace{10mm}
\textbf{Завдання 17} 

Задано програму, що містить код 1-2 класів, можливе успадкування, та клієнтський код.
У наведеному коді невеликі частини закриті. 
Також наведено вихід програми.
Необхідно відновити закриті частини, щоб отриманий код відповідав наведеному виходу програми.


\textbf{Завдання 18-19} 

Кілька запитань за завданнями для самостійної роботи.
\vspace{10mm}
Скоріш за все, що завдань буде трохи менше, а деякі попередні будуть переформульовані в стилі двох останніх.

%enditem

\end {large}}

%end

\newpage
%begin

{{
\begin {large}

\newpage
1. Що буде надруковано за виконання?

\begin{verbatim}
print('R', end=' ')
c=('a','b','c','d','e',0,1,2)
print(c[-10])
\end{verbatim}

%enditem


\newpage
2. Що буде надруковано за виконання?

\begin{verbatim}
print('R', end=' ')
b=['0','1','2','3','4','5','6','7','8','9']
print(b[:1:-2])
\end{verbatim}

%enditem



\newpage
3. Що буде надруковано за виконання?

\begin{verbatim}
print('R', end=' ')
h=['0','1','2','3','4']
h.remove(3)
print(h)
\end{verbatim}

%enditem



\newpage
4. Що буде надруковано за виконання?

\begin{verbatim}
print('R', end=' ')
c=[10,11,12,13,14]
del c[:2]
print(c)
\end{verbatim}

%enditem



\newpage
5. Укажіть ТИП значення виразу.\par
\begin{verbatim}
{x:0 for x in range(7)}
\end{verbatim}
%enditem


\newpage
6. Що буде надруковано за виконання?\par
\begin{verbatim}
try:
    s = {33:2, 25:5, 12:4, 25:7}
    s[70] = 3
    for x in s.items():
        print(*x, sep=' ', end=' ')
except:
    print('ERR')
\end{verbatim}
%enditem


\newpage
7. Нехай змінна \kw{s} позначає послідовність цілих, по якій можна ітерувати. 

Написати вираз, що перевіряє, чи правда, що послідовність  містить від’ємні числа.
%enditem


\newpage
8. Що буде надруковано за виконання?
%Оригинал был утерян. Требует отладки!
\begin{verbatim}
c = 0

class D:
    c = 1
    
    def __init__(self):
        self.c = 2
        c = 3
        D.c = 4

obj = D()
obj.c = 5
try:
    print(c, end=' ')
    print(obj.c, end=' ')
    print(D.c, end=' ')
    print(obj.c, end=' ')
    print(obj._D__c)
except:
    print('ERR')
\end{verbatim}
%enditem


\newpage
9. Що буде надруковано за виконання?
\begin{verbatim}
class C0:
    def b(self): print(1, end=' ')
    def h(self): print(2, end=' ')

class C1(C0):
    def h(self): print(3, end=' ')

try:
    C1().b()
    C1().h()
except:
    print('ERR')
\end{verbatim}
%enditem


\newpage
10. Що буде надруковано за виконання?
\begin{verbatim}
class A:
    b = 1

    def _f1(self): print(2, end=' ')
    def __f2(self): print(3, end=' ')

    def main(self):
        try: print(self.b, end=' ')
        except: print('E1', end=' ')

        try: self._f1()
        except: print('E2', end=' ')

        try: self.__f2()
        except: print('E3', end=' ')


class B(A):
    b = 4

    def _f1(self): print(5, end=' ')
    def __f2(self): print(6, end=' ')

try: B().main()
except: print('E4', end=' ')
\end{verbatim}
%enditem


\newpage
11. Змінні \kw{next} та \kw{prev} вузла списку позначають наступний та попередній за ним вузол відповідно. Починаючи з голови, у список записано послідовні цілі числа від~$-52$ до~$25$. Нехай змінна \kw{p} позначає вузол, що зберігає значення~$-4$.
Яке значення зберігається у вузлі \kw{p.prev.next.prev.next.next} ?
%enditem


\newpage
12. 
\begin{center}
	\adjustimage{max size={0.8\linewidth}{0.8\paperheight}}
	{BinTree/trees_exp/164.png}
\end{center}
Указати послідовність міток вершин дерева, зображеного на рис., що утвориться в результаті обходу дерева в прямому порядку. Мітки записувати через пробіл.\par Алгоритм починає роботу з кореня (на рис. це верхня вершина).
%enditem


\newpage
13. 
\begin{center}
	\adjustimage{max size={0.9\linewidth}{0.9\paperheight}}
	{Graphs/AdjList/180.png}
\end{center}
Для графа, зображеного на рис., вказати список суміжності вершини 8.\par
Формат оформлення відповіді:\par
список суміжності виписувати через пробіл на одному рядку, якщо пробілів десь буде більше одного - не важливо, упорядкування вершин у списку також значення не має.
%enditem


\newpage
14. Для графа на рисунку вказати послідовність, в якій алгоритм пошуку в глибину буде відмічати відвідані вершини, якщо списки суміжності вершин переглядаються алгоритмом за зростанням міток вершин та вершини відмічаються на першому вход\-женні у вершину.\par
Пошук розпочинається з вершини 2.\par

\begin{center}
	\adjustimage{max size={0.9\linewidth}{0.9\paperheight}}
	{Graphs/Traversals/164.png}
\end{center}
%enditem


\newpage
15. Для графа на рисунку з вершини 4 запущено алгоритм пошуку в ширину, що будує кістякове дерево. Списки суміжності вершин переглядаються алгоритмом за зростанням міток вершин.\par
\begin{center}
	\adjustimage{max size={0.9\linewidth}{0.9\paperheight}}
	{Graphs/Traversals/311.png}
\end{center}
\par Які з наведених ребер входять до складу отриманого кістякового дерева?\par
1) (3, 5)\par
2) (4, 6)\par
3) (5, 6)\par
4) (2, 3)\par
5) (1, 4)\par
6) (3, 4)\par
{\vskip 15pt} Номери відповідних ребер записати за зростанням через пробіл.


%enditem


\newpage
16. 

Рядки журналу через двокрапку містять ім'я користувача (ідентифікатор), час події,тип події, код події, наприклад
\begin{verbatim}
s="username : 2022-05-05 13:26 : ERROR : 404"
\end{verbatim}
у коді 
\begin{verbatim}
res = re.fullmatch('???', s) 
s = ''
code_str = ??? 
\end{verbatim}
чим треба замінити кожен з ???, щоб значенням code\_str був код події? 


\vspace{10mm}
\textbf{Завдання 17} 

Задано програму, що містить код 1-2 класів, можливе успадкування, та клієнтський код.
У наведеному коді невеликі частини закриті. 
Також наведено вихід програми.
Необхідно відновити закриті частини, щоб отриманий код відповідав наведеному виходу програми.


\textbf{Завдання 18-19} 

Кілька запитань за завданнями для самостійної роботи.
\vspace{10mm}
Скоріш за все, що завдань буде трохи менше, а деякі попередні будуть переформульовані в стилі двох останніх.

%enditem

\end {large}}

%end

\newpage
%begin

{{
\begin {large}

\newpage
1. Що буде надруковано за виконання?

\begin{verbatim}
print('R', end=' ')
g=('a','b','c','d','e',0,1)
print(g[4])
\end{verbatim}

%enditem


\newpage
2. Що буде надруковано за виконання?

\begin{verbatim}
print('R', end=' ')
g=(0,1,2,3,4,5,6,7,8,9)
print(g[13:0:-3])
\end{verbatim}

%enditem



\newpage
3. Що буде надруковано за виконання?

\begin{verbatim}
print('R', end=' ')
a=[10,11,12,13,14]
print(a.index('11'), end=' ')
\end{verbatim}

%enditem



\newpage
4. Що буде надруковано за виконання?

\begin{verbatim}
print('R', end=' ')
d=['0','1','2','3','4']
d[7:-7]=()
print(d)
\end{verbatim}

%enditem



\newpage
5. Укажіть ТИП значення виразу.\par
\begin{verbatim}
(2, 7, -4)
\end{verbatim}
%enditem


\newpage
6. Що буде надруковано за виконання?\par
\begin{verbatim}
try:
    t = {29:7, 53:3, 63:7, 66:3, 66:2}
    t[55] = 9
    for x, y in t.items():
        print(x, y, sep=' ', end=' ')
except:
    print('ERR')
\end{verbatim}
%enditem


\newpage
7. Записати ВИРАЗ, значенням якого є словник, ключами якого є цілі числа від 16 до 2010 (включно), що не діляться на 7, а ключам відповідають їх квадрати 
%enditem


\newpage
8. Що буде надруковано за виконання?
\begin{verbatim}
b = 0
n = 1

class C:
    b = 2
    
    def __init__(self):
        global b
        self.n = 3
        b = 4
        C.b = 5

obj = C()
try:
    print(b, end=' ')
    print(n, end=' ')
    print(C.b, end=' ')
    print(obj.n)
except:
    print('ERR')
\end{verbatim}
%enditem


\newpage
9. Що буде надруковано за виконання?
\begin{verbatim}
class B0:
    def a(self): print(1, end=' ')
    def g(self): print(2, end=' ')

class B1(B0):
    def a(self): print(3, end=' ')

try:
    B0().g()
    B1().a()
except:
    print('ERR')
\end{verbatim}
%enditem


\newpage
10. Що буде надруковано за виконання?
\begin{verbatim}
class A:
    c = 1

    def __h1(self): print(2, end=' ')
    def _h2(self): print(3, end=' ')

    def main(self):
        try: print(self.c, end=' ')
        except: print('E1', end=' ')

        try: self.__h1()
        except: print('E2', end=' ')

        try: self._h2()
        except: print('E3', end=' ')


class B(A):
    c = 4

    def __h1(self): print(5, end=' ')
    def _h2(self): print(6, end=' ')

try: B().main()
except: print('E4', end=' ')
\end{verbatim}
%enditem


\newpage
11. Змінні \kw{next} та \kw{prev} вузла списку позначають наступний та попередній за ним вузол відповідно. Починаючи з голови, у список записано послідовні цілі числа від~$-35$ до~$33$. Нехай змінна \kw{p} позначає вузол, що зберігає значення~$-6$.
Яке значення зберігається у вузлі \kw{p.next.prev.next.prev.prev} ?
%enditem


\newpage
12. 
\begin{center}
	\adjustimage{max size={0.8\linewidth}{0.8\paperheight}}
	{BinTree/trees_exp/0.png}
\end{center}
Указати послідовність міток вершин дерева, зображеного на рис., що утвориться в результаті обходу дерева в оберненому порядку. Мітки записувати через пробіл.\par Алгоритм починає роботу з кореня (на рис. це верхня вершина).
%enditem


\newpage
13. 
\begin{center}
	\adjustimage{max size={0.9\linewidth}{0.9\paperheight}}
	{Graphs/AdjList/298.png}
\end{center}
Для графа, зображеного на рис., вказати список суміжності вершини 7.\par
Формат оформлення відповіді:\par
список суміжності виписувати через пробіл на одному рядку, якщо пробілів десь буде більше одного - не важливо, упорядкування вершин у списку також значення не має.
%enditem


\newpage
14. Для графа на рисунку вказати послідовність, в якій алгоритм пошуку в глибину буде відмічати відвідані вершини, якщо списки суміжності вершин переглядаються алгоритмом за зростанням міток вершин та вершини відмічаються на першому вход\-женні у вершину.\par
Пошук розпочинається з вершини 6.\par

\begin{center}
	\adjustimage{max size={0.9\linewidth}{0.9\paperheight}}
	{Graphs/Traversals/61.png}
\end{center}
%enditem


\newpage
15. Для графа на рисунку з вершини 3 запущено алгоритм пошуку в ширину, що будує кістякове дерево. Списки суміжності вершин переглядаються алгоритмом за зростанням міток вершин.\par
\begin{center}
	\adjustimage{max size={0.9\linewidth}{0.9\paperheight}}
	{Graphs/Traversals/226.png}
\end{center}
\par Які з наведених ребер входять до складу отриманого кістякового дерева?\par
1) (1, 5)\par
2) (1, 2)\par
3) (0, 5)\par
4) (2, 5)\par
5) (0, 2)\par
6) (3, 4)\par
{\vskip 15pt} Номери відповідних ребер записати за зростанням через пробіл.


%enditem


\newpage
16. 

Рядки журналу через двокрапку містять ім'я користувача (ідентифікатор), час події,тип події, код події, наприклад
\begin{verbatim}
s="username : 2022-05-05 13:26 : ERROR : 404"
\end{verbatim}
у коді 
\begin{verbatim}
res = re.fullmatch('???', s) 
s = ''
code_str = ??? 
\end{verbatim}
чим треба замінити кожен з ???, щоб значенням code\_str був код події? 


\vspace{10mm}
\textbf{Завдання 17} 

Задано програму, що містить код 1-2 класів, можливе успадкування, та клієнтський код.
У наведеному коді невеликі частини закриті. 
Також наведено вихід програми.
Необхідно відновити закриті частини, щоб отриманий код відповідав наведеному виходу програми.


\textbf{Завдання 18-19} 

Кілька запитань за завданнями для самостійної роботи.
\vspace{10mm}
Скоріш за все, що завдань буде трохи менше, а деякі попередні будуть переформульовані в стилі двох останніх.

%enditem

\end {large}}

%end

\newpage
%begin

{{
\begin {large}

\newpage
1. Що буде надруковано за виконання?

\begin{verbatim}
print('R', end=' ')
e=('a','b','c','d','e',0,1,2,4)
print(e[-11])
\end{verbatim}

%enditem


\newpage
2. Що буде надруковано за виконання?

\begin{verbatim}
print('R', end=' ')
d='0123456789'
print(d[::3])
\end{verbatim}

%enditem



\newpage
3. Що буде надруковано за виконання?

\begin{verbatim}
print('R', end=' ')
e=['0','1','2','3','4']
e.extend(20)
print(e)
\end{verbatim}

%enditem



\newpage
4. Що буде надруковано за виконання?

\begin{verbatim}
print('R', end=' ')
e=['0','1','2','3','4']
e[-7:-3]='11'
print(e)
\end{verbatim}

%enditem



\newpage
5. Укажіть ТИП значення виразу.\par
\begin{verbatim}
[13,2,4]+[3,5]
\end{verbatim}
%enditem


\newpage
6. Що буде надруковано за виконання?\par
\begin{verbatim}
try:
    d = {65:8, 62:9, 79:9, 62:8}
    d[15] = 7
    for x in d.keys():
        print(x, sep=' ', end=' ')
except:
    print('ERR')
\end{verbatim}
%enditem


\newpage
7. Записати ВИРАЗ, значенням якого є множина, що складається з цілих чисел від 19 до 2020 (включно), що не діляться на жодне з чисел 2, 3.
%enditem


\newpage
8. Що буде надруковано за виконання?
\begin{verbatim}
c = 0
m = 1

class A:
    m = 2
    
    def __init__(self):
        self.c = 3
        m = 4
        A.m = 5

obj = A()
try:
    print(c, end=' ')
    print(m, end=' ')
    print(A.m, end=' ')
    print(obj.c)
except:
    print('ERR')
\end{verbatim}
%enditem


\newpage
9. Що буде надруковано за виконання?
\begin{verbatim}
class C0:
    def b(self): print(1, end=' ')
    def f(self): print(2, end=' ')

class C1(C0):
    def b(self): print(3, end=' ')

try:
    C1().f()
    C1().b()
except:
    print('ERR')
\end{verbatim}
%enditem


\newpage
10. Що буде надруковано за виконання?
\begin{verbatim}
class A:
    d = 1

    def _g1(self): print(2, end=' ')
    def _g2(self): print(3, end=' ')

    def main(self):
        try: print(self.d, end=' ')
        except: print('E1', end=' ')

        try: self._g1()
        except: print('E2', end=' ')

        try: self._g2()
        except: print('E3', end=' ')


class B(A):
    d = 4

    def _g1(self): print(5, end=' ')
    def _g2(self): print(6, end=' ')

try: B().main()
except: print('E4', end=' ')
\end{verbatim}
%enditem


\newpage
11. Змінні \kw{next} та \kw{prev} вузла списку позначають наступний та попередній за ним вузол відповідно. Починаючи з голови, у список записано послідовні цілі числа від~$-69$ до~$38$. Нехай змінна \kw{p} позначає вузол, що зберігає значення~$-8$.
Яке значення зберігається у вузлі \kw{p.prev.next.next.prev.next} ?
%enditem


\newpage
12. 
\begin{center}
	\adjustimage{max size={0.8\linewidth}{0.8\paperheight}}
	{BinTree/trees_exp/138.png}
\end{center}
Указати послідовність міток вершин дерева, зображеного на рис., що утвориться в результаті обходу дерева в прямому порядку. Мітки записувати через пробіл.\par Алгоритм починає роботу з кореня (на рис. це верхня вершина).
%enditem


\newpage
13. 
\begin{center}
	\adjustimage{max size={0.9\linewidth}{0.9\paperheight}}
	{Graphs/AdjList/216.png}
\end{center}
Для графа, зображеного на рис., вказати список суміжності вершини 1.\par
Формат оформлення відповіді:\par
список суміжності виписувати через пробіл на одному рядку, якщо пробілів десь буде більше одного - не важливо, упорядкування вершин у списку також значення не має.
%enditem


\newpage
14. Для графа на рисунку з вершини 3 запущено алгоритм пошуку в глибину, що будує кістякове дерево. Списки суміжності вершин переглядаються алгоритмом за зростанням міток вершин.\par
\begin{center}
	\adjustimage{max size={0.9\linewidth}{0.9\paperheight}}
	{Graphs/Traversals/343.png}
\end{center}
\par Які з наведених ребер входять до складу отриманого кістякового дерева?\par
1) (0, 1)\par
2) (0, 3)\par
3) (1, 4)\par
4) (2, 6)\par
5) (1, 3)\par
6) (0, 6)\par
{\vskip 15pt} Номери відповідних ребер записати за зростанням через пробіл.

%enditem


\newpage
15. Для графа на рисунку з вершини 3 запущено алгоритм пошуку в ширину, що будує кістякове дерево. Списки суміжності вершин переглядаються алгоритмом за зростанням міток вершин.\par
\begin{center}
	\adjustimage{max size={0.9\linewidth}{0.9\paperheight}}
	{Graphs/Traversals/212.png}
\end{center}
\par Які з наведених ребер входять до складу отриманого кістякового дерева?\par
1) (0, 5)\par
2) (1, 5)\par
3) (0, 6)\par
4) (1, 4)\par
5) (3, 6)\par
6) (0, 2)\par
{\vskip 15pt} Номери відповідних ребер записати за зростанням через пробіл.


%enditem


\newpage
16. 

Рядки журналу через двокрапку містять ім'я користувача (ідентифікатор), час події,тип події, код події, наприклад
\begin{verbatim}
s="username : 2022-05-05 13:26 : ERROR : 404"
\end{verbatim}
у коді 
\begin{verbatim}
res = re.fullmatch('???', s) 
s = ''
code_str = ??? 
\end{verbatim}
чим треба замінити кожен з ???, щоб значенням code\_str був код події? 


\vspace{10mm}
\textbf{Завдання 17} 

Задано програму, що містить код 1-2 класів, можливе успадкування, та клієнтський код.
У наведеному коді невеликі частини закриті. 
Також наведено вихід програми.
Необхідно відновити закриті частини, щоб отриманий код відповідав наведеному виходу програми.


\textbf{Завдання 18-19} 

Кілька запитань за завданнями для самостійної роботи.
\vspace{10mm}
Скоріш за все, що завдань буде трохи менше, а деякі попередні будуть переформульовані в стилі двох останніх.

%enditem

\end {large}}

%end

\newpage
%begin

{{
\begin {large}

\newpage
1. Що буде надруковано за виконання?

\begin{verbatim}
print('R', end=' ')
a=('a','b','c','d','e',0,1,2,3)
print(a[-1])
\end{verbatim}

%enditem


\newpage
2. Що буде надруковано за виконання?

\begin{verbatim}
print('R', end=' ')
e=[0,1,2,3,4,5,6,7,8,9]
print(e[-9:8:2])
\end{verbatim}

%enditem



\newpage
3. Що буде надруковано за виконання?

\begin{verbatim}
print('R', end=' ')
f=[10,11,12,13,14]
f.append('13')
print(f)
\end{verbatim}

%enditem



\newpage
4. Що буде надруковано за виконання?

\begin{verbatim}
print('R', end=' ')
a=[10,11,12,13,14]
del a[8:-9]
print(a)
\end{verbatim}

%enditem



\newpage
5. Укажіть ТИП значення виразу.\par
\begin{verbatim}
{y:13 for y in range(-7,7)}
\end{verbatim}
%enditem


\newpage
6. Що буде надруковано за виконання?\par
\begin{verbatim}
try:
    t = {50:7, 80:7, 62:5, 62:6}
    t[80] = 5
    for x in t.items():
        print(x, sep=' ', end=' ')
except:
    print('ERR')
\end{verbatim}
%enditem


\newpage
7. 
Записати ВИРАЗ, значенням якого є список, що складається з цілих чисел від 10 до 2005 (включно), причому спочатку за зростанням записано всі, що діляться на 3, а потім решту за зростанням.

%enditem


\newpage
8. Що буде надруковано за виконання?
\begin{verbatim}
b = 0
h = 1

class B:
    b = 2
    
    def __init__(self):
        global b
        self.h = 3
        b = 4
        B.b = 5

obj = B()
try:
    print(b, end=' ')
    print(h, end=' ')
    print(B.b, end=' ')
    print(obj.h)
except:
    print('ERR')
\end{verbatim}
%enditem


\newpage
9. Що буде надруковано за виконання?
\begin{verbatim}
class A0:
    def c(self): print(1, end=' ')
    def f(self): print(2, end=' ')

class A1(A0):
    def c(self): print(3, end=' ')

try:
    A1().f()
    A0().c()
except:
    print('ERR')
\end{verbatim}
%enditem


\newpage
10. Що буде надруковано за виконання?
\begin{verbatim}
class A:
    b = 1

    def __init__(self):
        a = 2
        a = 3

class B(A):
    b = 4
    c = 5

    def __init__(self):
        super().__init__()
        self.c = 6

try:
    obj = B()
    obj.b = obj.c + 10
    B.b = 13
except: print('EE', end=' ')

try: print(obj.b, end= ' ')
except: print('E1', end=' ')

try: print(A.b, end= ' ')
except: print('E2', end=' ')

try: print(B.b, end= ' ')
except: print('E3', end=' ')
\end{verbatim}
%enditem


\newpage
11. Змінні \kw{next} та \kw{prev} вузла списку позначають наступний та попередній за ним вузол відповідно. Починаючи з голови, у список записано послідовні цілі числа від~$-50$ до~$42$. Нехай змінна \kw{p} позначає вузол, що зберігає значення~$6$.
Яке значення зберігається у вузлі \kw{p.next.prev.prev.next.prev} ?
%enditem


\newpage
12. 
\begin{center}
	\adjustimage{max size={0.8\linewidth}{0.8\paperheight}}
	{BinTree/trees_exp/162.png}
\end{center}
Указати послідовність міток вершин дерева, зображеного на рис., що утвориться в результаті обходу дерева в оберненому порядку. Мітки записувати через пробіл.\par Алгоритм починає роботу з кореня (на рис. це верхня вершина).
%enditem


\newpage
13. 
\begin{center}
	\adjustimage{max size={0.9\linewidth}{0.9\paperheight}}
	{Graphs/AdjList/295.png}
\end{center}
Для графа, зображеного на рис., вказати список суміжності вершини 0.\par
Формат оформлення відповіді:\par
список суміжності виписувати через пробіл на одному рядку, якщо пробілів десь буде більше одного - не важливо, упорядкування вершин у списку також значення не має.
%enditem


\newpage
14. Для графа на рисунку з вершини 5 запущено алгоритм пошуку в глибину, що будує кістякове дерево. Списки суміжності вершин переглядаються алгоритмом за зростанням міток вершин.\par
\begin{center}
	\adjustimage{max size={0.9\linewidth}{0.9\paperheight}}
	{Graphs/Traversals/146.png}
\end{center}
\par Які з наведених ребер входять до складу отриманого кістякового дерева?\par
1) (1, 3)\par
2) (0, 5)\par
3) (1, 6)\par
4) (2, 5)\par
5) (2, 6)\par
6) (1, 4)\par
{\vskip 15pt} Номери відповідних ребер записати за зростанням через пробіл.

%enditem


\newpage
15. Для графа на рисунку з вершини 0 запущено алгоритм пошуку в ширину, що будує кістякове дерево. Списки суміжності вершин переглядаються алгоритмом за зростанням міток вершин.\par
\begin{center}
	\adjustimage{max size={0.9\linewidth}{0.9\paperheight}}
	{Graphs/Traversals/335.png}
\end{center}
\par Які з наведених ребер входять до складу отриманого кістякового дерева?\par
1) (2, 6)\par
2) (0, 3)\par
3) (1, 3)\par
4) (2, 5)\par
5) (0, 2)\par
6) (0, 1)\par
{\vskip 15pt} Номери відповідних ребер записати за зростанням через пробіл.


%enditem


\newpage
16. 

Рядки журналу через двокрапку містять ім'я користувача (ідентифікатор), час події,тип події, код події, наприклад
\begin{verbatim}
s="username : 2022-05-05 13:26 : ERROR : 404"
\end{verbatim}
у коді 
\begin{verbatim}
res = re.fullmatch('???', s) 
s = ''
code_str = ??? 
\end{verbatim}
чим треба замінити кожен з ???, щоб значенням code\_str був код події? 


\vspace{10mm}
\textbf{Завдання 17} 

Задано програму, що містить код 1-2 класів, можливе успадкування, та клієнтський код.
У наведеному коді невеликі частини закриті. 
Також наведено вихід програми.
Необхідно відновити закриті частини, щоб отриманий код відповідав наведеному виходу програми.


\textbf{Завдання 18-19} 

Кілька запитань за завданнями для самостійної роботи.
\vspace{10mm}
Скоріш за все, що завдань буде трохи менше, а деякі попередні будуть переформульовані в стилі двох останніх.

%enditem

\end {large}}

%end

\newpage
%begin

{{
\begin {large}

\newpage
1. Що буде надруковано за виконання?

\begin{verbatim}
print('R', end=' ')
h=('a','b','c','d','e',0,1,2,3,4)
print(h[-12])
\end{verbatim}

%enditem


\newpage
2. Що буде надруковано за виконання?

\begin{verbatim}
print('R', end=' ')
a=('0','1','2','3','4','5','6','7','8','9')
print(a[4:-13:2])
\end{verbatim}

%enditem



\newpage
3. Що буде надруковано за виконання?

\begin{verbatim}
print('R', end=' ')
b=[10,11,12,13,14]
print(b.index(9), end=' ')
\end{verbatim}

%enditem



\newpage
4. Що буде надруковано за виконання?

\begin{verbatim}
print('R', end=' ')
b=[10,11,12,13,14]
b[-7:]=[1,2]
print(b)
\end{verbatim}

%enditem



\newpage
5. Укажіть ТИП значення виразу.\par
\begin{verbatim}
{x for x in range(7)}
\end{verbatim}
%enditem


\newpage
6. Що буде надруковано за виконання?\par
\begin{verbatim}
try:
    d = {81:9, 72:8, 48:1, 81:6}
    d[55] = 1
    for x in d :
        print(x, sep=' ', end=' ')
except:
    print('ERR')
\end{verbatim}
%enditem


\newpage
7. Записати ВИРАЗ, значенням якого є словник, ключами якого є цілі числа від 18 до 2013 (включно), що діляться на 2, а ключам відповідають їх квадрати 
%enditem


\newpage
8. Що буде надруковано за виконання?
\begin{verbatim}
d = 0
k = 1

class A:
    k = 2
    
    def __init__(self):
        self.d = 3
        k = 4
        A.k = 5

obj = A()
try:
    print(d, end=' ')
    print(k, end=' ')
    print(A.k, end=' ')
    print(obj.d)
except:
    print('ERR')
\end{verbatim}
%enditem


\newpage
9. Що буде надруковано за виконання?
\begin{verbatim}
class D0:
    def d(self): print(1, end=' ')
    def g(self): print(2, end=' ')

class D1(D0):
    def g(self): print(3, end=' ')

try:
    D0().d()
    D1().g()
except:
    print('ERR')
\end{verbatim}
%enditem


\newpage
10. Що буде надруковано за виконання?
\begin{verbatim}
class A:
    b = 1

    def _f1(self): print(2, end=' ')
    def __f2(self): print(3, end=' ')

    def main(self):
        try: print(self.b, end=' ')
        except: print('E1', end=' ')

        try: self._f1()
        except: print('E2', end=' ')

        try: self.__f2()
        except: print('E3', end=' ')


class B(A):
    b = 4

    def _f1(self): print(5, end=' ')
    def __f2(self): print(6, end=' ')

try: B().main()
except: print('E4', end=' ')
\end{verbatim}
%enditem


\newpage
11. Змінні \kw{next} та \kw{prev} вузла списку позначають наступний та попередній за ним вузол відповідно. Починаючи з голови, у список записано послідовні цілі числа від~$-46$ до~$26$. Нехай змінна \kw{p} позначає вузол, що зберігає значення~$3$.
Яке значення зберігається у вузлі \kw{p.next.next.next.prev.next} ?
%enditem


\newpage
12. 
\begin{center}
	\adjustimage{max size={0.8\linewidth}{0.8\paperheight}}
	{BinTree/trees_exp/116.png}
\end{center}
Указати послідовність міток вершин дерева, зображеного на рис., що утвориться в результаті обходу дерева в оберненому порядку. Мітки записувати через пробіл.\par Алгоритм починає роботу з кореня (на рис. це верхня вершина).
%enditem


\newpage
13. 
\begin{center}
	\adjustimage{max size={0.9\linewidth}{0.9\paperheight}}
	{Graphs/AdjList/212.png}
\end{center}
Для графа, зображеного на рис., вказати список суміжності вершини 0.\par
Формат оформлення відповіді:\par
список суміжності виписувати через пробіл на одному рядку, якщо пробілів десь буде більше одного - не важливо, упорядкування вершин у списку також значення не має.
%enditem


\newpage
14. Для графа на рисунку вказати послідовність, в якій алгоритм пошуку в глибину буде відмічати відвідані вершини, якщо списки суміжності вершин переглядаються алгоритмом за зростанням міток вершин та вершини відмічаються на першому вход\-женні у вершину.\par
Пошук розпочинається з вершини 6.\par

\begin{center}
	\adjustimage{max size={0.9\linewidth}{0.9\paperheight}}
	{Graphs/Traversals/140.png}
\end{center}
%enditem


\newpage
15. Для графа на рисунку з вершини 1 запущено алгоритм пошуку в ширину, що будує кістякове дерево. Списки суміжності вершин переглядаються алгоритмом за зростанням міток вершин.\par
\begin{center}
	\adjustimage{max size={0.9\linewidth}{0.9\paperheight}}
	{Graphs/Traversals/112.png}
\end{center}
\par Які з наведених ребер входять до складу отриманого кістякового дерева?\par
1) (1, 2)\par
2) (0, 6)\par
3) (0, 1)\par
4) (2, 6)\par
5) (3, 6)\par
6) (5, 6)\par
{\vskip 15pt} Номери відповідних ребер записати за зростанням через пробіл.


%enditem


\newpage
16. 

Рядки журналу через двокрапку містять ім'я користувача (ідентифікатор), час події,тип події, код події, наприклад
\begin{verbatim}
s="username : 2022-05-05 13:26 : ERROR : 404"
\end{verbatim}
у коді 
\begin{verbatim}
res = re.fullmatch('???', s) 
s = ''
code_str = ??? 
\end{verbatim}
чим треба замінити кожен з ???, щоб значенням code\_str був код події? 


\vspace{10mm}
\textbf{Завдання 17} 

Задано програму, що містить код 1-2 класів, можливе успадкування, та клієнтський код.
У наведеному коді невеликі частини закриті. 
Також наведено вихід програми.
Необхідно відновити закриті частини, щоб отриманий код відповідав наведеному виходу програми.


\textbf{Завдання 18-19} 

Кілька запитань за завданнями для самостійної роботи.
\vspace{10mm}
Скоріш за все, що завдань буде трохи менше, а деякі попередні будуть переформульовані в стилі двох останніх.

%enditem

\end {large}}

%end

\newpage
%begin

{{
\begin {large}

\newpage
1. Що буде надруковано за виконання?

\begin{verbatim}
print('R', end=' ')
d=('a','b','c','d','e',0,1,2,4)
print(d[2])
\end{verbatim}

%enditem


\newpage
2. Що буде надруковано за виконання?

\begin{verbatim}
print('R', end=' ')
f='0123456789'
print(f[10:15:-2])
\end{verbatim}

%enditem



\newpage
3. Що буде надруковано за виконання?

\begin{verbatim}
print('R', end=' ')
c=['0','1','2','3','4']
c.append('11')
print(c)
\end{verbatim}

%enditem



\newpage
4. Що буде надруковано за виконання?

\begin{verbatim}
print('R', end=' ')
h=['0','1','2','3','4']
h[-4:]='13'
print(h)
\end{verbatim}

%enditem



\newpage
5. Укажіть ТИП значення виразу.\par
\begin{verbatim}
{x:x for x in range(7)}
\end{verbatim}
%enditem


\newpage
6. Що буде надруковано за виконання?\par
\begin{verbatim}
try:
    s = {25:1, 11:7, 19:4, 52:4, 37:5}
    s[19] = 9
    for x in s.values():
        print(x, sep=' ', end=' ')
except:
    print('ERR')
\end{verbatim}
%enditem


\newpage
7. Нехай змінна \kw{s} позначає дуже довгу послідовність цілих, по якій можна ітерувати. 

Написати вираз-генератор, що задає підпослідовність її елементів, що більші за задане число num. 

Обмеження: усі елементи шуканої підпослідовності одночасно в пам'яті розміщені бути не можуть. 
%enditem


\newpage
8. Що буде надруковано за виконання?
\begin{verbatim}
c = 0
j = 1

class D:
    c = 2
    
    def __init__(self):
        self.j = 3
        c = 4
        D.c = 5

obj = D()
try:
    print(c, end=' ')
    print(j, end=' ')
    print(D.c, end=' ')
    print(obj.j)
except:
    print('ERR')
\end{verbatim}
%enditem


\newpage
9. Що буде надруковано за виконання?
\begin{verbatim}
class B0:
    def a(self): print(1, end=' ')
    def h(self): print(2, end=' ')

class B1(B0):
    def h(self): print(3, end=' ')

try:
    B1().a()
    B1().h()
except:
    print('ERR')
\end{verbatim}
%enditem


\newpage
10. Що буде надруковано за виконання?
\begin{verbatim}
class A:
    c = 1

    def h1(self): print(2, end=' ')
    def _h2(self): print(3, end=' ')

    def main(self):
        try: print(self.c, end=' ')
        except: print('E1', end=' ')

        try: self.h1()
        except: print('E2', end=' ')

        try: self._h2()
        except: print('E3', end=' ')


class B(A):
    c = 4

    def h1(self): print(5, end=' ')
    def _h2(self): print(6, end=' ')

try: B().main()
except: print('E4', end=' ')
\end{verbatim}
%enditem


\newpage
11. Змінні \kw{next} та \kw{prev} вузла списку позначають наступний та попередній за ним вузол відповідно. Починаючи з голови, у список записано послідовні цілі числа від~$-62$ до~$16$. Нехай змінна \kw{p} позначає вузол, що зберігає значення~$4$.
Яке значення зберігається у вузлі \kw{p.prev.prev.prev.next.prev} ?
%enditem


\newpage
12. 
\begin{center}
	\adjustimage{max size={0.8\linewidth}{0.8\paperheight}}
	{BinTree/trees_exp/114.png}
\end{center}
Указати послідовність міток вершин дерева, зображеного на рис., що утвориться в результаті обходу дерева в прямому порядку. Мітки записувати через пробіл.\par Алгоритм починає роботу з кореня (на рис. це верхня вершина).
%enditem


\newpage
13. 
\begin{center}
	\adjustimage{max size={0.9\linewidth}{0.9\paperheight}}
	{Graphs/AdjList/348.png}
\end{center}
Для графа, зображеного на рис., вказати список суміжності вершини 3.\par
Формат оформлення відповіді:\par
список суміжності виписувати через пробіл на одному рядку, якщо пробілів десь буде більше одного - не важливо, упорядкування вершин у списку також значення не має.
%enditem


\newpage
14. Для графа на рисунку вказати послідовність, в якій алгоритм пошуку в глибину буде відмічати відвідані вершини, якщо списки суміжності вершин переглядаються алгоритмом за зростанням міток вершин та вершини відмічаються на першому вход\-женні у вершину.\par
Пошук розпочинається з вершини 1.\par

\begin{center}
	\adjustimage{max size={0.9\linewidth}{0.9\paperheight}}
	{Graphs/Traversals/64.png}
\end{center}
%enditem


\newpage
15. Для графа на рисунку з вершини 3 запущено алгоритм пошуку в ширину, що будує кістякове дерево. Списки суміжності вершин переглядаються алгоритмом за зростанням міток вершин.\par
\begin{center}
	\adjustimage{max size={0.9\linewidth}{0.9\paperheight}}
	{Graphs/Traversals/277.png}
\end{center}
\par Які з наведених ребер входять до складу отриманого кістякового дерева?\par
1) (1, 5)\par
2) (1, 6)\par
3) (0, 1)\par
4) (0, 5)\par
5) (0, 2)\par
6) (4, 6)\par
{\vskip 15pt} Номери відповідних ребер записати за зростанням через пробіл.


%enditem


\newpage
16. 

Рядки журналу через двокрапку містять ім'я користувача (ідентифікатор), час події,тип події, код події, наприклад
\begin{verbatim}
s="username : 2022-05-05 13:26 : ERROR : 404"
\end{verbatim}
у коді 
\begin{verbatim}
res = re.fullmatch('???', s) 
s = ''
code_str = ??? 
\end{verbatim}
чим треба замінити кожен з ???, щоб значенням code\_str був код події? 


\vspace{10mm}
\textbf{Завдання 17} 

Задано програму, що містить код 1-2 класів, можливе успадкування, та клієнтський код.
У наведеному коді невеликі частини закриті. 
Також наведено вихід програми.
Необхідно відновити закриті частини, щоб отриманий код відповідав наведеному виходу програми.


\textbf{Завдання 18-19} 

Кілька запитань за завданнями для самостійної роботи.
\vspace{10mm}
Скоріш за все, що завдань буде трохи менше, а деякі попередні будуть переформульовані в стилі двох останніх.

%enditem

\end {large}}

%end

\newpage
%begin

{{
\begin {large}

\newpage
1. Що буде надруковано за виконання?

\begin{verbatim}
print('R', end=' ')
h=('a','b','c','d','e',0,1,2)
print(h[-5])
\end{verbatim}

%enditem


\newpage
2. Що буде надруковано за виконання?

\begin{verbatim}
print('R', end=' ')
f=[0,1,2,3,4,5,6,7,8,9]
print(f[::3])
\end{verbatim}

%enditem



\newpage
3. Що буде надруковано за виконання?

\begin{verbatim}
print('R', end=' ')
c=['0','1','2','3','4']
c.extend([2,3])
print(c)
\end{verbatim}

%enditem



\newpage
4. Що буде надруковано за виконання?

\begin{verbatim}
print('R', end=' ')
e=[10,11,12,13,14]
e[-3:-4]=[4,5]
print(e)
\end{verbatim}

%enditem



\newpage
5. Укажіть ТИП значення виразу.\par
\begin{verbatim}
(11,2,3)
\end{verbatim}
%enditem


\newpage
6. Що буде надруковано за виконання?\par
\begin{verbatim}
try:
    d = {60:1, 54:1, 15:5, 42:1, 15:2}
    d[15] = 9
    for x in d.keys():
        print(x, sep=' ', end=' ')
except:
    print('ERR')
\end{verbatim}
%enditem


\newpage
7. Нехай змінна \kw{s} позначає дуже довгу послідовність, по якій можна ітерувати. 

Написати вираз-генератор, що задає підпослідовність її елементів, що не є числами типу \kw{float}.

Обмеження: усі елементи шуканої підпослідовності одночасно в пам'яті розміщені бути не можуть. 
%enditem


\newpage
8. Що буде надруковано за виконання?
%Оригинал был утерян. Требует отладки!
\begin{verbatim}
__d = 0

class C:
    __d = 1
    
    def __init__(self):
        self.__d = 2
        __d = 3
        C.__d = 4

obj = C()
C.__d = 5
try:
    print(__d, end=' ')
    print(obj.__d, end=' ')
    print(C.__d, end=' ')
    print(obj.__d, end=' ')
    print(obj._C__d)
except:
    print('ERR')
\end{verbatim}
%enditem


\newpage
9. Що буде надруковано за виконання?
\begin{verbatim}
class D0:
    def b(self): print(1, end=' ')
    def g(self): print(2, end=' ')

class D1(D0):
    def b(self): print(3, end=' ')

try:
    D1().g()
    D1().b()
except:
    print('ERR')
\end{verbatim}
%enditem


\newpage
10. Що буде надруковано за виконання?
\begin{verbatim}
class A:
    a = 1

    def _f1(self): print(2, end=' ')
    def f2(self): print(3, end=' ')

    def main(self):
        try: print(self.a, end=' ')
        except: print('E1', end=' ')

        try: self._f1()
        except: print('E2', end=' ')

        try: self.f2()
        except: print('E3', end=' ')


class B(A):
    a = 4

    def _f1(self): print(5, end=' ')
    def f2(self): print(6, end=' ')

try: B().main()
except: print('E4', end=' ')
\end{verbatim}
%enditem


\newpage
11. Змінні \kw{next} та \kw{prev} вузла списку позначають наступний та попередній за ним вузол відповідно. Починаючи з голови, у список записано послідовні цілі числа від~$-16$ до~$31$. Нехай змінна \kw{p} позначає вузол, що зберігає значення~$-5$.
Яке значення зберігається у вузлі \kw{p.next.prev.next.next.next} ?
%enditem


\newpage
12. 
\begin{center}
	\adjustimage{max size={0.8\linewidth}{0.8\paperheight}}
	{BinTree/trees_exp/122.png}
\end{center}
Указати послідовність міток вершин дерева, зображеного на рис., що утвориться в результаті обходу дерева в оберненому порядку. Мітки записувати через пробіл.\par Алгоритм починає роботу з кореня (на рис. це верхня вершина).
%enditem


\newpage
13. 
\begin{center}
	\adjustimage{max size={0.9\linewidth}{0.9\paperheight}}
	{Graphs/AdjList/207.png}
\end{center}
Для графа, зображеного на рис., вказати список суміжності вершини 3.\par
Формат оформлення відповіді:\par
список суміжності виписувати через пробіл на одному рядку, якщо пробілів десь буде більше одного - не важливо, упорядкування вершин у списку також значення не має.
%enditem


\newpage
14. Для графа на рисунку з вершини 6 запущено алгоритм пошуку в глибину, що будує кістякове дерево. Списки суміжності вершин переглядаються алгоритмом за зростанням міток вершин.\par
\begin{center}
	\adjustimage{max size={0.9\linewidth}{0.9\paperheight}}
	{Graphs/Traversals/306.png}
\end{center}
\par Які з наведених ребер входять до складу отриманого кістякового дерева?\par
1) (1, 6)\par
2) (0, 3)\par
3) (0, 5)\par
4) (3, 5)\par
5) (4, 6)\par
6) (0, 1)\par
{\vskip 15pt} Номери відповідних ребер записати за зростанням через пробіл.

%enditem


\newpage
15. Для графа на рисунку з вершини 3 запущено алгоритм пошуку в ширину, що будує кістякове дерево. Списки суміжності вершин переглядаються алгоритмом за зростанням міток вершин.\par
\begin{center}
	\adjustimage{max size={0.9\linewidth}{0.9\paperheight}}
	{Graphs/Traversals/288.png}
\end{center}
\par Які з наведених ребер входять до складу отриманого кістякового дерева?\par
1) (3, 6)\par
2) (0, 4)\par
3) (0, 5)\par
4) (0, 6)\par
5) (1, 2)\par
6) (1, 3)\par
{\vskip 15pt} Номери відповідних ребер записати за зростанням через пробіл.


%enditem


\newpage
16. 

Рядки журналу через двокрапку містять ім'я користувача (ідентифікатор), час події,тип події, код події, наприклад
\begin{verbatim}
s="username : 2022-05-05 13:26 : ERROR : 404"
\end{verbatim}
у коді 
\begin{verbatim}
res = re.fullmatch('???', s) 
s = ''
code_str = ??? 
\end{verbatim}
чим треба замінити кожен з ???, щоб значенням code\_str був код події? 


\vspace{10mm}
\textbf{Завдання 17} 

Задано програму, що містить код 1-2 класів, можливе успадкування, та клієнтський код.
У наведеному коді невеликі частини закриті. 
Також наведено вихід програми.
Необхідно відновити закриті частини, щоб отриманий код відповідав наведеному виходу програми.


\textbf{Завдання 18-19} 

Кілька запитань за завданнями для самостійної роботи.
\vspace{10mm}
Скоріш за все, що завдань буде трохи менше, а деякі попередні будуть переформульовані в стилі двох останніх.

%enditem

\end {large}}

%end

\newpage
%begin

{{
\begin {large}

\newpage
1. Що буде надруковано за виконання?

\begin{verbatim}
print('R', end=' ')
f=('a','b','c','d','e',0,1)
print(f[-9])
\end{verbatim}

%enditem


\newpage
2. Що буде надруковано за виконання?

\begin{verbatim}
print('R', end=' ')
h=['0','1','2','3','4','5','6','7','8','9']
print(h[:-10:-3])
\end{verbatim}

%enditem



\newpage
3. Що буде надруковано за виконання?

\begin{verbatim}
print('R', end=' ')
f=[10,11,12,13,14]
f.remove('12')
print(f)
\end{verbatim}

%enditem



\newpage
4. Що буде надруковано за виконання?

\begin{verbatim}
print('R', end=' ')
d=['0','1','2','3','4']
del d[8:-2]
print(d)
\end{verbatim}

%enditem



\newpage
5. Укажіть ТИП значення виразу.\par
\begin{verbatim}
{x*x:x for x in range(-7,7)}
\end{verbatim}
%enditem


\newpage
6. Що буде надруковано за виконання?\par
\begin{verbatim}
try:
    t = {10:9, 31:9, 28:8, 67:4, 67:8}
    t[28] = 9
    for x, y in t.items():
        print(x, y, sep=' ', end=' ')
except:
    print('ERR')
\end{verbatim}
%enditem


\newpage
7. Нехай змінна \kw{s} позначає послідовність цілих, по якій можна ітерувати.

Написати вираз, що перевіряє, чи правда, що всі елементи послідовності є точними квадратами.
%enditem


\newpage
8. Що буде надруковано за виконання?
\begin{verbatim}
a = 0
n = 1

class C:
    n = 2
    
    def __init__(self):
        global n
        self.a = 3
        n = 4
        C.n = 5

obj = C()
try:
    print(a, end=' ')
    print(n, end=' ')
    print(C.n, end=' ')
    print(obj.a)
except:
    print('ERR')
\end{verbatim}
%enditem


\newpage
9. Що буде надруковано за виконання?
\begin{verbatim}
class C0:
    def d(self): print(1, end=' ')
    def f(self): print(2, end=' ')

class C1(C0):
    def d(self): print(3, end=' ')

try:
    C1().f()
    C0().d()
except:
    print('ERR')
\end{verbatim}
%enditem


\newpage
10. Що буде надруковано за виконання?
\begin{verbatim}
class A:
    b = 1

    def h1(self): print(2, end=' ')
    def _h2(self): print(3, end=' ')

    def main(self):
        try: print(self.b, end=' ')
        except: print('E1', end=' ')

        try: self.h1()
        except: print('E2', end=' ')

        try: self._h2()
        except: print('E3', end=' ')


class B(A):
    b = 4

    def h1(self): print(5, end=' ')
    def _h2(self): print(6, end=' ')

try: B().main()
except: print('E4', end=' ')
\end{verbatim}
%enditem


\newpage
11. Змінні \kw{next} та \kw{prev} вузла списку позначають наступний та попередній за ним вузол відповідно. Починаючи з голови, у список записано послідовні цілі числа від~$-56$ до~$37$. Нехай змінна \kw{p} позначає вузол, що зберігає значення~$9$.
Яке значення зберігається у вузлі \kw{p.prev.next.prev.prev.prev} ?
%enditem


\newpage
12. 
\begin{center}
	\adjustimage{max size={0.8\linewidth}{0.8\paperheight}}
	{BinTree/trees_exp/157.png}
\end{center}
Указати послідовність міток вершин дерева, зображеного на рис., що утвориться в результаті обходу дерева в прямому порядку. Мітки записувати через пробіл.\par Алгоритм починає роботу з кореня (на рис. це верхня вершина).
%enditem


\newpage
13. 
\begin{center}
	\adjustimage{max size={0.9\linewidth}{0.9\paperheight}}
	{Graphs/AdjList/71.png}
\end{center}
Для графа, зображеного на рис., вказати список суміжності вершини 0.\par
Формат оформлення відповіді:\par
список суміжності виписувати через пробіл на одному рядку, якщо пробілів десь буде більше одного - не важливо, упорядкування вершин у списку також значення не має.
%enditem


\newpage
14. Для графа на рисунку вказати послідовність, в якій алгоритм пошуку в глибину буде відмічати відвідані вершини, якщо списки суміжності вершин переглядаються алгоритмом за зростанням міток вершин та вершини відмічаються на першому вход\-женні у вершину.\par
Пошук розпочинається з вершини 3.\par

\begin{center}
	\adjustimage{max size={0.9\linewidth}{0.9\paperheight}}
	{Graphs/Traversals/277.png}
\end{center}
%enditem


\newpage
15. Для графа на рисунку з вершини 6 запущено алгоритм пошуку в ширину, що будує кістякове дерево. Списки суміжності вершин переглядаються алгоритмом за зростанням міток вершин.\par
\begin{center}
	\adjustimage{max size={0.9\linewidth}{0.9\paperheight}}
	{Graphs/Traversals/91.png}
\end{center}
\par Які з наведених ребер входять до складу отриманого кістякового дерева?\par
1) (2, 4)\par
2) (0, 2)\par
3) (3, 5)\par
4) (1, 4)\par
5) (2, 5)\par
6) (3, 4)\par
{\vskip 15pt} Номери відповідних ребер записати за зростанням через пробіл.


%enditem


\newpage
16. 

Рядки журналу через двокрапку містять ім'я користувача (ідентифікатор), час події,тип події, код події, наприклад
\begin{verbatim}
s="username : 2022-05-05 13:26 : ERROR : 404"
\end{verbatim}
у коді 
\begin{verbatim}
res = re.fullmatch('???', s) 
s = ''
code_str = ??? 
\end{verbatim}
чим треба замінити кожен з ???, щоб значенням code\_str був код події? 


\vspace{10mm}
\textbf{Завдання 17} 

Задано програму, що містить код 1-2 класів, можливе успадкування, та клієнтський код.
У наведеному коді невеликі частини закриті. 
Також наведено вихід програми.
Необхідно відновити закриті частини, щоб отриманий код відповідав наведеному виходу програми.


\textbf{Завдання 18-19} 

Кілька запитань за завданнями для самостійної роботи.
\vspace{10mm}
Скоріш за все, що завдань буде трохи менше, а деякі попередні будуть переформульовані в стилі двох останніх.

%enditem

\end {large}}

%end

\newpage
%begin

{{
\begin {large}

\newpage
1. Що буде надруковано за виконання?

\begin{verbatim}
print('R', end=' ')
d=('a','b','c','d','e',0,1,2,3)
print(d[-3])
\end{verbatim}

%enditem


\newpage
2. Що буде надруковано за виконання?

\begin{verbatim}
print('R', end=' ')
a=(0,1,2,3,4,5,6,7,8,9)
print(a[-5::-1])
\end{verbatim}

%enditem



\newpage
3. Що буде надруковано за виконання?

\begin{verbatim}
print('R', end=' ')
a=['0','1','2','3','4']
print(a.pop(-5), end=' ')
print(a)
\end{verbatim}

%enditem



\newpage
4. Що буде надруковано за виконання?

\begin{verbatim}
print('R', end=' ')
f=['0','1','2','3','4']
f[7:]='12'
print(f)
\end{verbatim}

%enditem



\newpage
5. Укажіть ТИП значення виразу.\par
\begin{verbatim}
{y:None for y in range(-8,3)}
\end{verbatim}
%enditem


\newpage
6. Що буде надруковано за виконання?\par
\begin{verbatim}
try:
    s = {35:2, 65:1, 53:8, 65:4}
    s[19] = 4
    for x in s.items():
        print(x, sep=' ', end=' ')
except:
    print('ERR')
\end{verbatim}
%enditem


\newpage
7. Нехай змінна \kw{s} позначає послідовність цілих, по якій можна ітерувати. 

Написати вираз, що перевіряє, чи правда, що послідовність  містить число 13.
%enditem


\newpage
8. Що буде надруковано за виконання?
\begin{verbatim}
c = 0
j = 1

class C:
    c = 2
    
    def __init__(self):
        self.c = 3
        c = 4
        C.c = 5

obj = C()
try:
    print(c, end=' ')
    print(j, end=' ')
    print(C.c, end=' ')
    print(obj.c)
except:
    print('ERR')
\end{verbatim}
%enditem


\newpage
9. Що буде надруковано за виконання?
\begin{verbatim}
class B0:
    def a(self): print(1, end=' ')
    def h(self): print(2, end=' ')

class B1(B0):
    def h(self): print(3, end=' ')

try:
    B1().a()
    B1().h()
except:
    print('ERR')
\end{verbatim}
%enditem


\newpage
10. Що буде надруковано за виконання?
\begin{verbatim}
class A:
    c = 1

    def _g1(self): print(2, end=' ')
    def __g2(self): print(3, end=' ')

    def main(self):
        try: print(self.c, end=' ')
        except: print('E1', end=' ')

        try: self._g1()
        except: print('E2', end=' ')

        try: self.__g2()
        except: print('E3', end=' ')


class B(A):
    c = 4

    def _g1(self): print(5, end=' ')
    def __g2(self): print(6, end=' ')

try: B().main()
except: print('E4', end=' ')
\end{verbatim}
%enditem


\newpage
11. Змінні \kw{next} та \kw{prev} вузла списку позначають наступний та попередній за ним вузол відповідно. Починаючи з голови, у список записано послідовні цілі числа від~$-23$ до~$43$. Нехай змінна \kw{p} позначає вузол, що зберігає значення~$0$.
Яке значення зберігається у вузлі \kw{p.prev.next.prev.next.next} ?
%enditem


\newpage
12. 
\begin{center}
	\adjustimage{max size={0.8\linewidth}{0.8\paperheight}}
	{BinTree/trees_exp/175.png}
\end{center}
Указати послідовність міток вершин дерева, зображеного на рис., що утвориться в результаті обходу дерева в прямому порядку. Мітки записувати через пробіл.\par Алгоритм починає роботу з кореня (на рис. це верхня вершина).
%enditem


\newpage
13. 
\begin{center}
	\adjustimage{max size={0.9\linewidth}{0.9\paperheight}}
	{Graphs/AdjList/52.png}
\end{center}
Для графа, зображеного на рис., вказати список суміжності вершини 1.\par
Формат оформлення відповіді:\par
список суміжності виписувати через пробіл на одному рядку, якщо пробілів десь буде більше одного - не важливо, упорядкування вершин у списку також значення не має.
%enditem


\newpage
14. Для графа на рисунку з вершини 4 запущено алгоритм пошуку в глибину, що будує кістякове дерево. Списки суміжності вершин переглядаються алгоритмом за зростанням міток вершин.\par
\begin{center}
	\adjustimage{max size={0.9\linewidth}{0.9\paperheight}}
	{Graphs/Traversals/98.png}
\end{center}
\par Які з наведених ребер входять до складу отриманого кістякового дерева?\par
1) (0, 6)\par
2) (1, 2)\par
3) (4, 6)\par
4) (1, 4)\par
5) (2, 4)\par
6) (3, 4)\par
{\vskip 15pt} Номери відповідних ребер записати за зростанням через пробіл.

%enditem


\newpage
15. Для графа на рисунку з вершини 0 запущено алгоритм пошуку в ширину, що будує кістякове дерево. Списки суміжності вершин переглядаються алгоритмом за зростанням міток вершин.\par
\begin{center}
	\adjustimage{max size={0.9\linewidth}{0.9\paperheight}}
	{Graphs/Traversals/107.png}
\end{center}
\par Які з наведених ребер входять до складу отриманого кістякового дерева?\par
1) (3, 6)\par
2) (2, 5)\par
3) (2, 6)\par
4) (0, 1)\par
5) (3, 4)\par
6) (3, 5)\par
{\vskip 15pt} Номери відповідних ребер записати за зростанням через пробіл.


%enditem


\newpage
16. 

Рядки журналу через двокрапку містять ім'я користувача (ідентифікатор), час події,тип події, код події, наприклад
\begin{verbatim}
s="username : 2022-05-05 13:26 : ERROR : 404"
\end{verbatim}
у коді 
\begin{verbatim}
res = re.fullmatch('???', s) 
s = ''
code_str = ??? 
\end{verbatim}
чим треба замінити кожен з ???, щоб значенням code\_str був код події? 


\vspace{10mm}
\textbf{Завдання 17} 

Задано програму, що містить код 1-2 класів, можливе успадкування, та клієнтський код.
У наведеному коді невеликі частини закриті. 
Також наведено вихід програми.
Необхідно відновити закриті частини, щоб отриманий код відповідав наведеному виходу програми.


\textbf{Завдання 18-19} 

Кілька запитань за завданнями для самостійної роботи.
\vspace{10mm}
Скоріш за все, що завдань буде трохи менше, а деякі попередні будуть переформульовані в стилі двох останніх.

%enditem

\end {large}}

%end

\newpage
%begin

{{
\begin {large}

\newpage
1. Що буде надруковано за виконання?

\begin{verbatim}
print('R', end=' ')
e=('a','b','c','d','e',0,1,2,3,4)
print(e[3])
\end{verbatim}

%enditem


\newpage
2. Що буде надруковано за виконання?

\begin{verbatim}
print('R', end=' ')
c=[0,1,2,3,4,5,6,7,8,9]
print(c[8:7:-1])
\end{verbatim}

%enditem



\newpage
3. Що буде надруковано за виконання?

\begin{verbatim}
print('R', end=' ')
g=[10,11,12,13,14]
g+='12'
print(g)
\end{verbatim}

%enditem



\newpage
4. Що буде надруковано за виконання?

\begin{verbatim}
print('R', end=' ')
g=[10,11,12,13,14]
del g[:6]
print(g)
\end{verbatim}

%enditem



\newpage
5. Укажіть ТИП значення виразу.\par
\begin{verbatim}
{x for x in range(7) if el%3==1}
\end{verbatim}
%enditem


\newpage
6. Що буде надруковано за виконання?\par
\begin{verbatim}
try:
    s = {59:9, 61:6, 48:2, 34:5, 59:8}
    s[27] = 1
    for x in s :
        print(x, sep=' ', end=' ')
except:
    print('ERR')
\end{verbatim}
%enditem


\newpage
7. 
Записати ВИРАЗ, значенням якого є список, що складається з квадратів цілих чисел від 14 до 2013 (включно), що не діляться на жодне з чисел 2, 3.
%enditem


\newpage
8. Що буде надруковано за виконання?
\begin{verbatim}
b = 0
m = 1

class D:
    m = 2
    
    def __init__(self):
        global m
        self.m = 3
        m = 4
        D.m = 5

obj = D()
try:
    print(b, end=' ')
    print(m, end=' ')
    print(D.m, end=' ')
    print(obj.m)
except:
    print('ERR')
\end{verbatim}
%enditem


\newpage
9. Що буде надруковано за виконання?
\begin{verbatim}
class A0:
    def c(self): print(1, end=' ')
    def h(self): print(2, end=' ')

class A1(A0):
    def h(self): print(3, end=' ')

try:
    A0().c()
    A1().h()
except:
    print('ERR')
\end{verbatim}
%enditem


\newpage
10. Що буде надруковано за виконання?
\begin{verbatim}
class A:
    a = 1

    def __init__(self):
        b = 2
        a = 3

class B(A):
    b = 4
    c = 5

    def __init__(self):
        super().__init__()
        self.c = 6

try:
    obj = B()
    obj.c = obj.b + 10
    B.c = 13
except: print('EE', end=' ')

try: print(obj.c, end= ' ')
except: print('E1', end=' ')

try: print(A.c, end= ' ')
except: print('E2', end=' ')

try: print(B.c, end= ' ')
except: print('E3', end=' ')
\end{verbatim}
%enditem


\newpage
11. Змінні \kw{next} та \kw{prev} вузла списку позначають наступний та попередній за ним вузол відповідно. Починаючи з голови, у список записано послідовні цілі числа від~$-49$ до~$50$. Нехай змінна \kw{p} позначає вузол, що зберігає значення~$-1$.
Яке значення зберігається у вузлі \kw{p.next.prev.next.prev.prev} ?
%enditem


\newpage
12. 
\begin{center}
	\adjustimage{max size={0.8\linewidth}{0.8\paperheight}}
	{BinTree/trees_exp/84.png}
\end{center}
Указати послідовність міток вершин дерева, зображеного на рис., що утвориться в результаті обходу дерева в оберненому порядку. Мітки записувати через пробіл.\par Алгоритм починає роботу з кореня (на рис. це верхня вершина).
%enditem


\newpage
13. 
\begin{center}
	\adjustimage{max size={0.9\linewidth}{0.9\paperheight}}
	{Graphs/AdjList/128.png}
\end{center}
Для графа, зображеного на рис., вказати список суміжності вершини 3.\par
Формат оформлення відповіді:\par
список суміжності виписувати через пробіл на одному рядку, якщо пробілів десь буде більше одного - не важливо, упорядкування вершин у списку також значення не має.
%enditem


\newpage
14. Для графа на рисунку вказати послідовність, в якій алгоритм пошуку в глибину буде відмічати відвідані вершини, якщо списки суміжності вершин переглядаються алгоритмом за зростанням міток вершин та вершини відмічаються на першому вход\-женні у вершину.\par
Пошук розпочинається з вершини 5.\par

\begin{center}
	\adjustimage{max size={0.9\linewidth}{0.9\paperheight}}
	{Graphs/Traversals/262.png}
\end{center}
%enditem


\newpage
15. Для графа на рисунку з вершини 3 запущено алгоритм пошуку в ширину, що будує кістякове дерево. Списки суміжності вершин переглядаються алгоритмом за зростанням міток вершин.\par
\begin{center}
	\adjustimage{max size={0.9\linewidth}{0.9\paperheight}}
	{Graphs/Traversals/334.png}
\end{center}
\par Які з наведених ребер входять до складу отриманого кістякового дерева?\par
1) (2, 5)\par
2) (0, 1)\par
3) (0, 6)\par
4) (1, 2)\par
5) (2, 6)\par
6) (1, 4)\par
{\vskip 15pt} Номери відповідних ребер записати за зростанням через пробіл.


%enditem


\newpage
16. 

Рядки журналу через двокрапку містять ім'я користувача (ідентифікатор), час події,тип події, код події, наприклад
\begin{verbatim}
s="username : 2022-05-05 13:26 : ERROR : 404"
\end{verbatim}
у коді 
\begin{verbatim}
res = re.fullmatch('???', s) 
s = ''
code_str = ??? 
\end{verbatim}
чим треба замінити кожен з ???, щоб значенням code\_str був код події? 


\vspace{10mm}
\textbf{Завдання 17} 

Задано програму, що містить код 1-2 класів, можливе успадкування, та клієнтський код.
У наведеному коді невеликі частини закриті. 
Також наведено вихід програми.
Необхідно відновити закриті частини, щоб отриманий код відповідав наведеному виходу програми.


\textbf{Завдання 18-19} 

Кілька запитань за завданнями для самостійної роботи.
\vspace{10mm}
Скоріш за все, що завдань буде трохи менше, а деякі попередні будуть переформульовані в стилі двох останніх.

%enditem

\end {large}}

%end

\newpage
%begin

{{
\begin {large}

\newpage
1. Що буде надруковано за виконання?

\begin{verbatim}
print('R', end=' ')
g=('a','b','c','d','e',0,1,2,4)
print(g[-8])
\end{verbatim}

%enditem


\newpage
2. Що буде надруковано за виконання?

\begin{verbatim}
print('R', end=' ')
d='0123456789'
print(d[-10:-6:-3])
\end{verbatim}

%enditem



\newpage
3. Що буде надруковано за виконання?

\begin{verbatim}
print('R', end=' ')
b=[10,11,12,13,14]
b.insert(-7, '10')
print(b)
\end{verbatim}

%enditem



\newpage
4. Що буде надруковано за виконання?

\begin{verbatim}
print('R', end=' ')
c=[10,11,12,13,14]
c[-8:-6]=(14)
print(c)
\end{verbatim}

%enditem



\newpage
5. Укажіть ТИП значення виразу.\par
\begin{verbatim}
[x for x in range(7)]
\end{verbatim}
%enditem


\newpage
6. Що буде надруковано за виконання?\par
\begin{verbatim}
try:
    t = {90:2, 82:5, 72:8, 46:7, 90:0}
    t[46] = 0
    for x in t.values():
        print(x, sep=' ', end=' ')
except:
    print('ERR')
\end{verbatim}
%enditem


\newpage
7. Задано словник \kw{d} з цілими ключами та значеннями. Записати ВИРАЗ, значенням якого є множина, що складається  з тих значень, чиї ключі більші за 3.
%enditem


\newpage
8. Що буде надруковано за виконання?
\begin{verbatim}
a = 0
n = 1

class B:
    n = 2
    
    def __init__(self):
        global n
        self.a = 3
        n = 4
        B.n = 5

obj = B()
try:
    print(a, end=' ')
    print(n, end=' ')
    print(B.n, end=' ')
    print(obj.a)
except:
    print('ERR')
\end{verbatim}
%enditem


\newpage
9. Що буде надруковано за виконання?
\begin{verbatim}
class B0:
    def a(self): print(1, end=' ')
    def g(self): print(2, end=' ')

class B1(B0):
    def a(self): print(3, end=' ')

try:
    B0().g()
    B1().a()
except:
    print('ERR')
\end{verbatim}
%enditem


\newpage
10. Що буде надруковано за виконання?
\begin{verbatim}
class A:
    b = 1

    def __init__(self):
        c = 2
        c = 3

class B(A):
    b = 4
    a = 5

    def __init__(self):
        super().__init__()
        self.a = 6

try:
    obj = B()
    obj.a = obj.b + 10
    B.a = 13
except: print('EE', end=' ')

try: print(obj.a, end= ' ')
except: print('E1', end=' ')

try: print(A.a, end= ' ')
except: print('E2', end=' ')

try: print(B.a, end= ' ')
except: print('E3', end=' ')
\end{verbatim}
%enditem


\newpage
11. Змінні \kw{next} та \kw{prev} вузла списку позначають наступний та попередній за ним вузол відповідно. Починаючи з голови, у список записано послідовні цілі числа від~$-32$ до~$42$. Нехай змінна \kw{p} позначає вузол, що зберігає значення~$-4$.
Яке значення зберігається у вузлі \kw{p.next.next.next.prev.prev} ?
%enditem


\newpage
12. 
\begin{center}
	\adjustimage{max size={0.8\linewidth}{0.8\paperheight}}
	{BinTree/trees_exp/54.png}
\end{center}
Указати послідовність міток вершин дерева, зображеного на рис., що утвориться в результаті обходу дерева в прямому порядку. Мітки записувати через пробіл.\par Алгоритм починає роботу з кореня (на рис. це верхня вершина).
%enditem


\newpage
13. 
\begin{center}
	\adjustimage{max size={0.9\linewidth}{0.9\paperheight}}
	{Graphs/AdjList/149.png}
\end{center}
Для графа, зображеного на рис., вказати список суміжності вершини 0.\par
Формат оформлення відповіді:\par
список суміжності виписувати через пробіл на одному рядку, якщо пробілів десь буде більше одного - не важливо, упорядкування вершин у списку також значення не має.
%enditem


\newpage
14. Для графа на рисунку вказати послідовність, в якій алгоритм пошуку в глибину буде відмічати відвідані вершини, якщо списки суміжності вершин переглядаються алгоритмом за зростанням міток вершин та вершини відмічаються на першому вход\-женні у вершину.\par
Пошук розпочинається з вершини 1.\par

\begin{center}
	\adjustimage{max size={0.9\linewidth}{0.9\paperheight}}
	{Graphs/Traversals/74.png}
\end{center}
%enditem


\newpage
15. Для графа на рисунку з вершини 6 запущено алгоритм пошуку в ширину, що будує кістякове дерево. Списки суміжності вершин переглядаються алгоритмом за зростанням міток вершин.\par
\begin{center}
	\adjustimage{max size={0.9\linewidth}{0.9\paperheight}}
	{Graphs/Traversals/313.png}
\end{center}
\par Які з наведених ребер входять до складу отриманого кістякового дерева?\par
1) (5, 6)\par
2) (0, 5)\par
3) (2, 4)\par
4) (0, 4)\par
5) (3, 6)\par
6) (1, 6)\par
{\vskip 15pt} Номери відповідних ребер записати за зростанням через пробіл.


%enditem


\newpage
16. 

Рядки журналу через двокрапку містять ім'я користувача (ідентифікатор), час події,тип події, код події, наприклад
\begin{verbatim}
s="username : 2022-05-05 13:26 : ERROR : 404"
\end{verbatim}
у коді 
\begin{verbatim}
res = re.fullmatch('???', s) 
s = ''
code_str = ??? 
\end{verbatim}
чим треба замінити кожен з ???, щоб значенням code\_str був код події? 


\vspace{10mm}
\textbf{Завдання 17} 

Задано програму, що містить код 1-2 класів, можливе успадкування, та клієнтський код.
У наведеному коді невеликі частини закриті. 
Також наведено вихід програми.
Необхідно відновити закриті частини, щоб отриманий код відповідав наведеному виходу програми.


\textbf{Завдання 18-19} 

Кілька запитань за завданнями для самостійної роботи.
\vspace{10mm}
Скоріш за все, що завдань буде трохи менше, а деякі попередні будуть переформульовані в стилі двох останніх.

%enditem

\end {large}}

%end

\newpage
%begin

{{
\begin {large}

\newpage
1. Що буде надруковано за виконання?

\begin{verbatim}
print('R', end=' ')
b=('a','b','c','d','e',0,1,2,3)
print(b[13])
\end{verbatim}

%enditem


\newpage
2. Що буде надруковано за виконання?

\begin{verbatim}
print('R', end=' ')
g=('0','1','2','3','4','5','6','7','8','9')
print(g[-12::-2])
\end{verbatim}

%enditem



\newpage
3. Що буде надруковано за виконання?

\begin{verbatim}
print('R', end=' ')
d=['0','1','2','3','4']
d.remove('9')
print(d)
\end{verbatim}

%enditem



\newpage
4. Що буде надруковано за виконання?

\begin{verbatim}
print('R', end=' ')
a=['0','1','2','3','4']
a[5:7]=[2,3]
print(a)
\end{verbatim}

%enditem



\newpage
5. Укажіть ТИП значення виразу.\par
\begin{verbatim}
(2, 7, -4)
\end{verbatim}
%enditem


\newpage
6. Що буде надруковано за виконання?\par
\begin{verbatim}
try:
    d = {70:9, 90:7, 72:1, 72:6}
    d[72] = 0
    for x in d.items():
        print(*x, sep=' ', end=' ')
except:
    print('ERR')
\end{verbatim}
%enditem


\newpage
7. 
Записати ВИРАЗ, значенням якого є список, що складається з цілих чисел від 16 до 2015 (включно), причому спочатку за спаданням записано всі, що не діляться на 2, а потім решту за спаданням.

%enditem


\newpage
8. Що буде надруковано за виконання?
%Оригинал был утерян. Требует отладки!
\begin{verbatim}
a = 0

class A:
    a = 1
    
    def __init__(self):
        self.a = 2
        a = 3
        A.a = 4

obj = A()
A.a = 5
try:
    print(a, end=' ')
    print(obj.a, end=' ')
    print(A.a, end=' ')
    print(obj.a, end=' ')
    print(obj._A__a)
except:
    print('ERR')
\end{verbatim}
%enditem


\newpage
9. Що буде надруковано за виконання?
\begin{verbatim}
class C0:
    def c(self): print(1, end=' ')
    def f(self): print(2, end=' ')

class C1(C0):
    def c(self): print(3, end=' ')

try:
    C1().f()
    C1().c()
except:
    print('ERR')
\end{verbatim}
%enditem


\newpage
10. Що буде надруковано за виконання?
\begin{verbatim}
class A:
    c = 1

    def __init__(self):
        b = 2
        c = 3

class B(A):
    b = 4
    a = 5

    def __init__(self):
        super().__init__()
        self.a = 6

try:
    obj = B()
    obj.a = obj.c + 10
    B.a = 13
except: print('EE', end=' ')

try: print(obj.a, end= ' ')
except: print('E1', end=' ')

try: print(A.a, end= ' ')
except: print('E2', end=' ')

try: print(B.a, end= ' ')
except: print('E3', end=' ')
\end{verbatim}
%enditem


\newpage
11. Змінні \kw{next} та \kw{prev} вузла списку позначають наступний та попередній за ним вузол відповідно. Починаючи з голови, у список записано послідовні цілі числа від~$-64$ до~$33$. Нехай змінна \kw{p} позначає вузол, що зберігає значення~$-5$.
Яке значення зберігається у вузлі \kw{p.prev.prev.prev.next.next} ?
%enditem


\newpage
12. 
\begin{center}
	\adjustimage{max size={0.8\linewidth}{0.8\paperheight}}
	{BinTree/trees_exp/11.png}
\end{center}
Указати послідовність міток вершин дерева, зображеного на рис., що утвориться в результаті обходу дерева в оберненому порядку. Мітки записувати через пробіл.\par Алгоритм починає роботу з кореня (на рис. це верхня вершина).
%enditem


\newpage
13. 
\begin{center}
	\adjustimage{max size={0.9\linewidth}{0.9\paperheight}}
	{Graphs/AdjList/9.png}
\end{center}
Для графа, зображеного на рис., вказати список суміжності вершини 7.\par
Формат оформлення відповіді:\par
список суміжності виписувати через пробіл на одному рядку, якщо пробілів десь буде більше одного - не важливо, упорядкування вершин у списку також значення не має.
%enditem


\newpage
14. Для графа на рисунку з вершини 4 запущено алгоритм пошуку в глибину, що будує кістякове дерево. Списки суміжності вершин переглядаються алгоритмом за зростанням міток вершин.\par
\begin{center}
	\adjustimage{max size={0.9\linewidth}{0.9\paperheight}}
	{Graphs/Traversals/344.png}
\end{center}
\par Які з наведених ребер входять до складу отриманого кістякового дерева?\par
1) (3, 5)\par
2) (1, 4)\par
3) (0, 5)\par
4) (4, 6)\par
5) (0, 4)\par
6) (5, 6)\par
{\vskip 15pt} Номери відповідних ребер записати за зростанням через пробіл.

%enditem


\newpage
15. Для графа на рисунку з вершини 2 запущено алгоритм пошуку в ширину, що будує кістякове дерево. Списки суміжності вершин переглядаються алгоритмом за зростанням міток вершин.\par
\begin{center}
	\adjustimage{max size={0.9\linewidth}{0.9\paperheight}}
	{Graphs/Traversals/103.png}
\end{center}
\par Які з наведених ребер входять до складу отриманого кістякового дерева?\par
1) (0, 6)\par
2) (3, 4)\par
3) (2, 5)\par
4) (5, 6)\par
5) (2, 3)\par
6) (0, 1)\par
{\vskip 15pt} Номери відповідних ребер записати за зростанням через пробіл.


%enditem


\newpage
16. 

Рядки журналу через двокрапку містять ім'я користувача (ідентифікатор), час події,тип події, код події, наприклад
\begin{verbatim}
s="username : 2022-05-05 13:26 : ERROR : 404"
\end{verbatim}
у коді 
\begin{verbatim}
res = re.fullmatch('???', s) 
s = ''
code_str = ??? 
\end{verbatim}
чим треба замінити кожен з ???, щоб значенням code\_str був код події? 


\vspace{10mm}
\textbf{Завдання 17} 

Задано програму, що містить код 1-2 класів, можливе успадкування, та клієнтський код.
У наведеному коді невеликі частини закриті. 
Також наведено вихід програми.
Необхідно відновити закриті частини, щоб отриманий код відповідав наведеному виходу програми.


\textbf{Завдання 18-19} 

Кілька запитань за завданнями для самостійної роботи.
\vspace{10mm}
Скоріш за все, що завдань буде трохи менше, а деякі попередні будуть переформульовані в стилі двох останніх.

%enditem

\end {large}}

%end

\newpage
%begin

{{
\begin {large}

\newpage
1. Що буде надруковано за виконання?

\begin{verbatim}
print('R', end=' ')
a=('a','b','c','d','e',0,1)
print(a[-2])
\end{verbatim}

%enditem


\newpage
2. Що буде надруковано за виконання?

\begin{verbatim}
print('R', end=' ')
b=(0,1,2,3,4,5,6,7,8,9)
print(b[:-15:2])
\end{verbatim}

%enditem



\newpage
3. Що буде надруковано за виконання?

\begin{verbatim}
print('R', end=' ')
h=['0','1','2','3','4']
print(h.index('1'), end=' ')
\end{verbatim}

%enditem



\newpage
4. Що буде надруковано за виконання?

\begin{verbatim}
print('R', end=' ')
b=[10,11,12,13,14]
del b[3:]
print(b)
\end{verbatim}

%enditem



\newpage
5. Укажіть ТИП значення виразу.\par
\begin{verbatim}
[{1,2}, 5]
\end{verbatim}
%enditem


\newpage
6. Що буде надруковано за виконання?\par
\begin{verbatim}
try:
    d = {55:2, 71:2, 73:4, 12:4, 73:8}
    d[18] = 8
    for x in d :
        print(x, sep=' ', end=' ')
except:
    print('ERR')
\end{verbatim}
%enditem


\newpage
7. Нехай змінна \kw{s} позначає послідовність цілих, по якій можна ітерувати.

Написати вираз, що перевіряє, чи правда, що всі елементи послідовності є додатни\-ми.
%enditem


\newpage
8. Що буде надруковано за виконання?
\begin{verbatim}
d = 0
k = 1

class A:
    d = 2
    
    def __init__(self):
        self.k = 3
        d = 4
        A.d = 5

obj = A()
try:
    print(d, end=' ')
    print(k, end=' ')
    print(A.d, end=' ')
    print(obj.k)
except:
    print('ERR')
\end{verbatim}
%enditem


\newpage
9. Що буде надруковано за виконання?
\begin{verbatim}
class D0:
    def d(self): print(1, end=' ')
    def h(self): print(2, end=' ')

class D1(D0):
    def h(self): print(3, end=' ')

try:
    D1().d()
    D1().h()
except:
    print('ERR')
\end{verbatim}
%enditem


\newpage
10. Що буде надруковано за виконання?
\begin{verbatim}
class A:
    a = 1

    def __g1(self): print(2, end=' ')
    def _g2(self): print(3, end=' ')

    def main(self):
        try: print(self.a, end=' ')
        except: print('E1', end=' ')

        try: self.__g1()
        except: print('E2', end=' ')

        try: self._g2()
        except: print('E3', end=' ')


class B(A):
    a = 4

    def __g1(self): print(5, end=' ')
    def _g2(self): print(6, end=' ')

try: B().main()
except: print('E4', end=' ')
\end{verbatim}
%enditem


\newpage
11. Змінні \kw{next} та \kw{prev} вузла списку позначають наступний та попередній за ним вузол відповідно. Починаючи з голови, у список записано послідовні цілі числа від~$-30$ до~$27$. Нехай змінна \kw{p} позначає вузол, що зберігає значення~$-6$.
Яке значення зберігається у вузлі \kw{p.prev.next.prev.next.next} ?
%enditem


\newpage
12. 
\begin{center}
	\adjustimage{max size={0.8\linewidth}{0.8\paperheight}}
	{BinTree/trees_exp/95.png}
\end{center}
Указати послідовність міток вершин дерева, зображеного на рис., що утвориться в результаті обходу дерева в оберненому порядку. Мітки записувати через пробіл.\par Алгоритм починає роботу з кореня (на рис. це верхня вершина).
%enditem


\newpage
13. 
\begin{center}
	\adjustimage{max size={0.9\linewidth}{0.9\paperheight}}
	{Graphs/AdjList/308.png}
\end{center}
Для графа, зображеного на рис., вказати список суміжності вершини 0.\par
Формат оформлення відповіді:\par
список суміжності виписувати через пробіл на одному рядку, якщо пробілів десь буде більше одного - не важливо, упорядкування вершин у списку також значення не має.
%enditem


\newpage
14. Для графа на рисунку вказати послідовність, в якій алгоритм пошуку в глибину буде відмічати відвідані вершини, якщо списки суміжності вершин переглядаються алгоритмом за зростанням міток вершин та вершини відмічаються на першому вход\-женні у вершину.\par
Пошук розпочинається з вершини 4.\par

\begin{center}
	\adjustimage{max size={0.9\linewidth}{0.9\paperheight}}
	{Graphs/Traversals/225.png}
\end{center}
%enditem


\newpage
15. Для графа на рисунку з вершини 6 запущено алгоритм пошуку в ширину, що будує кістякове дерево. Списки суміжності вершин переглядаються алгоритмом за зростанням міток вершин.\par
\begin{center}
	\adjustimage{max size={0.9\linewidth}{0.9\paperheight}}
	{Graphs/Traversals/127.png}
\end{center}
\par Які з наведених ребер входять до складу отриманого кістякового дерева?\par
1) (0, 4)\par
2) (0, 1)\par
3) (1, 2)\par
4) (1, 3)\par
5) (0, 3)\par
6) (3, 4)\par
{\vskip 15pt} Номери відповідних ребер записати за зростанням через пробіл.


%enditem


\newpage
16. 

Рядки журналу через двокрапку містять ім'я користувача (ідентифікатор), час події,тип події, код події, наприклад
\begin{verbatim}
s="username : 2022-05-05 13:26 : ERROR : 404"
\end{verbatim}
у коді 
\begin{verbatim}
res = re.fullmatch('???', s) 
s = ''
code_str = ??? 
\end{verbatim}
чим треба замінити кожен з ???, щоб значенням code\_str був код події? 


\vspace{10mm}
\textbf{Завдання 17} 

Задано програму, що містить код 1-2 класів, можливе успадкування, та клієнтський код.
У наведеному коді невеликі частини закриті. 
Також наведено вихід програми.
Необхідно відновити закриті частини, щоб отриманий код відповідав наведеному виходу програми.


\textbf{Завдання 18-19} 

Кілька запитань за завданнями для самостійної роботи.
\vspace{10mm}
Скоріш за все, що завдань буде трохи менше, а деякі попередні будуть переформульовані в стилі двох останніх.

%enditem

\end {large}}

%end

\newpage
%begin

{{
\begin {large}

\newpage
1. Що буде надруковано за виконання?

\begin{verbatim}
print('R', end=' ')
c=('a','b','c','d','e',0,1,2)
print(c[6])
\end{verbatim}

%enditem


\newpage
2. Що буде надруковано за виконання?

\begin{verbatim}
print('R', end=' ')
e=['0','1','2','3','4','5','6','7','8','9']
print(e[3:-8:3])
\end{verbatim}

%enditem



\newpage
3. Що буде надруковано за виконання?

\begin{verbatim}
print('R', end=' ')
e=[10,11,12,13,14]
e.insert(-3, [1,0])
print(e)
\end{verbatim}

%enditem



\newpage
4. Що буде надруковано за виконання?

\begin{verbatim}
print('R', end=' ')
h=['0','1','2','3','4']
h[:2]=[13]
print(h)
\end{verbatim}

%enditem



\newpage
5. Укажіть ТИП значення виразу.\par
\begin{verbatim}
{x:0 for x in range(7)}
\end{verbatim}
%enditem


\newpage
6. Що буде надруковано за виконання?\par
\begin{verbatim}
try:
    s = {26:5, 39:3, 76:3, 26:5}
    s[43] = 6
    for x in s.items():
        print(*x, sep=' ', end=' ')
except:
    print('ERR')
\end{verbatim}
%enditem


\newpage
7. Нехай змінна \kw{s} позначає послідовність цілих, по якій можна ітерувати.

Написати вираз, що перевіряє, чи правда, що всі елементи послідовності більші за задане число num.
%enditem


\newpage
8. Що буде надруковано за виконання?
%Оригинал был утерян. Требует отладки!
\begin{verbatim}
_b = 0

class B:
    _b = 1
    
    def __init__(self):
        self._b = 2
        _b = 3
        B._b = 4

obj = B()
obj._b = 5
try:
    print(_b, end=' ')
    print(obj._b, end=' ')
    print(B._b, end=' ')
    print(obj._b, end=' ')
    print(obj._B__b)
except:
    print('ERR')
\end{verbatim}
%enditem


\newpage
9. Що буде надруковано за виконання?
\begin{verbatim}
class A0:
    def b(self): print(1, end=' ')
    def f(self): print(2, end=' ')

class A1(A0):
    def f(self): print(3, end=' ')

try:
    A1().b()
    A0().f()
except:
    print('ERR')
\end{verbatim}
%enditem


\newpage
10. Що буде надруковано за виконання?
\begin{verbatim}
class A:
    a = 1

    def __init__(self):
        b = 2
        a = 3

class B(A):
    b = 4
    c = 5

    def __init__(self):
        super().__init__()
        self.c = 6

try:
    obj = B()
    obj.b = obj.c + 10
    B.b = 13
except: print('EE', end=' ')

try: print(obj.b, end= ' ')
except: print('E1', end=' ')

try: print(A.b, end= ' ')
except: print('E2', end=' ')

try: print(B.b, end= ' ')
except: print('E3', end=' ')
\end{verbatim}
%enditem


\newpage
11. Змінні \kw{next} та \kw{prev} вузла списку позначають наступний та попередній за ним вузол відповідно. Починаючи з голови, у список записано послідовні цілі числа від~$-47$ до~$35$. Нехай змінна \kw{p} позначає вузол, що зберігає значення~$1$.
Яке значення зберігається у вузлі \kw{p.next.prev.next.prev.prev} ?
%enditem


\newpage
12. 
\begin{center}
	\adjustimage{max size={0.8\linewidth}{0.8\paperheight}}
	{BinTree/trees_exp/21.png}
\end{center}
Указати послідовність міток вершин дерева, зображеного на рис., що утвориться в результаті обходу дерева в прямому порядку. Мітки записувати через пробіл.\par Алгоритм починає роботу з кореня (на рис. це верхня вершина).
%enditem


\newpage
13. 
\begin{center}
	\adjustimage{max size={0.9\linewidth}{0.9\paperheight}}
	{Graphs/AdjList/256.png}
\end{center}
Для графа, зображеного на рис., вказати список суміжності вершини 1.\par
Формат оформлення відповіді:\par
список суміжності виписувати через пробіл на одному рядку, якщо пробілів десь буде більше одного - не важливо, упорядкування вершин у списку також значення не має.
%enditem


\newpage
14. Для графа на рисунку вказати послідовність, в якій алгоритм пошуку в глибину буде відмічати відвідані вершини, якщо списки суміжності вершин переглядаються алгоритмом за зростанням міток вершин та вершини відмічаються на першому вход\-женні у вершину.\par
Пошук розпочинається з вершини 1.\par

\begin{center}
	\adjustimage{max size={0.9\linewidth}{0.9\paperheight}}
	{Graphs/Traversals/241.png}
\end{center}
%enditem


\newpage
15. Для графа на рисунку з вершини 1 запущено алгоритм пошуку в ширину, що будує кістякове дерево. Списки суміжності вершин переглядаються алгоритмом за зростанням міток вершин.\par
\begin{center}
	\adjustimage{max size={0.9\linewidth}{0.9\paperheight}}
	{Graphs/Traversals/183.png}
\end{center}
\par Які з наведених ребер входять до складу отриманого кістякового дерева?\par
1) (0, 4)\par
2) (5, 6)\par
3) (1, 3)\par
4) (0, 6)\par
5) (1, 2)\par
6) (0, 3)\par
{\vskip 15pt} Номери відповідних ребер записати за зростанням через пробіл.


%enditem


\newpage
16. 

Рядки журналу через двокрапку містять ім'я користувача (ідентифікатор), час події,тип події, код події, наприклад
\begin{verbatim}
s="username : 2022-05-05 13:26 : ERROR : 404"
\end{verbatim}
у коді 
\begin{verbatim}
res = re.fullmatch('???', s) 
s = ''
code_str = ??? 
\end{verbatim}
чим треба замінити кожен з ???, щоб значенням code\_str був код події? 


\vspace{10mm}
\textbf{Завдання 17} 

Задано програму, що містить код 1-2 класів, можливе успадкування, та клієнтський код.
У наведеному коді невеликі частини закриті. 
Також наведено вихід програми.
Необхідно відновити закриті частини, щоб отриманий код відповідав наведеному виходу програми.


\textbf{Завдання 18-19} 

Кілька запитань за завданнями для самостійної роботи.
\vspace{10mm}
Скоріш за все, що завдань буде трохи менше, а деякі попередні будуть переформульовані в стилі двох останніх.

%enditem

\end {large}}

%end

\newpage
%begin

{{
\begin {large}

\newpage
1. Що буде надруковано за виконання?

\begin{verbatim}
print('R', end=' ')
b=('a','b','c','d','e',0,1,2,4)
print(b[-10])
\end{verbatim}

%enditem


\newpage
2. Що буде надруковано за виконання?

\begin{verbatim}
print('R', end=' ')
g=(0,1,2,3,4,5,6,7,8,9)
print(g[14:4:-1])
\end{verbatim}

%enditem



\newpage
3. Що буде надруковано за виконання?

\begin{verbatim}
print('R', end=' ')
d=[10,11,12,13,14]
print(d.pop(-1), end=' ')
print(d)
\end{verbatim}

%enditem



\newpage
4. Що буде надруковано за виконання?

\begin{verbatim}
print('R', end=' ')
d=[10,11,12,13,14]
d[-1:-5]=[1,0]
print(d)
\end{verbatim}

%enditem



\newpage
5. Укажіть ТИП значення виразу.\par
\begin{verbatim}
{21,2,4}
\end{verbatim}
%enditem


\newpage
6. Що буде надруковано за виконання?\par
\begin{verbatim}
try:
    t = {58:2, 31:2, 90:2, 90:9}
    t[58] = 9
    for x in t.keys():
        print(x, sep=' ', end=' ')
except:
    print('ERR')
\end{verbatim}
%enditem


\newpage
7. Нехай змінна \kw{s} позначає дуже довгу послідовність цілих, по якій можна ітерувати. 

Написати вираз-генератор, що задає підпослідовність її елементів, що є точними квадратами. 

Обмеження: усі елементи шуканої підпослідовності одночасно в пам'яті розміщені бути не можуть. 
%enditem


\newpage
8. Що буде надруковано за виконання?
%Оригинал был утерян. Требует отладки!
\begin{verbatim}
b = 0

class C:
    b = 1
    
    def __init__(self):
        self.b = 2
        b = 3
        C.b = 4

obj = C()
obj.b = 5
try:
    print(b, end=' ')
    print(obj.b, end=' ')
    print(C.b, end=' ')
    print(obj.b, end=' ')
    print(obj._C__b)
except:
    print('ERR')
\end{verbatim}
%enditem


\newpage
9. Що буде надруковано за виконання?
\begin{verbatim}
class C0:
    def a(self): print(1, end=' ')
    def g(self): print(2, end=' ')

class C1(C0):
    def a(self): print(3, end=' ')

try:
    C1().g()
    C0().a()
except:
    print('ERR')
\end{verbatim}
%enditem


\newpage
10. Що буде надруковано за виконання?
\begin{verbatim}
class A:
    d = 1

    def _h1(self): print(2, end=' ')
    def _h2(self): print(3, end=' ')

    def main(self):
        try: print(self.d, end=' ')
        except: print('E1', end=' ')

        try: self._h1()
        except: print('E2', end=' ')

        try: self._h2()
        except: print('E3', end=' ')


class B(A):
    d = 4

    def _h1(self): print(5, end=' ')
    def _h2(self): print(6, end=' ')

try: B().main()
except: print('E4', end=' ')
\end{verbatim}
%enditem


\newpage
11. Змінні \kw{next} та \kw{prev} вузла списку позначають наступний та попередній за ним вузол відповідно. Починаючи з голови, у список записано послідовні цілі числа від~$-19$ до~$26$. Нехай змінна \kw{p} позначає вузол, що зберігає значення~$-8$.
Яке значення зберігається у вузлі \kw{p.next.next.prev.next.next} ?
%enditem


\newpage
12. 
\begin{center}
	\adjustimage{max size={0.8\linewidth}{0.8\paperheight}}
	{BinTree/trees_exp/14.png}
\end{center}
Указати послідовність міток вершин дерева, зображеного на рис., що утвориться в результаті обходу дерева в прямому порядку. Мітки записувати через пробіл.\par Алгоритм починає роботу з кореня (на рис. це верхня вершина).
%enditem


\newpage
13. 
\begin{center}
	\adjustimage{max size={0.9\linewidth}{0.9\paperheight}}
	{Graphs/AdjList/151.png}
\end{center}
Для графа, зображеного на рис., вказати список суміжності вершини 1.\par
Формат оформлення відповіді:\par
список суміжності виписувати через пробіл на одному рядку, якщо пробілів десь буде більше одного - не важливо, упорядкування вершин у списку також значення не має.
%enditem


\newpage
14. Для графа на рисунку вказати послідовність, в якій алгоритм пошуку в глибину буде відмічати відвідані вершини, якщо списки суміжності вершин переглядаються алгоритмом за зростанням міток вершин та вершини відмічаються на першому вход\-женні у вершину.\par
Пошук розпочинається з вершини 4.\par

\begin{center}
	\adjustimage{max size={0.9\linewidth}{0.9\paperheight}}
	{Graphs/Traversals/278.png}
\end{center}
%enditem


\newpage
15. Для графа на рисунку з вершини 0 запущено алгоритм пошуку в ширину, що будує кістякове дерево. Списки суміжності вершин переглядаються алгоритмом за зростанням міток вершин.\par
\begin{center}
	\adjustimage{max size={0.9\linewidth}{0.9\paperheight}}
	{Graphs/Traversals/13.png}
\end{center}
\par Які з наведених ребер входять до складу отриманого кістякового дерева?\par
1) (3, 4)\par
2) (2, 5)\par
3) (2, 6)\par
4) (0, 4)\par
5) (1, 6)\par
6) (1, 3)\par
{\vskip 15pt} Номери відповідних ребер записати за зростанням через пробіл.


%enditem


\newpage
16. 

Рядки журналу через двокрапку містять ім'я користувача (ідентифікатор), час події,тип події, код події, наприклад
\begin{verbatim}
s="username : 2022-05-05 13:26 : ERROR : 404"
\end{verbatim}
у коді 
\begin{verbatim}
res = re.fullmatch('???', s) 
s = ''
code_str = ??? 
\end{verbatim}
чим треба замінити кожен з ???, щоб значенням code\_str був код події? 


\vspace{10mm}
\textbf{Завдання 17} 

Задано програму, що містить код 1-2 класів, можливе успадкування, та клієнтський код.
У наведеному коді невеликі частини закриті. 
Також наведено вихід програми.
Необхідно відновити закриті частини, щоб отриманий код відповідав наведеному виходу програми.


\textbf{Завдання 18-19} 

Кілька запитань за завданнями для самостійної роботи.
\vspace{10mm}
Скоріш за все, що завдань буде трохи менше, а деякі попередні будуть переформульовані в стилі двох останніх.

%enditem

\end {large}}

%end

\newpage
%begin

{{
\begin {large}

\newpage
1. Що буде надруковано за виконання?

\begin{verbatim}
print('R', end=' ')
g=('a','b','c','d','e',0,1,2)
print(g[1])
\end{verbatim}

%enditem


\newpage
2. Що буде надруковано за виконання?

\begin{verbatim}
print('R', end=' ')
a=('0','1','2','3','4','5','6','7','8','9')
print(a[5:-3:3])
\end{verbatim}

%enditem



\newpage
3. Що буде надруковано за виконання?

\begin{verbatim}
print('R', end=' ')
b=['0','1','2','3','4']
b.append('12')
print(b)
\end{verbatim}

%enditem



\newpage
4. Що буде надруковано за виконання?

\begin{verbatim}
print('R', end=' ')
f=['0','1','2','3','4']
f[2:-8]='10'
print(f)
\end{verbatim}

%enditem



\newpage
5. Укажіть ТИП значення виразу.\par
\begin{verbatim}
{1:3, 5:8}
\end{verbatim}
%enditem


\newpage
6. Що буде надруковано за виконання?\par
\begin{verbatim}
try:
    s = {37:5, 59:0, 55:1, 86:2, 86:5}
    s[37] = 6
    for x, y in s.items():
        print(x, y, sep=' ', end=' ')
except:
    print('ERR')
\end{verbatim}
%enditem


\newpage
7. Нехай змінна \kw{s} позначає послідовність цілих, по якій можна ітерувати. 

Написати вираз, що перевіряє, чи правда, що послідовність  містить хоча б одне число, що ділиться на 7.
%enditem


\newpage
8. Що буде надруковано за виконання?
%Оригинал был утерян. Требует отладки!
\begin{verbatim}
_d = 0

class D:
    _d = 1
    
    def __init__(self):
        self._d = 2
        _d = 3
        D._d = 4

obj = D()
D._d = 5
try:
    print(_d, end=' ')
    print(obj._d, end=' ')
    print(D._d, end=' ')
    print(obj._d, end=' ')
    print(obj._D__d)
except:
    print('ERR')
\end{verbatim}
%enditem


\newpage
9. Що буде надруковано за виконання?
\begin{verbatim}
class D0:
    def c(self): print(1, end=' ')
    def h(self): print(2, end=' ')

class D1(D0):
    def h(self): print(3, end=' ')

try:
    D0().c()
    D1().h()
except:
    print('ERR')
\end{verbatim}
%enditem


\newpage
10. Що буде надруковано за виконання?
\begin{verbatim}
class A:
    a = 1

    def _f1(self): print(2, end=' ')
    def f2(self): print(3, end=' ')

    def main(self):
        try: print(self.a, end=' ')
        except: print('E1', end=' ')

        try: self._f1()
        except: print('E2', end=' ')

        try: self.f2()
        except: print('E3', end=' ')


class B(A):
    a = 4

    def _f1(self): print(5, end=' ')
    def f2(self): print(6, end=' ')

try: B().main()
except: print('E4', end=' ')
\end{verbatim}
%enditem


\newpage
11. Змінні \kw{next} та \kw{prev} вузла списку позначають наступний та попередній за ним вузол відповідно. Починаючи з голови, у список записано послідовні цілі числа від~$-24$ до~$32$. Нехай змінна \kw{p} позначає вузол, що зберігає значення~$-7$.
Яке значення зберігається у вузлі \kw{p.prev.prev.next.prev.prev} ?
%enditem


\newpage
12. 
\begin{center}
	\adjustimage{max size={0.8\linewidth}{0.8\paperheight}}
	{BinTree/trees_exp/170.png}
\end{center}
Указати послідовність міток вершин дерева, зображеного на рис., що утвориться в результаті обходу дерева в оберненому порядку. Мітки записувати через пробіл.\par Алгоритм починає роботу з кореня (на рис. це верхня вершина).
%enditem


\newpage
13. 
\begin{center}
	\adjustimage{max size={0.9\linewidth}{0.9\paperheight}}
	{Graphs/AdjList/334.png}
\end{center}
Для графа, зображеного на рис., вказати список суміжності вершини 1.\par
Формат оформлення відповіді:\par
список суміжності виписувати через пробіл на одному рядку, якщо пробілів десь буде більше одного - не важливо, упорядкування вершин у списку також значення не має.
%enditem


\newpage
14. Для графа на рисунку з вершини 3 запущено алгоритм пошуку в глибину, що будує кістякове дерево. Списки суміжності вершин переглядаються алгоритмом за зростанням міток вершин.\par
\begin{center}
	\adjustimage{max size={0.9\linewidth}{0.9\paperheight}}
	{Graphs/Traversals/274.png}
\end{center}
\par Які з наведених ребер входять до складу отриманого кістякового дерева?\par
1) (4, 6)\par
2) (3, 6)\par
3) (2, 4)\par
4) (4, 5)\par
5) (3, 4)\par
6) (0, 4)\par
{\vskip 15pt} Номери відповідних ребер записати за зростанням через пробіл.

%enditem


\newpage
15. Для графа на рисунку з вершини 6 запущено алгоритм пошуку в ширину, що будує кістякове дерево. Списки суміжності вершин переглядаються алгоритмом за зростанням міток вершин.\par
\begin{center}
	\adjustimage{max size={0.9\linewidth}{0.9\paperheight}}
	{Graphs/Traversals/57.png}
\end{center}
\par Які з наведених ребер входять до складу отриманого кістякового дерева?\par
1) (2, 4)\par
2) (0, 2)\par
3) (0, 5)\par
4) (4, 5)\par
5) (1, 3)\par
6) (1, 2)\par
{\vskip 15pt} Номери відповідних ребер записати за зростанням через пробіл.


%enditem


\newpage
16. 

Рядки журналу через двокрапку містять ім'я користувача (ідентифікатор), час події,тип події, код події, наприклад
\begin{verbatim}
s="username : 2022-05-05 13:26 : ERROR : 404"
\end{verbatim}
у коді 
\begin{verbatim}
res = re.fullmatch('???', s) 
s = ''
code_str = ??? 
\end{verbatim}
чим треба замінити кожен з ???, щоб значенням code\_str був код події? 


\vspace{10mm}
\textbf{Завдання 17} 

Задано програму, що містить код 1-2 класів, можливе успадкування, та клієнтський код.
У наведеному коді невеликі частини закриті. 
Також наведено вихід програми.
Необхідно відновити закриті частини, щоб отриманий код відповідав наведеному виходу програми.


\textbf{Завдання 18-19} 

Кілька запитань за завданнями для самостійної роботи.
\vspace{10mm}
Скоріш за все, що завдань буде трохи менше, а деякі попередні будуть переформульовані в стилі двох останніх.

%enditem

\end {large}}

%end

\newpage
%begin

{{
\begin {large}

\newpage
1. Що буде надруковано за виконання?

\begin{verbatim}
print('R', end=' ')
h=('a','b','c','d','e',0,1)
print(h[12])
\end{verbatim}

%enditem


\newpage
2. Що буде надруковано за виконання?

\begin{verbatim}
print('R', end=' ')
b='0123456789'
print(b[-13::-2])
\end{verbatim}

%enditem



\newpage
3. Що буде надруковано за виконання?

\begin{verbatim}
print('R', end=' ')
g=['0','1','2','3','4']
g+='10'
print(g)
\end{verbatim}

%enditem



\newpage
4. Що буде надруковано за виконання?

\begin{verbatim}
print('R', end=' ')
a=[10,11,12,13,14]
del a[-9:8]
print(a)
\end{verbatim}

%enditem



\newpage
5. Укажіть ТИП значення виразу.\par
\begin{verbatim}
[7,2,3,4]
\end{verbatim}
%enditem


\newpage
6. Що буде надруковано за виконання?\par
\begin{verbatim}
try:
    d = {27:2, 23:2, 66:7, 46:8, 35:3}
    d[27] = 9
    for x in d.values():
        print(x, sep=' ', end=' ')
except:
    print('ERR')
\end{verbatim}
%enditem


\newpage
7. Нехай змінна \kw{s} позначає послідовність цілих, по якій можна ітерувати.

Написати вираз, що перевіряє, чи правда, що всі елементи послідовності менші за задане число num.
%enditem


\newpage
8. Що буде надруковано за виконання?
%Оригинал был утерян. Требует отладки!
\begin{verbatim}
__c = 0

class B:
    __c = 1
    
    def __init__(self):
        self.__c = 2
        __c = 3
        B.__c = 4

obj = B()
obj.__c = 5
try:
    print(__c, end=' ')
    print(obj.__c, end=' ')
    print(B.__c, end=' ')
    print(obj.__c, end=' ')
    print(obj._B__c)
except:
    print('ERR')
\end{verbatim}
%enditem


\newpage
9. Що буде надруковано за виконання?
\begin{verbatim}
class A0:
    def b(self): print(1, end=' ')
    def g(self): print(2, end=' ')

class A1(A0):
    def g(self): print(3, end=' ')

try:
    A1().b()
    A1().g()
except:
    print('ERR')
\end{verbatim}
%enditem


\newpage
10. Що буде надруковано за виконання?
\begin{verbatim}
class A:
    c = 1

    def __init__(self):
        a = 2
        a = 3

class B(A):
    c = 4
    b = 5

    def __init__(self):
        super().__init__()
        self.b = 6

try:
    obj = B()
    obj.a = obj.b + 10
    B.a = 13
except: print('EE', end=' ')

try: print(obj.a, end= ' ')
except: print('E1', end=' ')

try: print(A.a, end= ' ')
except: print('E2', end=' ')

try: print(B.a, end= ' ')
except: print('E3', end=' ')
\end{verbatim}
%enditem


\newpage
11. Змінні \kw{next} та \kw{prev} вузла списку позначають наступний та попередній за ним вузол відповідно. Починаючи з голови, у список записано послідовні цілі числа від~$-42$ до~$21$. Нехай змінна \kw{p} позначає вузол, що зберігає значення~$0$.
Яке значення зберігається у вузлі \kw{p.next.prev.prev.next.prev} ?
%enditem


\newpage
12. 
\begin{center}
	\adjustimage{max size={0.8\linewidth}{0.8\paperheight}}
	{BinTree/trees_exp/68.png}
\end{center}
Указати послідовність міток вершин дерева, зображеного на рис., що утвориться в результаті обходу дерева в прямому порядку. Мітки записувати через пробіл.\par Алгоритм починає роботу з кореня (на рис. це верхня вершина).
%enditem


\newpage
13. 
\begin{center}
	\adjustimage{max size={0.9\linewidth}{0.9\paperheight}}
	{Graphs/AdjList/284.png}
\end{center}
Для графа, зображеного на рис., вказати список суміжності вершини 8.\par
Формат оформлення відповіді:\par
список суміжності виписувати через пробіл на одному рядку, якщо пробілів десь буде більше одного - не важливо, упорядкування вершин у списку також значення не має.
%enditem


\newpage
14. Для графа на рисунку вказати послідовність, в якій алгоритм пошуку в глибину буде відмічати відвідані вершини, якщо списки суміжності вершин переглядаються алгоритмом за зростанням міток вершин та вершини відмічаються на першому вход\-женні у вершину.\par
Пошук розпочинається з вершини 0.\par

\begin{center}
	\adjustimage{max size={0.9\linewidth}{0.9\paperheight}}
	{Graphs/Traversals/118.png}
\end{center}
%enditem


\newpage
15. Для графа на рисунку з вершини 5 запущено алгоритм пошуку в ширину, що будує кістякове дерево. Списки суміжності вершин переглядаються алгоритмом за зростанням міток вершин.\par
\begin{center}
	\adjustimage{max size={0.9\linewidth}{0.9\paperheight}}
	{Graphs/Traversals/339.png}
\end{center}
\par Які з наведених ребер входять до складу отриманого кістякового дерева?\par
1) (5, 6)\par
2) (1, 3)\par
3) (0, 2)\par
4) (4, 5)\par
5) (2, 5)\par
6) (0, 1)\par
{\vskip 15pt} Номери відповідних ребер записати за зростанням через пробіл.


%enditem


\newpage
16. 

Рядки журналу через двокрапку містять ім'я користувача (ідентифікатор), час події,тип події, код події, наприклад
\begin{verbatim}
s="username : 2022-05-05 13:26 : ERROR : 404"
\end{verbatim}
у коді 
\begin{verbatim}
res = re.fullmatch('???', s) 
s = ''
code_str = ??? 
\end{verbatim}
чим треба замінити кожен з ???, щоб значенням code\_str був код події? 


\vspace{10mm}
\textbf{Завдання 17} 

Задано програму, що містить код 1-2 класів, можливе успадкування, та клієнтський код.
У наведеному коді невеликі частини закриті. 
Також наведено вихід програми.
Необхідно відновити закриті частини, щоб отриманий код відповідав наведеному виходу програми.


\textbf{Завдання 18-19} 

Кілька запитань за завданнями для самостійної роботи.
\vspace{10mm}
Скоріш за все, що завдань буде трохи менше, а деякі попередні будуть переформульовані в стилі двох останніх.

%enditem

\end {large}}

%end

\newpage
%begin

{{
\begin {large}

\newpage
1. Що буде надруковано за виконання?

\begin{verbatim}
print('R', end=' ')
d=('a','b','c','d','e',0,1,2,3)
print(d[-11])
\end{verbatim}

%enditem


\newpage
2. Що буде надруковано за виконання?

\begin{verbatim}
print('R', end=' ')
d=[0,1,2,3,4,5,6,7,8,9]
print(d[-15::2])
\end{verbatim}

%enditem



\newpage
3. Що буде надруковано за виконання?

\begin{verbatim}
print('R', end=' ')
a=[10,11,12,13,14]
a.extend([4,5])
print(a)
\end{verbatim}

%enditem



\newpage
4. Що буде надруковано за виконання?

\begin{verbatim}
print('R', end=' ')
e=['0','1','2','3','4']
e[1:]='12'
print(e)
\end{verbatim}

%enditem



\newpage
5. Укажіть ТИП значення виразу.\par
\begin{verbatim}
{y:y*y for y in range(-7,7)}
\end{verbatim}
%enditem


\newpage
6. Що буде надруковано за виконання?\par
\begin{verbatim}
try:
    t = {68:0, 39:4, 89:7, 68:5}
    t[68] = 2
    for x in t.items():
        print(x, sep=' ', end=' ')
except:
    print('ERR')
\end{verbatim}
%enditem


\newpage
7. Нехай змінна \kw{s} позначає послідовність цілих, по якій можна ітерувати.

Написати вираз, що перевіряє, чи правда, що всі елементи послідовності діляться на 11.
%enditem


\newpage
8. Що буде надруковано за виконання?
\begin{verbatim}
d = 0
h = 1

class D:
    d = 2
    
    def __init__(self):
        global d
        self.d = 3
        d = 4
        D.d = 5

obj = D()
try:
    print(d, end=' ')
    print(h, end=' ')
    print(D.d, end=' ')
    print(obj.d)
except:
    print('ERR')
\end{verbatim}
%enditem


\newpage
9. Що буде надруковано за виконання?
\begin{verbatim}
class B0:
    def d(self): print(1, end=' ')
    def f(self): print(2, end=' ')

class B1(B0):
    def d(self): print(3, end=' ')

try:
    B1().f()
    B1().d()
except:
    print('ERR')
\end{verbatim}
%enditem


\newpage
10. Що буде надруковано за виконання?
\begin{verbatim}
class A:
    c = 1

    def __init__(self):
        b = 2
        c = 3

class B(A):
    c = 4
    a = 5

    def __init__(self):
        super().__init__()
        self.a = 6

try:
    obj = B()
    obj.c = obj.a + 10
    B.c = 13
except: print('EE', end=' ')

try: print(obj.c, end= ' ')
except: print('E1', end=' ')

try: print(A.c, end= ' ')
except: print('E2', end=' ')

try: print(B.c, end= ' ')
except: print('E3', end=' ')
\end{verbatim}
%enditem


\newpage
11. Змінні \kw{next} та \kw{prev} вузла списку позначають наступний та попередній за ним вузол відповідно. Починаючи з голови, у список записано послідовні цілі числа від~$-70$ до~$20$. Нехай змінна \kw{p} позначає вузол, що зберігає значення~$4$.
Яке значення зберігається у вузлі \kw{p.prev.next.next.prev.next} ?
%enditem


\newpage
12. 
\begin{center}
	\adjustimage{max size={0.8\linewidth}{0.8\paperheight}}
	{BinTree/trees_exp/106.png}
\end{center}
Указати послідовність міток вершин дерева, зображеного на рис., що утвориться в результаті обходу дерева в оберненому порядку. Мітки записувати через пробіл.\par Алгоритм починає роботу з кореня (на рис. це верхня вершина).
%enditem


\newpage
13. 
\begin{center}
	\adjustimage{max size={0.9\linewidth}{0.9\paperheight}}
	{Graphs/AdjList/345.png}
\end{center}
Для графа, зображеного на рис., вказати список суміжності вершини 1.\par
Формат оформлення відповіді:\par
список суміжності виписувати через пробіл на одному рядку, якщо пробілів десь буде більше одного - не важливо, упорядкування вершин у списку також значення не має.
%enditem


\newpage
14. Для графа на рисунку вказати послідовність, в якій алгоритм пошуку в глибину буде відмічати відвідані вершини, якщо списки суміжності вершин переглядаються алгоритмом за зростанням міток вершин та вершини відмічаються на першому вход\-женні у вершину.\par
Пошук розпочинається з вершини 2.\par

\begin{center}
	\adjustimage{max size={0.9\linewidth}{0.9\paperheight}}
	{Graphs/Traversals/163.png}
\end{center}
%enditem


\newpage
15. Для графа на рисунку з вершини 0 запущено алгоритм пошуку в ширину, що будує кістякове дерево. Списки суміжності вершин переглядаються алгоритмом за зростанням міток вершин.\par
\begin{center}
	\adjustimage{max size={0.9\linewidth}{0.9\paperheight}}
	{Graphs/Traversals/0.png}
\end{center}
\par Які з наведених ребер входять до складу отриманого кістякового дерева?\par
1) (0, 5)\par
2) (2, 5)\par
3) (2, 6)\par
4) (4, 5)\par
5) (0, 3)\par
6) (0, 6)\par
{\vskip 15pt} Номери відповідних ребер записати за зростанням через пробіл.


%enditem


\newpage
16. 

Рядки журналу через двокрапку містять ім'я користувача (ідентифікатор), час події,тип події, код події, наприклад
\begin{verbatim}
s="username : 2022-05-05 13:26 : ERROR : 404"
\end{verbatim}
у коді 
\begin{verbatim}
res = re.fullmatch('???', s) 
s = ''
code_str = ??? 
\end{verbatim}
чим треба замінити кожен з ???, щоб значенням code\_str був код події? 


\vspace{10mm}
\textbf{Завдання 17} 

Задано програму, що містить код 1-2 класів, можливе успадкування, та клієнтський код.
У наведеному коді невеликі частини закриті. 
Також наведено вихід програми.
Необхідно відновити закриті частини, щоб отриманий код відповідав наведеному виходу програми.


\textbf{Завдання 18-19} 

Кілька запитань за завданнями для самостійної роботи.
\vspace{10mm}
Скоріш за все, що завдань буде трохи менше, а деякі попередні будуть переформульовані в стилі двох останніх.

%enditem

\end {large}}

%end

\newpage
%begin

{{
\begin {large}

\newpage
1. Що буде надруковано за виконання?

\begin{verbatim}
print('R', end=' ')
c=('a','b','c','d','e',0,1,2,3,4)
print(c[1])
\end{verbatim}

%enditem


\newpage
2. Що буде надруковано за виконання?

\begin{verbatim}
print('R', end=' ')
h=['0','1','2','3','4','5','6','7','8','9']
print(h[0:2:-3])
\end{verbatim}

%enditem



\newpage
3. Що буде надруковано за виконання?

\begin{verbatim}
print('R', end=' ')
c=[10,11,12,13,14]
print(c.index(15), end=' ')
\end{verbatim}

%enditem



\newpage
4. Що буде надруковано за виконання?

\begin{verbatim}
print('R', end=' ')
g=['0','1','2','3','4']
g[-6:-7]=[4,5]
print(g)
\end{verbatim}

%enditem



\newpage
5. Укажіть ТИП значення виразу.\par
\begin{verbatim}
{(1,4), 8}
\end{verbatim}
%enditem


\newpage
6. Що буде надруковано за виконання?\par
\begin{verbatim}
try:
    d = {62:7, 35:2, 54:6, 84:0, 62:1}
    d[67] = 5
    for x in d.keys():
        print(x, sep=' ', end=' ')
except:
    print('ERR')
\end{verbatim}
%enditem


\newpage
7. Нехай змінна \kw{s} позначає дуже довгу послідовність, по якій можна ітерувати. 

Написати вираз-генератор, що задає підпослідовність її елементів, що є числами типу \kw{int}.

Обмеження: усі елементи шуканої підпослідовності одночасно в пам'яті розміщені бути не можуть. 
%enditem


\newpage
8. Що буде надруковано за виконання?
%Оригинал был утерян. Требует отладки!
\begin{verbatim}
_a = 0

class A:
    _a = 1
    
    def __init__(self):
        self._a = 2
        _a = 3
        A._a = 4

obj = A()
A._a = 5
try:
    print(_a, end=' ')
    print(obj._a, end=' ')
    print(A._a, end=' ')
    print(obj._a, end=' ')
    print(obj._A__a)
except:
    print('ERR')
\end{verbatim}
%enditem


\newpage
9. Що буде надруковано за виконання?
\begin{verbatim}
class A0:
    def a(self): print(1, end=' ')
    def f(self): print(2, end=' ')

class A1(A0):
    def a(self): print(3, end=' ')

try:
    A1().f()
    A1().a()
except:
    print('ERR')
\end{verbatim}
%enditem


\newpage
10. Що буде надруковано за виконання?
\begin{verbatim}
class A:
    c = 1

    def __f1(self): print(2, end=' ')
    def _f2(self): print(3, end=' ')

    def main(self):
        try: print(self.c, end=' ')
        except: print('E1', end=' ')

        try: self.__f1()
        except: print('E2', end=' ')

        try: self._f2()
        except: print('E3', end=' ')


class B(A):
    c = 4

    def __f1(self): print(5, end=' ')
    def _f2(self): print(6, end=' ')

try: B().main()
except: print('E4', end=' ')
\end{verbatim}
%enditem


\newpage
11. Змінні \kw{next} та \kw{prev} вузла списку позначають наступний та попередній за ним вузол відповідно. Починаючи з голови, у список записано послідовні цілі числа від~$-17$ до~$34$. Нехай змінна \kw{p} позначає вузол, що зберігає значення~$6$.
Яке значення зберігається у вузлі \kw{p.prev.prev.next.next.next} ?
%enditem


\newpage
12. 
\begin{center}
	\adjustimage{max size={0.8\linewidth}{0.8\paperheight}}
	{BinTree/trees_exp/127.png}
\end{center}
Указати послідовність міток вершин дерева, зображеного на рис., що утвориться в результаті обходу дерева в оберненому порядку. Мітки записувати через пробіл.\par Алгоритм починає роботу з кореня (на рис. це верхня вершина).
%enditem


\newpage
13. 
\begin{center}
	\adjustimage{max size={0.9\linewidth}{0.9\paperheight}}
	{Graphs/AdjList/322.png}
\end{center}
Для графа, зображеного на рис., вказати список суміжності вершини 9.\par
Формат оформлення відповіді:\par
список суміжності виписувати через пробіл на одному рядку, якщо пробілів десь буде більше одного - не важливо, упорядкування вершин у списку також значення не має.
%enditem


\newpage
14. Для графа на рисунку вказати послідовність, в якій алгоритм пошуку в глибину буде відмічати відвідані вершини, якщо списки суміжності вершин переглядаються алгоритмом за зростанням міток вершин та вершини відмічаються на першому вход\-женні у вершину.\par
Пошук розпочинається з вершини 6.\par

\begin{center}
	\adjustimage{max size={0.9\linewidth}{0.9\paperheight}}
	{Graphs/Traversals/313.png}
\end{center}
%enditem


\newpage
15. Для графа на рисунку з вершини 5 запущено алгоритм пошуку в ширину, що будує кістякове дерево. Списки суміжності вершин переглядаються алгоритмом за зростанням міток вершин.\par
\begin{center}
	\adjustimage{max size={0.9\linewidth}{0.9\paperheight}}
	{Graphs/Traversals/315.png}
\end{center}
\par Які з наведених ребер входять до складу отриманого кістякового дерева?\par
1) (0, 4)\par
2) (2, 4)\par
3) (3, 6)\par
4) (1, 5)\par
5) (1, 2)\par
6) (1, 4)\par
{\vskip 15pt} Номери відповідних ребер записати за зростанням через пробіл.


%enditem


\newpage
16. 

Рядки журналу через двокрапку містять ім'я користувача (ідентифікатор), час події,тип події, код події, наприклад
\begin{verbatim}
s="username : 2022-05-05 13:26 : ERROR : 404"
\end{verbatim}
у коді 
\begin{verbatim}
res = re.fullmatch('???', s) 
s = ''
code_str = ??? 
\end{verbatim}
чим треба замінити кожен з ???, щоб значенням code\_str був код події? 


\vspace{10mm}
\textbf{Завдання 17} 

Задано програму, що містить код 1-2 класів, можливе успадкування, та клієнтський код.
У наведеному коді невеликі частини закриті. 
Також наведено вихід програми.
Необхідно відновити закриті частини, щоб отриманий код відповідав наведеному виходу програми.


\textbf{Завдання 18-19} 

Кілька запитань за завданнями для самостійної роботи.
\vspace{10mm}
Скоріш за все, що завдань буде трохи менше, а деякі попередні будуть переформульовані в стилі двох останніх.

%enditem

\end {large}}

%end

\newpage
%begin

{{
\begin {large}

\newpage
1. Що буде надруковано за виконання?

\begin{verbatim}
print('R', end=' ')
a=('a','b','c','d','e',0,1)
print(a[-10])
\end{verbatim}

%enditem


\newpage
2. Що буде надруковано за виконання?

\begin{verbatim}
print('R', end=' ')
c=('0','1','2','3','4','5','6','7','8','9')
print(c[:-14:-2])
\end{verbatim}

%enditem



\newpage
3. Що буде надруковано за виконання?

\begin{verbatim}
print('R', end=' ')
f=['0','1','2','3','4']
print(f.pop(5), end=' ')
print(f)
\end{verbatim}

%enditem



\newpage
4. Що буде надруковано за виконання?

\begin{verbatim}
print('R', end=' ')
b=[10,11,12,13,14]
b[:3]='10'
print(b)
\end{verbatim}

%enditem



\newpage
5. Укажіть ТИП значення виразу.\par
\begin{verbatim}
(11,2,3)
\end{verbatim}
%enditem


\newpage
6. Що буде надруковано за виконання?\par
\begin{verbatim}
try:
    t = {45:7, 13:9, 27:9, 41:3, 13:0}
    t[22] = 3
    for x in t.items():
        print(*x, sep=' ', end=' ')
except:
    print('ERR')
\end{verbatim}
%enditem


\newpage
7. 
 Задано числову послідовність \kw{s} довжини 6000. Записати ВИРАЗ, значенням якого є список, що складається з елементів послідовності, що діляться на 13, записаних в обернено\-му порядку.
%enditem


\newpage
8. Що буде надруковано за виконання?
%Оригинал был утерян. Требует отладки!
\begin{verbatim}
__b = 0

class D:
    __b = 1
    
    def __init__(self):
        self.__b = 2
        __b = 3
        D.__b = 4

obj = D()
D.__b = 5
try:
    print(__b, end=' ')
    print(obj.__b, end=' ')
    print(D.__b, end=' ')
    print(obj.__b, end=' ')
    print(obj._D__b)
except:
    print('ERR')
\end{verbatim}
%enditem


\newpage
9. Що буде надруковано за виконання?
\begin{verbatim}
class B0:
    def c(self): print(1, end=' ')
    def h(self): print(2, end=' ')

class B1(B0):
    def h(self): print(3, end=' ')

try:
    B1().c()
    B0().h()
except:
    print('ERR')
\end{verbatim}
%enditem


\newpage
10. Що буде надруковано за виконання?
\begin{verbatim}
class A:
    b = 1

    def _h1(self): print(2, end=' ')
    def __h2(self): print(3, end=' ')

    def main(self):
        try: print(self.b, end=' ')
        except: print('E1', end=' ')

        try: self._h1()
        except: print('E2', end=' ')

        try: self.__h2()
        except: print('E3', end=' ')


class B(A):
    b = 4

    def _h1(self): print(5, end=' ')
    def __h2(self): print(6, end=' ')

try: B().main()
except: print('E4', end=' ')
\end{verbatim}
%enditem


\newpage
11. Змінні \kw{next} та \kw{prev} вузла списку позначають наступний та попередній за ним вузол відповідно. Починаючи з голови, у список записано послідовні цілі числа від~$-51$ до~$16$. Нехай змінна \kw{p} позначає вузол, що зберігає значення~$2$.
Яке значення зберігається у вузлі \kw{p.next.next.prev.prev.prev} ?
%enditem


\newpage
12. 
\begin{center}
	\adjustimage{max size={0.8\linewidth}{0.8\paperheight}}
	{BinTree/trees_exp/28.png}
\end{center}
Указати послідовність міток вершин дерева, зображеного на рис., що утвориться в результаті обходу дерева в прямому порядку. Мітки записувати через пробіл.\par Алгоритм починає роботу з кореня (на рис. це верхня вершина).
%enditem


\newpage
13. 
\begin{center}
	\adjustimage{max size={0.9\linewidth}{0.9\paperheight}}
	{Graphs/AdjList/93.png}
\end{center}
Для графа, зображеного на рис., вказати список суміжності вершини 3.\par
Формат оформлення відповіді:\par
список суміжності виписувати через пробіл на одному рядку, якщо пробілів десь буде більше одного - не важливо, упорядкування вершин у списку також значення не має.
%enditem


\newpage
14. Для графа на рисунку вказати послідовність, в якій алгоритм пошуку в глибину буде відмічати відвідані вершини, якщо списки суміжності вершин переглядаються алгоритмом за зростанням міток вершин та вершини відмічаються на першому вход\-женні у вершину.\par
Пошук розпочинається з вершини 3.\par

\begin{center}
	\adjustimage{max size={0.9\linewidth}{0.9\paperheight}}
	{Graphs/Traversals/202.png}
\end{center}
%enditem


\newpage
15. Для графа на рисунку з вершини 6 запущено алгоритм пошуку в ширину, що будує кістякове дерево. Списки суміжності вершин переглядаються алгоритмом за зростанням міток вершин.\par
\begin{center}
	\adjustimage{max size={0.9\linewidth}{0.9\paperheight}}
	{Graphs/Traversals/119.png}
\end{center}
\par Які з наведених ребер входять до складу отриманого кістякового дерева?\par
1) (3, 4)\par
2) (0, 3)\par
3) (1, 3)\par
4) (0, 1)\par
5) (4, 6)\par
6) (2, 4)\par
{\vskip 15pt} Номери відповідних ребер записати за зростанням через пробіл.


%enditem


\newpage
16. 

Рядки журналу через двокрапку містять ім'я користувача (ідентифікатор), час події,тип події, код події, наприклад
\begin{verbatim}
s="username : 2022-05-05 13:26 : ERROR : 404"
\end{verbatim}
у коді 
\begin{verbatim}
res = re.fullmatch('???', s) 
s = ''
code_str = ??? 
\end{verbatim}
чим треба замінити кожен з ???, щоб значенням code\_str був код події? 


\vspace{10mm}
\textbf{Завдання 17} 

Задано програму, що містить код 1-2 класів, можливе успадкування, та клієнтський код.
У наведеному коді невеликі частини закриті. 
Також наведено вихід програми.
Необхідно відновити закриті частини, щоб отриманий код відповідав наведеному виходу програми.


\textbf{Завдання 18-19} 

Кілька запитань за завданнями для самостійної роботи.
\vspace{10mm}
Скоріш за все, що завдань буде трохи менше, а деякі попередні будуть переформульовані в стилі двох останніх.

%enditem

\end {large}}

%end

\newpage
%begin

{{
\begin {large}

\newpage
1. Що буде надруковано за виконання?

\begin{verbatim}
print('R', end=' ')
f=('a','b','c','d','e',0,1,2,4)
print(f[-6])
\end{verbatim}

%enditem


\newpage
2. Що буде надруковано за виконання?

\begin{verbatim}
print('R', end=' ')
e=['0','1','2','3','4','5','6','7','8','9']
print(e[-2:-9:-3])
\end{verbatim}

%enditem



\newpage
3. Що буде надруковано за виконання?

\begin{verbatim}
print('R', end=' ')
e=['0','1','2','3','4']
e+='11'
print(e)
\end{verbatim}

%enditem



\newpage
4. Що буде надруковано за виконання?

\begin{verbatim}
print('R', end=' ')
c=['0','1','2','3','4']
del c[:4]
print(c)
\end{verbatim}

%enditem



\newpage
5. Укажіть ТИП значення виразу.\par
\begin{verbatim}
{y:None for y in range(-8,3)}
\end{verbatim}
%enditem


\newpage
6. Що буде надруковано за виконання?\par
\begin{verbatim}
try:
    s = {42:6, 87:8, 34:2, 42:5}
    s[34] = 2
    for x in s.values():
        print(x, sep=' ', end=' ')
except:
    print('ERR')
\end{verbatim}
%enditem


\newpage
7. 
Записати ВИРАЗ, значенням якого є список, що складається з квадратних коренів цілих чисел від 13 до 2000 (включно), що не діляться на 5.
%enditem


\newpage
8. Що буде надруковано за виконання?
%Оригинал был утерян. Требует отладки!
\begin{verbatim}
d = 0

class A:
    d = 1
    
    def __init__(self):
        self.d = 2
        d = 3
        A.d = 4

obj = A()
obj.d = 5
try:
    print(d, end=' ')
    print(obj.d, end=' ')
    print(A.d, end=' ')
    print(obj.d, end=' ')
    print(obj._A__d)
except:
    print('ERR')
\end{verbatim}
%enditem


\newpage
9. Що буде надруковано за виконання?
\begin{verbatim}
class D0:
    def b(self): print(1, end=' ')
    def g(self): print(2, end=' ')

class D1(D0):
    def g(self): print(3, end=' ')

try:
    D1().b()
    D1().g()
except:
    print('ERR')
\end{verbatim}
%enditem


\newpage
10. Що буде надруковано за виконання?
\begin{verbatim}
class A:
    b = 1

    def __init__(self):
        c = 2
        c = 3

class B(A):
    c = 4
    a = 5

    def __init__(self):
        super().__init__()
        self.a = 6

try:
    obj = B()
    obj.c = obj.b + 10
    B.c = 13
except: print('EE', end=' ')

try: print(obj.c, end= ' ')
except: print('E1', end=' ')

try: print(A.c, end= ' ')
except: print('E2', end=' ')

try: print(B.c, end= ' ')
except: print('E3', end=' ')
\end{verbatim}
%enditem


\newpage
11. Змінні \kw{next} та \kw{prev} вузла списку позначають наступний та попередній за ним вузол відповідно. Починаючи з голови, у список записано послідовні цілі числа від~$-26$ до~$49$. Нехай змінна \kw{p} позначає вузол, що зберігає значення~$-1$.
Яке значення зберігається у вузлі \kw{p.next.next.next.next.prev} ?
%enditem


\newpage
12. 
\begin{center}
	\adjustimage{max size={0.8\linewidth}{0.8\paperheight}}
	{BinTree/trees_exp/141.png}
\end{center}
Указати послідовність міток вершин дерева, зображеного на рис., що утвориться в результаті обходу дерева в прямому порядку. Мітки записувати через пробіл.\par Алгоритм починає роботу з кореня (на рис. це верхня вершина).
%enditem


\newpage
13. 
\begin{center}
	\adjustimage{max size={0.9\linewidth}{0.9\paperheight}}
	{Graphs/AdjList/286.png}
\end{center}
Для графа, зображеного на рис., вказати список суміжності вершини 2.\par
Формат оформлення відповіді:\par
список суміжності виписувати через пробіл на одному рядку, якщо пробілів десь буде більше одного - не важливо, упорядкування вершин у списку також значення не має.
%enditem


\newpage
14. Для графа на рисунку з вершини 0 запущено алгоритм пошуку в глибину, що будує кістякове дерево. Списки суміжності вершин переглядаються алгоритмом за зростанням міток вершин.\par
\begin{center}
	\adjustimage{max size={0.9\linewidth}{0.9\paperheight}}
	{Graphs/Traversals/5.png}
\end{center}
\par Які з наведених ребер входять до складу отриманого кістякового дерева?\par
1) (2, 6)\par
2) (1, 2)\par
3) (1, 5)\par
4) (4, 6)\par
5) (3, 5)\par
6) (4, 5)\par
{\vskip 15pt} Номери відповідних ребер записати за зростанням через пробіл.

%enditem


\newpage
15. Для графа на рисунку з вершини 1 запущено алгоритм пошуку в ширину, що будує кістякове дерево. Списки суміжності вершин переглядаються алгоритмом за зростанням міток вершин.\par
\begin{center}
	\adjustimage{max size={0.9\linewidth}{0.9\paperheight}}
	{Graphs/Traversals/221.png}
\end{center}
\par Які з наведених ребер входять до складу отриманого кістякового дерева?\par
1) (3, 5)\par
2) (4, 5)\par
3) (5, 6)\par
4) (0, 3)\par
5) (0, 1)\par
6) (1, 5)\par
{\vskip 15pt} Номери відповідних ребер записати за зростанням через пробіл.


%enditem


\newpage
16. 

Рядки журналу через двокрапку містять ім'я користувача (ідентифікатор), час події,тип події, код події, наприклад
\begin{verbatim}
s="username : 2022-05-05 13:26 : ERROR : 404"
\end{verbatim}
у коді 
\begin{verbatim}
res = re.fullmatch('???', s) 
s = ''
code_str = ??? 
\end{verbatim}
чим треба замінити кожен з ???, щоб значенням code\_str був код події? 


\vspace{10mm}
\textbf{Завдання 17} 

Задано програму, що містить код 1-2 класів, можливе успадкування, та клієнтський код.
У наведеному коді невеликі частини закриті. 
Також наведено вихід програми.
Необхідно відновити закриті частини, щоб отриманий код відповідав наведеному виходу програми.


\textbf{Завдання 18-19} 

Кілька запитань за завданнями для самостійної роботи.
\vspace{10mm}
Скоріш за все, що завдань буде трохи менше, а деякі попередні будуть переформульовані в стилі двох останніх.

%enditem

\end {large}}

%end

\newpage
%begin

{{
\begin {large}

\newpage
1. Що буде надруковано за виконання?

\begin{verbatim}
print('R', end=' ')
e=('a','b','c','d','e',0,1,2)
print(e[-3])
\end{verbatim}

%enditem


\newpage
2. Що буде надруковано за виконання?

\begin{verbatim}
print('R', end=' ')
f='0123456789'
print(f[-5:-4:-1])
\end{verbatim}

%enditem



\newpage
3. Що буде надруковано за виконання?

\begin{verbatim}
print('R', end=' ')
h=[10,11,12,13,14]
h.append('13')
print(h)
\end{verbatim}

%enditem



\newpage
4. Що буде надруковано за виконання?

\begin{verbatim}
print('R', end=' ')
h=[10,11,12,13,14]
h[-2:-1]=(12)
print(h)
\end{verbatim}

%enditem



\newpage
5. Укажіть ТИП значення виразу.\par
\begin{verbatim}
{x*x:x for x in range(-7,7)}
\end{verbatim}
%enditem


\newpage
6. Що буде надруковано за виконання?\par
\begin{verbatim}
try:
    t = {40:6, 61:7, 47:5, 40:7}
    t[36] = 7
    for x in t :
        print(x, sep=' ', end=' ')
except:
    print('ERR')
\end{verbatim}
%enditem


\newpage
7. Нехай змінна \kw{s} позначає дуже довгу послідовність цілих, по якій можна ітерувати. 

Написати вираз-генератор, що задає підпослідовність її елементів, що є додатни\-ми. 

Обмеження: усі елементи шуканої підпослідовності одночасно в пам'яті розміщені бути не можуть. 
%enditem


\newpage
8. Що буде надруковано за виконання?
%Оригинал был утерян. Требует отладки!
\begin{verbatim}
__a = 0

class C:
    __a = 1
    
    def __init__(self):
        self.__a = 2
        __a = 3
        C.__a = 4

obj = C()
obj.__a = 5
try:
    print(__a, end=' ')
    print(obj.__a, end=' ')
    print(C.__a, end=' ')
    print(obj.__a, end=' ')
    print(obj._C__a)
except:
    print('ERR')
\end{verbatim}
%enditem


\newpage
9. Що буде надруковано за виконання?
\begin{verbatim}
class C0:
    def d(self): print(1, end=' ')
    def g(self): print(2, end=' ')

class C1(C0):
    def g(self): print(3, end=' ')

try:
    C0().d()
    C1().g()
except:
    print('ERR')
\end{verbatim}
%enditem


\newpage
10. Що буде надруковано за виконання?
\begin{verbatim}
class A:
    b = 1

    def __init__(self):
        a = 2
        b = 3

class B(A):
    b = 4
    c = 5

    def __init__(self):
        super().__init__()
        self.c = 6

try:
    obj = B()
    obj.a = obj.a + 10
    B.a = 13
except: print('EE', end=' ')

try: print(obj.a, end= ' ')
except: print('E1', end=' ')

try: print(A.a, end= ' ')
except: print('E2', end=' ')

try: print(B.a, end= ' ')
except: print('E3', end=' ')
\end{verbatim}
%enditem


\newpage
11. Змінні \kw{next} та \kw{prev} вузла списку позначають наступний та попередній за ним вузол відповідно. Починаючи з голови, у список записано послідовні цілі числа від~$-57$ до~$22$. Нехай змінна \kw{p} позначає вузол, що зберігає значення~$9$.
Яке значення зберігається у вузлі \kw{p.prev.prev.prev.prev.next} ?
%enditem


\newpage
12. 
\begin{center}
	\adjustimage{max size={0.8\linewidth}{0.8\paperheight}}
	{BinTree/trees_exp/39.png}
\end{center}
Указати послідовність міток вершин дерева, зображеного на рис., що утвориться в результаті обходу дерева в оберненому порядку. Мітки записувати через пробіл.\par Алгоритм починає роботу з кореня (на рис. це верхня вершина).
%enditem


\newpage
13. 
\begin{center}
	\adjustimage{max size={0.9\linewidth}{0.9\paperheight}}
	{Graphs/AdjList/349.png}
\end{center}
Для графа, зображеного на рис., вказати список суміжності вершини 1.\par
Формат оформлення відповіді:\par
список суміжності виписувати через пробіл на одному рядку, якщо пробілів десь буде більше одного - не важливо, упорядкування вершин у списку також значення не має.
%enditem


\newpage
14. Для графа на рисунку вказати послідовність, в якій алгоритм пошуку в глибину буде відмічати відвідані вершини, якщо списки суміжності вершин переглядаються алгоритмом за зростанням міток вершин та вершини відмічаються на першому вход\-женні у вершину.\par
Пошук розпочинається з вершини 2.\par

\begin{center}
	\adjustimage{max size={0.9\linewidth}{0.9\paperheight}}
	{Graphs/Traversals/244.png}
\end{center}
%enditem


\newpage
15. Для графа на рисунку з вершини 4 запущено алгоритм пошуку в ширину, що будує кістякове дерево. Списки суміжності вершин переглядаються алгоритмом за зростанням міток вершин.\par
\begin{center}
	\adjustimage{max size={0.9\linewidth}{0.9\paperheight}}
	{Graphs/Traversals/198.png}
\end{center}
\par Які з наведених ребер входять до складу отриманого кістякового дерева?\par
1) (0, 5)\par
2) (0, 1)\par
3) (3, 6)\par
4) (5, 6)\par
5) (1, 2)\par
6) (1, 3)\par
{\vskip 15pt} Номери відповідних ребер записати за зростанням через пробіл.


%enditem


\newpage
16. 

Рядки журналу через двокрапку містять ім'я користувача (ідентифікатор), час події,тип події, код події, наприклад
\begin{verbatim}
s="username : 2022-05-05 13:26 : ERROR : 404"
\end{verbatim}
у коді 
\begin{verbatim}
res = re.fullmatch('???', s) 
s = ''
code_str = ??? 
\end{verbatim}
чим треба замінити кожен з ???, щоб значенням code\_str був код події? 


\vspace{10mm}
\textbf{Завдання 17} 

Задано програму, що містить код 1-2 класів, можливе успадкування, та клієнтський код.
У наведеному коді невеликі частини закриті. 
Також наведено вихід програми.
Необхідно відновити закриті частини, щоб отриманий код відповідав наведеному виходу програми.


\textbf{Завдання 18-19} 

Кілька запитань за завданнями для самостійної роботи.
\vspace{10mm}
Скоріш за все, що завдань буде трохи менше, а деякі попередні будуть переформульовані в стилі двох останніх.

%enditem

\end {large}}

%end

\newpage
%begin

{{
\begin {large}

\newpage
1. Що буде надруковано за виконання?

\begin{verbatim}
print('R', end=' ')
f=('a','b','c','d','e',0,1,2,3)
print(f[0])
\end{verbatim}

%enditem


\newpage
2. Що буде надруковано за виконання?

\begin{verbatim}
print('R', end=' ')
h=(0,1,2,3,4,5,6,7,8,9)
print(h[:-13:3])
\end{verbatim}

%enditem



\newpage
3. Що буде надруковано за виконання?

\begin{verbatim}
print('R', end=' ')
h=[10,11,12,13,14]
h.extend([2,3])
print(h)
\end{verbatim}

%enditem



\newpage
4. Що буде надруковано за виконання?

\begin{verbatim}
print('R', end=' ')
e=['0','1','2','3','4']
e[4:]=[1,2]
print(e)
\end{verbatim}

%enditem



\newpage
5. Укажіть ТИП значення виразу.\par
\begin{verbatim}
{x:0 for x in range(7)}
\end{verbatim}
%enditem


\newpage
6. Що буде надруковано за виконання?\par
\begin{verbatim}
try:
    s = {25:1, 23:2, 82:5, 23:0}
    s[25] = 2
    for x, y in s.items():
        print(x, y, sep=' ', end=' ')
except:
    print('ERR')
\end{verbatim}
%enditem


\newpage
7. Записати ВИРАЗ, значенням якого є словник, ключами якого є цілі числа від 19 до 2007 (включно), що не діляться на 5, а ключам відповідають їх куби 
%enditem


\newpage
8. Що буде надруковано за виконання?
\begin{verbatim}
b = 0
k = 1

class A:
    k = 2
    
    def __init__(self):
        self.k = 3
        k = 4
        A.k = 5

obj = A()
try:
    print(b, end=' ')
    print(k, end=' ')
    print(A.k, end=' ')
    print(obj.k)
except:
    print('ERR')
\end{verbatim}
%enditem


\newpage
9. Що буде надруковано за виконання?
\begin{verbatim}
class B0:
    def c(self): print(1, end=' ')
    def h(self): print(2, end=' ')

class B1(B0):
    def c(self): print(3, end=' ')

try:
    B1().h()
    B1().c()
except:
    print('ERR')
\end{verbatim}
%enditem


\newpage
10. Що буде надруковано за виконання?
\begin{verbatim}
class A:
    d = 1

    def _g1(self): print(2, end=' ')
    def _g2(self): print(3, end=' ')

    def main(self):
        try: print(self.d, end=' ')
        except: print('E1', end=' ')

        try: self._g1()
        except: print('E2', end=' ')

        try: self._g2()
        except: print('E3', end=' ')


class B(A):
    d = 4

    def _g1(self): print(5, end=' ')
    def _g2(self): print(6, end=' ')

try: B().main()
except: print('E4', end=' ')
\end{verbatim}
%enditem


\newpage
11. Змінні \kw{next} та \kw{prev} вузла списку позначають наступний та попередній за ним вузол відповідно. Починаючи з голови, у список записано послідовні цілі числа від~$-67$ до~$44$. Нехай змінна \kw{p} позначає вузол, що зберігає значення~$7$.
Яке значення зберігається у вузлі \kw{p.prev.next.prev.next.next} ?
%enditem


\newpage
12. 
\begin{center}
	\adjustimage{max size={0.8\linewidth}{0.8\paperheight}}
	{BinTree/trees_exp/121.png}
\end{center}
Указати послідовність міток вершин дерева, зображеного на рис., що утвориться в результаті обходу дерева в прямому порядку. Мітки записувати через пробіл.\par Алгоритм починає роботу з кореня (на рис. це верхня вершина).
%enditem


\newpage
13. 
\begin{center}
	\adjustimage{max size={0.9\linewidth}{0.9\paperheight}}
	{Graphs/AdjList/73.png}
\end{center}
Для графа, зображеного на рис., вказати список суміжності вершини 1.\par
Формат оформлення відповіді:\par
список суміжності виписувати через пробіл на одному рядку, якщо пробілів десь буде більше одного - не важливо, упорядкування вершин у списку також значення не має.
%enditem


\newpage
14. Для графа на рисунку з вершини 5 запущено алгоритм пошуку в глибину, що будує кістякове дерево. Списки суміжності вершин переглядаються алгоритмом за зростанням міток вершин.\par
\begin{center}
	\adjustimage{max size={0.9\linewidth}{0.9\paperheight}}
	{Graphs/Traversals/40.png}
\end{center}
\par Які з наведених ребер входять до складу отриманого кістякового дерева?\par
1) (4, 5)\par
2) (3, 6)\par
3) (4, 6)\par
4) (0, 6)\par
5) (0, 2)\par
6) (1, 5)\par
{\vskip 15pt} Номери відповідних ребер записати за зростанням через пробіл.

%enditem


\newpage
15. Для графа на рисунку з вершини 5 запущено алгоритм пошуку в ширину, що будує кістякове дерево. Списки суміжності вершин переглядаються алгоритмом за зростанням міток вершин.\par
\begin{center}
	\adjustimage{max size={0.9\linewidth}{0.9\paperheight}}
	{Graphs/Traversals/20.png}
\end{center}
\par Які з наведених ребер входять до складу отриманого кістякового дерева?\par
1) (0, 2)\par
2) (5, 6)\par
3) (3, 6)\par
4) (4, 6)\par
5) (1, 5)\par
6) (3, 4)\par
{\vskip 15pt} Номери відповідних ребер записати за зростанням через пробіл.


%enditem


\newpage
16. 

Рядки журналу через двокрапку містять ім'я користувача (ідентифікатор), час події,тип події, код події, наприклад
\begin{verbatim}
s="username : 2022-05-05 13:26 : ERROR : 404"
\end{verbatim}
у коді 
\begin{verbatim}
res = re.fullmatch('???', s) 
s = ''
code_str = ??? 
\end{verbatim}
чим треба замінити кожен з ???, щоб значенням code\_str був код події? 


\vspace{10mm}
\textbf{Завдання 17} 

Задано програму, що містить код 1-2 класів, можливе успадкування, та клієнтський код.
У наведеному коді невеликі частини закриті. 
Також наведено вихід програми.
Необхідно відновити закриті частини, щоб отриманий код відповідав наведеному виходу програми.


\textbf{Завдання 18-19} 

Кілька запитань за завданнями для самостійної роботи.
\vspace{10mm}
Скоріш за все, що завдань буде трохи менше, а деякі попередні будуть переформульовані в стилі двох останніх.

%enditem

\end {large}}

%end

\newpage
%begin

{{
\begin {large}

\newpage
1. Що буде надруковано за виконання?

\begin{verbatim}
print('R', end=' ')
a=('a','b','c','d','e',0,1,2,3,4)
print(a[-13])
\end{verbatim}

%enditem


\newpage
2. Що буде надруковано за виконання?

\begin{verbatim}
print('R', end=' ')
e=[0,1,2,3,4,5,6,7,8,9]
print(e[::2])
\end{verbatim}

%enditem



\newpage
3. Що буде надруковано за виконання?

\begin{verbatim}
print('R', end=' ')
c=['0','1','2','3','4']
c.remove(4)
print(c)
\end{verbatim}

%enditem



\newpage
4. Що буде надруковано за виконання?

\begin{verbatim}
print('R', end=' ')
d=[10,11,12,13,14]
del d[0:0]
print(d)
\end{verbatim}

%enditem



\newpage
5. Укажіть ТИП значення виразу.\par
\begin{verbatim}
(2, 7, -4)
\end{verbatim}
%enditem


\newpage
6. Що буде надруковано за виконання?\par
\begin{verbatim}
try:
    d = {83:9, 66:0, 34:7, 29:0, 25:6}
    d[21] = 3
    for x in d.items():
        print(x, sep=' ', end=' ')
except:
    print('ERR')
\end{verbatim}
%enditem


\newpage
7. Нехай змінна \kw{s} позначає дуже довгу послідовність, по якій можна ітерувати. 

Написати вираз-генератор, що задає підпослідовність її елементів, що не є числами типу \kw{int}.

Обмеження: усі елементи шуканої підпослідовності одночасно в пам'яті розміщені бути не можуть. 
%enditem


\newpage
8. Що буде надруковано за виконання?
%Оригинал был утерян. Требует отладки!
\begin{verbatim}
c = 0

class B:
    c = 1
    
    def __init__(self):
        self.c = 2
        c = 3
        B.c = 4

obj = B()
B.c = 5
try:
    print(c, end=' ')
    print(obj.c, end=' ')
    print(B.c, end=' ')
    print(obj.c, end=' ')
    print(obj._B__c)
except:
    print('ERR')
\end{verbatim}
%enditem


\newpage
9. Що буде надруковано за виконання?
\begin{verbatim}
class C0:
    def a(self): print(1, end=' ')
    def f(self): print(2, end=' ')

class C1(C0):
    def f(self): print(3, end=' ')

try:
    C0().a()
    C1().f()
except:
    print('ERR')
\end{verbatim}
%enditem


\newpage
10. Що буде надруковано за виконання?
\begin{verbatim}
class A:
    a = 1

    def _g1(self): print(2, end=' ')
    def g2(self): print(3, end=' ')

    def main(self):
        try: print(self.a, end=' ')
        except: print('E1', end=' ')

        try: self._g1()
        except: print('E2', end=' ')

        try: self.g2()
        except: print('E3', end=' ')


class B(A):
    a = 4

    def _g1(self): print(5, end=' ')
    def g2(self): print(6, end=' ')

try: B().main()
except: print('E4', end=' ')
\end{verbatim}
%enditem


\newpage
11. Змінні \kw{next} та \kw{prev} вузла списку позначають наступний та попередній за ним вузол відповідно. Починаючи з голови, у список записано послідовні цілі числа від~$-59$ до~$48$. Нехай змінна \kw{p} позначає вузол, що зберігає значення~$5$.
Яке значення зберігається у вузлі \kw{p.next.prev.next.prev.prev} ?
%enditem


\newpage
12. 
\begin{center}
	\adjustimage{max size={0.8\linewidth}{0.8\paperheight}}
	{BinTree/trees_exp/96.png}
\end{center}
Указати послідовність міток вершин дерева, зображеного на рис., що утвориться в результаті обходу дерева в оберненому порядку. Мітки записувати через пробіл.\par Алгоритм починає роботу з кореня (на рис. це верхня вершина).
%enditem


\newpage
13. 
\begin{center}
	\adjustimage{max size={0.9\linewidth}{0.9\paperheight}}
	{Graphs/AdjList/294.png}
\end{center}
Для графа, зображеного на рис., вказати список суміжності вершини 0.\par
Формат оформлення відповіді:\par
список суміжності виписувати через пробіл на одному рядку, якщо пробілів десь буде більше одного - не важливо, упорядкування вершин у списку також значення не має.
%enditem


\newpage
14. Для графа на рисунку вказати послідовність, в якій алгоритм пошуку в глибину буде відмічати відвідані вершини, якщо списки суміжності вершин переглядаються алгоритмом за зростанням міток вершин та вершини відмічаються на першому вход\-женні у вершину.\par
Пошук розпочинається з вершини 6.\par

\begin{center}
	\adjustimage{max size={0.9\linewidth}{0.9\paperheight}}
	{Graphs/Traversals/50.png}
\end{center}
%enditem


\newpage
15. Для графа на рисунку з вершини 2 запущено алгоритм пошуку в ширину, що будує кістякове дерево. Списки суміжності вершин переглядаються алгоритмом за зростанням міток вершин.\par
\begin{center}
	\adjustimage{max size={0.9\linewidth}{0.9\paperheight}}
	{Graphs/Traversals/327.png}
\end{center}
\par Які з наведених ребер входять до складу отриманого кістякового дерева?\par
1) (0, 5)\par
2) (1, 5)\par
3) (2, 3)\par
4) (0, 6)\par
5) (0, 2)\par
6) (3, 6)\par
{\vskip 15pt} Номери відповідних ребер записати за зростанням через пробіл.


%enditem


\newpage
16. 

Рядки журналу через двокрапку містять ім'я користувача (ідентифікатор), час події,тип події, код події, наприклад
\begin{verbatim}
s="username : 2022-05-05 13:26 : ERROR : 404"
\end{verbatim}
у коді 
\begin{verbatim}
res = re.fullmatch('???', s) 
s = ''
code_str = ??? 
\end{verbatim}
чим треба замінити кожен з ???, щоб значенням code\_str був код події? 


\vspace{10mm}
\textbf{Завдання 17} 

Задано програму, що містить код 1-2 класів, можливе успадкування, та клієнтський код.
У наведеному коді невеликі частини закриті. 
Також наведено вихід програми.
Необхідно відновити закриті частини, щоб отриманий код відповідав наведеному виходу програми.


\textbf{Завдання 18-19} 

Кілька запитань за завданнями для самостійної роботи.
\vspace{10mm}
Скоріш за все, що завдань буде трохи менше, а деякі попередні будуть переформульовані в стилі двох останніх.

%enditem

\end {large}}

%end

\newpage
%begin

{{
\begin {large}

\newpage
1. Що буде надруковано за виконання?

\begin{verbatim}
print('R', end=' ')
g=('a','b','c','d','e',0,1,2,4)
print(g[-7])
\end{verbatim}

%enditem


\newpage
2. Що буде надруковано за виконання?

\begin{verbatim}
print('R', end=' ')
c=('0','1','2','3','4','5','6','7','8','9')
print(c[11:-1:-2])
\end{verbatim}

%enditem



\newpage
3. Що буде надруковано за виконання?

\begin{verbatim}
print('R', end=' ')
e=[10,11,12,13,14]
e.insert(-5, [1,0])
print(e)
\end{verbatim}

%enditem



\newpage
4. Що буде надруковано за виконання?

\begin{verbatim}
print('R', end=' ')
g=['0','1','2','3','4']
del g[6:]
print(g)
\end{verbatim}

%enditem



\newpage
5. Укажіть ТИП значення виразу.\par
\begin{verbatim}
[x for x in range(7)]
\end{verbatim}
%enditem


\newpage
6. Що буде надруковано за виконання?\par
\begin{verbatim}
try:
    d = {76:4, 56:0, 67:3, 56:0}
    d[56] = 7
    for x in d :
        print(x, sep=' ', end=' ')
except:
    print('ERR')
\end{verbatim}
%enditem


\newpage
7. Нехай змінна \kw{s} позначає дуже довгу послідовність, по якій можна ітерувати. 

Написати вираз-генератор, що задає підпослідовність її елементів, що є числами типу \kw{float}.

Обмеження: усі елементи шуканої підпослідовності одночасно в пам'яті розміщені бути не можуть. 
%enditem


\newpage
8. Що буде надруковано за виконання?
%Оригинал был утерян. Требует отладки!
\begin{verbatim}
_a = 0

class D:
    _a = 1
    
    def __init__(self):
        self._a = 2
        _a = 3
        D._a = 4

obj = D()
D._a = 5
try:
    print(_a, end=' ')
    print(obj._a, end=' ')
    print(D._a, end=' ')
    print(obj._a, end=' ')
    print(obj._D__a)
except:
    print('ERR')
\end{verbatim}
%enditem


\newpage
9. Що буде надруковано за виконання?
\begin{verbatim}
class A0:
    def d(self): print(1, end=' ')
    def f(self): print(2, end=' ')

class A1(A0):
    def d(self): print(3, end=' ')

try:
    A1().f()
    A1().d()
except:
    print('ERR')
\end{verbatim}
%enditem


\newpage
10. Що буде надруковано за виконання?
\begin{verbatim}
class A:
    a = 1

    def __init__(self):
        c = 2
        c = 3

class B(A):
    c = 4
    b = 5

    def __init__(self):
        super().__init__()
        self.b = 6

try:
    obj = B()
    obj.b = obj.c + 10
    B.b = 13
except: print('EE', end=' ')

try: print(obj.b, end= ' ')
except: print('E1', end=' ')

try: print(A.b, end= ' ')
except: print('E2', end=' ')

try: print(B.b, end= ' ')
except: print('E3', end=' ')
\end{verbatim}
%enditem


\newpage
11. Змінні \kw{next} та \kw{prev} вузла списку позначають наступний та попередній за ним вузол відповідно. Починаючи з голови, у список записано послідовні цілі числа від~$-60$ до~$28$. Нехай змінна \kw{p} позначає вузол, що зберігає значення~$3$.
Яке значення зберігається у вузлі \kw{p.next.prev.prev.prev.prev} ?
%enditem


\newpage
12. 
\begin{center}
	\adjustimage{max size={0.8\linewidth}{0.8\paperheight}}
	{BinTree/trees_exp/115.png}
\end{center}
Указати послідовність міток вершин дерева, зображеного на рис., що утвориться в результаті обходу дерева в прямому порядку. Мітки записувати через пробіл.\par Алгоритм починає роботу з кореня (на рис. це верхня вершина).
%enditem


\newpage
13. 
\begin{center}
	\adjustimage{max size={0.9\linewidth}{0.9\paperheight}}
	{Graphs/AdjList/132.png}
\end{center}
Для графа, зображеного на рис., вказати список суміжності вершини 1.\par
Формат оформлення відповіді:\par
список суміжності виписувати через пробіл на одному рядку, якщо пробілів десь буде більше одного - не важливо, упорядкування вершин у списку також значення не має.
%enditem


\newpage
14. Для графа на рисунку вказати послідовність, в якій алгоритм пошуку в глибину буде відмічати відвідані вершини, якщо списки суміжності вершин переглядаються алгоритмом за зростанням міток вершин та вершини відмічаються на першому вход\-женні у вершину.\par
Пошук розпочинається з вершини 3.\par

\begin{center}
	\adjustimage{max size={0.9\linewidth}{0.9\paperheight}}
	{Graphs/Traversals/102.png}
\end{center}
%enditem


\newpage
15. Для графа на рисунку з вершини 3 запущено алгоритм пошуку в ширину, що будує кістякове дерево. Списки суміжності вершин переглядаються алгоритмом за зростанням міток вершин.\par
\begin{center}
	\adjustimage{max size={0.9\linewidth}{0.9\paperheight}}
	{Graphs/Traversals/251.png}
\end{center}
\par Які з наведених ребер входять до складу отриманого кістякового дерева?\par
1) (0, 6)\par
2) (0, 2)\par
3) (0, 4)\par
4) (3, 5)\par
5) (4, 5)\par
6) (0, 5)\par
{\vskip 15pt} Номери відповідних ребер записати за зростанням через пробіл.


%enditem


\newpage
16. 

Рядки журналу через двокрапку містять ім'я користувача (ідентифікатор), час події,тип події, код події, наприклад
\begin{verbatim}
s="username : 2022-05-05 13:26 : ERROR : 404"
\end{verbatim}
у коді 
\begin{verbatim}
res = re.fullmatch('???', s) 
s = ''
code_str = ??? 
\end{verbatim}
чим треба замінити кожен з ???, щоб значенням code\_str був код події? 


\vspace{10mm}
\textbf{Завдання 17} 

Задано програму, що містить код 1-2 класів, можливе успадкування, та клієнтський код.
У наведеному коді невеликі частини закриті. 
Також наведено вихід програми.
Необхідно відновити закриті частини, щоб отриманий код відповідав наведеному виходу програми.


\textbf{Завдання 18-19} 

Кілька запитань за завданнями для самостійної роботи.
\vspace{10mm}
Скоріш за все, що завдань буде трохи менше, а деякі попередні будуть переформульовані в стилі двох останніх.

%enditem

\end {large}}

%end

\newpage
%begin

{{
\begin {large}

\newpage
1. Що буде надруковано за виконання?

\begin{verbatim}
print('R', end=' ')
c=('a','b','c','d','e',0,1,2)
print(c[-8])
\end{verbatim}

%enditem


\newpage
2. Що буде надруковано за виконання?

\begin{verbatim}
print('R', end=' ')
d='0123456789'
print(d[-2:-3:-3])
\end{verbatim}

%enditem



\newpage
3. Що буде надруковано за виконання?

\begin{verbatim}
print('R', end=' ')
b=['0','1','2','3','4']
b.append(18)
print(b)
\end{verbatim}

%enditem



\newpage
4. Що буде надруковано за виконання?

\begin{verbatim}
print('R', end=' ')
f=[10,11,12,13,14]
f[:5]=()
print(f)
\end{verbatim}

%enditem



\newpage
5. Укажіть ТИП значення виразу.\par
\begin{verbatim}
{1:3, 5:8}
\end{verbatim}
%enditem


\newpage
6. Що буде надруковано за виконання?\par
\begin{verbatim}
try:
    s = {18:2, 55:1, 30:1, 26:0, 26:7}
    s[50] = 9
    for x in s.items():
        print(x, sep=' ', end=' ')
except:
    print('ERR')
\end{verbatim}
%enditem


\newpage
7. Нехай змінна \kw{s} позначає послідовність цілих, по якій можна ітерувати.

Написати вираз, що перевіряє, чи правда, що всі елементи послідовності є цілими.
%enditem


\newpage
8. Що буде надруковано за виконання?
\begin{verbatim}
c = 0
j = 1

class B:
    j = 2
    
    def __init__(self):
        global j
        self.c = 3
        j = 4
        B.j = 5

obj = B()
try:
    print(c, end=' ')
    print(j, end=' ')
    print(B.j, end=' ')
    print(obj.c)
except:
    print('ERR')
\end{verbatim}
%enditem


\newpage
9. Що буде надруковано за виконання?
\begin{verbatim}
class D0:
    def b(self): print(1, end=' ')
    def h(self): print(2, end=' ')

class D1(D0):
    def h(self): print(3, end=' ')

try:
    D1().b()
    D0().h()
except:
    print('ERR')
\end{verbatim}
%enditem


\newpage
10. Що буде надруковано за виконання?
\begin{verbatim}
class A:
    b = 1

    def __init__(self):
        a = 2
        b = 3

class B(A):
    b = 4
    c = 5

    def __init__(self):
        super().__init__()
        self.c = 6

try:
    obj = B()
    obj.c = obj.b + 10
    B.c = 13
except: print('EE', end=' ')

try: print(obj.c, end= ' ')
except: print('E1', end=' ')

try: print(A.c, end= ' ')
except: print('E2', end=' ')

try: print(B.c, end= ' ')
except: print('E3', end=' ')
\end{verbatim}
%enditem


\newpage
11. Змінні \kw{next} та \kw{prev} вузла списку позначають наступний та попередній за ним вузол відповідно. Починаючи з голови, у список записано послідовні цілі числа від~$-41$ до~$24$. Нехай змінна \kw{p} позначає вузол, що зберігає значення~$-3$.
Яке значення зберігається у вузлі \kw{p.prev.next.next.next.next} ?
%enditem


\newpage
12. 
\begin{center}
	\adjustimage{max size={0.8\linewidth}{0.8\paperheight}}
	{BinTree/trees_exp/73.png}
\end{center}
Указати послідовність міток вершин дерева, зображеного на рис., що утвориться в результаті обходу дерева в оберненому порядку. Мітки записувати через пробіл.\par Алгоритм починає роботу з кореня (на рис. це верхня вершина).
%enditem


\newpage
13. 
\begin{center}
	\adjustimage{max size={0.9\linewidth}{0.9\paperheight}}
	{Graphs/AdjList/10.png}
\end{center}
Для графа, зображеного на рис., вказати список суміжності вершини 1.\par
Формат оформлення відповіді:\par
список суміжності виписувати через пробіл на одному рядку, якщо пробілів десь буде більше одного - не важливо, упорядкування вершин у списку також значення не має.
%enditem


\newpage
14. Для графа на рисунку вказати послідовність, в якій алгоритм пошуку в глибину буде відмічати відвідані вершини, якщо списки суміжності вершин переглядаються алгоритмом за зростанням міток вершин та вершини відмічаються на першому вход\-женні у вершину.\par
Пошук розпочинається з вершини 5.\par

\begin{center}
	\adjustimage{max size={0.9\linewidth}{0.9\paperheight}}
	{Graphs/Traversals/20.png}
\end{center}
%enditem


\newpage
15. Для графа на рисунку з вершини 2 запущено алгоритм пошуку в ширину, що будує кістякове дерево. Списки суміжності вершин переглядаються алгоритмом за зростанням міток вершин.\par
\begin{center}
	\adjustimage{max size={0.9\linewidth}{0.9\paperheight}}
	{Graphs/Traversals/89.png}
\end{center}
\par Які з наведених ребер входять до складу отриманого кістякового дерева?\par
1) (0, 1)\par
2) (1, 2)\par
3) (1, 3)\par
4) (0, 5)\par
5) (2, 5)\par
6) (0, 4)\par
{\vskip 15pt} Номери відповідних ребер записати за зростанням через пробіл.


%enditem


\newpage
16. 

Рядки журналу через двокрапку містять ім'я користувача (ідентифікатор), час події,тип події, код події, наприклад
\begin{verbatim}
s="username : 2022-05-05 13:26 : ERROR : 404"
\end{verbatim}
у коді 
\begin{verbatim}
res = re.fullmatch('???', s) 
s = ''
code_str = ??? 
\end{verbatim}
чим треба замінити кожен з ???, щоб значенням code\_str був код події? 


\vspace{10mm}
\textbf{Завдання 17} 

Задано програму, що містить код 1-2 класів, можливе успадкування, та клієнтський код.
У наведеному коді невеликі частини закриті. 
Також наведено вихід програми.
Необхідно відновити закриті частини, щоб отриманий код відповідав наведеному виходу програми.


\textbf{Завдання 18-19} 

Кілька запитань за завданнями для самостійної роботи.
\vspace{10mm}
Скоріш за все, що завдань буде трохи менше, а деякі попередні будуть переформульовані в стилі двох останніх.

%enditem

\end {large}}

%end

\newpage
%begin

{{
\begin {large}

\newpage
1. Що буде надруковано за виконання?

\begin{verbatim}
print('R', end=' ')
b=('a','b','c','d','e',0,1)
print(b[-4])
\end{verbatim}

%enditem


\newpage
2. Що буде надруковано за виконання?

\begin{verbatim}
print('R', end=' ')
f=(0,1,2,3,4,5,6,7,8,9)
print(f[-10::-1])
\end{verbatim}

%enditem



\newpage
3. Що буде надруковано за виконання?

\begin{verbatim}
print('R', end=' ')
d=['0','1','2','3','4']
print(d.index('1'), end=' ')
\end{verbatim}

%enditem



\newpage
4. Що буде надруковано за виконання?

\begin{verbatim}
print('R', end=' ')
b=[10,11,12,13,14]
b[-5:1]=[1,0]
print(b)
\end{verbatim}

%enditem



\newpage
5. Укажіть ТИП значення виразу.\par
\begin{verbatim}
{21,2,4}
\end{verbatim}
%enditem


\newpage
6. Що буде надруковано за виконання?\par
\begin{verbatim}
try:
    t = {80:0, 47:1, 35:8, 88:5, 47:1}
    t[77] = 9
    for x in t.values():
        print(x, sep=' ', end=' ')
except:
    print('ERR')
\end{verbatim}
%enditem


\newpage
7. Нехай змінна \kw{s} позначає дуже довгу послідовність, по якій можна ітерувати. 

Написати вираз-генератор, що задає підпослідовність її елементів, що не є рядками типу \kw{str}.

Обмеження: усі елементи шуканої підпослідовності одночасно в пам'яті розміщені бути не можуть. 
%enditem


\newpage
8. Що буде надруковано за виконання?
\begin{verbatim}
a = 0
h = 1

class C:
    a = 2
    
    def __init__(self):
        self.h = 3
        a = 4
        C.a = 5

obj = C()
try:
    print(a, end=' ')
    print(h, end=' ')
    print(C.a, end=' ')
    print(obj.h)
except:
    print('ERR')
\end{verbatim}
%enditem


\newpage
9. Що буде надруковано за виконання?
\begin{verbatim}
class B0:
    def c(self): print(1, end=' ')
    def g(self): print(2, end=' ')

class B1(B0):
    def g(self): print(3, end=' ')

try:
    B1().c()
    B1().g()
except:
    print('ERR')
\end{verbatim}
%enditem


\newpage
10. Що буде надруковано за виконання?
\begin{verbatim}
class A:
    c = 1

    def h1(self): print(2, end=' ')
    def _h2(self): print(3, end=' ')

    def main(self):
        try: print(self.c, end=' ')
        except: print('E1', end=' ')

        try: self.h1()
        except: print('E2', end=' ')

        try: self._h2()
        except: print('E3', end=' ')


class B(A):
    c = 4

    def h1(self): print(5, end=' ')
    def _h2(self): print(6, end=' ')

try: B().main()
except: print('E4', end=' ')
\end{verbatim}
%enditem


\newpage
11. Змінні \kw{next} та \kw{prev} вузла списку позначають наступний та попередній за ним вузол відповідно. Починаючи з голови, у список записано послідовні цілі числа від~$-27$ до~$17$. Нехай змінна \kw{p} позначає вузол, що зберігає значення~$-9$.
Яке значення зберігається у вузлі \kw{p.next.next.next.prev.prev} ?
%enditem


\newpage
12. 
\begin{center}
	\adjustimage{max size={0.8\linewidth}{0.8\paperheight}}
	{BinTree/trees_exp/135.png}
\end{center}
Указати послідовність міток вершин дерева, зображеного на рис., що утвориться в результаті обходу дерева в оберненому порядку. Мітки записувати через пробіл.\par Алгоритм починає роботу з кореня (на рис. це верхня вершина).
%enditem


\newpage
13. 
\begin{center}
	\adjustimage{max size={0.9\linewidth}{0.9\paperheight}}
	{Graphs/AdjList/76.png}
\end{center}
Для графа, зображеного на рис., вказати список суміжності вершини 2.\par
Формат оформлення відповіді:\par
список суміжності виписувати через пробіл на одному рядку, якщо пробілів десь буде більше одного - не важливо, упорядкування вершин у списку також значення не має.
%enditem


\newpage
14. Для графа на рисунку з вершини 5 запущено алгоритм пошуку в глибину, що будує кістякове дерево. Списки суміжності вершин переглядаються алгоритмом за зростанням міток вершин.\par
\begin{center}
	\adjustimage{max size={0.9\linewidth}{0.9\paperheight}}
	{Graphs/Traversals/338.png}
\end{center}
\par Які з наведених ребер входять до складу отриманого кістякового дерева?\par
1) (3, 6)\par
2) (1, 2)\par
3) (1, 3)\par
4) (0, 4)\par
5) (2, 3)\par
6) (0, 6)\par
{\vskip 15pt} Номери відповідних ребер записати за зростанням через пробіл.

%enditem


\newpage
15. Для графа на рисунку з вершини 5 запущено алгоритм пошуку в ширину, що будує кістякове дерево. Списки суміжності вершин переглядаються алгоритмом за зростанням міток вершин.\par
\begin{center}
	\adjustimage{max size={0.9\linewidth}{0.9\paperheight}}
	{Graphs/Traversals/207.png}
\end{center}
\par Які з наведених ребер входять до складу отриманого кістякового дерева?\par
1) (2, 6)\par
2) (3, 5)\par
3) (0, 2)\par
4) (1, 3)\par
5) (1, 4)\par
6) (1, 2)\par
{\vskip 15pt} Номери відповідних ребер записати за зростанням через пробіл.


%enditem


\newpage
16. 

Рядки журналу через двокрапку містять ім'я користувача (ідентифікатор), час події,тип події, код події, наприклад
\begin{verbatim}
s="username : 2022-05-05 13:26 : ERROR : 404"
\end{verbatim}
у коді 
\begin{verbatim}
res = re.fullmatch('???', s) 
s = ''
code_str = ??? 
\end{verbatim}
чим треба замінити кожен з ???, щоб значенням code\_str був код події? 


\vspace{10mm}
\textbf{Завдання 17} 

Задано програму, що містить код 1-2 класів, можливе успадкування, та клієнтський код.
У наведеному коді невеликі частини закриті. 
Також наведено вихід програми.
Необхідно відновити закриті частини, щоб отриманий код відповідав наведеному виходу програми.


\textbf{Завдання 18-19} 

Кілька запитань за завданнями для самостійної роботи.
\vspace{10mm}
Скоріш за все, що завдань буде трохи менше, а деякі попередні будуть переформульовані в стилі двох останніх.

%enditem

\end {large}}

%end

\newpage
%begin

{{
\begin {large}

\newpage
1. Що буде надруковано за виконання?

\begin{verbatim}
print('R', end=' ')
h=('a','b','c','d','e',0,1,2,3)
print(h[9])
\end{verbatim}

%enditem


\newpage
2. Що буде надруковано за виконання?

\begin{verbatim}
print('R', end=' ')
g=[0,1,2,3,4,5,6,7,8,9]
print(g[-9::2])
\end{verbatim}

%enditem



\newpage
3. Що буде надруковано за виконання?

\begin{verbatim}
print('R', end=' ')
a=[10,11,12,13,14]
a.insert(-6, [1,2])
print(a)
\end{verbatim}

%enditem



\newpage
4. Що буде надруковано за виконання?

\begin{verbatim}
print('R', end=' ')
a=['0','1','2','3','4']
a[3:3]='11'
print(a)
\end{verbatim}

%enditem



\newpage
5. Укажіть ТИП значення виразу.\par
\begin{verbatim}
[13,2,4]+[3,5]
\end{verbatim}
%enditem


\newpage
6. Що буде надруковано за виконання?\par
\begin{verbatim}
try:
    d = {46:1, 11:6, 34:9, 34:9}
    d[12] = 7
    for x, y in d.items():
        print(x, y, sep=' ', end=' ')
except:
    print('ERR')
\end{verbatim}
%enditem


\newpage
7. Нехай змінна \kw{s} позначає дуже довгу послідовність цілих, по якій можна ітерувати. 

Написати вираз-генератор, що задає підпослідовність її елементів, що діляться на 3. 

Обмеження: усі елементи шуканої підпослідовності одночасно в пам'яті розміщені бути не можуть. 
%enditem


\newpage
8. Що буде надруковано за виконання?
\begin{verbatim}
c = 0
n = 1

class D:
    c = 2
    
    def __init__(self):
        self.n = 3
        c = 4
        D.c = 5

obj = D()
try:
    print(c, end=' ')
    print(n, end=' ')
    print(D.c, end=' ')
    print(obj.n)
except:
    print('ERR')
\end{verbatim}
%enditem


\newpage
9. Що буде надруковано за виконання?
\begin{verbatim}
class A0:
    def d(self): print(1, end=' ')
    def f(self): print(2, end=' ')

class A1(A0):
    def d(self): print(3, end=' ')

try:
    A1().f()
    A0().d()
except:
    print('ERR')
\end{verbatim}
%enditem


\newpage
10. Що буде надруковано за виконання?
\begin{verbatim}
class A:
    b = 1

    def _f1(self): print(2, end=' ')
    def _f2(self): print(3, end=' ')

    def main(self):
        try: print(self.b, end=' ')
        except: print('E1', end=' ')

        try: self._f1()
        except: print('E2', end=' ')

        try: self._f2()
        except: print('E3', end=' ')


class B(A):
    b = 4

    def _f1(self): print(5, end=' ')
    def _f2(self): print(6, end=' ')

try: B().main()
except: print('E4', end=' ')
\end{verbatim}
%enditem


\newpage
11. Змінні \kw{next} та \kw{prev} вузла списку позначають наступний та попередній за ним вузол відповідно. Починаючи з голови, у список записано послідовні цілі числа від~$-29$ до~$47$. Нехай змінна \kw{p} позначає вузол, що зберігає значення~$-2$.
Яке значення зберігається у вузлі \kw{p.prev.prev.prev.next.next} ?
%enditem


\newpage
12. 
\begin{center}
	\adjustimage{max size={0.8\linewidth}{0.8\paperheight}}
	{BinTree/trees_exp/179.png}
\end{center}
Указати послідовність міток вершин дерева, зображеного на рис., що утвориться в результаті обходу дерева в прямому порядку. Мітки записувати через пробіл.\par Алгоритм починає роботу з кореня (на рис. це верхня вершина).
%enditem


\newpage
13. 
\begin{center}
	\adjustimage{max size={0.9\linewidth}{0.9\paperheight}}
	{Graphs/AdjList/183.png}
\end{center}
Для графа, зображеного на рис., вказати список суміжності вершини 1.\par
Формат оформлення відповіді:\par
список суміжності виписувати через пробіл на одному рядку, якщо пробілів десь буде більше одного - не важливо, упорядкування вершин у списку також значення не має.
%enditem


\newpage
14. Для графа на рисунку з вершини 2 запущено алгоритм пошуку в глибину, що будує кістякове дерево. Списки суміжності вершин переглядаються алгоритмом за зростанням міток вершин.\par
\begin{center}
	\adjustimage{max size={0.9\linewidth}{0.9\paperheight}}
	{Graphs/Traversals/264.png}
\end{center}
\par Які з наведених ребер входять до складу отриманого кістякового дерева?\par
1) (1, 3)\par
2) (0, 5)\par
3) (0, 2)\par
4) (2, 5)\par
5) (2, 4)\par
6) (5, 6)\par
{\vskip 15pt} Номери відповідних ребер записати за зростанням через пробіл.

%enditem


\newpage
15. Для графа на рисунку з вершини 5 запущено алгоритм пошуку в ширину, що будує кістякове дерево. Списки суміжності вершин переглядаються алгоритмом за зростанням міток вершин.\par
\begin{center}
	\adjustimage{max size={0.9\linewidth}{0.9\paperheight}}
	{Graphs/Traversals/193.png}
\end{center}
\par Які з наведених ребер входять до складу отриманого кістякового дерева?\par
1) (3, 6)\par
2) (0, 1)\par
3) (1, 4)\par
4) (3, 4)\par
5) (2, 6)\par
6) (1, 6)\par
{\vskip 15pt} Номери відповідних ребер записати за зростанням через пробіл.


%enditem


\newpage
16. 

Рядки журналу через двокрапку містять ім'я користувача (ідентифікатор), час події,тип події, код події, наприклад
\begin{verbatim}
s="username : 2022-05-05 13:26 : ERROR : 404"
\end{verbatim}
у коді 
\begin{verbatim}
res = re.fullmatch('???', s) 
s = ''
code_str = ??? 
\end{verbatim}
чим треба замінити кожен з ???, щоб значенням code\_str був код події? 


\vspace{10mm}
\textbf{Завдання 17} 

Задано програму, що містить код 1-2 класів, можливе успадкування, та клієнтський код.
У наведеному коді невеликі частини закриті. 
Також наведено вихід програми.
Необхідно відновити закриті частини, щоб отриманий код відповідав наведеному виходу програми.


\textbf{Завдання 18-19} 

Кілька запитань за завданнями для самостійної роботи.
\vspace{10mm}
Скоріш за все, що завдань буде трохи менше, а деякі попередні будуть переформульовані в стилі двох останніх.

%enditem

\end {large}}

%end

\newpage
%begin

{{
\begin {large}

\newpage
1. Що буде надруковано за виконання?

\begin{verbatim}
print('R', end=' ')
e=('a','b','c','d','e',0,1,2,3,4)
print(e[3])
\end{verbatim}

%enditem


\newpage
2. Що буде надруковано за виконання?

\begin{verbatim}
print('R', end=' ')
a=['0','1','2','3','4','5','6','7','8','9']
print(a[:-15:3])
\end{verbatim}

%enditem



\newpage
3. Що буде надруковано за виконання?

\begin{verbatim}
print('R', end=' ')
f=[10,11,12,13,14]
f.remove('10')
print(f)
\end{verbatim}

%enditem



\newpage
4. Що буде надруковано за виконання?

\begin{verbatim}
print('R', end=' ')
c=['0','1','2','3','4']
del c[6:-4]
print(c)
\end{verbatim}

%enditem



\newpage
5. Укажіть ТИП значення виразу.\par
\begin{verbatim}
{(1,4), 8}
\end{verbatim}
%enditem


\newpage
6. Що буде надруковано за виконання?\par
\begin{verbatim}
try:
    t = {16:9, 68:5, 43:9, 40:5, 68:8}
    t[72] = 0
    for x in t.keys():
        print(x, sep=' ', end=' ')
except:
    print('ERR')
\end{verbatim}
%enditem


\newpage
7. 
Записати ВИРАЗ, значенням якого є список, що складається з цілих чисел від 18 до 2007 (включно), що діляться на 3, записаних в обернено\-му порядку.
%enditem


\newpage
8. Що буде надруковано за виконання?
%Оригинал был утерян. Требует отладки!
\begin{verbatim}
_b = 0

class A:
    _b = 1
    
    def __init__(self):
        self._b = 2
        _b = 3
        A._b = 4

obj = A()
obj._b = 5
try:
    print(_b, end=' ')
    print(obj._b, end=' ')
    print(A._b, end=' ')
    print(obj._b, end=' ')
    print(obj._A__b)
except:
    print('ERR')
\end{verbatim}
%enditem


\newpage
9. Що буде надруковано за виконання?
\begin{verbatim}
class D0:
    def a(self): print(1, end=' ')
    def h(self): print(2, end=' ')

class D1(D0):
    def a(self): print(3, end=' ')

try:
    D1().h()
    D1().a()
except:
    print('ERR')
\end{verbatim}
%enditem


\newpage
10. Що буде надруковано за виконання?
\begin{verbatim}
class A:
    c = 1

    def __init__(self):
        a = 2
        a = 3

class B(A):
    a = 4
    b = 5

    def __init__(self):
        super().__init__()
        self.b = 6

try:
    obj = B()
    obj.a = obj.c + 10
    B.a = 13
except: print('EE', end=' ')

try: print(obj.a, end= ' ')
except: print('E1', end=' ')

try: print(A.a, end= ' ')
except: print('E2', end=' ')

try: print(B.a, end= ' ')
except: print('E3', end=' ')
\end{verbatim}
%enditem


\newpage
11. Змінні \kw{next} та \kw{prev} вузла списку позначають наступний та попередній за ним вузол відповідно. Починаючи з голови, у список записано послідовні цілі числа від~$-39$ до~$40$. Нехай змінна \kw{p} позначає вузол, що зберігає значення~$8$.
Яке значення зберігається у вузлі \kw{p.next.prev.next.prev.next} ?
%enditem


\newpage
12. 
\begin{center}
	\adjustimage{max size={0.8\linewidth}{0.8\paperheight}}
	{BinTree/trees_exp/43.png}
\end{center}
Указати послідовність міток вершин дерева, зображеного на рис., що утвориться в результаті обходу дерева в прямому порядку. Мітки записувати через пробіл.\par Алгоритм починає роботу з кореня (на рис. це верхня вершина).
%enditem


\newpage
13. 
\begin{center}
	\adjustimage{max size={0.9\linewidth}{0.9\paperheight}}
	{Graphs/AdjList/258.png}
\end{center}
Для графа, зображеного на рис., вказати список суміжності вершини 0.\par
Формат оформлення відповіді:\par
список суміжності виписувати через пробіл на одному рядку, якщо пробілів десь буде більше одного - не важливо, упорядкування вершин у списку також значення не має.
%enditem


\newpage
14. Для графа на рисунку вказати послідовність, в якій алгоритм пошуку в глибину буде відмічати відвідані вершини, якщо списки суміжності вершин переглядаються алгоритмом за зростанням міток вершин та вершини відмічаються на першому вход\-женні у вершину.\par
Пошук розпочинається з вершини 0.\par

\begin{center}
	\adjustimage{max size={0.9\linewidth}{0.9\paperheight}}
	{Graphs/Traversals/76.png}
\end{center}
%enditem


\newpage
15. Для графа на рисунку з вершини 0 запущено алгоритм пошуку в ширину, що будує кістякове дерево. Списки суміжності вершин переглядаються алгоритмом за зростанням міток вершин.\par
\begin{center}
	\adjustimage{max size={0.9\linewidth}{0.9\paperheight}}
	{Graphs/Traversals/179.png}
\end{center}
\par Які з наведених ребер входять до складу отриманого кістякового дерева?\par
1) (1, 6)\par
2) (0, 2)\par
3) (2, 4)\par
4) (4, 5)\par
5) (0, 4)\par
6) (0, 6)\par
{\vskip 15pt} Номери відповідних ребер записати за зростанням через пробіл.


%enditem


\newpage
16. 

Рядки журналу через двокрапку містять ім'я користувача (ідентифікатор), час події,тип події, код події, наприклад
\begin{verbatim}
s="username : 2022-05-05 13:26 : ERROR : 404"
\end{verbatim}
у коді 
\begin{verbatim}
res = re.fullmatch('???', s) 
s = ''
code_str = ??? 
\end{verbatim}
чим треба замінити кожен з ???, щоб значенням code\_str був код події? 


\vspace{10mm}
\textbf{Завдання 17} 

Задано програму, що містить код 1-2 класів, можливе успадкування, та клієнтський код.
У наведеному коді невеликі частини закриті. 
Також наведено вихід програми.
Необхідно відновити закриті частини, щоб отриманий код відповідав наведеному виходу програми.


\textbf{Завдання 18-19} 

Кілька запитань за завданнями для самостійної роботи.
\vspace{10mm}
Скоріш за все, що завдань буде трохи менше, а деякі попередні будуть переформульовані в стилі двох останніх.

%enditem

\end {large}}

%end

\newpage
%begin

{{
\begin {large}

\newpage
1. Що буде надруковано за виконання?

\begin{verbatim}
print('R', end=' ')
d=('a','b','c','d','e',0,1,2,4)
print(d[-9])
\end{verbatim}

%enditem


\newpage
2. Що буде надруковано за виконання?

\begin{verbatim}
print('R', end=' ')
b=[0,1,2,3,4,5,6,7,8,9]
print(b[10:-4:2])
\end{verbatim}

%enditem



\newpage
3. Що буде надруковано за виконання?

\begin{verbatim}
print('R', end=' ')
g=['0','1','2','3','4']
print(g.pop(-5), end=' ')
print(g)
\end{verbatim}

%enditem



\newpage
4. Що буде надруковано за виконання?

\begin{verbatim}
print('R', end=' ')
h=[10,11,12,13,14]
h[-4:]=[2,3]
print(h)
\end{verbatim}

%enditem



\newpage
5. Укажіть ТИП значення виразу.\par
\begin{verbatim}
{x:x for x in range(7)}
\end{verbatim}
%enditem


\newpage
6. Що буде надруковано за виконання?\par
\begin{verbatim}
try:
    s = {87:5, 64:6, 24:5, 35:7, 35:2}
    s[16] = 2
    for x in s.items():
        print(*x, sep=' ', end=' ')
except:
    print('ERR')
\end{verbatim}
%enditem


\newpage
7. Нехай змінна \kw{s} позначає послідовність цілих, по якій можна ітерувати. 

Написати вираз, що перевіряє, чи правда, що послідовність  містить числа, більші за 5.
%enditem


\newpage
8. Що буде надруковано за виконання?
\begin{verbatim}
b = 0
m = 1

class C:
    m = 2
    
    def __init__(self):
        global m
        self.b = 3
        m = 4
        C.m = 5

obj = C()
try:
    print(b, end=' ')
    print(m, end=' ')
    print(C.m, end=' ')
    print(obj.b)
except:
    print('ERR')
\end{verbatim}
%enditem


\newpage
9. Що буде надруковано за виконання?
\begin{verbatim}
class C0:
    def b(self): print(1, end=' ')
    def g(self): print(2, end=' ')

class C1(C0):
    def g(self): print(3, end=' ')

try:
    C0().b()
    C1().g()
except:
    print('ERR')
\end{verbatim}
%enditem


\newpage
10. Що буде надруковано за виконання?
\begin{verbatim}
class A:
    d = 1

    def __h1(self): print(2, end=' ')
    def _h2(self): print(3, end=' ')

    def main(self):
        try: print(self.d, end=' ')
        except: print('E1', end=' ')

        try: self.__h1()
        except: print('E2', end=' ')

        try: self._h2()
        except: print('E3', end=' ')


class B(A):
    d = 4

    def __h1(self): print(5, end=' ')
    def _h2(self): print(6, end=' ')

try: B().main()
except: print('E4', end=' ')
\end{verbatim}
%enditem


\newpage
11. Змінні \kw{next} та \kw{prev} вузла списку позначають наступний та попередній за ним вузол відповідно. Починаючи з голови, у список записано послідовні цілі числа від~$-25$ до~$18$. Нехай змінна \kw{p} позначає вузол, що зберігає значення~$4$.
Яке значення зберігається у вузлі \kw{p.prev.next.prev.next.prev} ?
%enditem


\newpage
12. 
\begin{center}
	\adjustimage{max size={0.8\linewidth}{0.8\paperheight}}
	{BinTree/trees_exp/132.png}
\end{center}
Указати послідовність міток вершин дерева, зображеного на рис., що утвориться в результаті обходу дерева в оберненому порядку. Мітки записувати через пробіл.\par Алгоритм починає роботу з кореня (на рис. це верхня вершина).
%enditem


\newpage
13. 
\begin{center}
	\adjustimage{max size={0.9\linewidth}{0.9\paperheight}}
	{Graphs/AdjList/134.png}
\end{center}
Для графа, зображеного на рис., вказати список суміжності вершини 8.\par
Формат оформлення відповіді:\par
список суміжності виписувати через пробіл на одному рядку, якщо пробілів десь буде більше одного - не важливо, упорядкування вершин у списку також значення не має.
%enditem


\newpage
14. Для графа на рисунку вказати послідовність, в якій алгоритм пошуку в глибину буде відмічати відвідані вершини, якщо списки суміжності вершин переглядаються алгоритмом за зростанням міток вершин та вершини відмічаються на першому вход\-женні у вершину.\par
Пошук розпочинається з вершини 4.\par

\begin{center}
	\adjustimage{max size={0.9\linewidth}{0.9\paperheight}}
	{Graphs/Traversals/239.png}
\end{center}
%enditem


\newpage
15. Для графа на рисунку з вершини 0 запущено алгоритм пошуку в ширину, що будує кістякове дерево. Списки суміжності вершин переглядаються алгоритмом за зростанням міток вершин.\par
\begin{center}
	\adjustimage{max size={0.9\linewidth}{0.9\paperheight}}
	{Graphs/Traversals/199.png}
\end{center}
\par Які з наведених ребер входять до складу отриманого кістякового дерева?\par
1) (2, 3)\par
2) (4, 5)\par
3) (2, 4)\par
4) (3, 5)\par
5) (3, 4)\par
6) (5, 6)\par
{\vskip 15pt} Номери відповідних ребер записати за зростанням через пробіл.


%enditem


\newpage
16. 

Рядки журналу через двокрапку містять ім'я користувача (ідентифікатор), час події,тип події, код події, наприклад
\begin{verbatim}
s="username : 2022-05-05 13:26 : ERROR : 404"
\end{verbatim}
у коді 
\begin{verbatim}
res = re.fullmatch('???', s) 
s = ''
code_str = ??? 
\end{verbatim}
чим треба замінити кожен з ???, щоб значенням code\_str був код події? 


\vspace{10mm}
\textbf{Завдання 17} 

Задано програму, що містить код 1-2 класів, можливе успадкування, та клієнтський код.
У наведеному коді невеликі частини закриті. 
Також наведено вихід програми.
Необхідно відновити закриті частини, щоб отриманий код відповідав наведеному виходу програми.


\textbf{Завдання 18-19} 

Кілька запитань за завданнями для самостійної роботи.
\vspace{10mm}
Скоріш за все, що завдань буде трохи менше, а деякі попередні будуть переформульовані в стилі двох останніх.

%enditem

\end {large}}

%end

\newpage
%begin

{{
\begin {large}

\newpage
1. Що буде надруковано за виконання?

\begin{verbatim}
print('R', end=' ')
d=('a','b','c','d','e',0,1,2,3)
print(d[-7])
\end{verbatim}

%enditem


\newpage
2. Що буде надруковано за виконання?

\begin{verbatim}
print('R', end=' ')
f=('0','1','2','3','4','5','6','7','8','9')
print(f[8:0:3])
\end{verbatim}

%enditem



\newpage
3. Що буде надруковано за виконання?

\begin{verbatim}
print('R', end=' ')
c=['0','1','2','3','4']
c.extend('10')
print(c)
\end{verbatim}

%enditem



\newpage
4. Що буде надруковано за виконання?

\begin{verbatim}
print('R', end=' ')
h=['0','1','2','3','4']
h[:0]=[15]
print(h)
\end{verbatim}

%enditem



\newpage
5. Укажіть ТИП значення виразу.\par
\begin{verbatim}
{x for x in range(7)}
\end{verbatim}
%enditem


\newpage
6. Що буде надруковано за виконання?\par
\begin{verbatim}
try:
    d = {37:4, 39:7, 75:1, 72:8, 72:6}
    d[77] = 4
    for x, y in d.items():
        print(x, y, sep=' ', end=' ')
except:
    print('ERR')
\end{verbatim}
%enditem


\newpage
7. Нехай змінна \kw{s} позначає послідовність цілих, по якій можна ітерувати. 

Написати вираз, що перевіряє, чи правда, що послідовність  містить від’ємні числа.
%enditem


\newpage
8. Що буде надруковано за виконання?
%Оригинал был утерян. Требует отладки!
\begin{verbatim}
__c = 0

class C:
    __c = 1
    
    def __init__(self):
        self.__c = 2
        __c = 3
        C.__c = 4

obj = C()
obj.__c = 5
try:
    print(__c, end=' ')
    print(obj.__c, end=' ')
    print(C.__c, end=' ')
    print(obj.__c, end=' ')
    print(obj._C__c)
except:
    print('ERR')
\end{verbatim}
%enditem


\newpage
9. Що буде надруковано за виконання?
\begin{verbatim}
class A0:
    def c(self): print(1, end=' ')
    def h(self): print(2, end=' ')

class A1(A0):
    def c(self): print(3, end=' ')

try:
    A1().h()
    A1().c()
except:
    print('ERR')
\end{verbatim}
%enditem


\newpage
10. Що буде надруковано за виконання?
\begin{verbatim}
class A:
    c = 1

    def _f1(self): print(2, end=' ')
    def __f2(self): print(3, end=' ')

    def main(self):
        try: print(self.c, end=' ')
        except: print('E1', end=' ')

        try: self._f1()
        except: print('E2', end=' ')

        try: self.__f2()
        except: print('E3', end=' ')


class B(A):
    c = 4

    def _f1(self): print(5, end=' ')
    def __f2(self): print(6, end=' ')

try: B().main()
except: print('E4', end=' ')
\end{verbatim}
%enditem


\newpage
11. Змінні \kw{next} та \kw{prev} вузла списку позначають наступний та попередній за ним вузол відповідно. Починаючи з голови, у список записано послідовні цілі числа від~$-22$ до~$25$. Нехай змінна \kw{p} позначає вузол, що зберігає значення~$-9$.
Яке значення зберігається у вузлі \kw{p.prev.next.next.next.prev} ?
%enditem


\newpage
12. 
\begin{center}
	\adjustimage{max size={0.8\linewidth}{0.8\paperheight}}
	{BinTree/trees_exp/10.png}
\end{center}
Указати послідовність міток вершин дерева, зображеного на рис., що утвориться в результаті обходу дерева в прямому порядку. Мітки записувати через пробіл.\par Алгоритм починає роботу з кореня (на рис. це верхня вершина).
%enditem


\newpage
13. 
\begin{center}
	\adjustimage{max size={0.9\linewidth}{0.9\paperheight}}
	{Graphs/AdjList/78.png}
\end{center}
Для графа, зображеного на рис., вказати список суміжності вершини 2.\par
Формат оформлення відповіді:\par
список суміжності виписувати через пробіл на одному рядку, якщо пробілів десь буде більше одного - не важливо, упорядкування вершин у списку також значення не має.
%enditem


\newpage
14. Для графа на рисунку вказати послідовність, в якій алгоритм пошуку в глибину буде відмічати відвідані вершини, якщо списки суміжності вершин переглядаються алгоритмом за зростанням міток вершин та вершини відмічаються на першому вход\-женні у вершину.\par
Пошук розпочинається з вершини 2.\par

\begin{center}
	\adjustimage{max size={0.9\linewidth}{0.9\paperheight}}
	{Graphs/Traversals/53.png}
\end{center}
%enditem


\newpage
15. Для графа на рисунку з вершини 3 запущено алгоритм пошуку в ширину, що будує кістякове дерево. Списки суміжності вершин переглядаються алгоритмом за зростанням міток вершин.\par
\begin{center}
	\adjustimage{max size={0.9\linewidth}{0.9\paperheight}}
	{Graphs/Traversals/52.png}
\end{center}
\par Які з наведених ребер входять до складу отриманого кістякового дерева?\par
1) (0, 1)\par
2) (4, 5)\par
3) (0, 4)\par
4) (3, 5)\par
5) (4, 6)\par
6) (2, 3)\par
{\vskip 15pt} Номери відповідних ребер записати за зростанням через пробіл.


%enditem


\newpage
16. 

Рядки журналу через двокрапку містять ім'я користувача (ідентифікатор), час події,тип події, код події, наприклад
\begin{verbatim}
s="username : 2022-05-05 13:26 : ERROR : 404"
\end{verbatim}
у коді 
\begin{verbatim}
res = re.fullmatch('???', s) 
s = ''
code_str = ??? 
\end{verbatim}
чим треба замінити кожен з ???, щоб значенням code\_str був код події? 


\vspace{10mm}
\textbf{Завдання 17} 

Задано програму, що містить код 1-2 класів, можливе успадкування, та клієнтський код.
У наведеному коді невеликі частини закриті. 
Також наведено вихід програми.
Необхідно відновити закриті частини, щоб отриманий код відповідав наведеному виходу програми.


\textbf{Завдання 18-19} 

Кілька запитань за завданнями для самостійної роботи.
\vspace{10mm}
Скоріш за все, що завдань буде трохи менше, а деякі попередні будуть переформульовані в стилі двох останніх.

%enditem

\end {large}}

%end

\newpage
%begin

{{
\begin {large}

\newpage
1. Що буде надруковано за виконання?

\begin{verbatim}
print('R', end=' ')
g=('a','b','c','d','e',0,1)
print(g[5])
\end{verbatim}

%enditem


\newpage
2. Що буде надруковано за виконання?

\begin{verbatim}
print('R', end=' ')
b=['0','1','2','3','4','5','6','7','8','9']
print(b[:7:-2])
\end{verbatim}

%enditem



\newpage
3. Що буде надруковано за виконання?

\begin{verbatim}
print('R', end=' ')
e=[10,11,12,13,14]
e+=[1,0]
print(e)
\end{verbatim}

%enditem



\newpage
4. Що буде надруковано за виконання?

\begin{verbatim}
print('R', end=' ')
c=[10,11,12,13,14]
c[-7:-9]='13'
print(c)
\end{verbatim}

%enditem



\newpage
5. Укажіть ТИП значення виразу.\par
\begin{verbatim}
{y:13 for y in range(-7,7)}
\end{verbatim}
%enditem


\newpage
6. Що буде надруковано за виконання?\par
\begin{verbatim}
try:
    s = {32:1, 43:4, 26:2, 38:2, 35:0}
    s[32] = 0
    for x in s :
        print(x, sep=' ', end=' ')
except:
    print('ERR')
\end{verbatim}
%enditem


\newpage
7. Нехай змінна \kw{s} позначає дуже довгу послідовність, по якій можна ітерувати. 

Написати вираз-генератор, що задає підпослідовність її елементів, що є числами типу \kw{float}.

Обмеження: усі елементи шуканої підпослідовності одночасно в пам'яті розміщені бути не можуть. 
%enditem


\newpage
8. Що буде надруковано за виконання?
\begin{verbatim}
d = 0
m = 1

class A:
    d = 2
    
    def __init__(self):
        global d
        self.m = 3
        d = 4
        A.d = 5

obj = A()
try:
    print(d, end=' ')
    print(m, end=' ')
    print(A.d, end=' ')
    print(obj.m)
except:
    print('ERR')
\end{verbatim}
%enditem


\newpage
9. Що буде надруковано за виконання?
\begin{verbatim}
class D0:
    def d(self): print(1, end=' ')
    def g(self): print(2, end=' ')

class D1(D0):
    def d(self): print(3, end=' ')

try:
    D1().g()
    D1().d()
except:
    print('ERR')
\end{verbatim}
%enditem


\newpage
10. Що буде надруковано за виконання?
\begin{verbatim}
class A:
    b = 1

    def __init__(self):
        c = 2
        c = 3

class B(A):
    b = 4
    a = 5

    def __init__(self):
        super().__init__()
        self.a = 6

try:
    obj = B()
    obj.b = obj.a + 10
    B.b = 13
except: print('EE', end=' ')

try: print(obj.b, end= ' ')
except: print('E1', end=' ')

try: print(A.b, end= ' ')
except: print('E2', end=' ')

try: print(B.b, end= ' ')
except: print('E3', end=' ')
\end{verbatim}
%enditem


\newpage
11. Змінні \kw{next} та \kw{prev} вузла списку позначають наступний та попередній за ним вузол відповідно. Починаючи з голови, у список записано послідовні цілі числа від~$-44$ до~$29$. Нехай змінна \kw{p} позначає вузол, що зберігає значення~$1$.
Яке значення зберігається у вузлі \kw{p.next.prev.prev.prev.next} ?
%enditem


\newpage
12. 
\begin{center}
	\adjustimage{max size={0.8\linewidth}{0.8\paperheight}}
	{BinTree/trees_exp/149.png}
\end{center}
Указати послідовність міток вершин дерева, зображеного на рис., що утвориться в результаті обходу дерева в оберненому порядку. Мітки записувати через пробіл.\par Алгоритм починає роботу з кореня (на рис. це верхня вершина).
%enditem


\newpage
13. 
\begin{center}
	\adjustimage{max size={0.9\linewidth}{0.9\paperheight}}
	{Graphs/AdjList/83.png}
\end{center}
Для графа, зображеного на рис., вказати список суміжності вершини 6.\par
Формат оформлення відповіді:\par
список суміжності виписувати через пробіл на одному рядку, якщо пробілів десь буде більше одного - не важливо, упорядкування вершин у списку також значення не має.
%enditem


\newpage
14. Для графа на рисунку вказати послідовність, в якій алгоритм пошуку в глибину буде відмічати відвідані вершини, якщо списки суміжності вершин переглядаються алгоритмом за зростанням міток вершин та вершини відмічаються на першому вход\-женні у вершину.\par
Пошук розпочинається з вершини 6.\par

\begin{center}
	\adjustimage{max size={0.9\linewidth}{0.9\paperheight}}
	{Graphs/Traversals/135.png}
\end{center}
%enditem


\newpage
15. Для графа на рисунку з вершини 6 запущено алгоритм пошуку в ширину, що будує кістякове дерево. Списки суміжності вершин переглядаються алгоритмом за зростанням міток вершин.\par
\begin{center}
	\adjustimage{max size={0.9\linewidth}{0.9\paperheight}}
	{Graphs/Traversals/87.png}
\end{center}
\par Які з наведених ребер входять до складу отриманого кістякового дерева?\par
1) (0, 6)\par
2) (4, 5)\par
3) (0, 5)\par
4) (0, 3)\par
5) (1, 4)\par
6) (3, 4)\par
{\vskip 15pt} Номери відповідних ребер записати за зростанням через пробіл.


%enditem


\newpage
16. 

Рядки журналу через двокрапку містять ім'я користувача (ідентифікатор), час події,тип події, код події, наприклад
\begin{verbatim}
s="username : 2022-05-05 13:26 : ERROR : 404"
\end{verbatim}
у коді 
\begin{verbatim}
res = re.fullmatch('???', s) 
s = ''
code_str = ??? 
\end{verbatim}
чим треба замінити кожен з ???, щоб значенням code\_str був код події? 


\vspace{10mm}
\textbf{Завдання 17} 

Задано програму, що містить код 1-2 класів, можливе успадкування, та клієнтський код.
У наведеному коді невеликі частини закриті. 
Також наведено вихід програми.
Необхідно відновити закриті частини, щоб отриманий код відповідав наведеному виходу програми.


\textbf{Завдання 18-19} 

Кілька запитань за завданнями для самостійної роботи.
\vspace{10mm}
Скоріш за все, що завдань буде трохи менше, а деякі попередні будуть переформульовані в стилі двох останніх.

%enditem

\end {large}}

%end

\newpage
%begin

{{
\begin {large}

\newpage
1. Що буде надруковано за виконання?

\begin{verbatim}
print('R', end=' ')
c=('a','b','c','d','e',0,1,2)
print(c[-5])
\end{verbatim}

%enditem


\newpage
2. Що буде надруковано за виконання?

\begin{verbatim}
print('R', end=' ')
a='0123456789'
print(a[::-1])
\end{verbatim}

%enditem



\newpage
3. Що буде надруковано за виконання?

\begin{verbatim}
print('R', end=' ')
a=[10,11,12,13,14]
print(a.index('10'), end=' ')
\end{verbatim}

%enditem



\newpage
4. Що буде надруковано за виконання?

\begin{verbatim}
print('R', end=' ')
g=[10,11,12,13,14]
del g[-8:1]
print(g)
\end{verbatim}

%enditem



\newpage
5. Укажіть ТИП значення виразу.\par
\begin{verbatim}
{y:y*y for y in range(-7,7)}
\end{verbatim}
%enditem


\newpage
6. Що буде надруковано за виконання?\par
\begin{verbatim}
try:
    t = {43:0, 70:6, 48:2, 70:7}
    t[70] = 7
    for x in t.keys():
        print(x, sep=' ', end=' ')
except:
    print('ERR')
\end{verbatim}
%enditem


\newpage
7. Нехай змінна \kw{s} позначає дуже довгу послідовність цілих, по якій можна ітерувати. 

Написати вираз-генератор, що задає підпослідовність її елементів, що є від'єм\-ни\-ми. 

Обмеження: усі елементи шуканої підпослідовності одночасно в пам'яті розміщені бути не можуть. 
%enditem


\newpage
8. Що буде надруковано за виконання?
%Оригинал был утерян. Требует отладки!
\begin{verbatim}
d = 0

class B:
    d = 1
    
    def __init__(self):
        self.d = 2
        d = 3
        B.d = 4

obj = B()
B.d = 5
try:
    print(d, end=' ')
    print(obj.d, end=' ')
    print(B.d, end=' ')
    print(obj.d, end=' ')
    print(obj._B__d)
except:
    print('ERR')
\end{verbatim}
%enditem


\newpage
9. Що буде надруковано за виконання?
\begin{verbatim}
class B0:
    def a(self): print(1, end=' ')
    def f(self): print(2, end=' ')

class B1(B0):
    def a(self): print(3, end=' ')

try:
    B1().f()
    B0().a()
except:
    print('ERR')
\end{verbatim}
%enditem


\newpage
10. Що буде надруковано за виконання?
\begin{verbatim}
class A:
    c = 1

    def __init__(self):
        b = 2
        c = 3

class B(A):
    b = 4
    a = 5

    def __init__(self):
        super().__init__()
        self.a = 6

try:
    obj = B()
    obj.a = obj.b + 10
    B.a = 13
except: print('EE', end=' ')

try: print(obj.a, end= ' ')
except: print('E1', end=' ')

try: print(A.a, end= ' ')
except: print('E2', end=' ')

try: print(B.a, end= ' ')
except: print('E3', end=' ')
\end{verbatim}
%enditem


\newpage
11. Змінні \kw{next} та \kw{prev} вузла списку позначають наступний та попередній за ним вузол відповідно. Починаючи з голови, у список записано послідовні цілі числа від~$-66$ до~$15$. Нехай змінна \kw{p} позначає вузол, що зберігає значення~$3$.
Яке значення зберігається у вузлі \kw{p.prev.next.next.next.next} ?
%enditem


\newpage
12. 
\begin{center}
	\adjustimage{max size={0.8\linewidth}{0.8\paperheight}}
	{BinTree/trees_exp/34.png}
\end{center}
Указати послідовність міток вершин дерева, зображеного на рис., що утвориться в результаті обходу дерева в оберненому порядку. Мітки записувати через пробіл.\par Алгоритм починає роботу з кореня (на рис. це верхня вершина).
%enditem


\newpage
13. 
\begin{center}
	\adjustimage{max size={0.9\linewidth}{0.9\paperheight}}
	{Graphs/AdjList/319.png}
\end{center}
Для графа, зображеного на рис., вказати список суміжності вершини 1.\par
Формат оформлення відповіді:\par
список суміжності виписувати через пробіл на одному рядку, якщо пробілів десь буде більше одного - не важливо, упорядкування вершин у списку також значення не має.
%enditem


\newpage
14. Для графа на рисунку вказати послідовність, в якій алгоритм пошуку в глибину буде відмічати відвідані вершини, якщо списки суміжності вершин переглядаються алгоритмом за зростанням міток вершин та вершини відмічаються на першому вход\-женні у вершину.\par
Пошук розпочинається з вершини 0.\par

\begin{center}
	\adjustimage{max size={0.9\linewidth}{0.9\paperheight}}
	{Graphs/Traversals/122.png}
\end{center}
%enditem


\newpage
15. Для графа на рисунку з вершини 2 запущено алгоритм пошуку в ширину, що будує кістякове дерево. Списки суміжності вершин переглядаються алгоритмом за зростанням міток вершин.\par
\begin{center}
	\adjustimage{max size={0.9\linewidth}{0.9\paperheight}}
	{Graphs/Traversals/177.png}
\end{center}
\par Які з наведених ребер входять до складу отриманого кістякового дерева?\par
1) (2, 6)\par
2) (0, 3)\par
3) (5, 6)\par
4) (3, 5)\par
5) (0, 6)\par
6) (0, 2)\par
{\vskip 15pt} Номери відповідних ребер записати за зростанням через пробіл.


%enditem


\newpage
16. 

Рядки журналу через двокрапку містять ім'я користувача (ідентифікатор), час події,тип події, код події, наприклад
\begin{verbatim}
s="username : 2022-05-05 13:26 : ERROR : 404"
\end{verbatim}
у коді 
\begin{verbatim}
res = re.fullmatch('???', s) 
s = ''
code_str = ??? 
\end{verbatim}
чим треба замінити кожен з ???, щоб значенням code\_str був код події? 


\vspace{10mm}
\textbf{Завдання 17} 

Задано програму, що містить код 1-2 класів, можливе успадкування, та клієнтський код.
У наведеному коді невеликі частини закриті. 
Також наведено вихід програми.
Необхідно відновити закриті частини, щоб отриманий код відповідав наведеному виходу програми.


\textbf{Завдання 18-19} 

Кілька запитань за завданнями для самостійної роботи.
\vspace{10mm}
Скоріш за все, що завдань буде трохи менше, а деякі попередні будуть переформульовані в стилі двох останніх.

%enditem

\end {large}}

%end

\newpage
%begin

{{
\begin {large}

\newpage
1. Що буде надруковано за виконання?

\begin{verbatim}
print('R', end=' ')
h=('a','b','c','d','e',0,1,2,3,4)
print(h[-14])
\end{verbatim}

%enditem


\newpage
2. Що буде надруковано за виконання?

\begin{verbatim}
print('R', end=' ')
h=(0,1,2,3,4,5,6,7,8,9)
print(h[9:-7:-3])
\end{verbatim}

%enditem



\newpage
3. Що буде надруковано за виконання?

\begin{verbatim}
print('R', end=' ')
d=['0','1','2','3','4']
d.remove('2')
print(d)
\end{verbatim}

%enditem



\newpage
4. Що буде надруковано за виконання?

\begin{verbatim}
print('R', end=' ')
f=['0','1','2','3','4']
f[:-6]=[15]
print(f)
\end{verbatim}

%enditem



\newpage
5. Укажіть ТИП значення виразу.\par
\begin{verbatim}
[7,2,3,4]
\end{verbatim}
%enditem


\newpage
6. Що буде надруковано за виконання?\par
\begin{verbatim}
try:
    s = {42:9, 68:2, 76:5, 36:1, 76:0}
    s[42] = 6
    for x in s.items():
        print(x, sep=' ', end=' ')
except:
    print('ERR')
\end{verbatim}
%enditem


\newpage
7. Нехай змінна \kw{s} позначає дуже довгу послідовність цілих, по якій можна ітерувати. 

Написати вираз-генератор, що задає підпослідовність її елементів, що менші за задане число num. 

Обмеження: усі елементи шуканої підпослідовності одночасно в пам'яті розміщені бути не можуть. 
%enditem


\newpage
8. Що буде надруковано за виконання?
\begin{verbatim}
a = 0
j = 1

class B:
    j = 2
    
    def __init__(self):
        self.a = 3
        j = 4
        B.j = 5

obj = B()
try:
    print(a, end=' ')
    print(j, end=' ')
    print(B.j, end=' ')
    print(obj.a)
except:
    print('ERR')
\end{verbatim}
%enditem


\newpage
9. Що буде надруковано за виконання?
\begin{verbatim}
class C0:
    def b(self): print(1, end=' ')
    def g(self): print(2, end=' ')

class C1(C0):
    def g(self): print(3, end=' ')

try:
    C0().b()
    C1().g()
except:
    print('ERR')
\end{verbatim}
%enditem


\newpage
10. Що буде надруковано за виконання?
\begin{verbatim}
class A:
    a = 1

    def __init__(self):
        c = 2
        c = 3

class B(A):
    c = 4
    b = 5

    def __init__(self):
        super().__init__()
        self.b = 6

try:
    obj = B()
    obj.c = obj.c + 10
    B.c = 13
except: print('EE', end=' ')

try: print(obj.c, end= ' ')
except: print('E1', end=' ')

try: print(A.c, end= ' ')
except: print('E2', end=' ')

try: print(B.c, end= ' ')
except: print('E3', end=' ')
\end{verbatim}
%enditem


\newpage
11. Змінні \kw{next} та \kw{prev} вузла списку позначають наступний та попередній за ним вузол відповідно. Починаючи з голови, у список записано послідовні цілі числа від~$-33$ до~$41$. Нехай змінна \kw{p} позначає вузол, що зберігає значення~$-3$.
Яке значення зберігається у вузлі \kw{p.next.prev.prev.prev.prev} ?
%enditem


\newpage
12. 
\begin{center}
	\adjustimage{max size={0.8\linewidth}{0.8\paperheight}}
	{BinTree/trees_exp/146.png}
\end{center}
Указати послідовність міток вершин дерева, зображеного на рис., що утвориться в результаті обходу дерева в прямому порядку. Мітки записувати через пробіл.\par Алгоритм починає роботу з кореня (на рис. це верхня вершина).
%enditem


\newpage
13. 
\begin{center}
	\adjustimage{max size={0.9\linewidth}{0.9\paperheight}}
	{Graphs/AdjList/8.png}
\end{center}
Для графа, зображеного на рис., вказати список суміжності вершини 7.\par
Формат оформлення відповіді:\par
список суміжності виписувати через пробіл на одному рядку, якщо пробілів десь буде більше одного - не важливо, упорядкування вершин у списку також значення не має.
%enditem


\newpage
14. Для графа на рисунку вказати послідовність, в якій алгоритм пошуку в глибину буде відмічати відвідані вершини, якщо списки суміжності вершин переглядаються алгоритмом за зростанням міток вершин та вершини відмічаються на першому вход\-женні у вершину.\par
Пошук розпочинається з вершини 5.\par

\begin{center}
	\adjustimage{max size={0.9\linewidth}{0.9\paperheight}}
	{Graphs/Traversals/106.png}
\end{center}
%enditem


\newpage
15. Для графа на рисунку з вершини 5 запущено алгоритм пошуку в ширину, що будує кістякове дерево. Списки суміжності вершин переглядаються алгоритмом за зростанням міток вершин.\par
\begin{center}
	\adjustimage{max size={0.9\linewidth}{0.9\paperheight}}
	{Graphs/Traversals/201.png}
\end{center}
\par Які з наведених ребер входять до складу отриманого кістякового дерева?\par
1) (4, 5)\par
2) (0, 3)\par
3) (2, 5)\par
4) (2, 3)\par
5) (0, 1)\par
6) (2, 4)\par
{\vskip 15pt} Номери відповідних ребер записати за зростанням через пробіл.


%enditem


\newpage
16. 

Рядки журналу через двокрапку містять ім'я користувача (ідентифікатор), час події,тип події, код події, наприклад
\begin{verbatim}
s="username : 2022-05-05 13:26 : ERROR : 404"
\end{verbatim}
у коді 
\begin{verbatim}
res = re.fullmatch('???', s) 
s = ''
code_str = ??? 
\end{verbatim}
чим треба замінити кожен з ???, щоб значенням code\_str був код події? 


\vspace{10mm}
\textbf{Завдання 17} 

Задано програму, що містить код 1-2 класів, можливе успадкування, та клієнтський код.
У наведеному коді невеликі частини закриті. 
Також наведено вихід програми.
Необхідно відновити закриті частини, щоб отриманий код відповідав наведеному виходу програми.


\textbf{Завдання 18-19} 

Кілька запитань за завданнями для самостійної роботи.
\vspace{10mm}
Скоріш за все, що завдань буде трохи менше, а деякі попередні будуть переформульовані в стилі двох останніх.

%enditem

\end {large}}

%end

\newpage
%begin

{{
\begin {large}

\newpage
1. Що буде надруковано за виконання?

\begin{verbatim}
print('R', end=' ')
e=('a','b','c','d','e',0,1)
print(e[-1])
\end{verbatim}

%enditem


\newpage
2. Що буде надруковано за виконання?

\begin{verbatim}
print('R', end=' ')
d='0123456789'
print(d[6:2:-1])
\end{verbatim}

%enditem



\newpage
3. Що буде надруковано за виконання?

\begin{verbatim}
print('R', end=' ')
b=[10,11,12,13,14]
b.insert(-1, '13')
print(b)
\end{verbatim}

%enditem



\newpage
4. Що буде надруковано за виконання?

\begin{verbatim}
print('R', end=' ')
b=['0','1','2','3','4']
del b[-3:-7]
print(b)
\end{verbatim}

%enditem



\newpage
5. Укажіть ТИП значення виразу.\par
\begin{verbatim}
[{1,2}, 5]
\end{verbatim}
%enditem


\newpage
6. Що буде надруковано за виконання?\par
\begin{verbatim}
try:
    t = {64:0, 50:3, 62:4, 14:8, 76:9}
    t[62] = 3
    for x in t.items():
        print(*x, sep=' ', end=' ')
except:
    print('ERR')
\end{verbatim}
%enditem


\newpage
7. Нехай змінна \kw{s} позначає дуже довгу послідовність, по якій можна ітерувати. 

Написати вираз-генератор, що задає підпослідовність її елементів, що не є рядками типу \kw{str}.

Обмеження: усі елементи шуканої підпослідовності одночасно в пам'яті розміщені бути не можуть. 
%enditem


\newpage
8. Що буде надруковано за виконання?
\begin{verbatim}
c = 0
h = 1

class C:
    h = 2
    
    def __init__(self):
        self.h = 3
        h = 4
        C.h = 5

obj = C()
try:
    print(c, end=' ')
    print(h, end=' ')
    print(C.h, end=' ')
    print(obj.h)
except:
    print('ERR')
\end{verbatim}
%enditem


\newpage
9. Що буде надруковано за виконання?
\begin{verbatim}
class A0:
    def a(self): print(1, end=' ')
    def h(self): print(2, end=' ')

class A1(A0):
    def a(self): print(3, end=' ')

try:
    A1().h()
    A0().a()
except:
    print('ERR')
\end{verbatim}
%enditem


\newpage
10. Що буде надруковано за виконання?
\begin{verbatim}
class A:
    b = 1

    def g1(self): print(2, end=' ')
    def _g2(self): print(3, end=' ')

    def main(self):
        try: print(self.b, end=' ')
        except: print('E1', end=' ')

        try: self.g1()
        except: print('E2', end=' ')

        try: self._g2()
        except: print('E3', end=' ')


class B(A):
    b = 4

    def g1(self): print(5, end=' ')
    def _g2(self): print(6, end=' ')

try: B().main()
except: print('E4', end=' ')
\end{verbatim}
%enditem


\newpage
11. Змінні \kw{next} та \kw{prev} вузла списку позначають наступний та попередній за ним вузол відповідно. Починаючи з голови, у список записано послідовні цілі числа від~$-43$ до~$36$. Нехай змінна \kw{p} позначає вузол, що зберігає значення~$0$.
Яке значення зберігається у вузлі \kw{p.next.next.next.prev.prev} ?
%enditem


\newpage
12. 
\begin{center}
	\adjustimage{max size={0.8\linewidth}{0.8\paperheight}}
	{BinTree/trees_exp/88.png}
\end{center}
Указати послідовність міток вершин дерева, зображеного на рис., що утвориться в результаті обходу дерева в прямому порядку. Мітки записувати через пробіл.\par Алгоритм починає роботу з кореня (на рис. це верхня вершина).
%enditem


\newpage
13. 
\begin{center}
	\adjustimage{max size={0.9\linewidth}{0.9\paperheight}}
	{Graphs/AdjList/312.png}
\end{center}
Для графа, зображеного на рис., вказати список суміжності вершини 0.\par
Формат оформлення відповіді:\par
список суміжності виписувати через пробіл на одному рядку, якщо пробілів десь буде більше одного - не важливо, упорядкування вершин у списку також значення не має.
%enditem


\newpage
14. Для графа на рисунку вказати послідовність, в якій алгоритм пошуку в глибину буде відмічати відвідані вершини, якщо списки суміжності вершин переглядаються алгоритмом за зростанням міток вершин та вершини відмічаються на першому вход\-женні у вершину.\par
Пошук розпочинається з вершини 3.\par

\begin{center}
	\adjustimage{max size={0.9\linewidth}{0.9\paperheight}}
	{Graphs/Traversals/125.png}
\end{center}
%enditem


\newpage
15. Для графа на рисунку з вершини 5 запущено алгоритм пошуку в ширину, що будує кістякове дерево. Списки суміжності вершин переглядаються алгоритмом за зростанням міток вершин.\par
\begin{center}
	\adjustimage{max size={0.9\linewidth}{0.9\paperheight}}
	{Graphs/Traversals/129.png}
\end{center}
\par Які з наведених ребер входять до складу отриманого кістякового дерева?\par
1) (0, 1)\par
2) (1, 5)\par
3) (4, 5)\par
4) (2, 5)\par
5) (4, 6)\par
6) (3, 4)\par
{\vskip 15pt} Номери відповідних ребер записати за зростанням через пробіл.


%enditem


\newpage
16. 

Рядки журналу через двокрапку містять ім'я користувача (ідентифікатор), час події,тип події, код події, наприклад
\begin{verbatim}
s="username : 2022-05-05 13:26 : ERROR : 404"
\end{verbatim}
у коді 
\begin{verbatim}
res = re.fullmatch('???', s) 
s = ''
code_str = ??? 
\end{verbatim}
чим треба замінити кожен з ???, щоб значенням code\_str був код події? 


\vspace{10mm}
\textbf{Завдання 17} 

Задано програму, що містить код 1-2 класів, можливе успадкування, та клієнтський код.
У наведеному коді невеликі частини закриті. 
Також наведено вихід програми.
Необхідно відновити закриті частини, щоб отриманий код відповідав наведеному виходу програми.


\textbf{Завдання 18-19} 

Кілька запитань за завданнями для самостійної роботи.
\vspace{10mm}
Скоріш за все, що завдань буде трохи менше, а деякі попередні будуть переформульовані в стилі двох останніх.

%enditem

\end {large}}

%end

\newpage
%begin

{{
\begin {large}

\newpage
1. Що буде надруковано за виконання?

\begin{verbatim}
print('R', end=' ')
a=('a','b','c','d','e',0,1,2,3,4)
print(a[2])
\end{verbatim}

%enditem


\newpage
2. Що буде надруковано за виконання?

\begin{verbatim}
print('R', end=' ')
e=['0','1','2','3','4','5','6','7','8','9']
print(e[::2])
\end{verbatim}

%enditem



\newpage
3. Що буде надруковано за виконання?

\begin{verbatim}
print('R', end=' ')
g=['0','1','2','3','4']
g.extend([1,2])
print(g)
\end{verbatim}

%enditem



\newpage
4. Що буде надруковано за виконання?

\begin{verbatim}
print('R', end=' ')
e=[10,11,12,13,14]
e[7:-5]=()
print(e)
\end{verbatim}

%enditem



\newpage
5. Укажіть ТИП значення виразу.\par
\begin{verbatim}
{x for x in range(7) if el%3==1}
\end{verbatim}
%enditem


\newpage
6. Що буде надруковано за виконання?\par
\begin{verbatim}
try:
    d = {26:6, 64:2, 82:2, 45:1, 64:6}
    d[82] = 9
    for x in d.values():
        print(x, sep=' ', end=' ')
except:
    print('ERR')
\end{verbatim}
%enditem


\newpage
7. Нехай змінна \kw{s} позначає дуже довгу послідовність цілих, по якій можна ітерувати. 

Написати вираз-генератор, що задає підпослідовність її елементів, що не діляться на жодне з чисел 2, 7, 5. 

Обмеження: усі елементи шуканої підпослідовності одночасно в пам'яті розміщені бути не можуть. 
%enditem


\newpage
8. Що буде надруковано за виконання?
%Оригинал был утерян. Требует отладки!
\begin{verbatim}
c = 0

class A:
    c = 1
    
    def __init__(self):
        self.c = 2
        c = 3
        A.c = 4

obj = A()
A.c = 5
try:
    print(c, end=' ')
    print(obj.c, end=' ')
    print(A.c, end=' ')
    print(obj.c, end=' ')
    print(obj._A__c)
except:
    print('ERR')
\end{verbatim}
%enditem


\newpage
9. Що буде надруковано за виконання?
\begin{verbatim}
class B0:
    def d(self): print(1, end=' ')
    def f(self): print(2, end=' ')

class B1(B0):
    def f(self): print(3, end=' ')

try:
    B1().d()
    B1().f()
except:
    print('ERR')
\end{verbatim}
%enditem


\newpage
10. Що буде надруковано за виконання?
\begin{verbatim}
class A:
    a = 1

    def __init__(self):
        b = 2
        a = 3

class B(A):
    a = 4
    c = 5

    def __init__(self):
        super().__init__()
        self.c = 6

try:
    obj = B()
    obj.b = obj.a + 10
    B.b = 13
except: print('EE', end=' ')

try: print(obj.b, end= ' ')
except: print('E1', end=' ')

try: print(A.b, end= ' ')
except: print('E2', end=' ')

try: print(B.b, end= ' ')
except: print('E3', end=' ')
\end{verbatim}
%enditem


\newpage
11. Змінні \kw{next} та \kw{prev} вузла списку позначають наступний та попередній за ним вузол відповідно. Починаючи з голови, у список записано послідовні цілі числа від~$-55$ до~$46$. Нехай змінна \kw{p} позначає вузол, що зберігає значення~$8$.
Яке значення зберігається у вузлі \kw{p.prev.prev.prev.next.next} ?
%enditem


\newpage
12. 
\begin{center}
	\adjustimage{max size={0.8\linewidth}{0.8\paperheight}}
	{BinTree/trees_exp/71.png}
\end{center}
Указати послідовність міток вершин дерева, зображеного на рис., що утвориться в результаті обходу дерева в оберненому порядку. Мітки записувати через пробіл.\par Алгоритм починає роботу з кореня (на рис. це верхня вершина).
%enditem


\newpage
13. 
\begin{center}
	\adjustimage{max size={0.9\linewidth}{0.9\paperheight}}
	{Graphs/AdjList/126.png}
\end{center}
Для графа, зображеного на рис., вказати список суміжності вершини 0.\par
Формат оформлення відповіді:\par
список суміжності виписувати через пробіл на одному рядку, якщо пробілів десь буде більше одного - не важливо, упорядкування вершин у списку також значення не має.
%enditem


\newpage
14. Для графа на рисунку вказати послідовність, в якій алгоритм пошуку в глибину буде відмічати відвідані вершини, якщо списки суміжності вершин переглядаються алгоритмом за зростанням міток вершин та вершини відмічаються на першому вход\-женні у вершину.\par
Пошук розпочинається з вершини 3.\par

\begin{center}
	\adjustimage{max size={0.9\linewidth}{0.9\paperheight}}
	{Graphs/Traversals/180.png}
\end{center}
%enditem


\newpage
15. Для графа на рисунку з вершини 5 запущено алгоритм пошуку в ширину, що будує кістякове дерево. Списки суміжності вершин переглядаються алгоритмом за зростанням міток вершин.\par
\begin{center}
	\adjustimage{max size={0.9\linewidth}{0.9\paperheight}}
	{Graphs/Traversals/40.png}
\end{center}
\par Які з наведених ребер входять до складу отриманого кістякового дерева?\par
1) (2, 3)\par
2) (4, 6)\par
3) (4, 5)\par
4) (0, 2)\par
5) (1, 2)\par
6) (0, 6)\par
{\vskip 15pt} Номери відповідних ребер записати за зростанням через пробіл.


%enditem


\newpage
16. 

Рядки журналу через двокрапку містять ім'я користувача (ідентифікатор), час події,тип події, код події, наприклад
\begin{verbatim}
s="username : 2022-05-05 13:26 : ERROR : 404"
\end{verbatim}
у коді 
\begin{verbatim}
res = re.fullmatch('???', s) 
s = ''
code_str = ??? 
\end{verbatim}
чим треба замінити кожен з ???, щоб значенням code\_str був код події? 


\vspace{10mm}
\textbf{Завдання 17} 

Задано програму, що містить код 1-2 класів, можливе успадкування, та клієнтський код.
У наведеному коді невеликі частини закриті. 
Також наведено вихід програми.
Необхідно відновити закриті частини, щоб отриманий код відповідав наведеному виходу програми.


\textbf{Завдання 18-19} 

Кілька запитань за завданнями для самостійної роботи.
\vspace{10mm}
Скоріш за все, що завдань буде трохи менше, а деякі попередні будуть переформульовані в стилі двох останніх.

%enditem

\end {large}}

%end

\newpage
%begin

{{
\begin {large}

\newpage
1. Що буде надруковано за виконання?

\begin{verbatim}
print('R', end=' ')
f=('a','b','c','d','e',0,1,2)
print(f[10])
\end{verbatim}

%enditem


\newpage
2. Що буде надруковано за виконання?

\begin{verbatim}
print('R', end=' ')
g=('0','1','2','3','4','5','6','7','8','9')
print(g[2:9:-3])
\end{verbatim}

%enditem



\newpage
3. Що буде надруковано за виконання?

\begin{verbatim}
print('R', end=' ')
h=[10,11,12,13,14]
print(h.pop(-1), end=' ')
print(h)
\end{verbatim}

%enditem



\newpage
4. Що буде надруковано за виконання?

\begin{verbatim}
print('R', end=' ')
a=[10,11,12,13,14]
a[8:-2]='11'
print(a)
\end{verbatim}

%enditem



\newpage
5. Укажіть ТИП значення виразу.\par
\begin{verbatim}
{(1,4), 8}
\end{verbatim}
%enditem


\newpage
6. Що буде надруковано за виконання?\par
\begin{verbatim}
try:
    s = {26:4, 54:3, 27:2, 27:3}
    s[51] = 6
    for x in s.values():
        print(x, sep=' ', end=' ')
except:
    print('ERR')
\end{verbatim}
%enditem


\newpage
7. 
Записати ВИРАЗ, значенням якого є список, що складається з квадратів цілих чисел від 17 до 2014 (включно), що не діляться на жодне з чисел 2, 3, записаних в обернено\-му порядку.
%enditem


\newpage
8. Що буде надруковано за виконання?
%Оригинал был утерян. Требует отладки!
\begin{verbatim}
__b = 0

class D:
    __b = 1
    
    def __init__(self):
        self.__b = 2
        __b = 3
        D.__b = 4

obj = D()
obj.__b = 5
try:
    print(__b, end=' ')
    print(obj.__b, end=' ')
    print(D.__b, end=' ')
    print(obj.__b, end=' ')
    print(obj._D__b)
except:
    print('ERR')
\end{verbatim}
%enditem


\newpage
9. Що буде надруковано за виконання?
\begin{verbatim}
class C0:
    def b(self): print(1, end=' ')
    def h(self): print(2, end=' ')

class C1(C0):
    def b(self): print(3, end=' ')

try:
    C1().h()
    C1().b()
except:
    print('ERR')
\end{verbatim}
%enditem


\newpage
10. Що буде надруковано за виконання?
\begin{verbatim}
class A:
    a = 1

    def __init__(self):
        b = 2
        b = 3

class B(A):
    a = 4
    c = 5

    def __init__(self):
        super().__init__()
        self.c = 6

try:
    obj = B()
    obj.a = obj.b + 10
    B.a = 13
except: print('EE', end=' ')

try: print(obj.a, end= ' ')
except: print('E1', end=' ')

try: print(A.a, end= ' ')
except: print('E2', end=' ')

try: print(B.a, end= ' ')
except: print('E3', end=' ')
\end{verbatim}
%enditem


\newpage
11. Змінні \kw{next} та \kw{prev} вузла списку позначають наступний та попередній за ним вузол відповідно. Починаючи з голови, у список записано послідовні цілі числа від~$-61$ до~$23$. Нехай змінна \kw{p} позначає вузол, що зберігає значення~$6$.
Яке значення зберігається у вузлі \kw{p.next.prev.next.next.prev} ?
%enditem


\newpage
12. 
\begin{center}
	\adjustimage{max size={0.8\linewidth}{0.8\paperheight}}
	{BinTree/trees_exp/147.png}
\end{center}
Указати послідовність міток вершин дерева, зображеного на рис., що утвориться в результаті обходу дерева в оберненому порядку. Мітки записувати через пробіл.\par Алгоритм починає роботу з кореня (на рис. це верхня вершина).
%enditem


\newpage
13. 
\begin{center}
	\adjustimage{max size={0.9\linewidth}{0.9\paperheight}}
	{Graphs/AdjList/250.png}
\end{center}
Для графа, зображеного на рис., вказати список суміжності вершини 1.\par
Формат оформлення відповіді:\par
список суміжності виписувати через пробіл на одному рядку, якщо пробілів десь буде більше одного - не важливо, упорядкування вершин у списку також значення не має.
%enditem


\newpage
14. Для графа на рисунку з вершини 6 запущено алгоритм пошуку в глибину, що будує кістякове дерево. Списки суміжності вершин переглядаються алгоритмом за зростанням міток вершин.\par
\begin{center}
	\adjustimage{max size={0.9\linewidth}{0.9\paperheight}}
	{Graphs/Traversals/240.png}
\end{center}
\par Які з наведених ребер входять до складу отриманого кістякового дерева?\par
1) (3, 5)\par
2) (0, 6)\par
3) (1, 5)\par
4) (1, 2)\par
5) (0, 2)\par
6) (0, 4)\par
{\vskip 15pt} Номери відповідних ребер записати за зростанням через пробіл.

%enditem


\newpage
15. Для графа на рисунку з вершини 1 запущено алгоритм пошуку в ширину, що будує кістякове дерево. Списки суміжності вершин переглядаються алгоритмом за зростанням міток вершин.\par
\begin{center}
	\adjustimage{max size={0.9\linewidth}{0.9\paperheight}}
	{Graphs/Traversals/204.png}
\end{center}
\par Які з наведених ребер входять до складу отриманого кістякового дерева?\par
1) (3, 6)\par
2) (1, 4)\par
3) (2, 5)\par
4) (4, 5)\par
5) (5, 6)\par
6) (2, 6)\par
{\vskip 15pt} Номери відповідних ребер записати за зростанням через пробіл.


%enditem


\newpage
16. 

Рядки журналу через двокрапку містять ім'я користувача (ідентифікатор), час події,тип події, код події, наприклад
\begin{verbatim}
s="username : 2022-05-05 13:26 : ERROR : 404"
\end{verbatim}
у коді 
\begin{verbatim}
res = re.fullmatch('???', s) 
s = ''
code_str = ??? 
\end{verbatim}
чим треба замінити кожен з ???, щоб значенням code\_str був код події? 


\vspace{10mm}
\textbf{Завдання 17} 

Задано програму, що містить код 1-2 класів, можливе успадкування, та клієнтський код.
У наведеному коді невеликі частини закриті. 
Також наведено вихід програми.
Необхідно відновити закриті частини, щоб отриманий код відповідав наведеному виходу програми.


\textbf{Завдання 18-19} 

Кілька запитань за завданнями для самостійної роботи.
\vspace{10mm}
Скоріш за все, що завдань буде трохи менше, а деякі попередні будуть переформульовані в стилі двох останніх.

%enditem

\end {large}}

%end

\newpage
%begin

{{
\begin {large}

\newpage
1. Що буде надруковано за виконання?

\begin{verbatim}
print('R', end=' ')
b=('a','b','c','d','e',0,1,2,3)
print(b[-12])
\end{verbatim}

%enditem


\newpage
2. Що буде надруковано за виконання?

\begin{verbatim}
print('R', end=' ')
c=[0,1,2,3,4,5,6,7,8,9]
print(c[:-5:-2])
\end{verbatim}

%enditem



\newpage
3. Що буде надруковано за виконання?

\begin{verbatim}
print('R', end=' ')
f=['0','1','2','3','4']
f+=[2,3]
print(f)
\end{verbatim}

%enditem



\newpage
4. Що буде надруковано за виконання?

\begin{verbatim}
print('R', end=' ')
d=['0','1','2','3','4']
d[-9:5]=[4,5]
print(d)
\end{verbatim}

%enditem



\newpage
5. Укажіть ТИП значення виразу.\par
\begin{verbatim}
[{1,2}, 5]
\end{verbatim}
%enditem


\newpage
6. Що буде надруковано за виконання?\par
\begin{verbatim}
try:
    t = {26:1, 30:9, 24:1, 41:6, 24:6}
    t[70] = 6
    for x in t.items():
        print(*x, sep=' ', end=' ')
except:
    print('ERR')
\end{verbatim}
%enditem


\newpage
7. Нехай змінна \kw{s} позначає послідовність цілих, по якій можна ітерувати. 

Написати вираз, що перевіряє, чи правда, що послідовність  містить числа, більші за 5.
%enditem


\newpage
8. Що буде надруковано за виконання?
%Оригинал был утерян. Требует отладки!
\begin{verbatim}
_a = 0

class C:
    _a = 1
    
    def __init__(self):
        self._a = 2
        _a = 3
        C._a = 4

obj = C()
obj._a = 5
try:
    print(_a, end=' ')
    print(obj._a, end=' ')
    print(C._a, end=' ')
    print(obj._a, end=' ')
    print(obj._C__a)
except:
    print('ERR')
\end{verbatim}
%enditem


\newpage
9. Що буде надруковано за виконання?
\begin{verbatim}
class D0:
    def c(self): print(1, end=' ')
    def f(self): print(2, end=' ')

class D1(D0):
    def f(self): print(3, end=' ')

try:
    D0().c()
    D1().f()
except:
    print('ERR')
\end{verbatim}
%enditem


\newpage
10. Що буде надруковано за виконання?
\begin{verbatim}
class A:
    b = 1

    def __init__(self):
        c = 2
        b = 3

class B(A):
    c = 4
    a = 5

    def __init__(self):
        super().__init__()
        self.a = 6

try:
    obj = B()
    obj.c = obj.a + 10
    B.c = 13
except: print('EE', end=' ')

try: print(obj.c, end= ' ')
except: print('E1', end=' ')

try: print(A.c, end= ' ')
except: print('E2', end=' ')

try: print(B.c, end= ' ')
except: print('E3', end=' ')
\end{verbatim}
%enditem


\newpage
11. Змінні \kw{next} та \kw{prev} вузла списку позначають наступний та попередній за ним вузол відповідно. Починаючи з голови, у список записано послідовні цілі числа від~$-15$ до~$19$. Нехай змінна \kw{p} позначає вузол, що зберігає значення~$-1$.
Яке значення зберігається у вузлі \kw{p.prev.next.prev.prev.next} ?
%enditem


\newpage
12. 
\begin{center}
	\adjustimage{max size={0.8\linewidth}{0.8\paperheight}}
	{BinTree/trees_exp/26.png}
\end{center}
Указати послідовність міток вершин дерева, зображеного на рис., що утвориться в результаті обходу дерева в прямому порядку. Мітки записувати через пробіл.\par Алгоритм починає роботу з кореня (на рис. це верхня вершина).
%enditem


\newpage
13. 
\begin{center}
	\adjustimage{max size={0.9\linewidth}{0.9\paperheight}}
	{Graphs/AdjList/174.png}
\end{center}
Для графа, зображеного на рис., вказати список суміжності вершини 5.\par
Формат оформлення відповіді:\par
список суміжності виписувати через пробіл на одному рядку, якщо пробілів десь буде більше одного - не важливо, упорядкування вершин у списку також значення не має.
%enditem


\newpage
14. Для графа на рисунку вказати послідовність, в якій алгоритм пошуку в глибину буде відмічати відвідані вершини, якщо списки суміжності вершин переглядаються алгоритмом за зростанням міток вершин та вершини відмічаються на першому вход\-женні у вершину.\par
Пошук розпочинається з вершини 3.\par

\begin{center}
	\adjustimage{max size={0.9\linewidth}{0.9\paperheight}}
	{Graphs/Traversals/333.png}
\end{center}
%enditem


\newpage
15. Для графа на рисунку з вершини 4 запущено алгоритм пошуку в ширину, що будує кістякове дерево. Списки суміжності вершин переглядаються алгоритмом за зростанням міток вершин.\par
\begin{center}
	\adjustimage{max size={0.9\linewidth}{0.9\paperheight}}
	{Graphs/Traversals/174.png}
\end{center}
\par Які з наведених ребер входять до складу отриманого кістякового дерева?\par
1) (1, 6)\par
2) (0, 5)\par
3) (4, 5)\par
4) (3, 5)\par
5) (2, 3)\par
6) (1, 2)\par
{\vskip 15pt} Номери відповідних ребер записати за зростанням через пробіл.


%enditem


\newpage
16. 

Рядки журналу через двокрапку містять ім'я користувача (ідентифікатор), час події,тип події, код події, наприклад
\begin{verbatim}
s="username : 2022-05-05 13:26 : ERROR : 404"
\end{verbatim}
у коді 
\begin{verbatim}
res = re.fullmatch('???', s) 
s = ''
code_str = ??? 
\end{verbatim}
чим треба замінити кожен з ???, щоб значенням code\_str був код події? 


\vspace{10mm}
\textbf{Завдання 17} 

Задано програму, що містить код 1-2 класів, можливе успадкування, та клієнтський код.
У наведеному коді невеликі частини закриті. 
Також наведено вихід програми.
Необхідно відновити закриті частини, щоб отриманий код відповідав наведеному виходу програми.


\textbf{Завдання 18-19} 

Кілька запитань за завданнями для самостійної роботи.
\vspace{10mm}
Скоріш за все, що завдань буде трохи менше, а деякі попередні будуть переформульовані в стилі двох останніх.

%enditem

\end {large}}

%end

\newpage
%begin

{{
\begin {large}

\newpage
1. Що буде надруковано за виконання?

\begin{verbatim}
print('R', end=' ')
h=('a','b','c','d','e',0,1,2,4)
print(h[6])
\end{verbatim}

%enditem


\newpage
2. Що буде надруковано за виконання?

\begin{verbatim}
print('R', end=' ')
b=(0,1,2,3,4,5,6,7,8,9)
print(b[:-10:3])
\end{verbatim}

%enditem



\newpage
3. Що буде надруковано за виконання?

\begin{verbatim}
print('R', end=' ')
a=['0','1','2','3','4']
a.append('11')
print(a)
\end{verbatim}

%enditem



\newpage
4. Що буде надруковано за виконання?

\begin{verbatim}
print('R', end=' ')
e=[10,11,12,13,14]
del e[-1:]
print(e)
\end{verbatim}

%enditem



\newpage
5. Укажіть ТИП значення виразу.\par
\begin{verbatim}
(2, 7, -4)
\end{verbatim}
%enditem


\newpage
6. Що буде надруковано за виконання?\par
\begin{verbatim}
try:
    d = {32:4, 11:4, 14:8, 21:1}
    d[86] = 3
    for x in d.keys():
        print(x, sep=' ', end=' ')
except:
    print('ERR')
\end{verbatim}
%enditem


\newpage
7. Нехай змінна \kw{s} позначає дуже довгу послідовність цілих, по якій можна ітерувати. 

Написати вираз-генератор, що задає підпослідовність її елементів, що не діляться на жодне з чисел 2, 7, 5. 

Обмеження: усі елементи шуканої підпослідовності одночасно в пам'яті розміщені бути не можуть. 
%enditem


\newpage
8. Що буде надруковано за виконання?
%Оригинал был утерян. Требует отладки!
\begin{verbatim}
__d = 0

class B:
    __d = 1
    
    def __init__(self):
        self.__d = 2
        __d = 3
        B.__d = 4

obj = B()
B.__d = 5
try:
    print(__d, end=' ')
    print(obj.__d, end=' ')
    print(B.__d, end=' ')
    print(obj.__d, end=' ')
    print(obj._B__d)
except:
    print('ERR')
\end{verbatim}
%enditem


\newpage
9. Що буде надруковано за виконання?
\begin{verbatim}
class B0:
    def a(self): print(1, end=' ')
    def g(self): print(2, end=' ')

class B1(B0):
    def a(self): print(3, end=' ')

try:
    B1().g()
    B0().a()
except:
    print('ERR')
\end{verbatim}
%enditem


\newpage
10. Що буде надруковано за виконання?
\begin{verbatim}
class A:
    c = 1

    def __init__(self):
        b = 2
        b = 3

class B(A):
    c = 4
    a = 5

    def __init__(self):
        super().__init__()
        self.a = 6

try:
    obj = B()
    obj.b = obj.c + 10
    B.b = 13
except: print('EE', end=' ')

try: print(obj.b, end= ' ')
except: print('E1', end=' ')

try: print(A.b, end= ' ')
except: print('E2', end=' ')

try: print(B.b, end= ' ')
except: print('E3', end=' ')
\end{verbatim}
%enditem


\newpage
11. Змінні \kw{next} та \kw{prev} вузла списку позначають наступний та попередній за ним вузол відповідно. Починаючи з голови, у список записано послідовні цілі числа від~$-38$ до~$39$. Нехай змінна \kw{p} позначає вузол, що зберігає значення~$-5$.
Яке значення зберігається у вузлі \kw{p.prev.prev.next.next.next} ?
%enditem


\newpage
12. 
\begin{center}
	\adjustimage{max size={0.8\linewidth}{0.8\paperheight}}
	{BinTree/trees_exp/144.png}
\end{center}
Указати послідовність міток вершин дерева, зображеного на рис., що утвориться в результаті обходу дерева в прямому порядку. Мітки записувати через пробіл.\par Алгоритм починає роботу з кореня (на рис. це верхня вершина).
%enditem


\newpage
13. 
\begin{center}
	\adjustimage{max size={0.9\linewidth}{0.9\paperheight}}
	{Graphs/AdjList/95.png}
\end{center}
Для графа, зображеного на рис., вказати список суміжності вершини 3.\par
Формат оформлення відповіді:\par
список суміжності виписувати через пробіл на одному рядку, якщо пробілів десь буде більше одного - не важливо, упорядкування вершин у списку також значення не має.
%enditem


\newpage
14. Для графа на рисунку з вершини 4 запущено алгоритм пошуку в глибину, що будує кістякове дерево. Списки суміжності вершин переглядаються алгоритмом за зростанням міток вершин.\par
\begin{center}
	\adjustimage{max size={0.9\linewidth}{0.9\paperheight}}
	{Graphs/Traversals/11.png}
\end{center}
\par Які з наведених ребер входять до складу отриманого кістякового дерева?\par
1) (1, 5)\par
2) (2, 4)\par
3) (3, 5)\par
4) (0, 6)\par
5) (1, 4)\par
6) (2, 6)\par
{\vskip 15pt} Номери відповідних ребер записати за зростанням через пробіл.

%enditem


\newpage
15. Для графа на рисунку з вершини 4 запущено алгоритм пошуку в ширину, що будує кістякове дерево. Списки суміжності вершин переглядаються алгоритмом за зростанням міток вершин.\par
\begin{center}
	\adjustimage{max size={0.9\linewidth}{0.9\paperheight}}
	{Graphs/Traversals/237.png}
\end{center}
\par Які з наведених ребер входять до складу отриманого кістякового дерева?\par
1) (2, 5)\par
2) (0, 4)\par
3) (2, 6)\par
4) (1, 5)\par
5) (4, 5)\par
6) (1, 6)\par
{\vskip 15pt} Номери відповідних ребер записати за зростанням через пробіл.


%enditem


\newpage
16. 

Рядки журналу через двокрапку містять ім'я користувача (ідентифікатор), час події,тип події, код події, наприклад
\begin{verbatim}
s="username : 2022-05-05 13:26 : ERROR : 404"
\end{verbatim}
у коді 
\begin{verbatim}
res = re.fullmatch('???', s) 
s = ''
code_str = ??? 
\end{verbatim}
чим треба замінити кожен з ???, щоб значенням code\_str був код події? 


\vspace{10mm}
\textbf{Завдання 17} 

Задано програму, що містить код 1-2 класів, можливе успадкування, та клієнтський код.
У наведеному коді невеликі частини закриті. 
Також наведено вихід програми.
Необхідно відновити закриті частини, щоб отриманий код відповідав наведеному виходу програми.


\textbf{Завдання 18-19} 

Кілька запитань за завданнями для самостійної роботи.
\vspace{10mm}
Скоріш за все, що завдань буде трохи менше, а деякі попередні будуть переформульовані в стилі двох останніх.

%enditem

\end {large}}

%end

\newpage
%begin

{{
\begin {large}

\newpage
1. Що буде надруковано за виконання?

\begin{verbatim}
print('R', end=' ')
b=('a','b','c','d','e',0,1)
print(b[-8])
\end{verbatim}

%enditem


\newpage
2. Що буде надруковано за виконання?

\begin{verbatim}
print('R', end=' ')
h=[0,1,2,3,4,5,6,7,8,9]
print(h[-6::3])
\end{verbatim}

%enditem



\newpage
3. Що буде надруковано за виконання?

\begin{verbatim}
print('R', end=' ')
f=[10,11,12,13,14]
f.insert(1, [4,5])
print(f)
\end{verbatim}

%enditem



\newpage
4. Що буде надруковано за виконання?

\begin{verbatim}
print('R', end=' ')
d=['0','1','2','3','4']
d[4:]=[1,0]
print(d)
\end{verbatim}

%enditem



\newpage
5. Укажіть ТИП значення виразу.\par
\begin{verbatim}
{21,2,4}
\end{verbatim}
%enditem


\newpage
6. Що буде надруковано за виконання?\par
\begin{verbatim}
try:
    t = {63:3, 82:4, 48:6, 58:0, 82:6}
    t[58] = 2
    for x, y in t.items():
        print(x, y, sep=' ', end=' ')
except:
    print('ERR')
\end{verbatim}
%enditem


\newpage
7. Нехай змінна \kw{s} позначає послідовність цілих, по якій можна ітерувати. 

Написати вираз, що перевіряє, чи правда, що послідовність  містить хоча б одне додатнє число.
%enditem


\newpage
8. Що буде надруковано за виконання?
%Оригинал был утерян. Требует отладки!
\begin{verbatim}
c = 0

class C:
    c = 1
    
    def __init__(self):
        self.c = 2
        c = 3
        C.c = 4

obj = C()
C.c = 5
try:
    print(c, end=' ')
    print(obj.c, end=' ')
    print(C.c, end=' ')
    print(obj.c, end=' ')
    print(obj._C__c)
except:
    print('ERR')
\end{verbatim}
%enditem


\newpage
9. Що буде надруковано за виконання?
\begin{verbatim}
class C0:
    def d(self): print(1, end=' ')
    def h(self): print(2, end=' ')

class C1(C0):
    def d(self): print(3, end=' ')

try:
    C1().h()
    C1().d()
except:
    print('ERR')
\end{verbatim}
%enditem


\newpage
10. Що буде надруковано за виконання?
\begin{verbatim}
class A:
    b = 1

    def __init__(self):
        a = 2
        b = 3

class B(A):
    a = 4
    c = 5

    def __init__(self):
        super().__init__()
        self.c = 6

try:
    obj = B()
    obj.a = obj.b + 10
    B.a = 13
except: print('EE', end=' ')

try: print(obj.a, end= ' ')
except: print('E1', end=' ')

try: print(A.a, end= ' ')
except: print('E2', end=' ')

try: print(B.a, end= ' ')
except: print('E3', end=' ')
\end{verbatim}
%enditem


\newpage
11. Змінні \kw{next} та \kw{prev} вузла списку позначають наступний та попередній за ним вузол відповідно. Починаючи з голови, у список записано послідовні цілі числа від~$-34$ до~$30$. Нехай змінна \kw{p} позначає вузол, що зберігає значення~$-8$.
Яке значення зберігається у вузлі \kw{p.next.next.prev.prev.prev} ?
%enditem


\newpage
12. 
\begin{center}
	\adjustimage{max size={0.8\linewidth}{0.8\paperheight}}
	{BinTree/trees_exp/108.png}
\end{center}
Указати послідовність міток вершин дерева, зображеного на рис., що утвориться в результаті обходу дерева в оберненому порядку. Мітки записувати через пробіл.\par Алгоритм починає роботу з кореня (на рис. це верхня вершина).
%enditem


\newpage
13. 
\begin{center}
	\adjustimage{max size={0.9\linewidth}{0.9\paperheight}}
	{Graphs/AdjList/173.png}
\end{center}
Для графа, зображеного на рис., вказати список суміжності вершини 8.\par
Формат оформлення відповіді:\par
список суміжності виписувати через пробіл на одному рядку, якщо пробілів десь буде більше одного - не важливо, упорядкування вершин у списку також значення не має.
%enditem


\newpage
14. Для графа на рисунку з вершини 0 запущено алгоритм пошуку в глибину, що будує кістякове дерево. Списки суміжності вершин переглядаються алгоритмом за зростанням міток вершин.\par
\begin{center}
	\adjustimage{max size={0.9\linewidth}{0.9\paperheight}}
	{Graphs/Traversals/273.png}
\end{center}
\par Які з наведених ребер входять до складу отриманого кістякового дерева?\par
1) (3, 4)\par
2) (2, 5)\par
3) (0, 1)\par
4) (0, 3)\par
5) (1, 3)\par
6) (5, 6)\par
{\vskip 15pt} Номери відповідних ребер записати за зростанням через пробіл.

%enditem


\newpage
15. Для графа на рисунку з вершини 0 запущено алгоритм пошуку в ширину, що будує кістякове дерево. Списки суміжності вершин переглядаються алгоритмом за зростанням міток вершин.\par
\begin{center}
	\adjustimage{max size={0.9\linewidth}{0.9\paperheight}}
	{Graphs/Traversals/329.png}
\end{center}
\par Які з наведених ребер входять до складу отриманого кістякового дерева?\par
1) (1, 5)\par
2) (1, 3)\par
3) (3, 5)\par
4) (2, 6)\par
5) (4, 5)\par
6) (0, 3)\par
{\vskip 15pt} Номери відповідних ребер записати за зростанням через пробіл.


%enditem


\newpage
16. 

Рядки журналу через двокрапку містять ім'я користувача (ідентифікатор), час події,тип події, код події, наприклад
\begin{verbatim}
s="username : 2022-05-05 13:26 : ERROR : 404"
\end{verbatim}
у коді 
\begin{verbatim}
res = re.fullmatch('???', s) 
s = ''
code_str = ??? 
\end{verbatim}
чим треба замінити кожен з ???, щоб значенням code\_str був код події? 


\vspace{10mm}
\textbf{Завдання 17} 

Задано програму, що містить код 1-2 класів, можливе успадкування, та клієнтський код.
У наведеному коді невеликі частини закриті. 
Також наведено вихід програми.
Необхідно відновити закриті частини, щоб отриманий код відповідав наведеному виходу програми.


\textbf{Завдання 18-19} 

Кілька запитань за завданнями для самостійної роботи.
\vspace{10mm}
Скоріш за все, що завдань буде трохи менше, а деякі попередні будуть переформульовані в стилі двох останніх.

%enditem

\end {large}}

%end

\newpage
%begin

{{
\begin {large}

\newpage
1. Що буде надруковано за виконання?

\begin{verbatim}
print('R', end=' ')
c=('a','b','c','d','e',0,1,2)
print(c[-12])
\end{verbatim}

%enditem


\newpage
2. Що буде надруковано за виконання?

\begin{verbatim}
print('R', end=' ')
e=('0','1','2','3','4','5','6','7','8','9')
print(e[13:1:-3])
\end{verbatim}

%enditem



\newpage
3. Що буде надруковано за виконання?

\begin{verbatim}
print('R', end=' ')
h=[10,11,12,13,14]
print(h.index('12'), end=' ')
\end{verbatim}

%enditem



\newpage
4. Що буде надруковано за виконання?

\begin{verbatim}
print('R', end=' ')
f=[10,11,12,13,14]
del f[:-1]
print(f)
\end{verbatim}

%enditem



\newpage
5. Укажіть ТИП значення виразу.\par
\begin{verbatim}
{y:13 for y in range(-7,7)}
\end{verbatim}
%enditem


\newpage
6. Що буде надруковано за виконання?\par
\begin{verbatim}
try:
    s = {68:1, 38:2, 27:9, 68:3}
    s[87] = 8
    for x in s.items():
        print(x, sep=' ', end=' ')
except:
    print('ERR')
\end{verbatim}
%enditem


\newpage
7. Записати ВИРАЗ, значенням якого є множина, що складається з цілих чисел від 17 до 2009 (включно), що діляться на 7.
%enditem


\newpage
8. Що буде надруковано за виконання?
\begin{verbatim}
d = 0
k = 1

class B:
    d = 2
    
    def __init__(self):
        global d
        self.d = 3
        d = 4
        B.d = 5

obj = B()
try:
    print(d, end=' ')
    print(k, end=' ')
    print(B.d, end=' ')
    print(obj.d)
except:
    print('ERR')
\end{verbatim}
%enditem


\newpage
9. Що буде надруковано за виконання?
\begin{verbatim}
class D0:
    def b(self): print(1, end=' ')
    def f(self): print(2, end=' ')

class D1(D0):
    def f(self): print(3, end=' ')

try:
    D0().b()
    D1().f()
except:
    print('ERR')
\end{verbatim}
%enditem


\newpage
10. Що буде надруковано за виконання?
\begin{verbatim}
class A:
    a = 1

    def __init__(self):
        c = 2
        c = 3

class B(A):
    c = 4
    b = 5

    def __init__(self):
        super().__init__()
        self.b = 6

try:
    obj = B()
    obj.c = obj.c + 10
    B.c = 13
except: print('EE', end=' ')

try: print(obj.c, end= ' ')
except: print('E1', end=' ')

try: print(A.c, end= ' ')
except: print('E2', end=' ')

try: print(B.c, end= ' ')
except: print('E3', end=' ')
\end{verbatim}
%enditem


\newpage
11. Змінні \kw{next} та \kw{prev} вузла списку позначають наступний та попередній за ним вузол відповідно. Починаючи з голови, у список записано послідовні цілі числа від~$-68$ до~$45$. Нехай змінна \kw{p} позначає вузол, що зберігає значення~$7$.
Яке значення зберігається у вузлі \kw{p.next.prev.next.prev.prev} ?
%enditem


\newpage
12. 
\begin{center}
	\adjustimage{max size={0.8\linewidth}{0.8\paperheight}}
	{BinTree/trees_exp/99.png}
\end{center}
Указати послідовність міток вершин дерева, зображеного на рис., що утвориться в результаті обходу дерева в прямому порядку. Мітки записувати через пробіл.\par Алгоритм починає роботу з кореня (на рис. це верхня вершина).
%enditem


\newpage
13. 
\begin{center}
	\adjustimage{max size={0.9\linewidth}{0.9\paperheight}}
	{Graphs/AdjList/23.png}
\end{center}
Для графа, зображеного на рис., вказати список суміжності вершини 3.\par
Формат оформлення відповіді:\par
список суміжності виписувати через пробіл на одному рядку, якщо пробілів десь буде більше одного - не важливо, упорядкування вершин у списку також значення не має.
%enditem


\newpage
14. Для графа на рисунку вказати послідовність, в якій алгоритм пошуку в глибину буде відмічати відвідані вершини, якщо списки суміжності вершин переглядаються алгоритмом за зростанням міток вершин та вершини відмічаються на першому вход\-женні у вершину.\par
Пошук розпочинається з вершини 3.\par

\begin{center}
	\adjustimage{max size={0.9\linewidth}{0.9\paperheight}}
	{Graphs/Traversals/21.png}
\end{center}
%enditem


\newpage
15. Для графа на рисунку з вершини 0 запущено алгоритм пошуку в ширину, що будує кістякове дерево. Списки суміжності вершин переглядаються алгоритмом за зростанням міток вершин.\par
\begin{center}
	\adjustimage{max size={0.9\linewidth}{0.9\paperheight}}
	{Graphs/Traversals/47.png}
\end{center}
\par Які з наведених ребер входять до складу отриманого кістякового дерева?\par
1) (3, 6)\par
2) (3, 5)\par
3) (2, 5)\par
4) (0, 5)\par
5) (0, 3)\par
6) (4, 6)\par
{\vskip 15pt} Номери відповідних ребер записати за зростанням через пробіл.


%enditem


\newpage
16. 

Рядки журналу через двокрапку містять ім'я користувача (ідентифікатор), час події,тип події, код події, наприклад
\begin{verbatim}
s="username : 2022-05-05 13:26 : ERROR : 404"
\end{verbatim}
у коді 
\begin{verbatim}
res = re.fullmatch('???', s) 
s = ''
code_str = ??? 
\end{verbatim}
чим треба замінити кожен з ???, щоб значенням code\_str був код події? 


\vspace{10mm}
\textbf{Завдання 17} 

Задано програму, що містить код 1-2 класів, можливе успадкування, та клієнтський код.
У наведеному коді невеликі частини закриті. 
Також наведено вихід програми.
Необхідно відновити закриті частини, щоб отриманий код відповідав наведеному виходу програми.


\textbf{Завдання 18-19} 

Кілька запитань за завданнями для самостійної роботи.
\vspace{10mm}
Скоріш за все, що завдань буде трохи менше, а деякі попередні будуть переформульовані в стилі двох останніх.

%enditem

\end {large}}

%end

\newpage
%begin

{{
\begin {large}

\newpage
1. Що буде надруковано за виконання?

\begin{verbatim}
print('R', end=' ')
f=('a','b','c','d','e',0,1,2,3,4)
print(f[-13])
\end{verbatim}

%enditem


\newpage
2. Що буде надруковано за виконання?

\begin{verbatim}
print('R', end=' ')
g=['0','1','2','3','4','5','6','7','8','9']
print(g[3:11:2])
\end{verbatim}

%enditem



\newpage
3. Що буде надруковано за виконання?

\begin{verbatim}
print('R', end=' ')
d=['0','1','2','3','4']
d.extend('12')
print(d)
\end{verbatim}

%enditem



\newpage
4. Що буде надруковано за виконання?

\begin{verbatim}
print('R', end=' ')
g=['0','1','2','3','4']
g[-5:-8]='13'
print(g)
\end{verbatim}

%enditem



\newpage
5. Укажіть ТИП значення виразу.\par
\begin{verbatim}
{x for x in range(7)}
\end{verbatim}
%enditem


\newpage
6. Що буде надруковано за виконання?\par
\begin{verbatim}
try:
    d = {25:7, 47:7, 13:3, 13:6}
    d[12] = 7
    for x in d :
        print(x, sep=' ', end=' ')
except:
    print('ERR')
\end{verbatim}
%enditem


\newpage
7. Нехай змінна \kw{s} позначає послідовність цілих, по якій можна ітерувати.

Написати вираз, що перевіряє, чи правда, що всі елементи послідовності є від’ємни\-ми.
%enditem


\newpage
8. Що буде надруковано за виконання?
%Оригинал был утерян. Требует отладки!
\begin{verbatim}
_d = 0

class D:
    _d = 1
    
    def __init__(self):
        self._d = 2
        _d = 3
        D._d = 4

obj = D()
obj._d = 5
try:
    print(_d, end=' ')
    print(obj._d, end=' ')
    print(D._d, end=' ')
    print(obj._d, end=' ')
    print(obj._D__d)
except:
    print('ERR')
\end{verbatim}
%enditem


\newpage
9. Що буде надруковано за виконання?
\begin{verbatim}
class A0:
    def c(self): print(1, end=' ')
    def g(self): print(2, end=' ')

class A1(A0):
    def g(self): print(3, end=' ')

try:
    A1().c()
    A1().g()
except:
    print('ERR')
\end{verbatim}
%enditem


\newpage
10. Що буде надруковано за виконання?
\begin{verbatim}
class A:
    c = 1

    def __init__(self):
        a = 2
        c = 3

class B(A):
    c = 4
    b = 5

    def __init__(self):
        super().__init__()
        self.b = 6

try:
    obj = B()
    obj.a = obj.b + 10
    B.a = 13
except: print('EE', end=' ')

try: print(obj.a, end= ' ')
except: print('E1', end=' ')

try: print(A.a, end= ' ')
except: print('E2', end=' ')

try: print(B.a, end= ' ')
except: print('E3', end=' ')
\end{verbatim}
%enditem


\newpage
11. Змінні \kw{next} та \kw{prev} вузла списку позначають наступний та попередній за ним вузол відповідно. Починаючи з голови, у список записано послідовні цілі числа від~$-58$ до~$38$. Нехай змінна \kw{p} позначає вузол, що зберігає значення~$-2$.
Яке значення зберігається у вузлі \kw{p.prev.next.prev.next.next} ?
%enditem


\newpage
12. 
\begin{center}
	\adjustimage{max size={0.8\linewidth}{0.8\paperheight}}
	{BinTree/trees_exp/118.png}
\end{center}
Указати послідовність міток вершин дерева, зображеного на рис., що утвориться в результаті обходу дерева в оберненому порядку. Мітки записувати через пробіл.\par Алгоритм починає роботу з кореня (на рис. це верхня вершина).
%enditem


\newpage
13. 
\begin{center}
	\adjustimage{max size={0.9\linewidth}{0.9\paperheight}}
	{Graphs/AdjList/101.png}
\end{center}
Для графа, зображеного на рис., вказати список суміжності вершини 3.\par
Формат оформлення відповіді:\par
список суміжності виписувати через пробіл на одному рядку, якщо пробілів десь буде більше одного - не важливо, упорядкування вершин у списку також значення не має.
%enditem


\newpage
14. Для графа на рисунку вказати послідовність, в якій алгоритм пошуку в глибину буде відмічати відвідані вершини, якщо списки суміжності вершин переглядаються алгоритмом за зростанням міток вершин та вершини відмічаються на першому вход\-женні у вершину.\par
Пошук розпочинається з вершини 4.\par

\begin{center}
	\adjustimage{max size={0.9\linewidth}{0.9\paperheight}}
	{Graphs/Traversals/302.png}
\end{center}
%enditem


\newpage
15. Для графа на рисунку з вершини 0 запущено алгоритм пошуку в ширину, що будує кістякове дерево. Списки суміжності вершин переглядаються алгоритмом за зростанням міток вершин.\par
\begin{center}
	\adjustimage{max size={0.9\linewidth}{0.9\paperheight}}
	{Graphs/Traversals/79.png}
\end{center}
\par Які з наведених ребер входять до складу отриманого кістякового дерева?\par
1) (1, 4)\par
2) (1, 6)\par
3) (0, 5)\par
4) (2, 6)\par
5) (0, 4)\par
6) (3, 5)\par
{\vskip 15pt} Номери відповідних ребер записати за зростанням через пробіл.


%enditem


\newpage
16. 

Рядки журналу через двокрапку містять ім'я користувача (ідентифікатор), час події,тип події, код події, наприклад
\begin{verbatim}
s="username : 2022-05-05 13:26 : ERROR : 404"
\end{verbatim}
у коді 
\begin{verbatim}
res = re.fullmatch('???', s) 
s = ''
code_str = ??? 
\end{verbatim}
чим треба замінити кожен з ???, щоб значенням code\_str був код події? 


\vspace{10mm}
\textbf{Завдання 17} 

Задано програму, що містить код 1-2 класів, можливе успадкування, та клієнтський код.
У наведеному коді невеликі частини закриті. 
Також наведено вихід програми.
Необхідно відновити закриті частини, щоб отриманий код відповідав наведеному виходу програми.


\textbf{Завдання 18-19} 

Кілька запитань за завданнями для самостійної роботи.
\vspace{10mm}
Скоріш за все, що завдань буде трохи менше, а деякі попередні будуть переформульовані в стилі двох останніх.

%enditem

\end {large}}

%end

\newpage
%begin

{{
\begin {large}

\newpage
1. Що буде надруковано за виконання?

\begin{verbatim}
print('R', end=' ')
e=('a','b','c','d','e',0,1,2,3)
print(e[-2])
\end{verbatim}

%enditem


\newpage
2. Що буде надруковано за виконання?

\begin{verbatim}
print('R', end=' ')
f='0123456789'
print(f[14::-1])
\end{verbatim}

%enditem



\newpage
3. Що буде надруковано за виконання?

\begin{verbatim}
print('R', end=' ')
e=['0','1','2','3','4']
e.remove(0)
print(e)
\end{verbatim}

%enditem



\newpage
4. Що буде надруковано за виконання?

\begin{verbatim}
print('R', end=' ')
c=[10,11,12,13,14]
del c[0:6]
print(c)
\end{verbatim}

%enditem



\newpage
5. Укажіть ТИП значення виразу.\par
\begin{verbatim}
[7,2,3,4]
\end{verbatim}
%enditem


\newpage
6. Що буде надруковано за виконання?\par
\begin{verbatim}
try:
    t = {39:7, 74:1, 47:8, 18:9}
    t[74] = 0
    for x, y in t.items():
        print(x, y, sep=' ', end=' ')
except:
    print('ERR')
\end{verbatim}
%enditem


\newpage
7. Нехай змінна \kw{s} позначає дуже довгу послідовність, по якій можна ітерувати. 

Написати вираз-генератор, що задає підпослідовність її елементів, що не є числами типу \kw{float}.

Обмеження: усі елементи шуканої підпослідовності одночасно в пам'яті розміщені бути не можуть. 
%enditem


\newpage
8. Що буде надруковано за виконання?
\begin{verbatim}
a = 0
n = 1

class A:
    a = 2
    
    def __init__(self):
        self.a = 3
        a = 4
        A.a = 5

obj = A()
try:
    print(a, end=' ')
    print(n, end=' ')
    print(A.a, end=' ')
    print(obj.a)
except:
    print('ERR')
\end{verbatim}
%enditem


\newpage
9. Що буде надруковано за виконання?
\begin{verbatim}
class C0:
    def c(self): print(1, end=' ')
    def f(self): print(2, end=' ')

class C1(C0):
    def c(self): print(3, end=' ')

try:
    C1().f()
    C1().c()
except:
    print('ERR')
\end{verbatim}
%enditem


\newpage
10. Що буде надруковано за виконання?
\begin{verbatim}
class A:
    d = 1

    def _g1(self): print(2, end=' ')
    def g2(self): print(3, end=' ')

    def main(self):
        try: print(self.d, end=' ')
        except: print('E1', end=' ')

        try: self._g1()
        except: print('E2', end=' ')

        try: self.g2()
        except: print('E3', end=' ')


class B(A):
    d = 4

    def _g1(self): print(5, end=' ')
    def g2(self): print(6, end=' ')

try: B().main()
except: print('E4', end=' ')
\end{verbatim}
%enditem


\newpage
11. Змінні \kw{next} та \kw{prev} вузла списку позначають наступний та попередній за ним вузол відповідно. Починаючи з голови, у список записано послідовні цілі числа від~$-45$ до~$41$. Нехай змінна \kw{p} позначає вузол, що зберігає значення~$5$.
Яке значення зберігається у вузлі \kw{p.next.next.next.prev.prev} ?
%enditem


\newpage
12. 
\begin{center}
	\adjustimage{max size={0.8\linewidth}{0.8\paperheight}}
	{BinTree/trees_exp/140.png}
\end{center}
Указати послідовність міток вершин дерева, зображеного на рис., що утвориться в результаті обходу дерева в оберненому порядку. Мітки записувати через пробіл.\par Алгоритм починає роботу з кореня (на рис. це верхня вершина).
%enditem


\newpage
13. 
\begin{center}
	\adjustimage{max size={0.9\linewidth}{0.9\paperheight}}
	{Graphs/AdjList/205.png}
\end{center}
Для графа, зображеного на рис., вказати список суміжності вершини 7.\par
Формат оформлення відповіді:\par
список суміжності виписувати через пробіл на одному рядку, якщо пробілів десь буде більше одного - не важливо, упорядкування вершин у списку також значення не має.
%enditem


\newpage
14. Для графа на рисунку вказати послідовність, в якій алгоритм пошуку в глибину буде відмічати відвідані вершини, якщо списки суміжності вершин переглядаються алгоритмом за зростанням міток вершин та вершини відмічаються на першому вход\-женні у вершину.\par
Пошук розпочинається з вершини 2.\par

\begin{center}
	\adjustimage{max size={0.9\linewidth}{0.9\paperheight}}
	{Graphs/Traversals/316.png}
\end{center}
%enditem


\newpage
15. Для графа на рисунку з вершини 1 запущено алгоритм пошуку в ширину, що будує кістякове дерево. Списки суміжності вершин переглядаються алгоритмом за зростанням міток вершин.\par
\begin{center}
	\adjustimage{max size={0.9\linewidth}{0.9\paperheight}}
	{Graphs/Traversals/37.png}
\end{center}
\par Які з наведених ребер входять до складу отриманого кістякового дерева?\par
1) (0, 5)\par
2) (1, 6)\par
3) (2, 6)\par
4) (1, 4)\par
5) (3, 4)\par
6) (2, 5)\par
{\vskip 15pt} Номери відповідних ребер записати за зростанням через пробіл.


%enditem


\newpage
16. 

Рядки журналу через двокрапку містять ім'я користувача (ідентифікатор), час події,тип події, код події, наприклад
\begin{verbatim}
s="username : 2022-05-05 13:26 : ERROR : 404"
\end{verbatim}
у коді 
\begin{verbatim}
res = re.fullmatch('???', s) 
s = ''
code_str = ??? 
\end{verbatim}
чим треба замінити кожен з ???, щоб значенням code\_str був код події? 


\vspace{10mm}
\textbf{Завдання 17} 

Задано програму, що містить код 1-2 класів, можливе успадкування, та клієнтський код.
У наведеному коді невеликі частини закриті. 
Також наведено вихід програми.
Необхідно відновити закриті частини, щоб отриманий код відповідав наведеному виходу програми.


\textbf{Завдання 18-19} 

Кілька запитань за завданнями для самостійної роботи.
\vspace{10mm}
Скоріш за все, що завдань буде трохи менше, а деякі попередні будуть переформульовані в стилі двох останніх.

%enditem

\end {large}}

%end

\newpage
%begin

{{
\begin {large}

\newpage
1. Що буде надруковано за виконання?

\begin{verbatim}
print('R', end=' ')
d=('a','b','c','d','e',0,1,2,4)
print(d[11])
\end{verbatim}

%enditem


\newpage
2. Що буде надруковано за виконання?

\begin{verbatim}
print('R', end=' ')
c=(0,1,2,3,4,5,6,7,8,9)
print(c[-1:6:-2])
\end{verbatim}

%enditem



\newpage
3. Що буде надруковано за виконання?

\begin{verbatim}
print('R', end=' ')
g=[10,11,12,13,14]
print(g.pop(1), end=' ')
print(g)
\end{verbatim}

%enditem



\newpage
4. Що буде надруковано за виконання?

\begin{verbatim}
print('R', end=' ')
a=['0','1','2','3','4']
a[-6:8]='10'
print(a)
\end{verbatim}

%enditem



\newpage
5. Укажіть ТИП значення виразу.\par
\begin{verbatim}
{x:x for x in range(7)}
\end{verbatim}
%enditem


\newpage
6. Що буде надруковано за виконання?\par
\begin{verbatim}
try:
    s = {10:9, 51:7, 55:3, 33:5, 51:0}
    s[85] = 2
    for x in s :
        print(x, sep=' ', end=' ')
except:
    print('ERR')
\end{verbatim}
%enditem


\newpage
7. Нехай змінна \kw{s} позначає дуже довгу послідовність, по якій можна ітерувати. 

Написати вираз-генератор, що задає підпослідовність її елементів, що є числами типу \kw{int}.

Обмеження: усі елементи шуканої підпослідовності одночасно в пам'яті розміщені бути не можуть. 
%enditem


\newpage
8. Що буде надруковано за виконання?
\begin{verbatim}
b = 0
j = 1

class D:
    j = 2
    
    def __init__(self):
        global j
        self.j = 3
        j = 4
        D.j = 5

obj = D()
try:
    print(b, end=' ')
    print(j, end=' ')
    print(D.j, end=' ')
    print(obj.j)
except:
    print('ERR')
\end{verbatim}
%enditem


\newpage
9. Що буде надруковано за виконання?
\begin{verbatim}
class A0:
    def d(self): print(1, end=' ')
    def g(self): print(2, end=' ')

class A1(A0):
    def g(self): print(3, end=' ')

try:
    A0().d()
    A1().g()
except:
    print('ERR')
\end{verbatim}
%enditem


\newpage
10. Що буде надруковано за виконання?
\begin{verbatim}
class A:
    a = 1

    def __init__(self):
        b = 2
        a = 3

class B(A):
    a = 4
    c = 5

    def __init__(self):
        super().__init__()
        self.c = 6

try:
    obj = B()
    obj.a = obj.c + 10
    B.a = 13
except: print('EE', end=' ')

try: print(obj.a, end= ' ')
except: print('E1', end=' ')

try: print(A.a, end= ' ')
except: print('E2', end=' ')

try: print(B.a, end= ' ')
except: print('E3', end=' ')
\end{verbatim}
%enditem


\newpage
11. Змінні \kw{next} та \kw{prev} вузла списку позначають наступний та попередній за ним вузол відповідно. Починаючи з голови, у список записано послідовні цілі числа від~$-28$ до~$37$. Нехай змінна \kw{p} позначає вузол, що зберігає значення~$9$.
Яке значення зберігається у вузлі \kw{p.prev.prev.prev.next.next} ?
%enditem


\newpage
12. 
\begin{center}
	\adjustimage{max size={0.8\linewidth}{0.8\paperheight}}
	{BinTree/trees_exp/66.png}
\end{center}
Указати послідовність міток вершин дерева, зображеного на рис., що утвориться в результаті обходу дерева в прямому порядку. Мітки записувати через пробіл.\par Алгоритм починає роботу з кореня (на рис. це верхня вершина).
%enditem


\newpage
13. 
\begin{center}
	\adjustimage{max size={0.9\linewidth}{0.9\paperheight}}
	{Graphs/AdjList/159.png}
\end{center}
Для графа, зображеного на рис., вказати список суміжності вершини 5.\par
Формат оформлення відповіді:\par
список суміжності виписувати через пробіл на одному рядку, якщо пробілів десь буде більше одного - не важливо, упорядкування вершин у списку також значення не має.
%enditem


\newpage
14. Для графа на рисунку вказати послідовність, в якій алгоритм пошуку в глибину буде відмічати відвідані вершини, якщо списки суміжності вершин переглядаються алгоритмом за зростанням міток вершин та вершини відмічаються на першому вход\-женні у вершину.\par
Пошук розпочинається з вершини 3.\par

\begin{center}
	\adjustimage{max size={0.9\linewidth}{0.9\paperheight}}
	{Graphs/Traversals/328.png}
\end{center}
%enditem


\newpage
15. Для графа на рисунку з вершини 0 запущено алгоритм пошуку в ширину, що будує кістякове дерево. Списки суміжності вершин переглядаються алгоритмом за зростанням міток вершин.\par
\begin{center}
	\adjustimage{max size={0.9\linewidth}{0.9\paperheight}}
	{Graphs/Traversals/115.png}
\end{center}
\par Які з наведених ребер входять до складу отриманого кістякового дерева?\par
1) (1, 5)\par
2) (2, 3)\par
3) (0, 2)\par
4) (4, 6)\par
5) (0, 3)\par
6) (3, 4)\par
{\vskip 15pt} Номери відповідних ребер записати за зростанням через пробіл.


%enditem


\newpage
16. 

Рядки журналу через двокрапку містять ім'я користувача (ідентифікатор), час події,тип події, код події, наприклад
\begin{verbatim}
s="username : 2022-05-05 13:26 : ERROR : 404"
\end{verbatim}
у коді 
\begin{verbatim}
res = re.fullmatch('???', s) 
s = ''
code_str = ??? 
\end{verbatim}
чим треба замінити кожен з ???, щоб значенням code\_str був код події? 


\vspace{10mm}
\textbf{Завдання 17} 

Задано програму, що містить код 1-2 класів, можливе успадкування, та клієнтський код.
У наведеному коді невеликі частини закриті. 
Також наведено вихід програми.
Необхідно відновити закриті частини, щоб отриманий код відповідав наведеному виходу програми.


\textbf{Завдання 18-19} 

Кілька запитань за завданнями для самостійної роботи.
\vspace{10mm}
Скоріш за все, що завдань буде трохи менше, а деякі попередні будуть переформульовані в стилі двох останніх.

%enditem

\end {large}}

%end

\newpage
%begin

{{
\begin {large}

\newpage
1. Що буде надруковано за виконання?

\begin{verbatim}
print('R', end=' ')
g=('a','b','c','d','e',0,1,2,3)
print(g[-1])
\end{verbatim}

%enditem


\newpage
2. Що буде надруковано за виконання?

\begin{verbatim}
print('R', end=' ')
a='0123456789'
print(a[4:-6:2])
\end{verbatim}

%enditem



\newpage
3. Що буде надруковано за виконання?

\begin{verbatim}
print('R', end=' ')
c=['0','1','2','3','4']
c.append(20)
print(c)
\end{verbatim}

%enditem



\newpage
4. Що буде надруковано за виконання?

\begin{verbatim}
print('R', end=' ')
b=[10,11,12,13,14]
b[:4]=[2,3]
print(b)
\end{verbatim}

%enditem



\newpage
5. Укажіть ТИП значення виразу.\par
\begin{verbatim}
(11,2,3)
\end{verbatim}
%enditem


\newpage
6. Що буде надруковано за виконання?\par
\begin{verbatim}
try:
    d = {79:2, 87:7, 83:6, 47:7, 83:3}
    d[79] = 9
    for x in d.items():
        print(*x, sep=' ', end=' ')
except:
    print('ERR')
\end{verbatim}
%enditem


\newpage
7. Нехай змінна \kw{s} позначає дуже довгу послідовність, по якій можна ітерувати. 

Написати вираз-генератор, що задає підпослідовність її елементів, що не є числами типу \kw{int}.

Обмеження: усі елементи шуканої підпослідовності одночасно в пам'яті розміщені бути не можуть. 
%enditem


\newpage
8. Що буде надруковано за виконання?
%Оригинал был утерян. Требует отладки!
\begin{verbatim}
_a = 0

class A:
    _a = 1
    
    def __init__(self):
        self._a = 2
        _a = 3
        A._a = 4

obj = A()
obj._a = 5
try:
    print(_a, end=' ')
    print(obj._a, end=' ')
    print(A._a, end=' ')
    print(obj._a, end=' ')
    print(obj._A__a)
except:
    print('ERR')
\end{verbatim}
%enditem


\newpage
9. Що буде надруковано за виконання?
\begin{verbatim}
class B0:
    def a(self): print(1, end=' ')
    def h(self): print(2, end=' ')

class B1(B0):
    def h(self): print(3, end=' ')

try:
    B1().a()
    B1().h()
except:
    print('ERR')
\end{verbatim}
%enditem


\newpage
10. Що буде надруковано за виконання?
\begin{verbatim}
class A:
    c = 1

    def __init__(self):
        a = 2
        a = 3

class B(A):
    a = 4
    b = 5

    def __init__(self):
        super().__init__()
        self.b = 6

try:
    obj = B()
    obj.b = obj.a + 10
    B.b = 13
except: print('EE', end=' ')

try: print(obj.b, end= ' ')
except: print('E1', end=' ')

try: print(A.b, end= ' ')
except: print('E2', end=' ')

try: print(B.b, end= ' ')
except: print('E3', end=' ')
\end{verbatim}
%enditem


\newpage
11. Змінні \kw{next} та \kw{prev} вузла списку позначають наступний та попередній за ним вузол відповідно. Починаючи з голови, у список записано послідовні цілі числа від~$-31$ до~$50$. Нехай змінна \kw{p} позначає вузол, що зберігає значення~$-7$.
Яке значення зберігається у вузлі \kw{p.prev.next.prev.prev.prev} ?
%enditem


\newpage
12. 
\begin{center}
	\adjustimage{max size={0.8\linewidth}{0.8\paperheight}}
	{BinTree/trees_exp/82.png}
\end{center}
Указати послідовність міток вершин дерева, зображеного на рис., що утвориться в результаті обходу дерева в прямому порядку. Мітки записувати через пробіл.\par Алгоритм починає роботу з кореня (на рис. це верхня вершина).
%enditem


\newpage
13. 
\begin{center}
	\adjustimage{max size={0.9\linewidth}{0.9\paperheight}}
	{Graphs/AdjList/269.png}
\end{center}
Для графа, зображеного на рис., вказати список суміжності вершини 8.\par
Формат оформлення відповіді:\par
список суміжності виписувати через пробіл на одному рядку, якщо пробілів десь буде більше одного - не важливо, упорядкування вершин у списку також значення не має.
%enditem


\newpage
14. Для графа на рисунку з вершини 5 запущено алгоритм пошуку в глибину, що будує кістякове дерево. Списки суміжності вершин переглядаються алгоритмом за зростанням міток вершин.\par
\begin{center}
	\adjustimage{max size={0.9\linewidth}{0.9\paperheight}}
	{Graphs/Traversals/339.png}
\end{center}
\par Які з наведених ребер входять до складу отриманого кістякового дерева?\par
1) (2, 5)\par
2) (2, 6)\par
3) (3, 4)\par
4) (1, 3)\par
5) (5, 6)\par
6) (0, 1)\par
{\vskip 15pt} Номери відповідних ребер записати за зростанням через пробіл.

%enditem


\newpage
15. Для графа на рисунку з вершини 2 запущено алгоритм пошуку в ширину, що будує кістякове дерево. Списки суміжності вершин переглядаються алгоритмом за зростанням міток вершин.\par
\begin{center}
	\adjustimage{max size={0.9\linewidth}{0.9\paperheight}}
	{Graphs/Traversals/310.png}
\end{center}
\par Які з наведених ребер входять до складу отриманого кістякового дерева?\par
1) (0, 5)\par
2) (2, 6)\par
3) (1, 2)\par
4) (0, 6)\par
5) (0, 4)\par
6) (3, 6)\par
{\vskip 15pt} Номери відповідних ребер записати за зростанням через пробіл.


%enditem


\newpage
16. 

Рядки журналу через двокрапку містять ім'я користувача (ідентифікатор), час події,тип події, код події, наприклад
\begin{verbatim}
s="username : 2022-05-05 13:26 : ERROR : 404"
\end{verbatim}
у коді 
\begin{verbatim}
res = re.fullmatch('???', s) 
s = ''
code_str = ??? 
\end{verbatim}
чим треба замінити кожен з ???, щоб значенням code\_str був код події? 


\vspace{10mm}
\textbf{Завдання 17} 

Задано програму, що містить код 1-2 класів, можливе успадкування, та клієнтський код.
У наведеному коді невеликі частини закриті. 
Також наведено вихід програми.
Необхідно відновити закриті частини, щоб отриманий код відповідав наведеному виходу програми.


\textbf{Завдання 18-19} 

Кілька запитань за завданнями для самостійної роботи.
\vspace{10mm}
Скоріш за все, що завдань буде трохи менше, а деякі попередні будуть переформульовані в стилі двох останніх.

%enditem

\end {large}}

%end

\newpage
%begin

{{
\begin {large}

\newpage
1. Що буде надруковано за виконання?

\begin{verbatim}
print('R', end=' ')
a=('a','b','c','d','e',0,1,2)
print(a[4])
\end{verbatim}

%enditem


\newpage
2. Що буде надруковано за виконання?

\begin{verbatim}
print('R', end=' ')
d=[0,1,2,3,4,5,6,7,8,9]
print(d[15:-14:-1])
\end{verbatim}

%enditem



\newpage
3. Що буде надруковано за виконання?

\begin{verbatim}
print('R', end=' ')
b=[10,11,12,13,14]
b+=15
print(b)
\end{verbatim}

%enditem



\newpage
4. Що буде надруковано за виконання?

\begin{verbatim}
print('R', end=' ')
h=['0','1','2','3','4']
h[5:]=[1,2]
print(h)
\end{verbatim}

%enditem



\newpage
5. Укажіть ТИП значення виразу.\par
\begin{verbatim}
{x*x:x for x in range(-7,7)}
\end{verbatim}
%enditem


\newpage
6. Що буде надруковано за виконання?\par
\begin{verbatim}
try:
    t = {19:4, 41:7, 27:9, 61:8, 41:2}
    t[19] = 2
    for x in t.items():
        print(x, sep=' ', end=' ')
except:
    print('ERR')
\end{verbatim}
%enditem


\newpage
7. 
 Задано числову послідовність \kw{s} довжини 6000. Записати ВИРАЗ, значенням якого є список, що складається з перших 2019 елементів послідовності, записаних в обернено\-му порядку.
%enditem


\newpage
8. Що буде надруковано за виконання?
%Оригинал был утерян. Требует отладки!
\begin{verbatim}
__b = 0

class B:
    __b = 1
    
    def __init__(self):
        self.__b = 2
        __b = 3
        B.__b = 4

obj = B()
B.__b = 5
try:
    print(__b, end=' ')
    print(obj.__b, end=' ')
    print(B.__b, end=' ')
    print(obj.__b, end=' ')
    print(obj._B__b)
except:
    print('ERR')
\end{verbatim}
%enditem


\newpage
9. Що буде надруковано за виконання?
\begin{verbatim}
class D0:
    def b(self): print(1, end=' ')
    def g(self): print(2, end=' ')

class D1(D0):
    def g(self): print(3, end=' ')

try:
    D1().b()
    D0().g()
except:
    print('ERR')
\end{verbatim}
%enditem


\newpage
10. Що буде надруковано за виконання?
\begin{verbatim}
class A:
    a = 1

    def __h1(self): print(2, end=' ')
    def _h2(self): print(3, end=' ')

    def main(self):
        try: print(self.a, end=' ')
        except: print('E1', end=' ')

        try: self.__h1()
        except: print('E2', end=' ')

        try: self._h2()
        except: print('E3', end=' ')


class B(A):
    a = 4

    def __h1(self): print(5, end=' ')
    def _h2(self): print(6, end=' ')

try: B().main()
except: print('E4', end=' ')
\end{verbatim}
%enditem


\newpage
11. Змінні \kw{next} та \kw{prev} вузла списку позначають наступний та попередній за ним вузол відповідно. Починаючи з голови, у список записано послідовні цілі числа від~$-20$ до~$31$. Нехай змінна \kw{p} позначає вузол, що зберігає значення~$2$.
Яке значення зберігається у вузлі \kw{p.next.prev.next.next.next} ?
%enditem


\newpage
12. 
\begin{center}
	\adjustimage{max size={0.8\linewidth}{0.8\paperheight}}
	{BinTree/trees_exp/169.png}
\end{center}
Указати послідовність міток вершин дерева, зображеного на рис., що утвориться в результаті обходу дерева в оберненому порядку. Мітки записувати через пробіл.\par Алгоритм починає роботу з кореня (на рис. це верхня вершина).
%enditem


\newpage
13. 
\begin{center}
	\adjustimage{max size={0.9\linewidth}{0.9\paperheight}}
	{Graphs/AdjList/112.png}
\end{center}
Для графа, зображеного на рис., вказати список суміжності вершини 6.\par
Формат оформлення відповіді:\par
список суміжності виписувати через пробіл на одному рядку, якщо пробілів десь буде більше одного - не важливо, упорядкування вершин у списку також значення не має.
%enditem


\newpage
14. Для графа на рисунку вказати послідовність, в якій алгоритм пошуку в глибину буде відмічати відвідані вершини, якщо списки суміжності вершин переглядаються алгоритмом за зростанням міток вершин та вершини відмічаються на першому вход\-женні у вершину.\par
Пошук розпочинається з вершини 0.\par

\begin{center}
	\adjustimage{max size={0.9\linewidth}{0.9\paperheight}}
	{Graphs/Traversals/148.png}
\end{center}
%enditem


\newpage
15. Для графа на рисунку з вершини 3 запущено алгоритм пошуку в ширину, що будує кістякове дерево. Списки суміжності вершин переглядаються алгоритмом за зростанням міток вершин.\par
\begin{center}
	\adjustimage{max size={0.9\linewidth}{0.9\paperheight}}
	{Graphs/Traversals/202.png}
\end{center}
\par Які з наведених ребер входять до складу отриманого кістякового дерева?\par
1) (3, 5)\par
2) (1, 6)\par
3) (3, 4)\par
4) (4, 6)\par
5) (1, 3)\par
6) (2, 3)\par
{\vskip 15pt} Номери відповідних ребер записати за зростанням через пробіл.


%enditem


\newpage
16. 

Рядки журналу через двокрапку містять ім'я користувача (ідентифікатор), час події,тип події, код події, наприклад
\begin{verbatim}
s="username : 2022-05-05 13:26 : ERROR : 404"
\end{verbatim}
у коді 
\begin{verbatim}
res = re.fullmatch('???', s) 
s = ''
code_str = ??? 
\end{verbatim}
чим треба замінити кожен з ???, щоб значенням code\_str був код події? 


\vspace{10mm}
\textbf{Завдання 17} 

Задано програму, що містить код 1-2 класів, можливе успадкування, та клієнтський код.
У наведеному коді невеликі частини закриті. 
Також наведено вихід програми.
Необхідно відновити закриті частини, щоб отриманий код відповідав наведеному виходу програми.


\textbf{Завдання 18-19} 

Кілька запитань за завданнями для самостійної роботи.
\vspace{10mm}
Скоріш за все, що завдань буде трохи менше, а деякі попередні будуть переформульовані в стилі двох останніх.

%enditem

\end {large}}

%end

\newpage
%begin

{{
\begin {large}

\newpage
1. Що буде надруковано за виконання?

\begin{verbatim}
print('R', end=' ')
e=('a','b','c','d','e',0,1)
print(e[-3])
\end{verbatim}

%enditem


\newpage
2. Що буде надруковано за виконання?

\begin{verbatim}
print('R', end=' ')
e=('0','1','2','3','4','5','6','7','8','9')
print(e[:4:-3])
\end{verbatim}

%enditem



\newpage
3. Що буде надруковано за виконання?

\begin{verbatim}
print('R', end=' ')
e=[10,11,12,13,14]
e.insert(-7, '13')
print(e)
\end{verbatim}

%enditem



\newpage
4. Що буде надруковано за виконання?

\begin{verbatim}
print('R', end=' ')
a=[10,11,12,13,14]
a[1:]=(14)
print(a)
\end{verbatim}

%enditem



\newpage
5. Укажіть ТИП значення виразу.\par
\begin{verbatim}
{x for x in range(7) if el%3==1}
\end{verbatim}
%enditem


\newpage
6. Що буде надруковано за виконання?\par
\begin{verbatim}
try:
    s = {89:2, 53:8, 80:8, 53:7}
    s[30] = 2
    for x in s.values():
        print(x, sep=' ', end=' ')
except:
    print('ERR')
\end{verbatim}
%enditem


\newpage
7. 
Записати ВИРАЗ, значенням якого є список, що складається з кубів цілих чисел від 16 до 2012 (включно), що не діляться на 5.
%enditem


\newpage
8. Що буде надруковано за виконання?
%Оригинал был утерян. Требует отладки!
\begin{verbatim}
a = 0

class A:
    a = 1
    
    def __init__(self):
        self.a = 2
        a = 3
        A.a = 4

obj = A()
A.a = 5
try:
    print(a, end=' ')
    print(obj.a, end=' ')
    print(A.a, end=' ')
    print(obj.a, end=' ')
    print(obj._A__a)
except:
    print('ERR')
\end{verbatim}
%enditem


\newpage
9. Що буде надруковано за виконання?
\begin{verbatim}
class C0:
    def b(self): print(1, end=' ')
    def f(self): print(2, end=' ')

class C1(C0):
    def f(self): print(3, end=' ')

try:
    C1().b()
    C1().f()
except:
    print('ERR')
\end{verbatim}
%enditem


\newpage
10. Що буде надруковано за виконання?
\begin{verbatim}
class A:
    c = 1

    def __init__(self):
        b = 2
        c = 3

class B(A):
    b = 4
    a = 5

    def __init__(self):
        super().__init__()
        self.a = 6

try:
    obj = B()
    obj.b = obj.c + 10
    B.b = 13
except: print('EE', end=' ')

try: print(obj.b, end= ' ')
except: print('E1', end=' ')

try: print(A.b, end= ' ')
except: print('E2', end=' ')

try: print(B.b, end= ' ')
except: print('E3', end=' ')
\end{verbatim}
%enditem


\newpage
11. Змінні \kw{next} та \kw{prev} вузла списку позначають наступний та попередній за ним вузол відповідно. Починаючи з голови, у список записано послідовні цілі числа від~$-22$ до~$34$. Нехай змінна \kw{p} позначає вузол, що зберігає значення~$-6$.
Яке значення зберігається у вузлі \kw{p.prev.next.prev.prev.prev} ?
%enditem


\newpage
12. 
\begin{center}
	\adjustimage{max size={0.8\linewidth}{0.8\paperheight}}
	{BinTree/trees_exp/48.png}
\end{center}
Указати послідовність міток вершин дерева, зображеного на рис., що утвориться в результаті обходу дерева в оберненому порядку. Мітки записувати через пробіл.\par Алгоритм починає роботу з кореня (на рис. це верхня вершина).
%enditem


\newpage
13. 
\begin{center}
	\adjustimage{max size={0.9\linewidth}{0.9\paperheight}}
	{Graphs/AdjList/289.png}
\end{center}
Для графа, зображеного на рис., вказати список суміжності вершини 1.\par
Формат оформлення відповіді:\par
список суміжності виписувати через пробіл на одному рядку, якщо пробілів десь буде більше одного - не важливо, упорядкування вершин у списку також значення не має.
%enditem


\newpage
14. Для графа на рисунку з вершини 4 запущено алгоритм пошуку в глибину, що будує кістякове дерево. Списки суміжності вершин переглядаються алгоритмом за зростанням міток вершин.\par
\begin{center}
	\adjustimage{max size={0.9\linewidth}{0.9\paperheight}}
	{Graphs/Traversals/39.png}
\end{center}
\par Які з наведених ребер входять до складу отриманого кістякового дерева?\par
1) (1, 4)\par
2) (2, 6)\par
3) (1, 6)\par
4) (0, 4)\par
5) (2, 5)\par
6) (2, 4)\par
{\vskip 15pt} Номери відповідних ребер записати за зростанням через пробіл.

%enditem


\newpage
15. Для графа на рисунку з вершини 6 запущено алгоритм пошуку в ширину, що будує кістякове дерево. Списки суміжності вершин переглядаються алгоритмом за зростанням міток вершин.\par
\begin{center}
	\adjustimage{max size={0.9\linewidth}{0.9\paperheight}}
	{Graphs/Traversals/42.png}
\end{center}
\par Які з наведених ребер входять до складу отриманого кістякового дерева?\par
1) (0, 2)\par
2) (1, 6)\par
3) (4, 5)\par
4) (3, 4)\par
5) (2, 5)\par
6) (0, 6)\par
{\vskip 15pt} Номери відповідних ребер записати за зростанням через пробіл.


%enditem


\newpage
16. 

Рядки журналу через двокрапку містять ім'я користувача (ідентифікатор), час події,тип події, код події, наприклад
\begin{verbatim}
s="username : 2022-05-05 13:26 : ERROR : 404"
\end{verbatim}
у коді 
\begin{verbatim}
res = re.fullmatch('???', s) 
s = ''
code_str = ??? 
\end{verbatim}
чим треба замінити кожен з ???, щоб значенням code\_str був код події? 


\vspace{10mm}
\textbf{Завдання 17} 

Задано програму, що містить код 1-2 класів, можливе успадкування, та клієнтський код.
У наведеному коді невеликі частини закриті. 
Також наведено вихід програми.
Необхідно відновити закриті частини, щоб отриманий код відповідав наведеному виходу програми.


\textbf{Завдання 18-19} 

Кілька запитань за завданнями для самостійної роботи.
\vspace{10mm}
Скоріш за все, що завдань буде трохи менше, а деякі попередні будуть переформульовані в стилі двох останніх.

%enditem

\end {large}}

%end

\newpage
%begin

{{
\begin {large}

\newpage
1. Що буде надруковано за виконання?

\begin{verbatim}
print('R', end=' ')
c=('a','b','c','d','e',0,1,2,4)
print(c[-13])
\end{verbatim}

%enditem


\newpage
2. Що буде надруковано за виконання?

\begin{verbatim}
print('R', end=' ')
f=(0,1,2,3,4,5,6,7,8,9)
print(f[-14:-2:-2])
\end{verbatim}

%enditem



\newpage
3. Що буде надруковано за виконання?

\begin{verbatim}
print('R', end=' ')
g=['0','1','2','3','4']
print(g.pop(), end=' ')
print(g)
\end{verbatim}

%enditem



\newpage
4. Що буде надруковано за виконання?

\begin{verbatim}
print('R', end=' ')
h=['0','1','2','3','4']
del h[-2:-3]
print(h)
\end{verbatim}

%enditem



\newpage
5. Укажіть ТИП значення виразу.\par
\begin{verbatim}
{y:None for y in range(-8,3)}
\end{verbatim}
%enditem


\newpage
6. Що буде надруковано за виконання?\par
\begin{verbatim}
try:
    d = {19:8, 35:6, 57:0, 32:4, 19:3}
    d[32] = 4
    for x in d.keys():
        print(x, sep=' ', end=' ')
except:
    print('ERR')
\end{verbatim}
%enditem


\newpage
7. Записати ВИРАЗ, значенням якого є множина, що складається з кубів цілих чисел від 17 до 2006 (включно), що не діляться на 5.
%enditem


\newpage
8. Що буде надруковано за виконання?
\begin{verbatim}
c = 0
h = 1

class D:
    c = 2
    
    def __init__(self):
        self.h = 3
        c = 4
        D.c = 5

obj = D()
try:
    print(c, end=' ')
    print(h, end=' ')
    print(D.c, end=' ')
    print(obj.h)
except:
    print('ERR')
\end{verbatim}
%enditem


\newpage
9. Що буде надруковано за виконання?
\begin{verbatim}
class A0:
    def c(self): print(1, end=' ')
    def h(self): print(2, end=' ')

class A1(A0):
    def h(self): print(3, end=' ')

try:
    A1().c()
    A1().h()
except:
    print('ERR')
\end{verbatim}
%enditem


\newpage
10. Що буде надруковано за виконання?
\begin{verbatim}
class A:
    b = 1

    def __init__(self):
        a = 2
        a = 3

class B(A):
    b = 4
    c = 5

    def __init__(self):
        super().__init__()
        self.c = 6

try:
    obj = B()
    obj.a = obj.b + 10
    B.a = 13
except: print('EE', end=' ')

try: print(obj.a, end= ' ')
except: print('E1', end=' ')

try: print(A.a, end= ' ')
except: print('E2', end=' ')

try: print(B.a, end= ' ')
except: print('E3', end=' ')
\end{verbatim}
%enditem


\newpage
11. Змінні \kw{next} та \kw{prev} вузла списку позначають наступний та попередній за ним вузол відповідно. Починаючи з голови, у список записано послідовні цілі числа від~$-16$ до~$26$. Нехай змінна \kw{p} позначає вузол, що зберігає значення~$-4$.
Яке значення зберігається у вузлі \kw{p.next.prev.next.next.next} ?
%enditem


\newpage
12. 
\begin{center}
	\adjustimage{max size={0.8\linewidth}{0.8\paperheight}}
	{BinTree/trees_exp/136.png}
\end{center}
Указати послідовність міток вершин дерева, зображеного на рис., що утвориться в результаті обходу дерева в прямому порядку. Мітки записувати через пробіл.\par Алгоритм починає роботу з кореня (на рис. це верхня вершина).
%enditem


\newpage
13. 
\begin{center}
	\adjustimage{max size={0.9\linewidth}{0.9\paperheight}}
	{Graphs/AdjList/215.png}
\end{center}
Для графа, зображеного на рис., вказати список суміжності вершини 7.\par
Формат оформлення відповіді:\par
список суміжності виписувати через пробіл на одному рядку, якщо пробілів десь буде більше одного - не важливо, упорядкування вершин у списку також значення не має.
%enditem


\newpage
14. Для графа на рисунку вказати послідовність, в якій алгоритм пошуку в глибину буде відмічати відвідані вершини, якщо списки суміжності вершин переглядаються алгоритмом за зростанням міток вершин та вершини відмічаються на першому вход\-женні у вершину.\par
Пошук розпочинається з вершини 1.\par

\begin{center}
	\adjustimage{max size={0.9\linewidth}{0.9\paperheight}}
	{Graphs/Traversals/287.png}
\end{center}
%enditem


\newpage
15. Для графа на рисунку з вершини 3 запущено алгоритм пошуку в ширину, що будує кістякове дерево. Списки суміжності вершин переглядаються алгоритмом за зростанням міток вершин.\par
\begin{center}
	\adjustimage{max size={0.9\linewidth}{0.9\paperheight}}
	{Graphs/Traversals/328.png}
\end{center}
\par Які з наведених ребер входять до складу отриманого кістякового дерева?\par
1) (0, 3)\par
2) (1, 6)\par
3) (2, 6)\par
4) (4, 6)\par
5) (3, 4)\par
6) (3, 6)\par
{\vskip 15pt} Номери відповідних ребер записати за зростанням через пробіл.


%enditem


\newpage
16. 

Рядки журналу через двокрапку містять ім'я користувача (ідентифікатор), час події,тип події, код події, наприклад
\begin{verbatim}
s="username : 2022-05-05 13:26 : ERROR : 404"
\end{verbatim}
у коді 
\begin{verbatim}
res = re.fullmatch('???', s) 
s = ''
code_str = ??? 
\end{verbatim}
чим треба замінити кожен з ???, щоб значенням code\_str був код події? 


\vspace{10mm}
\textbf{Завдання 17} 

Задано програму, що містить код 1-2 класів, можливе успадкування, та клієнтський код.
У наведеному коді невеликі частини закриті. 
Також наведено вихід програми.
Необхідно відновити закриті частини, щоб отриманий код відповідав наведеному виходу програми.


\textbf{Завдання 18-19} 

Кілька запитань за завданнями для самостійної роботи.
\vspace{10mm}
Скоріш за все, що завдань буде трохи менше, а деякі попередні будуть переформульовані в стилі двох останніх.

%enditem

\end {large}}

%end

\newpage
%begin

{{
\begin {large}

\newpage
1. Що буде надруковано за виконання?

\begin{verbatim}
print('R', end=' ')
h=('a','b','c','d','e',0,1,2,3,4)
print(h[-8])
\end{verbatim}

%enditem


\newpage
2. Що буде надруковано за виконання?

\begin{verbatim}
print('R', end=' ')
c=['0','1','2','3','4','5','6','7','8','9']
print(c[7::3])
\end{verbatim}

%enditem



\newpage
3. Що буде надруковано за виконання?

\begin{verbatim}
print('R', end=' ')
d=['0','1','2','3','4']
d.extend([4,5])
print(d)
\end{verbatim}

%enditem



\newpage
4. Що буде надруковано за виконання?

\begin{verbatim}
print('R', end=' ')
c=[10,11,12,13,14]
c[2:2]='12'
print(c)
\end{verbatim}

%enditem



\newpage
5. Укажіть ТИП значення виразу.\par
\begin{verbatim}
{y:y*y for y in range(-7,7)}
\end{verbatim}
%enditem


\newpage
6. Що буде надруковано за виконання?\par
\begin{verbatim}
try:
    d = {36:2, 85:1, 68:1, 85:6}
    d[36] = 2
    for x in d.values():
        print(x, sep=' ', end=' ')
except:
    print('ERR')
\end{verbatim}
%enditem


\newpage
7. Нехай змінна \kw{s} позначає дуже довгу послідовність цілих, по якій можна ітерувати. 

Написати вираз-генератор, що задає підпослідовність її елементів, що менші за задане число num. 

Обмеження: усі елементи шуканої підпослідовності одночасно в пам'яті розміщені бути не можуть. 
%enditem


\newpage
8. Що буде надруковано за виконання?
%Оригинал был утерян. Требует отладки!
\begin{verbatim}
__d = 0

class B:
    __d = 1
    
    def __init__(self):
        self.__d = 2
        __d = 3
        B.__d = 4

obj = B()
obj.__d = 5
try:
    print(__d, end=' ')
    print(obj.__d, end=' ')
    print(B.__d, end=' ')
    print(obj.__d, end=' ')
    print(obj._B__d)
except:
    print('ERR')
\end{verbatim}
%enditem


\newpage
9. Що буде надруковано за виконання?
\begin{verbatim}
class D0:
    def a(self): print(1, end=' ')
    def h(self): print(2, end=' ')

class D1(D0):
    def h(self): print(3, end=' ')

try:
    D0().a()
    D1().h()
except:
    print('ERR')
\end{verbatim}
%enditem


\newpage
10. Що буде надруковано за виконання?
\begin{verbatim}
class A:
    c = 1

    def _f1(self): print(2, end=' ')
    def _f2(self): print(3, end=' ')

    def main(self):
        try: print(self.c, end=' ')
        except: print('E1', end=' ')

        try: self._f1()
        except: print('E2', end=' ')

        try: self._f2()
        except: print('E3', end=' ')


class B(A):
    c = 4

    def _f1(self): print(5, end=' ')
    def _f2(self): print(6, end=' ')

try: B().main()
except: print('E4', end=' ')
\end{verbatim}
%enditem


\newpage
11. Змінні \kw{next} та \kw{prev} вузла списку позначають наступний та попередній за ним вузол відповідно. Починаючи з голови, у список записано послідовні цілі числа від~$-15$ до~$17$. Нехай змінна \kw{p} позначає вузол, що зберігає значення~$-5$.
Яке значення зберігається у вузлі \kw{p.prev.next.next.prev.next} ?
%enditem


\newpage
12. 
\begin{center}
	\adjustimage{max size={0.8\linewidth}{0.8\paperheight}}
	{BinTree/trees_exp/69.png}
\end{center}
Указати послідовність міток вершин дерева, зображеного на рис., що утвориться в результаті обходу дерева в оберненому порядку. Мітки записувати через пробіл.\par Алгоритм починає роботу з кореня (на рис. це верхня вершина).
%enditem


\newpage
13. 
\begin{center}
	\adjustimage{max size={0.9\linewidth}{0.9\paperheight}}
	{Graphs/AdjList/42.png}
\end{center}
Для графа, зображеного на рис., вказати список суміжності вершини 7.\par
Формат оформлення відповіді:\par
список суміжності виписувати через пробіл на одному рядку, якщо пробілів десь буде більше одного - не важливо, упорядкування вершин у списку також значення не має.
%enditem


\newpage
14. Для графа на рисунку вказати послідовність, в якій алгоритм пошуку в глибину буде відмічати відвідані вершини, якщо списки суміжності вершин переглядаються алгоритмом за зростанням міток вершин та вершини відмічаються на першому вход\-женні у вершину.\par
Пошук розпочинається з вершини 4.\par

\begin{center}
	\adjustimage{max size={0.9\linewidth}{0.9\paperheight}}
	{Graphs/Traversals/131.png}
\end{center}
%enditem


\newpage
15. Для графа на рисунку з вершини 3 запущено алгоритм пошуку в ширину, що будує кістякове дерево. Списки суміжності вершин переглядаються алгоритмом за зростанням міток вершин.\par
\begin{center}
	\adjustimage{max size={0.9\linewidth}{0.9\paperheight}}
	{Graphs/Traversals/319.png}
\end{center}
\par Які з наведених ребер входять до складу отриманого кістякового дерева?\par
1) (3, 5)\par
2) (3, 6)\par
3) (1, 4)\par
4) (0, 6)\par
5) (1, 2)\par
6) (1, 6)\par
{\vskip 15pt} Номери відповідних ребер записати за зростанням через пробіл.


%enditem


\newpage
16. 

Рядки журналу через двокрапку містять ім'я користувача (ідентифікатор), час події,тип події, код події, наприклад
\begin{verbatim}
s="username : 2022-05-05 13:26 : ERROR : 404"
\end{verbatim}
у коді 
\begin{verbatim}
res = re.fullmatch('???', s) 
s = ''
code_str = ??? 
\end{verbatim}
чим треба замінити кожен з ???, щоб значенням code\_str був код події? 


\vspace{10mm}
\textbf{Завдання 17} 

Задано програму, що містить код 1-2 класів, можливе успадкування, та клієнтський код.
У наведеному коді невеликі частини закриті. 
Також наведено вихід програми.
Необхідно відновити закриті частини, щоб отриманий код відповідав наведеному виходу програми.


\textbf{Завдання 18-19} 

Кілька запитань за завданнями для самостійної роботи.
\vspace{10mm}
Скоріш за все, що завдань буде трохи менше, а деякі попередні будуть переформульовані в стилі двох останніх.

%enditem

\end {large}}

%end

\newpage
%begin

{{
\begin {large}

\newpage
1. Що буде надруковано за виконання?

\begin{verbatim}
print('R', end=' ')
b=('a','b','c','d','e',0,1,2,4)
print(b[6])
\end{verbatim}

%enditem


\newpage
2. Що буде надруковано за виконання?

\begin{verbatim}
print('R', end=' ')
a=['0','1','2','3','4','5','6','7','8','9']
print(a[::2])
\end{verbatim}

%enditem



\newpage
3. Що буде надруковано за виконання?

\begin{verbatim}
print('R', end=' ')
c=[10,11,12,13,14]
c.remove(14)
print(c)
\end{verbatim}

%enditem



\newpage
4. Що буде надруковано за виконання?

\begin{verbatim}
print('R', end=' ')
b=['0','1','2','3','4']
b[:7]=(13)
print(b)
\end{verbatim}

%enditem



\newpage
5. Укажіть ТИП значення виразу.\par
\begin{verbatim}
{1:3, 5:8}
\end{verbatim}
%enditem


\newpage
6. Що буде надруковано за виконання?\par
\begin{verbatim}
try:
    t = {28:3, 34:5, 37:9, 37:2}
    t[66] = 1
    for x, y in t.items():
        print(x, y, sep=' ', end=' ')
except:
    print('ERR')
\end{verbatim}
%enditem


\newpage
7. Задано множину \kw{x}, що складається з цілих чисел. Записати ВИРАЗ, значенням якого є множина, що складається  з непарних елементів множини \kw{x}.
%enditem


\newpage
8. Що буде надруковано за виконання?
\begin{verbatim}
d = 0
m = 1

class A:
    m = 2
    
    def __init__(self):
        global m
        self.d = 3
        m = 4
        A.m = 5

obj = A()
try:
    print(d, end=' ')
    print(m, end=' ')
    print(A.m, end=' ')
    print(obj.d)
except:
    print('ERR')
\end{verbatim}
%enditem


\newpage
9. Що буде надруковано за виконання?
\begin{verbatim}
class B0:
    def d(self): print(1, end=' ')
    def f(self): print(2, end=' ')

class B1(B0):
    def d(self): print(3, end=' ')

try:
    B1().f()
    B0().d()
except:
    print('ERR')
\end{verbatim}
%enditem


\newpage
10. Що буде надруковано за виконання?
\begin{verbatim}
class A:
    d = 1

    def _f1(self): print(2, end=' ')
    def f2(self): print(3, end=' ')

    def main(self):
        try: print(self.d, end=' ')
        except: print('E1', end=' ')

        try: self._f1()
        except: print('E2', end=' ')

        try: self.f2()
        except: print('E3', end=' ')


class B(A):
    d = 4

    def _f1(self): print(5, end=' ')
    def f2(self): print(6, end=' ')

try: B().main()
except: print('E4', end=' ')
\end{verbatim}
%enditem


\newpage
11. Змінні \kw{next} та \kw{prev} вузла списку позначають наступний та попередній за ним вузол відповідно. Починаючи з голови, у список записано послідовні цілі числа від~$-21$ до~$30$. Нехай змінна \kw{p} позначає вузол, що зберігає значення~$0$.
Яке значення зберігається у вузлі \kw{p.next.prev.prev.next.prev} ?
%enditem


\newpage
12. 
\begin{center}
	\adjustimage{max size={0.8\linewidth}{0.8\paperheight}}
	{BinTree/trees_exp/113.png}
\end{center}
Указати послідовність міток вершин дерева, зображеного на рис., що утвориться в результаті обходу дерева в прямому порядку. Мітки записувати через пробіл.\par Алгоритм починає роботу з кореня (на рис. це верхня вершина).
%enditem


\newpage
13. 
\begin{center}
	\adjustimage{max size={0.9\linewidth}{0.9\paperheight}}
	{Graphs/AdjList/60.png}
\end{center}
Для графа, зображеного на рис., вказати список суміжності вершини 5.\par
Формат оформлення відповіді:\par
список суміжності виписувати через пробіл на одному рядку, якщо пробілів десь буде більше одного - не важливо, упорядкування вершин у списку також значення не має.
%enditem


\newpage
14. Для графа на рисунку вказати послідовність, в якій алгоритм пошуку в глибину буде відмічати відвідані вершини, якщо списки суміжності вершин переглядаються алгоритмом за зростанням міток вершин та вершини відмічаються на першому вход\-женні у вершину.\par
Пошук розпочинається з вершини 1.\par

\begin{center}
	\adjustimage{max size={0.9\linewidth}{0.9\paperheight}}
	{Graphs/Traversals/204.png}
\end{center}
%enditem


\newpage
15. Для графа на рисунку з вершини 3 запущено алгоритм пошуку в ширину, що будує кістякове дерево. Списки суміжності вершин переглядаються алгоритмом за зростанням міток вершин.\par
\begin{center}
	\adjustimage{max size={0.9\linewidth}{0.9\paperheight}}
	{Graphs/Traversals/136.png}
\end{center}
\par Які з наведених ребер входять до складу отриманого кістякового дерева?\par
1) (2, 3)\par
2) (0, 4)\par
3) (2, 5)\par
4) (3, 6)\par
5) (1, 2)\par
6) (0, 5)\par
{\vskip 15pt} Номери відповідних ребер записати за зростанням через пробіл.


%enditem


\newpage
16. 

Рядки журналу через двокрапку містять ім'я користувача (ідентифікатор), час події,тип події, код події, наприклад
\begin{verbatim}
s="username : 2022-05-05 13:26 : ERROR : 404"
\end{verbatim}
у коді 
\begin{verbatim}
res = re.fullmatch('???', s) 
s = ''
code_str = ??? 
\end{verbatim}
чим треба замінити кожен з ???, щоб значенням code\_str був код події? 


\vspace{10mm}
\textbf{Завдання 17} 

Задано програму, що містить код 1-2 класів, можливе успадкування, та клієнтський код.
У наведеному коді невеликі частини закриті. 
Також наведено вихід програми.
Необхідно відновити закриті частини, щоб отриманий код відповідав наведеному виходу програми.


\textbf{Завдання 18-19} 

Кілька запитань за завданнями для самостійної роботи.
\vspace{10mm}
Скоріш за все, що завдань буде трохи менше, а деякі попередні будуть переформульовані в стилі двох останніх.

%enditem

\end {large}}

%end

\newpage
%begin

{{
\begin {large}

\newpage
1. Що буде надруковано за виконання?

\begin{verbatim}
print('R', end=' ')
f=('a','b','c','d','e',0,1,2)
print(f[5])
\end{verbatim}

%enditem


\newpage
2. Що буде надруковано за виконання?

\begin{verbatim}
print('R', end=' ')
b=('0','1','2','3','4','5','6','7','8','9')
print(b[-7::-1])
\end{verbatim}

%enditem



\newpage
3. Що буде надруковано за виконання?

\begin{verbatim}
print('R', end=' ')
f=['0','1','2','3','4']
f+=[2,3]
print(f)
\end{verbatim}

%enditem



\newpage
4. Що буде надруковано за виконання?

\begin{verbatim}
print('R', end=' ')
f=['0','1','2','3','4']
del f[7:-7]
print(f)
\end{verbatim}

%enditem



\newpage
5. Укажіть ТИП значення виразу.\par
\begin{verbatim}
[x for x in range(7)]
\end{verbatim}
%enditem


\newpage
6. Що буде надруковано за виконання?\par
\begin{verbatim}
try:
    s = {54:1, 33:1, 69:5, 33:1}
    s[33] = 0
    for x in s :
        print(x, sep=' ', end=' ')
except:
    print('ERR')
\end{verbatim}
%enditem


\newpage
7. Нехай змінна \kw{s} позначає послідовність цілих, по якій можна ітерувати.

Написати вираз, що перевіряє, чи правда, що всі елементи послідовності діляться на 11.
%enditem


\newpage
8. Що буде надруковано за виконання?
%Оригинал был утерян. Требует отладки!
\begin{verbatim}
b = 0

class C:
    b = 1
    
    def __init__(self):
        self.b = 2
        b = 3
        C.b = 4

obj = C()
obj.b = 5
try:
    print(b, end=' ')
    print(obj.b, end=' ')
    print(C.b, end=' ')
    print(obj.b, end=' ')
    print(obj._C__b)
except:
    print('ERR')
\end{verbatim}
%enditem


\newpage
9. Що буде надруковано за виконання?
\begin{verbatim}
class A0:
    def a(self): print(1, end=' ')
    def g(self): print(2, end=' ')

class A1(A0):
    def g(self): print(3, end=' ')

try:
    A1().a()
    A1().g()
except:
    print('ERR')
\end{verbatim}
%enditem


\newpage
10. Що буде надруковано за виконання?
\begin{verbatim}
class A:
    a = 1

    def __init__(self):
        c = 2
        c = 3

class B(A):
    a = 4
    b = 5

    def __init__(self):
        super().__init__()
        self.b = 6

try:
    obj = B()
    obj.c = obj.c + 10
    B.c = 13
except: print('EE', end=' ')

try: print(obj.c, end= ' ')
except: print('E1', end=' ')

try: print(A.c, end= ' ')
except: print('E2', end=' ')

try: print(B.c, end= ' ')
except: print('E3', end=' ')
\end{verbatim}
%enditem


\newpage
11. Змінні \kw{next} та \kw{prev} вузла списку позначають наступний та попередній за ним вузол відповідно. Починаючи з голови, у список записано послідовні цілі числа від~$-48$ до~$20$. Нехай змінна \kw{p} позначає вузол, що зберігає значення~$2$.
Яке значення зберігається у вузлі \kw{p.prev.prev.next.next.next} ?
%enditem


\newpage
12. 
\begin{center}
	\adjustimage{max size={0.8\linewidth}{0.8\paperheight}}
	{BinTree/trees_exp/46.png}
\end{center}
Указати послідовність міток вершин дерева, зображеного на рис., що утвориться в результаті обходу дерева в оберненому порядку. Мітки записувати через пробіл.\par Алгоритм починає роботу з кореня (на рис. це верхня вершина).
%enditem


\newpage
13. 
\begin{center}
	\adjustimage{max size={0.9\linewidth}{0.9\paperheight}}
	{Graphs/AdjList/32.png}
\end{center}
Для графа, зображеного на рис., вказати список суміжності вершини 0.\par
Формат оформлення відповіді:\par
список суміжності виписувати через пробіл на одному рядку, якщо пробілів десь буде більше одного - не важливо, упорядкування вершин у списку також значення не має.
%enditem


\newpage
14. Для графа на рисунку з вершини 6 запущено алгоритм пошуку в глибину, що будує кістякове дерево. Списки суміжності вершин переглядаються алгоритмом за зростанням міток вершин.\par
\begin{center}
	\adjustimage{max size={0.9\linewidth}{0.9\paperheight}}
	{Graphs/Traversals/269.png}
\end{center}
\par Які з наведених ребер входять до складу отриманого кістякового дерева?\par
1) (5, 6)\par
2) (1, 6)\par
3) (0, 2)\par
4) (0, 4)\par
5) (2, 3)\par
6) (3, 6)\par
{\vskip 15pt} Номери відповідних ребер записати за зростанням через пробіл.

%enditem


\newpage
15. Для графа на рисунку з вершини 2 запущено алгоритм пошуку в ширину, що будує кістякове дерево. Списки суміжності вершин переглядаються алгоритмом за зростанням міток вершин.\par
\begin{center}
	\adjustimage{max size={0.9\linewidth}{0.9\paperheight}}
	{Graphs/Traversals/17.png}
\end{center}
\par Які з наведених ребер входять до складу отриманого кістякового дерева?\par
1) (2, 5)\par
2) (0, 3)\par
3) (2, 4)\par
4) (1, 6)\par
5) (2, 3)\par
6) (2, 6)\par
{\vskip 15pt} Номери відповідних ребер записати за зростанням через пробіл.


%enditem


\newpage
16. 

Рядки журналу через двокрапку містять ім'я користувача (ідентифікатор), час події,тип події, код події, наприклад
\begin{verbatim}
s="username : 2022-05-05 13:26 : ERROR : 404"
\end{verbatim}
у коді 
\begin{verbatim}
res = re.fullmatch('???', s) 
s = ''
code_str = ??? 
\end{verbatim}
чим треба замінити кожен з ???, щоб значенням code\_str був код події? 


\vspace{10mm}
\textbf{Завдання 17} 

Задано програму, що містить код 1-2 класів, можливе успадкування, та клієнтський код.
У наведеному коді невеликі частини закриті. 
Також наведено вихід програми.
Необхідно відновити закриті частини, щоб отриманий код відповідав наведеному виходу програми.


\textbf{Завдання 18-19} 

Кілька запитань за завданнями для самостійної роботи.
\vspace{10mm}
Скоріш за все, що завдань буде трохи менше, а деякі попередні будуть переформульовані в стилі двох останніх.

%enditem

\end {large}}

%end

\newpage
%begin

{{
\begin {large}

\newpage
1. Що буде надруковано за виконання?

\begin{verbatim}
print('R', end=' ')
d=('a','b','c','d','e',0,1,2,3)
print(d[7])
\end{verbatim}

%enditem


\newpage
2. Що буде надруковано за виконання?

\begin{verbatim}
print('R', end=' ')
g=[0,1,2,3,4,5,6,7,8,9]
print(g[-12:-9:-2])
\end{verbatim}

%enditem



\newpage
3. Що буде надруковано за виконання?

\begin{verbatim}
print('R', end=' ')
h=[10,11,12,13,14]
h.append('12')
print(h)
\end{verbatim}

%enditem



\newpage
4. Що буде надруковано за виконання?

\begin{verbatim}
print('R', end=' ')
d=[10,11,12,13,14]
del d[8:7]
print(d)
\end{verbatim}

%enditem



\newpage
5. Укажіть ТИП значення виразу.\par
\begin{verbatim}
{x:0 for x in range(7)}
\end{verbatim}
%enditem


\newpage
6. Що буде надруковано за виконання?\par
\begin{verbatim}
try:
    t = {14:5, 34:3, 84:8, 31:2}
    t[14] = 8
    for x in t.items():
        print(*x, sep=' ', end=' ')
except:
    print('ERR')
\end{verbatim}
%enditem


\newpage
7. Нехай змінна \kw{s} позначає дуже довгу послідовність цілих, по якій можна ітерувати. 

Написати вираз-генератор, що задає підпослідовність її елементів, що більші за задане число num. 

Обмеження: усі елементи шуканої підпослідовності одночасно в пам'яті розміщені бути не можуть. 
%enditem


\newpage
8. Що буде надруковано за виконання?
\begin{verbatim}
a = 0
n = 1

class C:
    a = 2
    
    def __init__(self):
        self.n = 3
        a = 4
        C.a = 5

obj = C()
try:
    print(a, end=' ')
    print(n, end=' ')
    print(C.a, end=' ')
    print(obj.n)
except:
    print('ERR')
\end{verbatim}
%enditem


\newpage
9. Що буде надруковано за виконання?
\begin{verbatim}
class D0:
    def c(self): print(1, end=' ')
    def h(self): print(2, end=' ')

class D1(D0):
    def h(self): print(3, end=' ')

try:
    D0().c()
    D1().h()
except:
    print('ERR')
\end{verbatim}
%enditem


\newpage
10. Що буде надруковано за виконання?
\begin{verbatim}
class A:
    b = 1

    def _h1(self): print(2, end=' ')
    def __h2(self): print(3, end=' ')

    def main(self):
        try: print(self.b, end=' ')
        except: print('E1', end=' ')

        try: self._h1()
        except: print('E2', end=' ')

        try: self.__h2()
        except: print('E3', end=' ')


class B(A):
    b = 4

    def _h1(self): print(5, end=' ')
    def __h2(self): print(6, end=' ')

try: B().main()
except: print('E4', end=' ')
\end{verbatim}
%enditem


\newpage
11. Змінні \kw{next} та \kw{prev} вузла списку позначають наступний та попередній за ним вузол відповідно. Починаючи з голови, у список записано послідовні цілі числа від~$-43$ до~$43$. Нехай змінна \kw{p} позначає вузол, що зберігає значення~$-4$.
Яке значення зберігається у вузлі \kw{p.next.next.prev.prev.prev} ?
%enditem


\newpage
12. 
\begin{center}
	\adjustimage{max size={0.8\linewidth}{0.8\paperheight}}
	{BinTree/trees_exp/172.png}
\end{center}
Указати послідовність міток вершин дерева, зображеного на рис., що утвориться в результаті обходу дерева в прямому порядку. Мітки записувати через пробіл.\par Алгоритм починає роботу з кореня (на рис. це верхня вершина).
%enditem


\newpage
13. 
\begin{center}
	\adjustimage{max size={0.9\linewidth}{0.9\paperheight}}
	{Graphs/AdjList/141.png}
\end{center}
Для графа, зображеного на рис., вказати список суміжності вершини 5.\par
Формат оформлення відповіді:\par
список суміжності виписувати через пробіл на одному рядку, якщо пробілів десь буде більше одного - не важливо, упорядкування вершин у списку також значення не має.
%enditem


\newpage
14. Для графа на рисунку з вершини 0 запущено алгоритм пошуку в глибину, що будує кістякове дерево. Списки суміжності вершин переглядаються алгоритмом за зростанням міток вершин.\par
\begin{center}
	\adjustimage{max size={0.9\linewidth}{0.9\paperheight}}
	{Graphs/Traversals/13.png}
\end{center}
\par Які з наведених ребер входять до складу отриманого кістякового дерева?\par
1) (4, 6)\par
2) (2, 6)\par
3) (1, 3)\par
4) (1, 6)\par
5) (3, 4)\par
6) (0, 3)\par
{\vskip 15pt} Номери відповідних ребер записати за зростанням через пробіл.

%enditem


\newpage
15. Для графа на рисунку з вершини 1 запущено алгоритм пошуку в ширину, що будує кістякове дерево. Списки суміжності вершин переглядаються алгоритмом за зростанням міток вершин.\par
\begin{center}
	\adjustimage{max size={0.9\linewidth}{0.9\paperheight}}
	{Graphs/Traversals/238.png}
\end{center}
\par Які з наведених ребер входять до складу отриманого кістякового дерева?\par
1) (1, 4)\par
2) (3, 6)\par
3) (0, 1)\par
4) (1, 6)\par
5) (0, 2)\par
6) (1, 5)\par
{\vskip 15pt} Номери відповідних ребер записати за зростанням через пробіл.


%enditem


\newpage
16. 

Рядки журналу через двокрапку містять ім'я користувача (ідентифікатор), час події,тип події, код події, наприклад
\begin{verbatim}
s="username : 2022-05-05 13:26 : ERROR : 404"
\end{verbatim}
у коді 
\begin{verbatim}
res = re.fullmatch('???', s) 
s = ''
code_str = ??? 
\end{verbatim}
чим треба замінити кожен з ???, щоб значенням code\_str був код події? 


\vspace{10mm}
\textbf{Завдання 17} 

Задано програму, що містить код 1-2 класів, можливе успадкування, та клієнтський код.
У наведеному коді невеликі частини закриті. 
Також наведено вихід програми.
Необхідно відновити закриті частини, щоб отриманий код відповідав наведеному виходу програми.


\textbf{Завдання 18-19} 

Кілька запитань за завданнями для самостійної роботи.
\vspace{10mm}
Скоріш за все, що завдань буде трохи менше, а деякі попередні будуть переформульовані в стилі двох останніх.

%enditem

\end {large}}

%end

\newpage
%begin

{{
\begin {large}

\newpage
1. Що буде надруковано за виконання?

\begin{verbatim}
print('R', end=' ')
a=('a','b','c','d','e',0,1)
print(a[-10])
\end{verbatim}

%enditem


\newpage
2. Що буде надруковано за виконання?

\begin{verbatim}
print('R', end=' ')
h=(0,1,2,3,4,5,6,7,8,9)
print(h[:-8:3])
\end{verbatim}

%enditem



\newpage
3. Що буде надруковано за виконання?

\begin{verbatim}
print('R', end=' ')
a=[10,11,12,13,14]
print(a.index(8), end=' ')
\end{verbatim}

%enditem



\newpage
4. Що буде надруковано за виконання?

\begin{verbatim}
print('R', end=' ')
g=['0','1','2','3','4']
g[-1:-6]=[12]
print(g)
\end{verbatim}

%enditem



\newpage
5. Укажіть ТИП значення виразу.\par
\begin{verbatim}
[13,2,4]+[3,5]
\end{verbatim}
%enditem


\newpage
6. Що буде надруковано за виконання?\par
\begin{verbatim}
try:
    s = {12:8, 71:9, 22:6, 12:3}
    s[35] = 2
    for x in s.items():
        print(x, sep=' ', end=' ')
except:
    print('ERR')
\end{verbatim}
%enditem


\newpage
7. Нехай змінна \kw{s} позначає послідовність цілих, по якій можна ітерувати.

Написати вираз, що перевіряє, чи правда, що всі елементи послідовності є від’ємни\-ми.
%enditem


\newpage
8. Що буде надруковано за виконання?
%Оригинал был утерян. Требует отладки!
\begin{verbatim}
_c = 0

class D:
    _c = 1
    
    def __init__(self):
        self._c = 2
        _c = 3
        D._c = 4

obj = D()
D._c = 5
try:
    print(_c, end=' ')
    print(obj._c, end=' ')
    print(D._c, end=' ')
    print(obj._c, end=' ')
    print(obj._D__c)
except:
    print('ERR')
\end{verbatim}
%enditem


\newpage
9. Що буде надруковано за виконання?
\begin{verbatim}
class B0:
    def d(self): print(1, end=' ')
    def g(self): print(2, end=' ')

class B1(B0):
    def g(self): print(3, end=' ')

try:
    B1().d()
    B1().g()
except:
    print('ERR')
\end{verbatim}
%enditem


\newpage
10. Що буде надруковано за виконання?
\begin{verbatim}
class A:
    b = 1

    def __init__(self):
        c = 2
        b = 3

class B(A):
    c = 4
    a = 5

    def __init__(self):
        super().__init__()
        self.a = 6

try:
    obj = B()
    obj.b = obj.a + 10
    B.b = 13
except: print('EE', end=' ')

try: print(obj.b, end= ' ')
except: print('E1', end=' ')

try: print(A.b, end= ' ')
except: print('E2', end=' ')

try: print(B.b, end= ' ')
except: print('E3', end=' ')
\end{verbatim}
%enditem


\newpage
11. Змінні \kw{next} та \kw{prev} вузла списку позначають наступний та попередній за ним вузол відповідно. Починаючи з голови, у список записано послідовні цілі числа від~$-44$ до~$29$. Нехай змінна \kw{p} позначає вузол, що зберігає значення~$8$.
Яке значення зберігається у вузлі \kw{p.prev.prev.prev.prev.prev} ?
%enditem


\newpage
12. 
\begin{center}
	\adjustimage{max size={0.8\linewidth}{0.8\paperheight}}
	{BinTree/trees_exp/81.png}
\end{center}
Указати послідовність міток вершин дерева, зображеного на рис., що утвориться в результаті обходу дерева в прямому порядку. Мітки записувати через пробіл.\par Алгоритм починає роботу з кореня (на рис. це верхня вершина).
%enditem


\newpage
13. 
\begin{center}
	\adjustimage{max size={0.9\linewidth}{0.9\paperheight}}
	{Graphs/AdjList/232.png}
\end{center}
Для графа, зображеного на рис., вказати список суміжності вершини 3.\par
Формат оформлення відповіді:\par
список суміжності виписувати через пробіл на одному рядку, якщо пробілів десь буде більше одного - не важливо, упорядкування вершин у списку також значення не має.
%enditem


\newpage
14. Для графа на рисунку вказати послідовність, в якій алгоритм пошуку в глибину буде відмічати відвідані вершини, якщо списки суміжності вершин переглядаються алгоритмом за зростанням міток вершин та вершини відмічаються на першому вход\-женні у вершину.\par
Пошук розпочинається з вершини 1.\par

\begin{center}
	\adjustimage{max size={0.9\linewidth}{0.9\paperheight}}
	{Graphs/Traversals/26.png}
\end{center}
%enditem


\newpage
15. Для графа на рисунку з вершини 5 запущено алгоритм пошуку в ширину, що будує кістякове дерево. Списки суміжності вершин переглядаються алгоритмом за зростанням міток вершин.\par
\begin{center}
	\adjustimage{max size={0.9\linewidth}{0.9\paperheight}}
	{Graphs/Traversals/95.png}
\end{center}
\par Які з наведених ребер входять до складу отриманого кістякового дерева?\par
1) (4, 6)\par
2) (3, 4)\par
3) (4, 5)\par
4) (5, 6)\par
5) (2, 5)\par
6) (1, 4)\par
{\vskip 15pt} Номери відповідних ребер записати за зростанням через пробіл.


%enditem


\newpage
16. 

Рядки журналу через двокрапку містять ім'я користувача (ідентифікатор), час події,тип події, код події, наприклад
\begin{verbatim}
s="username : 2022-05-05 13:26 : ERROR : 404"
\end{verbatim}
у коді 
\begin{verbatim}
res = re.fullmatch('???', s) 
s = ''
code_str = ??? 
\end{verbatim}
чим треба замінити кожен з ???, щоб значенням code\_str був код події? 


\vspace{10mm}
\textbf{Завдання 17} 

Задано програму, що містить код 1-2 класів, можливе успадкування, та клієнтський код.
У наведеному коді невеликі частини закриті. 
Також наведено вихід програми.
Необхідно відновити закриті частини, щоб отриманий код відповідав наведеному виходу програми.


\textbf{Завдання 18-19} 

Кілька запитань за завданнями для самостійної роботи.
\vspace{10mm}
Скоріш за все, що завдань буде трохи менше, а деякі попередні будуть переформульовані в стилі двох останніх.

%enditem

\end {large}}

%end

\newpage
%begin

{{
\begin {large}

\newpage
1. Що буде надруковано за виконання?

\begin{verbatim}
print('R', end=' ')
g=('a','b','c','d','e',0,1,2,3,4)
print(g[-11])
\end{verbatim}

%enditem


\newpage
2. Що буде надруковано за виконання?

\begin{verbatim}
print('R', end=' ')
d='0123456789'
print(d[-8:10:-3])
\end{verbatim}

%enditem



\newpage
3. Що буде надруковано за виконання?

\begin{verbatim}
print('R', end=' ')
b=['0','1','2','3','4']
b.remove('6')
print(b)
\end{verbatim}

%enditem



\newpage
4. Що буде надруковано за виконання?

\begin{verbatim}
print('R', end=' ')
e=[10,11,12,13,14]
e[:-4]='10'
print(e)
\end{verbatim}

%enditem



\newpage
5. Укажіть ТИП значення виразу.\par
\begin{verbatim}
{x:0 for x in range(7)}
\end{verbatim}
%enditem


\newpage
6. Що буде надруковано за виконання?\par
\begin{verbatim}
try:
    d = {31:4, 58:8, 46:2, 28:1}
    d[24] = 6
    for x in d.keys():
        print(x, sep=' ', end=' ')
except:
    print('ERR')
\end{verbatim}
%enditem


\newpage
7. Записати ВИРАЗ, значенням якого є словник, ключами якого є цілі числа від 13 до 2000 (включно), що не діляться на жодне з чисел 2, 3, а ключам відповідають їх квадратні корені  
%enditem


\newpage
8. Що буде надруковано за виконання?
\begin{verbatim}
b = 0
k = 1

class B:
    k = 2
    
    def __init__(self):
        global k
        self.b = 3
        k = 4
        B.k = 5

obj = B()
try:
    print(b, end=' ')
    print(k, end=' ')
    print(B.k, end=' ')
    print(obj.b)
except:
    print('ERR')
\end{verbatim}
%enditem


\newpage
9. Що буде надруковано за виконання?
\begin{verbatim}
class C0:
    def b(self): print(1, end=' ')
    def f(self): print(2, end=' ')

class C1(C0):
    def b(self): print(3, end=' ')

try:
    C1().f()
    C0().b()
except:
    print('ERR')
\end{verbatim}
%enditem


\newpage
10. Що буде надруковано за виконання?
\begin{verbatim}
class A:
    a = 1

    def g1(self): print(2, end=' ')
    def _g2(self): print(3, end=' ')

    def main(self):
        try: print(self.a, end=' ')
        except: print('E1', end=' ')

        try: self.g1()
        except: print('E2', end=' ')

        try: self._g2()
        except: print('E3', end=' ')


class B(A):
    a = 4

    def g1(self): print(5, end=' ')
    def _g2(self): print(6, end=' ')

try: B().main()
except: print('E4', end=' ')
\end{verbatim}
%enditem


\newpage
11. Змінні \kw{next} та \kw{prev} вузла списку позначають наступний та попередній за ним вузол відповідно. Починаючи з голови, у список записано послідовні цілі числа від~$-54$ до~$38$. Нехай змінна \kw{p} позначає вузол, що зберігає значення~$1$.
Яке значення зберігається у вузлі \kw{p.next.next.next.next.next} ?
%enditem


\newpage
12. 
\begin{center}
	\adjustimage{max size={0.8\linewidth}{0.8\paperheight}}
	{BinTree/trees_exp/5.png}
\end{center}
Указати послідовність міток вершин дерева, зображеного на рис., що утвориться в результаті обходу дерева в оберненому порядку. Мітки записувати через пробіл.\par Алгоритм починає роботу з кореня (на рис. це верхня вершина).
%enditem


\newpage
13. 
\begin{center}
	\adjustimage{max size={0.9\linewidth}{0.9\paperheight}}
	{Graphs/AdjList/225.png}
\end{center}
Для графа, зображеного на рис., вказати список суміжності вершини 1.\par
Формат оформлення відповіді:\par
список суміжності виписувати через пробіл на одному рядку, якщо пробілів десь буде більше одного - не важливо, упорядкування вершин у списку також значення не має.
%enditem


\newpage
14. Для графа на рисунку вказати послідовність, в якій алгоритм пошуку в глибину буде відмічати відвідані вершини, якщо списки суміжності вершин переглядаються алгоритмом за зростанням міток вершин та вершини відмічаються на першому вход\-женні у вершину.\par
Пошук розпочинається з вершини 2.\par

\begin{center}
	\adjustimage{max size={0.9\linewidth}{0.9\paperheight}}
	{Graphs/Traversals/188.png}
\end{center}
%enditem


\newpage
15. Для графа на рисунку з вершини 0 запущено алгоритм пошуку в ширину, що будує кістякове дерево. Списки суміжності вершин переглядаються алгоритмом за зростанням міток вершин.\par
\begin{center}
	\adjustimage{max size={0.9\linewidth}{0.9\paperheight}}
	{Graphs/Traversals/123.png}
\end{center}
\par Які з наведених ребер входять до складу отриманого кістякового дерева?\par
1) (1, 6)\par
2) (1, 2)\par
3) (0, 6)\par
4) (0, 5)\par
5) (3, 5)\par
6) (0, 1)\par
{\vskip 15pt} Номери відповідних ребер записати за зростанням через пробіл.


%enditem


\newpage
16. 

Рядки журналу через двокрапку містять ім'я користувача (ідентифікатор), час події,тип події, код події, наприклад
\begin{verbatim}
s="username : 2022-05-05 13:26 : ERROR : 404"
\end{verbatim}
у коді 
\begin{verbatim}
res = re.fullmatch('???', s) 
s = ''
code_str = ??? 
\end{verbatim}
чим треба замінити кожен з ???, щоб значенням code\_str був код події? 


\vspace{10mm}
\textbf{Завдання 17} 

Задано програму, що містить код 1-2 класів, можливе успадкування, та клієнтський код.
У наведеному коді невеликі частини закриті. 
Також наведено вихід програми.
Необхідно відновити закриті частини, щоб отриманий код відповідав наведеному виходу програми.


\textbf{Завдання 18-19} 

Кілька запитань за завданнями для самостійної роботи.
\vspace{10mm}
Скоріш за все, що завдань буде трохи менше, а деякі попередні будуть переформульовані в стилі двох останніх.

%enditem

\end {large}}

%end

\newpage
%begin

{{
\begin {large}

\newpage
1. Що буде надруковано за виконання?

\begin{verbatim}
print('R', end=' ')
h=('a','b','c','d','e',0,1,2,4)
print(h[-12])
\end{verbatim}

%enditem


\newpage
2. Що буде надруковано за виконання?

\begin{verbatim}
print('R', end=' ')
b=['0','1','2','3','4','5','6','7','8','9']
print(b[1:15:2])
\end{verbatim}

%enditem



\newpage
3. Що буде надруковано за виконання?

\begin{verbatim}
print('R', end=' ')
h=[10,11,12,13,14]
h.insert(5, [1,0])
print(h)
\end{verbatim}

%enditem



\newpage
4. Що буде надруковано за виконання?

\begin{verbatim}
print('R', end=' ')
h=['0','1','2','3','4']
del h[-9:]
print(h)
\end{verbatim}

%enditem



\newpage
5. Укажіть ТИП значення виразу.\par
\begin{verbatim}
(2, 7, -4)
\end{verbatim}
%enditem


\newpage
6. Що буде надруковано за виконання?\par
\begin{verbatim}
try:
    d = {28:0, 51:6, 20:7, 53:6, 12:8}
    d[28] = 1
    for x in d.items():
        print(x, sep=' ', end=' ')
except:
    print('ERR')
\end{verbatim}
%enditem


\newpage
7. Нехай змінна \kw{s} позначає дуже довгу послідовність, по якій можна ітерувати. 

Написати вираз-генератор, що задає підпослідовність її елементів, що є рядками типу \kw{str}.

Обмеження: усі елементи шуканої підпослідовності одночасно в пам'яті розміщені бути не можуть. 
%enditem


\newpage
8. Що буде надруковано за виконання?
\begin{verbatim}
d = 0
k = 1

class C:
    d = 2
    
    def __init__(self):
        self.k = 3
        d = 4
        C.d = 5

obj = C()
try:
    print(d, end=' ')
    print(k, end=' ')
    print(C.d, end=' ')
    print(obj.k)
except:
    print('ERR')
\end{verbatim}
%enditem


\newpage
9. Що буде надруковано за виконання?
\begin{verbatim}
class A0:
    def c(self): print(1, end=' ')
    def f(self): print(2, end=' ')

class A1(A0):
    def f(self): print(3, end=' ')

try:
    A1().c()
    A0().f()
except:
    print('ERR')
\end{verbatim}
%enditem


\newpage
10. Що буде надруковано за виконання?
\begin{verbatim}
class A:
    b = 1

    def __init__(self):
        c = 2
        b = 3

class B(A):
    c = 4
    a = 5

    def __init__(self):
        super().__init__()
        self.a = 6

try:
    obj = B()
    obj.c = obj.b + 10
    B.c = 13
except: print('EE', end=' ')

try: print(obj.c, end= ' ')
except: print('E1', end=' ')

try: print(A.c, end= ' ')
except: print('E2', end=' ')

try: print(B.c, end= ' ')
except: print('E3', end=' ')
\end{verbatim}
%enditem


\newpage
11. Змінні \kw{next} та \kw{prev} вузла списку позначають наступний та попередній за ним вузол відповідно. Починаючи з голови, у список записано послідовні цілі числа від~$-36$ до~$27$. Нехай змінна \kw{p} позначає вузол, що зберігає значення~$-6$.
Яке значення зберігається у вузлі \kw{p.prev.prev.prev.next.next} ?
%enditem


\newpage
12. 
\begin{center}
	\adjustimage{max size={0.8\linewidth}{0.8\paperheight}}
	{BinTree/trees_exp/165.png}
\end{center}
Указати послідовність міток вершин дерева, зображеного на рис., що утвориться в результаті обходу дерева в оберненому порядку. Мітки записувати через пробіл.\par Алгоритм починає роботу з кореня (на рис. це верхня вершина).
%enditem


\newpage
13. 
\begin{center}
	\adjustimage{max size={0.9\linewidth}{0.9\paperheight}}
	{Graphs/AdjList/27.png}
\end{center}
Для графа, зображеного на рис., вказати список суміжності вершини 0.\par
Формат оформлення відповіді:\par
список суміжності виписувати через пробіл на одному рядку, якщо пробілів десь буде більше одного - не важливо, упорядкування вершин у списку також значення не має.
%enditem


\newpage
14. Для графа на рисунку вказати послідовність, в якій алгоритм пошуку в глибину буде відмічати відвідані вершини, якщо списки суміжності вершин переглядаються алгоритмом за зростанням міток вершин та вершини відмічаються на першому вход\-женні у вершину.\par
Пошук розпочинається з вершини 5.\par

\begin{center}
	\adjustimage{max size={0.9\linewidth}{0.9\paperheight}}
	{Graphs/Traversals/307.png}
\end{center}
%enditem


\newpage
15. Для графа на рисунку з вершини 3 запущено алгоритм пошуку в ширину, що будує кістякове дерево. Списки суміжності вершин переглядаються алгоритмом за зростанням міток вершин.\par
\begin{center}
	\adjustimage{max size={0.9\linewidth}{0.9\paperheight}}
	{Graphs/Traversals/309.png}
\end{center}
\par Які з наведених ребер входять до складу отриманого кістякового дерева?\par
1) (2, 4)\par
2) (2, 5)\par
3) (1, 4)\par
4) (0, 4)\par
5) (0, 3)\par
6) (2, 6)\par
{\vskip 15pt} Номери відповідних ребер записати за зростанням через пробіл.


%enditem


\newpage
16. 

Рядки журналу через двокрапку містять ім'я користувача (ідентифікатор), час події,тип події, код події, наприклад
\begin{verbatim}
s="username : 2022-05-05 13:26 : ERROR : 404"
\end{verbatim}
у коді 
\begin{verbatim}
res = re.fullmatch('???', s) 
s = ''
code_str = ??? 
\end{verbatim}
чим треба замінити кожен з ???, щоб значенням code\_str був код події? 


\vspace{10mm}
\textbf{Завдання 17} 

Задано програму, що містить код 1-2 класів, можливе успадкування, та клієнтський код.
У наведеному коді невеликі частини закриті. 
Також наведено вихід програми.
Необхідно відновити закриті частини, щоб отриманий код відповідав наведеному виходу програми.


\textbf{Завдання 18-19} 

Кілька запитань за завданнями для самостійної роботи.
\vspace{10mm}
Скоріш за все, що завдань буде трохи менше, а деякі попередні будуть переформульовані в стилі двох останніх.

%enditem

\end {large}}

%end

\newpage
%begin

{{
\begin {large}

\newpage
1. Що буде надруковано за виконання?

\begin{verbatim}
print('R', end=' ')
e=('a','b','c','d','e',0,1)
print(e[-9])
\end{verbatim}

%enditem


\newpage
2. Що буде надруковано за виконання?

\begin{verbatim}
print('R', end=' ')
a=(0,1,2,3,4,5,6,7,8,9)
print(a[5:8:3])
\end{verbatim}

%enditem



\newpage
3. Що буде надруковано за виконання?

\begin{verbatim}
print('R', end=' ')
g=['0','1','2','3','4']
g.extend('10')
print(g)
\end{verbatim}

%enditem



\newpage
4. Що буде надруковано за виконання?

\begin{verbatim}
print('R', end=' ')
c=[10,11,12,13,14]
c[:2]='11'
print(c)
\end{verbatim}

%enditem



\newpage
5. Укажіть ТИП значення виразу.\par
\begin{verbatim}
{x for x in range(7)}
\end{verbatim}
%enditem


\newpage
6. Що буде надруковано за виконання?\par
\begin{verbatim}
try:
    s = {11:2, 88:2, 40:5, 11:5}
    s[11] = 4
    for x in s :
        print(x, sep=' ', end=' ')
except:
    print('ERR')
\end{verbatim}
%enditem


\newpage
7. Нехай змінна \kw{s} позначає послідовність цілих, по якій можна ітерувати. 

Написати вираз, що перевіряє, чи правда, що послідовність  містить число 13.
%enditem


\newpage
8. Що буде надруковано за виконання?
\begin{verbatim}
b = 0
n = 1

class B:
    n = 2
    
    def __init__(self):
        global n
        self.b = 3
        n = 4
        B.n = 5

obj = B()
try:
    print(b, end=' ')
    print(n, end=' ')
    print(B.n, end=' ')
    print(obj.b)
except:
    print('ERR')
\end{verbatim}
%enditem


\newpage
9. Що буде надруковано за виконання?
\begin{verbatim}
class C0:
    def d(self): print(1, end=' ')
    def g(self): print(2, end=' ')

class C1(C0):
    def d(self): print(3, end=' ')

try:
    C1().g()
    C1().d()
except:
    print('ERR')
\end{verbatim}
%enditem


\newpage
10. Що буде надруковано за виконання?
\begin{verbatim}
class A:
    a = 1

    def __init__(self):
        b = 2
        b = 3

class B(A):
    a = 4
    c = 5

    def __init__(self):
        super().__init__()
        self.c = 6

try:
    obj = B()
    obj.a = obj.c + 10
    B.a = 13
except: print('EE', end=' ')

try: print(obj.a, end= ' ')
except: print('E1', end=' ')

try: print(A.a, end= ' ')
except: print('E2', end=' ')

try: print(B.a, end= ' ')
except: print('E3', end=' ')
\end{verbatim}
%enditem


\newpage
11. Змінні \kw{next} та \kw{prev} вузла списку позначають наступний та попередній за ним вузол відповідно. Починаючи з голови, у список записано послідовні цілі числа від~$-39$ до~$39$. Нехай змінна \kw{p} позначає вузол, що зберігає значення~$-2$.
Яке значення зберігається у вузлі \kw{p.next.next.next.prev.prev} ?
%enditem


\newpage
12. 
\begin{center}
	\adjustimage{max size={0.8\linewidth}{0.8\paperheight}}
	{BinTree/trees_exp/142.png}
\end{center}
Указати послідовність міток вершин дерева, зображеного на рис., що утвориться в результаті обходу дерева в прямому порядку. Мітки записувати через пробіл.\par Алгоритм починає роботу з кореня (на рис. це верхня вершина).
%enditem


\newpage
13. 
\begin{center}
	\adjustimage{max size={0.9\linewidth}{0.9\paperheight}}
	{Graphs/AdjList/318.png}
\end{center}
Для графа, зображеного на рис., вказати список суміжності вершини 3.\par
Формат оформлення відповіді:\par
список суміжності виписувати через пробіл на одному рядку, якщо пробілів десь буде більше одного - не важливо, упорядкування вершин у списку також значення не має.
%enditem


\newpage
14. Для графа на рисунку вказати послідовність, в якій алгоритм пошуку в глибину буде відмічати відвідані вершини, якщо списки суміжності вершин переглядаються алгоритмом за зростанням міток вершин та вершини відмічаються на першому вход\-женні у вершину.\par
Пошук розпочинається з вершини 1.\par

\begin{center}
	\adjustimage{max size={0.9\linewidth}{0.9\paperheight}}
	{Graphs/Traversals/290.png}
\end{center}
%enditem


\newpage
15. Для графа на рисунку з вершини 0 запущено алгоритм пошуку в ширину, що будує кістякове дерево. Списки суміжності вершин переглядаються алгоритмом за зростанням міток вершин.\par
\begin{center}
	\adjustimage{max size={0.9\linewidth}{0.9\paperheight}}
	{Graphs/Traversals/62.png}
\end{center}
\par Які з наведених ребер входять до складу отриманого кістякового дерева?\par
1) (0, 4)\par
2) (5, 6)\par
3) (2, 6)\par
4) (2, 5)\par
5) (3, 6)\par
6) (0, 1)\par
{\vskip 15pt} Номери відповідних ребер записати за зростанням через пробіл.


%enditem


\newpage
16. 

Рядки журналу через двокрапку містять ім'я користувача (ідентифікатор), час події,тип події, код події, наприклад
\begin{verbatim}
s="username : 2022-05-05 13:26 : ERROR : 404"
\end{verbatim}
у коді 
\begin{verbatim}
res = re.fullmatch('???', s) 
s = ''
code_str = ??? 
\end{verbatim}
чим треба замінити кожен з ???, щоб значенням code\_str був код події? 


\vspace{10mm}
\textbf{Завдання 17} 

Задано програму, що містить код 1-2 класів, можливе успадкування, та клієнтський код.
У наведеному коді невеликі частини закриті. 
Також наведено вихід програми.
Необхідно відновити закриті частини, щоб отриманий код відповідав наведеному виходу програми.


\textbf{Завдання 18-19} 

Кілька запитань за завданнями для самостійної роботи.
\vspace{10mm}
Скоріш за все, що завдань буде трохи менше, а деякі попередні будуть переформульовані в стилі двох останніх.

%enditem

\end {large}}

%end

\newpage
%begin

{{
\begin {large}

\newpage
1. Що буде надруковано за виконання?

\begin{verbatim}
print('R', end=' ')
f=('a','b','c','d','e',0,1,2,3)
print(f[-5])
\end{verbatim}

%enditem


\newpage
2. Що буде надруковано за виконання?

\begin{verbatim}
print('R', end=' ')
g='0123456789'
print(g[-3:14:-2])
\end{verbatim}

%enditem



\newpage
3. Що буде надруковано за виконання?

\begin{verbatim}
print('R', end=' ')
f=['0','1','2','3','4']
print(f.pop(2), end=' ')
print(f)
\end{verbatim}

%enditem



\newpage
4. Що буде надруковано за виконання?

\begin{verbatim}
print('R', end=' ')
g=[10,11,12,13,14]
g[3:8]=[1,2]
print(g)
\end{verbatim}

%enditem



\newpage
5. Укажіть ТИП значення виразу.\par
\begin{verbatim}
{1:3, 5:8}
\end{verbatim}
%enditem


\newpage
6. Що буде надруковано за виконання?\par
\begin{verbatim}
try:
    t = {53:6, 90:3, 36:6, 53:0}
    t[90] = 0
    for x in t.values():
        print(x, sep=' ', end=' ')
except:
    print('ERR')
\end{verbatim}
%enditem


\newpage
7. Нехай змінна \kw{s} позначає послідовність цілих, по якій можна ітерувати. 

Написати вираз, що перевіряє, чи правда, що послідовність  містить хоча б одне число, що ділиться на 7.
%enditem


\newpage
8. Що буде надруковано за виконання?
\begin{verbatim}
c = 0
h = 1

class A:
    c = 2
    
    def __init__(self):
        self.c = 3
        c = 4
        A.c = 5

obj = A()
try:
    print(c, end=' ')
    print(h, end=' ')
    print(A.c, end=' ')
    print(obj.c)
except:
    print('ERR')
\end{verbatim}
%enditem


\newpage
9. Що буде надруковано за виконання?
\begin{verbatim}
class B0:
    def a(self): print(1, end=' ')
    def h(self): print(2, end=' ')

class B1(B0):
    def a(self): print(3, end=' ')

try:
    B0().h()
    B1().a()
except:
    print('ERR')
\end{verbatim}
%enditem


\newpage
10. Що буде надруковано за виконання?
\begin{verbatim}
class A:
    a = 1

    def __init__(self):
        c = 2
        c = 3

class B(A):
    c = 4
    b = 5

    def __init__(self):
        super().__init__()
        self.b = 6

try:
    obj = B()
    obj.b = obj.a + 10
    B.b = 13
except: print('EE', end=' ')

try: print(obj.b, end= ' ')
except: print('E1', end=' ')

try: print(A.b, end= ' ')
except: print('E2', end=' ')

try: print(B.b, end= ' ')
except: print('E3', end=' ')
\end{verbatim}
%enditem


\newpage
11. Змінні \kw{next} та \kw{prev} вузла списку позначають наступний та попередній за ним вузол відповідно. Починаючи з голови, у список записано послідовні цілі числа від~$-41$ до~$35$. Нехай змінна \kw{p} позначає вузол, що зберігає значення~$5$.
Яке значення зберігається у вузлі \kw{p.prev.prev.prev.prev.prev} ?
%enditem


\newpage
12. 
\begin{center}
	\adjustimage{max size={0.8\linewidth}{0.8\paperheight}}
	{BinTree/trees_exp/104.png}
\end{center}
Указати послідовність міток вершин дерева, зображеного на рис., що утвориться в результаті обходу дерева в прямому порядку. Мітки записувати через пробіл.\par Алгоритм починає роботу з кореня (на рис. це верхня вершина).
%enditem


\newpage
13. 
\begin{center}
	\adjustimage{max size={0.9\linewidth}{0.9\paperheight}}
	{Graphs/AdjList/344.png}
\end{center}
Для графа, зображеного на рис., вказати список суміжності вершини 1.\par
Формат оформлення відповіді:\par
список суміжності виписувати через пробіл на одному рядку, якщо пробілів десь буде більше одного - не важливо, упорядкування вершин у списку також значення не має.
%enditem


\newpage
14. Для графа на рисунку з вершини 2 запущено алгоритм пошуку в глибину, що будує кістякове дерево. Списки суміжності вершин переглядаються алгоритмом за зростанням міток вершин.\par
\begin{center}
	\adjustimage{max size={0.9\linewidth}{0.9\paperheight}}
	{Graphs/Traversals/134.png}
\end{center}
\par Які з наведених ребер входять до складу отриманого кістякового дерева?\par
1) (0, 4)\par
2) (3, 4)\par
3) (0, 5)\par
4) (2, 5)\par
5) (3, 5)\par
6) (2, 6)\par
{\vskip 15pt} Номери відповідних ребер записати за зростанням через пробіл.

%enditem


\newpage
15. Для графа на рисунку з вершини 5 запущено алгоритм пошуку в ширину, що будує кістякове дерево. Списки суміжності вершин переглядаються алгоритмом за зростанням міток вершин.\par
\begin{center}
	\adjustimage{max size={0.9\linewidth}{0.9\paperheight}}
	{Graphs/Traversals/73.png}
\end{center}
\par Які з наведених ребер входять до складу отриманого кістякового дерева?\par
1) (3, 6)\par
2) (1, 2)\par
3) (5, 6)\par
4) (0, 6)\par
5) (4, 5)\par
6) (0, 3)\par
{\vskip 15pt} Номери відповідних ребер записати за зростанням через пробіл.


%enditem


\newpage
16. 

Рядки журналу через двокрапку містять ім'я користувача (ідентифікатор), час події,тип події, код події, наприклад
\begin{verbatim}
s="username : 2022-05-05 13:26 : ERROR : 404"
\end{verbatim}
у коді 
\begin{verbatim}
res = re.fullmatch('???', s) 
s = ''
code_str = ??? 
\end{verbatim}
чим треба замінити кожен з ???, щоб значенням code\_str був код події? 


\vspace{10mm}
\textbf{Завдання 17} 

Задано програму, що містить код 1-2 класів, можливе успадкування, та клієнтський код.
У наведеному коді невеликі частини закриті. 
Також наведено вихід програми.
Необхідно відновити закриті частини, щоб отриманий код відповідав наведеному виходу програми.


\textbf{Завдання 18-19} 

Кілька запитань за завданнями для самостійної роботи.
\vspace{10mm}
Скоріш за все, що завдань буде трохи менше, а деякі попередні будуть переформульовані в стилі двох останніх.

%enditem

\end {large}}

%end

\newpage
%begin

{{
\begin {large}

\newpage
1. Що буде надруковано за виконання?

\begin{verbatim}
print('R', end=' ')
g=('a','b','c','d','e',0,1,2,3,4)
print(g[1])
\end{verbatim}

%enditem


\newpage
2. Що буде надруковано за виконання?

\begin{verbatim}
print('R', end=' ')
h=[0,1,2,3,4,5,6,7,8,9]
print(h[:12:-3])
\end{verbatim}

%enditem



\newpage
3. Що буде надруковано за виконання?

\begin{verbatim}
print('R', end=' ')
e=[10,11,12,13,14]
e.append([1,2])
print(e)
\end{verbatim}

%enditem



\newpage
4. Що буде надруковано за виконання?

\begin{verbatim}
print('R', end=' ')
a=['0','1','2','3','4']
a[-4:-9]=[2,3]
print(a)
\end{verbatim}

%enditem



\newpage
5. Укажіть ТИП значення виразу.\par
\begin{verbatim}
{(1,4), 8}
\end{verbatim}
%enditem


\newpage
6. Що буде надруковано за виконання?\par
\begin{verbatim}
try:
    s = {31:5, 70:6, 16:1, 22:1}
    s[70] = 7
    for x, y in s.items():
        print(x, y, sep=' ', end=' ')
except:
    print('ERR')
\end{verbatim}
%enditem


\newpage
7. Нехай змінна \kw{s} позначає дуже довгу послідовність, по якій можна ітерувати. 

Написати вираз-генератор, що задає підпослідовність її елементів, що є числами типу \kw{int}.

Обмеження: усі елементи шуканої підпослідовності одночасно в пам'яті розміщені бути не можуть. 
%enditem


\newpage
8. Що буде надруковано за виконання?
\begin{verbatim}
a = 0
m = 1

class D:
    m = 2
    
    def __init__(self):
        global m
        self.m = 3
        m = 4
        D.m = 5

obj = D()
try:
    print(a, end=' ')
    print(m, end=' ')
    print(D.m, end=' ')
    print(obj.m)
except:
    print('ERR')
\end{verbatim}
%enditem


\newpage
9. Що буде надруковано за виконання?
\begin{verbatim}
class D0:
    def b(self): print(1, end=' ')
    def h(self): print(2, end=' ')

class D1(D0):
    def b(self): print(3, end=' ')

try:
    D1().h()
    D1().b()
except:
    print('ERR')
\end{verbatim}
%enditem


\newpage
10. Що буде надруковано за виконання?
\begin{verbatim}
class A:
    c = 1

    def __init__(self):
        b = 2
        c = 3

class B(A):
    c = 4
    a = 5

    def __init__(self):
        super().__init__()
        self.a = 6

try:
    obj = B()
    obj.b = obj.a + 10
    B.b = 13
except: print('EE', end=' ')

try: print(obj.b, end= ' ')
except: print('E1', end=' ')

try: print(A.b, end= ' ')
except: print('E2', end=' ')

try: print(B.b, end= ' ')
except: print('E3', end=' ')
\end{verbatim}
%enditem


\newpage
11. Змінні \kw{next} та \kw{prev} вузла списку позначають наступний та попередній за ним вузол відповідно. Починаючи з голови, у список записано послідовні цілі числа від~$-55$ до~$36$. Нехай змінна \kw{p} позначає вузол, що зберігає значення~$-7$.
Яке значення зберігається у вузлі \kw{p.next.next.next.next.next} ?
%enditem


\newpage
12. 
\begin{center}
	\adjustimage{max size={0.8\linewidth}{0.8\paperheight}}
	{BinTree/trees_exp/19.png}
\end{center}
Указати послідовність міток вершин дерева, зображеного на рис., що утвориться в результаті обходу дерева в оберненому порядку. Мітки записувати через пробіл.\par Алгоритм починає роботу з кореня (на рис. це верхня вершина).
%enditem


\newpage
13. 
\begin{center}
	\adjustimage{max size={0.9\linewidth}{0.9\paperheight}}
	{Graphs/AdjList/236.png}
\end{center}
Для графа, зображеного на рис., вказати список суміжності вершини 0.\par
Формат оформлення відповіді:\par
список суміжності виписувати через пробіл на одному рядку, якщо пробілів десь буде більше одного - не важливо, упорядкування вершин у списку також значення не має.
%enditem


\newpage
14. Для графа на рисунку вказати послідовність, в якій алгоритм пошуку в глибину буде відмічати відвідані вершини, якщо списки суміжності вершин переглядаються алгоритмом за зростанням міток вершин та вершини відмічаються на першому вход\-женні у вершину.\par
Пошук розпочинається з вершини 2.\par

\begin{center}
	\adjustimage{max size={0.9\linewidth}{0.9\paperheight}}
	{Graphs/Traversals/105.png}
\end{center}
%enditem


\newpage
15. Для графа на рисунку з вершини 1 запущено алгоритм пошуку в ширину, що будує кістякове дерево. Списки суміжності вершин переглядаються алгоритмом за зростанням міток вершин.\par
\begin{center}
	\adjustimage{max size={0.9\linewidth}{0.9\paperheight}}
	{Graphs/Traversals/314.png}
\end{center}
\par Які з наведених ребер входять до складу отриманого кістякового дерева?\par
1) (3, 6)\par
2) (2, 4)\par
3) (0, 1)\par
4) (2, 3)\par
5) (0, 2)\par
6) (2, 5)\par
{\vskip 15pt} Номери відповідних ребер записати за зростанням через пробіл.


%enditem


\newpage
16. 

Рядки журналу через двокрапку містять ім'я користувача (ідентифікатор), час події,тип події, код події, наприклад
\begin{verbatim}
s="username : 2022-05-05 13:26 : ERROR : 404"
\end{verbatim}
у коді 
\begin{verbatim}
res = re.fullmatch('???', s) 
s = ''
code_str = ??? 
\end{verbatim}
чим треба замінити кожен з ???, щоб значенням code\_str був код події? 


\vspace{10mm}
\textbf{Завдання 17} 

Задано програму, що містить код 1-2 класів, можливе успадкування, та клієнтський код.
У наведеному коді невеликі частини закриті. 
Також наведено вихід програми.
Необхідно відновити закриті частини, щоб отриманий код відповідав наведеному виходу програми.


\textbf{Завдання 18-19} 

Кілька запитань за завданнями для самостійної роботи.
\vspace{10mm}
Скоріш за все, що завдань буде трохи менше, а деякі попередні будуть переформульовані в стилі двох останніх.

%enditem

\end {large}}

%end

\newpage
%begin

{{
\begin {large}

\newpage
1. Що буде надруковано за виконання?

\begin{verbatim}
print('R', end=' ')
a=('a','b','c','d','e',0,1,2)
print(a[4])
\end{verbatim}

%enditem


\newpage
2. Що буде надруковано за виконання?

\begin{verbatim}
print('R', end=' ')
d=('0','1','2','3','4','5','6','7','8','9')
print(d[:-12:-1])
\end{verbatim}

%enditem



\newpage
3. Що буде надруковано за виконання?

\begin{verbatim}
print('R', end=' ')
d=[10,11,12,13,14]
d+='11'
print(d)
\end{verbatim}

%enditem



\newpage
4. Що буде надруковано за виконання?

\begin{verbatim}
print('R', end=' ')
e=[10,11,12,13,14]
e[-8:]='13'
print(e)
\end{verbatim}

%enditem



\newpage
5. Укажіть ТИП значення виразу.\par
\begin{verbatim}
[x for x in range(7)]
\end{verbatim}
%enditem


\newpage
6. Що буде надруковано за виконання?\par
\begin{verbatim}
try:
    d = {56:6, 21:9, 42:2, 42:2}
    d[21] = 2
    for x in d.items():
        print(*x, sep=' ', end=' ')
except:
    print('ERR')
\end{verbatim}
%enditem


\newpage
7. 
 Задано числову послідовність \kw{s} довжини 6000. Записати ВИРАЗ, значенням якого є список, що складається з першої половини послідовності, записаної в оберненому порядку.
%enditem


\newpage
8. Що буде надруковано за виконання?
%Оригинал был утерян. Требует отладки!
\begin{verbatim}
_b = 0

class B:
    _b = 1
    
    def __init__(self):
        self._b = 2
        _b = 3
        B._b = 4

obj = B()
obj._b = 5
try:
    print(_b, end=' ')
    print(obj._b, end=' ')
    print(B._b, end=' ')
    print(obj._b, end=' ')
    print(obj._B__b)
except:
    print('ERR')
\end{verbatim}
%enditem


\newpage
9. Що буде надруковано за виконання?
\begin{verbatim}
class B0:
    def b(self): print(1, end=' ')
    def f(self): print(2, end=' ')

class B1(B0):
    def f(self): print(3, end=' ')

try:
    B0().b()
    B1().f()
except:
    print('ERR')
\end{verbatim}
%enditem


\newpage
10. Що буде надруковано за виконання?
\begin{verbatim}
class A:
    b = 1

    def __init__(self):
        a = 2
        a = 3

class B(A):
    a = 4
    c = 5

    def __init__(self):
        super().__init__()
        self.c = 6

try:
    obj = B()
    obj.c = obj.b + 10
    B.c = 13
except: print('EE', end=' ')

try: print(obj.c, end= ' ')
except: print('E1', end=' ')

try: print(A.c, end= ' ')
except: print('E2', end=' ')

try: print(B.c, end= ' ')
except: print('E3', end=' ')
\end{verbatim}
%enditem


\newpage
11. Змінні \kw{next} та \kw{prev} вузла списку позначають наступний та попередній за ним вузол відповідно. Починаючи з голови, у список записано послідовні цілі числа від~$-38$ до~$23$. Нехай змінна \kw{p} позначає вузол, що зберігає значення~$6$.
Яке значення зберігається у вузлі \kw{p.next.next.prev.next.prev} ?
%enditem


\newpage
12. 
\begin{center}
	\adjustimage{max size={0.8\linewidth}{0.8\paperheight}}
	{BinTree/trees_exp/89.png}
\end{center}
Указати послідовність міток вершин дерева, зображеного на рис., що утвориться в результаті обходу дерева в прямому порядку. Мітки записувати через пробіл.\par Алгоритм починає роботу з кореня (на рис. це верхня вершина).
%enditem


\newpage
13. 
\begin{center}
	\adjustimage{max size={0.9\linewidth}{0.9\paperheight}}
	{Graphs/AdjList/296.png}
\end{center}
Для графа, зображеного на рис., вказати список суміжності вершини 6.\par
Формат оформлення відповіді:\par
список суміжності виписувати через пробіл на одному рядку, якщо пробілів десь буде більше одного - не важливо, упорядкування вершин у списку також значення не має.
%enditem


\newpage
14. Для графа на рисунку з вершини 5 запущено алгоритм пошуку в глибину, що будує кістякове дерево. Списки суміжності вершин переглядаються алгоритмом за зростанням міток вершин.\par
\begin{center}
	\adjustimage{max size={0.9\linewidth}{0.9\paperheight}}
	{Graphs/Traversals/31.png}
\end{center}
\par Які з наведених ребер входять до складу отриманого кістякового дерева?\par
1) (2, 3)\par
2) (1, 5)\par
3) (0, 5)\par
4) (3, 4)\par
5) (2, 6)\par
6) (0, 6)\par
{\vskip 15pt} Номери відповідних ребер записати за зростанням через пробіл.

%enditem


\newpage
15. Для графа на рисунку з вершини 4 запущено алгоритм пошуку в ширину, що будує кістякове дерево. Списки суміжності вершин переглядаються алгоритмом за зростанням міток вершин.\par
\begin{center}
	\adjustimage{max size={0.9\linewidth}{0.9\paperheight}}
	{Graphs/Traversals/165.png}
\end{center}
\par Які з наведених ребер входять до складу отриманого кістякового дерева?\par
1) (1, 3)\par
2) (1, 4)\par
3) (5, 6)\par
4) (2, 6)\par
5) (3, 5)\par
6) (2, 3)\par
{\vskip 15pt} Номери відповідних ребер записати за зростанням через пробіл.


%enditem


\newpage
16. 

Рядки журналу через двокрапку містять ім'я користувача (ідентифікатор), час події,тип події, код події, наприклад
\begin{verbatim}
s="username : 2022-05-05 13:26 : ERROR : 404"
\end{verbatim}
у коді 
\begin{verbatim}
res = re.fullmatch('???', s) 
s = ''
code_str = ??? 
\end{verbatim}
чим треба замінити кожен з ???, щоб значенням code\_str був код події? 


\vspace{10mm}
\textbf{Завдання 17} 

Задано програму, що містить код 1-2 класів, можливе успадкування, та клієнтський код.
У наведеному коді невеликі частини закриті. 
Також наведено вихід програми.
Необхідно відновити закриті частини, щоб отриманий код відповідав наведеному виходу програми.


\textbf{Завдання 18-19} 

Кілька запитань за завданнями для самостійної роботи.
\vspace{10mm}
Скоріш за все, що завдань буде трохи менше, а деякі попередні будуть переформульовані в стилі двох останніх.

%enditem

\end {large}}

%end

\newpage
%begin

{{
\begin {large}

\newpage
1. Що буде надруковано за виконання?

\begin{verbatim}
print('R', end=' ')
b=('a','b','c','d','e',0,1,2,3,4)
print(b[-7])
\end{verbatim}

%enditem


\newpage
2. Що буде надруковано за виконання?

\begin{verbatim}
print('R', end=' ')
c=['0','1','2','3','4','5','6','7','8','9']
print(c[-15:3:-2])
\end{verbatim}

%enditem



\newpage
3. Що буде надруковано за виконання?

\begin{verbatim}
print('R', end=' ')
b=['0','1','2','3','4']
print(b.index('7'), end=' ')
\end{verbatim}

%enditem



\newpage
4. Що буде надруковано за виконання?

\begin{verbatim}
print('R', end=' ')
d=['0','1','2','3','4']
del d[6:3]
print(d)
\end{verbatim}

%enditem



\newpage
5. Укажіть ТИП значення виразу.\par
\begin{verbatim}
{x for x in range(7) if el%3==1}
\end{verbatim}
%enditem


\newpage
6. Що буде надруковано за виконання?\par
\begin{verbatim}
try:
    t = {72:7, 36:3, 44:0, 49:5, 44:3}
    t[47] = 3
    for x in t.keys():
        print(x, sep=' ', end=' ')
except:
    print('ERR')
\end{verbatim}
%enditem


\newpage
7. Нехай змінна \kw{s} позначає послідовність цілих, по якій можна ітерувати. 

Написати вираз, що перевіряє, чи правда, що послідовність  містить число 35.
%enditem


\newpage
8. Що буде надруковано за виконання?
\begin{verbatim}
a = 0
j = 1

class D:
    a = 2
    
    def __init__(self):
        global a
        self.a = 3
        a = 4
        D.a = 5

obj = D()
try:
    print(a, end=' ')
    print(j, end=' ')
    print(D.a, end=' ')
    print(obj.a)
except:
    print('ERR')
\end{verbatim}
%enditem


\newpage
9. Що буде надруковано за виконання?
\begin{verbatim}
class A0:
    def a(self): print(1, end=' ')
    def g(self): print(2, end=' ')

class A1(A0):
    def a(self): print(3, end=' ')

try:
    A1().g()
    A0().a()
except:
    print('ERR')
\end{verbatim}
%enditem


\newpage
10. Що буде надруковано за виконання?
\begin{verbatim}
class A:
    b = 1

    def _g1(self): print(2, end=' ')
    def __g2(self): print(3, end=' ')

    def main(self):
        try: print(self.b, end=' ')
        except: print('E1', end=' ')

        try: self._g1()
        except: print('E2', end=' ')

        try: self.__g2()
        except: print('E3', end=' ')


class B(A):
    b = 4

    def _g1(self): print(5, end=' ')
    def __g2(self): print(6, end=' ')

try: B().main()
except: print('E4', end=' ')
\end{verbatim}
%enditem


\newpage
11. Змінні \kw{next} та \kw{prev} вузла списку позначають наступний та попередній за ним вузол відповідно. Починаючи з голови, у список записано послідовні цілі числа від~$-37$ до~$22$. Нехай змінна \kw{p} позначає вузол, що зберігає значення~$-8$.
Яке значення зберігається у вузлі \kw{p.prev.prev.next.prev.next} ?
%enditem


\newpage
12. 
\begin{center}
	\adjustimage{max size={0.8\linewidth}{0.8\paperheight}}
	{BinTree/trees_exp/13.png}
\end{center}
Указати послідовність міток вершин дерева, зображеного на рис., що утвориться в результаті обходу дерева в оберненому порядку. Мітки записувати через пробіл.\par Алгоритм починає роботу з кореня (на рис. це верхня вершина).
%enditem


\newpage
13. 
\begin{center}
	\adjustimage{max size={0.9\linewidth}{0.9\paperheight}}
	{Graphs/AdjList/235.png}
\end{center}
Для графа, зображеного на рис., вказати список суміжності вершини 0.\par
Формат оформлення відповіді:\par
список суміжності виписувати через пробіл на одному рядку, якщо пробілів десь буде більше одного - не важливо, упорядкування вершин у списку також значення не має.
%enditem


\newpage
14. Для графа на рисунку з вершини 0 запущено алгоритм пошуку в глибину, що будує кістякове дерево. Списки суміжності вершин переглядаються алгоритмом за зростанням міток вершин.\par
\begin{center}
	\adjustimage{max size={0.9\linewidth}{0.9\paperheight}}
	{Graphs/Traversals/49.png}
\end{center}
\par Які з наведених ребер входять до складу отриманого кістякового дерева?\par
1) (2, 6)\par
2) (4, 5)\par
3) (0, 5)\par
4) (1, 6)\par
5) (1, 4)\par
6) (3, 4)\par
{\vskip 15pt} Номери відповідних ребер записати за зростанням через пробіл.

%enditem


\newpage
15. Для графа на рисунку з вершини 0 запущено алгоритм пошуку в ширину, що будує кістякове дерево. Списки суміжності вершин переглядаються алгоритмом за зростанням міток вершин.\par
\begin{center}
	\adjustimage{max size={0.9\linewidth}{0.9\paperheight}}
	{Graphs/Traversals/261.png}
\end{center}
\par Які з наведених ребер входять до складу отриманого кістякового дерева?\par
1) (1, 4)\par
2) (2, 3)\par
3) (1, 5)\par
4) (5, 6)\par
5) (4, 5)\par
6) (0, 6)\par
{\vskip 15pt} Номери відповідних ребер записати за зростанням через пробіл.


%enditem


\newpage
16. 

Рядки журналу через двокрапку містять ім'я користувача (ідентифікатор), час події,тип події, код події, наприклад
\begin{verbatim}
s="username : 2022-05-05 13:26 : ERROR : 404"
\end{verbatim}
у коді 
\begin{verbatim}
res = re.fullmatch('???', s) 
s = ''
code_str = ??? 
\end{verbatim}
чим треба замінити кожен з ???, щоб значенням code\_str був код події? 


\vspace{10mm}
\textbf{Завдання 17} 

Задано програму, що містить код 1-2 класів, можливе успадкування, та клієнтський код.
У наведеному коді невеликі частини закриті. 
Також наведено вихід програми.
Необхідно відновити закриті частини, щоб отриманий код відповідав наведеному виходу програми.


\textbf{Завдання 18-19} 

Кілька запитань за завданнями для самостійної роботи.
\vspace{10mm}
Скоріш за все, що завдань буде трохи менше, а деякі попередні будуть переформульовані в стилі двох останніх.

%enditem

\end {large}}

%end

\newpage
%begin

{{
\begin {large}

\newpage
1. Що буде надруковано за виконання?

\begin{verbatim}
print('R', end=' ')
c=('a','b','c','d','e',0,1)
print(c[-4])
\end{verbatim}

%enditem


\newpage
2. Що буде надруковано за виконання?

\begin{verbatim}
print('R', end=' ')
e=[0,1,2,3,4,5,6,7,8,9]
print(e[-11::3])
\end{verbatim}

%enditem



\newpage
3. Що буде надруковано за виконання?

\begin{verbatim}
print('R', end=' ')
c=['0','1','2','3','4']
c+='11'
print(c)
\end{verbatim}

%enditem



\newpage
4. Що буде надруковано за виконання?

\begin{verbatim}
print('R', end=' ')
b=['0','1','2','3','4']
del b[5:-5]
print(b)
\end{verbatim}

%enditem



\newpage
5. Укажіть ТИП значення виразу.\par
\begin{verbatim}
{y:13 for y in range(-7,7)}
\end{verbatim}
%enditem


\newpage
6. Що буде надруковано за виконання?\par
\begin{verbatim}
try:
    t = {55:7, 41:1, 68:9, 52:9, 17:2}
    t[68] = 7
    for x, y in t.items():
        print(x, y, sep=' ', end=' ')
except:
    print('ERR')
\end{verbatim}
%enditem


\newpage
7. Нехай змінна \kw{s} позначає дуже довгу послідовність цілих, по якій можна ітерувати. 

Написати вираз-генератор, що задає підпослідовність її елементів, що є додатни\-ми. 

Обмеження: усі елементи шуканої підпослідовності одночасно в пам'яті розміщені бути не можуть. 
%enditem


\newpage
8. Що буде надруковано за виконання?
\begin{verbatim}
d = 0
j = 1

class B:
    j = 2
    
    def __init__(self):
        self.j = 3
        j = 4
        B.j = 5

obj = B()
try:
    print(d, end=' ')
    print(j, end=' ')
    print(B.j, end=' ')
    print(obj.j)
except:
    print('ERR')
\end{verbatim}
%enditem


\newpage
9. Що буде надруковано за виконання?
\begin{verbatim}
class D0:
    def c(self): print(1, end=' ')
    def g(self): print(2, end=' ')

class D1(D0):
    def g(self): print(3, end=' ')

try:
    D1().c()
    D1().g()
except:
    print('ERR')
\end{verbatim}
%enditem


\newpage
10. Що буде надруковано за виконання?
\begin{verbatim}
class A:
    c = 1

    def __init__(self):
        a = 2
        c = 3

class B(A):
    c = 4
    b = 5

    def __init__(self):
        super().__init__()
        self.b = 6

try:
    obj = B()
    obj.a = obj.c + 10
    B.a = 13
except: print('EE', end=' ')

try: print(obj.a, end= ' ')
except: print('E1', end=' ')

try: print(A.a, end= ' ')
except: print('E2', end=' ')

try: print(B.a, end= ' ')
except: print('E3', end=' ')
\end{verbatim}
%enditem


\newpage
11. Змінні \kw{next} та \kw{prev} вузла списку позначають наступний та попередній за ним вузол відповідно. Починаючи з голови, у список записано послідовні цілі числа від~$-64$ до~$21$. Нехай змінна \kw{p} позначає вузол, що зберігає значення~$3$.
Яке значення зберігається у вузлі \kw{p.next.next.next.prev.prev} ?
%enditem


\newpage
12. 
\begin{center}
	\adjustimage{max size={0.8\linewidth}{0.8\paperheight}}
	{BinTree/trees_exp/78.png}
\end{center}
Указати послідовність міток вершин дерева, зображеного на рис., що утвориться в результаті обходу дерева в оберненому порядку. Мітки записувати через пробіл.\par Алгоритм починає роботу з кореня (на рис. це верхня вершина).
%enditem


\newpage
13. 
\begin{center}
	\adjustimage{max size={0.9\linewidth}{0.9\paperheight}}
	{Graphs/AdjList/329.png}
\end{center}
Для графа, зображеного на рис., вказати список суміжності вершини 6.\par
Формат оформлення відповіді:\par
список суміжності виписувати через пробіл на одному рядку, якщо пробілів десь буде більше одного - не важливо, упорядкування вершин у списку також значення не має.
%enditem


\newpage
14. Для графа на рисунку з вершини 4 запущено алгоритм пошуку в глибину, що будує кістякове дерево. Списки суміжності вершин переглядаються алгоритмом за зростанням міток вершин.\par
\begin{center}
	\adjustimage{max size={0.9\linewidth}{0.9\paperheight}}
	{Graphs/Traversals/323.png}
\end{center}
\par Які з наведених ребер входять до складу отриманого кістякового дерева?\par
1) (5, 6)\par
2) (0, 2)\par
3) (1, 5)\par
4) (1, 4)\par
5) (0, 5)\par
6) (4, 6)\par
{\vskip 15pt} Номери відповідних ребер записати за зростанням через пробіл.

%enditem


\newpage
15. Для графа на рисунку з вершини 0 запущено алгоритм пошуку в ширину, що будує кістякове дерево. Списки суміжності вершин переглядаються алгоритмом за зростанням міток вершин.\par
\begin{center}
	\adjustimage{max size={0.9\linewidth}{0.9\paperheight}}
	{Graphs/Traversals/122.png}
\end{center}
\par Які з наведених ребер входять до складу отриманого кістякового дерева?\par
1) (4, 5)\par
2) (3, 6)\par
3) (1, 3)\par
4) (2, 5)\par
5) (1, 5)\par
6) (1, 6)\par
{\vskip 15pt} Номери відповідних ребер записати за зростанням через пробіл.


%enditem


\newpage
16. 

Рядки журналу через двокрапку містять ім'я користувача (ідентифікатор), час події,тип події, код події, наприклад
\begin{verbatim}
s="username : 2022-05-05 13:26 : ERROR : 404"
\end{verbatim}
у коді 
\begin{verbatim}
res = re.fullmatch('???', s) 
s = ''
code_str = ??? 
\end{verbatim}
чим треба замінити кожен з ???, щоб значенням code\_str був код події? 


\vspace{10mm}
\textbf{Завдання 17} 

Задано програму, що містить код 1-2 класів, можливе успадкування, та клієнтський код.
У наведеному коді невеликі частини закриті. 
Також наведено вихід програми.
Необхідно відновити закриті частини, щоб отриманий код відповідав наведеному виходу програми.


\textbf{Завдання 18-19} 

Кілька запитань за завданнями для самостійної роботи.
\vspace{10mm}
Скоріш за все, що завдань буде трохи менше, а деякі попередні будуть переформульовані в стилі двох останніх.

%enditem

\end {large}}

%end

\newpage
%begin

{{
\begin {large}

\newpage
1. Що буде надруковано за виконання?

\begin{verbatim}
print('R', end=' ')
d=('a','b','c','d','e',0,1,2,4)
print(d[8])
\end{verbatim}

%enditem


\newpage
2. Що буде надруковано за виконання?

\begin{verbatim}
print('R', end=' ')
f='0123456789'
print(f[0:-11:-1])
\end{verbatim}

%enditem



\newpage
3. Що буде надруковано за виконання?

\begin{verbatim}
print('R', end=' ')
a=[10,11,12,13,14]
print(a.pop(4), end=' ')
print(a)
\end{verbatim}

%enditem



\newpage
4. Що буде надруковано за виконання?

\begin{verbatim}
print('R', end=' ')
f=[10,11,12,13,14]
f[1:-1]='12'
print(f)
\end{verbatim}

%enditem



\newpage
5. Укажіть ТИП значення виразу.\par
\begin{verbatim}
{y:None for y in range(-8,3)}
\end{verbatim}
%enditem


\newpage
6. Що буде надруковано за виконання?\par
\begin{verbatim}
try:
    s = {52:0, 55:9, 21:9, 36:5, 52:0}
    s[67] = 5
    for x in s.items():
        print(x, sep=' ', end=' ')
except:
    print('ERR')
\end{verbatim}
%enditem


\newpage
7. Нехай змінна \kw{s} позначає дуже довгу послідовність цілих, по якій можна ітерувати. 

Написати вираз-генератор, що задає підпослідовність її елементів, що є точними квадратами. 

Обмеження: усі елементи шуканої підпослідовності одночасно в пам'яті розміщені бути не можуть. 
%enditem


\newpage
8. Що буде надруковано за виконання?
\begin{verbatim}
c = 0
h = 1

class C:
    c = 2
    
    def __init__(self):
        self.c = 3
        c = 4
        C.c = 5

obj = C()
try:
    print(c, end=' ')
    print(h, end=' ')
    print(C.c, end=' ')
    print(obj.c)
except:
    print('ERR')
\end{verbatim}
%enditem


\newpage
9. Що буде надруковано за виконання?
\begin{verbatim}
class C0:
    def d(self): print(1, end=' ')
    def h(self): print(2, end=' ')

class C1(C0):
    def d(self): print(3, end=' ')

try:
    C1().h()
    C1().d()
except:
    print('ERR')
\end{verbatim}
%enditem


\newpage
10. Що буде надруковано за виконання?
\begin{verbatim}
class A:
    d = 1

    def __h1(self): print(2, end=' ')
    def _h2(self): print(3, end=' ')

    def main(self):
        try: print(self.d, end=' ')
        except: print('E1', end=' ')

        try: self.__h1()
        except: print('E2', end=' ')

        try: self._h2()
        except: print('E3', end=' ')


class B(A):
    d = 4

    def __h1(self): print(5, end=' ')
    def _h2(self): print(6, end=' ')

try: B().main()
except: print('E4', end=' ')
\end{verbatim}
%enditem


\newpage
11. Змінні \kw{next} та \kw{prev} вузла списку позначають наступний та попередній за ним вузол відповідно. Починаючи з голови, у список записано послідовні цілі числа від~$-40$ до~$33$. Нехай змінна \kw{p} позначає вузол, що зберігає значення~$-3$.
Яке значення зберігається у вузлі \kw{p.prev.prev.prev.next.next} ?
%enditem


\newpage
12. 
\begin{center}
	\adjustimage{max size={0.8\linewidth}{0.8\paperheight}}
	{BinTree/trees_exp/49.png}
\end{center}
Указати послідовність міток вершин дерева, зображеного на рис., що утвориться в результаті обходу дерева в прямому порядку. Мітки записувати через пробіл.\par Алгоритм починає роботу з кореня (на рис. це верхня вершина).
%enditem


\newpage
13. 
\begin{center}
	\adjustimage{max size={0.9\linewidth}{0.9\paperheight}}
	{Graphs/AdjList/29.png}
\end{center}
Для графа, зображеного на рис., вказати список суміжності вершини 2.\par
Формат оформлення відповіді:\par
список суміжності виписувати через пробіл на одному рядку, якщо пробілів десь буде більше одного - не важливо, упорядкування вершин у списку також значення не має.
%enditem


\newpage
14. Для графа на рисунку вказати послідовність, в якій алгоритм пошуку в глибину буде відмічати відвідані вершини, якщо списки суміжності вершин переглядаються алгоритмом за зростанням міток вершин та вершини відмічаються на першому вход\-женні у вершину.\par
Пошук розпочинається з вершини 4.\par

\begin{center}
	\adjustimage{max size={0.9\linewidth}{0.9\paperheight}}
	{Graphs/Traversals/260.png}
\end{center}
%enditem


\newpage
15. Для графа на рисунку з вершини 6 запущено алгоритм пошуку в ширину, що будує кістякове дерево. Списки суміжності вершин переглядаються алгоритмом за зростанням міток вершин.\par
\begin{center}
	\adjustimage{max size={0.9\linewidth}{0.9\paperheight}}
	{Graphs/Traversals/67.png}
\end{center}
\par Які з наведених ребер входять до складу отриманого кістякового дерева?\par
1) (1, 5)\par
2) (0, 5)\par
3) (3, 6)\par
4) (3, 4)\par
5) (1, 6)\par
6) (1, 4)\par
{\vskip 15pt} Номери відповідних ребер записати за зростанням через пробіл.


%enditem


\newpage
16. 

Рядки журналу через двокрапку містять ім'я користувача (ідентифікатор), час події,тип події, код події, наприклад
\begin{verbatim}
s="username : 2022-05-05 13:26 : ERROR : 404"
\end{verbatim}
у коді 
\begin{verbatim}
res = re.fullmatch('???', s) 
s = ''
code_str = ??? 
\end{verbatim}
чим треба замінити кожен з ???, щоб значенням code\_str був код події? 


\vspace{10mm}
\textbf{Завдання 17} 

Задано програму, що містить код 1-2 класів, можливе успадкування, та клієнтський код.
У наведеному коді невеликі частини закриті. 
Також наведено вихід програми.
Необхідно відновити закриті частини, щоб отриманий код відповідав наведеному виходу програми.


\textbf{Завдання 18-19} 

Кілька запитань за завданнями для самостійної роботи.
\vspace{10mm}
Скоріш за все, що завдань буде трохи менше, а деякі попередні будуть переформульовані в стилі двох останніх.

%enditem

\end {large}}

%end

\newpage
%begin

{{
\begin {large}

\newpage
1. Що буде надруковано за виконання?

\begin{verbatim}
print('R', end=' ')
g=('a','b','c','d','e',0,1,2,3)
print(g[2])
\end{verbatim}

%enditem


\newpage
2. Що буде надруковано за виконання?

\begin{verbatim}
print('R', end=' ')
h=(0,1,2,3,4,5,6,7,8,9)
print(h[::2])
\end{verbatim}

%enditem



\newpage
3. Що буде надруковано за виконання?

\begin{verbatim}
print('R', end=' ')
a=[10,11,12,13,14]
a.insert(-1, [1,0])
print(a)
\end{verbatim}

%enditem



\newpage
4. Що буде надруковано за виконання?

\begin{verbatim}
print('R', end=' ')
h=['0','1','2','3','4']
h[:1]=[1,0]
print(h)
\end{verbatim}

%enditem



\newpage
5. Укажіть ТИП значення виразу.\par
\begin{verbatim}
(11,2,3)
\end{verbatim}
%enditem


\newpage
6. Що буде надруковано за виконання?\par
\begin{verbatim}
try:
    d = {65:8, 69:2, 72:5, 90:6, 47:8}
    d[80] = 3
    for x in d.values():
        print(x, sep=' ', end=' ')
except:
    print('ERR')
\end{verbatim}
%enditem


\newpage
7. Записати ВИРАЗ, значенням якого є множина, що складається з цілих чисел від 15 до 2002 (включно), що діляться на 3.
%enditem


\newpage
8. Що буде надруковано за виконання?
\begin{verbatim}
b = 0
m = 1

class A:
    m = 2
    
    def __init__(self):
        global m
        self.m = 3
        m = 4
        A.m = 5

obj = A()
try:
    print(b, end=' ')
    print(m, end=' ')
    print(A.m, end=' ')
    print(obj.m)
except:
    print('ERR')
\end{verbatim}
%enditem


\newpage
9. Що буде надруковано за виконання?
\begin{verbatim}
class A0:
    def d(self): print(1, end=' ')
    def f(self): print(2, end=' ')

class A1(A0):
    def d(self): print(3, end=' ')

try:
    A0().f()
    A1().d()
except:
    print('ERR')
\end{verbatim}
%enditem


\newpage
10. Що буде надруковано за виконання?
\begin{verbatim}
class A:
    a = 1

    def _f1(self): print(2, end=' ')
    def _f2(self): print(3, end=' ')

    def main(self):
        try: print(self.a, end=' ')
        except: print('E1', end=' ')

        try: self._f1()
        except: print('E2', end=' ')

        try: self._f2()
        except: print('E3', end=' ')


class B(A):
    a = 4

    def _f1(self): print(5, end=' ')
    def _f2(self): print(6, end=' ')

try: B().main()
except: print('E4', end=' ')
\end{verbatim}
%enditem


\newpage
11. Змінні \kw{next} та \kw{prev} вузла списку позначають наступний та попередній за ним вузол відповідно. Починаючи з голови, у список записано послідовні цілі числа від~$-33$ до~$32$. Нехай змінна \kw{p} позначає вузол, що зберігає значення~$-9$.
Яке значення зберігається у вузлі \kw{p.next.prev.next.prev.next} ?
%enditem


\newpage
12. 
\begin{center}
	\adjustimage{max size={0.8\linewidth}{0.8\paperheight}}
	{BinTree/trees_exp/57.png}
\end{center}
Указати послідовність міток вершин дерева, зображеного на рис., що утвориться в результаті обходу дерева в прямому порядку. Мітки записувати через пробіл.\par Алгоритм починає роботу з кореня (на рис. це верхня вершина).
%enditem


\newpage
13. 
\begin{center}
	\adjustimage{max size={0.9\linewidth}{0.9\paperheight}}
	{Graphs/AdjList/15.png}
\end{center}
Для графа, зображеного на рис., вказати список суміжності вершини 1.\par
Формат оформлення відповіді:\par
список суміжності виписувати через пробіл на одному рядку, якщо пробілів десь буде більше одного - не важливо, упорядкування вершин у списку також значення не має.
%enditem


\newpage
14. Для графа на рисунку вказати послідовність, в якій алгоритм пошуку в глибину буде відмічати відвідані вершини, якщо списки суміжності вершин переглядаються алгоритмом за зростанням міток вершин та вершини відмічаються на першому вход\-женні у вершину.\par
Пошук розпочинається з вершини 0.\par

\begin{center}
	\adjustimage{max size={0.9\linewidth}{0.9\paperheight}}
	{Graphs/Traversals/279.png}
\end{center}
%enditem


\newpage
15. Для графа на рисунку з вершини 3 запущено алгоритм пошуку в ширину, що будує кістякове дерево. Списки суміжності вершин переглядаються алгоритмом за зростанням міток вершин.\par
\begin{center}
	\adjustimage{max size={0.9\linewidth}{0.9\paperheight}}
	{Graphs/Traversals/24.png}
\end{center}
\par Які з наведених ребер входять до складу отриманого кістякового дерева?\par
1) (5, 6)\par
2) (2, 5)\par
3) (1, 6)\par
4) (4, 6)\par
5) (0, 6)\par
6) (1, 3)\par
{\vskip 15pt} Номери відповідних ребер записати за зростанням через пробіл.


%enditem


\newpage
16. 

Рядки журналу через двокрапку містять ім'я користувача (ідентифікатор), час події,тип події, код події, наприклад
\begin{verbatim}
s="username : 2022-05-05 13:26 : ERROR : 404"
\end{verbatim}
у коді 
\begin{verbatim}
res = re.fullmatch('???', s) 
s = ''
code_str = ??? 
\end{verbatim}
чим треба замінити кожен з ???, щоб значенням code\_str був код події? 


\vspace{10mm}
\textbf{Завдання 17} 

Задано програму, що містить код 1-2 класів, можливе успадкування, та клієнтський код.
У наведеному коді невеликі частини закриті. 
Також наведено вихід програми.
Необхідно відновити закриті частини, щоб отриманий код відповідав наведеному виходу програми.


\textbf{Завдання 18-19} 

Кілька запитань за завданнями для самостійної роботи.
\vspace{10mm}
Скоріш за все, що завдань буде трохи менше, а деякі попередні будуть переформульовані в стилі двох останніх.

%enditem

\end {large}}

%end

\newpage
%begin

{{
\begin {large}

\newpage
1. Що буде надруковано за виконання?

\begin{verbatim}
print('R', end=' ')
a=('a','b','c','d','e',0,1,2)
print(a[-6])
\end{verbatim}

%enditem


\newpage
2. Що буде надруковано за виконання?

\begin{verbatim}
print('R', end=' ')
c=('0','1','2','3','4','5','6','7','8','9')
print(c[12::-3])
\end{verbatim}

%enditem



\newpage
3. Що буде надруковано за виконання?

\begin{verbatim}
print('R', end=' ')
g=['0','1','2','3','4']
g.append('13')
print(g)
\end{verbatim}

%enditem



\newpage
4. Що буде надруковано за виконання?

\begin{verbatim}
print('R', end=' ')
e=[10,11,12,13,14]
del e[-6:]
print(e)
\end{verbatim}

%enditem



\newpage
5. Укажіть ТИП значення виразу.\par
\begin{verbatim}
[{1,2}, 5]
\end{verbatim}
%enditem


\newpage
6. Що буде надруковано за виконання?\par
\begin{verbatim}
try:
    d = {78:2, 38:2, 19:5, 19:0}
    d[51] = 2
    for x in d.items():
        print(*x, sep=' ', end=' ')
except:
    print('ERR')
\end{verbatim}
%enditem


\newpage
7. Нехай змінна \kw{s} позначає послідовність цілих, по якій можна ітерувати. 

Написати вираз, що перевіряє, чи правда, що послідовність  містить число 35.
%enditem


\newpage
8. Що буде надруковано за виконання?
%Оригинал был утерян. Требует отладки!
\begin{verbatim}
__a = 0

class D:
    __a = 1
    
    def __init__(self):
        self.__a = 2
        __a = 3
        D.__a = 4

obj = D()
D.__a = 5
try:
    print(__a, end=' ')
    print(obj.__a, end=' ')
    print(D.__a, end=' ')
    print(obj.__a, end=' ')
    print(obj._D__a)
except:
    print('ERR')
\end{verbatim}
%enditem


\newpage
9. Що буде надруковано за виконання?
\begin{verbatim}
class D0:
    def b(self): print(1, end=' ')
    def h(self): print(2, end=' ')

class D1(D0):
    def h(self): print(3, end=' ')

try:
    D1().b()
    D1().h()
except:
    print('ERR')
\end{verbatim}
%enditem


\newpage
10. Що буде надруковано за виконання?
\begin{verbatim}
class A:
    c = 1

    def f1(self): print(2, end=' ')
    def _f2(self): print(3, end=' ')

    def main(self):
        try: print(self.c, end=' ')
        except: print('E1', end=' ')

        try: self.f1()
        except: print('E2', end=' ')

        try: self._f2()
        except: print('E3', end=' ')


class B(A):
    c = 4

    def f1(self): print(5, end=' ')
    def _f2(self): print(6, end=' ')

try: B().main()
except: print('E4', end=' ')
\end{verbatim}
%enditem


\newpage
11. Змінні \kw{next} та \kw{prev} вузла списку позначають наступний та попередній за ним вузол відповідно. Починаючи з голови, у список записано послідовні цілі числа від~$-34$ до~$19$. Нехай змінна \kw{p} позначає вузол, що зберігає значення~$7$.
Яке значення зберігається у вузлі \kw{p.prev.next.prev.next.prev} ?
%enditem


\newpage
12. 
\begin{center}
	\adjustimage{max size={0.8\linewidth}{0.8\paperheight}}
	{BinTree/trees_exp/47.png}
\end{center}
Указати послідовність міток вершин дерева, зображеного на рис., що утвориться в результаті обходу дерева в оберненому порядку. Мітки записувати через пробіл.\par Алгоритм починає роботу з кореня (на рис. це верхня вершина).
%enditem


\newpage
13. 
\begin{center}
	\adjustimage{max size={0.9\linewidth}{0.9\paperheight}}
	{Graphs/AdjList/184.png}
\end{center}
Для графа, зображеного на рис., вказати список суміжності вершини 6.\par
Формат оформлення відповіді:\par
список суміжності виписувати через пробіл на одному рядку, якщо пробілів десь буде більше одного - не важливо, упорядкування вершин у списку також значення не має.
%enditem


\newpage
14. Для графа на рисунку вказати послідовність, в якій алгоритм пошуку в глибину буде відмічати відвідані вершини, якщо списки суміжності вершин переглядаються алгоритмом за зростанням міток вершин та вершини відмічаються на першому вход\-женні у вершину.\par
Пошук розпочинається з вершини 0.\par

\begin{center}
	\adjustimage{max size={0.9\linewidth}{0.9\paperheight}}
	{Graphs/Traversals/0.png}
\end{center}
%enditem


\newpage
15. Для графа на рисунку з вершини 4 запущено алгоритм пошуку в ширину, що будує кістякове дерево. Списки суміжності вершин переглядаються алгоритмом за зростанням міток вершин.\par
\begin{center}
	\adjustimage{max size={0.9\linewidth}{0.9\paperheight}}
	{Graphs/Traversals/131.png}
\end{center}
\par Які з наведених ребер входять до складу отриманого кістякового дерева?\par
1) (3, 4)\par
2) (2, 6)\par
3) (3, 6)\par
4) (4, 6)\par
5) (2, 5)\par
6) (2, 3)\par
{\vskip 15pt} Номери відповідних ребер записати за зростанням через пробіл.


%enditem


\newpage
16. 

Рядки журналу через двокрапку містять ім'я користувача (ідентифікатор), час події,тип події, код події, наприклад
\begin{verbatim}
s="username : 2022-05-05 13:26 : ERROR : 404"
\end{verbatim}
у коді 
\begin{verbatim}
res = re.fullmatch('???', s) 
s = ''
code_str = ??? 
\end{verbatim}
чим треба замінити кожен з ???, щоб значенням code\_str був код події? 


\vspace{10mm}
\textbf{Завдання 17} 

Задано програму, що містить код 1-2 класів, можливе успадкування, та клієнтський код.
У наведеному коді невеликі частини закриті. 
Також наведено вихід програми.
Необхідно відновити закриті частини, щоб отриманий код відповідав наведеному виходу програми.


\textbf{Завдання 18-19} 

Кілька запитань за завданнями для самостійної роботи.
\vspace{10mm}
Скоріш за все, що завдань буде трохи менше, а деякі попередні будуть переформульовані в стилі двох останніх.

%enditem

\end {large}}

%end

\newpage
%begin

{{
\begin {large}

\newpage
1. Що буде надруковано за виконання?

\begin{verbatim}
print('R', end=' ')
h=('a','b','c','d','e',0,1,2,4)
print(h[12])
\end{verbatim}

%enditem


\newpage
2. Що буде надруковано за виконання?

\begin{verbatim}
print('R', end=' ')
f=['0','1','2','3','4','5','6','7','8','9']
print(f[:5:-1])
\end{verbatim}

%enditem



\newpage
3. Що буде надруковано за виконання?

\begin{verbatim}
print('R', end=' ')
f=[10,11,12,13,14]
f.extend([4,5])
print(f)
\end{verbatim}

%enditem



\newpage
4. Що буде надруковано за виконання?

\begin{verbatim}
print('R', end=' ')
a=['0','1','2','3','4']
a[-3:-2]=()
print(a)
\end{verbatim}

%enditem



\newpage
5. Укажіть ТИП значення виразу.\par
\begin{verbatim}
{y:y*y for y in range(-7,7)}
\end{verbatim}
%enditem


\newpage
6. Що буде надруковано за виконання?\par
\begin{verbatim}
try:
    t = {19:4, 16:0, 65:0, 77:6, 65:4}
    t[48] = 7
    for x in t :
        print(x, sep=' ', end=' ')
except:
    print('ERR')
\end{verbatim}
%enditem


\newpage
7. Записати ВИРАЗ, значенням якого є словник, ключами якого є цілі числа від 15 до 2005 (включно), що діляться на 7, а ключам відповідають їх квадрати 
%enditem


\newpage
8. Що буде надруковано за виконання?
%Оригинал был утерян. Требует отладки!
\begin{verbatim}
d = 0

class A:
    d = 1
    
    def __init__(self):
        self.d = 2
        d = 3
        A.d = 4

obj = A()
obj.d = 5
try:
    print(d, end=' ')
    print(obj.d, end=' ')
    print(A.d, end=' ')
    print(obj.d, end=' ')
    print(obj._A__d)
except:
    print('ERR')
\end{verbatim}
%enditem


\newpage
9. Що буде надруковано за виконання?
\begin{verbatim}
class C0:
    def a(self): print(1, end=' ')
    def f(self): print(2, end=' ')

class C1(C0):
    def a(self): print(3, end=' ')

try:
    C1().f()
    C1().a()
except:
    print('ERR')
\end{verbatim}
%enditem


\newpage
10. Що буде надруковано за виконання?
\begin{verbatim}
class A:
    c = 1

    def _h1(self): print(2, end=' ')
    def h2(self): print(3, end=' ')

    def main(self):
        try: print(self.c, end=' ')
        except: print('E1', end=' ')

        try: self._h1()
        except: print('E2', end=' ')

        try: self.h2()
        except: print('E3', end=' ')


class B(A):
    c = 4

    def _h1(self): print(5, end=' ')
    def h2(self): print(6, end=' ')

try: B().main()
except: print('E4', end=' ')
\end{verbatim}
%enditem


\newpage
11. Змінні \kw{next} та \kw{prev} вузла списку позначають наступний та попередній за ним вузол відповідно. Починаючи з голови, у список записано послідовні цілі числа від~$-42$ до~$18$. Нехай змінна \kw{p} позначає вузол, що зберігає значення~$9$.
Яке значення зберігається у вузлі \kw{p.next.prev.next.prev.prev} ?
%enditem


\newpage
12. 
\begin{center}
	\adjustimage{max size={0.8\linewidth}{0.8\paperheight}}
	{BinTree/trees_exp/63.png}
\end{center}
Указати послідовність міток вершин дерева, зображеного на рис., що утвориться в результаті обходу дерева в оберненому порядку. Мітки записувати через пробіл.\par Алгоритм починає роботу з кореня (на рис. це верхня вершина).
%enditem


\newpage
13. 
\begin{center}
	\adjustimage{max size={0.9\linewidth}{0.9\paperheight}}
	{Graphs/AdjList/35.png}
\end{center}
Для графа, зображеного на рис., вказати список суміжності вершини 2.\par
Формат оформлення відповіді:\par
список суміжності виписувати через пробіл на одному рядку, якщо пробілів десь буде більше одного - не важливо, упорядкування вершин у списку також значення не має.
%enditem


\newpage
14. Для графа на рисунку вказати послідовність, в якій алгоритм пошуку в глибину буде відмічати відвідані вершини, якщо списки суміжності вершин переглядаються алгоритмом за зростанням міток вершин та вершини відмічаються на першому вход\-женні у вершину.\par
Пошук розпочинається з вершини 0.\par

\begin{center}
	\adjustimage{max size={0.9\linewidth}{0.9\paperheight}}
	{Graphs/Traversals/123.png}
\end{center}
%enditem


\newpage
15. Для графа на рисунку з вершини 4 запущено алгоритм пошуку в ширину, що будує кістякове дерево. Списки суміжності вершин переглядаються алгоритмом за зростанням міток вершин.\par
\begin{center}
	\adjustimage{max size={0.9\linewidth}{0.9\paperheight}}
	{Graphs/Traversals/151.png}
\end{center}
\par Які з наведених ребер входять до складу отриманого кістякового дерева?\par
1) (5, 6)\par
2) (2, 5)\par
3) (4, 5)\par
4) (3, 4)\par
5) (0, 4)\par
6) (2, 3)\par
{\vskip 15pt} Номери відповідних ребер записати за зростанням через пробіл.


%enditem


\newpage
16. 

Рядки журналу через двокрапку містять ім'я користувача (ідентифікатор), час події,тип події, код події, наприклад
\begin{verbatim}
s="username : 2022-05-05 13:26 : ERROR : 404"
\end{verbatim}
у коді 
\begin{verbatim}
res = re.fullmatch('???', s) 
s = ''
code_str = ??? 
\end{verbatim}
чим треба замінити кожен з ???, щоб значенням code\_str був код події? 


\vspace{10mm}
\textbf{Завдання 17} 

Задано програму, що містить код 1-2 класів, можливе успадкування, та клієнтський код.
У наведеному коді невеликі частини закриті. 
Також наведено вихід програми.
Необхідно відновити закриті частини, щоб отриманий код відповідав наведеному виходу програми.


\textbf{Завдання 18-19} 

Кілька запитань за завданнями для самостійної роботи.
\vspace{10mm}
Скоріш за все, що завдань буде трохи менше, а деякі попередні будуть переформульовані в стилі двох останніх.

%enditem

\end {large}}

%end

\newpage
%begin

{{
\begin {large}

\newpage
1. Що буде надруковано за виконання?

\begin{verbatim}
print('R', end=' ')
f=('a','b','c','d','e',0,1,2,3,4)
print(f[-5])
\end{verbatim}

%enditem


\newpage
2. Що буде надруковано за виконання?

\begin{verbatim}
print('R', end=' ')
g=(0,1,2,3,4,5,6,7,8,9)
print(g[-13:13:-2])
\end{verbatim}

%enditem



\newpage
3. Що буде надруковано за виконання?

\begin{verbatim}
print('R', end=' ')
b=['0','1','2','3','4']
b.remove('8')
print(b)
\end{verbatim}

%enditem



\newpage
4. Що буде надруковано за виконання?

\begin{verbatim}
print('R', end=' ')
f=[10,11,12,13,14]
f[:6]=[4,5]
print(f)
\end{verbatim}

%enditem



\newpage
5. Укажіть ТИП значення виразу.\par
\begin{verbatim}
{x*x:x for x in range(-7,7)}
\end{verbatim}
%enditem


\newpage
6. Що буде надруковано за виконання?\par
\begin{verbatim}
try:
    s = {73:5, 38:4, 26:1, 38:1}
    s[40] = 3
    for x in s.keys():
        print(x, sep=' ', end=' ')
except:
    print('ERR')
\end{verbatim}
%enditem


\newpage
7. Нехай змінна \kw{s} позначає послідовність цілих, по якій можна ітерувати.

Написати вираз, що перевіряє, чи правда, що всі елементи послідовності є додатни\-ми.
%enditem


\newpage
8. Що буде надруковано за виконання?
%Оригинал был утерян. Требует отладки!
\begin{verbatim}
__c = 0

class C:
    __c = 1
    
    def __init__(self):
        self.__c = 2
        __c = 3
        C.__c = 4

obj = C()
C.__c = 5
try:
    print(__c, end=' ')
    print(obj.__c, end=' ')
    print(C.__c, end=' ')
    print(obj.__c, end=' ')
    print(obj._C__c)
except:
    print('ERR')
\end{verbatim}
%enditem


\newpage
9. Що буде надруковано за виконання?
\begin{verbatim}
class B0:
    def c(self): print(1, end=' ')
    def g(self): print(2, end=' ')

class B1(B0):
    def c(self): print(3, end=' ')

try:
    B1().g()
    B0().c()
except:
    print('ERR')
\end{verbatim}
%enditem


\newpage
10. Що буде надруковано за виконання?
\begin{verbatim}
class A:
    b = 1

    def __init__(self):
        a = 2
        a = 3

class B(A):
    b = 4
    c = 5

    def __init__(self):
        super().__init__()
        self.c = 6

try:
    obj = B()
    obj.b = obj.a + 10
    B.b = 13
except: print('EE', end=' ')

try: print(obj.b, end= ' ')
except: print('E1', end=' ')

try: print(A.b, end= ' ')
except: print('E2', end=' ')

try: print(B.b, end= ' ')
except: print('E3', end=' ')
\end{verbatim}
%enditem


\newpage
11. Змінні \kw{next} та \kw{prev} вузла списку позначають наступний та попередній за ним вузол відповідно. Починаючи з голови, у список записано послідовні цілі числа від~$-28$ до~$47$. Нехай змінна \kw{p} позначає вузол, що зберігає значення~$-1$.
Яке значення зберігається у вузлі \kw{p.prev.next.prev.next.next} ?
%enditem


\newpage
12. 
\begin{center}
	\adjustimage{max size={0.8\linewidth}{0.8\paperheight}}
	{BinTree/trees_exp/92.png}
\end{center}
Указати послідовність міток вершин дерева, зображеного на рис., що утвориться в результаті обходу дерева в прямому порядку. Мітки записувати через пробіл.\par Алгоритм починає роботу з кореня (на рис. це верхня вершина).
%enditem


\newpage
13. 
\begin{center}
	\adjustimage{max size={0.9\linewidth}{0.9\paperheight}}
	{Graphs/AdjList/154.png}
\end{center}
Для графа, зображеного на рис., вказати список суміжності вершини 2.\par
Формат оформлення відповіді:\par
список суміжності виписувати через пробіл на одному рядку, якщо пробілів десь буде більше одного - не важливо, упорядкування вершин у списку також значення не має.
%enditem


\newpage
14. Для графа на рисунку вказати послідовність, в якій алгоритм пошуку в глибину буде відмічати відвідані вершини, якщо списки суміжності вершин переглядаються алгоритмом за зростанням міток вершин та вершини відмічаються на першому вход\-женні у вершину.\par
Пошук розпочинається з вершини 4.\par

\begin{center}
	\adjustimage{max size={0.9\linewidth}{0.9\paperheight}}
	{Graphs/Traversals/231.png}
\end{center}
%enditem


\newpage
15. Для графа на рисунку з вершини 6 запущено алгоритм пошуку в ширину, що будує кістякове дерево. Списки суміжності вершин переглядаються алгоритмом за зростанням міток вершин.\par
\begin{center}
	\adjustimage{max size={0.9\linewidth}{0.9\paperheight}}
	{Graphs/Traversals/300.png}
\end{center}
\par Які з наведених ребер входять до складу отриманого кістякового дерева?\par
1) (0, 1)\par
2) (3, 5)\par
3) (3, 6)\par
4) (0, 2)\par
5) (0, 3)\par
6) (2, 4)\par
{\vskip 15pt} Номери відповідних ребер записати за зростанням через пробіл.


%enditem


\newpage
16. 

Рядки журналу через двокрапку містять ім'я користувача (ідентифікатор), час події,тип події, код події, наприклад
\begin{verbatim}
s="username : 2022-05-05 13:26 : ERROR : 404"
\end{verbatim}
у коді 
\begin{verbatim}
res = re.fullmatch('???', s) 
s = ''
code_str = ??? 
\end{verbatim}
чим треба замінити кожен з ???, щоб значенням code\_str був код події? 


\vspace{10mm}
\textbf{Завдання 17} 

Задано програму, що містить код 1-2 класів, можливе успадкування, та клієнтський код.
У наведеному коді невеликі частини закриті. 
Також наведено вихід програми.
Необхідно відновити закриті частини, щоб отриманий код відповідав наведеному виходу програми.


\textbf{Завдання 18-19} 

Кілька запитань за завданнями для самостійної роботи.
\vspace{10mm}
Скоріш за все, що завдань буде трохи менше, а деякі попередні будуть переформульовані в стилі двох останніх.

%enditem

\end {large}}

%end

\newpage
%begin

{{
\begin {large}

\newpage
1. Що буде надруковано за виконання?

\begin{verbatim}
print('R', end=' ')
c=('a','b','c','d','e',0,1,2,3)
print(c[-10])
\end{verbatim}

%enditem


\newpage
2. Що буде надруковано за виконання?

\begin{verbatim}
print('R', end=' ')
d='0123456789'
print(d[-4:-9:-3])
\end{verbatim}

%enditem



\newpage
3. Що буде надруковано за виконання?

\begin{verbatim}
print('R', end=' ')
e=['0','1','2','3','4']
print(e.index(2), end=' ')
\end{verbatim}

%enditem



\newpage
4. Що буде надруковано за виконання?

\begin{verbatim}
print('R', end=' ')
c=[10,11,12,13,14]
del c[-7:]
print(c)
\end{verbatim}

%enditem



\newpage
5. Укажіть ТИП значення виразу.\par
\begin{verbatim}
{x:x for x in range(7)}
\end{verbatim}
%enditem


\newpage
6. Що буде надруковано за виконання?\par
\begin{verbatim}
try:
    s = {85:1, 65:6, 57:1, 11:9, 11:5}
    s[40] = 6
    for x in s.keys():
        print(x, sep=' ', end=' ')
except:
    print('ERR')
\end{verbatim}
%enditem


\newpage
7. Задано словник \kw{d} з цілими ключами та значеннями. Записати ВИРАЗ, значенням якого є множина, що складається  з тих ключів словника, яким відповідають значення, що діляться на 7.
%enditem


\newpage
8. Що буде надруковано за виконання?
%Оригинал был утерян. Требует отладки!
\begin{verbatim}
c = 0

class C:
    c = 1
    
    def __init__(self):
        self.c = 2
        c = 3
        C.c = 4

obj = C()
C.c = 5
try:
    print(c, end=' ')
    print(obj.c, end=' ')
    print(C.c, end=' ')
    print(obj.c, end=' ')
    print(obj._C__c)
except:
    print('ERR')
\end{verbatim}
%enditem


\newpage
9. Що буде надруковано за виконання?
\begin{verbatim}
class A0:
    def c(self): print(1, end=' ')
    def h(self): print(2, end=' ')

class A1(A0):
    def h(self): print(3, end=' ')

try:
    A1().c()
    A0().h()
except:
    print('ERR')
\end{verbatim}
%enditem


\newpage
10. Що буде надруковано за виконання?
\begin{verbatim}
class A:
    a = 1

    def __init__(self):
        b = 2
        a = 3

class B(A):
    b = 4
    c = 5

    def __init__(self):
        super().__init__()
        self.c = 6

try:
    obj = B()
    obj.c = obj.a + 10
    B.c = 13
except: print('EE', end=' ')

try: print(obj.c, end= ' ')
except: print('E1', end=' ')

try: print(A.c, end= ' ')
except: print('E2', end=' ')

try: print(B.c, end= ' ')
except: print('E3', end=' ')
\end{verbatim}
%enditem


\newpage
11. Змінні \kw{next} та \kw{prev} вузла списку позначають наступний та попередній за ним вузол відповідно. Починаючи з голови, у список записано послідовні цілі числа від~$-53$ до~$46$. Нехай змінна \kw{p} позначає вузол, що зберігає значення~$4$.
Яке значення зберігається у вузлі \kw{p.prev.next.prev.next.prev} ?
%enditem


\newpage
12. 
\begin{center}
	\adjustimage{max size={0.8\linewidth}{0.8\paperheight}}
	{BinTree/trees_exp/79.png}
\end{center}
Указати послідовність міток вершин дерева, зображеного на рис., що утвориться в результаті обходу дерева в оберненому порядку. Мітки записувати через пробіл.\par Алгоритм починає роботу з кореня (на рис. це верхня вершина).
%enditem


\newpage
13. 
\begin{center}
	\adjustimage{max size={0.9\linewidth}{0.9\paperheight}}
	{Graphs/AdjList/231.png}
\end{center}
Для графа, зображеного на рис., вказати список суміжності вершини 5.\par
Формат оформлення відповіді:\par
список суміжності виписувати через пробіл на одному рядку, якщо пробілів десь буде більше одного - не важливо, упорядкування вершин у списку також значення не має.
%enditem


\newpage
14. Для графа на рисунку вказати послідовність, в якій алгоритм пошуку в глибину буде відмічати відвідані вершини, якщо списки суміжності вершин переглядаються алгоритмом за зростанням міток вершин та вершини відмічаються на першому вход\-женні у вершину.\par
Пошук розпочинається з вершини 1.\par

\begin{center}
	\adjustimage{max size={0.9\linewidth}{0.9\paperheight}}
	{Graphs/Traversals/171.png}
\end{center}
%enditem


\newpage
15. Для графа на рисунку з вершини 2 запущено алгоритм пошуку в ширину, що будує кістякове дерево. Списки суміжності вершин переглядаються алгоритмом за зростанням міток вершин.\par
\begin{center}
	\adjustimage{max size={0.9\linewidth}{0.9\paperheight}}
	{Graphs/Traversals/164.png}
\end{center}
\par Які з наведених ребер входять до складу отриманого кістякового дерева?\par
1) (1, 2)\par
2) (2, 3)\par
3) (2, 5)\par
4) (1, 5)\par
5) (4, 6)\par
6) (2, 6)\par
{\vskip 15pt} Номери відповідних ребер записати за зростанням через пробіл.


%enditem


\newpage
16. 

Рядки журналу через двокрапку містять ім'я користувача (ідентифікатор), час події,тип події, код події, наприклад
\begin{verbatim}
s="username : 2022-05-05 13:26 : ERROR : 404"
\end{verbatim}
у коді 
\begin{verbatim}
res = re.fullmatch('???', s) 
s = ''
code_str = ??? 
\end{verbatim}
чим треба замінити кожен з ???, щоб значенням code\_str був код події? 


\vspace{10mm}
\textbf{Завдання 17} 

Задано програму, що містить код 1-2 класів, можливе успадкування, та клієнтський код.
У наведеному коді невеликі частини закриті. 
Також наведено вихід програми.
Необхідно відновити закриті частини, щоб отриманий код відповідав наведеному виходу програми.


\textbf{Завдання 18-19} 

Кілька запитань за завданнями для самостійної роботи.
\vspace{10mm}
Скоріш за все, що завдань буде трохи менше, а деякі попередні будуть переформульовані в стилі двох останніх.

%enditem

\end {large}}

%end

\newpage
%begin

{{
\begin {large}

\newpage
1. Що буде надруковано за виконання?

\begin{verbatim}
print('R', end=' ')
d=('a','b','c','d','e',0,1,2)
print(d[-13])
\end{verbatim}

%enditem


\newpage
2. Що буде надруковано за виконання?

\begin{verbatim}
print('R', end=' ')
e=('0','1','2','3','4','5','6','7','8','9')
print(e[12::3])
\end{verbatim}

%enditem



\newpage
3. Що буде надруковано за виконання?

\begin{verbatim}
print('R', end=' ')
h=[10,11,12,13,14]
print(h.pop(-3), end=' ')
print(h)
\end{verbatim}

%enditem



\newpage
4. Що буде надруковано за виконання?

\begin{verbatim}
print('R', end=' ')
b=['0','1','2','3','4']
b[-5:0]='12'
print(b)
\end{verbatim}

%enditem



\newpage
5. Укажіть ТИП значення виразу.\par
\begin{verbatim}
{21,2,4}
\end{verbatim}
%enditem


\newpage
6. Що буде надруковано за виконання?\par
\begin{verbatim}
try:
    d = {14:3, 78:4, 62:9, 87:7, 78:0}
    d[14] = 2
    for x in d.values():
        print(x, sep=' ', end=' ')
except:
    print('ERR')
\end{verbatim}
%enditem


\newpage
7. Записати ВИРАЗ, значенням якого є множина, що складається з квадратів цілих чисел від 10 до 2002 (включно), що не діляться на 2.
%enditem


\newpage
8. Що буде надруковано за виконання?
%Оригинал был утерян. Требует отладки!
\begin{verbatim}
_a = 0

class D:
    _a = 1
    
    def __init__(self):
        self._a = 2
        _a = 3
        D._a = 4

obj = D()
obj._a = 5
try:
    print(_a, end=' ')
    print(obj._a, end=' ')
    print(D._a, end=' ')
    print(obj._a, end=' ')
    print(obj._D__a)
except:
    print('ERR')
\end{verbatim}
%enditem


\newpage
9. Що буде надруковано за виконання?
\begin{verbatim}
class D0:
    def a(self): print(1, end=' ')
    def f(self): print(2, end=' ')

class D1(D0):
    def f(self): print(3, end=' ')

try:
    D1().a()
    D1().f()
except:
    print('ERR')
\end{verbatim}
%enditem


\newpage
10. Що буде надруковано за виконання?
\begin{verbatim}
class A:
    d = 1

    def g1(self): print(2, end=' ')
    def _g2(self): print(3, end=' ')

    def main(self):
        try: print(self.d, end=' ')
        except: print('E1', end=' ')

        try: self.g1()
        except: print('E2', end=' ')

        try: self._g2()
        except: print('E3', end=' ')


class B(A):
    d = 4

    def g1(self): print(5, end=' ')
    def _g2(self): print(6, end=' ')

try: B().main()
except: print('E4', end=' ')
\end{verbatim}
%enditem


\newpage
11. Змінні \kw{next} та \kw{prev} вузла списку позначають наступний та попередній за ним вузол відповідно. Починаючи з голови, у список записано послідовні цілі числа від~$-50$ до~$49$. Нехай змінна \kw{p} позначає вузол, що зберігає значення~$2$.
Яке значення зберігається у вузлі \kw{p.next.prev.next.prev.next} ?
%enditem


\newpage
12. 
\begin{center}
	\adjustimage{max size={0.8\linewidth}{0.8\paperheight}}
	{BinTree/trees_exp/4.png}
\end{center}
Указати послідовність міток вершин дерева, зображеного на рис., що утвориться в результаті обходу дерева в прямому порядку. Мітки записувати через пробіл.\par Алгоритм починає роботу з кореня (на рис. це верхня вершина).
%enditem


\newpage
13. 
\begin{center}
	\adjustimage{max size={0.9\linewidth}{0.9\paperheight}}
	{Graphs/AdjList/155.png}
\end{center}
Для графа, зображеного на рис., вказати список суміжності вершини 3.\par
Формат оформлення відповіді:\par
список суміжності виписувати через пробіл на одному рядку, якщо пробілів десь буде більше одного - не важливо, упорядкування вершин у списку також значення не має.
%enditem


\newpage
14. Для графа на рисунку вказати послідовність, в якій алгоритм пошуку в глибину буде відмічати відвідані вершини, якщо списки суміжності вершин переглядаються алгоритмом за зростанням міток вершин та вершини відмічаються на першому вход\-женні у вершину.\par
Пошук розпочинається з вершини 6.\par

\begin{center}
	\adjustimage{max size={0.9\linewidth}{0.9\paperheight}}
	{Graphs/Traversals/144.png}
\end{center}
%enditem


\newpage
15. Для графа на рисунку з вершини 4 запущено алгоритм пошуку в ширину, що будує кістякове дерево. Списки суміжності вершин переглядаються алгоритмом за зростанням міток вершин.\par
\begin{center}
	\adjustimage{max size={0.9\linewidth}{0.9\paperheight}}
	{Graphs/Traversals/189.png}
\end{center}
\par Які з наведених ребер входять до складу отриманого кістякового дерева?\par
1) (0, 1)\par
2) (0, 2)\par
3) (3, 6)\par
4) (0, 5)\par
5) (0, 4)\par
6) (2, 4)\par
{\vskip 15pt} Номери відповідних ребер записати за зростанням через пробіл.


%enditem


\newpage
16. 

Рядки журналу через двокрапку містять ім'я користувача (ідентифікатор), час події,тип події, код події, наприклад
\begin{verbatim}
s="username : 2022-05-05 13:26 : ERROR : 404"
\end{verbatim}
у коді 
\begin{verbatim}
res = re.fullmatch('???', s) 
s = ''
code_str = ??? 
\end{verbatim}
чим треба замінити кожен з ???, щоб значенням code\_str був код події? 


\vspace{10mm}
\textbf{Завдання 17} 

Задано програму, що містить код 1-2 класів, можливе успадкування, та клієнтський код.
У наведеному коді невеликі частини закриті. 
Також наведено вихід програми.
Необхідно відновити закриті частини, щоб отриманий код відповідав наведеному виходу програми.


\textbf{Завдання 18-19} 

Кілька запитань за завданнями для самостійної роботи.
\vspace{10mm}
Скоріш за все, що завдань буде трохи менше, а деякі попередні будуть переформульовані в стилі двох останніх.

%enditem

\end {large}}

%end

\newpage
%begin

{{
\begin {large}

\newpage
1. Що буде надруковано за виконання?

\begin{verbatim}
print('R', end=' ')
e=('a','b','c','d','e',0,1)
print(e[-2])
\end{verbatim}

%enditem


\newpage
2. Що буде надруковано за виконання?

\begin{verbatim}
print('R', end=' ')
b=[0,1,2,3,4,5,6,7,8,9]
print(b[2:-5:2])
\end{verbatim}

%enditem



\newpage
3. Що буде надруковано за виконання?

\begin{verbatim}
print('R', end=' ')
c=['0','1','2','3','4']
print(c.index('9'), end=' ')
\end{verbatim}

%enditem



\newpage
4. Що буде надруковано за виконання?

\begin{verbatim}
print('R', end=' ')
d=[10,11,12,13,14]
d[0:-3]='13'
print(d)
\end{verbatim}

%enditem



\newpage
5. Укажіть ТИП значення виразу.\par
\begin{verbatim}
[7,2,3,4]
\end{verbatim}
%enditem


\newpage
6. Що буде надруковано за виконання?\par
\begin{verbatim}
try:
    t = {61:1, 56:5, 50:6, 73:1, 56:7}
    t[59] = 4
    for x in t.items():
        print(*x, sep=' ', end=' ')
except:
    print('ERR')
\end{verbatim}
%enditem


\newpage
7. Нехай змінна \kw{s} позначає дуже довгу послідовність, по якій можна ітерувати. 

Написати вираз-генератор, що задає підпослідовність її елементів, що є рядками типу \kw{str}.

Обмеження: усі елементи шуканої підпослідовності одночасно в пам'яті розміщені бути не можуть. 
%enditem


\newpage
8. Що буде надруковано за виконання?
%Оригинал был утерян. Требует отладки!
\begin{verbatim}
__b = 0

class B:
    __b = 1
    
    def __init__(self):
        self.__b = 2
        __b = 3
        B.__b = 4

obj = B()
obj.__b = 5
try:
    print(__b, end=' ')
    print(obj.__b, end=' ')
    print(B.__b, end=' ')
    print(obj.__b, end=' ')
    print(obj._B__b)
except:
    print('ERR')
\end{verbatim}
%enditem


\newpage
9. Що буде надруковано за виконання?
\begin{verbatim}
class C0:
    def d(self): print(1, end=' ')
    def g(self): print(2, end=' ')

class C1(C0):
    def d(self): print(3, end=' ')

try:
    C1().g()
    C1().d()
except:
    print('ERR')
\end{verbatim}
%enditem


\newpage
10. Що буде надруковано за виконання?
\begin{verbatim}
class A:
    b = 1

    def _f1(self): print(2, end=' ')
    def __f2(self): print(3, end=' ')

    def main(self):
        try: print(self.b, end=' ')
        except: print('E1', end=' ')

        try: self._f1()
        except: print('E2', end=' ')

        try: self.__f2()
        except: print('E3', end=' ')


class B(A):
    b = 4

    def _f1(self): print(5, end=' ')
    def __f2(self): print(6, end=' ')

try: B().main()
except: print('E4', end=' ')
\end{verbatim}
%enditem


\newpage
11. Змінні \kw{next} та \kw{prev} вузла списку позначають наступний та попередній за ним вузол відповідно. Починаючи з голови, у список записано послідовні цілі числа від~$-66$ до~$25$. Нехай змінна \kw{p} позначає вузол, що зберігає значення~$1$.
Яке значення зберігається у вузлі \kw{p.next.next.prev.next.prev} ?
%enditem


\newpage
12. 
\begin{center}
	\adjustimage{max size={0.8\linewidth}{0.8\paperheight}}
	{BinTree/trees_exp/83.png}
\end{center}
Указати послідовність міток вершин дерева, зображеного на рис., що утвориться в результаті обходу дерева в оберненому порядку. Мітки записувати через пробіл.\par Алгоритм починає роботу з кореня (на рис. це верхня вершина).
%enditem


\newpage
13. 
\begin{center}
	\adjustimage{max size={0.9\linewidth}{0.9\paperheight}}
	{Graphs/AdjList/241.png}
\end{center}
Для графа, зображеного на рис., вказати список суміжності вершини 5.\par
Формат оформлення відповіді:\par
список суміжності виписувати через пробіл на одному рядку, якщо пробілів десь буде більше одного - не важливо, упорядкування вершин у списку також значення не має.
%enditem


\newpage
14. Для графа на рисунку вказати послідовність, в якій алгоритм пошуку в глибину буде відмічати відвідані вершини, якщо списки суміжності вершин переглядаються алгоритмом за зростанням міток вершин та вершини відмічаються на першому вход\-женні у вершину.\par
Пошук розпочинається з вершини 2.\par

\begin{center}
	\adjustimage{max size={0.9\linewidth}{0.9\paperheight}}
	{Graphs/Traversals/349.png}
\end{center}
%enditem


\newpage
15. Для графа на рисунку з вершини 5 запущено алгоритм пошуку в ширину, що будує кістякове дерево. Списки суміжності вершин переглядаються алгоритмом за зростанням міток вершин.\par
\begin{center}
	\adjustimage{max size={0.9\linewidth}{0.9\paperheight}}
	{Graphs/Traversals/262.png}
\end{center}
\par Які з наведених ребер входять до складу отриманого кістякового дерева?\par
1) (0, 4)\par
2) (2, 6)\par
3) (1, 6)\par
4) (2, 4)\par
5) (0, 3)\par
6) (1, 4)\par
{\vskip 15pt} Номери відповідних ребер записати за зростанням через пробіл.


%enditem


\newpage
16. 

Рядки журналу через двокрапку містять ім'я користувача (ідентифікатор), час події,тип події, код події, наприклад
\begin{verbatim}
s="username : 2022-05-05 13:26 : ERROR : 404"
\end{verbatim}
у коді 
\begin{verbatim}
res = re.fullmatch('???', s) 
s = ''
code_str = ??? 
\end{verbatim}
чим треба замінити кожен з ???, щоб значенням code\_str був код події? 


\vspace{10mm}
\textbf{Завдання 17} 

Задано програму, що містить код 1-2 класів, можливе успадкування, та клієнтський код.
У наведеному коді невеликі частини закриті. 
Також наведено вихід програми.
Необхідно відновити закриті частини, щоб отриманий код відповідав наведеному виходу програми.


\textbf{Завдання 18-19} 

Кілька запитань за завданнями для самостійної роботи.
\vspace{10mm}
Скоріш за все, що завдань буде трохи менше, а деякі попередні будуть переформульовані в стилі двох останніх.

%enditem

\end {large}}

%end

\newpage
%begin

{{
\begin {large}

\newpage
1. Що буде надруковано за виконання?

\begin{verbatim}
print('R', end=' ')
b=('a','b','c','d','e',0,1,2)
print(b[0])
\end{verbatim}

%enditem


\newpage
2. Що буде надруковано за виконання?

\begin{verbatim}
print('R', end=' ')
a=['0','1','2','3','4','5','6','7','8','9']
print(a[:-12:-2])
\end{verbatim}

%enditem



\newpage
3. Що буде надруковано за виконання?

\begin{verbatim}
print('R', end=' ')
d=[10,11,12,13,14]
d.insert(-6, [2,3])
print(d)
\end{verbatim}

%enditem



\newpage
4. Що буде надруковано за виконання?

\begin{verbatim}
print('R', end=' ')
g=['0','1','2','3','4']
g[2:]=[1,2]
print(g)
\end{verbatim}

%enditem



\newpage
5. Укажіть ТИП значення виразу.\par
\begin{verbatim}
[13,2,4]+[3,5]
\end{verbatim}
%enditem


\newpage
6. Що буде надруковано за виконання?\par
\begin{verbatim}
try:
    s = {38:2, 40:9, 48:1, 48:7}
    s[48] = 4
    for x in s :
        print(x, sep=' ', end=' ')
except:
    print('ERR')
\end{verbatim}
%enditem


\newpage
7. Нехай змінна \kw{s} позначає дуже довгу послідовність цілих, по якій можна ітерувати. 

Написати вираз-генератор, що задає підпослідовність її елементів, що діляться на 3. 

Обмеження: усі елементи шуканої підпослідовності одночасно в пам'яті розміщені бути не можуть. 
%enditem


\newpage
8. Що буде надруковано за виконання?
\begin{verbatim}
c = 0
k = 1

class A:
    c = 2
    
    def __init__(self):
        global c
        self.k = 3
        c = 4
        A.c = 5

obj = A()
try:
    print(c, end=' ')
    print(k, end=' ')
    print(A.c, end=' ')
    print(obj.k)
except:
    print('ERR')
\end{verbatim}
%enditem


\newpage
9. Що буде надруковано за виконання?
\begin{verbatim}
class B0:
    def b(self): print(1, end=' ')
    def f(self): print(2, end=' ')

class B1(B0):
    def b(self): print(3, end=' ')

try:
    B0().f()
    B1().b()
except:
    print('ERR')
\end{verbatim}
%enditem


\newpage
10. Що буде надруковано за виконання?
\begin{verbatim}
class A:
    a = 1

    def _g1(self): print(2, end=' ')
    def _g2(self): print(3, end=' ')

    def main(self):
        try: print(self.a, end=' ')
        except: print('E1', end=' ')

        try: self._g1()
        except: print('E2', end=' ')

        try: self._g2()
        except: print('E3', end=' ')


class B(A):
    a = 4

    def _g1(self): print(5, end=' ')
    def _g2(self): print(6, end=' ')

try: B().main()
except: print('E4', end=' ')
\end{verbatim}
%enditem


\newpage
11. Змінні \kw{next} та \kw{prev} вузла списку позначають наступний та попередній за ним вузол відповідно. Починаючи з голови, у список записано послідовні цілі числа від~$-27$ до~$28$. Нехай змінна \kw{p} позначає вузол, що зберігає значення~$-3$.
Яке значення зберігається у вузлі \kw{p.prev.prev.next.prev.next} ?
%enditem


\newpage
12. 
\begin{center}
	\adjustimage{max size={0.8\linewidth}{0.8\paperheight}}
	{BinTree/trees_exp/52.png}
\end{center}
Указати послідовність міток вершин дерева, зображеного на рис., що утвориться в результаті обходу дерева в прямому порядку. Мітки записувати через пробіл.\par Алгоритм починає роботу з кореня (на рис. це верхня вершина).
%enditem


\newpage
13. 
\begin{center}
	\adjustimage{max size={0.9\linewidth}{0.9\paperheight}}
	{Graphs/AdjList/339.png}
\end{center}
Для графа, зображеного на рис., вказати список суміжності вершини 3.\par
Формат оформлення відповіді:\par
список суміжності виписувати через пробіл на одному рядку, якщо пробілів десь буде більше одного - не важливо, упорядкування вершин у списку також значення не має.
%enditem


\newpage
14. Для графа на рисунку вказати послідовність, в якій алгоритм пошуку в глибину буде відмічати відвідані вершини, якщо списки суміжності вершин переглядаються алгоритмом за зростанням міток вершин та вершини відмічаються на першому вход\-женні у вершину.\par
Пошук розпочинається з вершини 4.\par

\begin{center}
	\adjustimage{max size={0.9\linewidth}{0.9\paperheight}}
	{Graphs/Traversals/326.png}
\end{center}
%enditem


\newpage
15. Для графа на рисунку з вершини 3 запущено алгоритм пошуку в ширину, що будує кістякове дерево. Списки суміжності вершин переглядаються алгоритмом за зростанням міток вершин.\par
\begin{center}
	\adjustimage{max size={0.9\linewidth}{0.9\paperheight}}
	{Graphs/Traversals/332.png}
\end{center}
\par Які з наведених ребер входять до складу отриманого кістякового дерева?\par
1) (4, 6)\par
2) (0, 6)\par
3) (1, 5)\par
4) (0, 4)\par
5) (5, 6)\par
6) (2, 3)\par
{\vskip 15pt} Номери відповідних ребер записати за зростанням через пробіл.


%enditem


\newpage
16. 

Рядки журналу через двокрапку містять ім'я користувача (ідентифікатор), час події,тип події, код події, наприклад
\begin{verbatim}
s="username : 2022-05-05 13:26 : ERROR : 404"
\end{verbatim}
у коді 
\begin{verbatim}
res = re.fullmatch('???', s) 
s = ''
code_str = ??? 
\end{verbatim}
чим треба замінити кожен з ???, щоб значенням code\_str був код події? 


\vspace{10mm}
\textbf{Завдання 17} 

Задано програму, що містить код 1-2 класів, можливе успадкування, та клієнтський код.
У наведеному коді невеликі частини закриті. 
Також наведено вихід програми.
Необхідно відновити закриті частини, щоб отриманий код відповідав наведеному виходу програми.


\textbf{Завдання 18-19} 

Кілька запитань за завданнями для самостійної роботи.
\vspace{10mm}
Скоріш за все, що завдань буде трохи менше, а деякі попередні будуть переформульовані в стилі двох останніх.

%enditem

\end {large}}

%end

\newpage
%begin

{{
\begin {large}

\newpage
1. Що буде надруковано за виконання?

\begin{verbatim}
print('R', end=' ')
b=('a','b','c','d','e',0,1,2,3)
print(b[-3])
\end{verbatim}

%enditem


\newpage
2. Що буде надруковано за виконання?

\begin{verbatim}
print('R', end=' ')
e=(0,1,2,3,4,5,6,7,8,9)
print(e[-9:5:2])
\end{verbatim}

%enditem



\newpage
3. Що буде надруковано за виконання?

\begin{verbatim}
print('R', end=' ')
a=[10,11,12,13,14]
a.extend(17)
print(a)
\end{verbatim}

%enditem



\newpage
4. Що буде надруковано за виконання?

\begin{verbatim}
print('R', end=' ')
a=['0','1','2','3','4']
del a[4:-8]
print(a)
\end{verbatim}

%enditem



\newpage
5. Укажіть ТИП значення виразу.\par
\begin{verbatim}
{y:y*y for y in range(-7,7)}
\end{verbatim}
%enditem


\newpage
6. Що буде надруковано за виконання?\par
\begin{verbatim}
try:
    d = {50:6, 18:8, 77:1, 80:8, 50:7}
    d[31] = 1
    for x in d.items():
        print(x, sep=' ', end=' ')
except:
    print('ERR')
\end{verbatim}
%enditem


\newpage
7. Нехай змінна \kw{s} позначає послідовність цілих, по якій можна ітерувати.

Написати вираз, що перевіряє, чи правда, що всі елементи послідовності є цілими.
%enditem


\newpage
8. Що буде надруковано за виконання?
%Оригинал был утерян. Требует отладки!
\begin{verbatim}
_d = 0

class A:
    _d = 1
    
    def __init__(self):
        self._d = 2
        _d = 3
        A._d = 4

obj = A()
A._d = 5
try:
    print(_d, end=' ')
    print(obj._d, end=' ')
    print(A._d, end=' ')
    print(obj._d, end=' ')
    print(obj._A__d)
except:
    print('ERR')
\end{verbatim}
%enditem


\newpage
9. Що буде надруковано за виконання?
\begin{verbatim}
class D0:
    def c(self): print(1, end=' ')
    def g(self): print(2, end=' ')

class D1(D0):
    def g(self): print(3, end=' ')

try:
    D1().c()
    D1().g()
except:
    print('ERR')
\end{verbatim}
%enditem


\newpage
10. Що буде надруковано за виконання?
\begin{verbatim}
class A:
    c = 1

    def __init__(self):
        a = 2
        a = 3

class B(A):
    c = 4
    b = 5

    def __init__(self):
        super().__init__()
        self.b = 6

try:
    obj = B()
    obj.b = obj.c + 10
    B.b = 13
except: print('EE', end=' ')

try: print(obj.b, end= ' ')
except: print('E1', end=' ')

try: print(A.b, end= ' ')
except: print('E2', end=' ')

try: print(B.b, end= ' ')
except: print('E3', end=' ')
\end{verbatim}
%enditem


\newpage
11. Змінні \kw{next} та \kw{prev} вузла списку позначають наступний та попередній за ним вузол відповідно. Починаючи з голови, у список записано послідовні цілі числа від~$-68$ до~$42$. Нехай змінна \kw{p} позначає вузол, що зберігає значення~$6$.
Яке значення зберігається у вузлі \kw{p.prev.prev.prev.next.prev} ?
%enditem


\newpage
12. 
\begin{center}
	\adjustimage{max size={0.8\linewidth}{0.8\paperheight}}
	{BinTree/trees_exp/158.png}
\end{center}
Указати послідовність міток вершин дерева, зображеного на рис., що утвориться в результаті обходу дерева в оберненому порядку. Мітки записувати через пробіл.\par Алгоритм починає роботу з кореня (на рис. це верхня вершина).
%enditem


\newpage
13. 
\begin{center}
	\adjustimage{max size={0.9\linewidth}{0.9\paperheight}}
	{Graphs/AdjList/313.png}
\end{center}
Для графа, зображеного на рис., вказати список суміжності вершини 1.\par
Формат оформлення відповіді:\par
список суміжності виписувати через пробіл на одному рядку, якщо пробілів десь буде більше одного - не важливо, упорядкування вершин у списку також значення не має.
%enditem


\newpage
14. Для графа на рисунку вказати послідовність, в якій алгоритм пошуку в глибину буде відмічати відвідані вершини, якщо списки суміжності вершин переглядаються алгоритмом за зростанням міток вершин та вершини відмічаються на першому вход\-женні у вершину.\par
Пошук розпочинається з вершини 3.\par

\begin{center}
	\adjustimage{max size={0.9\linewidth}{0.9\paperheight}}
	{Graphs/Traversals/29.png}
\end{center}
%enditem


\newpage
15. Для графа на рисунку з вершини 2 запущено алгоритм пошуку в ширину, що будує кістякове дерево. Списки суміжності вершин переглядаються алгоритмом за зростанням міток вершин.\par
\begin{center}
	\adjustimage{max size={0.9\linewidth}{0.9\paperheight}}
	{Graphs/Traversals/321.png}
\end{center}
\par Які з наведених ребер входять до складу отриманого кістякового дерева?\par
1) (1, 2)\par
2) (2, 5)\par
3) (2, 6)\par
4) (2, 3)\par
5) (0, 3)\par
6) (5, 6)\par
{\vskip 15pt} Номери відповідних ребер записати за зростанням через пробіл.


%enditem


\newpage
16. 

Рядки журналу через двокрапку містять ім'я користувача (ідентифікатор), час події,тип події, код події, наприклад
\begin{verbatim}
s="username : 2022-05-05 13:26 : ERROR : 404"
\end{verbatim}
у коді 
\begin{verbatim}
res = re.fullmatch('???', s) 
s = ''
code_str = ??? 
\end{verbatim}
чим треба замінити кожен з ???, щоб значенням code\_str був код події? 


\vspace{10mm}
\textbf{Завдання 17} 

Задано програму, що містить код 1-2 класів, можливе успадкування, та клієнтський код.
У наведеному коді невеликі частини закриті. 
Також наведено вихід програми.
Необхідно відновити закриті частини, щоб отриманий код відповідав наведеному виходу програми.


\textbf{Завдання 18-19} 

Кілька запитань за завданнями для самостійної роботи.
\vspace{10mm}
Скоріш за все, що завдань буде трохи менше, а деякі попередні будуть переформульовані в стилі двох останніх.

%enditem

\end {large}}

%end

\newpage
%begin

{{
\begin {large}

\newpage
1. Що буде надруковано за виконання?

\begin{verbatim}
print('R', end=' ')
a=('a','b','c','d','e',0,1,2,3,4)
print(a[3])
\end{verbatim}

%enditem


\newpage
2. Що буде надруковано за виконання?

\begin{verbatim}
print('R', end=' ')
c=[0,1,2,3,4,5,6,7,8,9]
print(c[7:1:-3])
\end{verbatim}

%enditem



\newpage
3. Що буде надруковано за виконання?

\begin{verbatim}
print('R', end=' ')
c=['0','1','2','3','4']
c.remove('4')
print(c)
\end{verbatim}

%enditem



\newpage
4. Що буде надруковано за виконання?

\begin{verbatim}
print('R', end=' ')
h=[10,11,12,13,14]
h[:4]=[15]
print(h)
\end{verbatim}

%enditem



\newpage
5. Укажіть ТИП значення виразу.\par
\begin{verbatim}
(11,2,3)
\end{verbatim}
%enditem


\newpage
6. Що буде надруковано за виконання?\par
\begin{verbatim}
try:
    t = {79:1, 90:9, 27:2, 27:2}
    t[90] = 6
    for x, y in t.items():
        print(x, y, sep=' ', end=' ')
except:
    print('ERR')
\end{verbatim}
%enditem


\newpage
7. Нехай змінна \kw{s} позначає послідовність цілих, по якій можна ітерувати. 

Написати вираз, що перевіряє, чи правда, що послідовність  містить число 13.
%enditem


\newpage
8. Що буде надруковано за виконання?
%Оригинал был утерян. Требует отладки!
\begin{verbatim}
c = 0

class D:
    c = 1
    
    def __init__(self):
        self.c = 2
        c = 3
        D.c = 4

obj = D()
obj.c = 5
try:
    print(c, end=' ')
    print(obj.c, end=' ')
    print(D.c, end=' ')
    print(obj.c, end=' ')
    print(obj._D__c)
except:
    print('ERR')
\end{verbatim}
%enditem


\newpage
9. Що буде надруковано за виконання?
\begin{verbatim}
class C0:
    def a(self): print(1, end=' ')
    def h(self): print(2, end=' ')

class C1(C0):
    def a(self): print(3, end=' ')

try:
    C0().h()
    C1().a()
except:
    print('ERR')
\end{verbatim}
%enditem


\newpage
10. Що буде надруковано за виконання?
\begin{verbatim}
class A:
    a = 1

    def _h1(self): print(2, end=' ')
    def h2(self): print(3, end=' ')

    def main(self):
        try: print(self.a, end=' ')
        except: print('E1', end=' ')

        try: self._h1()
        except: print('E2', end=' ')

        try: self.h2()
        except: print('E3', end=' ')


class B(A):
    a = 4

    def _h1(self): print(5, end=' ')
    def h2(self): print(6, end=' ')

try: B().main()
except: print('E4', end=' ')
\end{verbatim}
%enditem


\newpage
11. Змінні \kw{next} та \kw{prev} вузла списку позначають наступний та попередній за ним вузол відповідно. Починаючи з голови, у список записано послідовні цілі числа від~$-26$ до~$15$. Нехай змінна \kw{p} позначає вузол, що зберігає значення~$-7$.
Яке значення зберігається у вузлі \kw{p.next.next.next.prev.next} ?
%enditem


\newpage
12. 
\begin{center}
	\adjustimage{max size={0.8\linewidth}{0.8\paperheight}}
	{BinTree/trees_exp/40.png}
\end{center}
Указати послідовність міток вершин дерева, зображеного на рис., що утвориться в результаті обходу дерева в прямому порядку. Мітки записувати через пробіл.\par Алгоритм починає роботу з кореня (на рис. це верхня вершина).
%enditem


\newpage
13. 
\begin{center}
	\adjustimage{max size={0.9\linewidth}{0.9\paperheight}}
	{Graphs/AdjList/185.png}
\end{center}
Для графа, зображеного на рис., вказати список суміжності вершини 1.\par
Формат оформлення відповіді:\par
список суміжності виписувати через пробіл на одному рядку, якщо пробілів десь буде більше одного - не важливо, упорядкування вершин у списку також значення не має.
%enditem


\newpage
14. Для графа на рисунку вказати послідовність, в якій алгоритм пошуку в глибину буде відмічати відвідані вершини, якщо списки суміжності вершин переглядаються алгоритмом за зростанням міток вершин та вершини відмічаються на першому вход\-женні у вершину.\par
Пошук розпочинається з вершини 3.\par

\begin{center}
	\adjustimage{max size={0.9\linewidth}{0.9\paperheight}}
	{Graphs/Traversals/24.png}
\end{center}
%enditem


\newpage
15. Для графа на рисунку з вершини 1 запущено алгоритм пошуку в ширину, що будує кістякове дерево. Списки суміжності вершин переглядаються алгоритмом за зростанням міток вершин.\par
\begin{center}
	\adjustimage{max size={0.9\linewidth}{0.9\paperheight}}
	{Graphs/Traversals/340.png}
\end{center}
\par Які з наведених ребер входять до складу отриманого кістякового дерева?\par
1) (1, 5)\par
2) (2, 4)\par
3) (1, 2)\par
4) (2, 5)\par
5) (5, 6)\par
6) (2, 3)\par
{\vskip 15pt} Номери відповідних ребер записати за зростанням через пробіл.


%enditem


\newpage
16. 

Рядки журналу через двокрапку містять ім'я користувача (ідентифікатор), час події,тип події, код події, наприклад
\begin{verbatim}
s="username : 2022-05-05 13:26 : ERROR : 404"
\end{verbatim}
у коді 
\begin{verbatim}
res = re.fullmatch('???', s) 
s = ''
code_str = ??? 
\end{verbatim}
чим треба замінити кожен з ???, щоб значенням code\_str був код події? 


\vspace{10mm}
\textbf{Завдання 17} 

Задано програму, що містить код 1-2 класів, можливе успадкування, та клієнтський код.
У наведеному коді невеликі частини закриті. 
Також наведено вихід програми.
Необхідно відновити закриті частини, щоб отриманий код відповідав наведеному виходу програми.


\textbf{Завдання 18-19} 

Кілька запитань за завданнями для самостійної роботи.
\vspace{10mm}
Скоріш за все, що завдань буде трохи менше, а деякі попередні будуть переформульовані в стилі двох останніх.

%enditem

\end {large}}

%end

\newpage
%begin

{{
\begin {large}

\newpage
1. Що буде надруковано за виконання?

\begin{verbatim}
print('R', end=' ')
h=('a','b','c','d','e',0,1)
print(h[-13])
\end{verbatim}

%enditem


\newpage
2. Що буде надруковано за виконання?

\begin{verbatim}
print('R', end=' ')
b='0123456789'
print(b[::3])
\end{verbatim}

%enditem



\newpage
3. Що буде надруковано за виконання?

\begin{verbatim}
print('R', end=' ')
b=[10,11,12,13,14]
b+='10'
print(b)
\end{verbatim}

%enditem



\newpage
4. Що буде надруковано за виконання?

\begin{verbatim}
print('R', end=' ')
e=[10,11,12,13,14]
del e[-2:5]
print(e)
\end{verbatim}

%enditem



\newpage
5. Укажіть ТИП значення виразу.\par
\begin{verbatim}
{y:None for y in range(-8,3)}
\end{verbatim}
%enditem


\newpage
6. Що буде надруковано за виконання?\par
\begin{verbatim}
try:
    t = {75:9, 69:3, 43:2, 69:0}
    t[41] = 0
    for x in t.values():
        print(x, sep=' ', end=' ')
except:
    print('ERR')
\end{verbatim}
%enditem


\newpage
7. Нехай змінна \kw{s} позначає послідовність цілих, по якій можна ітерувати.

Написати вираз, що перевіряє, чи правда, що всі елементи послідовності більші за задане число num.
%enditem


\newpage
8. Що буде надруковано за виконання?
\begin{verbatim}
d = 0
n = 1

class C:
    n = 2
    
    def __init__(self):
        self.d = 3
        n = 4
        C.n = 5

obj = C()
try:
    print(d, end=' ')
    print(n, end=' ')
    print(C.n, end=' ')
    print(obj.d)
except:
    print('ERR')
\end{verbatim}
%enditem


\newpage
9. Що буде надруковано за виконання?
\begin{verbatim}
class B0:
    def d(self): print(1, end=' ')
    def g(self): print(2, end=' ')

class B1(B0):
    def d(self): print(3, end=' ')

try:
    B1().g()
    B0().d()
except:
    print('ERR')
\end{verbatim}
%enditem


\newpage
10. Що буде надруковано за виконання?
\begin{verbatim}
class A:
    c = 1

    def __init__(self):
        b = 2
        c = 3

class B(A):
    b = 4
    a = 5

    def __init__(self):
        super().__init__()
        self.a = 6

try:
    obj = B()
    obj.b = obj.a + 10
    B.b = 13
except: print('EE', end=' ')

try: print(obj.b, end= ' ')
except: print('E1', end=' ')

try: print(A.b, end= ' ')
except: print('E2', end=' ')

try: print(B.b, end= ' ')
except: print('E3', end=' ')
\end{verbatim}
%enditem


\newpage
11. Змінні \kw{next} та \kw{prev} вузла списку позначають наступний та попередній за ним вузол відповідно. Починаючи з голови, у список записано послідовні цілі числа від~$-69$ до~$16$. Нехай змінна \kw{p} позначає вузол, що зберігає значення~$0$.
Яке значення зберігається у вузлі \kw{p.prev.next.prev.next.next} ?
%enditem


\newpage
12. 
\begin{center}
	\adjustimage{max size={0.8\linewidth}{0.8\paperheight}}
	{BinTree/trees_exp/87.png}
\end{center}
Указати послідовність міток вершин дерева, зображеного на рис., що утвориться в результаті обходу дерева в оберненому порядку. Мітки записувати через пробіл.\par Алгоритм починає роботу з кореня (на рис. це верхня вершина).
%enditem


\newpage
13. 
\begin{center}
	\adjustimage{max size={0.9\linewidth}{0.9\paperheight}}
	{Graphs/AdjList/18.png}
\end{center}
Для графа, зображеного на рис., вказати список суміжності вершини 7.\par
Формат оформлення відповіді:\par
список суміжності виписувати через пробіл на одному рядку, якщо пробілів десь буде більше одного - не важливо, упорядкування вершин у списку також значення не має.
%enditem


\newpage
14. Для графа на рисунку вказати послідовність, в якій алгоритм пошуку в глибину буде відмічати відвідані вершини, якщо списки суміжності вершин переглядаються алгоритмом за зростанням міток вершин та вершини відмічаються на першому вход\-женні у вершину.\par
Пошук розпочинається з вершини 4.\par

\begin{center}
	\adjustimage{max size={0.9\linewidth}{0.9\paperheight}}
	{Graphs/Traversals/292.png}
\end{center}
%enditem


\newpage
15. Для графа на рисунку з вершини 1 запущено алгоритм пошуку в ширину, що будує кістякове дерево. Списки суміжності вершин переглядаються алгоритмом за зростанням міток вершин.\par
\begin{center}
	\adjustimage{max size={0.9\linewidth}{0.9\paperheight}}
	{Graphs/Traversals/324.png}
\end{center}
\par Які з наведених ребер входять до складу отриманого кістякового дерева?\par
1) (1, 6)\par
2) (1, 4)\par
3) (1, 5)\par
4) (1, 3)\par
5) (3, 4)\par
6) (0, 5)\par
{\vskip 15pt} Номери відповідних ребер записати за зростанням через пробіл.


%enditem


\newpage
16. 

Рядки журналу через двокрапку містять ім'я користувача (ідентифікатор), час події,тип події, код події, наприклад
\begin{verbatim}
s="username : 2022-05-05 13:26 : ERROR : 404"
\end{verbatim}
у коді 
\begin{verbatim}
res = re.fullmatch('???', s) 
s = ''
code_str = ??? 
\end{verbatim}
чим треба замінити кожен з ???, щоб значенням code\_str був код події? 


\vspace{10mm}
\textbf{Завдання 17} 

Задано програму, що містить код 1-2 класів, можливе успадкування, та клієнтський код.
У наведеному коді невеликі частини закриті. 
Також наведено вихід програми.
Необхідно відновити закриті частини, щоб отриманий код відповідав наведеному виходу програми.


\textbf{Завдання 18-19} 

Кілька запитань за завданнями для самостійної роботи.
\vspace{10mm}
Скоріш за все, що завдань буде трохи менше, а деякі попередні будуть переформульовані в стилі двох останніх.

%enditem

\end {large}}

%end

\newpage
%begin

{{
\begin {large}

\newpage
1. Що буде надруковано за виконання?

\begin{verbatim}
print('R', end=' ')
g=('a','b','c','d','e',0,1,2,4)
print(g[14])
\end{verbatim}

%enditem


\newpage
2. Що буде надруковано за виконання?

\begin{verbatim}
print('R', end=' ')
h=('0','1','2','3','4','5','6','7','8','9')
print(h[11::-1])
\end{verbatim}

%enditem



\newpage
3. Що буде надруковано за виконання?

\begin{verbatim}
print('R', end=' ')
g=['0','1','2','3','4']
g.append('12')
print(g)
\end{verbatim}

%enditem



\newpage
4. Що буде надруковано за виконання?

\begin{verbatim}
print('R', end=' ')
g=['0','1','2','3','4']
g[-1:5]=[1,0]
print(g)
\end{verbatim}

%enditem



\newpage
5. Укажіть ТИП значення виразу.\par
\begin{verbatim}
{x for x in range(7) if el%3==1}
\end{verbatim}
%enditem


\newpage
6. Що буде надруковано за виконання?\par
\begin{verbatim}
try:
    s = {54:6, 24:0, 80:3, 42:7, 54:3}
    s[39] = 5
    for x, y in s.items():
        print(x, y, sep=' ', end=' ')
except:
    print('ERR')
\end{verbatim}
%enditem


\newpage
7. Записати ВИРАЗ, значенням якого є словник, ключами якого є цілі числа від 14 до 2019 (включно), що не діляться на жодне з чисел 2, 3, а ключам відповідають їх квадратні корені  
%enditem


\newpage
8. Що буде надруковано за виконання?
\begin{verbatim}
b = 0
m = 1

class B:
    b = 2
    
    def __init__(self):
        global b
        self.m = 3
        b = 4
        B.b = 5

obj = B()
try:
    print(b, end=' ')
    print(m, end=' ')
    print(B.b, end=' ')
    print(obj.m)
except:
    print('ERR')
\end{verbatim}
%enditem


\newpage
9. Що буде надруковано за виконання?
\begin{verbatim}
class A0:
    def b(self): print(1, end=' ')
    def h(self): print(2, end=' ')

class A1(A0):
    def h(self): print(3, end=' ')

try:
    A1().b()
    A1().h()
except:
    print('ERR')
\end{verbatim}
%enditem


\newpage
10. Що буде надруковано за виконання?
\begin{verbatim}
class A:
    a = 1

    def __init__(self):
        c = 2
        a = 3

class B(A):
    c = 4
    b = 5

    def __init__(self):
        super().__init__()
        self.b = 6

try:
    obj = B()
    obj.c = obj.b + 10
    B.c = 13
except: print('EE', end=' ')

try: print(obj.c, end= ' ')
except: print('E1', end=' ')

try: print(A.c, end= ' ')
except: print('E2', end=' ')

try: print(B.c, end= ' ')
except: print('E3', end=' ')
\end{verbatim}
%enditem


\newpage
11. Змінні \kw{next} та \kw{prev} вузла списку позначають наступний та попередній за ним вузол відповідно. Починаючи з голови, у список записано послідовні цілі числа від~$-51$ до~$24$. Нехай змінна \kw{p} позначає вузол, що зберігає значення~$-6$.
Яке значення зберігається у вузлі \kw{p.next.prev.next.prev.prev} ?
%enditem


\newpage
12. 
\begin{center}
	\adjustimage{max size={0.8\linewidth}{0.8\paperheight}}
	{BinTree/trees_exp/98.png}
\end{center}
Указати послідовність міток вершин дерева, зображеного на рис., що утвориться в результаті обходу дерева в прямому порядку. Мітки записувати через пробіл.\par Алгоритм починає роботу з кореня (на рис. це верхня вершина).
%enditem


\newpage
13. 
\begin{center}
	\adjustimage{max size={0.9\linewidth}{0.9\paperheight}}
	{Graphs/AdjList/290.png}
\end{center}
Для графа, зображеного на рис., вказати список суміжності вершини 1.\par
Формат оформлення відповіді:\par
список суміжності виписувати через пробіл на одному рядку, якщо пробілів десь буде більше одного - не важливо, упорядкування вершин у списку також значення не має.
%enditem


\newpage
14. Для графа на рисунку з вершини 1 запущено алгоритм пошуку в глибину, що будує кістякове дерево. Списки суміжності вершин переглядаються алгоритмом за зростанням міток вершин.\par
\begin{center}
	\adjustimage{max size={0.9\linewidth}{0.9\paperheight}}
	{Graphs/Traversals/183.png}
\end{center}
\par Які з наведених ребер входять до складу отриманого кістякового дерева?\par
1) (0, 4)\par
2) (4, 6)\par
3) (0, 6)\par
4) (2, 5)\par
5) (0, 3)\par
6) (1, 3)\par
{\vskip 15pt} Номери відповідних ребер записати за зростанням через пробіл.

%enditem


\newpage
15. Для графа на рисунку з вершини 2 запущено алгоритм пошуку в ширину, що будує кістякове дерево. Списки суміжності вершин переглядаються алгоритмом за зростанням міток вершин.\par
\begin{center}
	\adjustimage{max size={0.9\linewidth}{0.9\paperheight}}
	{Graphs/Traversals/41.png}
\end{center}
\par Які з наведених ребер входять до складу отриманого кістякового дерева?\par
1) (0, 3)\par
2) (2, 3)\par
3) (0, 2)\par
4) (0, 6)\par
5) (0, 5)\par
6) (1, 5)\par
{\vskip 15pt} Номери відповідних ребер записати за зростанням через пробіл.


%enditem


\newpage
16. 

Рядки журналу через двокрапку містять ім'я користувача (ідентифікатор), час події,тип події, код події, наприклад
\begin{verbatim}
s="username : 2022-05-05 13:26 : ERROR : 404"
\end{verbatim}
у коді 
\begin{verbatim}
res = re.fullmatch('???', s) 
s = ''
code_str = ??? 
\end{verbatim}
чим треба замінити кожен з ???, щоб значенням code\_str був код події? 


\vspace{10mm}
\textbf{Завдання 17} 

Задано програму, що містить код 1-2 класів, можливе успадкування, та клієнтський код.
У наведеному коді невеликі частини закриті. 
Також наведено вихід програми.
Необхідно відновити закриті частини, щоб отриманий код відповідав наведеному виходу програми.


\textbf{Завдання 18-19} 

Кілька запитань за завданнями для самостійної роботи.
\vspace{10mm}
Скоріш за все, що завдань буде трохи менше, а деякі попередні будуть переформульовані в стилі двох останніх.

%enditem

\end {large}}

%end

\newpage
%begin

{{
\begin {large}

\newpage
1. Що буде надруковано за виконання?

\begin{verbatim}
print('R', end=' ')
c=('a','b','c','d','e',0,1,2)
print(c[-7])
\end{verbatim}

%enditem


\newpage
2. Що буде надруковано за виконання?

\begin{verbatim}
print('R', end=' ')
g=(0,1,2,3,4,5,6,7,8,9)
print(g[-3:6:-1])
\end{verbatim}

%enditem



\newpage
3. Що буде надруковано за виконання?

\begin{verbatim}
print('R', end=' ')
h=[10,11,12,13,14]
h.append([1,2])
print(h)
\end{verbatim}

%enditem



\newpage
4. Що буде надруковано за виконання?

\begin{verbatim}
print('R', end=' ')
c=['0','1','2','3','4']
c[5:-6]=[2,3]
print(c)
\end{verbatim}

%enditem



\newpage
5. Укажіть ТИП значення виразу.\par
\begin{verbatim}
(2, 7, -4)
\end{verbatim}
%enditem


\newpage
6. Що буде надруковано за виконання?\par
\begin{verbatim}
try:
    d = {43:0, 74:8, 75:1, 75:4}
    d[65] = 1
    for x in d :
        print(x, sep=' ', end=' ')
except:
    print('ERR')
\end{verbatim}
%enditem


\newpage
7. Нехай змінна \kw{s} позначає дуже довгу послідовність цілих, по якій можна ітерувати. 

Написати вираз-генератор, що задає підпослідовність її елементів, що є від'єм\-ни\-ми. 

Обмеження: усі елементи шуканої підпослідовності одночасно в пам'яті розміщені бути не можуть. 
%enditem


\newpage
8. Що буде надруковано за виконання?
%Оригинал был утерян. Требует отладки!
\begin{verbatim}
__b = 0

class C:
    __b = 1
    
    def __init__(self):
        self.__b = 2
        __b = 3
        C.__b = 4

obj = C()
C.__b = 5
try:
    print(__b, end=' ')
    print(obj.__b, end=' ')
    print(C.__b, end=' ')
    print(obj.__b, end=' ')
    print(obj._C__b)
except:
    print('ERR')
\end{verbatim}
%enditem


\newpage
9. Що буде надруковано за виконання?
\begin{verbatim}
class D0:
    def a(self): print(1, end=' ')
    def f(self): print(2, end=' ')

class D1(D0):
    def a(self): print(3, end=' ')

try:
    D1().f()
    D1().a()
except:
    print('ERR')
\end{verbatim}
%enditem


\newpage
10. Що буде надруковано за виконання?
\begin{verbatim}
class A:
    c = 1

    def __h1(self): print(2, end=' ')
    def _h2(self): print(3, end=' ')

    def main(self):
        try: print(self.c, end=' ')
        except: print('E1', end=' ')

        try: self.__h1()
        except: print('E2', end=' ')

        try: self._h2()
        except: print('E3', end=' ')


class B(A):
    c = 4

    def __h1(self): print(5, end=' ')
    def _h2(self): print(6, end=' ')

try: B().main()
except: print('E4', end=' ')
\end{verbatim}
%enditem


\newpage
11. Змінні \kw{next} та \kw{prev} вузла списку позначають наступний та попередній за ним вузол відповідно. Починаючи з голови, у список записано послідовні цілі числа від~$-25$ до~$48$. Нехай змінна \kw{p} позначає вузол, що зберігає значення~$-5$.
Яке значення зберігається у вузлі \kw{p.prev.prev.prev.next.next} ?
%enditem


\newpage
12. 
\begin{center}
	\adjustimage{max size={0.8\linewidth}{0.8\paperheight}}
	{BinTree/trees_exp/67.png}
\end{center}
Указати послідовність міток вершин дерева, зображеного на рис., що утвориться в результаті обходу дерева в оберненому порядку. Мітки записувати через пробіл.\par Алгоритм починає роботу з кореня (на рис. це верхня вершина).
%enditem


\newpage
13. 
\begin{center}
	\adjustimage{max size={0.9\linewidth}{0.9\paperheight}}
	{Graphs/AdjList/85.png}
\end{center}
Для графа, зображеного на рис., вказати список суміжності вершини 0.\par
Формат оформлення відповіді:\par
список суміжності виписувати через пробіл на одному рядку, якщо пробілів десь буде більше одного - не важливо, упорядкування вершин у списку також значення не має.
%enditem


\newpage
14. Для графа на рисунку з вершини 5 запущено алгоритм пошуку в глибину, що будує кістякове дерево. Списки суміжності вершин переглядаються алгоритмом за зростанням міток вершин.\par
\begin{center}
	\adjustimage{max size={0.9\linewidth}{0.9\paperheight}}
	{Graphs/Traversals/86.png}
\end{center}
\par Які з наведених ребер входять до складу отриманого кістякового дерева?\par
1) (0, 1)\par
2) (0, 4)\par
3) (3, 5)\par
4) (1, 4)\par
5) (2, 5)\par
6) (0, 2)\par
{\vskip 15pt} Номери відповідних ребер записати за зростанням через пробіл.

%enditem


\newpage
15. Для графа на рисунку з вершини 3 запущено алгоритм пошуку в ширину, що будує кістякове дерево. Списки суміжності вершин переглядаються алгоритмом за зростанням міток вершин.\par
\begin{center}
	\adjustimage{max size={0.9\linewidth}{0.9\paperheight}}
	{Graphs/Traversals/230.png}
\end{center}
\par Які з наведених ребер входять до складу отриманого кістякового дерева?\par
1) (0, 6)\par
2) (3, 5)\par
3) (1, 3)\par
4) (1, 2)\par
5) (0, 5)\par
6) (0, 4)\par
{\vskip 15pt} Номери відповідних ребер записати за зростанням через пробіл.


%enditem


\newpage
16. 

Рядки журналу через двокрапку містять ім'я користувача (ідентифікатор), час події,тип події, код події, наприклад
\begin{verbatim}
s="username : 2022-05-05 13:26 : ERROR : 404"
\end{verbatim}
у коді 
\begin{verbatim}
res = re.fullmatch('???', s) 
s = ''
code_str = ??? 
\end{verbatim}
чим треба замінити кожен з ???, щоб значенням code\_str був код події? 


\vspace{10mm}
\textbf{Завдання 17} 

Задано програму, що містить код 1-2 класів, можливе успадкування, та клієнтський код.
У наведеному коді невеликі частини закриті. 
Також наведено вихід програми.
Необхідно відновити закриті частини, щоб отриманий код відповідав наведеному виходу програми.


\textbf{Завдання 18-19} 

Кілька запитань за завданнями для самостійної роботи.
\vspace{10mm}
Скоріш за все, що завдань буде трохи менше, а деякі попередні будуть переформульовані в стилі двох останніх.

%enditem

\end {large}}

%end

\newpage
%begin

{{
\begin {large}

\newpage
1. Що буде надруковано за виконання?

\begin{verbatim}
print('R', end=' ')
e=('a','b','c','d','e',0,1,2,3)
print(e[-9])
\end{verbatim}

%enditem


\newpage
2. Що буде надруковано за виконання?

\begin{verbatim}
print('R', end=' ')
f=[0,1,2,3,4,5,6,7,8,9]
print(f[:3:2])
\end{verbatim}

%enditem



\newpage
3. Що буде надруковано за виконання?

\begin{verbatim}
print('R', end=' ')
d=['0','1','2','3','4']
d.extend('10')
print(d)
\end{verbatim}

%enditem



\newpage
4. Що буде надруковано за виконання?

\begin{verbatim}
print('R', end=' ')
d=[10,11,12,13,14]
del d[-8:-2]
print(d)
\end{verbatim}

%enditem



\newpage
5. Укажіть ТИП значення виразу.\par
\begin{verbatim}
[13,2,4]+[3,5]
\end{verbatim}
%enditem


\newpage
6. Що буде надруковано за виконання?\par
\begin{verbatim}
try:
    t = {77:4, 67:8, 16:6, 20:7, 16:5}
    t[67] = 6
    for x in t.items():
        print(x, sep=' ', end=' ')
except:
    print('ERR')
\end{verbatim}
%enditem


\newpage
7. Нехай змінна \kw{s} позначає дуже довгу послідовність, по якій можна ітерувати. 

Написати вираз-генератор, що задає підпослідовність її елементів, що не є числами типу \kw{float}.

Обмеження: усі елементи шуканої підпослідовності одночасно в пам'яті розміщені бути не можуть. 
%enditem


\newpage
8. Що буде надруковано за виконання?
\begin{verbatim}
a = 0
j = 1

class D:
    j = 2
    
    def __init__(self):
        self.a = 3
        j = 4
        D.j = 5

obj = D()
try:
    print(a, end=' ')
    print(j, end=' ')
    print(D.j, end=' ')
    print(obj.a)
except:
    print('ERR')
\end{verbatim}
%enditem


\newpage
9. Що буде надруковано за виконання?
\begin{verbatim}
class C0:
    def c(self): print(1, end=' ')
    def f(self): print(2, end=' ')

class C1(C0):
    def f(self): print(3, end=' ')

try:
    C1().c()
    C1().f()
except:
    print('ERR')
\end{verbatim}
%enditem


\newpage
10. Що буде надруковано за виконання?
\begin{verbatim}
class A:
    b = 1

    def __init__(self):
        c = 2
        c = 3

class B(A):
    b = 4
    a = 5

    def __init__(self):
        super().__init__()
        self.a = 6

try:
    obj = B()
    obj.c = obj.a + 10
    B.c = 13
except: print('EE', end=' ')

try: print(obj.c, end= ' ')
except: print('E1', end=' ')

try: print(A.c, end= ' ')
except: print('E2', end=' ')

try: print(B.c, end= ' ')
except: print('E3', end=' ')
\end{verbatim}
%enditem


\newpage
11. Змінні \kw{next} та \kw{prev} вузла списку позначають наступний та попередній за ним вузол відповідно. Починаючи з голови, у список записано послідовні цілі числа від~$-61$ до~$44$. Нехай змінна \kw{p} позначає вузол, що зберігає значення~$-1$.
Яке значення зберігається у вузлі \kw{p.next.next.next.prev.prev} ?
%enditem


\newpage
12. 
\begin{center}
	\adjustimage{max size={0.8\linewidth}{0.8\paperheight}}
	{BinTree/trees_exp/56.png}
\end{center}
Указати послідовність міток вершин дерева, зображеного на рис., що утвориться в результаті обходу дерева в прямому порядку. Мітки записувати через пробіл.\par Алгоритм починає роботу з кореня (на рис. це верхня вершина).
%enditem


\newpage
13. 
\begin{center}
	\adjustimage{max size={0.9\linewidth}{0.9\paperheight}}
	{Graphs/AdjList/2.png}
\end{center}
Для графа, зображеного на рис., вказати список суміжності вершини 6.\par
Формат оформлення відповіді:\par
список суміжності виписувати через пробіл на одному рядку, якщо пробілів десь буде більше одного - не важливо, упорядкування вершин у списку також значення не має.
%enditem


\newpage
14. Для графа на рисунку вказати послідовність, в якій алгоритм пошуку в глибину буде відмічати відвідані вершини, якщо списки суміжності вершин переглядаються алгоритмом за зростанням міток вершин та вершини відмічаються на першому вход\-женні у вершину.\par
Пошук розпочинається з вершини 6.\par

\begin{center}
	\adjustimage{max size={0.9\linewidth}{0.9\paperheight}}
	{Graphs/Traversals/266.png}
\end{center}
%enditem


\newpage
15. Для графа на рисунку з вершини 4 запущено алгоритм пошуку в ширину, що будує кістякове дерево. Списки суміжності вершин переглядаються алгоритмом за зростанням міток вершин.\par
\begin{center}
	\adjustimage{max size={0.9\linewidth}{0.9\paperheight}}
	{Graphs/Traversals/88.png}
\end{center}
\par Які з наведених ребер входять до складу отриманого кістякового дерева?\par
1) (3, 6)\par
2) (5, 6)\par
3) (2, 3)\par
4) (1, 5)\par
5) (0, 1)\par
6) (3, 5)\par
{\vskip 15pt} Номери відповідних ребер записати за зростанням через пробіл.


%enditem


\newpage
16. 

Рядки журналу через двокрапку містять ім'я користувача (ідентифікатор), час події,тип події, код події, наприклад
\begin{verbatim}
s="username : 2022-05-05 13:26 : ERROR : 404"
\end{verbatim}
у коді 
\begin{verbatim}
res = re.fullmatch('???', s) 
s = ''
code_str = ??? 
\end{verbatim}
чим треба замінити кожен з ???, щоб значенням code\_str був код події? 


\vspace{10mm}
\textbf{Завдання 17} 

Задано програму, що містить код 1-2 класів, можливе успадкування, та клієнтський код.
У наведеному коді невеликі частини закриті. 
Також наведено вихід програми.
Необхідно відновити закриті частини, щоб отриманий код відповідав наведеному виходу програми.


\textbf{Завдання 18-19} 

Кілька запитань за завданнями для самостійної роботи.
\vspace{10mm}
Скоріш за все, що завдань буде трохи менше, а деякі попередні будуть переформульовані в стилі двох останніх.

%enditem

\end {large}}

%end

\newpage
%begin

{{
\begin {large}

\newpage
1. Що буде надруковано за виконання?

\begin{verbatim}
print('R', end=' ')
d=('a','b','c','d','e',0,1)
print(d[-8])
\end{verbatim}

%enditem


\newpage
2. Що буде надруковано за виконання?

\begin{verbatim}
print('R', end=' ')
d=['0','1','2','3','4','5','6','7','8','9']
print(d[13::-3])
\end{verbatim}

%enditem



\newpage
3. Що буде надруковано за виконання?

\begin{verbatim}
print('R', end=' ')
e=['0','1','2','3','4']
e.remove(1)
print(e)
\end{verbatim}

%enditem



\newpage
4. Що буде надруковано за виконання?

\begin{verbatim}
print('R', end=' ')
f=['0','1','2','3','4']
f[:8]=(14)
print(f)
\end{verbatim}

%enditem



\newpage
5. Укажіть ТИП значення виразу.\par
\begin{verbatim}
{21,2,4}
\end{verbatim}
%enditem


\newpage
6. Що буде надруковано за виконання?\par
\begin{verbatim}
try:
    s = {25:6, 44:3, 55:7, 55:8}
    s[55] = 3
    for x in s.items():
        print(*x, sep=' ', end=' ')
except:
    print('ERR')
\end{verbatim}
%enditem


\newpage
7. Нехай змінна \kw{s} позначає послідовність цілих, по якій можна ітерувати. 

Написати вираз, що перевіряє, чи правда, що послідовність  містить хоча б одне додатнє число.
%enditem


\newpage
8. Що буде надруковано за виконання?
%Оригинал был утерян. Требует отладки!
\begin{verbatim}
_a = 0

class B:
    _a = 1
    
    def __init__(self):
        self._a = 2
        _a = 3
        B._a = 4

obj = B()
obj._a = 5
try:
    print(_a, end=' ')
    print(obj._a, end=' ')
    print(B._a, end=' ')
    print(obj._a, end=' ')
    print(obj._B__a)
except:
    print('ERR')
\end{verbatim}
%enditem


\newpage
9. Що буде надруковано за виконання?
\begin{verbatim}
class B0:
    def b(self): print(1, end=' ')
    def g(self): print(2, end=' ')

class B1(B0):
    def g(self): print(3, end=' ')

try:
    B0().b()
    B1().g()
except:
    print('ERR')
\end{verbatim}
%enditem


\newpage
10. Що буде надруковано за виконання?
\begin{verbatim}
class A:
    c = 1

    def __init__(self):
        b = 2
        c = 3

class B(A):
    c = 4
    a = 5

    def __init__(self):
        super().__init__()
        self.a = 6

try:
    obj = B()
    obj.b = obj.a + 10
    B.b = 13
except: print('EE', end=' ')

try: print(obj.b, end= ' ')
except: print('E1', end=' ')

try: print(A.b, end= ' ')
except: print('E2', end=' ')

try: print(B.b, end= ' ')
except: print('E3', end=' ')
\end{verbatim}
%enditem


\newpage
11. Змінні \kw{next} та \kw{prev} вузла списку позначають наступний та попередній за ним вузол відповідно. Починаючи з голови, у список записано послідовні цілі числа від~$-65$ до~$40$. Нехай змінна \kw{p} позначає вузол, що зберігає значення~$7$.
Яке значення зберігається у вузлі \kw{p.next.prev.prev.next.prev} ?
%enditem


\newpage
12. 
\begin{center}
	\adjustimage{max size={0.8\linewidth}{0.8\paperheight}}
	{BinTree/trees_exp/23.png}
\end{center}
Указати послідовність міток вершин дерева, зображеного на рис., що утвориться в результаті обходу дерева в оберненому порядку. Мітки записувати через пробіл.\par Алгоритм починає роботу з кореня (на рис. це верхня вершина).
%enditem


\newpage
13. 
\begin{center}
	\adjustimage{max size={0.9\linewidth}{0.9\paperheight}}
	{Graphs/AdjList/299.png}
\end{center}
Для графа, зображеного на рис., вказати список суміжності вершини 3.\par
Формат оформлення відповіді:\par
список суміжності виписувати через пробіл на одному рядку, якщо пробілів десь буде більше одного - не важливо, упорядкування вершин у списку також значення не має.
%enditem


\newpage
14. Для графа на рисунку вказати послідовність, в якій алгоритм пошуку в глибину буде відмічати відвідані вершини, якщо списки суміжності вершин переглядаються алгоритмом за зростанням міток вершин та вершини відмічаються на першому вход\-женні у вершину.\par
Пошук розпочинається з вершини 2.\par

\begin{center}
	\adjustimage{max size={0.9\linewidth}{0.9\paperheight}}
	{Graphs/Traversals/97.png}
\end{center}
%enditem


\newpage
15. Для графа на рисунку з вершини 0 запущено алгоритм пошуку в ширину, що будує кістякове дерево. Списки суміжності вершин переглядаються алгоритмом за зростанням міток вершин.\par
\begin{center}
	\adjustimage{max size={0.9\linewidth}{0.9\paperheight}}
	{Graphs/Traversals/83.png}
\end{center}
\par Які з наведених ребер входять до складу отриманого кістякового дерева?\par
1) (0, 4)\par
2) (2, 3)\par
3) (1, 2)\par
4) (4, 5)\par
5) (3, 5)\par
6) (3, 4)\par
{\vskip 15pt} Номери відповідних ребер записати за зростанням через пробіл.


%enditem


\newpage
16. 

Рядки журналу через двокрапку містять ім'я користувача (ідентифікатор), час події,тип події, код події, наприклад
\begin{verbatim}
s="username : 2022-05-05 13:26 : ERROR : 404"
\end{verbatim}
у коді 
\begin{verbatim}
res = re.fullmatch('???', s) 
s = ''
code_str = ??? 
\end{verbatim}
чим треба замінити кожен з ???, щоб значенням code\_str був код події? 


\vspace{10mm}
\textbf{Завдання 17} 

Задано програму, що містить код 1-2 класів, можливе успадкування, та клієнтський код.
У наведеному коді невеликі частини закриті. 
Також наведено вихід програми.
Необхідно відновити закриті частини, щоб отриманий код відповідав наведеному виходу програми.


\textbf{Завдання 18-19} 

Кілька запитань за завданнями для самостійної роботи.
\vspace{10mm}
Скоріш за все, що завдань буде трохи менше, а деякі попередні будуть переформульовані в стилі двох останніх.

%enditem

\end {large}}

%end

\newpage
%begin

{{
\begin {large}

\newpage
1. Що буде надруковано за виконання?

\begin{verbatim}
print('R', end=' ')
f=('a','b','c','d','e',0,1,2,3,4)
print(f[5])
\end{verbatim}

%enditem


\newpage
2. Що буде надруковано за виконання?

\begin{verbatim}
print('R', end=' ')
a=('0','1','2','3','4','5','6','7','8','9')
print(a[::3])
\end{verbatim}

%enditem



\newpage
3. Що буде надруковано за виконання?

\begin{verbatim}
print('R', end=' ')
f=[10,11,12,13,14]
f+='12'
print(f)
\end{verbatim}

%enditem



\newpage
4. Що буде надруковано за виконання?

\begin{verbatim}
print('R', end=' ')
b=[10,11,12,13,14]
b[-7:]='10'
print(b)
\end{verbatim}

%enditem



\newpage
5. Укажіть ТИП значення виразу.\par
\begin{verbatim}
[{1,2}, 5]
\end{verbatim}
%enditem


\newpage
6. Що буде надруковано за виконання?\par
\begin{verbatim}
try:
    d = {75:1, 56:1, 63:1, 47:2, 47:9}
    d[64] = 1
    for x in d.keys():
        print(x, sep=' ', end=' ')
except:
    print('ERR')
\end{verbatim}
%enditem


\newpage
7. Записати ВИРАЗ, значенням якого є множина, що складається з кубів цілих чисел від 20 до 2011 (включно), що не діляться на жодне з чисел 2, 3.
%enditem


\newpage
8. Що буде надруковано за виконання?
\begin{verbatim}
c = 0
h = 1

class A:
    h = 2
    
    def __init__(self):
        global h
        self.h = 3
        h = 4
        A.h = 5

obj = A()
try:
    print(c, end=' ')
    print(h, end=' ')
    print(A.h, end=' ')
    print(obj.h)
except:
    print('ERR')
\end{verbatim}
%enditem


\newpage
9. Що буде надруковано за виконання?
\begin{verbatim}
class A0:
    def d(self): print(1, end=' ')
    def h(self): print(2, end=' ')

class A1(A0):
    def d(self): print(3, end=' ')

try:
    A1().h()
    A0().d()
except:
    print('ERR')
\end{verbatim}
%enditem


\newpage
10. Що буде надруковано за виконання?
\begin{verbatim}
class A:
    d = 1

    def f1(self): print(2, end=' ')
    def _f2(self): print(3, end=' ')

    def main(self):
        try: print(self.d, end=' ')
        except: print('E1', end=' ')

        try: self.f1()
        except: print('E2', end=' ')

        try: self._f2()
        except: print('E3', end=' ')


class B(A):
    d = 4

    def f1(self): print(5, end=' ')
    def _f2(self): print(6, end=' ')

try: B().main()
except: print('E4', end=' ')
\end{verbatim}
%enditem


\newpage
11. Змінні \kw{next} та \kw{prev} вузла списку позначають наступний та попередній за ним вузол відповідно. Починаючи з голови, у список записано послідовні цілі числа від~$-62$ до~$45$. Нехай змінна \kw{p} позначає вузол, що зберігає значення~$-2$.
Яке значення зберігається у вузлі \kw{p.prev.next.next.prev.next} ?
%enditem


\newpage
12. 
\begin{center}
	\adjustimage{max size={0.8\linewidth}{0.8\paperheight}}
	{BinTree/trees_exp/6.png}
\end{center}
Указати послідовність міток вершин дерева, зображеного на рис., що утвориться в результаті обходу дерева в прямому порядку. Мітки записувати через пробіл.\par Алгоритм починає роботу з кореня (на рис. це верхня вершина).
%enditem


\newpage
13. 
\begin{center}
	\adjustimage{max size={0.9\linewidth}{0.9\paperheight}}
	{Graphs/AdjList/47.png}
\end{center}
Для графа, зображеного на рис., вказати список суміжності вершини 1.\par
Формат оформлення відповіді:\par
список суміжності виписувати через пробіл на одному рядку, якщо пробілів десь буде більше одного - не важливо, упорядкування вершин у списку також значення не має.
%enditem


\newpage
14. Для графа на рисунку вказати послідовність, в якій алгоритм пошуку в глибину буде відмічати відвідані вершини, якщо списки суміжності вершин переглядаються алгоритмом за зростанням міток вершин та вершини відмічаються на першому вход\-женні у вершину.\par
Пошук розпочинається з вершини 2.\par

\begin{center}
	\adjustimage{max size={0.9\linewidth}{0.9\paperheight}}
	{Graphs/Traversals/321.png}
\end{center}
%enditem


\newpage
15. Для графа на рисунку з вершини 2 запущено алгоритм пошуку в ширину, що будує кістякове дерево. Списки суміжності вершин переглядаються алгоритмом за зростанням міток вершин.\par
\begin{center}
	\adjustimage{max size={0.9\linewidth}{0.9\paperheight}}
	{Graphs/Traversals/113.png}
\end{center}
\par Які з наведених ребер входять до складу отриманого кістякового дерева?\par
1) (2, 6)\par
2) (3, 4)\par
3) (5, 6)\par
4) (3, 6)\par
5) (1, 4)\par
6) (2, 4)\par
{\vskip 15pt} Номери відповідних ребер записати за зростанням через пробіл.


%enditem


\newpage
16. 

Рядки журналу через двокрапку містять ім'я користувача (ідентифікатор), час події,тип події, код події, наприклад
\begin{verbatim}
s="username : 2022-05-05 13:26 : ERROR : 404"
\end{verbatim}
у коді 
\begin{verbatim}
res = re.fullmatch('???', s) 
s = ''
code_str = ??? 
\end{verbatim}
чим треба замінити кожен з ???, щоб значенням code\_str був код події? 


\vspace{10mm}
\textbf{Завдання 17} 

Задано програму, що містить код 1-2 класів, можливе успадкування, та клієнтський код.
У наведеному коді невеликі частини закриті. 
Також наведено вихід програми.
Необхідно відновити закриті частини, щоб отриманий код відповідав наведеному виходу програми.


\textbf{Завдання 18-19} 

Кілька запитань за завданнями для самостійної роботи.
\vspace{10mm}
Скоріш за все, що завдань буде трохи менше, а деякі попередні будуть переформульовані в стилі двох останніх.

%enditem

\end {large}}

%end

\newpage
%begin

{{
\begin {large}

\newpage
1. Що буде надруковано за виконання?

\begin{verbatim}
print('R', end=' ')
c=('a','b','c','d','e',0,1,2,4)
print(c[-12])
\end{verbatim}

%enditem


\newpage
2. Що буде надруковано за виконання?

\begin{verbatim}
print('R', end=' ')
b='0123456789'
print(b[-5:0:-2])
\end{verbatim}

%enditem



\newpage
3. Що буде надруковано за виконання?

\begin{verbatim}
print('R', end=' ')
g=[10,11,12,13,14]
print(g.pop(-4), end=' ')
print(g)
\end{verbatim}

%enditem



\newpage
4. Що буде надруковано за виконання?

\begin{verbatim}
print('R', end=' ')
c=[10,11,12,13,14]
c[6:-3]='11'
print(c)
\end{verbatim}

%enditem



\newpage
5. Укажіть ТИП значення виразу.\par
\begin{verbatim}
{x:0 for x in range(7)}
\end{verbatim}
%enditem


\newpage
6. Що буде надруковано за виконання?\par
\begin{verbatim}
try:
    d = {39:5, 22:5, 70:3, 70:1}
    d[66] = 8
    for x, y in d.items():
        print(x, y, sep=' ', end=' ')
except:
    print('ERR')
\end{verbatim}
%enditem


\newpage
7. Нехай змінна \kw{s} позначає дуже довгу послідовність, по якій можна ітерувати. 

Написати вираз-генератор, що задає підпослідовність її елементів, що не є рядками типу \kw{str}.

Обмеження: усі елементи шуканої підпослідовності одночасно в пам'яті розміщені бути не можуть. 
%enditem


\newpage
8. Що буде надруковано за виконання?
%Оригинал был утерян. Требует отладки!
\begin{verbatim}
d = 0

class A:
    d = 1
    
    def __init__(self):
        self.d = 2
        d = 3
        A.d = 4

obj = A()
A.d = 5
try:
    print(d, end=' ')
    print(obj.d, end=' ')
    print(A.d, end=' ')
    print(obj.d, end=' ')
    print(obj._A__d)
except:
    print('ERR')
\end{verbatim}
%enditem


\newpage
9. Що буде надруковано за виконання?
\begin{verbatim}
class D0:
    def b(self): print(1, end=' ')
    def g(self): print(2, end=' ')

class D1(D0):
    def b(self): print(3, end=' ')

try:
    D1().g()
    D1().b()
except:
    print('ERR')
\end{verbatim}
%enditem


\newpage
10. Що буде надруковано за виконання?
\begin{verbatim}
class A:
    b = 1

    def _g1(self): print(2, end=' ')
    def _g2(self): print(3, end=' ')

    def main(self):
        try: print(self.b, end=' ')
        except: print('E1', end=' ')

        try: self._g1()
        except: print('E2', end=' ')

        try: self._g2()
        except: print('E3', end=' ')


class B(A):
    b = 4

    def _g1(self): print(5, end=' ')
    def _g2(self): print(6, end=' ')

try: B().main()
except: print('E4', end=' ')
\end{verbatim}
%enditem


\newpage
11. Змінні \kw{next} та \kw{prev} вузла списку позначають наступний та попередній за ним вузол відповідно. Починаючи з голови, у список записано послідовні цілі числа від~$-63$ до~$39$. Нехай змінна \kw{p} позначає вузол, що зберігає значення~$-4$.
Яке значення зберігається у вузлі \kw{p.prev.prev.prev.prev.prev} ?
%enditem


\newpage
12. 
\begin{center}
	\adjustimage{max size={0.8\linewidth}{0.8\paperheight}}
	{BinTree/trees_exp/31.png}
\end{center}
Указати послідовність міток вершин дерева, зображеного на рис., що утвориться в результаті обходу дерева в прямому порядку. Мітки записувати через пробіл.\par Алгоритм починає роботу з кореня (на рис. це верхня вершина).
%enditem


\newpage
13. 
\begin{center}
	\adjustimage{max size={0.9\linewidth}{0.9\paperheight}}
	{Graphs/AdjList/200.png}
\end{center}
Для графа, зображеного на рис., вказати список суміжності вершини 1.\par
Формат оформлення відповіді:\par
список суміжності виписувати через пробіл на одному рядку, якщо пробілів десь буде більше одного - не важливо, упорядкування вершин у списку також значення не має.
%enditem


\newpage
14. Для графа на рисунку з вершини 1 запущено алгоритм пошуку в глибину, що будує кістякове дерево. Списки суміжності вершин переглядаються алгоритмом за зростанням міток вершин.\par
\begin{center}
	\adjustimage{max size={0.9\linewidth}{0.9\paperheight}}
	{Graphs/Traversals/268.png}
\end{center}
\par Які з наведених ребер входять до складу отриманого кістякового дерева?\par
1) (0, 6)\par
2) (0, 2)\par
3) (1, 3)\par
4) (4, 6)\par
5) (3, 5)\par
6) (2, 4)\par
{\vskip 15pt} Номери відповідних ребер записати за зростанням через пробіл.

%enditem


\newpage
15. Для графа на рисунку з вершини 4 запущено алгоритм пошуку в ширину, що будує кістякове дерево. Списки суміжності вершин переглядаються алгоритмом за зростанням міток вершин.\par
\begin{center}
	\adjustimage{max size={0.9\linewidth}{0.9\paperheight}}
	{Graphs/Traversals/295.png}
\end{center}
\par Які з наведених ребер входять до складу отриманого кістякового дерева?\par
1) (0, 5)\par
2) (1, 6)\par
3) (4, 6)\par
4) (2, 4)\par
5) (0, 1)\par
6) (3, 4)\par
{\vskip 15pt} Номери відповідних ребер записати за зростанням через пробіл.


%enditem


\newpage
16. 

Рядки журналу через двокрапку містять ім'я користувача (ідентифікатор), час події,тип події, код події, наприклад
\begin{verbatim}
s="username : 2022-05-05 13:26 : ERROR : 404"
\end{verbatim}
у коді 
\begin{verbatim}
res = re.fullmatch('???', s) 
s = ''
code_str = ??? 
\end{verbatim}
чим треба замінити кожен з ???, щоб значенням code\_str був код події? 


\vspace{10mm}
\textbf{Завдання 17} 

Задано програму, що містить код 1-2 класів, можливе успадкування, та клієнтський код.
У наведеному коді невеликі частини закриті. 
Також наведено вихід програми.
Необхідно відновити закриті частини, щоб отриманий код відповідав наведеному виходу програми.


\textbf{Завдання 18-19} 

Кілька запитань за завданнями для самостійної роботи.
\vspace{10mm}
Скоріш за все, що завдань буде трохи менше, а деякі попередні будуть переформульовані в стилі двох останніх.

%enditem

\end {large}}

%end

\newpage
%begin

{{
\begin {large}

\newpage
1. Що буде надруковано за виконання?

\begin{verbatim}
print('R', end=' ')
f=('a','b','c','d','e',0,1,2,4)
print(f[11])
\end{verbatim}

%enditem


\newpage
2. Що буде надруковано за виконання?

\begin{verbatim}
print('R', end=' ')
g=['0','1','2','3','4','5','6','7','8','9']
print(g[8::-2])
\end{verbatim}

%enditem



\newpage
3. Що буде надруковано за виконання?

\begin{verbatim}
print('R', end=' ')
f=['0','1','2','3','4']
print(f.index('3'), end=' ')
\end{verbatim}

%enditem



\newpage
4. Що буде надруковано за виконання?

\begin{verbatim}
print('R', end=' ')
e=['0','1','2','3','4']
e[-6:-4]=[4,5]
print(e)
\end{verbatim}

%enditem



\newpage
5. Укажіть ТИП значення виразу.\par
\begin{verbatim}
{(1,4), 8}
\end{verbatim}
%enditem


\newpage
6. Що буде надруковано за виконання?\par
\begin{verbatim}
try:
    t = {34:2, 13:9, 22:0, 34:8}
    t[69] = 2
    for x in t.keys():
        print(x, sep=' ', end=' ')
except:
    print('ERR')
\end{verbatim}
%enditem


\newpage
7. 
Записати ВИРАЗ, значенням якого є список, що складається з цілих чисел від 10 до 2008 (включно), що діляться на 2.
%enditem


\newpage
8. Що буде надруковано за виконання?
\begin{verbatim}
d = 0
k = 1

class C:
    d = 2
    
    def __init__(self):
        self.d = 3
        d = 4
        C.d = 5

obj = C()
try:
    print(d, end=' ')
    print(k, end=' ')
    print(C.d, end=' ')
    print(obj.d)
except:
    print('ERR')
\end{verbatim}
%enditem


\newpage
9. Що буде надруковано за виконання?
\begin{verbatim}
class B0:
    def d(self): print(1, end=' ')
    def f(self): print(2, end=' ')

class B1(B0):
    def d(self): print(3, end=' ')

try:
    B1().f()
    B1().d()
except:
    print('ERR')
\end{verbatim}
%enditem


\newpage
10. Що буде надруковано за виконання?
\begin{verbatim}
class A:
    b = 1

    def __init__(self):
        c = 2
        c = 3

class B(A):
    c = 4
    a = 5

    def __init__(self):
        super().__init__()
        self.a = 6

try:
    obj = B()
    obj.c = obj.b + 10
    B.c = 13
except: print('EE', end=' ')

try: print(obj.c, end= ' ')
except: print('E1', end=' ')

try: print(A.c, end= ' ')
except: print('E2', end=' ')

try: print(B.c, end= ' ')
except: print('E3', end=' ')
\end{verbatim}
%enditem


\newpage
11. Змінні \kw{next} та \kw{prev} вузла списку позначають наступний та попередній за ним вузол відповідно. Починаючи з голови, у список записано послідовні цілі числа від~$-20$ до~$16$. Нехай змінна \kw{p} позначає вузол, що зберігає значення~$9$.
Яке значення зберігається у вузлі \kw{p.next.next.next.next.next} ?
%enditem


\newpage
12. 
\begin{center}
	\adjustimage{max size={0.8\linewidth}{0.8\paperheight}}
	{BinTree/trees_exp/152.png}
\end{center}
Указати послідовність міток вершин дерева, зображеного на рис., що утвориться в результаті обходу дерева в оберненому порядку. Мітки записувати через пробіл.\par Алгоритм починає роботу з кореня (на рис. це верхня вершина).
%enditem


\newpage
13. 
\begin{center}
	\adjustimage{max size={0.9\linewidth}{0.9\paperheight}}
	{Graphs/AdjList/237.png}
\end{center}
Для графа, зображеного на рис., вказати список суміжності вершини 8.\par
Формат оформлення відповіді:\par
список суміжності виписувати через пробіл на одному рядку, якщо пробілів десь буде більше одного - не важливо, упорядкування вершин у списку також значення не має.
%enditem


\newpage
14. Для графа на рисунку вказати послідовність, в якій алгоритм пошуку в глибину буде відмічати відвідані вершини, якщо списки суміжності вершин переглядаються алгоритмом за зростанням міток вершин та вершини відмічаються на першому вход\-женні у вершину.\par
Пошук розпочинається з вершини 3.\par

\begin{center}
	\adjustimage{max size={0.9\linewidth}{0.9\paperheight}}
	{Graphs/Traversals/348.png}
\end{center}
%enditem


\newpage
15. Для графа на рисунку з вершини 0 запущено алгоритм пошуку в ширину, що будує кістякове дерево. Списки суміжності вершин переглядаються алгоритмом за зростанням міток вершин.\par
\begin{center}
	\adjustimage{max size={0.9\linewidth}{0.9\paperheight}}
	{Graphs/Traversals/44.png}
\end{center}
\par Які з наведених ребер входять до складу отриманого кістякового дерева?\par
1) (0, 2)\par
2) (1, 4)\par
3) (4, 6)\par
4) (4, 5)\par
5) (1, 5)\par
6) (2, 5)\par
{\vskip 15pt} Номери відповідних ребер записати за зростанням через пробіл.


%enditem


\newpage
16. 

Рядки журналу через двокрапку містять ім'я користувача (ідентифікатор), час події,тип події, код події, наприклад
\begin{verbatim}
s="username : 2022-05-05 13:26 : ERROR : 404"
\end{verbatim}
у коді 
\begin{verbatim}
res = re.fullmatch('???', s) 
s = ''
code_str = ??? 
\end{verbatim}
чим треба замінити кожен з ???, щоб значенням code\_str був код події? 


\vspace{10mm}
\textbf{Завдання 17} 

Задано програму, що містить код 1-2 класів, можливе успадкування, та клієнтський код.
У наведеному коді невеликі частини закриті. 
Також наведено вихід програми.
Необхідно відновити закриті частини, щоб отриманий код відповідав наведеному виходу програми.


\textbf{Завдання 18-19} 

Кілька запитань за завданнями для самостійної роботи.
\vspace{10mm}
Скоріш за все, що завдань буде трохи менше, а деякі попередні будуть переформульовані в стилі двох останніх.

%enditem

\end {large}}

%end

\newpage
%begin

{{
\begin {large}

\newpage
1. Що буде надруковано за виконання?

\begin{verbatim}
print('R', end=' ')
g=('a','b','c','d','e',0,1)
print(g[-7])
\end{verbatim}

%enditem


\newpage
2. Що буде надруковано за виконання?

\begin{verbatim}
print('R', end=' ')
a=[0,1,2,3,4,5,6,7,8,9]
print(a[4:-8:-1])
\end{verbatim}

%enditem



\newpage
3. Що буде надруковано за виконання?

\begin{verbatim}
print('R', end=' ')
b=['0','1','2','3','4']
b.insert(6, [2,3])
print(b)
\end{verbatim}

%enditem



\newpage
4. Що буде надруковано за виконання?

\begin{verbatim}
print('R', end=' ')
f=[10,11,12,13,14]
del f[4:0]
print(f)
\end{verbatim}

%enditem



\newpage
5. Укажіть ТИП значення виразу.\par
\begin{verbatim}
{x for x in range(7)}
\end{verbatim}
%enditem


\newpage
6. Що буде надруковано за виконання?\par
\begin{verbatim}
try:
    s = {82:1, 41:4, 34:1, 66:6, 34:9}
    s[39] = 0
    for x in s.items():
        print(*x, sep=' ', end=' ')
except:
    print('ERR')
\end{verbatim}
%enditem


\newpage
7. Нехай змінна \kw{s} позначає послідовність цілих, по якій можна ітерувати. 

Написати вираз, що перевіряє, чи правда, що послідовність  містить хоча б одне число, що ділиться на 7.
%enditem


\newpage
8. Що буде надруковано за виконання?
\begin{verbatim}
b = 0
n = 1

class D:
    b = 2
    
    def __init__(self):
        global b
        self.b = 3
        b = 4
        D.b = 5

obj = D()
try:
    print(b, end=' ')
    print(n, end=' ')
    print(D.b, end=' ')
    print(obj.b)
except:
    print('ERR')
\end{verbatim}
%enditem


\newpage
9. Що буде надруковано за виконання?
\begin{verbatim}
class A0:
    def a(self): print(1, end=' ')
    def h(self): print(2, end=' ')

class A1(A0):
    def a(self): print(3, end=' ')

try:
    A0().h()
    A1().a()
except:
    print('ERR')
\end{verbatim}
%enditem


\newpage
10. Що буде надруковано за виконання?
\begin{verbatim}
class A:
    a = 1

    def _f1(self): print(2, end=' ')
    def __f2(self): print(3, end=' ')

    def main(self):
        try: print(self.a, end=' ')
        except: print('E1', end=' ')

        try: self._f1()
        except: print('E2', end=' ')

        try: self.__f2()
        except: print('E3', end=' ')


class B(A):
    a = 4

    def _f1(self): print(5, end=' ')
    def __f2(self): print(6, end=' ')

try: B().main()
except: print('E4', end=' ')
\end{verbatim}
%enditem


\newpage
11. Змінні \kw{next} та \kw{prev} вузла списку позначають наступний та попередній за ним вузол відповідно. Починаючи з голови, у список записано послідовні цілі числа від~$-57$ до~$47$. Нехай змінна \kw{p} позначає вузол, що зберігає значення~$5$.
Яке значення зберігається у вузлі \kw{p.prev.next.next.next.prev} ?
%enditem


\newpage
12. 
\begin{center}
	\adjustimage{max size={0.8\linewidth}{0.8\paperheight}}
	{BinTree/trees_exp/151.png}
\end{center}
Указати послідовність міток вершин дерева, зображеного на рис., що утвориться в результаті обходу дерева в прямому порядку. Мітки записувати через пробіл.\par Алгоритм починає роботу з кореня (на рис. це верхня вершина).
%enditem


\newpage
13. 
\begin{center}
	\adjustimage{max size={0.9\linewidth}{0.9\paperheight}}
	{Graphs/AdjList/115.png}
\end{center}
Для графа, зображеного на рис., вказати список суміжності вершини 1.\par
Формат оформлення відповіді:\par
список суміжності виписувати через пробіл на одному рядку, якщо пробілів десь буде більше одного - не важливо, упорядкування вершин у списку також значення не має.
%enditem


\newpage
14. Для графа на рисунку з вершини 6 запущено алгоритм пошуку в глибину, що будує кістякове дерево. Списки суміжності вершин переглядаються алгоритмом за зростанням міток вершин.\par
\begin{center}
	\adjustimage{max size={0.9\linewidth}{0.9\paperheight}}
	{Graphs/Traversals/214.png}
\end{center}
\par Які з наведених ребер входять до складу отриманого кістякового дерева?\par
1) (1, 6)\par
2) (3, 6)\par
3) (2, 4)\par
4) (5, 6)\par
5) (0, 6)\par
6) (4, 5)\par
{\vskip 15pt} Номери відповідних ребер записати за зростанням через пробіл.

%enditem


\newpage
15. Для графа на рисунку з вершини 5 запущено алгоритм пошуку в ширину, що будує кістякове дерево. Списки суміжності вершин переглядаються алгоритмом за зростанням міток вершин.\par
\begin{center}
	\adjustimage{max size={0.9\linewidth}{0.9\paperheight}}
	{Graphs/Traversals/51.png}
\end{center}
\par Які з наведених ребер входять до складу отриманого кістякового дерева?\par
1) (4, 5)\par
2) (3, 5)\par
3) (1, 2)\par
4) (0, 1)\par
5) (4, 6)\par
6) (3, 4)\par
{\vskip 15pt} Номери відповідних ребер записати за зростанням через пробіл.


%enditem


\newpage
16. 

Рядки журналу через двокрапку містять ім'я користувача (ідентифікатор), час події,тип події, код події, наприклад
\begin{verbatim}
s="username : 2022-05-05 13:26 : ERROR : 404"
\end{verbatim}
у коді 
\begin{verbatim}
res = re.fullmatch('???', s) 
s = ''
code_str = ??? 
\end{verbatim}
чим треба замінити кожен з ???, щоб значенням code\_str був код події? 


\vspace{10mm}
\textbf{Завдання 17} 

Задано програму, що містить код 1-2 класів, можливе успадкування, та клієнтський код.
У наведеному коді невеликі частини закриті. 
Також наведено вихід програми.
Необхідно відновити закриті частини, щоб отриманий код відповідав наведеному виходу програми.


\textbf{Завдання 18-19} 

Кілька запитань за завданнями для самостійної роботи.
\vspace{10mm}
Скоріш за все, що завдань буде трохи менше, а деякі попередні будуть переформульовані в стилі двох останніх.

%enditem

\end {large}}

%end

\newpage
%begin

{{
\begin {large}

\newpage
1. Що буде надруковано за виконання?

\begin{verbatim}
print('R', end=' ')
b=('a','b','c','d','e',0,1,2,3,4)
print(b[6])
\end{verbatim}

%enditem


\newpage
2. Що буде надруковано за виконання?

\begin{verbatim}
print('R', end=' ')
h=(0,1,2,3,4,5,6,7,8,9)
print(h[-1:-1:3])
\end{verbatim}

%enditem



\newpage
3. Що буде надруковано за виконання?

\begin{verbatim}
print('R', end=' ')
c=[10,11,12,13,14]
c+=[1,0]
print(c)
\end{verbatim}

%enditem



\newpage
4. Що буде надруковано за виконання?

\begin{verbatim}
print('R', end=' ')
b=['0','1','2','3','4']
b[1:]=()
print(b)
\end{verbatim}

%enditem



\newpage
5. Укажіть ТИП значення виразу.\par
\begin{verbatim}
{1:3, 5:8}
\end{verbatim}
%enditem


\newpage
6. Що буде надруковано за виконання?\par
\begin{verbatim}
try:
    d = {31:6, 41:8, 10:9, 40:1, 31:9}
    d[83] = 5
    for x in d.values():
        print(x, sep=' ', end=' ')
except:
    print('ERR')
\end{verbatim}
%enditem


\newpage
7. 
Записати ВИРАЗ, значенням якого є список, що складається з квадратів цілих чисел від 15 до 2006 (включно), що не діляться на 3.
%enditem


\newpage
8. Що буде надруковано за виконання?
\begin{verbatim}
a = 0
k = 1

class B:
    k = 2
    
    def __init__(self):
        self.k = 3
        k = 4
        B.k = 5

obj = B()
try:
    print(a, end=' ')
    print(k, end=' ')
    print(B.k, end=' ')
    print(obj.k)
except:
    print('ERR')
\end{verbatim}
%enditem


\newpage
9. Що буде надруковано за виконання?
\begin{verbatim}
class C0:
    def c(self): print(1, end=' ')
    def g(self): print(2, end=' ')

class C1(C0):
    def c(self): print(3, end=' ')

try:
    C1().g()
    C0().c()
except:
    print('ERR')
\end{verbatim}
%enditem


\newpage
10. Що буде надруковано за виконання?
\begin{verbatim}
class A:
    c = 1

    def _h1(self): print(2, end=' ')
    def h2(self): print(3, end=' ')

    def main(self):
        try: print(self.c, end=' ')
        except: print('E1', end=' ')

        try: self._h1()
        except: print('E2', end=' ')

        try: self.h2()
        except: print('E3', end=' ')


class B(A):
    c = 4

    def _h1(self): print(5, end=' ')
    def h2(self): print(6, end=' ')

try: B().main()
except: print('E4', end=' ')
\end{verbatim}
%enditem


\newpage
11. Змінні \kw{next} та \kw{prev} вузла списку позначають наступний та попередній за ним вузол відповідно. Починаючи з голови, у список записано послідовні цілі числа від~$-49$ до~$38$. Нехай змінна \kw{p} позначає вузол, що зберігає значення~$-8$.
Яке значення зберігається у вузлі \kw{p.next.prev.prev.prev.next} ?
%enditem


\newpage
12. 
\begin{center}
	\adjustimage{max size={0.8\linewidth}{0.8\paperheight}}
	{BinTree/trees_exp/24.png}
\end{center}
Указати послідовність міток вершин дерева, зображеного на рис., що утвориться в результаті обходу дерева в оберненому порядку. Мітки записувати через пробіл.\par Алгоритм починає роботу з кореня (на рис. це верхня вершина).
%enditem


\newpage
13. 
\begin{center}
	\adjustimage{max size={0.9\linewidth}{0.9\paperheight}}
	{Graphs/AdjList/104.png}
\end{center}
Для графа, зображеного на рис., вказати список суміжності вершини 3.\par
Формат оформлення відповіді:\par
список суміжності виписувати через пробіл на одному рядку, якщо пробілів десь буде більше одного - не важливо, упорядкування вершин у списку також значення не має.
%enditem


\newpage
14. Для графа на рисунку вказати послідовність, в якій алгоритм пошуку в глибину буде відмічати відвідані вершини, якщо списки суміжності вершин переглядаються алгоритмом за зростанням міток вершин та вершини відмічаються на першому вход\-женні у вершину.\par
Пошук розпочинається з вершини 2.\par

\begin{center}
	\adjustimage{max size={0.9\linewidth}{0.9\paperheight}}
	{Graphs/Traversals/141.png}
\end{center}
%enditem


\newpage
15. Для графа на рисунку з вершини 1 запущено алгоритм пошуку в ширину, що будує кістякове дерево. Списки суміжності вершин переглядаються алгоритмом за зростанням міток вершин.\par
\begin{center}
	\adjustimage{max size={0.9\linewidth}{0.9\paperheight}}
	{Graphs/Traversals/56.png}
\end{center}
\par Які з наведених ребер входять до складу отриманого кістякового дерева?\par
1) (0, 3)\par
2) (0, 2)\par
3) (2, 5)\par
4) (3, 4)\par
5) (1, 2)\par
6) (0, 5)\par
{\vskip 15pt} Номери відповідних ребер записати за зростанням через пробіл.


%enditem


\newpage
16. 

Рядки журналу через двокрапку містять ім'я користувача (ідентифікатор), час події,тип події, код події, наприклад
\begin{verbatim}
s="username : 2022-05-05 13:26 : ERROR : 404"
\end{verbatim}
у коді 
\begin{verbatim}
res = re.fullmatch('???', s) 
s = ''
code_str = ??? 
\end{verbatim}
чим треба замінити кожен з ???, щоб значенням code\_str був код події? 


\vspace{10mm}
\textbf{Завдання 17} 

Задано програму, що містить код 1-2 класів, можливе успадкування, та клієнтський код.
У наведеному коді невеликі частини закриті. 
Також наведено вихід програми.
Необхідно відновити закриті частини, щоб отриманий код відповідав наведеному виходу програми.


\textbf{Завдання 18-19} 

Кілька запитань за завданнями для самостійної роботи.
\vspace{10mm}
Скоріш за все, що завдань буде трохи менше, а деякі попередні будуть переформульовані в стилі двох останніх.

%enditem

\end {large}}

%end

\newpage
%begin

{{
\begin {large}

\newpage
1. Що буде надруковано за виконання?

\begin{verbatim}
print('R', end=' ')
e=('a','b','c','d','e',0,1,2)
print(e[2])
\end{verbatim}

%enditem


\newpage
2. Що буде надруковано за виконання?

\begin{verbatim}
print('R', end=' ')
c=('0','1','2','3','4','5','6','7','8','9')
print(c[:8:-3])
\end{verbatim}

%enditem



\newpage
3. Що буде надруковано за виконання?

\begin{verbatim}
print('R', end=' ')
e=['0','1','2','3','4']
print(e.pop(-6), end=' ')
print(e)
\end{verbatim}

%enditem



\newpage
4. Що буде надруковано за виконання?

\begin{verbatim}
print('R', end=' ')
d=[10,11,12,13,14]
d[:4]=[1,0]
print(d)
\end{verbatim}

%enditem



\newpage
5. Укажіть ТИП значення виразу.\par
\begin{verbatim}
{x:x for x in range(7)}
\end{verbatim}
%enditem


\newpage
6. Що буде надруковано за виконання?\par
\begin{verbatim}
try:
    t = {52:3, 48:2, 73:4, 41:6, 43:9}
    t[73] = 3
    for x in t :
        print(x, sep=' ', end=' ')
except:
    print('ERR')
\end{verbatim}
%enditem


\newpage
7. 
Записати ВИРАЗ, значенням якого є список, що складається з кубів цілих чисел від 20 до 2010 (включно), що діляться на 7, записаних в обернено\-му порядку.
%enditem


\newpage
8. Що буде надруковано за виконання?
%Оригинал был утерян. Требует отладки!
\begin{verbatim}
_c = 0

class C:
    _c = 1
    
    def __init__(self):
        self._c = 2
        _c = 3
        C._c = 4

obj = C()
C._c = 5
try:
    print(_c, end=' ')
    print(obj._c, end=' ')
    print(C._c, end=' ')
    print(obj._c, end=' ')
    print(obj._C__c)
except:
    print('ERR')
\end{verbatim}
%enditem


\newpage
9. Що буде надруковано за виконання?
\begin{verbatim}
class D0:
    def b(self): print(1, end=' ')
    def h(self): print(2, end=' ')

class D1(D0):
    def b(self): print(3, end=' ')

try:
    D1().h()
    D0().b()
except:
    print('ERR')
\end{verbatim}
%enditem


\newpage
10. Що буде надруковано за виконання?
\begin{verbatim}
class A:
    a = 1

    def __init__(self):
        b = 2
        b = 3

class B(A):
    b = 4
    c = 5

    def __init__(self):
        super().__init__()
        self.c = 6

try:
    obj = B()
    obj.c = obj.a + 10
    B.c = 13
except: print('EE', end=' ')

try: print(obj.c, end= ' ')
except: print('E1', end=' ')

try: print(A.c, end= ' ')
except: print('E2', end=' ')

try: print(B.c, end= ' ')
except: print('E3', end=' ')
\end{verbatim}
%enditem


\newpage
11. Змінні \kw{next} та \kw{prev} вузла списку позначають наступний та попередній за ним вузол відповідно. Починаючи з голови, у список записано послідовні цілі числа від~$-29$ до~$37$. Нехай змінна \kw{p} позначає вузол, що зберігає значення~$8$.
Яке значення зберігається у вузлі \kw{p.prev.next.next.prev.prev} ?
%enditem


\newpage
12. 
\begin{center}
	\adjustimage{max size={0.8\linewidth}{0.8\paperheight}}
	{BinTree/trees_exp/119.png}
\end{center}
Указати послідовність міток вершин дерева, зображеного на рис., що утвориться в результаті обходу дерева в оберненому порядку. Мітки записувати через пробіл.\par Алгоритм починає роботу з кореня (на рис. це верхня вершина).
%enditem


\newpage
13. 
\begin{center}
	\adjustimage{max size={0.9\linewidth}{0.9\paperheight}}
	{Graphs/AdjList/267.png}
\end{center}
Для графа, зображеного на рис., вказати список суміжності вершини 0.\par
Формат оформлення відповіді:\par
список суміжності виписувати через пробіл на одному рядку, якщо пробілів десь буде більше одного - не важливо, упорядкування вершин у списку також значення не має.
%enditem


\newpage
14. Для графа на рисунку вказати послідовність, в якій алгоритм пошуку в глибину буде відмічати відвідані вершини, якщо списки суміжності вершин переглядаються алгоритмом за зростанням міток вершин та вершини відмічаються на першому вход\-женні у вершину.\par
Пошук розпочинається з вершини 1.\par

\begin{center}
	\adjustimage{max size={0.9\linewidth}{0.9\paperheight}}
	{Graphs/Traversals/109.png}
\end{center}
%enditem


\newpage
15. Для графа на рисунку з вершини 2 запущено алгоритм пошуку в ширину, що будує кістякове дерево. Списки суміжності вершин переглядаються алгоритмом за зростанням міток вершин.\par
\begin{center}
	\adjustimage{max size={0.9\linewidth}{0.9\paperheight}}
	{Graphs/Traversals/249.png}
\end{center}
\par Які з наведених ребер входять до складу отриманого кістякового дерева?\par
1) (3, 4)\par
2) (0, 5)\par
3) (0, 4)\par
4) (1, 6)\par
5) (0, 6)\par
6) (2, 3)\par
{\vskip 15pt} Номери відповідних ребер записати за зростанням через пробіл.


%enditem


\newpage
16. 

Рядки журналу через двокрапку містять ім'я користувача (ідентифікатор), час події,тип події, код події, наприклад
\begin{verbatim}
s="username : 2022-05-05 13:26 : ERROR : 404"
\end{verbatim}
у коді 
\begin{verbatim}
res = re.fullmatch('???', s) 
s = ''
code_str = ??? 
\end{verbatim}
чим треба замінити кожен з ???, щоб значенням code\_str був код події? 


\vspace{10mm}
\textbf{Завдання 17} 

Задано програму, що містить код 1-2 класів, можливе успадкування, та клієнтський код.
У наведеному коді невеликі частини закриті. 
Також наведено вихід програми.
Необхідно відновити закриті частини, щоб отриманий код відповідав наведеному виходу програми.


\textbf{Завдання 18-19} 

Кілька запитань за завданнями для самостійної роботи.
\vspace{10mm}
Скоріш за все, що завдань буде трохи менше, а деякі попередні будуть переформульовані в стилі двох останніх.

%enditem

\end {large}}

%end

\newpage
%begin

{{
\begin {large}

\newpage
1. Що буде надруковано за виконання?

\begin{verbatim}
print('R', end=' ')
h=('a','b','c','d','e',0,1,2,3)
print(h[-6])
\end{verbatim}

%enditem


\newpage
2. Що буде надруковано за виконання?

\begin{verbatim}
print('R', end=' ')
e='0123456789'
print(e[3:-11:2])
\end{verbatim}

%enditem



\newpage
3. Що буде надруковано за виконання?

\begin{verbatim}
print('R', end=' ')
d=[10,11,12,13,14]
d.remove(11)
print(d)
\end{verbatim}

%enditem



\newpage
4. Що буде надруковано за виконання?

\begin{verbatim}
print('R', end=' ')
g=['0','1','2','3','4']
del g[8:6]
print(g)
\end{verbatim}

%enditem



\newpage
5. Укажіть ТИП значення виразу.\par
\begin{verbatim}
{y:13 for y in range(-7,7)}
\end{verbatim}
%enditem


\newpage
6. Що буде надруковано за виконання?\par
\begin{verbatim}
try:
    s = {75:4, 53:4, 71:1, 71:4}
    s[85] = 8
    for x in s.items():
        print(x, sep=' ', end=' ')
except:
    print('ERR')
\end{verbatim}
%enditem


\newpage
7. 
Записати ВИРАЗ, значенням якого є список, що складається з цілих чисел від 18 до 2000 (включно), причому спочатку за зростанням записано всі, що діляться на 3, а потім решту за спаданням.

%enditem


\newpage
8. Що буде надруковано за виконання?
%Оригинал был утерян. Требует отладки!
\begin{verbatim}
d = 0

class D:
    d = 1
    
    def __init__(self):
        self.d = 2
        d = 3
        D.d = 4

obj = D()
obj.d = 5
try:
    print(d, end=' ')
    print(obj.d, end=' ')
    print(D.d, end=' ')
    print(obj.d, end=' ')
    print(obj._D__d)
except:
    print('ERR')
\end{verbatim}
%enditem


\newpage
9. Що буде надруковано за виконання?
\begin{verbatim}
class C0:
    def d(self): print(1, end=' ')
    def f(self): print(2, end=' ')

class C1(C0):
    def f(self): print(3, end=' ')

try:
    C0().d()
    C1().f()
except:
    print('ERR')
\end{verbatim}
%enditem


\newpage
10. Що буде надруковано за виконання?
\begin{verbatim}
class A:
    b = 1

    def __init__(self):
        a = 2
        b = 3

class B(A):
    b = 4
    c = 5

    def __init__(self):
        super().__init__()
        self.c = 6

try:
    obj = B()
    obj.b = obj.c + 10
    B.b = 13
except: print('EE', end=' ')

try: print(obj.b, end= ' ')
except: print('E1', end=' ')

try: print(A.b, end= ' ')
except: print('E2', end=' ')

try: print(B.b, end= ' ')
except: print('E3', end=' ')
\end{verbatim}
%enditem


\newpage
11. Змінні \kw{next} та \kw{prev} вузла списку позначають наступний та попередній за ним вузол відповідно. Починаючи з голови, у список записано послідовні цілі числа від~$-46$ до~$36$. Нехай змінна \kw{p} позначає вузол, що зберігає значення~$-9$.
Яке значення зберігається у вузлі \kw{p.next.prev.prev.next.next} ?
%enditem


\newpage
12. 
\begin{center}
	\adjustimage{max size={0.8\linewidth}{0.8\paperheight}}
	{BinTree/trees_exp/41.png}
\end{center}
Указати послідовність міток вершин дерева, зображеного на рис., що утвориться в результаті обходу дерева в прямому порядку. Мітки записувати через пробіл.\par Алгоритм починає роботу з кореня (на рис. це верхня вершина).
%enditem


\newpage
13. 
\begin{center}
	\adjustimage{max size={0.9\linewidth}{0.9\paperheight}}
	{Graphs/AdjList/262.png}
\end{center}
Для графа, зображеного на рис., вказати список суміжності вершини 1.\par
Формат оформлення відповіді:\par
список суміжності виписувати через пробіл на одному рядку, якщо пробілів десь буде більше одного - не важливо, упорядкування вершин у списку також значення не має.
%enditem


\newpage
14. Для графа на рисунку вказати послідовність, в якій алгоритм пошуку в глибину буде відмічати відвідані вершини, якщо списки суміжності вершин переглядаються алгоритмом за зростанням міток вершин та вершини відмічаються на першому вход\-женні у вершину.\par
Пошук розпочинається з вершини 0.\par

\begin{center}
	\adjustimage{max size={0.9\linewidth}{0.9\paperheight}}
	{Graphs/Traversals/115.png}
\end{center}
%enditem


\newpage
15. Для графа на рисунку з вершини 2 запущено алгоритм пошуку в ширину, що будує кістякове дерево. Списки суміжності вершин переглядаються алгоритмом за зростанням міток вершин.\par
\begin{center}
	\adjustimage{max size={0.9\linewidth}{0.9\paperheight}}
	{Graphs/Traversals/14.png}
\end{center}
\par Які з наведених ребер входять до складу отриманого кістякового дерева?\par
1) (3, 4)\par
2) (2, 4)\par
3) (4, 6)\par
4) (2, 5)\par
5) (0, 5)\par
6) (3, 5)\par
{\vskip 15pt} Номери відповідних ребер записати за зростанням через пробіл.


%enditem


\newpage
16. 

Рядки журналу через двокрапку містять ім'я користувача (ідентифікатор), час події,тип події, код події, наприклад
\begin{verbatim}
s="username : 2022-05-05 13:26 : ERROR : 404"
\end{verbatim}
у коді 
\begin{verbatim}
res = re.fullmatch('???', s) 
s = ''
code_str = ??? 
\end{verbatim}
чим треба замінити кожен з ???, щоб значенням code\_str був код події? 


\vspace{10mm}
\textbf{Завдання 17} 

Задано програму, що містить код 1-2 класів, можливе успадкування, та клієнтський код.
У наведеному коді невеликі частини закриті. 
Також наведено вихід програми.
Необхідно відновити закриті частини, щоб отриманий код відповідав наведеному виходу програми.


\textbf{Завдання 18-19} 

Кілька запитань за завданнями для самостійної роботи.
\vspace{10mm}
Скоріш за все, що завдань буде трохи менше, а деякі попередні будуть переформульовані в стилі двох останніх.

%enditem

\end {large}}

%end

\newpage
%begin

{{
\begin {large}

\newpage
1. Що буде надруковано за виконання?

\begin{verbatim}
print('R', end=' ')
d=('a','b','c','d','e',0,1,2,3)
print(d[-11])
\end{verbatim}

%enditem


\newpage
2. Що буде надруковано за виконання?

\begin{verbatim}
print('R', end=' ')
f=[0,1,2,3,4,5,6,7,8,9]
print(f[:14:-3])
\end{verbatim}

%enditem



\newpage
3. Що буде надруковано за виконання?

\begin{verbatim}
print('R', end=' ')
h=[10,11,12,13,14]
print(h.index(9), end=' ')
\end{verbatim}

%enditem



\newpage
4. Що буде надруковано за виконання?

\begin{verbatim}
print('R', end=' ')
a=['0','1','2','3','4']
del a[-9:-5]
print(a)
\end{verbatim}

%enditem



\newpage
5. Укажіть ТИП значення виразу.\par
\begin{verbatim}
[x for x in range(7)]
\end{verbatim}
%enditem


\newpage
6. Що буде надруковано за виконання?\par
\begin{verbatim}
try:
    d = {77:9, 16:9, 84:5, 16:9}
    d[58] = 5
    for x in d.items():
        print(x, sep=' ', end=' ')
except:
    print('ERR')
\end{verbatim}
%enditem


\newpage
7. Нехай змінна \kw{s} позначає послідовність цілих, по якій можна ітерувати. 

Написати вираз, що перевіряє, чи правда, що послідовність  містить від’ємні числа.
%enditem


\newpage
8. Що буде надруковано за виконання?
%Оригинал был утерян. Требует отладки!
\begin{verbatim}
__b = 0

class B:
    __b = 1
    
    def __init__(self):
        self.__b = 2
        __b = 3
        B.__b = 4

obj = B()
B.__b = 5
try:
    print(__b, end=' ')
    print(obj.__b, end=' ')
    print(B.__b, end=' ')
    print(obj.__b, end=' ')
    print(obj._B__b)
except:
    print('ERR')
\end{verbatim}
%enditem


\newpage
9. Що буде надруковано за виконання?
\begin{verbatim}
class A0:
    def a(self): print(1, end=' ')
    def f(self): print(2, end=' ')

class A1(A0):
    def a(self): print(3, end=' ')

try:
    A1().f()
    A1().a()
except:
    print('ERR')
\end{verbatim}
%enditem


\newpage
10. Що буде надруковано за виконання?
\begin{verbatim}
class A:
    a = 1

    def __init__(self):
        c = 2
        a = 3

class B(A):
    c = 4
    b = 5

    def __init__(self):
        super().__init__()
        self.b = 6

try:
    obj = B()
    obj.a = obj.c + 10
    B.a = 13
except: print('EE', end=' ')

try: print(obj.a, end= ' ')
except: print('E1', end=' ')

try: print(A.a, end= ' ')
except: print('E2', end=' ')

try: print(B.a, end= ' ')
except: print('E3', end=' ')
\end{verbatim}
%enditem


\newpage
11. Змінні \kw{next} та \kw{prev} вузла списку позначають наступний та попередній за ним вузол відповідно. Починаючи з голови, у список записано послідовні цілі числа від~$-56$ до~$24$. Нехай змінна \kw{p} позначає вузол, що зберігає значення~$4$.
Яке значення зберігається у вузлі \kw{p.next.prev.next.next.next} ?
%enditem


\newpage
12. 
\begin{center}
	\adjustimage{max size={0.8\linewidth}{0.8\paperheight}}
	{BinTree/trees_exp/120.png}
\end{center}
Указати послідовність міток вершин дерева, зображеного на рис., що утвориться в результаті обходу дерева в оберненому порядку. Мітки записувати через пробіл.\par Алгоритм починає роботу з кореня (на рис. це верхня вершина).
%enditem


\newpage
13. 
\begin{center}
	\adjustimage{max size={0.9\linewidth}{0.9\paperheight}}
	{Graphs/AdjList/90.png}
\end{center}
Для графа, зображеного на рис., вказати список суміжності вершини 1.\par
Формат оформлення відповіді:\par
список суміжності виписувати через пробіл на одному рядку, якщо пробілів десь буде більше одного - не важливо, упорядкування вершин у списку також значення не має.
%enditem


\newpage
14. Для графа на рисунку вказати послідовність, в якій алгоритм пошуку в глибину буде відмічати відвідані вершини, якщо списки суміжності вершин переглядаються алгоритмом за зростанням міток вершин та вершини відмічаються на першому вход\-женні у вершину.\par
Пошук розпочинається з вершини 6.\par

\begin{center}
	\adjustimage{max size={0.9\linewidth}{0.9\paperheight}}
	{Graphs/Traversals/87.png}
\end{center}
%enditem


\newpage
15. Для графа на рисунку з вершини 3 запущено алгоритм пошуку в ширину, що будує кістякове дерево. Списки суміжності вершин переглядаються алгоритмом за зростанням міток вершин.\par
\begin{center}
	\adjustimage{max size={0.9\linewidth}{0.9\paperheight}}
	{Graphs/Traversals/29.png}
\end{center}
\par Які з наведених ребер входять до складу отриманого кістякового дерева?\par
1) (5, 6)\par
2) (3, 4)\par
3) (1, 6)\par
4) (2, 5)\par
5) (0, 3)\par
6) (4, 6)\par
{\vskip 15pt} Номери відповідних ребер записати за зростанням через пробіл.


%enditem


\newpage
16. 

Рядки журналу через двокрапку містять ім'я користувача (ідентифікатор), час події,тип події, код події, наприклад
\begin{verbatim}
s="username : 2022-05-05 13:26 : ERROR : 404"
\end{verbatim}
у коді 
\begin{verbatim}
res = re.fullmatch('???', s) 
s = ''
code_str = ??? 
\end{verbatim}
чим треба замінити кожен з ???, щоб значенням code\_str був код події? 


\vspace{10mm}
\textbf{Завдання 17} 

Задано програму, що містить код 1-2 класів, можливе успадкування, та клієнтський код.
У наведеному коді невеликі частини закриті. 
Також наведено вихід програми.
Необхідно відновити закриті частини, щоб отриманий код відповідав наведеному виходу програми.


\textbf{Завдання 18-19} 

Кілька запитань за завданнями для самостійної роботи.
\vspace{10mm}
Скоріш за все, що завдань буде трохи менше, а деякі попередні будуть переформульовані в стилі двох останніх.

%enditem

\end {large}}

%end

\newpage
%begin

{{
\begin {large}

\newpage
1. Що буде надруковано за виконання?

\begin{verbatim}
print('R', end=' ')
a=('a','b','c','d','e',0,1,2)
print(a[-9])
\end{verbatim}

%enditem


\newpage
2. Що буде надруковано за виконання?

\begin{verbatim}
print('R', end=' ')
d=['0','1','2','3','4','5','6','7','8','9']
print(d[-12:4:2])
\end{verbatim}

%enditem



\newpage
3. Що буде надруковано за виконання?

\begin{verbatim}
print('R', end=' ')
a=['0','1','2','3','4']
a.insert(3, [4,5])
print(a)
\end{verbatim}

%enditem



\newpage
4. Що буде надруковано за виконання?

\begin{verbatim}
print('R', end=' ')
h=[10,11,12,13,14]
h[2:]=[1,2]
print(h)
\end{verbatim}

%enditem



\newpage
5. Укажіть ТИП значення виразу.\par
\begin{verbatim}
{x*x:x for x in range(-7,7)}
\end{verbatim}
%enditem


\newpage
6. Що буде надруковано за виконання?\par
\begin{verbatim}
try:
    t = {53:5, 19:8, 13:0, 53:2}
    t[13] = 8
    for x in t :
        print(x, sep=' ', end=' ')
except:
    print('ERR')
\end{verbatim}
%enditem


\newpage
7. Нехай змінна \kw{s} позначає дуже довгу послідовність цілих, по якій можна ітерувати. 

Написати вираз-генератор, що задає підпослідовність її елементів, що є від'єм\-ни\-ми. 

Обмеження: усі елементи шуканої підпослідовності одночасно в пам'яті розміщені бути не можуть. 
%enditem


\newpage
8. Що буде надруковано за виконання?
%Оригинал был утерян. Требует отладки!
\begin{verbatim}
__a = 0

class A:
    __a = 1
    
    def __init__(self):
        self.__a = 2
        __a = 3
        A.__a = 4

obj = A()
obj.__a = 5
try:
    print(__a, end=' ')
    print(obj.__a, end=' ')
    print(A.__a, end=' ')
    print(obj.__a, end=' ')
    print(obj._A__a)
except:
    print('ERR')
\end{verbatim}
%enditem


\newpage
9. Що буде надруковано за виконання?
\begin{verbatim}
class B0:
    def c(self): print(1, end=' ')
    def g(self): print(2, end=' ')

class B1(B0):
    def c(self): print(3, end=' ')

try:
    B1().g()
    B1().c()
except:
    print('ERR')
\end{verbatim}
%enditem


\newpage
10. Що буде надруковано за виконання?
\begin{verbatim}
class A:
    c = 1

    def __init__(self):
        a = 2
        a = 3

class B(A):
    c = 4
    b = 5

    def __init__(self):
        super().__init__()
        self.b = 6

try:
    obj = B()
    obj.a = obj.b + 10
    B.a = 13
except: print('EE', end=' ')

try: print(obj.a, end= ' ')
except: print('E1', end=' ')

try: print(A.a, end= ' ')
except: print('E2', end=' ')

try: print(B.a, end= ' ')
except: print('E3', end=' ')
\end{verbatim}
%enditem


\newpage
11. Змінні \kw{next} та \kw{prev} вузла списку позначають наступний та попередній за ним вузол відповідно. Починаючи з голови, у список записано послідовні цілі числа від~$-19$ до~$30$. Нехай змінна \kw{p} позначає вузол, що зберігає значення~$3$.
Яке значення зберігається у вузлі \kw{p.prev.next.prev.prev.prev} ?
%enditem


\newpage
12. 
\begin{center}
	\adjustimage{max size={0.8\linewidth}{0.8\paperheight}}
	{BinTree/trees_exp/15.png}
\end{center}
Указати послідовність міток вершин дерева, зображеного на рис., що утвориться в результаті обходу дерева в прямому порядку. Мітки записувати через пробіл.\par Алгоритм починає роботу з кореня (на рис. це верхня вершина).
%enditem


\newpage
13. 
\begin{center}
	\adjustimage{max size={0.9\linewidth}{0.9\paperheight}}
	{Graphs/AdjList/266.png}
\end{center}
Для графа, зображеного на рис., вказати список суміжності вершини 7.\par
Формат оформлення відповіді:\par
список суміжності виписувати через пробіл на одному рядку, якщо пробілів десь буде більше одного - не важливо, упорядкування вершин у списку також значення не має.
%enditem


\newpage
14. Для графа на рисунку з вершини 3 запущено алгоритм пошуку в глибину, що будує кістякове дерево. Списки суміжності вершин переглядаються алгоритмом за зростанням міток вершин.\par
\begin{center}
	\adjustimage{max size={0.9\linewidth}{0.9\paperheight}}
	{Graphs/Traversals/93.png}
\end{center}
\par Які з наведених ребер входять до складу отриманого кістякового дерева?\par
1) (1, 3)\par
2) (1, 5)\par
3) (0, 2)\par
4) (3, 5)\par
5) (3, 4)\par
6) (2, 4)\par
{\vskip 15pt} Номери відповідних ребер записати за зростанням через пробіл.

%enditem


\newpage
15. Для графа на рисунку з вершини 5 запущено алгоритм пошуку в ширину, що будує кістякове дерево. Списки суміжності вершин переглядаються алгоритмом за зростанням міток вершин.\par
\begin{center}
	\adjustimage{max size={0.9\linewidth}{0.9\paperheight}}
	{Graphs/Traversals/297.png}
\end{center}
\par Які з наведених ребер входять до складу отриманого кістякового дерева?\par
1) (3, 4)\par
2) (0, 6)\par
3) (5, 6)\par
4) (2, 3)\par
5) (4, 5)\par
6) (1, 5)\par
{\vskip 15pt} Номери відповідних ребер записати за зростанням через пробіл.


%enditem


\newpage
16. 

Рядки журналу через двокрапку містять ім'я користувача (ідентифікатор), час події,тип події, код події, наприклад
\begin{verbatim}
s="username : 2022-05-05 13:26 : ERROR : 404"
\end{verbatim}
у коді 
\begin{verbatim}
res = re.fullmatch('???', s) 
s = ''
code_str = ??? 
\end{verbatim}
чим треба замінити кожен з ???, щоб значенням code\_str був код події? 


\vspace{10mm}
\textbf{Завдання 17} 

Задано програму, що містить код 1-2 класів, можливе успадкування, та клієнтський код.
У наведеному коді невеликі частини закриті. 
Також наведено вихід програми.
Необхідно відновити закриті частини, щоб отриманий код відповідав наведеному виходу програми.


\textbf{Завдання 18-19} 

Кілька запитань за завданнями для самостійної роботи.
\vspace{10mm}
Скоріш за все, що завдань буде трохи менше, а деякі попередні будуть переформульовані в стилі двох останніх.

%enditem

\end {large}}

%end

\newpage
%begin

{{
\begin {large}

\newpage
1. Що буде надруковано за виконання?

\begin{verbatim}
print('R', end=' ')
h=('a','b','c','d','e',0,1,2,4)
print(h[-1])
\end{verbatim}

%enditem


\newpage
2. Що буде надруковано за виконання?

\begin{verbatim}
print('R', end=' ')
c=(0,1,2,3,4,5,6,7,8,9)
print(c[:15:-2])
\end{verbatim}

%enditem



\newpage
3. Що буде надруковано за виконання?

\begin{verbatim}
print('R', end=' ')
e=['0','1','2','3','4']
e.append(18)
print(e)
\end{verbatim}

%enditem



\newpage
4. Що буде надруковано за виконання?

\begin{verbatim}
print('R', end=' ')
e=[10,11,12,13,14]
e[:-7]=[2,3]
print(e)
\end{verbatim}

%enditem



\newpage
5. Укажіть ТИП значення виразу.\par
\begin{verbatim}
[7,2,3,4]
\end{verbatim}
%enditem


\newpage
6. Що буде надруковано за виконання?\par
\begin{verbatim}
try:
    s = {68:9, 31:4, 11:7, 11:3}
    s[11] = 6
    for x in s.values():
        print(x, sep=' ', end=' ')
except:
    print('ERR')
\end{verbatim}
%enditem


\newpage
7. Нехай змінна \kw{s} позначає дуже довгу послідовність, по якій можна ітерувати. 

Написати вираз-генератор, що задає підпослідовність її елементів, що не є числами типу \kw{int}.

Обмеження: усі елементи шуканої підпослідовності одночасно в пам'яті розміщені бути не можуть. 
%enditem


\newpage
8. Що буде надруковано за виконання?
%Оригинал был утерян. Требует отладки!
\begin{verbatim}
c = 0

class B:
    c = 1
    
    def __init__(self):
        self.c = 2
        c = 3
        B.c = 4

obj = B()
obj.c = 5
try:
    print(c, end=' ')
    print(obj.c, end=' ')
    print(B.c, end=' ')
    print(obj.c, end=' ')
    print(obj._B__c)
except:
    print('ERR')
\end{verbatim}
%enditem


\newpage
9. Що буде надруковано за виконання?
\begin{verbatim}
class A0:
    def c(self): print(1, end=' ')
    def h(self): print(2, end=' ')

class A1(A0):
    def c(self): print(3, end=' ')

try:
    A0().h()
    A1().c()
except:
    print('ERR')
\end{verbatim}
%enditem


\newpage
10. Що буде надруковано за виконання?
\begin{verbatim}
class A:
    c = 1

    def __init__(self):
        b = 2
        b = 3

class B(A):
    b = 4
    a = 5

    def __init__(self):
        super().__init__()
        self.a = 6

try:
    obj = B()
    obj.a = obj.c + 10
    B.a = 13
except: print('EE', end=' ')

try: print(obj.a, end= ' ')
except: print('E1', end=' ')

try: print(A.a, end= ' ')
except: print('E2', end=' ')

try: print(B.a, end= ' ')
except: print('E3', end=' ')
\end{verbatim}
%enditem


\newpage
11. Змінні \kw{next} та \kw{prev} вузла списку позначають наступний та попередній за ним вузол відповідно. Починаючи з голови, у список записано послідовні цілі числа від~$-47$ до~$43$. Нехай змінна \kw{p} позначає вузол, що зберігає значення~$-6$.
Яке значення зберігається у вузлі \kw{p.prev.prev.prev.prev.next} ?
%enditem


\newpage
12. 
\begin{center}
	\adjustimage{max size={0.8\linewidth}{0.8\paperheight}}
	{BinTree/trees_exp/177.png}
\end{center}
Указати послідовність міток вершин дерева, зображеного на рис., що утвориться в результаті обходу дерева в оберненому порядку. Мітки записувати через пробіл.\par Алгоритм починає роботу з кореня (на рис. це верхня вершина).
%enditem


\newpage
13. 
\begin{center}
	\adjustimage{max size={0.9\linewidth}{0.9\paperheight}}
	{Graphs/AdjList/152.png}
\end{center}
Для графа, зображеного на рис., вказати список суміжності вершини 7.\par
Формат оформлення відповіді:\par
список суміжності виписувати через пробіл на одному рядку, якщо пробілів десь буде більше одного - не важливо, упорядкування вершин у списку також значення не має.
%enditem


\newpage
14. Для графа на рисунку вказати послідовність, в якій алгоритм пошуку в глибину буде відмічати відвідані вершини, якщо списки суміжності вершин переглядаються алгоритмом за зростанням міток вершин та вершини відмічаються на першому вход\-женні у вершину.\par
Пошук розпочинається з вершини 4.\par

\begin{center}
	\adjustimage{max size={0.9\linewidth}{0.9\paperheight}}
	{Graphs/Traversals/151.png}
\end{center}
%enditem


\newpage
15. Для графа на рисунку з вершини 6 запущено алгоритм пошуку в ширину, що будує кістякове дерево. Списки суміжності вершин переглядаються алгоритмом за зростанням міток вершин.\par
\begin{center}
	\adjustimage{max size={0.9\linewidth}{0.9\paperheight}}
	{Graphs/Traversals/23.png}
\end{center}
\par Які з наведених ребер входять до складу отриманого кістякового дерева?\par
1) (2, 5)\par
2) (3, 4)\par
3) (1, 2)\par
4) (3, 6)\par
5) (1, 4)\par
6) (4, 5)\par
{\vskip 15pt} Номери відповідних ребер записати за зростанням через пробіл.


%enditem


\newpage
16. 

Рядки журналу через двокрапку містять ім'я користувача (ідентифікатор), час події,тип події, код події, наприклад
\begin{verbatim}
s="username : 2022-05-05 13:26 : ERROR : 404"
\end{verbatim}
у коді 
\begin{verbatim}
res = re.fullmatch('???', s) 
s = ''
code_str = ??? 
\end{verbatim}
чим треба замінити кожен з ???, щоб значенням code\_str був код події? 


\vspace{10mm}
\textbf{Завдання 17} 

Задано програму, що містить код 1-2 класів, можливе успадкування, та клієнтський код.
У наведеному коді невеликі частини закриті. 
Також наведено вихід програми.
Необхідно відновити закриті частини, щоб отриманий код відповідав наведеному виходу програми.


\textbf{Завдання 18-19} 

Кілька запитань за завданнями для самостійної роботи.
\vspace{10mm}
Скоріш за все, що завдань буде трохи менше, а деякі попередні будуть переформульовані в стилі двох останніх.

%enditem

\end {large}}

%end

\newpage
%begin

{{
\begin {large}

\newpage
1. Що буде надруковано за виконання?

\begin{verbatim}
print('R', end=' ')
d=('a','b','c','d','e',0,1,2,3,4)
print(d[-4])
\end{verbatim}

%enditem


\newpage
2. Що буде надруковано за виконання?

\begin{verbatim}
print('R', end=' ')
e=('0','1','2','3','4','5','6','7','8','9')
print(e[-14::-1])
\end{verbatim}

%enditem



\newpage
3. Що буде надруковано за виконання?

\begin{verbatim}
print('R', end=' ')
a=[10,11,12,13,14]
a.extend('11')
print(a)
\end{verbatim}

%enditem



\newpage
4. Що буде надруковано за виконання?

\begin{verbatim}
print('R', end=' ')
f=['0','1','2','3','4']
f[:3]=()
print(f)
\end{verbatim}

%enditem



\newpage
5. Укажіть ТИП значення виразу.\par
\begin{verbatim}
{(1,4), 8}
\end{verbatim}
%enditem


\newpage
6. Що буде надруковано за виконання?\par
\begin{verbatim}
try:
    s = {89:2, 82:4, 86:2, 82:5}
    s[89] = 9
    for x in s.items():
        print(*x, sep=' ', end=' ')
except:
    print('ERR')
\end{verbatim}
%enditem


\newpage
7. Нехай змінна \kw{s} позначає дуже довгу послідовність, по якій можна ітерувати. 

Написати вираз-генератор, що задає підпослідовність її елементів, що є числами типу \kw{float}.

Обмеження: усі елементи шуканої підпослідовності одночасно в пам'яті розміщені бути не можуть. 
%enditem


\newpage
8. Що буде надруковано за виконання?
\begin{verbatim}
a = 0
j = 1

class C:
    a = 2
    
    def __init__(self):
        global a
        self.j = 3
        a = 4
        C.a = 5

obj = C()
try:
    print(a, end=' ')
    print(j, end=' ')
    print(C.a, end=' ')
    print(obj.j)
except:
    print('ERR')
\end{verbatim}
%enditem


\newpage
9. Що буде надруковано за виконання?
\begin{verbatim}
class D0:
    def b(self): print(1, end=' ')
    def g(self): print(2, end=' ')

class D1(D0):
    def b(self): print(3, end=' ')

try:
    D1().g()
    D1().b()
except:
    print('ERR')
\end{verbatim}
%enditem


\newpage
10. Що буде надруковано за виконання?
\begin{verbatim}
class A:
    d = 1

    def __g1(self): print(2, end=' ')
    def _g2(self): print(3, end=' ')

    def main(self):
        try: print(self.d, end=' ')
        except: print('E1', end=' ')

        try: self.__g1()
        except: print('E2', end=' ')

        try: self._g2()
        except: print('E3', end=' ')


class B(A):
    d = 4

    def __g1(self): print(5, end=' ')
    def _g2(self): print(6, end=' ')

try: B().main()
except: print('E4', end=' ')
\end{verbatim}
%enditem


\newpage
11. Змінні \kw{next} та \kw{prev} вузла списку позначають наступний та попередній за ним вузол відповідно. Починаючи з голови, у список записано послідовні цілі числа від~$-60$ до~$27$. Нехай змінна \kw{p} позначає вузол, що зберігає значення~$9$.
Яке значення зберігається у вузлі \kw{p.next.next.next.next.prev} ?
%enditem


\newpage
12. 
\begin{center}
	\adjustimage{max size={0.8\linewidth}{0.8\paperheight}}
	{BinTree/trees_exp/173.png}
\end{center}
Указати послідовність міток вершин дерева, зображеного на рис., що утвориться в результаті обходу дерева в прямому порядку. Мітки записувати через пробіл.\par Алгоритм починає роботу з кореня (на рис. це верхня вершина).
%enditem


\newpage
13. 
\begin{center}
	\adjustimage{max size={0.9\linewidth}{0.9\paperheight}}
	{Graphs/AdjList/30.png}
\end{center}
Для графа, зображеного на рис., вказати список суміжності вершини 1.\par
Формат оформлення відповіді:\par
список суміжності виписувати через пробіл на одному рядку, якщо пробілів десь буде більше одного - не важливо, упорядкування вершин у списку також значення не має.
%enditem


\newpage
14. Для графа на рисунку вказати послідовність, в якій алгоритм пошуку в глибину буде відмічати відвідані вершини, якщо списки суміжності вершин переглядаються алгоритмом за зростанням міток вершин та вершини відмічаються на першому вход\-женні у вершину.\par
Пошук розпочинається з вершини 4.\par

\begin{center}
	\adjustimage{max size={0.9\linewidth}{0.9\paperheight}}
	{Graphs/Traversals/85.png}
\end{center}
%enditem


\newpage
15. Для графа на рисунку з вершини 1 запущено алгоритм пошуку в ширину, що будує кістякове дерево. Списки суміжності вершин переглядаються алгоритмом за зростанням міток вершин.\par
\begin{center}
	\adjustimage{max size={0.9\linewidth}{0.9\paperheight}}
	{Graphs/Traversals/18.png}
\end{center}
\par Які з наведених ребер входять до складу отриманого кістякового дерева?\par
1) (2, 3)\par
2) (3, 4)\par
3) (4, 6)\par
4) (1, 6)\par
5) (2, 4)\par
6) (5, 6)\par
{\vskip 15pt} Номери відповідних ребер записати за зростанням через пробіл.


%enditem


\newpage
16. 

Рядки журналу через двокрапку містять ім'я користувача (ідентифікатор), час події,тип події, код події, наприклад
\begin{verbatim}
s="username : 2022-05-05 13:26 : ERROR : 404"
\end{verbatim}
у коді 
\begin{verbatim}
res = re.fullmatch('???', s) 
s = ''
code_str = ??? 
\end{verbatim}
чим треба замінити кожен з ???, щоб значенням code\_str був код події? 


\vspace{10mm}
\textbf{Завдання 17} 

Задано програму, що містить код 1-2 класів, можливе успадкування, та клієнтський код.
У наведеному коді невеликі частини закриті. 
Також наведено вихід програми.
Необхідно відновити закриті частини, щоб отриманий код відповідав наведеному виходу програми.


\textbf{Завдання 18-19} 

Кілька запитань за завданнями для самостійної роботи.
\vspace{10mm}
Скоріш за все, що завдань буде трохи менше, а деякі попередні будуть переформульовані в стилі двох останніх.

%enditem

\end {large}}

%end

\newpage
%begin

{{
\begin {large}

\newpage
1. Що буде надруковано за виконання?

\begin{verbatim}
print('R', end=' ')
a=('a','b','c','d','e',0,1)
print(a[-11])
\end{verbatim}

%enditem


\newpage
2. Що буде надруковано за виконання?

\begin{verbatim}
print('R', end=' ')
g='0123456789'
print(g[-13:12:3])
\end{verbatim}

%enditem



\newpage
3. Що буде надруковано за виконання?

\begin{verbatim}
print('R', end=' ')
d=['0','1','2','3','4']
print(d.index(4), end=' ')
\end{verbatim}

%enditem



\newpage
4. Що буде надруковано за виконання?

\begin{verbatim}
print('R', end=' ')
c=['0','1','2','3','4']
del c[-3:]
print(c)
\end{verbatim}

%enditem



\newpage
5. Укажіть ТИП значення виразу.\par
\begin{verbatim}
[{1,2}, 5]
\end{verbatim}
%enditem


\newpage
6. Що буде надруковано за виконання?\par
\begin{verbatim}
try:
    d = {56:1, 28:2, 27:4, 76:4, 27:8}
    d[80] = 4
    for x in d.keys():
        print(x, sep=' ', end=' ')
except:
    print('ERR')
\end{verbatim}
%enditem


\newpage
7. Нехай змінна \kw{s} позначає дуже довгу послідовність цілих, по якій можна ітерувати. 

Написати вираз-генератор, що задає підпослідовність її елементів, що є точними квадратами. 

Обмеження: усі елементи шуканої підпослідовності одночасно в пам'яті розміщені бути не можуть. 
%enditem


\newpage
8. Що буде надруковано за виконання?
\begin{verbatim}
d = 0
n = 1

class A:
    n = 2
    
    def __init__(self):
        self.d = 3
        n = 4
        A.n = 5

obj = A()
try:
    print(d, end=' ')
    print(n, end=' ')
    print(A.n, end=' ')
    print(obj.d)
except:
    print('ERR')
\end{verbatim}
%enditem


\newpage
9. Що буде надруковано за виконання?
\begin{verbatim}
class B0:
    def d(self): print(1, end=' ')
    def h(self): print(2, end=' ')

class B1(B0):
    def d(self): print(3, end=' ')

try:
    B1().h()
    B0().d()
except:
    print('ERR')
\end{verbatim}
%enditem


\newpage
10. Що буде надруковано за виконання?
\begin{verbatim}
class A:
    c = 1

    def __init__(self):
        a = 2
        c = 3

class B(A):
    c = 4
    b = 5

    def __init__(self):
        super().__init__()
        self.b = 6

try:
    obj = B()
    obj.b = obj.a + 10
    B.b = 13
except: print('EE', end=' ')

try: print(obj.b, end= ' ')
except: print('E1', end=' ')

try: print(A.b, end= ' ')
except: print('E2', end=' ')

try: print(B.b, end= ' ')
except: print('E3', end=' ')
\end{verbatim}
%enditem


\newpage
11. Змінні \kw{next} та \kw{prev} вузла списку позначають наступний та попередній за ним вузол відповідно. Починаючи з голови, у список записано послідовні цілі числа від~$-70$ до~$42$. Нехай змінна \kw{p} позначає вузол, що зберігає значення~$-7$.
Яке значення зберігається у вузлі \kw{p.prev.prev.next.next.next} ?
%enditem


\newpage
12. 
\begin{center}
	\adjustimage{max size={0.8\linewidth}{0.8\paperheight}}
	{BinTree/trees_exp/74.png}
\end{center}
Указати послідовність міток вершин дерева, зображеного на рис., що утвориться в результаті обходу дерева в оберненому порядку. Мітки записувати через пробіл.\par Алгоритм починає роботу з кореня (на рис. це верхня вершина).
%enditem


\newpage
13. 
\begin{center}
	\adjustimage{max size={0.9\linewidth}{0.9\paperheight}}
	{Graphs/AdjList/164.png}
\end{center}
Для графа, зображеного на рис., вказати список суміжності вершини 3.\par
Формат оформлення відповіді:\par
список суміжності виписувати через пробіл на одному рядку, якщо пробілів десь буде більше одного - не важливо, упорядкування вершин у списку також значення не має.
%enditem


\newpage
14. Для графа на рисунку вказати послідовність, в якій алгоритм пошуку в глибину буде відмічати відвідані вершини, якщо списки суміжності вершин переглядаються алгоритмом за зростанням міток вершин та вершини відмічаються на першому вход\-женні у вершину.\par
Пошук розпочинається з вершини 0.\par

\begin{center}
	\adjustimage{max size={0.9\linewidth}{0.9\paperheight}}
	{Graphs/Traversals/199.png}
\end{center}
%enditem


\newpage
15. Для графа на рисунку з вершини 5 запущено алгоритм пошуку в ширину, що будує кістякове дерево. Списки суміжності вершин переглядаються алгоритмом за зростанням міток вершин.\par
\begin{center}
	\adjustimage{max size={0.9\linewidth}{0.9\paperheight}}
	{Graphs/Traversals/197.png}
\end{center}
\par Які з наведених ребер входять до складу отриманого кістякового дерева?\par
1) (2, 5)\par
2) (3, 4)\par
3) (1, 3)\par
4) (4, 5)\par
5) (1, 5)\par
6) (1, 4)\par
{\vskip 15pt} Номери відповідних ребер записати за зростанням через пробіл.


%enditem


\newpage
16. 

Рядки журналу через двокрапку містять ім'я користувача (ідентифікатор), час події,тип події, код події, наприклад
\begin{verbatim}
s="username : 2022-05-05 13:26 : ERROR : 404"
\end{verbatim}
у коді 
\begin{verbatim}
res = re.fullmatch('???', s) 
s = ''
code_str = ??? 
\end{verbatim}
чим треба замінити кожен з ???, щоб значенням code\_str був код події? 


\vspace{10mm}
\textbf{Завдання 17} 

Задано програму, що містить код 1-2 класів, можливе успадкування, та клієнтський код.
У наведеному коді невеликі частини закриті. 
Також наведено вихід програми.
Необхідно відновити закриті частини, щоб отриманий код відповідав наведеному виходу програми.


\textbf{Завдання 18-19} 

Кілька запитань за завданнями для самостійної роботи.
\vspace{10mm}
Скоріш за все, що завдань буде трохи менше, а деякі попередні будуть переформульовані в стилі двох останніх.

%enditem

\end {large}}

%end

\newpage
%begin

{{
\begin {large}

\newpage
1. Що буде надруковано за виконання?

\begin{verbatim}
print('R', end=' ')
f=('a','b','c','d','e',0,1,2)
print(f[4])
\end{verbatim}

%enditem


\newpage
2. Що буде надруковано за виконання?

\begin{verbatim}
print('R', end=' ')
h=[0,1,2,3,4,5,6,7,8,9]
print(h[-8:13:2])
\end{verbatim}

%enditem



\newpage
3. Що буде надруковано за виконання?

\begin{verbatim}
print('R', end=' ')
h=[10,11,12,13,14]
h+=[1,2]
print(h)
\end{verbatim}

%enditem



\newpage
4. Що буде надруковано за виконання?

\begin{verbatim}
print('R', end=' ')
g=[10,11,12,13,14]
g[0:]=[4,5]
print(g)
\end{verbatim}

%enditem



\newpage
5. Укажіть ТИП значення виразу.\par
\begin{verbatim}
{x for x in range(7)}
\end{verbatim}
%enditem


\newpage
6. Що буде надруковано за виконання?\par
\begin{verbatim}
try:
    t = {81:4, 20:2, 80:0, 80:3}
    t[28] = 9
    for x, y in t.items():
        print(x, y, sep=' ', end=' ')
except:
    print('ERR')
\end{verbatim}
%enditem


\newpage
7. Нехай змінна \kw{s} позначає послідовність цілих, по якій можна ітерувати.

Написати вираз, що перевіряє, чи правда, що всі елементи послідовності менші за задане число num.
%enditem


\newpage
8. Що буде надруковано за виконання?
%Оригинал был утерян. Требует отладки!
\begin{verbatim}
_d = 0

class D:
    _d = 1
    
    def __init__(self):
        self._d = 2
        _d = 3
        D._d = 4

obj = D()
D._d = 5
try:
    print(_d, end=' ')
    print(obj._d, end=' ')
    print(D._d, end=' ')
    print(obj._d, end=' ')
    print(obj._D__d)
except:
    print('ERR')
\end{verbatim}
%enditem


\newpage
9. Що буде надруковано за виконання?
\begin{verbatim}
class C0:
    def a(self): print(1, end=' ')
    def f(self): print(2, end=' ')

class C1(C0):
    def f(self): print(3, end=' ')

try:
    C1().a()
    C1().f()
except:
    print('ERR')
\end{verbatim}
%enditem


\newpage
10. Що буде надруковано за виконання?
\begin{verbatim}
class A:
    b = 1

    def _f1(self): print(2, end=' ')
    def _f2(self): print(3, end=' ')

    def main(self):
        try: print(self.b, end=' ')
        except: print('E1', end=' ')

        try: self._f1()
        except: print('E2', end=' ')

        try: self._f2()
        except: print('E3', end=' ')


class B(A):
    b = 4

    def _f1(self): print(5, end=' ')
    def _f2(self): print(6, end=' ')

try: B().main()
except: print('E4', end=' ')
\end{verbatim}
%enditem


\newpage
11. Змінні \kw{next} та \kw{prev} вузла списку позначають наступний та попередній за ним вузол відповідно. Починаючи з голови, у список записано послідовні цілі числа від~$-17$ до~$40$. Нехай змінна \kw{p} позначає вузол, що зберігає значення~$-2$.
Яке значення зберігається у вузлі \kw{p.next.next.prev.prev.prev} ?
%enditem


\newpage
12. 
\begin{center}
	\adjustimage{max size={0.8\linewidth}{0.8\paperheight}}
	{BinTree/trees_exp/153.png}
\end{center}
Указати послідовність міток вершин дерева, зображеного на рис., що утвориться в результаті обходу дерева в прямому порядку. Мітки записувати через пробіл.\par Алгоритм починає роботу з кореня (на рис. це верхня вершина).
%enditem


\newpage
13. 
\begin{center}
	\adjustimage{max size={0.9\linewidth}{0.9\paperheight}}
	{Graphs/AdjList/175.png}
\end{center}
Для графа, зображеного на рис., вказати список суміжності вершини 1.\par
Формат оформлення відповіді:\par
список суміжності виписувати через пробіл на одному рядку, якщо пробілів десь буде більше одного - не важливо, упорядкування вершин у списку також значення не має.
%enditem


\newpage
14. Для графа на рисунку з вершини 1 запущено алгоритм пошуку в глибину, що будує кістякове дерево. Списки суміжності вершин переглядаються алгоритмом за зростанням міток вершин.\par
\begin{center}
	\adjustimage{max size={0.9\linewidth}{0.9\paperheight}}
	{Graphs/Traversals/229.png}
\end{center}
\par Які з наведених ребер входять до складу отриманого кістякового дерева?\par
1) (4, 5)\par
2) (3, 5)\par
3) (1, 2)\par
4) (5, 6)\par
5) (2, 5)\par
6) (0, 6)\par
{\vskip 15pt} Номери відповідних ребер записати за зростанням через пробіл.

%enditem


\newpage
15. Для графа на рисунку з вершини 2 запущено алгоритм пошуку в ширину, що будує кістякове дерево. Списки суміжності вершин переглядаються алгоритмом за зростанням міток вершин.\par
\begin{center}
	\adjustimage{max size={0.9\linewidth}{0.9\paperheight}}
	{Graphs/Traversals/264.png}
\end{center}
\par Які з наведених ребер входять до складу отриманого кістякового дерева?\par
1) (1, 6)\par
2) (1, 3)\par
3) (4, 5)\par
4) (5, 6)\par
5) (0, 4)\par
6) (0, 5)\par
{\vskip 15pt} Номери відповідних ребер записати за зростанням через пробіл.


%enditem


\newpage
16. 

Рядки журналу через двокрапку містять ім'я користувача (ідентифікатор), час події,тип події, код події, наприклад
\begin{verbatim}
s="username : 2022-05-05 13:26 : ERROR : 404"
\end{verbatim}
у коді 
\begin{verbatim}
res = re.fullmatch('???', s) 
s = ''
code_str = ??? 
\end{verbatim}
чим треба замінити кожен з ???, щоб значенням code\_str був код події? 


\vspace{10mm}
\textbf{Завдання 17} 

Задано програму, що містить код 1-2 класів, можливе успадкування, та клієнтський код.
У наведеному коді невеликі частини закриті. 
Також наведено вихід програми.
Необхідно відновити закриті частини, щоб отриманий код відповідав наведеному виходу програми.


\textbf{Завдання 18-19} 

Кілька запитань за завданнями для самостійної роботи.
\vspace{10mm}
Скоріш за все, що завдань буде трохи менше, а деякі попередні будуть переформульовані в стилі двох останніх.

%enditem

\end {large}}

%end

\newpage
%begin

{{
\begin {large}

\newpage
1. Що буде надруковано за виконання?

\begin{verbatim}
print('R', end=' ')
e=('a','b','c','d','e',0,1,2,3)
print(e[-6])
\end{verbatim}

%enditem


\newpage
2. Що буде надруковано за виконання?

\begin{verbatim}
print('R', end=' ')
d='0123456789'
print(d[15:10:-3])
\end{verbatim}

%enditem



\newpage
3. Що буде надруковано за виконання?

\begin{verbatim}
print('R', end=' ')
f=['0','1','2','3','4']
f.extend('13')
print(f)
\end{verbatim}

%enditem



\newpage
4. Що буде надруковано за виконання?

\begin{verbatim}
print('R', end=' ')
h=[10,11,12,13,14]
h[:7]='13'
print(h)
\end{verbatim}

%enditem



\newpage
5. Укажіть ТИП значення виразу.\par
\begin{verbatim}
{y:13 for y in range(-7,7)}
\end{verbatim}
%enditem


\newpage
6. Що буде надруковано за виконання?\par
\begin{verbatim}
try:
    d = {10:9, 24:6, 37:8, 10:2}
    d[77] = 1
    for x in d.values():
        print(x, sep=' ', end=' ')
except:
    print('ERR')
\end{verbatim}
%enditem


\newpage
7. Нехай змінна \kw{s} позначає послідовність цілих, по якій можна ітерувати. 

Написати вираз, що перевіряє, чи правда, що послідовність  містить числа, більші за 5.
%enditem


\newpage
8. Що буде надруковано за виконання?
\begin{verbatim}
c = 0
m = 1

class B:
    c = 2
    
    def __init__(self):
        self.m = 3
        c = 4
        B.c = 5

obj = B()
try:
    print(c, end=' ')
    print(m, end=' ')
    print(B.c, end=' ')
    print(obj.m)
except:
    print('ERR')
\end{verbatim}
%enditem


\newpage
9. Що буде надруковано за виконання?
\begin{verbatim}
class A0:
    def c(self): print(1, end=' ')
    def g(self): print(2, end=' ')

class A1(A0):
    def c(self): print(3, end=' ')

try:
    A1().g()
    A0().c()
except:
    print('ERR')
\end{verbatim}
%enditem


\newpage
10. Що буде надруковано за виконання?
\begin{verbatim}
class A:
    b = 1

    def __init__(self):
        c = 2
        b = 3

class B(A):
    b = 4
    a = 5

    def __init__(self):
        super().__init__()
        self.a = 6

try:
    obj = B()
    obj.c = obj.b + 10
    B.c = 13
except: print('EE', end=' ')

try: print(obj.c, end= ' ')
except: print('E1', end=' ')

try: print(A.c, end= ' ')
except: print('E2', end=' ')

try: print(B.c, end= ' ')
except: print('E3', end=' ')
\end{verbatim}
%enditem


\newpage
11. Змінні \kw{next} та \kw{prev} вузла списку позначають наступний та попередній за ним вузол відповідно. Починаючи з голови, у список записано послідовні цілі числа від~$-30$ до~$28$. Нехай змінна \kw{p} позначає вузол, що зберігає значення~$7$.
Яке значення зберігається у вузлі \kw{p.next.prev.next.prev.next} ?
%enditem


\newpage
12. 
\begin{center}
	\adjustimage{max size={0.8\linewidth}{0.8\paperheight}}
	{BinTree/trees_exp/130.png}
\end{center}
Указати послідовність міток вершин дерева, зображеного на рис., що утвориться в результаті обходу дерева в оберненому порядку. Мітки записувати через пробіл.\par Алгоритм починає роботу з кореня (на рис. це верхня вершина).
%enditem


\newpage
13. 
\begin{center}
	\adjustimage{max size={0.9\linewidth}{0.9\paperheight}}
	{Graphs/AdjList/1.png}
\end{center}
Для графа, зображеного на рис., вказати список суміжності вершини 1.\par
Формат оформлення відповіді:\par
список суміжності виписувати через пробіл на одному рядку, якщо пробілів десь буде більше одного - не важливо, упорядкування вершин у списку також значення не має.
%enditem


\newpage
14. Для графа на рисунку вказати послідовність, в якій алгоритм пошуку в глибину буде відмічати відвідані вершини, якщо списки суміжності вершин переглядаються алгоритмом за зростанням міток вершин та вершини відмічаються на першому вход\-женні у вершину.\par
Пошук розпочинається з вершини 4.\par

\begin{center}
	\adjustimage{max size={0.9\linewidth}{0.9\paperheight}}
	{Graphs/Traversals/186.png}
\end{center}
%enditem


\newpage
15. Для графа на рисунку з вершини 2 запущено алгоритм пошуку в ширину, що будує кістякове дерево. Списки суміжності вершин переглядаються алгоритмом за зростанням міток вершин.\par
\begin{center}
	\adjustimage{max size={0.9\linewidth}{0.9\paperheight}}
	{Graphs/Traversals/133.png}
\end{center}
\par Які з наведених ребер входять до складу отриманого кістякового дерева?\par
1) (0, 1)\par
2) (1, 2)\par
3) (4, 6)\par
4) (4, 5)\par
5) (5, 6)\par
6) (0, 4)\par
{\vskip 15pt} Номери відповідних ребер записати за зростанням через пробіл.


%enditem


\newpage
16. 

Рядки журналу через двокрапку містять ім'я користувача (ідентифікатор), час події,тип події, код події, наприклад
\begin{verbatim}
s="username : 2022-05-05 13:26 : ERROR : 404"
\end{verbatim}
у коді 
\begin{verbatim}
res = re.fullmatch('???', s) 
s = ''
code_str = ??? 
\end{verbatim}
чим треба замінити кожен з ???, щоб значенням code\_str був код події? 


\vspace{10mm}
\textbf{Завдання 17} 

Задано програму, що містить код 1-2 класів, можливе успадкування, та клієнтський код.
У наведеному коді невеликі частини закриті. 
Також наведено вихід програми.
Необхідно відновити закриті частини, щоб отриманий код відповідав наведеному виходу програми.


\textbf{Завдання 18-19} 

Кілька запитань за завданнями для самостійної роботи.
\vspace{10mm}
Скоріш за все, що завдань буде трохи менше, а деякі попередні будуть переформульовані в стилі двох останніх.

%enditem

\end {large}}

%end

\newpage
%begin

{{
\begin {large}

\newpage
1. Що буде надруковано за виконання?

\begin{verbatim}
print('R', end=' ')
b=('a','b','c','d','e',0,1,2,3,4)
print(b[-9])
\end{verbatim}

%enditem


\newpage
2. Що буде надруковано за виконання?

\begin{verbatim}
print('R', end=' ')
a=(0,1,2,3,4,5,6,7,8,9)
print(a[::-2])
\end{verbatim}

%enditem



\newpage
3. Що буде надруковано за виконання?

\begin{verbatim}
print('R', end=' ')
c=[10,11,12,13,14]
print(c.pop(5), end=' ')
print(c)
\end{verbatim}

%enditem



\newpage
4. Що буде надруковано за виконання?

\begin{verbatim}
print('R', end=' ')
a=['0','1','2','3','4']
del a[7:1]
print(a)
\end{verbatim}

%enditem



\newpage
5. Укажіть ТИП значення виразу.\par
\begin{verbatim}
{y:None for y in range(-8,3)}
\end{verbatim}
%enditem


\newpage
6. Що буде надруковано за виконання?\par
\begin{verbatim}
try:
    s = {23:2, 74:8, 44:6, 29:7, 29:1}
    s[23] = 3
    for x in s.items():
        print(x, sep=' ', end=' ')
except:
    print('ERR')
\end{verbatim}
%enditem


\newpage
7. Нехай змінна \kw{s} позначає послідовність цілих, по якій можна ітерувати. 

Написати вираз, що перевіряє, чи правда, що послідовність  містить число 13.
%enditem


\newpage
8. Що буде надруковано за виконання?
\begin{verbatim}
b = 0
h = 1

class D:
    h = 2
    
    def __init__(self):
        global h
        self.b = 3
        h = 4
        D.h = 5

obj = D()
try:
    print(b, end=' ')
    print(h, end=' ')
    print(D.h, end=' ')
    print(obj.b)
except:
    print('ERR')
\end{verbatim}
%enditem


\newpage
9. Що буде надруковано за виконання?
\begin{verbatim}
class C0:
    def d(self): print(1, end=' ')
    def h(self): print(2, end=' ')

class C1(C0):
    def h(self): print(3, end=' ')

try:
    C0().d()
    C1().h()
except:
    print('ERR')
\end{verbatim}
%enditem


\newpage
10. Що буде надруковано за виконання?
\begin{verbatim}
class A:
    b = 1

    def _g1(self): print(2, end=' ')
    def __g2(self): print(3, end=' ')

    def main(self):
        try: print(self.b, end=' ')
        except: print('E1', end=' ')

        try: self._g1()
        except: print('E2', end=' ')

        try: self.__g2()
        except: print('E3', end=' ')


class B(A):
    b = 4

    def _g1(self): print(5, end=' ')
    def __g2(self): print(6, end=' ')

try: B().main()
except: print('E4', end=' ')
\end{verbatim}
%enditem


\newpage
11. Змінні \kw{next} та \kw{prev} вузла списку позначають наступний та попередній за ним вузол відповідно. Починаючи з голови, у список записано послідовні цілі числа від~$-58$ до~$17$. Нехай змінна \kw{p} позначає вузол, що зберігає значення~$8$.
Яке значення зберігається у вузлі \kw{p.prev.next.prev.next.prev} ?
%enditem


\newpage
12. 
\begin{center}
	\adjustimage{max size={0.8\linewidth}{0.8\paperheight}}
	{BinTree/trees_exp/16.png}
\end{center}
Указати послідовність міток вершин дерева, зображеного на рис., що утвориться в результаті обходу дерева в прямому порядку. Мітки записувати через пробіл.\par Алгоритм починає роботу з кореня (на рис. це верхня вершина).
%enditem


\newpage
13. 
\begin{center}
	\adjustimage{max size={0.9\linewidth}{0.9\paperheight}}
	{Graphs/AdjList/160.png}
\end{center}
Для графа, зображеного на рис., вказати список суміжності вершини 9.\par
Формат оформлення відповіді:\par
список суміжності виписувати через пробіл на одному рядку, якщо пробілів десь буде більше одного - не важливо, упорядкування вершин у списку також значення не має.
%enditem


\newpage
14. Для графа на рисунку вказати послідовність, в якій алгоритм пошуку в глибину буде відмічати відвідані вершини, якщо списки суміжності вершин переглядаються алгоритмом за зростанням міток вершин та вершини відмічаються на першому вход\-женні у вершину.\par
Пошук розпочинається з вершини 3.\par

\begin{center}
	\adjustimage{max size={0.9\linewidth}{0.9\paperheight}}
	{Graphs/Traversals/3.png}
\end{center}
%enditem


\newpage
15. Для графа на рисунку з вершини 0 запущено алгоритм пошуку в ширину, що будує кістякове дерево. Списки суміжності вершин переглядаються алгоритмом за зростанням міток вершин.\par
\begin{center}
	\adjustimage{max size={0.9\linewidth}{0.9\paperheight}}
	{Graphs/Traversals/305.png}
\end{center}
\par Які з наведених ребер входять до складу отриманого кістякового дерева?\par
1) (1, 6)\par
2) (5, 6)\par
3) (4, 5)\par
4) (2, 6)\par
5) (2, 5)\par
6) (0, 4)\par
{\vskip 15pt} Номери відповідних ребер записати за зростанням через пробіл.


%enditem


\newpage
16. 

Рядки журналу через двокрапку містять ім'я користувача (ідентифікатор), час події,тип події, код події, наприклад
\begin{verbatim}
s="username : 2022-05-05 13:26 : ERROR : 404"
\end{verbatim}
у коді 
\begin{verbatim}
res = re.fullmatch('???', s) 
s = ''
code_str = ??? 
\end{verbatim}
чим треба замінити кожен з ???, щоб значенням code\_str був код події? 


\vspace{10mm}
\textbf{Завдання 17} 

Задано програму, що містить код 1-2 класів, можливе успадкування, та клієнтський код.
У наведеному коді невеликі частини закриті. 
Також наведено вихід програми.
Необхідно відновити закриті частини, щоб отриманий код відповідав наведеному виходу програми.


\textbf{Завдання 18-19} 

Кілька запитань за завданнями для самостійної роботи.
\vspace{10mm}
Скоріш за все, що завдань буде трохи менше, а деякі попередні будуть переформульовані в стилі двох останніх.

%enditem

\end {large}}

%end

\newpage
%begin

{{
\begin {large}

\newpage
1. Що буде надруковано за виконання?

\begin{verbatim}
print('R', end=' ')
c=('a','b','c','d','e',0,1)
print(c[13])
\end{verbatim}

%enditem


\newpage
2. Що буде надруковано за виконання?

\begin{verbatim}
print('R', end=' ')
f=['0','1','2','3','4','5','6','7','8','9']
print(f[0:9:-1])
\end{verbatim}

%enditem



\newpage
3. Що буде надруковано за виконання?

\begin{verbatim}
print('R', end=' ')
b=['0','1','2','3','4']
b.append([4,5])
print(b)
\end{verbatim}

%enditem



\newpage
4. Що буде надруковано за виконання?

\begin{verbatim}
print('R', end=' ')
b=[10,11,12,13,14]
b[-2:-1]=[12]
print(b)
\end{verbatim}

%enditem



\newpage
5. Укажіть ТИП значення виразу.\par
\begin{verbatim}
{x:x for x in range(7)}
\end{verbatim}
%enditem


\newpage
6. Що буде надруковано за виконання?\par
\begin{verbatim}
try:
    t = {80:1, 52:9, 38:1, 23:1, 89:4}
    t[38] = 3
    for x in t :
        print(x, sep=' ', end=' ')
except:
    print('ERR')
\end{verbatim}
%enditem


\newpage
7. Нехай змінна \kw{s} позначає послідовність цілих, по якій можна ітерувати.

Написати вираз, що перевіряє, чи правда, що всі елементи послідовності не діляться на 8.
%enditem


\newpage
8. Що буде надруковано за виконання?
\begin{verbatim}
b = 0
h = 1

class D:
    h = 2
    
    def __init__(self):
        self.b = 3
        h = 4
        D.h = 5

obj = D()
try:
    print(b, end=' ')
    print(h, end=' ')
    print(D.h, end=' ')
    print(obj.b)
except:
    print('ERR')
\end{verbatim}
%enditem


\newpage
9. Що буде надруковано за виконання?
\begin{verbatim}
class D0:
    def a(self): print(1, end=' ')
    def f(self): print(2, end=' ')

class D1(D0):
    def f(self): print(3, end=' ')

try:
    D1().a()
    D1().f()
except:
    print('ERR')
\end{verbatim}
%enditem


\newpage
10. Що буде надруковано за виконання?
\begin{verbatim}
class A:
    a = 1

    def _h1(self): print(2, end=' ')
    def h2(self): print(3, end=' ')

    def main(self):
        try: print(self.a, end=' ')
        except: print('E1', end=' ')

        try: self._h1()
        except: print('E2', end=' ')

        try: self.h2()
        except: print('E3', end=' ')


class B(A):
    a = 4

    def _h1(self): print(5, end=' ')
    def h2(self): print(6, end=' ')

try: B().main()
except: print('E4', end=' ')
\end{verbatim}
%enditem


\newpage
11. Змінні \kw{next} та \kw{prev} вузла списку позначають наступний та попередній за ним вузол відповідно. Починаючи з голови, у список записано послідовні цілі числа від~$-24$ до~$34$. Нехай змінна \kw{p} позначає вузол, що зберігає значення~$1$.
Яке значення зберігається у вузлі \kw{p.next.next.prev.next.prev} ?
%enditem


\newpage
12. 
\begin{center}
	\adjustimage{max size={0.8\linewidth}{0.8\paperheight}}
	{BinTree/trees_exp/86.png}
\end{center}
Указати послідовність міток вершин дерева, зображеного на рис., що утвориться в результаті обходу дерева в прямому порядку. Мітки записувати через пробіл.\par Алгоритм починає роботу з кореня (на рис. це верхня вершина).
%enditem


\newpage
13. 
\begin{center}
	\adjustimage{max size={0.9\linewidth}{0.9\paperheight}}
	{Graphs/AdjList/219.png}
\end{center}
Для графа, зображеного на рис., вказати список суміжності вершини 3.\par
Формат оформлення відповіді:\par
список суміжності виписувати через пробіл на одному рядку, якщо пробілів десь буде більше одного - не важливо, упорядкування вершин у списку також значення не має.
%enditem


\newpage
14. Для графа на рисунку вказати послідовність, в якій алгоритм пошуку в глибину буде відмічати відвідані вершини, якщо списки суміжності вершин переглядаються алгоритмом за зростанням міток вершин та вершини відмічаються на першому вход\-женні у вершину.\par
Пошук розпочинається з вершини 2.\par

\begin{center}
	\adjustimage{max size={0.9\linewidth}{0.9\paperheight}}
	{Graphs/Traversals/114.png}
\end{center}
%enditem


\newpage
15. Для графа на рисунку з вершини 2 запущено алгоритм пошуку в ширину, що будує кістякове дерево. Списки суміжності вершин переглядаються алгоритмом за зростанням міток вершин.\par
\begin{center}
	\adjustimage{max size={0.9\linewidth}{0.9\paperheight}}
	{Graphs/Traversals/66.png}
\end{center}
\par Які з наведених ребер входять до складу отриманого кістякового дерева?\par
1) (1, 2)\par
2) (0, 6)\par
3) (2, 5)\par
4) (3, 5)\par
5) (5, 6)\par
6) (1, 5)\par
{\vskip 15pt} Номери відповідних ребер записати за зростанням через пробіл.


%enditem


\newpage
16. 

Рядки журналу через двокрапку містять ім'я користувача (ідентифікатор), час події,тип події, код події, наприклад
\begin{verbatim}
s="username : 2022-05-05 13:26 : ERROR : 404"
\end{verbatim}
у коді 
\begin{verbatim}
res = re.fullmatch('???', s) 
s = ''
code_str = ??? 
\end{verbatim}
чим треба замінити кожен з ???, щоб значенням code\_str був код події? 


\vspace{10mm}
\textbf{Завдання 17} 

Задано програму, що містить код 1-2 класів, можливе успадкування, та клієнтський код.
У наведеному коді невеликі частини закриті. 
Також наведено вихід програми.
Необхідно відновити закриті частини, щоб отриманий код відповідав наведеному виходу програми.


\textbf{Завдання 18-19} 

Кілька запитань за завданнями для самостійної роботи.
\vspace{10mm}
Скоріш за все, що завдань буде трохи менше, а деякі попередні будуть переформульовані в стилі двох останніх.

%enditem

\end {large}}

%end

\newpage
%begin

{{
\begin {large}

\newpage
1. Що буде надруковано за виконання?

\begin{verbatim}
print('R', end=' ')
g=('a','b','c','d','e',0,1,2,4)
print(g[1])
\end{verbatim}

%enditem


\newpage
2. Що буде надруковано за виконання?

\begin{verbatim}
print('R', end=' ')
b=('0','1','2','3','4','5','6','7','8','9')
print(b[-2::3])
\end{verbatim}

%enditem



\newpage
3. Що буде надруковано за виконання?

\begin{verbatim}
print('R', end=' ')
g=[10,11,12,13,14]
g.remove(13)
print(g)
\end{verbatim}

%enditem



\newpage
4. Що буде надруковано за виконання?

\begin{verbatim}
print('R', end=' ')
d=['0','1','2','3','4']
d[-4:-9]=(13)
print(d)
\end{verbatim}

%enditem



\newpage
5. Укажіть ТИП значення виразу.\par
\begin{verbatim}
{21,2,4}
\end{verbatim}
%enditem


\newpage
6. Що буде надруковано за виконання?\par
\begin{verbatim}
try:
    t = {71:4, 57:1, 87:2, 87:8}
    t[81] = 9
    for x in t.items():
        print(*x, sep=' ', end=' ')
except:
    print('ERR')
\end{verbatim}
%enditem


\newpage
7. Нехай змінна \kw{s} позначає дуже довгу послідовність, по якій можна ітерувати. 

Написати вираз-генератор, що задає підпослідовність її елементів, що є числами типу \kw{int}.

Обмеження: усі елементи шуканої підпослідовності одночасно в пам'яті розміщені бути не можуть. 
%enditem


\newpage
8. Що буде надруковано за виконання?
%Оригинал был утерян. Требует отладки!
\begin{verbatim}
_b = 0

class C:
    _b = 1
    
    def __init__(self):
        self._b = 2
        _b = 3
        C._b = 4

obj = C()
C._b = 5
try:
    print(_b, end=' ')
    print(obj._b, end=' ')
    print(C._b, end=' ')
    print(obj._b, end=' ')
    print(obj._C__b)
except:
    print('ERR')
\end{verbatim}
%enditem


\newpage
9. Що буде надруковано за виконання?
\begin{verbatim}
class B0:
    def b(self): print(1, end=' ')
    def g(self): print(2, end=' ')

class B1(B0):
    def g(self): print(3, end=' ')

try:
    B1().b()
    B1().g()
except:
    print('ERR')
\end{verbatim}
%enditem


\newpage
10. Що буде надруковано за виконання?
\begin{verbatim}
class A:
    a = 1

    def __init__(self):
        b = 2
        b = 3

class B(A):
    b = 4
    c = 5

    def __init__(self):
        super().__init__()
        self.c = 6

try:
    obj = B()
    obj.a = obj.c + 10
    B.a = 13
except: print('EE', end=' ')

try: print(obj.a, end= ' ')
except: print('E1', end=' ')

try: print(A.a, end= ' ')
except: print('E2', end=' ')

try: print(B.a, end= ' ')
except: print('E3', end=' ')
\end{verbatim}
%enditem


\newpage
11. Змінні \kw{next} та \kw{prev} вузла списку позначають наступний та попередній за ним вузол відповідно. Починаючи з голови, у список записано послідовні цілі числа від~$-18$ до~$19$. Нехай змінна \kw{p} позначає вузол, що зберігає значення~$4$.
Яке значення зберігається у вузлі \kw{p.prev.prev.next.prev.next} ?
%enditem


\newpage
12. 
\begin{center}
	\adjustimage{max size={0.8\linewidth}{0.8\paperheight}}
	{BinTree/trees_exp/128.png}
\end{center}
Указати послідовність міток вершин дерева, зображеного на рис., що утвориться в результаті обходу дерева в оберненому порядку. Мітки записувати через пробіл.\par Алгоритм починає роботу з кореня (на рис. це верхня вершина).
%enditem


\newpage
13. 
\begin{center}
	\adjustimage{max size={0.9\linewidth}{0.9\paperheight}}
	{Graphs/AdjList/61.png}
\end{center}
Для графа, зображеного на рис., вказати список суміжності вершини 8.\par
Формат оформлення відповіді:\par
список суміжності виписувати через пробіл на одному рядку, якщо пробілів десь буде більше одного - не важливо, упорядкування вершин у списку також значення не має.
%enditem


\newpage
14. Для графа на рисунку з вершини 4 запущено алгоритм пошуку в глибину, що будує кістякове дерево. Списки суміжності вершин переглядаються алгоритмом за зростанням міток вершин.\par
\begin{center}
	\adjustimage{max size={0.9\linewidth}{0.9\paperheight}}
	{Graphs/Traversals/198.png}
\end{center}
\par Які з наведених ребер входять до складу отриманого кістякового дерева?\par
1) (0, 1)\par
2) (3, 6)\par
3) (3, 4)\par
4) (1, 3)\par
5) (0, 6)\par
6) (5, 6)\par
{\vskip 15pt} Номери відповідних ребер записати за зростанням через пробіл.

%enditem


\newpage
15. Для графа на рисунку з вершини 2 запущено алгоритм пошуку в ширину, що будує кістякове дерево. Списки суміжності вершин переглядаються алгоритмом за зростанням міток вершин.\par
\begin{center}
	\adjustimage{max size={0.9\linewidth}{0.9\paperheight}}
	{Graphs/Traversals/349.png}
\end{center}
\par Які з наведених ребер входять до складу отриманого кістякового дерева?\par
1) (2, 6)\par
2) (3, 4)\par
3) (1, 2)\par
4) (4, 6)\par
5) (1, 4)\par
6) (3, 5)\par
{\vskip 15pt} Номери відповідних ребер записати за зростанням через пробіл.


%enditem


\newpage
16. 

Рядки журналу через двокрапку містять ім'я користувача (ідентифікатор), час події,тип події, код події, наприклад
\begin{verbatim}
s="username : 2022-05-05 13:26 : ERROR : 404"
\end{verbatim}
у коді 
\begin{verbatim}
res = re.fullmatch('???', s) 
s = ''
code_str = ??? 
\end{verbatim}
чим треба замінити кожен з ???, щоб значенням code\_str був код події? 


\vspace{10mm}
\textbf{Завдання 17} 

Задано програму, що містить код 1-2 класів, можливе успадкування, та клієнтський код.
У наведеному коді невеликі частини закриті. 
Також наведено вихід програми.
Необхідно відновити закриті частини, щоб отриманий код відповідав наведеному виходу програми.


\textbf{Завдання 18-19} 

Кілька запитань за завданнями для самостійної роботи.
\vspace{10mm}
Скоріш за все, що завдань буде трохи менше, а деякі попередні будуть переформульовані в стилі двох останніх.

%enditem

\end {large}}

%end

\newpage
%begin

{{
\begin {large}

\newpage
1. Що буде надруковано за виконання?

\begin{verbatim}
print('R', end=' ')
b=('a','b','c','d','e',0,1)
print(b[-7])
\end{verbatim}

%enditem


\newpage
2. Що буде надруковано за виконання?

\begin{verbatim}
print('R', end=' ')
d='0123456789'
print(d[6::3])
\end{verbatim}

%enditem



\newpage
3. Що буде надруковано за виконання?

\begin{verbatim}
print('R', end=' ')
e=[10,11,12,13,14]
e.insert(7, '12')
print(e)
\end{verbatim}

%enditem



\newpage
4. Що буде надруковано за виконання?

\begin{verbatim}
print('R', end=' ')
c=[10,11,12,13,14]
c[3:2]='12'
print(c)
\end{verbatim}

%enditem



\newpage
5. Укажіть ТИП значення виразу.\par
\begin{verbatim}
[13,2,4]+[3,5]
\end{verbatim}
%enditem


\newpage
6. Що буде надруковано за виконання?\par
\begin{verbatim}
try:
    d = {27:1, 26:4, 63:7, 50:2, 44:9}
    d[50] = 6
    for x, y in d.items():
        print(x, y, sep=' ', end=' ')
except:
    print('ERR')
\end{verbatim}
%enditem


\newpage
7. Записати ВИРАЗ, значенням якого є множина, що складається з цілих чисел від 14 до 2001 (включно), що не діляться на 5.
%enditem


\newpage
8. Що буде надруковано за виконання?
\begin{verbatim}
c = 0
k = 1

class B:
    c = 2
    
    def __init__(self):
        global c
        self.k = 3
        c = 4
        B.c = 5

obj = B()
try:
    print(c, end=' ')
    print(k, end=' ')
    print(B.c, end=' ')
    print(obj.k)
except:
    print('ERR')
\end{verbatim}
%enditem


\newpage
9. Що буде надруковано за виконання?
\begin{verbatim}
class A0:
    def a(self): print(1, end=' ')
    def h(self): print(2, end=' ')

class A1(A0):
    def h(self): print(3, end=' ')

try:
    A0().a()
    A1().h()
except:
    print('ERR')
\end{verbatim}
%enditem


\newpage
10. Що буде надруковано за виконання?
\begin{verbatim}
class A:
    d = 1

    def f1(self): print(2, end=' ')
    def _f2(self): print(3, end=' ')

    def main(self):
        try: print(self.d, end=' ')
        except: print('E1', end=' ')

        try: self.f1()
        except: print('E2', end=' ')

        try: self._f2()
        except: print('E3', end=' ')


class B(A):
    d = 4

    def f1(self): print(5, end=' ')
    def _f2(self): print(6, end=' ')

try: B().main()
except: print('E4', end=' ')
\end{verbatim}
%enditem


\newpage
11. Змінні \kw{next} та \kw{prev} вузла списку позначають наступний та попередній за ним вузол відповідно. Починаючи з голови, у список записано послідовні цілі числа від~$-45$ до~$29$. Нехай змінна \kw{p} позначає вузол, що зберігає значення~$2$.
Яке значення зберігається у вузлі \kw{p.next.prev.prev.next.next} ?
%enditem


\newpage
12. 
\begin{center}
	\adjustimage{max size={0.8\linewidth}{0.8\paperheight}}
	{BinTree/trees_exp/22.png}
\end{center}
Указати послідовність міток вершин дерева, зображеного на рис., що утвориться в результаті обходу дерева в прямому порядку. Мітки записувати через пробіл.\par Алгоритм починає роботу з кореня (на рис. це верхня вершина).
%enditem


\newpage
13. 
\begin{center}
	\adjustimage{max size={0.9\linewidth}{0.9\paperheight}}
	{Graphs/AdjList/248.png}
\end{center}
Для графа, зображеного на рис., вказати список суміжності вершини 1.\par
Формат оформлення відповіді:\par
список суміжності виписувати через пробіл на одному рядку, якщо пробілів десь буде більше одного - не важливо, упорядкування вершин у списку також значення не має.
%enditem


\newpage
14. Для графа на рисунку з вершини 5 запущено алгоритм пошуку в глибину, що будує кістякове дерево. Списки суміжності вершин переглядаються алгоритмом за зростанням міток вершин.\par
\begin{center}
	\adjustimage{max size={0.9\linewidth}{0.9\paperheight}}
	{Graphs/Traversals/153.png}
\end{center}
\par Які з наведених ребер входять до складу отриманого кістякового дерева?\par
1) (1, 2)\par
2) (2, 5)\par
3) (0, 1)\par
4) (5, 6)\par
5) (0, 2)\par
6) (4, 5)\par
{\vskip 15pt} Номери відповідних ребер записати за зростанням через пробіл.

%enditem


\newpage
15. Для графа на рисунку з вершини 1 запущено алгоритм пошуку в ширину, що будує кістякове дерево. Списки суміжності вершин переглядаються алгоритмом за зростанням міток вершин.\par
\begin{center}
	\adjustimage{max size={0.9\linewidth}{0.9\paperheight}}
	{Graphs/Traversals/70.png}
\end{center}
\par Які з наведених ребер входять до складу отриманого кістякового дерева?\par
1) (1, 3)\par
2) (2, 6)\par
3) (0, 2)\par
4) (5, 6)\par
5) (0, 6)\par
6) (1, 2)\par
{\vskip 15pt} Номери відповідних ребер записати за зростанням через пробіл.


%enditem


\newpage
16. 

Рядки журналу через двокрапку містять ім'я користувача (ідентифікатор), час події,тип події, код події, наприклад
\begin{verbatim}
s="username : 2022-05-05 13:26 : ERROR : 404"
\end{verbatim}
у коді 
\begin{verbatim}
res = re.fullmatch('???', s) 
s = ''
code_str = ??? 
\end{verbatim}
чим треба замінити кожен з ???, щоб значенням code\_str був код події? 


\vspace{10mm}
\textbf{Завдання 17} 

Задано програму, що містить код 1-2 класів, можливе успадкування, та клієнтський код.
У наведеному коді невеликі частини закриті. 
Також наведено вихід програми.
Необхідно відновити закриті частини, щоб отриманий код відповідав наведеному виходу програми.


\textbf{Завдання 18-19} 

Кілька запитань за завданнями для самостійної роботи.
\vspace{10mm}
Скоріш за все, що завдань буде трохи менше, а деякі попередні будуть переформульовані в стилі двох останніх.

%enditem

\end {large}}

%end

\newpage
%begin

{{
\begin {large}

\newpage
1. Що буде надруковано за виконання?

\begin{verbatim}
print('R', end=' ')
f=('a','b','c','d','e',0,1,2,4)
print(f[3])
\end{verbatim}

%enditem


\newpage
2. Що буде надруковано за виконання?

\begin{verbatim}
print('R', end=' ')
g=['0','1','2','3','4','5','6','7','8','9']
print(g[10:7:-3])
\end{verbatim}

%enditem



\newpage
3. Що буде надруковано за виконання?

\begin{verbatim}
print('R', end=' ')
a=['0','1','2','3','4']
a.insert(0, '13')
print(a)
\end{verbatim}

%enditem



\newpage
4. Що буде надруковано за виконання?

\begin{verbatim}
print('R', end=' ')
d=['0','1','2','3','4']
del d[-5:-8]
print(d)
\end{verbatim}

%enditem



\newpage
5. Укажіть ТИП значення виразу.\par
\begin{verbatim}
(2, 7, -4)
\end{verbatim}
%enditem


\newpage
6. Що буде надруковано за виконання?\par
\begin{verbatim}
try:
    s = {15:5, 53:9, 16:0, 16:0}
    s[80] = 6
    for x in s.keys():
        print(x, sep=' ', end=' ')
except:
    print('ERR')
\end{verbatim}
%enditem


\newpage
7. Нехай змінна \kw{s} позначає дуже довгу послідовність, по якій можна ітерувати. 

Написати вираз-генератор, що задає підпослідовність її елементів, що є рядками типу \kw{str}.

Обмеження: усі елементи шуканої підпослідовності одночасно в пам'яті розміщені бути не можуть. 
%enditem


\newpage
8. Що буде надруковано за виконання?
\begin{verbatim}
d = 0
j = 1

class A:
    j = 2
    
    def __init__(self):
        self.j = 3
        j = 4
        A.j = 5

obj = A()
try:
    print(d, end=' ')
    print(j, end=' ')
    print(A.j, end=' ')
    print(obj.j)
except:
    print('ERR')
\end{verbatim}
%enditem


\newpage
9. Що буде надруковано за виконання?
\begin{verbatim}
class C0:
    def b(self): print(1, end=' ')
    def f(self): print(2, end=' ')

class C1(C0):
    def f(self): print(3, end=' ')

try:
    C1().b()
    C1().f()
except:
    print('ERR')
\end{verbatim}
%enditem


\newpage
10. Що буде надруковано за виконання?
\begin{verbatim}
class A:
    c = 1

    def __g1(self): print(2, end=' ')
    def _g2(self): print(3, end=' ')

    def main(self):
        try: print(self.c, end=' ')
        except: print('E1', end=' ')

        try: self.__g1()
        except: print('E2', end=' ')

        try: self._g2()
        except: print('E3', end=' ')


class B(A):
    c = 4

    def __g1(self): print(5, end=' ')
    def _g2(self): print(6, end=' ')

try: B().main()
except: print('E4', end=' ')
\end{verbatim}
%enditem


\newpage
11. Змінні \kw{next} та \kw{prev} вузла списку позначають наступний та попередній за ним вузол відповідно. Починаючи з голови, у список записано послідовні цілі числа від~$-23$ до~$21$. Нехай змінна \kw{p} позначає вузол, що зберігає значення~$-9$.
Яке значення зберігається у вузлі \kw{p.prev.next.next.prev.prev} ?
%enditem


\newpage
12. 
\begin{center}
	\adjustimage{max size={0.8\linewidth}{0.8\paperheight}}
	{BinTree/trees_exp/62.png}
\end{center}
Указати послідовність міток вершин дерева, зображеного на рис., що утвориться в результаті обходу дерева в оберненому порядку. Мітки записувати через пробіл.\par Алгоритм починає роботу з кореня (на рис. це верхня вершина).
%enditem


\newpage
13. 
\begin{center}
	\adjustimage{max size={0.9\linewidth}{0.9\paperheight}}
	{Graphs/AdjList/69.png}
\end{center}
Для графа, зображеного на рис., вказати список суміжності вершини 3.\par
Формат оформлення відповіді:\par
список суміжності виписувати через пробіл на одному рядку, якщо пробілів десь буде більше одного - не важливо, упорядкування вершин у списку також значення не має.
%enditem


\newpage
14. Для графа на рисунку вказати послідовність, в якій алгоритм пошуку в глибину буде відмічати відвідані вершини, якщо списки суміжності вершин переглядаються алгоритмом за зростанням міток вершин та вершини відмічаються на першому вход\-женні у вершину.\par
Пошук розпочинається з вершини 3.\par

\begin{center}
	\adjustimage{max size={0.9\linewidth}{0.9\paperheight}}
	{Graphs/Traversals/280.png}
\end{center}
%enditem


\newpage
15. Для графа на рисунку з вершини 3 запущено алгоритм пошуку в ширину, що будує кістякове дерево. Списки суміжності вершин переглядаються алгоритмом за зростанням міток вершин.\par
\begin{center}
	\adjustimage{max size={0.9\linewidth}{0.9\paperheight}}
	{Graphs/Traversals/280.png}
\end{center}
\par Які з наведених ребер входять до складу отриманого кістякового дерева?\par
1) (2, 5)\par
2) (3, 4)\par
3) (4, 5)\par
4) (2, 3)\par
5) (3, 5)\par
6) (0, 6)\par
{\vskip 15pt} Номери відповідних ребер записати за зростанням через пробіл.


%enditem


\newpage
16. 

Рядки журналу через двокрапку містять ім'я користувача (ідентифікатор), час події,тип події, код події, наприклад
\begin{verbatim}
s="username : 2022-05-05 13:26 : ERROR : 404"
\end{verbatim}
у коді 
\begin{verbatim}
res = re.fullmatch('???', s) 
s = ''
code_str = ??? 
\end{verbatim}
чим треба замінити кожен з ???, щоб значенням code\_str був код події? 


\vspace{10mm}
\textbf{Завдання 17} 

Задано програму, що містить код 1-2 класів, можливе успадкування, та клієнтський код.
У наведеному коді невеликі частини закриті. 
Також наведено вихід програми.
Необхідно відновити закриті частини, щоб отриманий код відповідав наведеному виходу програми.


\textbf{Завдання 18-19} 

Кілька запитань за завданнями для самостійної роботи.
\vspace{10mm}
Скоріш за все, що завдань буде трохи менше, а деякі попередні будуть переформульовані в стилі двох останніх.

%enditem

\end {large}}

%end

\newpage
%begin

{{
\begin {large}

\newpage
1. Що буде надруковано за виконання?

\begin{verbatim}
print('R', end=' ')
e=('a','b','c','d','e',0,1,2,3,4)
print(e[-12])
\end{verbatim}

%enditem


\newpage
2. Що буде надруковано за виконання?

\begin{verbatim}
print('R', end=' ')
c=(0,1,2,3,4,5,6,7,8,9)
print(c[:-6:2])
\end{verbatim}

%enditem



\newpage
3. Що буде надруковано за виконання?

\begin{verbatim}
print('R', end=' ')
b=['0','1','2','3','4']
b.extend([2,3])
print(b)
\end{verbatim}

%enditem



\newpage
4. Що буде надруковано за виконання?

\begin{verbatim}
print('R', end=' ')
f=['0','1','2','3','4']
f[:5]='10'
print(f)
\end{verbatim}

%enditem



\newpage
5. Укажіть ТИП значення виразу.\par
\begin{verbatim}
(11,2,3)
\end{verbatim}
%enditem


\newpage
6. Що буде надруковано за виконання?\par
\begin{verbatim}
try:
    d = {51:6, 18:2, 56:1, 56:6}
    d[36] = 3
    for x, y in d.items():
        print(x, y, sep=' ', end=' ')
except:
    print('ERR')
\end{verbatim}
%enditem


\newpage
7. Записати ВИРАЗ, значенням якого є множина, що складається з квадратних корен\-ів цілих чисел від 13 до 2007 (включно), що діляться на 3.
%enditem


\newpage
8. Що буде надруковано за виконання?
%Оригинал был утерян. Требует отладки!
\begin{verbatim}
a = 0

class A:
    a = 1
    
    def __init__(self):
        self.a = 2
        a = 3
        A.a = 4

obj = A()
obj.a = 5
try:
    print(a, end=' ')
    print(obj.a, end=' ')
    print(A.a, end=' ')
    print(obj.a, end=' ')
    print(obj._A__a)
except:
    print('ERR')
\end{verbatim}
%enditem


\newpage
9. Що буде надруковано за виконання?
\begin{verbatim}
class B0:
    def c(self): print(1, end=' ')
    def f(self): print(2, end=' ')

class B1(B0):
    def f(self): print(3, end=' ')

try:
    B1().c()
    B0().f()
except:
    print('ERR')
\end{verbatim}
%enditem


\newpage
10. Що буде надруковано за виконання?
\begin{verbatim}
class A:
    c = 1

    def _h1(self): print(2, end=' ')
    def __h2(self): print(3, end=' ')

    def main(self):
        try: print(self.c, end=' ')
        except: print('E1', end=' ')

        try: self._h1()
        except: print('E2', end=' ')

        try: self.__h2()
        except: print('E3', end=' ')


class B(A):
    c = 4

    def _h1(self): print(5, end=' ')
    def __h2(self): print(6, end=' ')

try: B().main()
except: print('E4', end=' ')
\end{verbatim}
%enditem


\newpage
11. Змінні \kw{next} та \kw{prev} вузла списку позначають наступний та попередній за ним вузол відповідно. Починаючи з голови, у список записано послідовні цілі числа від~$-32$ до~$22$. Нехай змінна \kw{p} позначає вузол, що зберігає значення~$-8$.
Яке значення зберігається у вузлі \kw{p.prev.next.prev.next.prev} ?
%enditem


\newpage
12. 
\begin{center}
	\adjustimage{max size={0.8\linewidth}{0.8\paperheight}}
	{BinTree/trees_exp/168.png}
\end{center}
Указати послідовність міток вершин дерева, зображеного на рис., що утвориться в результаті обходу дерева в оберненому порядку. Мітки записувати через пробіл.\par Алгоритм починає роботу з кореня (на рис. це верхня вершина).
%enditem


\newpage
13. 
\begin{center}
	\adjustimage{max size={0.9\linewidth}{0.9\paperheight}}
	{Graphs/AdjList/309.png}
\end{center}
Для графа, зображеного на рис., вказати список суміжності вершини 7.\par
Формат оформлення відповіді:\par
список суміжності виписувати через пробіл на одному рядку, якщо пробілів десь буде більше одного - не важливо, упорядкування вершин у списку також значення не має.
%enditem


\newpage
14. Для графа на рисунку вказати послідовність, в якій алгоритм пошуку в глибину буде відмічати відвідані вершини, якщо списки суміжності вершин переглядаються алгоритмом за зростанням міток вершин та вершини відмічаються на першому вход\-женні у вершину.\par
Пошук розпочинається з вершини 1.\par

\begin{center}
	\adjustimage{max size={0.9\linewidth}{0.9\paperheight}}
	{Graphs/Traversals/112.png}
\end{center}
%enditem


\newpage
15. Для графа на рисунку з вершини 4 запущено алгоритм пошуку в ширину, що будує кістякове дерево. Списки суміжності вершин переглядаються алгоритмом за зростанням міток вершин.\par
\begin{center}
	\adjustimage{max size={0.9\linewidth}{0.9\paperheight}}
	{Graphs/Traversals/6.png}
\end{center}
\par Які з наведених ребер входять до складу отриманого кістякового дерева?\par
1) (1, 2)\par
2) (1, 4)\par
3) (1, 3)\par
4) (2, 6)\par
5) (0, 1)\par
6) (0, 6)\par
{\vskip 15pt} Номери відповідних ребер записати за зростанням через пробіл.


%enditem


\newpage
16. 

Рядки журналу через двокрапку містять ім'я користувача (ідентифікатор), час події,тип події, код події, наприклад
\begin{verbatim}
s="username : 2022-05-05 13:26 : ERROR : 404"
\end{verbatim}
у коді 
\begin{verbatim}
res = re.fullmatch('???', s) 
s = ''
code_str = ??? 
\end{verbatim}
чим треба замінити кожен з ???, щоб значенням code\_str був код події? 


\vspace{10mm}
\textbf{Завдання 17} 

Задано програму, що містить код 1-2 класів, можливе успадкування, та клієнтський код.
У наведеному коді невеликі частини закриті. 
Також наведено вихід програми.
Необхідно відновити закриті частини, щоб отриманий код відповідав наведеному виходу програми.


\textbf{Завдання 18-19} 

Кілька запитань за завданнями для самостійної роботи.
\vspace{10mm}
Скоріш за все, що завдань буде трохи менше, а деякі попередні будуть переформульовані в стилі двох останніх.

%enditem

\end {large}}

%end

\newpage
%begin

{{
\begin {large}

\newpage
1. Що буде надруковано за виконання?

\begin{verbatim}
print('R', end=' ')
h=('a','b','c','d','e',0,1,2,3)
print(h[-2])
\end{verbatim}

%enditem


\newpage
2. Що буде надруковано за виконання?

\begin{verbatim}
print('R', end=' ')
h=[0,1,2,3,4,5,6,7,8,9]
print(h[-6:11:-2])
\end{verbatim}

%enditem



\newpage
3. Що буде надруковано за виконання?

\begin{verbatim}
print('R', end=' ')
f=[10,11,12,13,14]
print(f.index(15), end=' ')
\end{verbatim}

%enditem



\newpage
4. Що буде надруковано за виконання?

\begin{verbatim}
print('R', end=' ')
g=[10,11,12,13,14]
g[:-7]='11'
print(g)
\end{verbatim}

%enditem



\newpage
5. Укажіть ТИП значення виразу.\par
\begin{verbatim}
[x for x in range(7)]
\end{verbatim}
%enditem


\newpage
6. Що буде надруковано за виконання?\par
\begin{verbatim}
try:
    t = {38:7, 63:0, 45:0, 82:0, 45:0}
    t[18] = 0
    for x in t.items():
        print(x, sep=' ', end=' ')
except:
    print('ERR')
\end{verbatim}
%enditem


\newpage
7. Нехай змінна \kw{s} позначає дуже довгу послідовність цілих, по якій можна ітерувати. 

Написати вираз-генератор, що задає підпослідовність її елементів, що є додатни\-ми. 

Обмеження: усі елементи шуканої підпослідовності одночасно в пам'яті розміщені бути не можуть. 
%enditem


\newpage
8. Що буде надруковано за виконання?
%Оригинал был утерян. Требует отладки!
\begin{verbatim}
__a = 0

class C:
    __a = 1
    
    def __init__(self):
        self.__a = 2
        __a = 3
        C.__a = 4

obj = C()
C.__a = 5
try:
    print(__a, end=' ')
    print(obj.__a, end=' ')
    print(C.__a, end=' ')
    print(obj.__a, end=' ')
    print(obj._C__a)
except:
    print('ERR')
\end{verbatim}
%enditem


\newpage
9. Що буде надруковано за виконання?
\begin{verbatim}
class D0:
    def d(self): print(1, end=' ')
    def g(self): print(2, end=' ')

class D1(D0):
    def g(self): print(3, end=' ')

try:
    D1().d()
    D1().g()
except:
    print('ERR')
\end{verbatim}
%enditem


\newpage
10. Що буде надруковано за виконання?
\begin{verbatim}
class A:
    b = 1

    def __init__(self):
        a = 2
        b = 3

class B(A):
    a = 4
    c = 5

    def __init__(self):
        super().__init__()
        self.c = 6

try:
    obj = B()
    obj.b = obj.a + 10
    B.b = 13
except: print('EE', end=' ')

try: print(obj.b, end= ' ')
except: print('E1', end=' ')

try: print(A.b, end= ' ')
except: print('E2', end=' ')

try: print(B.b, end= ' ')
except: print('E3', end=' ')
\end{verbatim}
%enditem


\newpage
11. Змінні \kw{next} та \kw{prev} вузла списку позначають наступний та попередній за ним вузол відповідно. Починаючи з голови, у список записано послідовні цілі числа від~$-31$ до~$49$. Нехай змінна \kw{p} позначає вузол, що зберігає значення~$-3$.
Яке значення зберігається у вузлі \kw{p.next.prev.next.prev.next} ?
%enditem


\newpage
12. 
\begin{center}
	\adjustimage{max size={0.8\linewidth}{0.8\paperheight}}
	{BinTree/trees_exp/55.png}
\end{center}
Указати послідовність міток вершин дерева, зображеного на рис., що утвориться в результаті обходу дерева в прямому порядку. Мітки записувати через пробіл.\par Алгоритм починає роботу з кореня (на рис. це верхня вершина).
%enditem


\newpage
13. 
\begin{center}
	\adjustimage{max size={0.9\linewidth}{0.9\paperheight}}
	{Graphs/AdjList/260.png}
\end{center}
Для графа, зображеного на рис., вказати список суміжності вершини 8.\par
Формат оформлення відповіді:\par
список суміжності виписувати через пробіл на одному рядку, якщо пробілів десь буде більше одного - не важливо, упорядкування вершин у списку також значення не має.
%enditem


\newpage
14. Для графа на рисунку з вершини 0 запущено алгоритм пошуку в глибину, що будує кістякове дерево. Списки суміжності вершин переглядаються алгоритмом за зростанням міток вершин.\par
\begin{center}
	\adjustimage{max size={0.9\linewidth}{0.9\paperheight}}
	{Graphs/Traversals/110.png}
\end{center}
\par Які з наведених ребер входять до складу отриманого кістякового дерева?\par
1) (2, 6)\par
2) (3, 6)\par
3) (4, 5)\par
4) (1, 5)\par
5) (2, 5)\par
6) (5, 6)\par
{\vskip 15pt} Номери відповідних ребер записати за зростанням через пробіл.

%enditem


\newpage
15. Для графа на рисунку з вершини 1 запущено алгоритм пошуку в ширину, що будує кістякове дерево. Списки суміжності вершин переглядаються алгоритмом за зростанням міток вершин.\par
\begin{center}
	\adjustimage{max size={0.9\linewidth}{0.9\paperheight}}
	{Graphs/Traversals/109.png}
\end{center}
\par Які з наведених ребер входять до складу отриманого кістякового дерева?\par
1) (0, 5)\par
2) (3, 6)\par
3) (0, 6)\par
4) (0, 3)\par
5) (0, 4)\par
6) (0, 2)\par
{\vskip 15pt} Номери відповідних ребер записати за зростанням через пробіл.


%enditem


\newpage
16. 

Рядки журналу через двокрапку містять ім'я користувача (ідентифікатор), час події,тип події, код події, наприклад
\begin{verbatim}
s="username : 2022-05-05 13:26 : ERROR : 404"
\end{verbatim}
у коді 
\begin{verbatim}
res = re.fullmatch('???', s) 
s = ''
code_str = ??? 
\end{verbatim}
чим треба замінити кожен з ???, щоб значенням code\_str був код події? 


\vspace{10mm}
\textbf{Завдання 17} 

Задано програму, що містить код 1-2 класів, можливе успадкування, та клієнтський код.
У наведеному коді невеликі частини закриті. 
Також наведено вихід програми.
Необхідно відновити закриті частини, щоб отриманий код відповідав наведеному виходу програми.


\textbf{Завдання 18-19} 

Кілька запитань за завданнями для самостійної роботи.
\vspace{10mm}
Скоріш за все, що завдань буде трохи менше, а деякі попередні будуть переформульовані в стилі двох останніх.

%enditem

\end {large}}

%end

\newpage
%begin

{{
\begin {large}

\newpage
1. Що буде надруковано за виконання?

\begin{verbatim}
print('R', end=' ')
g=('a','b','c','d','e',0,1,2)
print(g[0])
\end{verbatim}

%enditem


\newpage
2. Що буде надруковано за виконання?

\begin{verbatim}
print('R', end=' ')
b=('0','1','2','3','4','5','6','7','8','9')
print(b[-11:-7:-1])
\end{verbatim}

%enditem



\newpage
3. Що буде надруковано за виконання?

\begin{verbatim}
print('R', end=' ')
g=[10,11,12,13,14]
print(g.pop(-2), end=' ')
print(g)
\end{verbatim}

%enditem



\newpage
4. Що буде надруковано за виконання?

\begin{verbatim}
print('R', end=' ')
h=[10,11,12,13,14]
del h[5:6]
print(h)
\end{verbatim}

%enditem



\newpage
5. Укажіть ТИП значення виразу.\par
\begin{verbatim}
[7,2,3,4]
\end{verbatim}
%enditem


\newpage
6. Що буде надруковано за виконання?\par
\begin{verbatim}
try:
    s = {16:7, 51:9, 41:8, 16:1}
    s[51] = 1
    for x in s :
        print(x, sep=' ', end=' ')
except:
    print('ERR')
\end{verbatim}
%enditem


\newpage
7. 
Записати ВИРАЗ, значенням якого є список, що складається з квадратних коренів цілих чисел від 17 до 2018 (включно), що не діляться на жодне з чисел 2, 3.
%enditem


\newpage
8. Що буде надруковано за виконання?
\begin{verbatim}
a = 0
n = 1

class C:
    a = 2
    
    def __init__(self):
        global a
        self.a = 3
        a = 4
        C.a = 5

obj = C()
try:
    print(a, end=' ')
    print(n, end=' ')
    print(C.a, end=' ')
    print(obj.a)
except:
    print('ERR')
\end{verbatim}
%enditem


\newpage
9. Що буде надруковано за виконання?
\begin{verbatim}
class D0:
    def d(self): print(1, end=' ')
    def h(self): print(2, end=' ')

class D1(D0):
    def h(self): print(3, end=' ')

try:
    D1().d()
    D1().h()
except:
    print('ERR')
\end{verbatim}
%enditem


\newpage
10. Що буде надруковано за виконання?
\begin{verbatim}
class A:
    a = 1

    def __init__(self):
        c = 2
        c = 3

class B(A):
    a = 4
    b = 5

    def __init__(self):
        super().__init__()
        self.b = 6

try:
    obj = B()
    obj.b = obj.c + 10
    B.b = 13
except: print('EE', end=' ')

try: print(obj.b, end= ' ')
except: print('E1', end=' ')

try: print(A.b, end= ' ')
except: print('E2', end=' ')

try: print(B.b, end= ' ')
except: print('E3', end=' ')
\end{verbatim}
%enditem


\newpage
11. Змінні \kw{next} та \kw{prev} вузла списку позначають наступний та попередній за ним вузол відповідно. Починаючи з голови, у список записано послідовні цілі числа від~$-59$ до~$48$. Нехай змінна \kw{p} позначає вузол, що зберігає значення~$-4$.
Яке значення зберігається у вузлі \kw{p.prev.next.prev.next.next} ?
%enditem


\newpage
12. 
\begin{center}
	\adjustimage{max size={0.8\linewidth}{0.8\paperheight}}
	{BinTree/trees_exp/7.png}
\end{center}
Указати послідовність міток вершин дерева, зображеного на рис., що утвориться в результаті обходу дерева в прямому порядку. Мітки записувати через пробіл.\par Алгоритм починає роботу з кореня (на рис. це верхня вершина).
%enditem


\newpage
13. 
\begin{center}
	\adjustimage{max size={0.9\linewidth}{0.9\paperheight}}
	{Graphs/AdjList/310.png}
\end{center}
Для графа, зображеного на рис., вказати список суміжності вершини 3.\par
Формат оформлення відповіді:\par
список суміжності виписувати через пробіл на одному рядку, якщо пробілів десь буде більше одного - не важливо, упорядкування вершин у списку також значення не має.
%enditem


\newpage
14. Для графа на рисунку з вершини 2 запущено алгоритм пошуку в глибину, що будує кістякове дерево. Списки суміжності вершин переглядаються алгоритмом за зростанням міток вершин.\par
\begin{center}
	\adjustimage{max size={0.9\linewidth}{0.9\paperheight}}
	{Graphs/Traversals/89.png}
\end{center}
\par Які з наведених ребер входять до складу отриманого кістякового дерева?\par
1) (1, 3)\par
2) (0, 1)\par
3) (1, 5)\par
4) (0, 4)\par
5) (0, 3)\par
6) (0, 6)\par
{\vskip 15pt} Номери відповідних ребер записати за зростанням через пробіл.

%enditem


\newpage
15. Для графа на рисунку з вершини 2 запущено алгоритм пошуку в ширину, що будує кістякове дерево. Списки суміжності вершин переглядаються алгоритмом за зростанням міток вершин.\par
\begin{center}
	\adjustimage{max size={0.9\linewidth}{0.9\paperheight}}
	{Graphs/Traversals/234.png}
\end{center}
\par Які з наведених ребер входять до складу отриманого кістякового дерева?\par
1) (1, 3)\par
2) (1, 4)\par
3) (0, 3)\par
4) (0, 4)\par
5) (5, 6)\par
6) (2, 6)\par
{\vskip 15pt} Номери відповідних ребер записати за зростанням через пробіл.


%enditem


\newpage
16. 

Рядки журналу через двокрапку містять ім'я користувача (ідентифікатор), час події,тип події, код події, наприклад
\begin{verbatim}
s="username : 2022-05-05 13:26 : ERROR : 404"
\end{verbatim}
у коді 
\begin{verbatim}
res = re.fullmatch('???', s) 
s = ''
code_str = ??? 
\end{verbatim}
чим треба замінити кожен з ???, щоб значенням code\_str був код події? 


\vspace{10mm}
\textbf{Завдання 17} 

Задано програму, що містить код 1-2 класів, можливе успадкування, та клієнтський код.
У наведеному коді невеликі частини закриті. 
Також наведено вихід програми.
Необхідно відновити закриті частини, щоб отриманий код відповідав наведеному виходу програми.


\textbf{Завдання 18-19} 

Кілька запитань за завданнями для самостійної роботи.
\vspace{10mm}
Скоріш за все, що завдань буде трохи менше, а деякі попередні будуть переформульовані в стилі двох останніх.

%enditem

\end {large}}

%end

\newpage
%begin

{{
\begin {large}

\newpage
1. Що буде надруковано за виконання?

\begin{verbatim}
print('R', end=' ')
c=('a','b','c','d','e',0,1,2,4)
print(c[-5])
\end{verbatim}

%enditem


\newpage
2. Що буде надруковано за виконання?

\begin{verbatim}
print('R', end=' ')
f=('0','1','2','3','4','5','6','7','8','9')
print(f[:-4:3])
\end{verbatim}

%enditem



\newpage
3. Що буде надруковано за виконання?

\begin{verbatim}
print('R', end=' ')
h=['0','1','2','3','4']
h.append('11')
print(h)
\end{verbatim}

%enditem



\newpage
4. Що буде надруковано за виконання?

\begin{verbatim}
print('R', end=' ')
e=['0','1','2','3','4']
e[-9:-3]=()
print(e)
\end{verbatim}

%enditem



\newpage
5. Укажіть ТИП значення виразу.\par
\begin{verbatim}
{x for x in range(7) if el%3==1}
\end{verbatim}
%enditem


\newpage
6. Що буде надруковано за виконання?\par
\begin{verbatim}
try:
    s = {84:9, 50:0, 22:1, 22:1}
    s[88] = 1
    for x in s.values():
        print(x, sep=' ', end=' ')
except:
    print('ERR')
\end{verbatim}
%enditem


\newpage
7. Нехай змінна \kw{s} позначає послідовність цілих, по якій можна ітерувати.

Написати вираз, що перевіряє, чи правда, що всі елементи послідовності є точними квадратами.
%enditem


\newpage
8. Що буде надруковано за виконання?
%Оригинал был утерян. Требует отладки!
\begin{verbatim}
d = 0

class A:
    d = 1
    
    def __init__(self):
        self.d = 2
        d = 3
        A.d = 4

obj = A()
obj.d = 5
try:
    print(d, end=' ')
    print(obj.d, end=' ')
    print(A.d, end=' ')
    print(obj.d, end=' ')
    print(obj._A__d)
except:
    print('ERR')
\end{verbatim}
%enditem


\newpage
9. Що буде надруковано за виконання?
\begin{verbatim}
class B0:
    def b(self): print(1, end=' ')
    def h(self): print(2, end=' ')

class B1(B0):
    def b(self): print(3, end=' ')

try:
    B1().h()
    B1().b()
except:
    print('ERR')
\end{verbatim}
%enditem


\newpage
10. Що буде надруковано за виконання?
\begin{verbatim}
class A:
    b = 1

    def _f1(self): print(2, end=' ')
    def _f2(self): print(3, end=' ')

    def main(self):
        try: print(self.b, end=' ')
        except: print('E1', end=' ')

        try: self._f1()
        except: print('E2', end=' ')

        try: self._f2()
        except: print('E3', end=' ')


class B(A):
    b = 4

    def _f1(self): print(5, end=' ')
    def _f2(self): print(6, end=' ')

try: B().main()
except: print('E4', end=' ')
\end{verbatim}
%enditem


\newpage
11. Змінні \kw{next} та \kw{prev} вузла списку позначають наступний та попередній за ним вузол відповідно. Починаючи з голови, у список записано послідовні цілі числа від~$-67$ до~$35$. Нехай змінна \kw{p} позначає вузол, що зберігає значення~$-5$.
Яке значення зберігається у вузлі \kw{p.next.prev.next.prev.prev} ?
%enditem


\newpage
12. 
\begin{center}
	\adjustimage{max size={0.8\linewidth}{0.8\paperheight}}
	{BinTree/trees_exp/97.png}
\end{center}
Указати послідовність міток вершин дерева, зображеного на рис., що утвориться в результаті обходу дерева в оберненому порядку. Мітки записувати через пробіл.\par Алгоритм починає роботу з кореня (на рис. це верхня вершина).
%enditem


\newpage
13. 
\begin{center}
	\adjustimage{max size={0.9\linewidth}{0.9\paperheight}}
	{Graphs/AdjList/129.png}
\end{center}
Для графа, зображеного на рис., вказати список суміжності вершини 1.\par
Формат оформлення відповіді:\par
список суміжності виписувати через пробіл на одному рядку, якщо пробілів десь буде більше одного - не важливо, упорядкування вершин у списку також значення не має.
%enditem


\newpage
14. Для графа на рисунку вказати послідовність, в якій алгоритм пошуку в глибину буде відмічати відвідані вершини, якщо списки суміжності вершин переглядаються алгоритмом за зростанням міток вершин та вершини відмічаються на першому вход\-женні у вершину.\par
Пошук розпочинається з вершини 2.\par

\begin{center}
	\adjustimage{max size={0.9\linewidth}{0.9\paperheight}}
	{Graphs/Traversals/126.png}
\end{center}
%enditem


\newpage
15. Для графа на рисунку з вершини 5 запущено алгоритм пошуку в ширину, що будує кістякове дерево. Списки суміжності вершин переглядаються алгоритмом за зростанням міток вершин.\par
\begin{center}
	\adjustimage{max size={0.9\linewidth}{0.9\paperheight}}
	{Graphs/Traversals/147.png}
\end{center}
\par Які з наведених ребер входять до складу отриманого кістякового дерева?\par
1) (1, 3)\par
2) (3, 6)\par
3) (2, 4)\par
4) (4, 5)\par
5) (1, 4)\par
6) (0, 2)\par
{\vskip 15pt} Номери відповідних ребер записати за зростанням через пробіл.


%enditem


\newpage
16. 

Рядки журналу через двокрапку містять ім'я користувача (ідентифікатор), час події,тип події, код події, наприклад
\begin{verbatim}
s="username : 2022-05-05 13:26 : ERROR : 404"
\end{verbatim}
у коді 
\begin{verbatim}
res = re.fullmatch('???', s) 
s = ''
code_str = ??? 
\end{verbatim}
чим треба замінити кожен з ???, щоб значенням code\_str був код події? 


\vspace{10mm}
\textbf{Завдання 17} 

Задано програму, що містить код 1-2 класів, можливе успадкування, та клієнтський код.
У наведеному коді невеликі частини закриті. 
Також наведено вихід програми.
Необхідно відновити закриті частини, щоб отриманий код відповідав наведеному виходу програми.


\textbf{Завдання 18-19} 

Кілька запитань за завданнями для самостійної роботи.
\vspace{10mm}
Скоріш за все, що завдань буде трохи менше, а деякі попередні будуть переформульовані в стилі двох останніх.

%enditem

\end {large}}

%end

\newpage
%begin

{{
\begin {large}

\newpage
1. Що буде надруковано за виконання?

\begin{verbatim}
print('R', end=' ')
a=('a','b','c','d','e',0,1,2)
print(a[-11])
\end{verbatim}

%enditem


\newpage
2. Що буде надруковано за виконання?

\begin{verbatim}
print('R', end=' ')
e=['0','1','2','3','4','5','6','7','8','9']
print(e[-10:-13:2])
\end{verbatim}

%enditem



\newpage
3. Що буде надруковано за виконання?

\begin{verbatim}
print('R', end=' ')
d=['0','1','2','3','4']
d.remove('5')
print(d)
\end{verbatim}

%enditem



\newpage
4. Що буде надруковано за виконання?

\begin{verbatim}
print('R', end=' ')
a=['0','1','2','3','4']
a[2:-2]=[1,0]
print(a)
\end{verbatim}

%enditem



\newpage
5. Укажіть ТИП значення виразу.\par
\begin{verbatim}
{x*x:x for x in range(-7,7)}
\end{verbatim}
%enditem


\newpage
6. Що буде надруковано за виконання?\par
\begin{verbatim}
try:
    t = {37:4, 84:3, 46:1, 59:8, 84:8}
    t[17] = 5
    for x in t.items():
        print(*x, sep=' ', end=' ')
except:
    print('ERR')
\end{verbatim}
%enditem


\newpage
7. 
 Задано числову послідовність \kw{s} довжини 6000. Записати ВИРАЗ, значенням якого є список, що складається з другої половини послідовності, записаної через один елемент.
%enditem


\newpage
8. Що буде надруковано за виконання?
%Оригинал был утерян. Требует отладки!
\begin{verbatim}
_b = 0

class B:
    _b = 1
    
    def __init__(self):
        self._b = 2
        _b = 3
        B._b = 4

obj = B()
B._b = 5
try:
    print(_b, end=' ')
    print(obj._b, end=' ')
    print(B._b, end=' ')
    print(obj._b, end=' ')
    print(obj._B__b)
except:
    print('ERR')
\end{verbatim}
%enditem


\newpage
9. Що буде надруковано за виконання?
\begin{verbatim}
class A0:
    def c(self): print(1, end=' ')
    def g(self): print(2, end=' ')

class A1(A0):
    def g(self): print(3, end=' ')

try:
    A1().c()
    A0().g()
except:
    print('ERR')
\end{verbatim}
%enditem


\newpage
10. Що буде надруковано за виконання?
\begin{verbatim}
class A:
    b = 1

    def __init__(self):
        c = 2
        b = 3

class B(A):
    c = 4
    a = 5

    def __init__(self):
        super().__init__()
        self.a = 6

try:
    obj = B()
    obj.a = obj.c + 10
    B.a = 13
except: print('EE', end=' ')

try: print(obj.a, end= ' ')
except: print('E1', end=' ')

try: print(A.a, end= ' ')
except: print('E2', end=' ')

try: print(B.a, end= ' ')
except: print('E3', end=' ')
\end{verbatim}
%enditem


\newpage
11. Змінні \kw{next} та \kw{prev} вузла списку позначають наступний та попередній за ним вузол відповідно. Починаючи з голови, у список записано послідовні цілі числа від~$-35$ до~$46$. Нехай змінна \kw{p} позначає вузол, що зберігає значення~$3$.
Яке значення зберігається у вузлі \kw{p.next.next.prev.prev.next} ?
%enditem


\newpage
12. 
\begin{center}
	\adjustimage{max size={0.8\linewidth}{0.8\paperheight}}
	{BinTree/trees_exp/160.png}
\end{center}
Указати послідовність міток вершин дерева, зображеного на рис., що утвориться в результаті обходу дерева в прямому порядку. Мітки записувати через пробіл.\par Алгоритм починає роботу з кореня (на рис. це верхня вершина).
%enditem


\newpage
13. 
\begin{center}
	\adjustimage{max size={0.9\linewidth}{0.9\paperheight}}
	{Graphs/AdjList/38.png}
\end{center}
Для графа, зображеного на рис., вказати список суміжності вершини 3.\par
Формат оформлення відповіді:\par
список суміжності виписувати через пробіл на одному рядку, якщо пробілів десь буде більше одного - не важливо, упорядкування вершин у списку також значення не має.
%enditem


\newpage
14. Для графа на рисунку з вершини 3 запущено алгоритм пошуку в глибину, що будує кістякове дерево. Списки суміжності вершин переглядаються алгоритмом за зростанням міток вершин.\par
\begin{center}
	\adjustimage{max size={0.9\linewidth}{0.9\paperheight}}
	{Graphs/Traversals/145.png}
\end{center}
\par Які з наведених ребер входять до складу отриманого кістякового дерева?\par
1) (1, 4)\par
2) (5, 6)\par
3) (0, 3)\par
4) (4, 5)\par
5) (1, 3)\par
6) (0, 1)\par
{\vskip 15pt} Номери відповідних ребер записати за зростанням через пробіл.

%enditem


\newpage
15. Для графа на рисунку з вершини 1 запущено алгоритм пошуку в ширину, що будує кістякове дерево. Списки суміжності вершин переглядаються алгоритмом за зростанням міток вершин.\par
\begin{center}
	\adjustimage{max size={0.9\linewidth}{0.9\paperheight}}
	{Graphs/Traversals/220.png}
\end{center}
\par Які з наведених ребер входять до складу отриманого кістякового дерева?\par
1) (1, 6)\par
2) (2, 3)\par
3) (1, 3)\par
4) (4, 5)\par
5) (2, 4)\par
6) (0, 5)\par
{\vskip 15pt} Номери відповідних ребер записати за зростанням через пробіл.


%enditem


\newpage
16. 

Рядки журналу через двокрапку містять ім'я користувача (ідентифікатор), час події,тип події, код події, наприклад
\begin{verbatim}
s="username : 2022-05-05 13:26 : ERROR : 404"
\end{verbatim}
у коді 
\begin{verbatim}
res = re.fullmatch('???', s) 
s = ''
code_str = ??? 
\end{verbatim}
чим треба замінити кожен з ???, щоб значенням code\_str був код події? 


\vspace{10mm}
\textbf{Завдання 17} 

Задано програму, що містить код 1-2 класів, можливе успадкування, та клієнтський код.
У наведеному коді невеликі частини закриті. 
Також наведено вихід програми.
Необхідно відновити закриті частини, щоб отриманий код відповідав наведеному виходу програми.


\textbf{Завдання 18-19} 

Кілька запитань за завданнями для самостійної роботи.
\vspace{10mm}
Скоріш за все, що завдань буде трохи менше, а деякі попередні будуть переформульовані в стилі двох останніх.

%enditem

\end {large}}

%end

\newpage
%begin

{{
\begin {large}

\newpage
1. Що буде надруковано за виконання?

\begin{verbatim}
print('R', end=' ')
d=('a','b','c','d','e',0,1)
print(d[8])
\end{verbatim}

%enditem


\newpage
2. Що буде надруковано за виконання?

\begin{verbatim}
print('R', end=' ')
a=(0,1,2,3,4,5,6,7,8,9)
print(a[5::-3])
\end{verbatim}

%enditem



\newpage
3. Що буде надруковано за виконання?

\begin{verbatim}
print('R', end=' ')
c=[10,11,12,13,14]
c+=16
print(c)
\end{verbatim}

%enditem



\newpage
4. Що буде надруковано за виконання?

\begin{verbatim}
print('R', end=' ')
b=[10,11,12,13,14]
del b[-8:2]
print(b)
\end{verbatim}

%enditem



\newpage
5. Укажіть ТИП значення виразу.\par
\begin{verbatim}
{y:y*y for y in range(-7,7)}
\end{verbatim}
%enditem


\newpage
6. Що буде надруковано за виконання?\par
\begin{verbatim}
try:
    d = {28:7, 16:8, 86:9, 38:9}
    d[38] = 2
    for x in d.keys():
        print(x, sep=' ', end=' ')
except:
    print('ERR')
\end{verbatim}
%enditem


\newpage
7. Нехай змінна \kw{s} позначає дуже довгу послідовність цілих, по якій можна ітерувати. 

Написати вираз-генератор, що задає підпослідовність її елементів, що не діляться на жодне з чисел 2, 7, 5. 

Обмеження: усі елементи шуканої підпослідовності одночасно в пам'яті розміщені бути не можуть. 
%enditem


\newpage
8. Що буде надруковано за виконання?
\begin{verbatim}
a = 0
m = 1

class C:
    a = 2
    
    def __init__(self):
        self.a = 3
        a = 4
        C.a = 5

obj = C()
try:
    print(a, end=' ')
    print(m, end=' ')
    print(C.a, end=' ')
    print(obj.a)
except:
    print('ERR')
\end{verbatim}
%enditem


\newpage
9. Що буде надруковано за виконання?
\begin{verbatim}
class C0:
    def a(self): print(1, end=' ')
    def f(self): print(2, end=' ')

class C1(C0):
    def a(self): print(3, end=' ')

try:
    C0().f()
    C1().a()
except:
    print('ERR')
\end{verbatim}
%enditem


\newpage
10. Що буде надруковано за виконання?
\begin{verbatim}
class A:
    a = 1

    def _h1(self): print(2, end=' ')
    def h2(self): print(3, end=' ')

    def main(self):
        try: print(self.a, end=' ')
        except: print('E1', end=' ')

        try: self._h1()
        except: print('E2', end=' ')

        try: self.h2()
        except: print('E3', end=' ')


class B(A):
    a = 4

    def _h1(self): print(5, end=' ')
    def h2(self): print(6, end=' ')

try: B().main()
except: print('E4', end=' ')
\end{verbatim}
%enditem


\newpage
11. Змінні \kw{next} та \kw{prev} вузла списку позначають наступний та попередній за ним вузол відповідно. Починаючи з голови, у список записано послідовні цілі числа від~$-52$ до~$44$. Нехай змінна \kw{p} позначає вузол, що зберігає значення~$0$.
Яке значення зберігається у вузлі \kw{p.prev.prev.next.next.prev} ?
%enditem


\newpage
12. 
\begin{center}
	\adjustimage{max size={0.8\linewidth}{0.8\paperheight}}
	{BinTree/trees_exp/133.png}
\end{center}
Указати послідовність міток вершин дерева, зображеного на рис., що утвориться в результаті обходу дерева в оберненому порядку. Мітки записувати через пробіл.\par Алгоритм починає роботу з кореня (на рис. це верхня вершина).
%enditem


\newpage
13. 
\begin{center}
	\adjustimage{max size={0.9\linewidth}{0.9\paperheight}}
	{Graphs/AdjList/31.png}
\end{center}
Для графа, зображеного на рис., вказати список суміжності вершини 9.\par
Формат оформлення відповіді:\par
список суміжності виписувати через пробіл на одному рядку, якщо пробілів десь буде більше одного - не важливо, упорядкування вершин у списку також значення не має.
%enditem


\newpage
14. Для графа на рисунку вказати послідовність, в якій алгоритм пошуку в глибину буде відмічати відвідані вершини, якщо списки суміжності вершин переглядаються алгоритмом за зростанням міток вершин та вершини відмічаються на першому вход\-женні у вершину.\par
Пошук розпочинається з вершини 6.\par

\begin{center}
	\adjustimage{max size={0.9\linewidth}{0.9\paperheight}}
	{Graphs/Traversals/286.png}
\end{center}
%enditem


\newpage
15. Для графа на рисунку з вершини 4 запущено алгоритм пошуку в ширину, що будує кістякове дерево. Списки суміжності вершин переглядаються алгоритмом за зростанням міток вершин.\par
\begin{center}
	\adjustimage{max size={0.9\linewidth}{0.9\paperheight}}
	{Graphs/Traversals/175.png}
\end{center}
\par Які з наведених ребер входять до складу отриманого кістякового дерева?\par
1) (0, 3)\par
2) (3, 6)\par
3) (2, 4)\par
4) (0, 6)\par
5) (3, 5)\par
6) (0, 1)\par
{\vskip 15pt} Номери відповідних ребер записати за зростанням через пробіл.


%enditem


\newpage
16. 

Рядки журналу через двокрапку містять ім'я користувача (ідентифікатор), час події,тип події, код події, наприклад
\begin{verbatim}
s="username : 2022-05-05 13:26 : ERROR : 404"
\end{verbatim}
у коді 
\begin{verbatim}
res = re.fullmatch('???', s) 
s = ''
code_str = ??? 
\end{verbatim}
чим треба замінити кожен з ???, щоб значенням code\_str був код події? 


\vspace{10mm}
\textbf{Завдання 17} 

Задано програму, що містить код 1-2 класів, можливе успадкування, та клієнтський код.
У наведеному коді невеликі частини закриті. 
Також наведено вихід програми.
Необхідно відновити закриті частини, щоб отриманий код відповідав наведеному виходу програми.


\textbf{Завдання 18-19} 

Кілька запитань за завданнями для самостійної роботи.
\vspace{10mm}
Скоріш за все, що завдань буде трохи менше, а деякі попередні будуть переформульовані в стилі двох останніх.

%enditem

\end {large}}

%end

\newpage
%begin

{{
\begin {large}

\newpage
1. Що буде надруковано за виконання?

\begin{verbatim}
print('R', end=' ')
h=('a','b','c','d','e',0,1,2,3,4)
print(h[10])
\end{verbatim}

%enditem


\newpage
2. Що буде надруковано за виконання?

\begin{verbatim}
print('R', end=' ')
a='0123456789'
print(a[::-2])
\end{verbatim}

%enditem



\newpage
3. Що буде надруковано за виконання?

\begin{verbatim}
print('R', end=' ')
c=['0','1','2','3','4']
c.remove('6')
print(c)
\end{verbatim}

%enditem



\newpage
4. Що буде надруковано за виконання?

\begin{verbatim}
print('R', end=' ')
c=['0','1','2','3','4']
c[-4:]=[4,5]
print(c)
\end{verbatim}

%enditem



\newpage
5. Укажіть ТИП значення виразу.\par
\begin{verbatim}
{1:3, 5:8}
\end{verbatim}
%enditem


\newpage
6. Що буде надруковано за виконання?\par
\begin{verbatim}
try:
    t = {52:7, 62:5, 25:9, 70:6}
    t[62] = 5
    for x in t.items():
        print(*x, sep=' ', end=' ')
except:
    print('ERR')
\end{verbatim}
%enditem


\newpage
7. Нехай змінна \kw{s} позначає послідовність цілих, по якій можна ітерувати. 

Написати вираз, що перевіряє, чи правда, що послідовність  містить хоча б одне число, що ділиться на 7.
%enditem


\newpage
8. Що буде надруковано за виконання?
\begin{verbatim}
c = 0
h = 1

class A:
    h = 2
    
    def __init__(self):
        global h
        self.h = 3
        h = 4
        A.h = 5

obj = A()
try:
    print(c, end=' ')
    print(h, end=' ')
    print(A.h, end=' ')
    print(obj.h)
except:
    print('ERR')
\end{verbatim}
%enditem


\newpage
9. Що буде надруковано за виконання?
\begin{verbatim}
class D0:
    def b(self): print(1, end=' ')
    def f(self): print(2, end=' ')

class D1(D0):
    def f(self): print(3, end=' ')

try:
    D1().b()
    D0().f()
except:
    print('ERR')
\end{verbatim}
%enditem


\newpage
10. Що буде надруковано за виконання?
\begin{verbatim}
class A:
    c = 1

    def __init__(self):
        b = 2
        b = 3

class B(A):
    c = 4
    a = 5

    def __init__(self):
        super().__init__()
        self.a = 6

try:
    obj = B()
    obj.b = obj.b + 10
    B.b = 13
except: print('EE', end=' ')

try: print(obj.b, end= ' ')
except: print('E1', end=' ')

try: print(A.b, end= ' ')
except: print('E2', end=' ')

try: print(B.b, end= ' ')
except: print('E3', end=' ')
\end{verbatim}
%enditem


\newpage
11. Змінні \kw{next} та \kw{prev} вузла списку позначають наступний та попередній за ним вузол відповідно. Починаючи з голови, у список записано послідовні цілі числа від~$-36$ до~$31$. Нехай змінна \kw{p} позначає вузол, що зберігає значення~$6$.
Яке значення зберігається у вузлі \kw{p.prev.next.next.prev.next} ?
%enditem


\newpage
12. 
\begin{center}
	\adjustimage{max size={0.8\linewidth}{0.8\paperheight}}
	{BinTree/trees_exp/35.png}
\end{center}
Указати послідовність міток вершин дерева, зображеного на рис., що утвориться в результаті обходу дерева в оберненому порядку. Мітки записувати через пробіл.\par Алгоритм починає роботу з кореня (на рис. це верхня вершина).
%enditem


\newpage
13. 
\begin{center}
	\adjustimage{max size={0.9\linewidth}{0.9\paperheight}}
	{Graphs/AdjList/14.png}
\end{center}
Для графа, зображеного на рис., вказати список суміжності вершини 5.\par
Формат оформлення відповіді:\par
список суміжності виписувати через пробіл на одному рядку, якщо пробілів десь буде більше одного - не важливо, упорядкування вершин у списку також значення не має.
%enditem


\newpage
14. Для графа на рисунку вказати послідовність, в якій алгоритм пошуку в глибину буде відмічати відвідані вершини, якщо списки суміжності вершин переглядаються алгоритмом за зростанням міток вершин та вершини відмічаються на першому вход\-женні у вершину.\par
Пошук розпочинається з вершини 6.\par

\begin{center}
	\adjustimage{max size={0.9\linewidth}{0.9\paperheight}}
	{Graphs/Traversals/143.png}
\end{center}
%enditem


\newpage
15. Для графа на рисунку з вершини 6 запущено алгоритм пошуку в ширину, що будує кістякове дерево. Списки суміжності вершин переглядаються алгоритмом за зростанням міток вершин.\par
\begin{center}
	\adjustimage{max size={0.9\linewidth}{0.9\paperheight}}
	{Graphs/Traversals/150.png}
\end{center}
\par Які з наведених ребер входять до складу отриманого кістякового дерева?\par
1) (3, 6)\par
2) (0, 5)\par
3) (1, 4)\par
4) (0, 6)\par
5) (2, 6)\par
6) (4, 6)\par
{\vskip 15pt} Номери відповідних ребер записати за зростанням через пробіл.


%enditem


\newpage
16. 

Рядки журналу через двокрапку містять ім'я користувача (ідентифікатор), час події,тип події, код події, наприклад
\begin{verbatim}
s="username : 2022-05-05 13:26 : ERROR : 404"
\end{verbatim}
у коді 
\begin{verbatim}
res = re.fullmatch('???', s) 
s = ''
code_str = ??? 
\end{verbatim}
чим треба замінити кожен з ???, щоб значенням code\_str був код події? 


\vspace{10mm}
\textbf{Завдання 17} 

Задано програму, що містить код 1-2 класів, можливе успадкування, та клієнтський код.
У наведеному коді невеликі частини закриті. 
Також наведено вихід програми.
Необхідно відновити закриті частини, щоб отриманий код відповідав наведеному виходу програми.


\textbf{Завдання 18-19} 

Кілька запитань за завданнями для самостійної роботи.
\vspace{10mm}
Скоріш за все, що завдань буде трохи менше, а деякі попередні будуть переформульовані в стилі двох останніх.

%enditem

\end {large}}

%end

\newpage
%begin

{{
\begin {large}

\newpage
1. Що буде надруковано за виконання?

\begin{verbatim}
print('R', end=' ')
f=('a','b','c','d','e',0,1,2,3)
print(f[-12])
\end{verbatim}

%enditem


\newpage
2. Що буде надруковано за виконання?

\begin{verbatim}
print('R', end=' ')
e=[0,1,2,3,4,5,6,7,8,9]
print(e[14:-14:-1])
\end{verbatim}

%enditem



\newpage
3. Що буде надруковано за виконання?

\begin{verbatim}
print('R', end=' ')
b=[10,11,12,13,14]
b.extend('10')
print(b)
\end{verbatim}

%enditem



\newpage
4. Що буде надруковано за виконання?

\begin{verbatim}
print('R', end=' ')
a=[10,11,12,13,14]
del a[-6:-9]
print(a)
\end{verbatim}

%enditem



\newpage
5. Укажіть ТИП значення виразу.\par
\begin{verbatim}
{x:0 for x in range(7)}
\end{verbatim}
%enditem


\newpage
6. Що буде надруковано за виконання?\par
\begin{verbatim}
try:
    s = {55:6, 19:9, 22:3, 53:1, 22:2}
    s[11] = 0
    for x in s :
        print(x, sep=' ', end=' ')
except:
    print('ERR')
\end{verbatim}
%enditem


\newpage
7. Нехай змінна \kw{s} позначає дуже довгу послідовність, по якій можна ітерувати. 

Написати вираз-генератор, що задає підпослідовність її елементів, що є числами типу \kw{float}.

Обмеження: усі елементи шуканої підпослідовності одночасно в пам'яті розміщені бути не можуть. 
%enditem


\newpage
8. Що буде надруковано за виконання?
\begin{verbatim}
d = 0
n = 1

class B:
    d = 2
    
    def __init__(self):
        global d
        self.d = 3
        d = 4
        B.d = 5

obj = B()
try:
    print(d, end=' ')
    print(n, end=' ')
    print(B.d, end=' ')
    print(obj.d)
except:
    print('ERR')
\end{verbatim}
%enditem


\newpage
9. Що буде надруковано за виконання?
\begin{verbatim}
class C0:
    def a(self): print(1, end=' ')
    def h(self): print(2, end=' ')

class C1(C0):
    def a(self): print(3, end=' ')

try:
    C0().h()
    C1().a()
except:
    print('ERR')
\end{verbatim}
%enditem


\newpage
10. Що буде надруковано за виконання?
\begin{verbatim}
class A:
    d = 1

    def g1(self): print(2, end=' ')
    def _g2(self): print(3, end=' ')

    def main(self):
        try: print(self.d, end=' ')
        except: print('E1', end=' ')

        try: self.g1()
        except: print('E2', end=' ')

        try: self._g2()
        except: print('E3', end=' ')


class B(A):
    d = 4

    def g1(self): print(5, end=' ')
    def _g2(self): print(6, end=' ')

try: B().main()
except: print('E4', end=' ')
\end{verbatim}
%enditem


\newpage
11. Змінні \kw{next} та \kw{prev} вузла списку позначають наступний та попередній за ним вузол відповідно. Починаючи з голови, у список записано послідовні цілі числа від~$-47$ до~$25$. Нехай змінна \kw{p} позначає вузол, що зберігає значення~$-1$.
Яке значення зберігається у вузлі \kw{p.next.prev.prev.next.prev} ?
%enditem


\newpage
12. 
\begin{center}
	\adjustimage{max size={0.8\linewidth}{0.8\paperheight}}
	{BinTree/trees_exp/155.png}
\end{center}
Указати послідовність міток вершин дерева, зображеного на рис., що утвориться в результаті обходу дерева в прямому порядку. Мітки записувати через пробіл.\par Алгоритм починає роботу з кореня (на рис. це верхня вершина).
%enditem


\newpage
13. 
\begin{center}
	\adjustimage{max size={0.9\linewidth}{0.9\paperheight}}
	{Graphs/AdjList/281.png}
\end{center}
Для графа, зображеного на рис., вказати список суміжності вершини 2.\par
Формат оформлення відповіді:\par
список суміжності виписувати через пробіл на одному рядку, якщо пробілів десь буде більше одного - не важливо, упорядкування вершин у списку також значення не має.
%enditem


\newpage
14. Для графа на рисунку з вершини 0 запущено алгоритм пошуку в глибину, що будує кістякове дерево. Списки суміжності вершин переглядаються алгоритмом за зростанням міток вершин.\par
\begin{center}
	\adjustimage{max size={0.9\linewidth}{0.9\paperheight}}
	{Graphs/Traversals/235.png}
\end{center}
\par Які з наведених ребер входять до складу отриманого кістякового дерева?\par
1) (3, 5)\par
2) (2, 3)\par
3) (3, 4)\par
4) (0, 2)\par
5) (1, 2)\par
6) (1, 5)\par
{\vskip 15pt} Номери відповідних ребер записати за зростанням через пробіл.

%enditem


\newpage
15. Для графа на рисунку з вершини 4 запущено алгоритм пошуку в ширину, що будує кістякове дерево. Списки суміжності вершин переглядаються алгоритмом за зростанням міток вершин.\par
\begin{center}
	\adjustimage{max size={0.9\linewidth}{0.9\paperheight}}
	{Graphs/Traversals/8.png}
\end{center}
\par Які з наведених ребер входять до складу отриманого кістякового дерева?\par
1) (4, 6)\par
2) (5, 6)\par
3) (1, 2)\par
4) (2, 3)\par
5) (0, 3)\par
6) (0, 1)\par
{\vskip 15pt} Номери відповідних ребер записати за зростанням через пробіл.


%enditem


\newpage
16. 

Рядки журналу через двокрапку містять ім'я користувача (ідентифікатор), час події,тип події, код події, наприклад
\begin{verbatim}
s="username : 2022-05-05 13:26 : ERROR : 404"
\end{verbatim}
у коді 
\begin{verbatim}
res = re.fullmatch('???', s) 
s = ''
code_str = ??? 
\end{verbatim}
чим треба замінити кожен з ???, щоб значенням code\_str був код події? 


\vspace{10mm}
\textbf{Завдання 17} 

Задано програму, що містить код 1-2 класів, можливе успадкування, та клієнтський код.
У наведеному коді невеликі частини закриті. 
Також наведено вихід програми.
Необхідно відновити закриті частини, щоб отриманий код відповідав наведеному виходу програми.


\textbf{Завдання 18-19} 

Кілька запитань за завданнями для самостійної роботи.
\vspace{10mm}
Скоріш за все, що завдань буде трохи менше, а деякі попередні будуть переформульовані в стилі двох останніх.

%enditem

\end {large}}

%end

\newpage
%begin

{{
\begin {large}

\newpage
1. Що буде надруковано за виконання?

\begin{verbatim}
print('R', end=' ')
e=('a','b','c','d','e',0,1,2,3)
print(e[1])
\end{verbatim}

%enditem


\newpage
2. Що буде надруковано за виконання?

\begin{verbatim}
print('R', end=' ')
c='0123456789'
print(c[:-10:-3])
\end{verbatim}

%enditem



\newpage
3. Що буде надруковано за виконання?

\begin{verbatim}
print('R', end=' ')
a=['0','1','2','3','4']
a+=[1,2]
print(a)
\end{verbatim}

%enditem



\newpage
4. Що буде надруковано за виконання?

\begin{verbatim}
print('R', end=' ')
g=[10,11,12,13,14]
g[-1:]=[2,3]
print(g)
\end{verbatim}

%enditem



\newpage
5. Укажіть ТИП значення виразу.\par
\begin{verbatim}
(11,2,3)
\end{verbatim}
%enditem


\newpage
6. Що буде надруковано за виконання?\par
\begin{verbatim}
try:
    d = {64:4, 67:8, 83:7, 50:8, 83:5}
    d[38] = 7
    for x in d.items():
        print(x, sep=' ', end=' ')
except:
    print('ERR')
\end{verbatim}
%enditem


\newpage
7. Нехай змінна \kw{s} позначає дуже довгу послідовність цілих, по якій можна ітерувати. 

Написати вираз-генератор, що задає підпослідовність її елементів, що більші за задане число num. 

Обмеження: усі елементи шуканої підпослідовності одночасно в пам'яті розміщені бути не можуть. 
%enditem


\newpage
8. Що буде надруковано за виконання?
%Оригинал был утерян. Требует отладки!
\begin{verbatim}
__c = 0

class D:
    __c = 1
    
    def __init__(self):
        self.__c = 2
        __c = 3
        D.__c = 4

obj = D()
obj.__c = 5
try:
    print(__c, end=' ')
    print(obj.__c, end=' ')
    print(D.__c, end=' ')
    print(obj.__c, end=' ')
    print(obj._D__c)
except:
    print('ERR')
\end{verbatim}
%enditem


\newpage
9. Що буде надруковано за виконання?
\begin{verbatim}
class A0:
    def c(self): print(1, end=' ')
    def g(self): print(2, end=' ')

class A1(A0):
    def c(self): print(3, end=' ')

try:
    A1().g()
    A1().c()
except:
    print('ERR')
\end{verbatim}
%enditem


\newpage
10. Що буде надруковано за виконання?
\begin{verbatim}
class A:
    c = 1

    def __init__(self):
        a = 2
        c = 3

class B(A):
    a = 4
    b = 5

    def __init__(self):
        super().__init__()
        self.b = 6

try:
    obj = B()
    obj.a = obj.c + 10
    B.a = 13
except: print('EE', end=' ')

try: print(obj.a, end= ' ')
except: print('E1', end=' ')

try: print(A.a, end= ' ')
except: print('E2', end=' ')

try: print(B.a, end= ' ')
except: print('E3', end=' ')
\end{verbatim}
%enditem


\newpage
11. Змінні \kw{next} та \kw{prev} вузла списку позначають наступний та попередній за ним вузол відповідно. Починаючи з голови, у список записано послідовні цілі числа від~$-23$ до~$32$. Нехай змінна \kw{p} позначає вузол, що зберігає значення~$5$.
Яке значення зберігається у вузлі \kw{p.next.next.prev.prev.prev} ?
%enditem


\newpage
12. 
\begin{center}
	\adjustimage{max size={0.8\linewidth}{0.8\paperheight}}
	{BinTree/trees_exp/154.png}
\end{center}
Указати послідовність міток вершин дерева, зображеного на рис., що утвориться в результаті обходу дерева в прямому порядку. Мітки записувати через пробіл.\par Алгоритм починає роботу з кореня (на рис. це верхня вершина).
%enditem


\newpage
13. 
\begin{center}
	\adjustimage{max size={0.9\linewidth}{0.9\paperheight}}
	{Graphs/AdjList/270.png}
\end{center}
Для графа, зображеного на рис., вказати список суміжності вершини 7.\par
Формат оформлення відповіді:\par
список суміжності виписувати через пробіл на одному рядку, якщо пробілів десь буде більше одного - не важливо, упорядкування вершин у списку також значення не має.
%enditem


\newpage
14. Для графа на рисунку вказати послідовність, в якій алгоритм пошуку в глибину буде відмічати відвідані вершини, якщо списки суміжності вершин переглядаються алгоритмом за зростанням міток вершин та вершини відмічаються на першому вход\-женні у вершину.\par
Пошук розпочинається з вершини 4.\par

\begin{center}
	\adjustimage{max size={0.9\linewidth}{0.9\paperheight}}
	{Graphs/Traversals/296.png}
\end{center}
%enditem


\newpage
15. Для графа на рисунку з вершини 2 запущено алгоритм пошуку в ширину, що будує кістякове дерево. Списки суміжності вершин переглядаються алгоритмом за зростанням міток вершин.\par
\begin{center}
	\adjustimage{max size={0.9\linewidth}{0.9\paperheight}}
	{Graphs/Traversals/188.png}
\end{center}
\par Які з наведених ребер входять до складу отриманого кістякового дерева?\par
1) (4, 5)\par
2) (3, 5)\par
3) (0, 4)\par
4) (2, 4)\par
5) (4, 6)\par
6) (0, 1)\par
{\vskip 15pt} Номери відповідних ребер записати за зростанням через пробіл.


%enditem


\newpage
16. 

Рядки журналу через двокрапку містять ім'я користувача (ідентифікатор), час події,тип події, код події, наприклад
\begin{verbatim}
s="username : 2022-05-05 13:26 : ERROR : 404"
\end{verbatim}
у коді 
\begin{verbatim}
res = re.fullmatch('???', s) 
s = ''
code_str = ??? 
\end{verbatim}
чим треба замінити кожен з ???, щоб значенням code\_str був код події? 


\vspace{10mm}
\textbf{Завдання 17} 

Задано програму, що містить код 1-2 класів, можливе успадкування, та клієнтський код.
У наведеному коді невеликі частини закриті. 
Також наведено вихід програми.
Необхідно відновити закриті частини, щоб отриманий код відповідав наведеному виходу програми.


\textbf{Завдання 18-19} 

Кілька запитань за завданнями для самостійної роботи.
\vspace{10mm}
Скоріш за все, що завдань буде трохи менше, а деякі попередні будуть переформульовані в стилі двох останніх.

%enditem

\end {large}}

%end

\newpage
%begin

{{
\begin {large}

\newpage
1. Що буде надруковано за виконання?

\begin{verbatim}
print('R', end=' ')
b=('a','b','c','d','e',0,1,2,3,4)
print(b[-13])
\end{verbatim}

%enditem


\newpage
2. Що буде надруковано за виконання?

\begin{verbatim}
print('R', end=' ')
f=['0','1','2','3','4','5','6','7','8','9']
print(f[9::3])
\end{verbatim}

%enditem



\newpage
3. Що буде надруковано за виконання?

\begin{verbatim}
print('R', end=' ')
g=[10,11,12,13,14]
print(g.index(12), end=' ')
\end{verbatim}

%enditem



\newpage
4. Що буде надруковано за виконання?

\begin{verbatim}
print('R', end=' ')
d=['0','1','2','3','4']
del d[0:]
print(d)
\end{verbatim}

%enditem



\newpage
5. Укажіть ТИП значення виразу.\par
\begin{verbatim}
[x for x in range(7)]
\end{verbatim}
%enditem


\newpage
6. Що буде надруковано за виконання?\par
\begin{verbatim}
try:
    t = {53:0, 66:0, 58:9, 19:0}
    t[22] = 9
    for x, y in t.items():
        print(x, y, sep=' ', end=' ')
except:
    print('ERR')
\end{verbatim}
%enditem


\newpage
7. Нехай змінна \kw{s} позначає дуже довгу послідовність, по якій можна ітерувати. 

Написати вираз-генератор, що задає підпослідовність її елементів, що не є числами типу \kw{float}.

Обмеження: усі елементи шуканої підпослідовності одночасно в пам'яті розміщені бути не можуть. 
%enditem


\newpage
8. Що буде надруковано за виконання?
%Оригинал был утерян. Требует отладки!
\begin{verbatim}
_d = 0

class D:
    _d = 1
    
    def __init__(self):
        self._d = 2
        _d = 3
        D._d = 4

obj = D()
D._d = 5
try:
    print(_d, end=' ')
    print(obj._d, end=' ')
    print(D._d, end=' ')
    print(obj._d, end=' ')
    print(obj._D__d)
except:
    print('ERR')
\end{verbatim}
%enditem


\newpage
9. Що буде надруковано за виконання?
\begin{verbatim}
class B0:
    def d(self): print(1, end=' ')
    def g(self): print(2, end=' ')

class B1(B0):
    def g(self): print(3, end=' ')

try:
    B1().d()
    B1().g()
except:
    print('ERR')
\end{verbatim}
%enditem


\newpage
10. Що буде надруковано за виконання?
\begin{verbatim}
class A:
    c = 1

    def __f1(self): print(2, end=' ')
    def _f2(self): print(3, end=' ')

    def main(self):
        try: print(self.c, end=' ')
        except: print('E1', end=' ')

        try: self.__f1()
        except: print('E2', end=' ')

        try: self._f2()
        except: print('E3', end=' ')


class B(A):
    c = 4

    def __f1(self): print(5, end=' ')
    def _f2(self): print(6, end=' ')

try: B().main()
except: print('E4', end=' ')
\end{verbatim}
%enditem


\newpage
11. Змінні \kw{next} та \kw{prev} вузла списку позначають наступний та попередній за ним вузол відповідно. Починаючи з голови, у список записано послідовні цілі числа від~$-53$ до~$50$. Нехай змінна \kw{p} позначає вузол, що зберігає значення~$-8$.
Яке значення зберігається у вузлі \kw{p.prev.prev.next.next.next} ?
%enditem


\newpage
12. 
\begin{center}
	\adjustimage{max size={0.8\linewidth}{0.8\paperheight}}
	{BinTree/trees_exp/37.png}
\end{center}
Указати послідовність міток вершин дерева, зображеного на рис., що утвориться в результаті обходу дерева в оберненому порядку. Мітки записувати через пробіл.\par Алгоритм починає роботу з кореня (на рис. це верхня вершина).
%enditem


\newpage
13. 
\begin{center}
	\adjustimage{max size={0.9\linewidth}{0.9\paperheight}}
	{Graphs/AdjList/238.png}
\end{center}
Для графа, зображеного на рис., вказати список суміжності вершини 2.\par
Формат оформлення відповіді:\par
список суміжності виписувати через пробіл на одному рядку, якщо пробілів десь буде більше одного - не важливо, упорядкування вершин у списку також значення не має.
%enditem


\newpage
14. Для графа на рисунку з вершини 0 запущено алгоритм пошуку в глибину, що будує кістякове дерево. Списки суміжності вершин переглядаються алгоритмом за зростанням міток вершин.\par
\begin{center}
	\adjustimage{max size={0.9\linewidth}{0.9\paperheight}}
	{Graphs/Traversals/283.png}
\end{center}
\par Які з наведених ребер входять до складу отриманого кістякового дерева?\par
1) (2, 4)\par
2) (1, 4)\par
3) (1, 5)\par
4) (0, 4)\par
5) (0, 2)\par
6) (0, 1)\par
{\vskip 15pt} Номери відповідних ребер записати за зростанням через пробіл.

%enditem


\newpage
15. Для графа на рисунку з вершини 6 запущено алгоритм пошуку в ширину, що будує кістякове дерево. Списки суміжності вершин переглядаються алгоритмом за зростанням міток вершин.\par
\begin{center}
	\adjustimage{max size={0.9\linewidth}{0.9\paperheight}}
	{Graphs/Traversals/108.png}
\end{center}
\par Які з наведених ребер входять до складу отриманого кістякового дерева?\par
1) (2, 4)\par
2) (1, 6)\par
3) (4, 5)\par
4) (1, 4)\par
5) (3, 4)\par
6) (3, 5)\par
{\vskip 15pt} Номери відповідних ребер записати за зростанням через пробіл.


%enditem


\newpage
16. 

Рядки журналу через двокрапку містять ім'я користувача (ідентифікатор), час події,тип події, код події, наприклад
\begin{verbatim}
s="username : 2022-05-05 13:26 : ERROR : 404"
\end{verbatim}
у коді 
\begin{verbatim}
res = re.fullmatch('???', s) 
s = ''
code_str = ??? 
\end{verbatim}
чим треба замінити кожен з ???, щоб значенням code\_str був код події? 


\vspace{10mm}
\textbf{Завдання 17} 

Задано програму, що містить код 1-2 класів, можливе успадкування, та клієнтський код.
У наведеному коді невеликі частини закриті. 
Також наведено вихід програми.
Необхідно відновити закриті частини, щоб отриманий код відповідав наведеному виходу програми.


\textbf{Завдання 18-19} 

Кілька запитань за завданнями для самостійної роботи.
\vspace{10mm}
Скоріш за все, що завдань буде трохи менше, а деякі попередні будуть переформульовані в стилі двох останніх.

%enditem

\end {large}}

%end

\newpage
%begin

{{
\begin {large}

\newpage
1. Що буде надруковано за виконання?

\begin{verbatim}
print('R', end=' ')
a=('a','b','c','d','e',0,1)
print(a[5])
\end{verbatim}

%enditem


\newpage
2. Що буде надруковано за виконання?

\begin{verbatim}
print('R', end=' ')
d=('0','1','2','3','4','5','6','7','8','9')
print(d[-7:-3:-1])
\end{verbatim}

%enditem



\newpage
3. Що буде надруковано за виконання?

\begin{verbatim}
print('R', end=' ')
f=['0','1','2','3','4']
f.append([1,0])
print(f)
\end{verbatim}

%enditem



\newpage
4. Що буде надруковано за виконання?

\begin{verbatim}
print('R', end=' ')
b=[10,11,12,13,14]
b[1:-1]='13'
print(b)
\end{verbatim}

%enditem



\newpage
5. Укажіть ТИП значення виразу.\par
\begin{verbatim}
{x for x in range(7) if el%3==1}
\end{verbatim}
%enditem


\newpage
6. Що буде надруковано за виконання?\par
\begin{verbatim}
try:
    s = {23:2, 24:7, 48:8, 23:9}
    s[24] = 6
    for x in s.values():
        print(x, sep=' ', end=' ')
except:
    print('ERR')
\end{verbatim}
%enditem


\newpage
7. 
Записати ВИРАЗ, значенням якого є список, що складається з квадратних коренів цілих чисел від 18 до 2017 (включно), що не діляться на жодне з чисел 2, 3.
%enditem


\newpage
8. Що буде надруковано за виконання?
%Оригинал был утерян. Требует отладки!
\begin{verbatim}
__a = 0

class A:
    __a = 1
    
    def __init__(self):
        self.__a = 2
        __a = 3
        A.__a = 4

obj = A()
obj.__a = 5
try:
    print(__a, end=' ')
    print(obj.__a, end=' ')
    print(A.__a, end=' ')
    print(obj.__a, end=' ')
    print(obj._A__a)
except:
    print('ERR')
\end{verbatim}
%enditem


\newpage
9. Що буде надруковано за виконання?
\begin{verbatim}
class D0:
    def a(self): print(1, end=' ')
    def f(self): print(2, end=' ')

class D1(D0):
    def f(self): print(3, end=' ')

try:
    D0().a()
    D1().f()
except:
    print('ERR')
\end{verbatim}
%enditem


\newpage
10. Що буде надруковано за виконання?
\begin{verbatim}
class A:
    b = 1

    def _h1(self): print(2, end=' ')
    def _h2(self): print(3, end=' ')

    def main(self):
        try: print(self.b, end=' ')
        except: print('E1', end=' ')

        try: self._h1()
        except: print('E2', end=' ')

        try: self._h2()
        except: print('E3', end=' ')


class B(A):
    b = 4

    def _h1(self): print(5, end=' ')
    def _h2(self): print(6, end=' ')

try: B().main()
except: print('E4', end=' ')
\end{verbatim}
%enditem


\newpage
11. Змінні \kw{next} та \kw{prev} вузла списку позначають наступний та попередній за ним вузол відповідно. Починаючи з голови, у список записано послідовні цілі числа від~$-46$ до~$26$. Нехай змінна \kw{p} позначає вузол, що зберігає значення~$-5$.
Яке значення зберігається у вузлі \kw{p.next.next.prev.next.prev} ?
%enditem


\newpage
12. 
\begin{center}
	\adjustimage{max size={0.8\linewidth}{0.8\paperheight}}
	{BinTree/trees_exp/110.png}
\end{center}
Указати послідовність міток вершин дерева, зображеного на рис., що утвориться в результаті обходу дерева в прямому порядку. Мітки записувати через пробіл.\par Алгоритм починає роботу з кореня (на рис. це верхня вершина).
%enditem


\newpage
13. 
\begin{center}
	\adjustimage{max size={0.9\linewidth}{0.9\paperheight}}
	{Graphs/AdjList/136.png}
\end{center}
Для графа, зображеного на рис., вказати список суміжності вершини 0.\par
Формат оформлення відповіді:\par
список суміжності виписувати через пробіл на одному рядку, якщо пробілів десь буде більше одного - не важливо, упорядкування вершин у списку також значення не має.
%enditem


\newpage
14. Для графа на рисунку з вершини 3 запущено алгоритм пошуку в глибину, що будує кістякове дерево. Списки суміжності вершин переглядаються алгоритмом за зростанням міток вершин.\par
\begin{center}
	\adjustimage{max size={0.9\linewidth}{0.9\paperheight}}
	{Graphs/Traversals/243.png}
\end{center}
\par Які з наведених ребер входять до складу отриманого кістякового дерева?\par
1) (4, 6)\par
2) (0, 3)\par
3) (3, 4)\par
4) (0, 1)\par
5) (0, 6)\par
6) (0, 4)\par
{\vskip 15pt} Номери відповідних ребер записати за зростанням через пробіл.

%enditem


\newpage
15. Для графа на рисунку з вершини 1 запущено алгоритм пошуку в ширину, що будує кістякове дерево. Списки суміжності вершин переглядаються алгоритмом за зростанням міток вершин.\par
\begin{center}
	\adjustimage{max size={0.9\linewidth}{0.9\paperheight}}
	{Graphs/Traversals/242.png}
\end{center}
\par Які з наведених ребер входять до складу отриманого кістякового дерева?\par
1) (1, 6)\par
2) (1, 4)\par
3) (2, 5)\par
4) (4, 5)\par
5) (0, 4)\par
6) (3, 4)\par
{\vskip 15pt} Номери відповідних ребер записати за зростанням через пробіл.


%enditem


\newpage
16. 

Рядки журналу через двокрапку містять ім'я користувача (ідентифікатор), час події,тип події, код події, наприклад
\begin{verbatim}
s="username : 2022-05-05 13:26 : ERROR : 404"
\end{verbatim}
у коді 
\begin{verbatim}
res = re.fullmatch('???', s) 
s = ''
code_str = ??? 
\end{verbatim}
чим треба замінити кожен з ???, щоб значенням code\_str був код події? 


\vspace{10mm}
\textbf{Завдання 17} 

Задано програму, що містить код 1-2 класів, можливе успадкування, та клієнтський код.
У наведеному коді невеликі частини закриті. 
Також наведено вихід програми.
Необхідно відновити закриті частини, щоб отриманий код відповідав наведеному виходу програми.


\textbf{Завдання 18-19} 

Кілька запитань за завданнями для самостійної роботи.
\vspace{10mm}
Скоріш за все, що завдань буде трохи менше, а деякі попередні будуть переформульовані в стилі двох останніх.

%enditem

\end {large}}

%end

\newpage
%begin

{{
\begin {large}

\newpage
1. Що буде надруковано за виконання?

\begin{verbatim}
print('R', end=' ')
c=('a','b','c','d','e',0,1,2,4)
print(c[-4])
\end{verbatim}

%enditem


\newpage
2. Що буде надруковано за виконання?

\begin{verbatim}
print('R', end=' ')
g=(0,1,2,3,4,5,6,7,8,9)
print(g[-15:-2:2])
\end{verbatim}

%enditem



\newpage
3. Що буде надруковано за виконання?

\begin{verbatim}
print('R', end=' ')
d=[10,11,12,13,14]
d.insert(-2, 19)
print(d)
\end{verbatim}

%enditem



\newpage
4. Що буде надруковано за виконання?

\begin{verbatim}
print('R', end=' ')
f=['0','1','2','3','4']
f[3:-8]=(13)
print(f)
\end{verbatim}

%enditem



\newpage
5. Укажіть ТИП значення виразу.\par
\begin{verbatim}
{21,2,4}
\end{verbatim}
%enditem


\newpage
6. Що буде надруковано за виконання?\par
\begin{verbatim}
try:
    d = {16:8, 18:8, 25:5, 20:2, 16:8}
    d[16] = 8
    for x in d.keys():
        print(x, sep=' ', end=' ')
except:
    print('ERR')
\end{verbatim}
%enditem


\newpage
7. Нехай змінна \kw{s} позначає послідовність цілих, по якій можна ітерувати. 

Написати вираз, що перевіряє, чи правда, що послідовність  містить хоча б одне додатнє число.
%enditem


\newpage
8. Що буде надруковано за виконання?
\begin{verbatim}
b = 0
m = 1

class D:
    m = 2
    
    def __init__(self):
        self.m = 3
        m = 4
        D.m = 5

obj = D()
try:
    print(b, end=' ')
    print(m, end=' ')
    print(D.m, end=' ')
    print(obj.m)
except:
    print('ERR')
\end{verbatim}
%enditem


\newpage
9. Що буде надруковано за виконання?
\begin{verbatim}
class C0:
    def d(self): print(1, end=' ')
    def h(self): print(2, end=' ')

class C1(C0):
    def h(self): print(3, end=' ')

try:
    C1().d()
    C0().h()
except:
    print('ERR')
\end{verbatim}
%enditem


\newpage
10. Що буде надруковано за виконання?
\begin{verbatim}
class A:
    b = 1

    def __init__(self):
        a = 2
        a = 3

class B(A):
    b = 4
    c = 5

    def __init__(self):
        super().__init__()
        self.c = 6

try:
    obj = B()
    obj.b = obj.a + 10
    B.b = 13
except: print('EE', end=' ')

try: print(obj.b, end= ' ')
except: print('E1', end=' ')

try: print(A.b, end= ' ')
except: print('E2', end=' ')

try: print(B.b, end= ' ')
except: print('E3', end=' ')
\end{verbatim}
%enditem


\newpage
11. Змінні \kw{next} та \kw{prev} вузла списку позначають наступний та попередній за ним вузол відповідно. Починаючи з голови, у список записано послідовні цілі числа від~$-45$ до~$41$. Нехай змінна \kw{p} позначає вузол, що зберігає значення~$3$.
Яке значення зберігається у вузлі \kw{p.prev.prev.next.prev.next} ?
%enditem


\newpage
12. 
\begin{center}
	\adjustimage{max size={0.8\linewidth}{0.8\paperheight}}
	{BinTree/trees_exp/137.png}
\end{center}
Указати послідовність міток вершин дерева, зображеного на рис., що утвориться в результаті обходу дерева в оберненому порядку. Мітки записувати через пробіл.\par Алгоритм починає роботу з кореня (на рис. це верхня вершина).
%enditem


\newpage
13. 
\begin{center}
	\adjustimage{max size={0.9\linewidth}{0.9\paperheight}}
	{Graphs/AdjList/86.png}
\end{center}
Для графа, зображеного на рис., вказати список суміжності вершини 9.\par
Формат оформлення відповіді:\par
список суміжності виписувати через пробіл на одному рядку, якщо пробілів десь буде більше одного - не важливо, упорядкування вершин у списку також значення не має.
%enditem


\newpage
14. Для графа на рисунку вказати послідовність, в якій алгоритм пошуку в глибину буде відмічати відвідані вершини, якщо списки суміжності вершин переглядаються алгоритмом за зростанням міток вершин та вершини відмічаються на першому вход\-женні у вершину.\par
Пошук розпочинається з вершини 4.\par

\begin{center}
	\adjustimage{max size={0.9\linewidth}{0.9\paperheight}}
	{Graphs/Traversals/217.png}
\end{center}
%enditem


\newpage
15. Для графа на рисунку з вершини 0 запущено алгоритм пошуку в ширину, що будує кістякове дерево. Списки суміжності вершин переглядаються алгоритмом за зростанням міток вершин.\par
\begin{center}
	\adjustimage{max size={0.9\linewidth}{0.9\paperheight}}
	{Graphs/Traversals/77.png}
\end{center}
\par Які з наведених ребер входять до складу отриманого кістякового дерева?\par
1) (0, 1)\par
2) (3, 5)\par
3) (4, 6)\par
4) (1, 6)\par
5) (1, 5)\par
6) (0, 2)\par
{\vskip 15pt} Номери відповідних ребер записати за зростанням через пробіл.


%enditem


\newpage
16. 

Рядки журналу через двокрапку містять ім'я користувача (ідентифікатор), час події,тип події, код події, наприклад
\begin{verbatim}
s="username : 2022-05-05 13:26 : ERROR : 404"
\end{verbatim}
у коді 
\begin{verbatim}
res = re.fullmatch('???', s) 
s = ''
code_str = ??? 
\end{verbatim}
чим треба замінити кожен з ???, щоб значенням code\_str був код події? 


\vspace{10mm}
\textbf{Завдання 17} 

Задано програму, що містить код 1-2 класів, можливе успадкування, та клієнтський код.
У наведеному коді невеликі частини закриті. 
Також наведено вихід програми.
Необхідно відновити закриті частини, щоб отриманий код відповідав наведеному виходу програми.


\textbf{Завдання 18-19} 

Кілька запитань за завданнями для самостійної роботи.
\vspace{10mm}
Скоріш за все, що завдань буде трохи менше, а деякі попередні будуть переформульовані в стилі двох останніх.

%enditem

\end {large}}

%end

\newpage
%begin

{{
\begin {large}

\newpage
1. Що буде надруковано за виконання?

\begin{verbatim}
print('R', end=' ')
d=('a','b','c','d','e',0,1,2)
print(d[-14])
\end{verbatim}

%enditem


\newpage
2. Що буде надруковано за виконання?

\begin{verbatim}
print('R', end=' ')
b=[0,1,2,3,4,5,6,7,8,9]
print(b[1::-2])
\end{verbatim}

%enditem



\newpage
3. Що буде надруковано за виконання?

\begin{verbatim}
print('R', end=' ')
e=['0','1','2','3','4']
print(e.pop(0), end=' ')
print(e)
\end{verbatim}

%enditem



\newpage
4. Що буде надруковано за виконання?

\begin{verbatim}
print('R', end=' ')
e=[10,11,12,13,14]
e[-2:4]='11'
print(e)
\end{verbatim}

%enditem



\newpage
5. Укажіть ТИП значення виразу.\par
\begin{verbatim}
{(1,4), 8}
\end{verbatim}
%enditem


\newpage
6. Що буде надруковано за виконання?\par
\begin{verbatim}
try:
    s = {46:7, 28:6, 88:9, 88:9}
    s[14] = 2
    for x in s.keys():
        print(x, sep=' ', end=' ')
except:
    print('ERR')
\end{verbatim}
%enditem


\newpage
7. Нехай змінна \kw{s} позначає дуже довгу послідовність цілих, по якій можна ітерувати. 

Написати вираз-генератор, що задає підпослідовність її елементів, що менші за задане число num. 

Обмеження: усі елементи шуканої підпослідовності одночасно в пам'яті розміщені бути не можуть. 
%enditem


\newpage
8. Що буде надруковано за виконання?
%Оригинал был утерян. Требует отладки!
\begin{verbatim}
b = 0

class B:
    b = 1
    
    def __init__(self):
        self.b = 2
        b = 3
        B.b = 4

obj = B()
obj.b = 5
try:
    print(b, end=' ')
    print(obj.b, end=' ')
    print(B.b, end=' ')
    print(obj.b, end=' ')
    print(obj._B__b)
except:
    print('ERR')
\end{verbatim}
%enditem


\newpage
9. Що буде надруковано за виконання?
\begin{verbatim}
class B0:
    def c(self): print(1, end=' ')
    def h(self): print(2, end=' ')

class B1(B0):
    def h(self): print(3, end=' ')

try:
    B1().c()
    B1().h()
except:
    print('ERR')
\end{verbatim}
%enditem


\newpage
10. Що буде надруковано за виконання?
\begin{verbatim}
class A:
    a = 1

    def __g1(self): print(2, end=' ')
    def _g2(self): print(3, end=' ')

    def main(self):
        try: print(self.a, end=' ')
        except: print('E1', end=' ')

        try: self.__g1()
        except: print('E2', end=' ')

        try: self._g2()
        except: print('E3', end=' ')


class B(A):
    a = 4

    def __g1(self): print(5, end=' ')
    def _g2(self): print(6, end=' ')

try: B().main()
except: print('E4', end=' ')
\end{verbatim}
%enditem


\newpage
11. Змінні \kw{next} та \kw{prev} вузла списку позначають наступний та попередній за ним вузол відповідно. Починаючи з голови, у список записано послідовні цілі числа від~$-17$ до~$23$. Нехай змінна \kw{p} позначає вузол, що зберігає значення~$9$.
Яке значення зберігається у вузлі \kw{p.next.prev.prev.next.prev} ?
%enditem


\newpage
12. 
\begin{center}
	\adjustimage{max size={0.8\linewidth}{0.8\paperheight}}
	{BinTree/trees_exp/167.png}
\end{center}
Указати послідовність міток вершин дерева, зображеного на рис., що утвориться в результаті обходу дерева в оберненому порядку. Мітки записувати через пробіл.\par Алгоритм починає роботу з кореня (на рис. це верхня вершина).
%enditem


\newpage
13. 
\begin{center}
	\adjustimage{max size={0.9\linewidth}{0.9\paperheight}}
	{Graphs/AdjList/36.png}
\end{center}
Для графа, зображеного на рис., вказати список суміжності вершини 1.\par
Формат оформлення відповіді:\par
список суміжності виписувати через пробіл на одному рядку, якщо пробілів десь буде більше одного - не важливо, упорядкування вершин у списку також значення не має.
%enditem


\newpage
14. Для графа на рисунку з вершини 1 запущено алгоритм пошуку в глибину, що будує кістякове дерево. Списки суміжності вершин переглядаються алгоритмом за зростанням міток вершин.\par
\begin{center}
	\adjustimage{max size={0.9\linewidth}{0.9\paperheight}}
	{Graphs/Traversals/70.png}
\end{center}
\par Які з наведених ребер входять до складу отриманого кістякового дерева?\par
1) (1, 3)\par
2) (2, 6)\par
3) (0, 6)\par
4) (2, 3)\par
5) (2, 4)\par
6) (3, 4)\par
{\vskip 15pt} Номери відповідних ребер записати за зростанням через пробіл.

%enditem


\newpage
15. Для графа на рисунку з вершини 5 запущено алгоритм пошуку в ширину, що будує кістякове дерево. Списки суміжності вершин переглядаються алгоритмом за зростанням міток вершин.\par
\begin{center}
	\adjustimage{max size={0.9\linewidth}{0.9\paperheight}}
	{Graphs/Traversals/276.png}
\end{center}
\par Які з наведених ребер входять до складу отриманого кістякового дерева?\par
1) (0, 5)\par
2) (4, 6)\par
3) (1, 5)\par
4) (1, 6)\par
5) (0, 6)\par
6) (2, 3)\par
{\vskip 15pt} Номери відповідних ребер записати за зростанням через пробіл.


%enditem


\newpage
16. 

Рядки журналу через двокрапку містять ім'я користувача (ідентифікатор), час події,тип події, код події, наприклад
\begin{verbatim}
s="username : 2022-05-05 13:26 : ERROR : 404"
\end{verbatim}
у коді 
\begin{verbatim}
res = re.fullmatch('???', s) 
s = ''
code_str = ??? 
\end{verbatim}
чим треба замінити кожен з ???, щоб значенням code\_str був код події? 


\vspace{10mm}
\textbf{Завдання 17} 

Задано програму, що містить код 1-2 класів, можливе успадкування, та клієнтський код.
У наведеному коді невеликі частини закриті. 
Також наведено вихід програми.
Необхідно відновити закриті частини, щоб отриманий код відповідав наведеному виходу програми.


\textbf{Завдання 18-19} 

Кілька запитань за завданнями для самостійної роботи.
\vspace{10mm}
Скоріш за все, що завдань буде трохи менше, а деякі попередні будуть переформульовані в стилі двох останніх.

%enditem

\end {large}}

%end

\newpage
%begin

{{
\begin {large}

\newpage
1. Що буде надруковано за виконання?

\begin{verbatim}
print('R', end=' ')
g=('a','b','c','d','e',0,1)
print(g[-8])
\end{verbatim}

%enditem


\newpage
2. Що буде надруковано за виконання?

\begin{verbatim}
print('R', end=' ')
h='0123456789'
print(h[-4:2:2])
\end{verbatim}

%enditem



\newpage
3. Що буде надруковано за виконання?

\begin{verbatim}
print('R', end=' ')
h=[10,11,12,13,14]
print(h.pop(3), end=' ')
print(h)
\end{verbatim}

%enditem



\newpage
4. Що буде надруковано за виконання?

\begin{verbatim}
print('R', end=' ')
h=['0','1','2','3','4']
h[:0]=[15]
print(h)
\end{verbatim}

%enditem



\newpage
5. Укажіть ТИП значення виразу.\par
\begin{verbatim}
{x for x in range(7)}
\end{verbatim}
%enditem


\newpage
6. Що буде надруковано за виконання?\par
\begin{verbatim}
try:
    t = {30:1, 47:2, 82:6, 47:9}
    t[73] = 9
    for x in t.items():
        print(*x, sep=' ', end=' ')
except:
    print('ERR')
\end{verbatim}
%enditem


\newpage
7. 
Записати ВИРАЗ, значенням якого є список, що складається з цілих чисел від 12 до 2004 (включно), причому спочатку за спаданням записано всі, що не діляться на 5, а потім решту за зростанням.

%enditem


\newpage
8. Що буде надруковано за виконання?
%Оригинал был утерян. Требует отладки!
\begin{verbatim}
__c = 0

class C:
    __c = 1
    
    def __init__(self):
        self.__c = 2
        __c = 3
        C.__c = 4

obj = C()
C.__c = 5
try:
    print(__c, end=' ')
    print(obj.__c, end=' ')
    print(C.__c, end=' ')
    print(obj.__c, end=' ')
    print(obj._C__c)
except:
    print('ERR')
\end{verbatim}
%enditem


\newpage
9. Що буде надруковано за виконання?
\begin{verbatim}
class A0:
    def b(self): print(1, end=' ')
    def f(self): print(2, end=' ')

class A1(A0):
    def f(self): print(3, end=' ')

try:
    A1().b()
    A1().f()
except:
    print('ERR')
\end{verbatim}
%enditem


\newpage
10. Що буде надруковано за виконання?
\begin{verbatim}
class A:
    a = 1

    def __init__(self):
        c = 2
        c = 3

class B(A):
    c = 4
    b = 5

    def __init__(self):
        super().__init__()
        self.b = 6

try:
    obj = B()
    obj.c = obj.c + 10
    B.c = 13
except: print('EE', end=' ')

try: print(obj.c, end= ' ')
except: print('E1', end=' ')

try: print(A.c, end= ' ')
except: print('E2', end=' ')

try: print(B.c, end= ' ')
except: print('E3', end=' ')
\end{verbatim}
%enditem


\newpage
11. Змінні \kw{next} та \kw{prev} вузла списку позначають наступний та попередній за ним вузол відповідно. Починаючи з голови, у список записано послідовні цілі числа від~$-20$ до~$15$. Нехай змінна \kw{p} позначає вузол, що зберігає значення~$4$.
Яке значення зберігається у вузлі \kw{p.prev.next.next.prev.next} ?
%enditem


\newpage
12. 
\begin{center}
	\adjustimage{max size={0.8\linewidth}{0.8\paperheight}}
	{BinTree/trees_exp/163.png}
\end{center}
Указати послідовність міток вершин дерева, зображеного на рис., що утвориться в результаті обходу дерева в прямому порядку. Мітки записувати через пробіл.\par Алгоритм починає роботу з кореня (на рис. це верхня вершина).
%enditem


\newpage
13. 
\begin{center}
	\adjustimage{max size={0.9\linewidth}{0.9\paperheight}}
	{Graphs/AdjList/240.png}
\end{center}
Для графа, зображеного на рис., вказати список суміжності вершини 0.\par
Формат оформлення відповіді:\par
список суміжності виписувати через пробіл на одному рядку, якщо пробілів десь буде більше одного - не важливо, упорядкування вершин у списку також значення не має.
%enditem


\newpage
14. Для графа на рисунку вказати послідовність, в якій алгоритм пошуку в глибину буде відмічати відвідані вершини, якщо списки суміжності вершин переглядаються алгоритмом за зростанням міток вершин та вершини відмічаються на першому вход\-женні у вершину.\par
Пошук розпочинається з вершини 3.\par

\begin{center}
	\adjustimage{max size={0.9\linewidth}{0.9\paperheight}}
	{Graphs/Traversals/35.png}
\end{center}
%enditem


\newpage
15. Для графа на рисунку з вершини 5 запущено алгоритм пошуку в ширину, що будує кістякове дерево. Списки суміжності вершин переглядаються алгоритмом за зростанням міток вершин.\par
\begin{center}
	\adjustimage{max size={0.9\linewidth}{0.9\paperheight}}
	{Graphs/Traversals/205.png}
\end{center}
\par Які з наведених ребер входять до складу отриманого кістякового дерева?\par
1) (0, 2)\par
2) (0, 3)\par
3) (0, 5)\par
4) (1, 4)\par
5) (2, 5)\par
6) (3, 5)\par
{\vskip 15pt} Номери відповідних ребер записати за зростанням через пробіл.


%enditem


\newpage
16. 

Рядки журналу через двокрапку містять ім'я користувача (ідентифікатор), час події,тип події, код події, наприклад
\begin{verbatim}
s="username : 2022-05-05 13:26 : ERROR : 404"
\end{verbatim}
у коді 
\begin{verbatim}
res = re.fullmatch('???', s) 
s = ''
code_str = ??? 
\end{verbatim}
чим треба замінити кожен з ???, щоб значенням code\_str був код події? 


\vspace{10mm}
\textbf{Завдання 17} 

Задано програму, що містить код 1-2 класів, можливе успадкування, та клієнтський код.
У наведеному коді невеликі частини закриті. 
Також наведено вихід програми.
Необхідно відновити закриті частини, щоб отриманий код відповідав наведеному виходу програми.


\textbf{Завдання 18-19} 

Кілька запитань за завданнями для самостійної роботи.
\vspace{10mm}
Скоріш за все, що завдань буде трохи менше, а деякі попередні будуть переформульовані в стилі двох останніх.

%enditem

\end {large}}

%end

\newpage
%begin

{{
\begin {large}

\newpage
1. Що буде надруковано за виконання?

\begin{verbatim}
print('R', end=' ')
h=('a','b','c','d','e',0,1,2,3,4)
print(h[-1])
\end{verbatim}

%enditem


\newpage
2. Що буде надруковано за виконання?

\begin{verbatim}
print('R', end=' ')
f=(0,1,2,3,4,5,6,7,8,9)
print(f[-7:-15:-3])
\end{verbatim}

%enditem



\newpage
3. Що буде надруковано за виконання?

\begin{verbatim}
print('R', end=' ')
a=['0','1','2','3','4']
a.remove('7')
print(a)
\end{verbatim}

%enditem



\newpage
4. Що буде надруковано за виконання?

\begin{verbatim}
print('R', end=' ')
a=[10,11,12,13,14]
del a[6:-5]
print(a)
\end{verbatim}

%enditem



\newpage
5. Укажіть ТИП значення виразу.\par
\begin{verbatim}
{x*x:x for x in range(-7,7)}
\end{verbatim}
%enditem


\newpage
6. Що буде надруковано за виконання?\par
\begin{verbatim}
try:
    d = {10:2, 32:7, 74:0, 32:8}
    d[20] = 7
    for x in d.values():
        print(x, sep=' ', end=' ')
except:
    print('ERR')
\end{verbatim}
%enditem


\newpage
7. Нехай змінна \kw{s} позначає дуже довгу послідовність цілих, по якій можна ітерувати. 

Написати вираз-генератор, що задає підпослідовність її елементів, що діляться на 3. 

Обмеження: усі елементи шуканої підпослідовності одночасно в пам'яті розміщені бути не можуть. 
%enditem


\newpage
8. Що буде надруковано за виконання?
\begin{verbatim}
a = 0
j = 1

class D:
    a = 2
    
    def __init__(self):
        global a
        self.j = 3
        a = 4
        D.a = 5

obj = D()
try:
    print(a, end=' ')
    print(j, end=' ')
    print(D.a, end=' ')
    print(obj.j)
except:
    print('ERR')
\end{verbatim}
%enditem


\newpage
9. Що буде надруковано за виконання?
\begin{verbatim}
class D0:
    def c(self): print(1, end=' ')
    def g(self): print(2, end=' ')

class D1(D0):
    def c(self): print(3, end=' ')

try:
    D1().g()
    D1().c()
except:
    print('ERR')
\end{verbatim}
%enditem


\newpage
10. Що буде надруковано за виконання?
\begin{verbatim}
class A:
    a = 1

    def __init__(self):
        b = 2
        a = 3

class B(A):
    a = 4
    c = 5

    def __init__(self):
        super().__init__()
        self.c = 6

try:
    obj = B()
    obj.b = obj.a + 10
    B.b = 13
except: print('EE', end=' ')

try: print(obj.b, end= ' ')
except: print('E1', end=' ')

try: print(A.b, end= ' ')
except: print('E2', end=' ')

try: print(B.b, end= ' ')
except: print('E3', end=' ')
\end{verbatim}
%enditem


\newpage
11. Змінні \kw{next} та \kw{prev} вузла списку позначають наступний та попередній за ним вузол відповідно. Починаючи з голови, у список записано послідовні цілі числа від~$-63$ до~$20$. Нехай змінна \kw{p} позначає вузол, що зберігає значення~$-4$.
Яке значення зберігається у вузлі \kw{p.next.prev.next.next.prev} ?
%enditem


\newpage
12. 
\begin{center}
	\adjustimage{max size={0.8\linewidth}{0.8\paperheight}}
	{BinTree/trees_exp/124.png}
\end{center}
Указати послідовність міток вершин дерева, зображеного на рис., що утвориться в результаті обходу дерева в оберненому порядку. Мітки записувати через пробіл.\par Алгоритм починає роботу з кореня (на рис. це верхня вершина).
%enditem


\newpage
13. 
\begin{center}
	\adjustimage{max size={0.9\linewidth}{0.9\paperheight}}
	{Graphs/AdjList/45.png}
\end{center}
Для графа, зображеного на рис., вказати список суміжності вершини 7.\par
Формат оформлення відповіді:\par
список суміжності виписувати через пробіл на одному рядку, якщо пробілів десь буде більше одного - не важливо, упорядкування вершин у списку також значення не має.
%enditem


\newpage
14. Для графа на рисунку вказати послідовність, в якій алгоритм пошуку в глибину буде відмічати відвідані вершини, якщо списки суміжності вершин переглядаються алгоритмом за зростанням міток вершин та вершини відмічаються на першому вход\-женні у вершину.\par
Пошук розпочинається з вершини 0.\par

\begin{center}
	\adjustimage{max size={0.9\linewidth}{0.9\paperheight}}
	{Graphs/Traversals/75.png}
\end{center}
%enditem


\newpage
15. Для графа на рисунку з вершини 5 запущено алгоритм пошуку в ширину, що будує кістякове дерево. Списки суміжності вершин переглядаються алгоритмом за зростанням міток вершин.\par
\begin{center}
	\adjustimage{max size={0.9\linewidth}{0.9\paperheight}}
	{Graphs/Traversals/257.png}
\end{center}
\par Які з наведених ребер входять до складу отриманого кістякового дерева?\par
1) (1, 4)\par
2) (0, 5)\par
3) (3, 5)\par
4) (1, 3)\par
5) (0, 1)\par
6) (5, 6)\par
{\vskip 15pt} Номери відповідних ребер записати за зростанням через пробіл.


%enditem


\newpage
16. 

Рядки журналу через двокрапку містять ім'я користувача (ідентифікатор), час події,тип події, код події, наприклад
\begin{verbatim}
s="username : 2022-05-05 13:26 : ERROR : 404"
\end{verbatim}
у коді 
\begin{verbatim}
res = re.fullmatch('???', s) 
s = ''
code_str = ??? 
\end{verbatim}
чим треба замінити кожен з ???, щоб значенням code\_str був код події? 


\vspace{10mm}
\textbf{Завдання 17} 

Задано програму, що містить код 1-2 класів, можливе успадкування, та клієнтський код.
У наведеному коді невеликі частини закриті. 
Також наведено вихід програми.
Необхідно відновити закриті частини, щоб отриманий код відповідав наведеному виходу програми.


\textbf{Завдання 18-19} 

Кілька запитань за завданнями для самостійної роботи.
\vspace{10mm}
Скоріш за все, що завдань буде трохи менше, а деякі попередні будуть переформульовані в стилі двох останніх.

%enditem

\end {large}}

%end

\newpage
%begin

{{
\begin {large}

\newpage
1. Що буде надруковано за виконання?

\begin{verbatim}
print('R', end=' ')
d=('a','b','c','d','e',0,1,2,3)
print(d[0])
\end{verbatim}

%enditem


\newpage
2. Що буде надруковано за виконання?

\begin{verbatim}
print('R', end=' ')
d=('0','1','2','3','4','5','6','7','8','9')
print(d[9:3:3])
\end{verbatim}

%enditem



\newpage
3. Що буде надруковано за виконання?

\begin{verbatim}
print('R', end=' ')
c=[10,11,12,13,14]
c.insert(-4, '11')
print(c)
\end{verbatim}

%enditem



\newpage
4. Що буде надруковано за виконання?

\begin{verbatim}
print('R', end=' ')
h=['0','1','2','3','4']
h[-5:7]=[1,2]
print(h)
\end{verbatim}

%enditem



\newpage
5. Укажіть ТИП значення виразу.\par
\begin{verbatim}
{y:13 for y in range(-7,7)}
\end{verbatim}
%enditem


\newpage
6. Що буде надруковано за виконання?\par
\begin{verbatim}
try:
    t = {48:0, 71:8, 21:5, 25:5}
    t[71] = 1
    for x in t.items():
        print(x, sep=' ', end=' ')
except:
    print('ERR')
\end{verbatim}
%enditem


\newpage
7. Нехай змінна \kw{s} позначає послідовність цілих, по якій можна ітерувати. 

Написати вираз, що перевіряє, чи правда, що послідовність  містить числа, більші за 5.
%enditem


\newpage
8. Що буде надруковано за виконання?
%Оригинал был утерян. Требует отладки!
\begin{verbatim}
d = 0

class C:
    d = 1
    
    def __init__(self):
        self.d = 2
        d = 3
        C.d = 4

obj = C()
obj.d = 5
try:
    print(d, end=' ')
    print(obj.d, end=' ')
    print(C.d, end=' ')
    print(obj.d, end=' ')
    print(obj._C__d)
except:
    print('ERR')
\end{verbatim}
%enditem


\newpage
9. Що буде надруковано за виконання?
\begin{verbatim}
class C0:
    def a(self): print(1, end=' ')
    def g(self): print(2, end=' ')

class C1(C0):
    def g(self): print(3, end=' ')

try:
    C1().a()
    C1().g()
except:
    print('ERR')
\end{verbatim}
%enditem


\newpage
10. Що буде надруковано за виконання?
\begin{verbatim}
class A:
    d = 1

    def _h1(self): print(2, end=' ')
    def __h2(self): print(3, end=' ')

    def main(self):
        try: print(self.d, end=' ')
        except: print('E1', end=' ')

        try: self._h1()
        except: print('E2', end=' ')

        try: self.__h2()
        except: print('E3', end=' ')


class B(A):
    d = 4

    def _h1(self): print(5, end=' ')
    def __h2(self): print(6, end=' ')

try: B().main()
except: print('E4', end=' ')
\end{verbatim}
%enditem


\newpage
11. Змінні \kw{next} та \kw{prev} вузла списку позначають наступний та попередній за ним вузол відповідно. Починаючи з голови, у список записано послідовні цілі числа від~$-35$ до~$18$. Нехай змінна \kw{p} позначає вузол, що зберігає значення~$8$.
Яке значення зберігається у вузлі \kw{p.prev.next.prev.prev.next} ?
%enditem


\newpage
12. 
\begin{center}
	\adjustimage{max size={0.8\linewidth}{0.8\paperheight}}
	{BinTree/trees_exp/36.png}
\end{center}
Указати послідовність міток вершин дерева, зображеного на рис., що утвориться в результаті обходу дерева в прямому порядку. Мітки записувати через пробіл.\par Алгоритм починає роботу з кореня (на рис. це верхня вершина).
%enditem


\newpage
13. 
\begin{center}
	\adjustimage{max size={0.9\linewidth}{0.9\paperheight}}
	{Graphs/AdjList/245.png}
\end{center}
Для графа, зображеного на рис., вказати список суміжності вершини 3.\par
Формат оформлення відповіді:\par
список суміжності виписувати через пробіл на одному рядку, якщо пробілів десь буде більше одного - не важливо, упорядкування вершин у списку також значення не має.
%enditem


\newpage
14. Для графа на рисунку з вершини 0 запущено алгоритм пошуку в глибину, що будує кістякове дерево. Списки суміжності вершин переглядаються алгоритмом за зростанням міток вершин.\par
\begin{center}
	\adjustimage{max size={0.9\linewidth}{0.9\paperheight}}
	{Graphs/Traversals/65.png}
\end{center}
\par Які з наведених ребер входять до складу отриманого кістякового дерева?\par
1) (0, 6)\par
2) (2, 3)\par
3) (4, 6)\par
4) (0, 4)\par
5) (4, 5)\par
6) (3, 5)\par
{\vskip 15pt} Номери відповідних ребер записати за зростанням через пробіл.

%enditem


\newpage
15. Для графа на рисунку з вершини 4 запущено алгоритм пошуку в ширину, що будує кістякове дерево. Списки суміжності вершин переглядаються алгоритмом за зростанням міток вершин.\par
\begin{center}
	\adjustimage{max size={0.9\linewidth}{0.9\paperheight}}
	{Graphs/Traversals/121.png}
\end{center}
\par Які з наведених ребер входять до складу отриманого кістякового дерева?\par
1) (0, 3)\par
2) (1, 3)\par
3) (4, 6)\par
4) (2, 4)\par
5) (0, 2)\par
6) (0, 4)\par
{\vskip 15pt} Номери відповідних ребер записати за зростанням через пробіл.


%enditem


\newpage
16. 

Рядки журналу через двокрапку містять ім'я користувача (ідентифікатор), час події,тип події, код події, наприклад
\begin{verbatim}
s="username : 2022-05-05 13:26 : ERROR : 404"
\end{verbatim}
у коді 
\begin{verbatim}
res = re.fullmatch('???', s) 
s = ''
code_str = ??? 
\end{verbatim}
чим треба замінити кожен з ???, щоб значенням code\_str був код події? 


\vspace{10mm}
\textbf{Завдання 17} 

Задано програму, що містить код 1-2 класів, можливе успадкування, та клієнтський код.
У наведеному коді невеликі частини закриті. 
Також наведено вихід програми.
Необхідно відновити закриті частини, щоб отриманий код відповідав наведеному виходу програми.


\textbf{Завдання 18-19} 

Кілька запитань за завданнями для самостійної роботи.
\vspace{10mm}
Скоріш за все, що завдань буде трохи менше, а деякі попередні будуть переформульовані в стилі двох останніх.

%enditem

\end {large}}

%end

\newpage
%begin

{{
\begin {large}

\newpage
1. Що буде надруковано за виконання?

\begin{verbatim}
print('R', end=' ')
g=('a','b','c','d','e',0,1,2,4)
print(g[-10])
\end{verbatim}

%enditem


\newpage
2. Що буде надруковано за виконання?

\begin{verbatim}
print('R', end=' ')
b=[0,1,2,3,4,5,6,7,8,9]
print(b[::-2])
\end{verbatim}

%enditem



\newpage
3. Що буде надруковано за виконання?

\begin{verbatim}
print('R', end=' ')
b=['0','1','2','3','4']
b.extend('13')
print(b)
\end{verbatim}

%enditem



\newpage
4. Що буде надруковано за виконання?

\begin{verbatim}
print('R', end=' ')
c=['0','1','2','3','4']
del c[-3:]
print(c)
\end{verbatim}

%enditem



\newpage
5. Укажіть ТИП значення виразу.\par
\begin{verbatim}
{1:3, 5:8}
\end{verbatim}
%enditem


\newpage
6. Що буде надруковано за виконання?\par
\begin{verbatim}
try:
    s = {58:3, 69:8, 78:4, 58:6}
    s[85] = 1
    for x, y in s.items():
        print(x, y, sep=' ', end=' ')
except:
    print('ERR')
\end{verbatim}
%enditem


\newpage
7. Нехай змінна \kw{s} позначає дуже довгу послідовність, по якій можна ітерувати. 

Написати вираз-генератор, що задає підпослідовність її елементів, що не є числами типу \kw{int}.

Обмеження: усі елементи шуканої підпослідовності одночасно в пам'яті розміщені бути не можуть. 
%enditem


\newpage
8. Що буде надруковано за виконання?
\begin{verbatim}
c = 0
k = 1

class A:
    k = 2
    
    def __init__(self):
        self.c = 3
        k = 4
        A.k = 5

obj = A()
try:
    print(c, end=' ')
    print(k, end=' ')
    print(A.k, end=' ')
    print(obj.c)
except:
    print('ERR')
\end{verbatim}
%enditem


\newpage
9. Що буде надруковано за виконання?
\begin{verbatim}
class B0:
    def d(self): print(1, end=' ')
    def f(self): print(2, end=' ')

class B1(B0):
    def d(self): print(3, end=' ')

try:
    B1().f()
    B0().d()
except:
    print('ERR')
\end{verbatim}
%enditem


\newpage
10. Що буде надруковано за виконання?
\begin{verbatim}
class A:
    d = 1

    def f1(self): print(2, end=' ')
    def _f2(self): print(3, end=' ')

    def main(self):
        try: print(self.d, end=' ')
        except: print('E1', end=' ')

        try: self.f1()
        except: print('E2', end=' ')

        try: self._f2()
        except: print('E3', end=' ')


class B(A):
    d = 4

    def f1(self): print(5, end=' ')
    def _f2(self): print(6, end=' ')

try: B().main()
except: print('E4', end=' ')
\end{verbatim}
%enditem


\newpage
11. Змінні \kw{next} та \kw{prev} вузла списку позначають наступний та попередній за ним вузол відповідно. Починаючи з голови, у список записано послідовні цілі числа від~$-44$ до~$33$. Нехай змінна \kw{p} позначає вузол, що зберігає значення~$2$.
Яке значення зберігається у вузлі \kw{p.next.next.next.prev.next} ?
%enditem


\newpage
12. 
\begin{center}
	\adjustimage{max size={0.8\linewidth}{0.8\paperheight}}
	{BinTree/trees_exp/3.png}
\end{center}
Указати послідовність міток вершин дерева, зображеного на рис., що утвориться в результаті обходу дерева в оберненому порядку. Мітки записувати через пробіл.\par Алгоритм починає роботу з кореня (на рис. це верхня вершина).
%enditem


\newpage
13. 
\begin{center}
	\adjustimage{max size={0.9\linewidth}{0.9\paperheight}}
	{Graphs/AdjList/121.png}
\end{center}
Для графа, зображеного на рис., вказати список суміжності вершини 7.\par
Формат оформлення відповіді:\par
список суміжності виписувати через пробіл на одному рядку, якщо пробілів десь буде більше одного - не важливо, упорядкування вершин у списку також значення не має.
%enditem


\newpage
14. Для графа на рисунку з вершини 5 запущено алгоритм пошуку в глибину, що будує кістякове дерево. Списки суміжності вершин переглядаються алгоритмом за зростанням міток вершин.\par
\begin{center}
	\adjustimage{max size={0.9\linewidth}{0.9\paperheight}}
	{Graphs/Traversals/224.png}
\end{center}
\par Які з наведених ребер входять до складу отриманого кістякового дерева?\par
1) (1, 2)\par
2) (0, 5)\par
3) (2, 6)\par
4) (2, 3)\par
5) (0, 2)\par
6) (5, 6)\par
{\vskip 15pt} Номери відповідних ребер записати за зростанням через пробіл.

%enditem


\newpage
15. Для графа на рисунку з вершини 6 запущено алгоритм пошуку в ширину, що будує кістякове дерево. Списки суміжності вершин переглядаються алгоритмом за зростанням міток вершин.\par
\begin{center}
	\adjustimage{max size={0.9\linewidth}{0.9\paperheight}}
	{Graphs/Traversals/286.png}
\end{center}
\par Які з наведених ребер входять до складу отриманого кістякового дерева?\par
1) (0, 1)\par
2) (0, 5)\par
3) (2, 5)\par
4) (4, 6)\par
5) (1, 6)\par
6) (1, 3)\par
{\vskip 15pt} Номери відповідних ребер записати за зростанням через пробіл.


%enditem


\newpage
16. 

Рядки журналу через двокрапку містять ім'я користувача (ідентифікатор), час події,тип події, код події, наприклад
\begin{verbatim}
s="username : 2022-05-05 13:26 : ERROR : 404"
\end{verbatim}
у коді 
\begin{verbatim}
res = re.fullmatch('???', s) 
s = ''
code_str = ??? 
\end{verbatim}
чим треба замінити кожен з ???, щоб значенням code\_str був код події? 


\vspace{10mm}
\textbf{Завдання 17} 

Задано програму, що містить код 1-2 класів, можливе успадкування, та клієнтський код.
У наведеному коді невеликі частини закриті. 
Також наведено вихід програми.
Необхідно відновити закриті частини, щоб отриманий код відповідав наведеному виходу програми.


\textbf{Завдання 18-19} 

Кілька запитань за завданнями для самостійної роботи.
\vspace{10mm}
Скоріш за все, що завдань буде трохи менше, а деякі попередні будуть переформульовані в стилі двох останніх.

%enditem

\end {large}}

%end

\newpage
%begin

{{
\begin {large}

\newpage
1. Що буде надруковано за виконання?

\begin{verbatim}
print('R', end=' ')
e=('a','b','c','d','e',0,1,2)
print(e[-8])
\end{verbatim}

%enditem


\newpage
2. Що буде надруковано за виконання?

\begin{verbatim}
print('R', end=' ')
c=['0','1','2','3','4','5','6','7','8','9']
print(c[:-6:-1])
\end{verbatim}

%enditem



\newpage
3. Що буде надруковано за виконання?

\begin{verbatim}
print('R', end=' ')
g=[10,11,12,13,14]
print(g.index(17), end=' ')
\end{verbatim}

%enditem



\newpage
4. Що буде надруковано за виконання?

\begin{verbatim}
print('R', end=' ')
e=[10,11,12,13,14]
e[:1]='10'
print(e)
\end{verbatim}

%enditem



\newpage
5. Укажіть ТИП значення виразу.\par
\begin{verbatim}
[13,2,4]+[3,5]
\end{verbatim}
%enditem


\newpage
6. Що буде надруковано за виконання?\par
\begin{verbatim}
try:
    d = {87:4, 47:6, 88:4, 87:5}
    d[33] = 2
    for x in d :
        print(x, sep=' ', end=' ')
except:
    print('ERR')
\end{verbatim}
%enditem


\newpage
7. 
Записати ВИРАЗ, значенням якого є список, що складається з кубів цілих чисел від 12 до 2004 (включно), що діляться на 7.
%enditem


\newpage
8. Що буде надруковано за виконання?
%Оригинал был утерян. Требует отладки!
\begin{verbatim}
_a = 0

class A:
    _a = 1
    
    def __init__(self):
        self._a = 2
        _a = 3
        A._a = 4

obj = A()
A._a = 5
try:
    print(_a, end=' ')
    print(obj._a, end=' ')
    print(A._a, end=' ')
    print(obj._a, end=' ')
    print(obj._A__a)
except:
    print('ERR')
\end{verbatim}
%enditem


\newpage
9. Що буде надруковано за виконання?
\begin{verbatim}
class A0:
    def b(self): print(1, end=' ')
    def h(self): print(2, end=' ')

class A1(A0):
    def b(self): print(3, end=' ')

try:
    A0().h()
    A1().b()
except:
    print('ERR')
\end{verbatim}
%enditem


\newpage
10. Що буде надруковано за виконання?
\begin{verbatim}
class A:
    a = 1

    def _g1(self): print(2, end=' ')
    def g2(self): print(3, end=' ')

    def main(self):
        try: print(self.a, end=' ')
        except: print('E1', end=' ')

        try: self._g1()
        except: print('E2', end=' ')

        try: self.g2()
        except: print('E3', end=' ')


class B(A):
    a = 4

    def _g1(self): print(5, end=' ')
    def g2(self): print(6, end=' ')

try: B().main()
except: print('E4', end=' ')
\end{verbatim}
%enditem


\newpage
11. Змінні \kw{next} та \kw{prev} вузла списку позначають наступний та попередній за ним вузол відповідно. Починаючи з голови, у список записано послідовні цілі числа від~$-21$ до~$45$. Нехай змінна \kw{p} позначає вузол, що зберігає значення~$-1$.
Яке значення зберігається у вузлі \kw{p.prev.prev.prev.next.prev} ?
%enditem


\newpage
12. 
\begin{center}
	\adjustimage{max size={0.8\linewidth}{0.8\paperheight}}
	{BinTree/trees_exp/91.png}
\end{center}
Указати послідовність міток вершин дерева, зображеного на рис., що утвориться в результаті обходу дерева в прямому порядку. Мітки записувати через пробіл.\par Алгоритм починає роботу з кореня (на рис. це верхня вершина).
%enditem


\newpage
13. 
\begin{center}
	\adjustimage{max size={0.9\linewidth}{0.9\paperheight}}
	{Graphs/AdjList/106.png}
\end{center}
Для графа, зображеного на рис., вказати список суміжності вершини 1.\par
Формат оформлення відповіді:\par
список суміжності виписувати через пробіл на одному рядку, якщо пробілів десь буде більше одного - не важливо, упорядкування вершин у списку також значення не має.
%enditem


\newpage
14. Для графа на рисунку вказати послідовність, в якій алгоритм пошуку в глибину буде відмічати відвідані вершини, якщо списки суміжності вершин переглядаються алгоритмом за зростанням міток вершин та вершини відмічаються на першому вход\-женні у вершину.\par
Пошук розпочинається з вершини 4.\par

\begin{center}
	\adjustimage{max size={0.9\linewidth}{0.9\paperheight}}
	{Graphs/Traversals/63.png}
\end{center}
%enditem


\newpage
15. Для графа на рисунку з вершини 0 запущено алгоритм пошуку в ширину, що будує кістякове дерево. Списки суміжності вершин переглядаються алгоритмом за зростанням міток вершин.\par
\begin{center}
	\adjustimage{max size={0.9\linewidth}{0.9\paperheight}}
	{Graphs/Traversals/218.png}
\end{center}
\par Які з наведених ребер входять до складу отриманого кістякового дерева?\par
1) (1, 4)\par
2) (0, 6)\par
3) (3, 6)\par
4) (2, 6)\par
5) (2, 5)\par
6) (0, 2)\par
{\vskip 15pt} Номери відповідних ребер записати за зростанням через пробіл.


%enditem


\newpage
16. 

Рядки журналу через двокрапку містять ім'я користувача (ідентифікатор), час події,тип події, код події, наприклад
\begin{verbatim}
s="username : 2022-05-05 13:26 : ERROR : 404"
\end{verbatim}
у коді 
\begin{verbatim}
res = re.fullmatch('???', s) 
s = ''
code_str = ??? 
\end{verbatim}
чим треба замінити кожен з ???, щоб значенням code\_str був код події? 


\vspace{10mm}
\textbf{Завдання 17} 

Задано програму, що містить код 1-2 класів, можливе успадкування, та клієнтський код.
У наведеному коді невеликі частини закриті. 
Також наведено вихід програми.
Необхідно відновити закриті частини, щоб отриманий код відповідав наведеному виходу програми.


\textbf{Завдання 18-19} 

Кілька запитань за завданнями для самостійної роботи.
\vspace{10mm}
Скоріш за все, що завдань буде трохи менше, а деякі попередні будуть переформульовані в стилі двох останніх.

%enditem

\end {large}}

%end

\newpage
%begin

{{
\begin {large}

\newpage
1. Що буде надруковано за виконання?

\begin{verbatim}
print('R', end=' ')
b=('a','b','c','d','e',0,1)
print(b[-3])
\end{verbatim}

%enditem


\newpage
2. Що буде надруковано за виконання?

\begin{verbatim}
print('R', end=' ')
e=('0','1','2','3','4','5','6','7','8','9')
print(e[:-9:3])
\end{verbatim}

%enditem



\newpage
3. Що буде надруковано за виконання?

\begin{verbatim}
print('R', end=' ')
d=['0','1','2','3','4']
d.append([1,2])
print(d)
\end{verbatim}

%enditem



\newpage
4. Що буде надруковано за виконання?

\begin{verbatim}
print('R', end=' ')
b=[10,11,12,13,14]
b[7:-4]='12'
print(b)
\end{verbatim}

%enditem



\newpage
5. Укажіть ТИП значення виразу.\par
\begin{verbatim}
{x:x for x in range(7)}
\end{verbatim}
%enditem


\newpage
6. Що буде надруковано за виконання?\par
\begin{verbatim}
try:
    t = {45:4, 20:3, 24:1, 59:1}
    t[24] = 9
    for x in t :
        print(x, sep=' ', end=' ')
except:
    print('ERR')
\end{verbatim}
%enditem


\newpage
7. Нехай змінна \kw{s} позначає послідовність цілих, по якій можна ітерувати. 

Написати вираз, що перевіряє, чи правда, що послідовність  містить від’ємні числа.
%enditem


\newpage
8. Що буде надруковано за виконання?
%Оригинал был утерян. Требует отладки!
\begin{verbatim}
_b = 0

class B:
    _b = 1
    
    def __init__(self):
        self._b = 2
        _b = 3
        B._b = 4

obj = B()
B._b = 5
try:
    print(_b, end=' ')
    print(obj._b, end=' ')
    print(B._b, end=' ')
    print(obj._b, end=' ')
    print(obj._B__b)
except:
    print('ERR')
\end{verbatim}
%enditem


\newpage
9. Що буде надруковано за виконання?
\begin{verbatim}
class D0:
    def a(self): print(1, end=' ')
    def f(self): print(2, end=' ')

class D1(D0):
    def a(self): print(3, end=' ')

try:
    D1().f()
    D0().a()
except:
    print('ERR')
\end{verbatim}
%enditem


\newpage
10. Що буде надруковано за виконання?
\begin{verbatim}
class A:
    a = 1

    def __init__(self):
        c = 2
        a = 3

class B(A):
    a = 4
    b = 5

    def __init__(self):
        super().__init__()
        self.b = 6

try:
    obj = B()
    obj.a = obj.b + 10
    B.a = 13
except: print('EE', end=' ')

try: print(obj.a, end= ' ')
except: print('E1', end=' ')

try: print(A.a, end= ' ')
except: print('E2', end=' ')

try: print(B.a, end= ' ')
except: print('E3', end=' ')
\end{verbatim}
%enditem


\newpage
11. Змінні \kw{next} та \kw{prev} вузла списку позначають наступний та попередній за ним вузол відповідно. Починаючи з голови, у список записано послідовні цілі числа від~$-43$ до~$32$. Нехай змінна \kw{p} позначає вузол, що зберігає значення~$0$.
Яке значення зберігається у вузлі \kw{p.next.next.next.next.prev} ?
%enditem


\newpage
12. 
\begin{center}
	\adjustimage{max size={0.8\linewidth}{0.8\paperheight}}
	{BinTree/trees_exp/108.png}
\end{center}
Указати послідовність міток вершин дерева, зображеного на рис., що утвориться в результаті обходу дерева в оберненому порядку. Мітки записувати через пробіл.\par Алгоритм починає роботу з кореня (на рис. це верхня вершина).
%enditem


\newpage
13. 
\begin{center}
	\adjustimage{max size={0.9\linewidth}{0.9\paperheight}}
	{Graphs/AdjList/307.png}
\end{center}
Для графа, зображеного на рис., вказати список суміжності вершини 1.\par
Формат оформлення відповіді:\par
список суміжності виписувати через пробіл на одному рядку, якщо пробілів десь буде більше одного - не важливо, упорядкування вершин у списку також значення не має.
%enditem


\newpage
14. Для графа на рисунку з вершини 1 запущено алгоритм пошуку в глибину, що будує кістякове дерево. Списки суміжності вершин переглядаються алгоритмом за зростанням міток вершин.\par
\begin{center}
	\adjustimage{max size={0.9\linewidth}{0.9\paperheight}}
	{Graphs/Traversals/92.png}
\end{center}
\par Які з наведених ребер входять до складу отриманого кістякового дерева?\par
1) (2, 4)\par
2) (2, 3)\par
3) (4, 5)\par
4) (0, 5)\par
5) (0, 3)\par
6) (5, 6)\par
{\vskip 15pt} Номери відповідних ребер записати за зростанням через пробіл.

%enditem


\newpage
15. Для графа на рисунку з вершини 0 запущено алгоритм пошуку в ширину, що будує кістякове дерево. Списки суміжності вершин переглядаються алгоритмом за зростанням міток вершин.\par
\begin{center}
	\adjustimage{max size={0.9\linewidth}{0.9\paperheight}}
	{Graphs/Traversals/137.png}
\end{center}
\par Які з наведених ребер входять до складу отриманого кістякового дерева?\par
1) (1, 2)\par
2) (5, 6)\par
3) (0, 3)\par
4) (1, 3)\par
5) (2, 4)\par
6) (1, 6)\par
{\vskip 15pt} Номери відповідних ребер записати за зростанням через пробіл.


%enditem


\newpage
16. 

Рядки журналу через двокрапку містять ім'я користувача (ідентифікатор), час події,тип події, код події, наприклад
\begin{verbatim}
s="username : 2022-05-05 13:26 : ERROR : 404"
\end{verbatim}
у коді 
\begin{verbatim}
res = re.fullmatch('???', s) 
s = ''
code_str = ??? 
\end{verbatim}
чим треба замінити кожен з ???, щоб значенням code\_str був код події? 


\vspace{10mm}
\textbf{Завдання 17} 

Задано програму, що містить код 1-2 класів, можливе успадкування, та клієнтський код.
У наведеному коді невеликі частини закриті. 
Також наведено вихід програми.
Необхідно відновити закриті частини, щоб отриманий код відповідав наведеному виходу програми.


\textbf{Завдання 18-19} 

Кілька запитань за завданнями для самостійної роботи.
\vspace{10mm}
Скоріш за все, що завдань буде трохи менше, а деякі попередні будуть переформульовані в стилі двох останніх.

%enditem

\end {large}}

%end

\newpage
%begin

{{
\begin {large}

\newpage
1. Що буде надруковано за виконання?

\begin{verbatim}
print('R', end=' ')
c=('a','b','c','d','e',0,1,2,3)
print(c[6])
\end{verbatim}

%enditem


\newpage
2. Що буде надруковано за виконання?

\begin{verbatim}
print('R', end=' ')
g=(0,1,2,3,4,5,6,7,8,9)
print(g[-6:-1:-2])
\end{verbatim}

%enditem



\newpage
3. Що буде надруковано за виконання?

\begin{verbatim}
print('R', end=' ')
h=[10,11,12,13,14]
h+='12'
print(h)
\end{verbatim}

%enditem



\newpage
4. Що буде надруковано за виконання?

\begin{verbatim}
print('R', end=' ')
d=['0','1','2','3','4']
del d[8:3]
print(d)
\end{verbatim}

%enditem



\newpage
5. Укажіть ТИП значення виразу.\par
\begin{verbatim}
{y:None for y in range(-8,3)}
\end{verbatim}
%enditem


\newpage
6. Що буде надруковано за виконання?\par
\begin{verbatim}
try:
    s = {90:1, 55:9, 46:5, 26:4, 41:1}
    s[75] = 6
    for x, y in s.items():
        print(x, y, sep=' ', end=' ')
except:
    print('ERR')
\end{verbatim}
%enditem


\newpage
7. Задано словник \kw{d} з цілими ключами та значеннями. Записати ВИРАЗ, значенням якого є множина, що складається  з тих значень, чиї ключі діляться на 4.
%enditem


\newpage
8. Що буде надруковано за виконання?
%Оригинал был утерян. Требует отладки!
\begin{verbatim}
c = 0

class D:
    c = 1
    
    def __init__(self):
        self.c = 2
        c = 3
        D.c = 4

obj = D()
obj.c = 5
try:
    print(c, end=' ')
    print(obj.c, end=' ')
    print(D.c, end=' ')
    print(obj.c, end=' ')
    print(obj._D__c)
except:
    print('ERR')
\end{verbatim}
%enditem


\newpage
9. Що буде надруковано за виконання?
\begin{verbatim}
class B0:
    def c(self): print(1, end=' ')
    def h(self): print(2, end=' ')

class B1(B0):
    def c(self): print(3, end=' ')

try:
    B0().h()
    B1().c()
except:
    print('ERR')
\end{verbatim}
%enditem


\newpage
10. Що буде надруковано за виконання?
\begin{verbatim}
class A:
    c = 1

    def __init__(self):
        b = 2
        b = 3

class B(A):
    b = 4
    a = 5

    def __init__(self):
        super().__init__()
        self.a = 6

try:
    obj = B()
    obj.c = obj.c + 10
    B.c = 13
except: print('EE', end=' ')

try: print(obj.c, end= ' ')
except: print('E1', end=' ')

try: print(A.c, end= ' ')
except: print('E2', end=' ')

try: print(B.c, end= ' ')
except: print('E3', end=' ')
\end{verbatim}
%enditem


\newpage
11. Змінні \kw{next} та \kw{prev} вузла списку позначають наступний та попередній за ним вузол відповідно. Починаючи з голови, у список записано послідовні цілі числа від~$-67$ до~$47$. Нехай змінна \kw{p} позначає вузол, що зберігає значення~$6$.
Яке значення зберігається у вузлі \kw{p.prev.prev.prev.prev.next} ?
%enditem


\newpage
12. 
\begin{center}
	\adjustimage{max size={0.8\linewidth}{0.8\paperheight}}
	{BinTree/trees_exp/154.png}
\end{center}
Указати послідовність міток вершин дерева, зображеного на рис., що утвориться в результаті обходу дерева в прямому порядку. Мітки записувати через пробіл.\par Алгоритм починає роботу з кореня (на рис. це верхня вершина).
%enditem


\newpage
13. 
\begin{center}
	\adjustimage{max size={0.9\linewidth}{0.9\paperheight}}
	{Graphs/AdjList/268.png}
\end{center}
Для графа, зображеного на рис., вказати список суміжності вершини 3.\par
Формат оформлення відповіді:\par
список суміжності виписувати через пробіл на одному рядку, якщо пробілів десь буде більше одного - не важливо, упорядкування вершин у списку також значення не має.
%enditem


\newpage
14. Для графа на рисунку з вершини 6 запущено алгоритм пошуку в глибину, що будує кістякове дерево. Списки суміжності вершин переглядаються алгоритмом за зростанням міток вершин.\par
\begin{center}
	\adjustimage{max size={0.9\linewidth}{0.9\paperheight}}
	{Graphs/Traversals/300.png}
\end{center}
\par Які з наведених ребер входять до складу отриманого кістякового дерева?\par
1) (3, 6)\par
2) (0, 5)\par
3) (0, 2)\par
4) (2, 5)\par
5) (0, 1)\par
6) (5, 6)\par
{\vskip 15pt} Номери відповідних ребер записати за зростанням через пробіл.

%enditem


\newpage
15. Для графа на рисунку з вершини 4 запущено алгоритм пошуку в ширину, що будує кістякове дерево. Списки суміжності вершин переглядаються алгоритмом за зростанням міток вершин.\par
\begin{center}
	\adjustimage{max size={0.9\linewidth}{0.9\paperheight}}
	{Graphs/Traversals/63.png}
\end{center}
\par Які з наведених ребер входять до складу отриманого кістякового дерева?\par
1) (1, 3)\par
2) (1, 2)\par
3) (0, 2)\par
4) (1, 4)\par
5) (5, 6)\par
6) (2, 5)\par
{\vskip 15pt} Номери відповідних ребер записати за зростанням через пробіл.


%enditem


\newpage
16. 

Рядки журналу через двокрапку містять ім'я користувача (ідентифікатор), час події,тип події, код події, наприклад
\begin{verbatim}
s="username : 2022-05-05 13:26 : ERROR : 404"
\end{verbatim}
у коді 
\begin{verbatim}
res = re.fullmatch('???', s) 
s = ''
code_str = ??? 
\end{verbatim}
чим треба замінити кожен з ???, щоб значенням code\_str був код події? 


\vspace{10mm}
\textbf{Завдання 17} 

Задано програму, що містить код 1-2 класів, можливе успадкування, та клієнтський код.
У наведеному коді невеликі частини закриті. 
Також наведено вихід програми.
Необхідно відновити закриті частини, щоб отриманий код відповідав наведеному виходу програми.


\textbf{Завдання 18-19} 

Кілька запитань за завданнями для самостійної роботи.
\vspace{10mm}
Скоріш за все, що завдань буде трохи менше, а деякі попередні будуть переформульовані в стилі двох останніх.

%enditem

\end {large}}

%end

\newpage
%begin

{{
\begin {large}

\newpage
1. Що буде надруковано за виконання?

\begin{verbatim}
print('R', end=' ')
a=('a','b','c','d','e',0,1,2)
print(a[-7])
\end{verbatim}

%enditem


\newpage
2. Що буде надруковано за виконання?

\begin{verbatim}
print('R', end=' ')
h=[0,1,2,3,4,5,6,7,8,9]
print(h[::-3])
\end{verbatim}

%enditem



\newpage
3. Що буде надруковано за виконання?

\begin{verbatim}
print('R', end=' ')
e=[10,11,12,13,14]
e.insert(4, [4,5])
print(e)
\end{verbatim}

%enditem



\newpage
4. Що буде надруковано за виконання?

\begin{verbatim}
print('R', end=' ')
f=['0','1','2','3','4']
f[-7:]='11'
print(f)
\end{verbatim}

%enditem



\newpage
5. Укажіть ТИП значення виразу.\par
\begin{verbatim}
(2, 7, -4)
\end{verbatim}
%enditem


\newpage
6. Що буде надруковано за виконання?\par
\begin{verbatim}
try:
    d = {34:8, 45:0, 52:1, 34:7}
    d[20] = 3
    for x in d.keys():
        print(x, sep=' ', end=' ')
except:
    print('ERR')
\end{verbatim}
%enditem


\newpage
7. Нехай змінна \kw{s} позначає послідовність цілих, по якій можна ітерувати.

Написати вираз, що перевіряє, чи правда, що всі елементи послідовності є від’ємни\-ми.
%enditem


\newpage
8. Що буде надруковано за виконання?
%Оригинал был утерян. Требует отладки!
\begin{verbatim}
__c = 0

class A:
    __c = 1
    
    def __init__(self):
        self.__c = 2
        __c = 3
        A.__c = 4

obj = A()
obj.__c = 5
try:
    print(__c, end=' ')
    print(obj.__c, end=' ')
    print(A.__c, end=' ')
    print(obj.__c, end=' ')
    print(obj._A__c)
except:
    print('ERR')
\end{verbatim}
%enditem


\newpage
9. Що буде надруковано за виконання?
\begin{verbatim}
class C0:
    def d(self): print(1, end=' ')
    def g(self): print(2, end=' ')

class C1(C0):
    def d(self): print(3, end=' ')

try:
    C1().g()
    C1().d()
except:
    print('ERR')
\end{verbatim}
%enditem


\newpage
10. Що буде надруковано за виконання?
\begin{verbatim}
class A:
    b = 1

    def __g1(self): print(2, end=' ')
    def _g2(self): print(3, end=' ')

    def main(self):
        try: print(self.b, end=' ')
        except: print('E1', end=' ')

        try: self.__g1()
        except: print('E2', end=' ')

        try: self._g2()
        except: print('E3', end=' ')


class B(A):
    b = 4

    def __g1(self): print(5, end=' ')
    def _g2(self): print(6, end=' ')

try: B().main()
except: print('E4', end=' ')
\end{verbatim}
%enditem


\newpage
11. Змінні \kw{next} та \kw{prev} вузла списку позначають наступний та попередній за ним вузол відповідно. Починаючи з голови, у список записано послідовні цілі числа від~$-64$ до~$28$. Нехай змінна \kw{p} позначає вузол, що зберігає значення~$-6$.
Яке значення зберігається у вузлі \kw{p.prev.next.next.prev.next} ?
%enditem


\newpage
12. 
\begin{center}
	\adjustimage{max size={0.8\linewidth}{0.8\paperheight}}
	{BinTree/trees_exp/129.png}
\end{center}
Указати послідовність міток вершин дерева, зображеного на рис., що утвориться в результаті обходу дерева в оберненому порядку. Мітки записувати через пробіл.\par Алгоритм починає роботу з кореня (на рис. це верхня вершина).
%enditem


\newpage
13. 
\begin{center}
	\adjustimage{max size={0.9\linewidth}{0.9\paperheight}}
	{Graphs/AdjList/107.png}
\end{center}
Для графа, зображеного на рис., вказати список суміжності вершини 1.\par
Формат оформлення відповіді:\par
список суміжності виписувати через пробіл на одному рядку, якщо пробілів десь буде більше одного - не важливо, упорядкування вершин у списку також значення не має.
%enditem


\newpage
14. Для графа на рисунку вказати послідовність, в якій алгоритм пошуку в глибину буде відмічати відвідані вершини, якщо списки суміжності вершин переглядаються алгоритмом за зростанням міток вершин та вершини відмічаються на першому вход\-женні у вершину.\par
Пошук розпочинається з вершини 1.\par

\begin{center}
	\adjustimage{max size={0.9\linewidth}{0.9\paperheight}}
	{Graphs/Traversals/220.png}
\end{center}
%enditem


\newpage
15. Для графа на рисунку з вершини 3 запущено алгоритм пошуку в ширину, що будує кістякове дерево. Списки суміжності вершин переглядаються алгоритмом за зростанням міток вершин.\par
\begin{center}
	\adjustimage{max size={0.9\linewidth}{0.9\paperheight}}
	{Graphs/Traversals/298.png}
\end{center}
\par Які з наведених ребер входять до складу отриманого кістякового дерева?\par
1) (4, 6)\par
2) (4, 5)\par
3) (0, 5)\par
4) (0, 2)\par
5) (1, 2)\par
6) (1, 6)\par
{\vskip 15pt} Номери відповідних ребер записати за зростанням через пробіл.


%enditem


\newpage
16. 

Рядки журналу через двокрапку містять ім'я користувача (ідентифікатор), час події,тип події, код події, наприклад
\begin{verbatim}
s="username : 2022-05-05 13:26 : ERROR : 404"
\end{verbatim}
у коді 
\begin{verbatim}
res = re.fullmatch('???', s) 
s = ''
code_str = ??? 
\end{verbatim}
чим треба замінити кожен з ???, щоб значенням code\_str був код події? 


\vspace{10mm}
\textbf{Завдання 17} 

Задано програму, що містить код 1-2 класів, можливе успадкування, та клієнтський код.
У наведеному коді невеликі частини закриті. 
Також наведено вихід програми.
Необхідно відновити закриті частини, щоб отриманий код відповідав наведеному виходу програми.


\textbf{Завдання 18-19} 

Кілька запитань за завданнями для самостійної роботи.
\vspace{10mm}
Скоріш за все, що завдань буде трохи менше, а деякі попередні будуть переформульовані в стилі двох останніх.

%enditem

\end {large}}

%end

\newpage
%begin

{{
\begin {large}

\newpage
1. Що буде надруковано за виконання?

\begin{verbatim}
print('R', end=' ')
f=('a','b','c','d','e',0,1,2,4)
print(f[7])
\end{verbatim}

%enditem


\newpage
2. Що буде надруковано за виконання?

\begin{verbatim}
print('R', end=' ')
a='0123456789'
print(a[1::-1])
\end{verbatim}

%enditem



\newpage
3. Що буде надруковано за виконання?

\begin{verbatim}
print('R', end=' ')
f=['0','1','2','3','4']
print(f.index(3), end=' ')
\end{verbatim}

%enditem



\newpage
4. Що буде надруковано за виконання?

\begin{verbatim}
print('R', end=' ')
g=[10,11,12,13,14]
g[4:]=[1,0]
print(g)
\end{verbatim}

%enditem



\newpage
5. Укажіть ТИП значення виразу.\par
\begin{verbatim}
{x:0 for x in range(7)}
\end{verbatim}
%enditem


\newpage
6. Що буде надруковано за виконання?\par
\begin{verbatim}
try:
    d = {36:5, 19:9, 73:5, 85:4, 33:5}
    d[18] = 4
    for x in d.items():
        print(x, sep=' ', end=' ')
except:
    print('ERR')
\end{verbatim}
%enditem


\newpage
7. Нехай змінна \kw{s} позначає дуже довгу послідовність цілих, по якій можна ітерувати. 

Написати вираз-генератор, що задає підпослідовність її елементів, що є додатни\-ми. 

Обмеження: усі елементи шуканої підпослідовності одночасно в пам'яті розміщені бути не можуть. 
%enditem


\newpage
8. Що буде надруковано за виконання?
\begin{verbatim}
d = 0
h = 1

class C:
    d = 2
    
    def __init__(self):
        global d
        self.h = 3
        d = 4
        C.d = 5

obj = C()
try:
    print(d, end=' ')
    print(h, end=' ')
    print(C.d, end=' ')
    print(obj.h)
except:
    print('ERR')
\end{verbatim}
%enditem


\newpage
9. Що буде надруковано за виконання?
\begin{verbatim}
class A0:
    def b(self): print(1, end=' ')
    def h(self): print(2, end=' ')

class A1(A0):
    def h(self): print(3, end=' ')

try:
    A1().b()
    A1().h()
except:
    print('ERR')
\end{verbatim}
%enditem


\newpage
10. Що буде надруковано за виконання?
\begin{verbatim}
class A:
    c = 1

    def f1(self): print(2, end=' ')
    def _f2(self): print(3, end=' ')

    def main(self):
        try: print(self.c, end=' ')
        except: print('E1', end=' ')

        try: self.f1()
        except: print('E2', end=' ')

        try: self._f2()
        except: print('E3', end=' ')


class B(A):
    c = 4

    def f1(self): print(5, end=' ')
    def _f2(self): print(6, end=' ')

try: B().main()
except: print('E4', end=' ')
\end{verbatim}
%enditem


\newpage
11. Змінні \kw{next} та \kw{prev} вузла списку позначають наступний та попередній за ним вузол відповідно. Починаючи з голови, у список записано послідовні цілі числа від~$-55$ до~$36$. Нехай змінна \kw{p} позначає вузол, що зберігає значення~$-3$.
Яке значення зберігається у вузлі \kw{p.next.prev.prev.next.prev} ?
%enditem


\newpage
12. 
\begin{center}
	\adjustimage{max size={0.8\linewidth}{0.8\paperheight}}
	{BinTree/trees_exp/84.png}
\end{center}
Указати послідовність міток вершин дерева, зображеного на рис., що утвориться в результаті обходу дерева в прямому порядку. Мітки записувати через пробіл.\par Алгоритм починає роботу з кореня (на рис. це верхня вершина).
%enditem


\newpage
13. 
\begin{center}
	\adjustimage{max size={0.9\linewidth}{0.9\paperheight}}
	{Graphs/AdjList/206.png}
\end{center}
Для графа, зображеного на рис., вказати список суміжності вершини 2.\par
Формат оформлення відповіді:\par
список суміжності виписувати через пробіл на одному рядку, якщо пробілів десь буде більше одного - не важливо, упорядкування вершин у списку також значення не має.
%enditem


\newpage
14. Для графа на рисунку з вершини 6 запущено алгоритм пошуку в глибину, що будує кістякове дерево. Списки суміжності вершин переглядаються алгоритмом за зростанням міток вершин.\par
\begin{center}
	\adjustimage{max size={0.9\linewidth}{0.9\paperheight}}
	{Graphs/Traversals/138.png}
\end{center}
\par Які з наведених ребер входять до складу отриманого кістякового дерева?\par
1) (2, 4)\par
2) (2, 6)\par
3) (3, 4)\par
4) (0, 3)\par
5) (0, 4)\par
6) (1, 4)\par
{\vskip 15pt} Номери відповідних ребер записати за зростанням через пробіл.

%enditem


\newpage
15. Для графа на рисунку з вершини 3 запущено алгоритм пошуку в ширину, що будує кістякове дерево. Списки суміжності вершин переглядаються алгоритмом за зростанням міток вершин.\par
\begin{center}
	\adjustimage{max size={0.9\linewidth}{0.9\paperheight}}
	{Graphs/Traversals/159.png}
\end{center}
\par Які з наведених ребер входять до складу отриманого кістякового дерева?\par
1) (0, 3)\par
2) (2, 5)\par
3) (3, 5)\par
4) (3, 4)\par
5) (0, 6)\par
6) (5, 6)\par
{\vskip 15pt} Номери відповідних ребер записати за зростанням через пробіл.


%enditem


\newpage
16. 

Рядки журналу через двокрапку містять ім'я користувача (ідентифікатор), час події,тип події, код події, наприклад
\begin{verbatim}
s="username : 2022-05-05 13:26 : ERROR : 404"
\end{verbatim}
у коді 
\begin{verbatim}
res = re.fullmatch('???', s) 
s = ''
code_str = ??? 
\end{verbatim}
чим треба замінити кожен з ???, щоб значенням code\_str був код події? 


\vspace{10mm}
\textbf{Завдання 17} 

Задано програму, що містить код 1-2 класів, можливе успадкування, та клієнтський код.
У наведеному коді невеликі частини закриті. 
Також наведено вихід програми.
Необхідно відновити закриті частини, щоб отриманий код відповідав наведеному виходу програми.


\textbf{Завдання 18-19} 

Кілька запитань за завданнями для самостійної роботи.
\vspace{10mm}
Скоріш за все, що завдань буде трохи менше, а деякі попередні будуть переформульовані в стилі двох останніх.

%enditem

\end {large}}

%end

\newpage
%begin

{{
\begin {large}

\newpage
1. Що буде надруковано за виконання?

\begin{verbatim}
print('R', end=' ')
b=('a','b','c','d','e',0,1,2,3,4)
print(b[9])
\end{verbatim}

%enditem


\newpage
2. Що буде надруковано за виконання?

\begin{verbatim}
print('R', end=' ')
c=['0','1','2','3','4','5','6','7','8','9']
print(c[8:1:2])
\end{verbatim}

%enditem



\newpage
3. Що буде надруковано за виконання?

\begin{verbatim}
print('R', end=' ')
h=['0','1','2','3','4']
print(h.pop(1), end=' ')
print(h)
\end{verbatim}

%enditem



\newpage
4. Що буде надруковано за виконання?

\begin{verbatim}
print('R', end=' ')
d=['0','1','2','3','4']
del d[:-6]
print(d)
\end{verbatim}

%enditem



\newpage
5. Укажіть ТИП значення виразу.\par
\begin{verbatim}
[{1,2}, 5]
\end{verbatim}
%enditem


\newpage
6. Що буде надруковано за виконання?\par
\begin{verbatim}
try:
    s = {90:0, 64:4, 86:7, 90:1}
    s[37] = 1
    for x in s.items():
        print(*x, sep=' ', end=' ')
except:
    print('ERR')
\end{verbatim}
%enditem


\newpage
7. Задано множину \kw{x}, що складається з цілих чисел. Записати ВИРАЗ, значенням якого є множина, що складається  з елементів множини \kw{x}, що не діляться на 7.
%enditem


\newpage
8. Що буде надруковано за виконання?
%Оригинал был утерян. Требует отладки!
\begin{verbatim}
__d = 0

class B:
    __d = 1
    
    def __init__(self):
        self.__d = 2
        __d = 3
        B.__d = 4

obj = B()
B.__d = 5
try:
    print(__d, end=' ')
    print(obj.__d, end=' ')
    print(B.__d, end=' ')
    print(obj.__d, end=' ')
    print(obj._B__d)
except:
    print('ERR')
\end{verbatim}
%enditem


\newpage
9. Що буде надруковано за виконання?
\begin{verbatim}
class D0:
    def b(self): print(1, end=' ')
    def g(self): print(2, end=' ')

class D1(D0):
    def b(self): print(3, end=' ')

try:
    D1().g()
    D0().b()
except:
    print('ERR')
\end{verbatim}
%enditem


\newpage
10. Що буде надруковано за виконання?
\begin{verbatim}
class A:
    a = 1

    def __init__(self):
        b = 2
        a = 3

class B(A):
    a = 4
    c = 5

    def __init__(self):
        super().__init__()
        self.c = 6

try:
    obj = B()
    obj.a = obj.b + 10
    B.a = 13
except: print('EE', end=' ')

try: print(obj.a, end= ' ')
except: print('E1', end=' ')

try: print(A.a, end= ' ')
except: print('E2', end=' ')

try: print(B.a, end= ' ')
except: print('E3', end=' ')
\end{verbatim}
%enditem


\newpage
11. Змінні \kw{next} та \kw{prev} вузла списку позначають наступний та попередній за ним вузол відповідно. Починаючи з голови, у список записано послідовні цілі числа від~$-51$ до~$25$. Нехай змінна \kw{p} позначає вузол, що зберігає значення~$-2$.
Яке значення зберігається у вузлі \kw{p.next.prev.prev.prev.next} ?
%enditem


\newpage
12. 
\begin{center}
	\adjustimage{max size={0.8\linewidth}{0.8\paperheight}}
	{BinTree/trees_exp/62.png}
\end{center}
Указати послідовність міток вершин дерева, зображеного на рис., що утвориться в результаті обходу дерева в оберненому порядку. Мітки записувати через пробіл.\par Алгоритм починає роботу з кореня (на рис. це верхня вершина).
%enditem


\newpage
13. 
\begin{center}
	\adjustimage{max size={0.9\linewidth}{0.9\paperheight}}
	{Graphs/AdjList/332.png}
\end{center}
Для графа, зображеного на рис., вказати список суміжності вершини 7.\par
Формат оформлення відповіді:\par
список суміжності виписувати через пробіл на одному рядку, якщо пробілів десь буде більше одного - не важливо, упорядкування вершин у списку також значення не має.
%enditem


\newpage
14. Для графа на рисунку вказати послідовність, в якій алгоритм пошуку в глибину буде відмічати відвідані вершини, якщо списки суміжності вершин переглядаються алгоритмом за зростанням міток вершин та вершини відмічаються на першому вход\-женні у вершину.\par
Пошук розпочинається з вершини 3.\par

\begin{center}
	\adjustimage{max size={0.9\linewidth}{0.9\paperheight}}
	{Graphs/Traversals/68.png}
\end{center}
%enditem


\newpage
15. Для графа на рисунку з вершини 2 запущено алгоритм пошуку в ширину, що будує кістякове дерево. Списки суміжності вершин переглядаються алгоритмом за зростанням міток вершин.\par
\begin{center}
	\adjustimage{max size={0.9\linewidth}{0.9\paperheight}}
	{Graphs/Traversals/303.png}
\end{center}
\par Які з наведених ребер входять до складу отриманого кістякового дерева?\par
1) (1, 3)\par
2) (1, 4)\par
3) (4, 6)\par
4) (0, 4)\par
5) (1, 2)\par
6) (2, 6)\par
{\vskip 15pt} Номери відповідних ребер записати за зростанням через пробіл.


%enditem


\newpage
16. 

Рядки журналу через двокрапку містять ім'я користувача (ідентифікатор), час події,тип події, код події, наприклад
\begin{verbatim}
s="username : 2022-05-05 13:26 : ERROR : 404"
\end{verbatim}
у коді 
\begin{verbatim}
res = re.fullmatch('???', s) 
s = ''
code_str = ??? 
\end{verbatim}
чим треба замінити кожен з ???, щоб значенням code\_str був код події? 


\vspace{10mm}
\textbf{Завдання 17} 

Задано програму, що містить код 1-2 класів, можливе успадкування, та клієнтський код.
У наведеному коді невеликі частини закриті. 
Також наведено вихід програми.
Необхідно відновити закриті частини, щоб отриманий код відповідав наведеному виходу програми.


\textbf{Завдання 18-19} 

Кілька запитань за завданнями для самостійної роботи.
\vspace{10mm}
Скоріш за все, що завдань буде трохи менше, а деякі попередні будуть переформульовані в стилі двох останніх.

%enditem

\end {large}}

%end

\newpage
%begin

{{
\begin {large}

\newpage
1. Що буде надруковано за виконання?

\begin{verbatim}
print('R', end=' ')
g=('a','b','c','d','e',0,1,2,3)
print(g[3])
\end{verbatim}

%enditem


\newpage
2. Що буде надруковано за виконання?

\begin{verbatim}
print('R', end=' ')
b='0123456789'
print(b[3:-5:3])
\end{verbatim}

%enditem



\newpage
3. Що буде надруковано за виконання?

\begin{verbatim}
print('R', end=' ')
b=[10,11,12,13,14]
b.append('10')
print(b)
\end{verbatim}

%enditem



\newpage
4. Що буде надруковано за виконання?

\begin{verbatim}
print('R', end=' ')
g=[10,11,12,13,14]
g[-4:8]=(12)
print(g)
\end{verbatim}

%enditem



\newpage
5. Укажіть ТИП значення виразу.\par
\begin{verbatim}
{y:y*y for y in range(-7,7)}
\end{verbatim}
%enditem


\newpage
6. Що буде надруковано за виконання?\par
\begin{verbatim}
try:
    t = {65:8, 75:1, 52:4, 68:6, 65:9}
    t[42] = 2
    for x in t.values():
        print(x, sep=' ', end=' ')
except:
    print('ERR')
\end{verbatim}
%enditem


\newpage
7. Нехай змінна \kw{s} позначає послідовність цілих, по якій можна ітерувати.

Написати вираз, що перевіряє, чи правда, що всі елементи послідовності діляться на 11.
%enditem


\newpage
8. Що буде надруковано за виконання?
%Оригинал был утерян. Требует отладки!
\begin{verbatim}
_a = 0

class D:
    _a = 1
    
    def __init__(self):
        self._a = 2
        _a = 3
        D._a = 4

obj = D()
D._a = 5
try:
    print(_a, end=' ')
    print(obj._a, end=' ')
    print(D._a, end=' ')
    print(obj._a, end=' ')
    print(obj._D__a)
except:
    print('ERR')
\end{verbatim}
%enditem


\newpage
9. Що буде надруковано за виконання?
\begin{verbatim}
class A0:
    def c(self): print(1, end=' ')
    def f(self): print(2, end=' ')

class A1(A0):
    def c(self): print(3, end=' ')

try:
    A1().f()
    A1().c()
except:
    print('ERR')
\end{verbatim}
%enditem


\newpage
10. Що буде надруковано за виконання?
\begin{verbatim}
class A:
    c = 1

    def __init__(self):
        a = 2
        a = 3

class B(A):
    a = 4
    b = 5

    def __init__(self):
        super().__init__()
        self.b = 6

try:
    obj = B()
    obj.a = obj.c + 10
    B.a = 13
except: print('EE', end=' ')

try: print(obj.a, end= ' ')
except: print('E1', end=' ')

try: print(A.a, end= ' ')
except: print('E2', end=' ')

try: print(B.a, end= ' ')
except: print('E3', end=' ')
\end{verbatim}
%enditem


\newpage
11. Змінні \kw{next} та \kw{prev} вузла списку позначають наступний та попередній за ним вузол відповідно. Починаючи з голови, у список записано послідовні цілі числа від~$-29$ до~$24$. Нехай змінна \kw{p} позначає вузол, що зберігає значення~$-9$.
Яке значення зберігається у вузлі \kw{p.prev.next.next.next.prev} ?
%enditem


\newpage
12. 
\begin{center}
	\adjustimage{max size={0.8\linewidth}{0.8\paperheight}}
	{BinTree/trees_exp/110.png}
\end{center}
Указати послідовність міток вершин дерева, зображеного на рис., що утвориться в результаті обходу дерева в прямому порядку. Мітки записувати через пробіл.\par Алгоритм починає роботу з кореня (на рис. це верхня вершина).
%enditem


\newpage
13. 
\begin{center}
	\adjustimage{max size={0.9\linewidth}{0.9\paperheight}}
	{Graphs/AdjList/146.png}
\end{center}
Для графа, зображеного на рис., вказати список суміжності вершини 1.\par
Формат оформлення відповіді:\par
список суміжності виписувати через пробіл на одному рядку, якщо пробілів десь буде більше одного - не важливо, упорядкування вершин у списку також значення не має.
%enditem


\newpage
14. Для графа на рисунку вказати послідовність, в якій алгоритм пошуку в глибину буде відмічати відвідані вершини, якщо списки суміжності вершин переглядаються алгоритмом за зростанням міток вершин та вершини відмічаються на першому вход\-женні у вершину.\par
Пошук розпочинається з вершини 1.\par

\begin{center}
	\adjustimage{max size={0.9\linewidth}{0.9\paperheight}}
	{Graphs/Traversals/252.png}
\end{center}
%enditem


\newpage
15. Для графа на рисунку з вершини 6 запущено алгоритм пошуку в ширину, що будує кістякове дерево. Списки суміжності вершин переглядаються алгоритмом за зростанням міток вершин.\par
\begin{center}
	\adjustimage{max size={0.9\linewidth}{0.9\paperheight}}
	{Graphs/Traversals/294.png}
\end{center}
\par Які з наведених ребер входять до складу отриманого кістякового дерева?\par
1) (1, 3)\par
2) (2, 3)\par
3) (0, 4)\par
4) (2, 5)\par
5) (1, 2)\par
6) (0, 2)\par
{\vskip 15pt} Номери відповідних ребер записати за зростанням через пробіл.


%enditem


\newpage
16. 

Рядки журналу через двокрапку містять ім'я користувача (ідентифікатор), час події,тип події, код події, наприклад
\begin{verbatim}
s="username : 2022-05-05 13:26 : ERROR : 404"
\end{verbatim}
у коді 
\begin{verbatim}
res = re.fullmatch('???', s) 
s = ''
code_str = ??? 
\end{verbatim}
чим треба замінити кожен з ???, щоб значенням code\_str був код події? 


\vspace{10mm}
\textbf{Завдання 17} 

Задано програму, що містить код 1-2 класів, можливе успадкування, та клієнтський код.
У наведеному коді невеликі частини закриті. 
Також наведено вихід програми.
Необхідно відновити закриті частини, щоб отриманий код відповідав наведеному виходу програми.


\textbf{Завдання 18-19} 

Кілька запитань за завданнями для самостійної роботи.
\vspace{10mm}
Скоріш за все, що завдань буде трохи менше, а деякі попередні будуть переформульовані в стилі двох останніх.

%enditem

\end {large}}

%end

\newpage
%begin

{{
\begin {large}

\newpage
1. Що буде надруковано за виконання?

\begin{verbatim}
print('R', end=' ')
a=('a','b','c','d','e',0,1,2,3,4)
print(a[-1])
\end{verbatim}

%enditem


\newpage
2. Що буде надруковано за виконання?

\begin{verbatim}
print('R', end=' ')
f=(0,1,2,3,4,5,6,7,8,9)
print(f[4:-7:2])
\end{verbatim}

%enditem



\newpage
3. Що буде надруковано за виконання?

\begin{verbatim}
print('R', end=' ')
e=['0','1','2','3','4']
e.extend([1,0])
print(e)
\end{verbatim}

%enditem



\newpage
4. Що буде надруковано за виконання?

\begin{verbatim}
print('R', end=' ')
h=[10,11,12,13,14]
h[3:-6]=[14]
print(h)
\end{verbatim}

%enditem



\newpage
5. Укажіть ТИП значення виразу.\par
\begin{verbatim}
[7,2,3,4]
\end{verbatim}
%enditem


\newpage
6. Що буде надруковано за виконання?\par
\begin{verbatim}
try:
    t = {63:4, 50:2, 62:9, 62:2}
    t[62] = 9
    for x, y in t.items():
        print(x, y, sep=' ', end=' ')
except:
    print('ERR')
\end{verbatim}
%enditem


\newpage
7. Нехай змінна \kw{s} позначає дуже довгу послідовність, по якій можна ітерувати. 

Написати вираз-генератор, що задає підпослідовність її елементів, що не є рядками типу \kw{str}.

Обмеження: усі елементи шуканої підпослідовності одночасно в пам'яті розміщені бути не можуть. 
%enditem


\newpage
8. Що буде надруковано за виконання?
%Оригинал был утерян. Требует отладки!
\begin{verbatim}
b = 0

class C:
    b = 1
    
    def __init__(self):
        self.b = 2
        b = 3
        C.b = 4

obj = C()
obj.b = 5
try:
    print(b, end=' ')
    print(obj.b, end=' ')
    print(C.b, end=' ')
    print(obj.b, end=' ')
    print(obj._C__b)
except:
    print('ERR')
\end{verbatim}
%enditem


\newpage
9. Що буде надруковано за виконання?
\begin{verbatim}
class B0:
    def a(self): print(1, end=' ')
    def h(self): print(2, end=' ')

class B1(B0):
    def h(self): print(3, end=' ')

try:
    B0().a()
    B1().h()
except:
    print('ERR')
\end{verbatim}
%enditem


\newpage
10. Що буде надруковано за виконання?
\begin{verbatim}
class A:
    b = 1

    def __init__(self):
        c = 2
        c = 3

class B(A):
    c = 4
    a = 5

    def __init__(self):
        super().__init__()
        self.a = 6

try:
    obj = B()
    obj.b = obj.b + 10
    B.b = 13
except: print('EE', end=' ')

try: print(obj.b, end= ' ')
except: print('E1', end=' ')

try: print(A.b, end= ' ')
except: print('E2', end=' ')

try: print(B.b, end= ' ')
except: print('E3', end=' ')
\end{verbatim}
%enditem


\newpage
11. Змінні \kw{next} та \kw{prev} вузла списку позначають наступний та попередній за ним вузол відповідно. Починаючи з голови, у список записано послідовні цілі числа від~$-28$ до~$35$. Нехай змінна \kw{p} позначає вузол, що зберігає значення~$1$.
Яке значення зберігається у вузлі \kw{p.prev.prev.prev.prev.prev} ?
%enditem


\newpage
12. 
\begin{center}
	\adjustimage{max size={0.8\linewidth}{0.8\paperheight}}
	{BinTree/trees_exp/60.png}
\end{center}
Указати послідовність міток вершин дерева, зображеного на рис., що утвориться в результаті обходу дерева в оберненому порядку. Мітки записувати через пробіл.\par Алгоритм починає роботу з кореня (на рис. це верхня вершина).
%enditem


\newpage
13. 
\begin{center}
	\adjustimage{max size={0.9\linewidth}{0.9\paperheight}}
	{Graphs/AdjList/96.png}
\end{center}
Для графа, зображеного на рис., вказати список суміжності вершини 5.\par
Формат оформлення відповіді:\par
список суміжності виписувати через пробіл на одному рядку, якщо пробілів десь буде більше одного - не важливо, упорядкування вершин у списку також значення не має.
%enditem


\newpage
14. Для графа на рисунку вказати послідовність, в якій алгоритм пошуку в глибину буде відмічати відвідані вершини, якщо списки суміжності вершин переглядаються алгоритмом за зростанням міток вершин та вершини відмічаються на першому вход\-женні у вершину.\par
Пошук розпочинається з вершини 4.\par

\begin{center}
	\adjustimage{max size={0.9\linewidth}{0.9\paperheight}}
	{Graphs/Traversals/317.png}
\end{center}
%enditem


\newpage
15. Для графа на рисунку з вершини 6 запущено алгоритм пошуку в ширину, що будує кістякове дерево. Списки суміжності вершин переглядаються алгоритмом за зростанням міток вершин.\par
\begin{center}
	\adjustimage{max size={0.9\linewidth}{0.9\paperheight}}
	{Graphs/Traversals/33.png}
\end{center}
\par Які з наведених ребер входять до складу отриманого кістякового дерева?\par
1) (5, 6)\par
2) (2, 5)\par
3) (0, 2)\par
4) (3, 5)\par
5) (2, 4)\par
6) (1, 2)\par
{\vskip 15pt} Номери відповідних ребер записати за зростанням через пробіл.


%enditem


\newpage
16. 

Рядки журналу через двокрапку містять ім'я користувача (ідентифікатор), час події,тип події, код події, наприклад
\begin{verbatim}
s="username : 2022-05-05 13:26 : ERROR : 404"
\end{verbatim}
у коді 
\begin{verbatim}
res = re.fullmatch('???', s) 
s = ''
code_str = ??? 
\end{verbatim}
чим треба замінити кожен з ???, щоб значенням code\_str був код події? 


\vspace{10mm}
\textbf{Завдання 17} 

Задано програму, що містить код 1-2 класів, можливе успадкування, та клієнтський код.
У наведеному коді невеликі частини закриті. 
Також наведено вихід програми.
Необхідно відновити закриті частини, щоб отриманий код відповідав наведеному виходу програми.


\textbf{Завдання 18-19} 

Кілька запитань за завданнями для самостійної роботи.
\vspace{10mm}
Скоріш за все, що завдань буде трохи менше, а деякі попередні будуть переформульовані в стилі двох останніх.

%enditem

\end {large}}

%end

\newpage
%begin

{{
\begin {large}

\newpage
1. Що буде надруковано за виконання?

\begin{verbatim}
print('R', end=' ')
e=('a','b','c','d','e',0,1,2,4)
print(e[-9])
\end{verbatim}

%enditem


\newpage
2. Що буде надруковано за виконання?

\begin{verbatim}
print('R', end=' ')
g=['0','1','2','3','4','5','6','7','8','9']
print(g[:12:-3])
\end{verbatim}

%enditem



\newpage
3. Що буде надруковано за виконання?

\begin{verbatim}
print('R', end=' ')
c=[10,11,12,13,14]
c.remove('11')
print(c)
\end{verbatim}

%enditem



\newpage
4. Що буде надруковано за виконання?

\begin{verbatim}
print('R', end=' ')
e=['0','1','2','3','4']
del e[:-7]
print(e)
\end{verbatim}

%enditem



\newpage
5. Укажіть ТИП значення виразу.\par
\begin{verbatim}
(2, 7, -4)
\end{verbatim}
%enditem


\newpage
6. Що буде надруковано за виконання?\par
\begin{verbatim}
try:
    d = {48:3, 44:5, 17:9, 84:0}
    d[48] = 8
    for x in d.items():
        print(*x, sep=' ', end=' ')
except:
    print('ERR')
\end{verbatim}
%enditem


\newpage
7. Нехай змінна \kw{s} позначає послідовність цілих, по якій можна ітерувати. 

Написати вираз, що перевіряє, чи правда, що послідовність  містить число 35.
%enditem


\newpage
8. Що буде надруковано за виконання?
%Оригинал был утерян. Требует отладки!
\begin{verbatim}
_c = 0

class C:
    _c = 1
    
    def __init__(self):
        self._c = 2
        _c = 3
        C._c = 4

obj = C()
C._c = 5
try:
    print(_c, end=' ')
    print(obj._c, end=' ')
    print(C._c, end=' ')
    print(obj._c, end=' ')
    print(obj._C__c)
except:
    print('ERR')
\end{verbatim}
%enditem


\newpage
9. Що буде надруковано за виконання?
\begin{verbatim}
class C0:
    def d(self): print(1, end=' ')
    def g(self): print(2, end=' ')

class C1(C0):
    def g(self): print(3, end=' ')

try:
    C1().d()
    C1().g()
except:
    print('ERR')
\end{verbatim}
%enditem


\newpage
10. Що буде надруковано за виконання?
\begin{verbatim}
class A:
    b = 1

    def __init__(self):
        a = 2
        b = 3

class B(A):
    b = 4
    c = 5

    def __init__(self):
        super().__init__()
        self.c = 6

try:
    obj = B()
    obj.c = obj.a + 10
    B.c = 13
except: print('EE', end=' ')

try: print(obj.c, end= ' ')
except: print('E1', end=' ')

try: print(A.c, end= ' ')
except: print('E2', end=' ')

try: print(B.c, end= ' ')
except: print('E3', end=' ')
\end{verbatim}
%enditem


\newpage
11. Змінні \kw{next} та \kw{prev} вузла списку позначають наступний та попередній за ним вузол відповідно. Починаючи з голови, у список записано послідовні цілі числа від~$-62$ до~$37$. Нехай змінна \kw{p} позначає вузол, що зберігає значення~$5$.
Яке значення зберігається у вузлі \kw{p.next.next.next.next.next} ?
%enditem


\newpage
12. 
\begin{center}
	\adjustimage{max size={0.8\linewidth}{0.8\paperheight}}
	{BinTree/trees_exp/54.png}
\end{center}
Указати послідовність міток вершин дерева, зображеного на рис., що утвориться в результаті обходу дерева в прямому порядку. Мітки записувати через пробіл.\par Алгоритм починає роботу з кореня (на рис. це верхня вершина).
%enditem


\newpage
13. 
\begin{center}
	\adjustimage{max size={0.9\linewidth}{0.9\paperheight}}
	{Graphs/AdjList/327.png}
\end{center}
Для графа, зображеного на рис., вказати список суміжності вершини 1.\par
Формат оформлення відповіді:\par
список суміжності виписувати через пробіл на одному рядку, якщо пробілів десь буде більше одного - не важливо, упорядкування вершин у списку також значення не має.
%enditem


\newpage
14. Для графа на рисунку вказати послідовність, в якій алгоритм пошуку в глибину буде відмічати відвідані вершини, якщо списки суміжності вершин переглядаються алгоритмом за зростанням міток вершин та вершини відмічаються на першому вход\-женні у вершину.\par
Пошук розпочинається з вершини 4.\par

\begin{center}
	\adjustimage{max size={0.9\linewidth}{0.9\paperheight}}
	{Graphs/Traversals/175.png}
\end{center}
%enditem


\newpage
15. Для графа на рисунку з вершини 1 запущено алгоритм пошуку в ширину, що будує кістякове дерево. Списки суміжності вершин переглядаються алгоритмом за зростанням міток вершин.\par
\begin{center}
	\adjustimage{max size={0.9\linewidth}{0.9\paperheight}}
	{Graphs/Traversals/347.png}
\end{center}
\par Які з наведених ребер входять до складу отриманого кістякового дерева?\par
1) (1, 6)\par
2) (4, 6)\par
3) (1, 3)\par
4) (0, 5)\par
5) (3, 5)\par
6) (0, 6)\par
{\vskip 15pt} Номери відповідних ребер записати за зростанням через пробіл.


%enditem


\newpage
16. 

Рядки журналу через двокрапку містять ім'я користувача (ідентифікатор), час події,тип події, код події, наприклад
\begin{verbatim}
s="username : 2022-05-05 13:26 : ERROR : 404"
\end{verbatim}
у коді 
\begin{verbatim}
res = re.fullmatch('???', s) 
s = ''
code_str = ??? 
\end{verbatim}
чим треба замінити кожен з ???, щоб значенням code\_str був код події? 


\vspace{10mm}
\textbf{Завдання 17} 

Задано програму, що містить код 1-2 класів, можливе успадкування, та клієнтський код.
У наведеному коді невеликі частини закриті. 
Також наведено вихід програми.
Необхідно відновити закриті частини, щоб отриманий код відповідав наведеному виходу програми.


\textbf{Завдання 18-19} 

Кілька запитань за завданнями для самостійної роботи.
\vspace{10mm}
Скоріш за все, що завдань буде трохи менше, а деякі попередні будуть переформульовані в стилі двох останніх.

%enditem

\end {large}}

%end

\newpage
%begin

{{
\begin {large}

\newpage
1. Що буде надруковано за виконання?

\begin{verbatim}
print('R', end=' ')
f=('a','b','c','d','e',0,1,2)
print(f[2])
\end{verbatim}

%enditem


\newpage
2. Що буде надруковано за виконання?

\begin{verbatim}
print('R', end=' ')
e=[0,1,2,3,4,5,6,7,8,9]
print(e[-11::-1])
\end{verbatim}

%enditem



\newpage
3. Що буде надруковано за виконання?

\begin{verbatim}
print('R', end=' ')
a=['0','1','2','3','4']
a+=[2,3]
print(a)
\end{verbatim}

%enditem



\newpage
4. Що буде надруковано за виконання?

\begin{verbatim}
print('R', end=' ')
a=['0','1','2','3','4']
a[-9:-9]=[1,2]
print(a)
\end{verbatim}

%enditem



\newpage
5. Укажіть ТИП значення виразу.\par
\begin{verbatim}
{x*x:x for x in range(-7,7)}
\end{verbatim}
%enditem


\newpage
6. Що буде надруковано за виконання?\par
\begin{verbatim}
try:
    s = {12:6, 78:8, 55:1, 55:2}
    s[78] = 4
    for x in s.values():
        print(x, sep=' ', end=' ')
except:
    print('ERR')
\end{verbatim}
%enditem


\newpage
7. Нехай змінна \kw{s} позначає дуже довгу послідовність, по якій можна ітерувати. 

Написати вираз-генератор, що задає підпослідовність її елементів, що є числами типу \kw{float}.

Обмеження: усі елементи шуканої підпослідовності одночасно в пам'яті розміщені бути не можуть. 
%enditem


\newpage
8. Що буде надруковано за виконання?
%Оригинал был утерян. Требует отладки!
\begin{verbatim}
b = 0

class A:
    b = 1
    
    def __init__(self):
        self.b = 2
        b = 3
        A.b = 4

obj = A()
obj.b = 5
try:
    print(b, end=' ')
    print(obj.b, end=' ')
    print(A.b, end=' ')
    print(obj.b, end=' ')
    print(obj._A__b)
except:
    print('ERR')
\end{verbatim}
%enditem


\newpage
9. Що буде надруковано за виконання?
\begin{verbatim}
class B0:
    def b(self): print(1, end=' ')
    def f(self): print(2, end=' ')

class B1(B0):
    def f(self): print(3, end=' ')

try:
    B1().b()
    B1().f()
except:
    print('ERR')
\end{verbatim}
%enditem


\newpage
10. Що буде надруковано за виконання?
\begin{verbatim}
class A:
    a = 1

    def _h1(self): print(2, end=' ')
    def _h2(self): print(3, end=' ')

    def main(self):
        try: print(self.a, end=' ')
        except: print('E1', end=' ')

        try: self._h1()
        except: print('E2', end=' ')

        try: self._h2()
        except: print('E3', end=' ')


class B(A):
    a = 4

    def _h1(self): print(5, end=' ')
    def _h2(self): print(6, end=' ')

try: B().main()
except: print('E4', end=' ')
\end{verbatim}
%enditem


\newpage
11. Змінні \kw{next} та \kw{prev} вузла списку позначають наступний та попередній за ним вузол відповідно. Починаючи з голови, у список записано послідовні цілі числа від~$-39$ до~$20$. Нехай змінна \kw{p} позначає вузол, що зберігає значення~$7$.
Яке значення зберігається у вузлі \kw{p.prev.prev.next.prev.prev} ?
%enditem


\newpage
12. 
\begin{center}
	\adjustimage{max size={0.8\linewidth}{0.8\paperheight}}
	{BinTree/trees_exp/41.png}
\end{center}
Указати послідовність міток вершин дерева, зображеного на рис., що утвориться в результаті обходу дерева в оберненому порядку. Мітки записувати через пробіл.\par Алгоритм починає роботу з кореня (на рис. це верхня вершина).
%enditem


\newpage
13. 
\begin{center}
	\adjustimage{max size={0.9\linewidth}{0.9\paperheight}}
	{Graphs/AdjList/172.png}
\end{center}
Для графа, зображеного на рис., вказати список суміжності вершини 1.\par
Формат оформлення відповіді:\par
список суміжності виписувати через пробіл на одному рядку, якщо пробілів десь буде більше одного - не важливо, упорядкування вершин у списку також значення не має.
%enditem


\newpage
14. Для графа на рисунку вказати послідовність, в якій алгоритм пошуку в глибину буде відмічати відвідані вершини, якщо списки суміжності вершин переглядаються алгоритмом за зростанням міток вершин та вершини відмічаються на першому вход\-женні у вершину.\par
Пошук розпочинається з вершини 0.\par

\begin{center}
	\adjustimage{max size={0.9\linewidth}{0.9\paperheight}}
	{Graphs/Traversals/179.png}
\end{center}
%enditem


\newpage
15. Для графа на рисунку з вершини 2 запущено алгоритм пошуку в ширину, що будує кістякове дерево. Списки суміжності вершин переглядаються алгоритмом за зростанням міток вершин.\par
\begin{center}
	\adjustimage{max size={0.9\linewidth}{0.9\paperheight}}
	{Graphs/Traversals/244.png}
\end{center}
\par Які з наведених ребер входять до складу отриманого кістякового дерева?\par
1) (2, 4)\par
2) (2, 5)\par
3) (0, 2)\par
4) (0, 1)\par
5) (5, 6)\par
6) (3, 4)\par
{\vskip 15pt} Номери відповідних ребер записати за зростанням через пробіл.


%enditem


\newpage
16. 

Рядки журналу через двокрапку містять ім'я користувача (ідентифікатор), час події,тип події, код події, наприклад
\begin{verbatim}
s="username : 2022-05-05 13:26 : ERROR : 404"
\end{verbatim}
у коді 
\begin{verbatim}
res = re.fullmatch('???', s) 
s = ''
code_str = ??? 
\end{verbatim}
чим треба замінити кожен з ???, щоб значенням code\_str був код події? 


\vspace{10mm}
\textbf{Завдання 17} 

Задано програму, що містить код 1-2 класів, можливе успадкування, та клієнтський код.
У наведеному коді невеликі частини закриті. 
Також наведено вихід програми.
Необхідно відновити закриті частини, щоб отриманий код відповідав наведеному виходу програми.


\textbf{Завдання 18-19} 

Кілька запитань за завданнями для самостійної роботи.
\vspace{10mm}
Скоріш за все, що завдань буде трохи менше, а деякі попередні будуть переформульовані в стилі двох останніх.

%enditem

\end {large}}

%end

\newpage
%begin

{{
\begin {large}

\newpage
1. Що буде надруковано за виконання?

\begin{verbatim}
print('R', end=' ')
c=('a','b','c','d','e',0,1)
print(c[-2])
\end{verbatim}

%enditem


\newpage
2. Що буде надруковано за виконання?

\begin{verbatim}
print('R', end=' ')
d=('0','1','2','3','4','5','6','7','8','9')
print(d[-5:4:-2])
\end{verbatim}

%enditem



\newpage
3. Що буде надруковано за виконання?

\begin{verbatim}
print('R', end=' ')
g=[10,11,12,13,14]
g+=[1,2]
print(g)
\end{verbatim}

%enditem



\newpage
4. Що буде надруковано за виконання?

\begin{verbatim}
print('R', end=' ')
f=[10,11,12,13,14]
f[6:-4]=()
print(f)
\end{verbatim}

%enditem



\newpage
5. Укажіть ТИП значення виразу.\par
\begin{verbatim}
{x for x in range(7) if el%3==1}
\end{verbatim}
%enditem


\newpage
6. Що буде надруковано за виконання?\par
\begin{verbatim}
try:
    t = {34:2, 21:6, 67:7, 68:2}
    t[24] = 5
    for x in t :
        print(x, sep=' ', end=' ')
except:
    print('ERR')
\end{verbatim}
%enditem


\newpage
7. 
Записати ВИРАЗ, значенням якого є список, що складається з квадратів цілих чисел від 19 до 2020 (включно), що не діляться на 3, записаних в обернено\-му порядку.
%enditem


\newpage
8. Що буде надруковано за виконання?
%Оригинал был утерян. Требует отладки!
\begin{verbatim}
__a = 0

class D:
    __a = 1
    
    def __init__(self):
        self.__a = 2
        __a = 3
        D.__a = 4

obj = D()
obj.__a = 5
try:
    print(__a, end=' ')
    print(obj.__a, end=' ')
    print(D.__a, end=' ')
    print(obj.__a, end=' ')
    print(obj._D__a)
except:
    print('ERR')
\end{verbatim}
%enditem


\newpage
9. Що буде надруковано за виконання?
\begin{verbatim}
class A0:
    def a(self): print(1, end=' ')
    def h(self): print(2, end=' ')

class A1(A0):
    def a(self): print(3, end=' ')

try:
    A0().h()
    A1().a()
except:
    print('ERR')
\end{verbatim}
%enditem


\newpage
10. Що буде надруковано за виконання?
\begin{verbatim}
class A:
    c = 1

    def _f1(self): print(2, end=' ')
    def f2(self): print(3, end=' ')

    def main(self):
        try: print(self.c, end=' ')
        except: print('E1', end=' ')

        try: self._f1()
        except: print('E2', end=' ')

        try: self.f2()
        except: print('E3', end=' ')


class B(A):
    c = 4

    def _f1(self): print(5, end=' ')
    def f2(self): print(6, end=' ')

try: B().main()
except: print('E4', end=' ')
\end{verbatim}
%enditem


\newpage
11. Змінні \kw{next} та \kw{prev} вузла списку позначають наступний та попередній за ним вузол відповідно. Починаючи з голови, у список записано послідовні цілі числа від~$-66$ до~$41$. Нехай змінна \kw{p} позначає вузол, що зберігає значення~$-7$.
Яке значення зберігається у вузлі \kw{p.next.next.prev.next.next} ?
%enditem


\newpage
12. 
\begin{center}
	\adjustimage{max size={0.8\linewidth}{0.8\paperheight}}
	{BinTree/trees_exp/134.png}
\end{center}
Указати послідовність міток вершин дерева, зображеного на рис., що утвориться в результаті обходу дерева в прямому порядку. Мітки записувати через пробіл.\par Алгоритм починає роботу з кореня (на рис. це верхня вершина).
%enditem


\newpage
13. 
\begin{center}
	\adjustimage{max size={0.9\linewidth}{0.9\paperheight}}
	{Graphs/AdjList/275.png}
\end{center}
Для графа, зображеного на рис., вказати список суміжності вершини 1.\par
Формат оформлення відповіді:\par
список суміжності виписувати через пробіл на одному рядку, якщо пробілів десь буде більше одного - не важливо, упорядкування вершин у списку також значення не має.
%enditem


\newpage
14. Для графа на рисунку вказати послідовність, в якій алгоритм пошуку в глибину буде відмічати відвідані вершини, якщо списки суміжності вершин переглядаються алгоритмом за зростанням міток вершин та вершини відмічаються на першому вход\-женні у вершину.\par
Пошук розпочинається з вершини 0.\par

\begin{center}
	\adjustimage{max size={0.9\linewidth}{0.9\paperheight}}
	{Graphs/Traversals/233.png}
\end{center}
%enditem


\newpage
15. Для графа на рисунку з вершини 3 запущено алгоритм пошуку в ширину, що будує кістякове дерево. Списки суміжності вершин переглядаються алгоритмом за зростанням міток вершин.\par
\begin{center}
	\adjustimage{max size={0.9\linewidth}{0.9\paperheight}}
	{Graphs/Traversals/58.png}
\end{center}
\par Які з наведених ребер входять до складу отриманого кістякового дерева?\par
1) (0, 6)\par
2) (2, 5)\par
3) (5, 6)\par
4) (1, 2)\par
5) (1, 3)\par
6) (3, 6)\par
{\vskip 15pt} Номери відповідних ребер записати за зростанням через пробіл.


%enditem


\newpage
16. 

Рядки журналу через двокрапку містять ім'я користувача (ідентифікатор), час події,тип події, код події, наприклад
\begin{verbatim}
s="username : 2022-05-05 13:26 : ERROR : 404"
\end{verbatim}
у коді 
\begin{verbatim}
res = re.fullmatch('???', s) 
s = ''
code_str = ??? 
\end{verbatim}
чим треба замінити кожен з ???, щоб значенням code\_str був код події? 


\vspace{10mm}
\textbf{Завдання 17} 

Задано програму, що містить код 1-2 класів, можливе успадкування, та клієнтський код.
У наведеному коді невеликі частини закриті. 
Також наведено вихід програми.
Необхідно відновити закриті частини, щоб отриманий код відповідав наведеному виходу програми.


\textbf{Завдання 18-19} 

Кілька запитань за завданнями для самостійної роботи.
\vspace{10mm}
Скоріш за все, що завдань буде трохи менше, а деякі попередні будуть переформульовані в стилі двох останніх.

%enditem

\end {large}}

%end

\newpage
%begin

{{
\begin {large}

\newpage
1. Що буде надруковано за виконання?

\begin{verbatim}
print('R', end=' ')
h=('a','b','c','d','e',0,1,2,3,4)
print(h[-10])
\end{verbatim}

%enditem


\newpage
2. Що буде надруковано за виконання?

\begin{verbatim}
print('R', end=' ')
h=['0','1','2','3','4','5','6','7','8','9']
print(h[::-2])
\end{verbatim}

%enditem



\newpage
3. Що буде надруковано за виконання?

\begin{verbatim}
print('R', end=' ')
f=[10,11,12,13,14]
print(f.index(16), end=' ')
\end{verbatim}

%enditem



\newpage
4. Що буде надруковано за виконання?

\begin{verbatim}
print('R', end=' ')
b=['0','1','2','3','4']
del b[-6:3]
print(b)
\end{verbatim}

%enditem



\newpage
5. Укажіть ТИП значення виразу.\par
\begin{verbatim}
{1:3, 5:8}
\end{verbatim}
%enditem


\newpage
6. Що буде надруковано за виконання?\par
\begin{verbatim}
try:
    s = {75:6, 13:2, 47:4, 75:6}
    s[15] = 0
    for x in s.items():
        print(x, sep=' ', end=' ')
except:
    print('ERR')
\end{verbatim}
%enditem


\newpage
7. Нехай змінна \kw{s} позначає дуже довгу послідовність цілих, по якій можна ітерувати. 

Написати вираз-генератор, що задає підпослідовність її елементів, що є від'єм\-ни\-ми. 

Обмеження: усі елементи шуканої підпослідовності одночасно в пам'яті розміщені бути не можуть. 
%enditem


\newpage
8. Що буде надруковано за виконання?
\begin{verbatim}
b = 0
m = 1

class B:
    m = 2
    
    def __init__(self):
        self.b = 3
        m = 4
        B.m = 5

obj = B()
try:
    print(b, end=' ')
    print(m, end=' ')
    print(B.m, end=' ')
    print(obj.b)
except:
    print('ERR')
\end{verbatim}
%enditem


\newpage
9. Що буде надруковано за виконання?
\begin{verbatim}
class D0:
    def d(self): print(1, end=' ')
    def f(self): print(2, end=' ')

class D1(D0):
    def f(self): print(3, end=' ')

try:
    D1().d()
    D1().f()
except:
    print('ERR')
\end{verbatim}
%enditem


\newpage
10. Що буде надруковано за виконання?
\begin{verbatim}
class A:
    b = 1

    def _h1(self): print(2, end=' ')
    def __h2(self): print(3, end=' ')

    def main(self):
        try: print(self.b, end=' ')
        except: print('E1', end=' ')

        try: self._h1()
        except: print('E2', end=' ')

        try: self.__h2()
        except: print('E3', end=' ')


class B(A):
    b = 4

    def _h1(self): print(5, end=' ')
    def __h2(self): print(6, end=' ')

try: B().main()
except: print('E4', end=' ')
\end{verbatim}
%enditem


\newpage
11. Змінні \kw{next} та \kw{prev} вузла списку позначають наступний та попередній за ним вузол відповідно. Починаючи з голови, у список записано послідовні цілі числа від~$-48$ до~$22$. Нехай змінна \kw{p} позначає вузол, що зберігає значення~$-5$.
Яке значення зберігається у вузлі \kw{p.prev.next.prev.prev.prev} ?
%enditem


\newpage
12. 
\begin{center}
	\adjustimage{max size={0.8\linewidth}{0.8\paperheight}}
	{BinTree/trees_exp/141.png}
\end{center}
Указати послідовність міток вершин дерева, зображеного на рис., що утвориться в результаті обходу дерева в оберненому порядку. Мітки записувати через пробіл.\par Алгоритм починає роботу з кореня (на рис. це верхня вершина).
%enditem


\newpage
13. 
\begin{center}
	\adjustimage{max size={0.9\linewidth}{0.9\paperheight}}
	{Graphs/AdjList/63.png}
\end{center}
Для графа, зображеного на рис., вказати список суміжності вершини 0.\par
Формат оформлення відповіді:\par
список суміжності виписувати через пробіл на одному рядку, якщо пробілів десь буде більше одного - не важливо, упорядкування вершин у списку також значення не має.
%enditem


\newpage
14. Для графа на рисунку вказати послідовність, в якій алгоритм пошуку в глибину буде відмічати відвідані вершини, якщо списки суміжності вершин переглядаються алгоритмом за зростанням міток вершин та вершини відмічаються на першому вход\-женні у вершину.\par
Пошук розпочинається з вершини 2.\par

\begin{center}
	\adjustimage{max size={0.9\linewidth}{0.9\paperheight}}
	{Graphs/Traversals/17.png}
\end{center}
%enditem


\newpage
15. Для графа на рисунку з вершини 2 запущено алгоритм пошуку в ширину, що будує кістякове дерево. Списки суміжності вершин переглядаються алгоритмом за зростанням міток вершин.\par
\begin{center}
	\adjustimage{max size={0.9\linewidth}{0.9\paperheight}}
	{Graphs/Traversals/34.png}
\end{center}
\par Які з наведених ребер входять до складу отриманого кістякового дерева?\par
1) (1, 5)\par
2) (5, 6)\par
3) (0, 3)\par
4) (3, 5)\par
5) (3, 4)\par
6) (4, 5)\par
{\vskip 15pt} Номери відповідних ребер записати за зростанням через пробіл.


%enditem


\newpage
16. 

Рядки журналу через двокрапку містять ім'я користувача (ідентифікатор), час події,тип події, код події, наприклад
\begin{verbatim}
s="username : 2022-05-05 13:26 : ERROR : 404"
\end{verbatim}
у коді 
\begin{verbatim}
res = re.fullmatch('???', s) 
s = ''
code_str = ??? 
\end{verbatim}
чим треба замінити кожен з ???, щоб значенням code\_str був код події? 


\vspace{10mm}
\textbf{Завдання 17} 

Задано програму, що містить код 1-2 класів, можливе успадкування, та клієнтський код.
У наведеному коді невеликі частини закриті. 
Також наведено вихід програми.
Необхідно відновити закриті частини, щоб отриманий код відповідав наведеному виходу програми.


\textbf{Завдання 18-19} 

Кілька запитань за завданнями для самостійної роботи.
\vspace{10mm}
Скоріш за все, що завдань буде трохи менше, а деякі попередні будуть переформульовані в стилі двох останніх.

%enditem

\end {large}}

%end

\newpage
%begin

{{
\begin {large}

\newpage
1. Що буде надруковано за виконання?

\begin{verbatim}
print('R', end=' ')
d=('a','b','c','d','e',0,1)
print(d[-8])
\end{verbatim}

%enditem


\newpage
2. Що буде надруковано за виконання?

\begin{verbatim}
print('R', end=' ')
a='0123456789'
print(a[:14:2])
\end{verbatim}

%enditem



\newpage
3. Що буде надруковано за виконання?

\begin{verbatim}
print('R', end=' ')
d=['0','1','2','3','4']
d.append('12')
print(d)
\end{verbatim}

%enditem



\newpage
4. Що буде надруковано за виконання?

\begin{verbatim}
print('R', end=' ')
c=[10,11,12,13,14]
c[5:7]='13'
print(c)
\end{verbatim}

%enditem



\newpage
5. Укажіть ТИП значення виразу.\par
\begin{verbatim}
{y:None for y in range(-8,3)}
\end{verbatim}
%enditem


\newpage
6. Що буде надруковано за виконання?\par
\begin{verbatim}
try:
    d = {30:1, 71:2, 10:0, 10:6}
    d[37] = 4
    for x in d.keys():
        print(x, sep=' ', end=' ')
except:
    print('ERR')
\end{verbatim}
%enditem


\newpage
7. Нехай змінна \kw{s} позначає дуже довгу послідовність, по якій можна ітерувати. 

Написати вираз-генератор, що задає підпослідовність її елементів, що не є числами типу \kw{float}.

Обмеження: усі елементи шуканої підпослідовності одночасно в пам'яті розміщені бути не можуть. 
%enditem


\newpage
8. Що буде надруковано за виконання?
\begin{verbatim}
a = 0
n = 1

class C:
    a = 2
    
    def __init__(self):
        self.a = 3
        a = 4
        C.a = 5

obj = C()
try:
    print(a, end=' ')
    print(n, end=' ')
    print(C.a, end=' ')
    print(obj.a)
except:
    print('ERR')
\end{verbatim}
%enditem


\newpage
9. Що буде надруковано за виконання?
\begin{verbatim}
class C0:
    def c(self): print(1, end=' ')
    def g(self): print(2, end=' ')

class C1(C0):
    def g(self): print(3, end=' ')

try:
    C1().c()
    C0().g()
except:
    print('ERR')
\end{verbatim}
%enditem


\newpage
10. Що буде надруковано за виконання?
\begin{verbatim}
class A:
    a = 1

    def __init__(self):
        b = 2
        a = 3

class B(A):
    b = 4
    c = 5

    def __init__(self):
        super().__init__()
        self.c = 6

try:
    obj = B()
    obj.b = obj.c + 10
    B.b = 13
except: print('EE', end=' ')

try: print(obj.b, end= ' ')
except: print('E1', end=' ')

try: print(A.b, end= ' ')
except: print('E2', end=' ')

try: print(B.b, end= ' ')
except: print('E3', end=' ')
\end{verbatim}
%enditem


\newpage
11. Змінні \kw{next} та \kw{prev} вузла списку позначають наступний та попередній за ним вузол відповідно. Починаючи з голови, у список записано послідовні цілі числа від~$-56$ до~$29$. Нехай змінна \kw{p} позначає вузол, що зберігає значення~$-4$.
Яке значення зберігається у вузлі \kw{p.next.prev.next.next.next} ?
%enditem


\newpage
12. 
\begin{center}
	\adjustimage{max size={0.8\linewidth}{0.8\paperheight}}
	{BinTree/trees_exp/5.png}
\end{center}
Указати послідовність міток вершин дерева, зображеного на рис., що утвориться в результаті обходу дерева в прямому порядку. Мітки записувати через пробіл.\par Алгоритм починає роботу з кореня (на рис. це верхня вершина).
%enditem


\newpage
13. 
\begin{center}
	\adjustimage{max size={0.9\linewidth}{0.9\paperheight}}
	{Graphs/AdjList/41.png}
\end{center}
Для графа, зображеного на рис., вказати список суміжності вершини 0.\par
Формат оформлення відповіді:\par
список суміжності виписувати через пробіл на одному рядку, якщо пробілів десь буде більше одного - не важливо, упорядкування вершин у списку також значення не має.
%enditem


\newpage
14. Для графа на рисунку з вершини 3 запущено алгоритм пошуку в глибину, що будує кістякове дерево. Списки суміжності вершин переглядаються алгоритмом за зростанням міток вершин.\par
\begin{center}
	\adjustimage{max size={0.9\linewidth}{0.9\paperheight}}
	{Graphs/Traversals/345.png}
\end{center}
\par Які з наведених ребер входять до складу отриманого кістякового дерева?\par
1) (2, 6)\par
2) (2, 4)\par
3) (0, 1)\par
4) (0, 4)\par
5) (0, 6)\par
6) (1, 4)\par
{\vskip 15pt} Номери відповідних ребер записати за зростанням через пробіл.

%enditem


\newpage
15. Для графа на рисунку з вершини 1 запущено алгоритм пошуку в ширину, що будує кістякове дерево. Списки суміжності вершин переглядаються алгоритмом за зростанням міток вершин.\par
\begin{center}
	\adjustimage{max size={0.9\linewidth}{0.9\paperheight}}
	{Graphs/Traversals/111.png}
\end{center}
\par Які з наведених ребер входять до складу отриманого кістякового дерева?\par
1) (0, 4)\par
2) (1, 3)\par
3) (2, 6)\par
4) (0, 2)\par
5) (5, 6)\par
6) (2, 4)\par
{\vskip 15pt} Номери відповідних ребер записати за зростанням через пробіл.


%enditem


\newpage
16. 

Рядки журналу через двокрапку містять ім'я користувача (ідентифікатор), час події,тип події, код події, наприклад
\begin{verbatim}
s="username : 2022-05-05 13:26 : ERROR : 404"
\end{verbatim}
у коді 
\begin{verbatim}
res = re.fullmatch('???', s) 
s = ''
code_str = ??? 
\end{verbatim}
чим треба замінити кожен з ???, щоб значенням code\_str був код події? 


\vspace{10mm}
\textbf{Завдання 17} 

Задано програму, що містить код 1-2 класів, можливе успадкування, та клієнтський код.
У наведеному коді невеликі частини закриті. 
Також наведено вихід програми.
Необхідно відновити закриті частини, щоб отриманий код відповідав наведеному виходу програми.


\textbf{Завдання 18-19} 

Кілька запитань за завданнями для самостійної роботи.
\vspace{10mm}
Скоріш за все, що завдань буде трохи менше, а деякі попередні будуть переформульовані в стилі двох останніх.

%enditem

\end {large}}

%end

\newpage
%begin

{{
\begin {large}

\newpage
1. Що буде надруковано за виконання?

\begin{verbatim}
print('R', end=' ')
b=('a','b','c','d','e',0,1,2,4)
print(b[4])
\end{verbatim}

%enditem


\newpage
2. Що буде надруковано за виконання?

\begin{verbatim}
print('R', end=' ')
d=('0','1','2','3','4','5','6','7','8','9')
print(d[11:-11:-3])
\end{verbatim}

%enditem



\newpage
3. Що буде надруковано за виконання?

\begin{verbatim}
print('R', end=' ')
g=['0','1','2','3','4']
g.remove('3')
print(g)
\end{verbatim}

%enditem



\newpage
4. Що буде надруковано за виконання?

\begin{verbatim}
print('R', end=' ')
h=[10,11,12,13,14]
h[-5:]='10'
print(h)
\end{verbatim}

%enditem



\newpage
5. Укажіть ТИП значення виразу.\par
\begin{verbatim}
{x for x in range(7)}
\end{verbatim}
%enditem


\newpage
6. Що буде надруковано за виконання?\par
\begin{verbatim}
try:
    s = {28:2, 83:1, 71:9, 45:8}
    s[72] = 4
    for x, y in s.items():
        print(x, y, sep=' ', end=' ')
except:
    print('ERR')
\end{verbatim}
%enditem


\newpage
7. Записати ВИРАЗ, значенням якого є множина, що складається з цілих чисел від 16 до 2017 (включно), що не діляться на 7.
%enditem


\newpage
8. Що буде надруковано за виконання?
%Оригинал был утерян. Требует отладки!
\begin{verbatim}
__d = 0

class B:
    __d = 1
    
    def __init__(self):
        self.__d = 2
        __d = 3
        B.__d = 4

obj = B()
B.__d = 5
try:
    print(__d, end=' ')
    print(obj.__d, end=' ')
    print(B.__d, end=' ')
    print(obj.__d, end=' ')
    print(obj._B__d)
except:
    print('ERR')
\end{verbatim}
%enditem


\newpage
9. Що буде надруковано за виконання?
\begin{verbatim}
class C0:
    def a(self): print(1, end=' ')
    def h(self): print(2, end=' ')

class C1(C0):
    def h(self): print(3, end=' ')

try:
    C1().a()
    C1().h()
except:
    print('ERR')
\end{verbatim}
%enditem


\newpage
10. Що буде надруковано за виконання?
\begin{verbatim}
class A:
    d = 1

    def _g1(self): print(2, end=' ')
    def __g2(self): print(3, end=' ')

    def main(self):
        try: print(self.d, end=' ')
        except: print('E1', end=' ')

        try: self._g1()
        except: print('E2', end=' ')

        try: self.__g2()
        except: print('E3', end=' ')


class B(A):
    d = 4

    def _g1(self): print(5, end=' ')
    def __g2(self): print(6, end=' ')

try: B().main()
except: print('E4', end=' ')
\end{verbatim}
%enditem


\newpage
11. Змінні \kw{next} та \kw{prev} вузла списку позначають наступний та попередній за ним вузол відповідно. Починаючи з голови, у список записано послідовні цілі числа від~$-50$ до~$33$. Нехай змінна \kw{p} позначає вузол, що зберігає значення~$-8$.
Яке значення зберігається у вузлі \kw{p.next.next.next.next.prev} ?
%enditem


\newpage
12. 
\begin{center}
	\adjustimage{max size={0.8\linewidth}{0.8\paperheight}}
	{BinTree/trees_exp/158.png}
\end{center}
Указати послідовність міток вершин дерева, зображеного на рис., що утвориться в результаті обходу дерева в оберненому порядку. Мітки записувати через пробіл.\par Алгоритм починає роботу з кореня (на рис. це верхня вершина).
%enditem


\newpage
13. 
\begin{center}
	\adjustimage{max size={0.9\linewidth}{0.9\paperheight}}
	{Graphs/AdjList/7.png}
\end{center}
Для графа, зображеного на рис., вказати список суміжності вершини 3.\par
Формат оформлення відповіді:\par
список суміжності виписувати через пробіл на одному рядку, якщо пробілів десь буде більше одного - не важливо, упорядкування вершин у списку також значення не має.
%enditem


\newpage
14. Для графа на рисунку вказати послідовність, в якій алгоритм пошуку в глибину буде відмічати відвідані вершини, якщо списки суміжності вершин переглядаються алгоритмом за зростанням міток вершин та вершини відмічаються на першому вход\-женні у вершину.\par
Пошук розпочинається з вершини 6.\par

\begin{center}
	\adjustimage{max size={0.9\linewidth}{0.9\paperheight}}
	{Graphs/Traversals/313.png}
\end{center}
%enditem


\newpage
15. Для графа на рисунку з вершини 5 запущено алгоритм пошуку в ширину, що будує кістякове дерево. Списки суміжності вершин переглядаються алгоритмом за зростанням міток вершин.\par
\begin{center}
	\adjustimage{max size={0.9\linewidth}{0.9\paperheight}}
	{Graphs/Traversals/106.png}
\end{center}
\par Які з наведених ребер входять до складу отриманого кістякового дерева?\par
1) (2, 6)\par
2) (1, 6)\par
3) (5, 6)\par
4) (0, 5)\par
5) (2, 4)\par
6) (0, 6)\par
{\vskip 15pt} Номери відповідних ребер записати за зростанням через пробіл.


%enditem


\newpage
16. 

Рядки журналу через двокрапку містять ім'я користувача (ідентифікатор), час події,тип події, код події, наприклад
\begin{verbatim}
s="username : 2022-05-05 13:26 : ERROR : 404"
\end{verbatim}
у коді 
\begin{verbatim}
res = re.fullmatch('???', s) 
s = ''
code_str = ??? 
\end{verbatim}
чим треба замінити кожен з ???, щоб значенням code\_str був код події? 


\vspace{10mm}
\textbf{Завдання 17} 

Задано програму, що містить код 1-2 класів, можливе успадкування, та клієнтський код.
У наведеному коді невеликі частини закриті. 
Також наведено вихід програми.
Необхідно відновити закриті частини, щоб отриманий код відповідав наведеному виходу програми.


\textbf{Завдання 18-19} 

Кілька запитань за завданнями для самостійної роботи.
\vspace{10mm}
Скоріш за все, що завдань буде трохи менше, а деякі попередні будуть переформульовані в стилі двох останніх.

%enditem

\end {large}}

%end

\newpage
%begin

{{
\begin {large}

\newpage
1. Що буде надруковано за виконання?

\begin{verbatim}
print('R', end=' ')
a=('a','b','c','d','e',0,1,2,3)
print(a[14])
\end{verbatim}

%enditem


\newpage
2. Що буде надруковано за виконання?

\begin{verbatim}
print('R', end=' ')
a=(0,1,2,3,4,5,6,7,8,9)
print(a[-4::3])
\end{verbatim}

%enditem



\newpage
3. Що буде надруковано за виконання?

\begin{verbatim}
print('R', end=' ')
h=[10,11,12,13,14]
h.extend([2,3])
print(h)
\end{verbatim}

%enditem



\newpage
4. Що буде надруковано за виконання?

\begin{verbatim}
print('R', end=' ')
e=['0','1','2','3','4']
del e[2:-8]
print(e)
\end{verbatim}

%enditem



\newpage
5. Укажіть ТИП значення виразу.\par
\begin{verbatim}
[x for x in range(7)]
\end{verbatim}
%enditem


\newpage
6. Що буде надруковано за виконання?\par
\begin{verbatim}
try:
    t = {84:7, 25:8, 51:0, 47:5}
    t[84] = 8
    for x in t :
        print(x, sep=' ', end=' ')
except:
    print('ERR')
\end{verbatim}
%enditem


\newpage
7. 
Записати ВИРАЗ, значенням якого є список, що складається з цілих чисел від 19 до 2001 (включно), причому спочатку за зростанням записано всі, що діляться на 2, а потім решту за спаданням.

%enditem


\newpage
8. Що буде надруковано за виконання?
%Оригинал был утерян. Требует отладки!
\begin{verbatim}
d = 0

class D:
    d = 1
    
    def __init__(self):
        self.d = 2
        d = 3
        D.d = 4

obj = D()
obj.d = 5
try:
    print(d, end=' ')
    print(obj.d, end=' ')
    print(D.d, end=' ')
    print(obj.d, end=' ')
    print(obj._D__d)
except:
    print('ERR')
\end{verbatim}
%enditem


\newpage
9. Що буде надруковано за виконання?
\begin{verbatim}
class D0:
    def b(self): print(1, end=' ')
    def g(self): print(2, end=' ')

class D1(D0):
    def b(self): print(3, end=' ')

try:
    D0().g()
    D1().b()
except:
    print('ERR')
\end{verbatim}
%enditem


\newpage
10. Що буде надруковано за виконання?
\begin{verbatim}
class A:
    b = 1

    def __init__(self):
        c = 2
        c = 3

class B(A):
    b = 4
    a = 5

    def __init__(self):
        super().__init__()
        self.a = 6

try:
    obj = B()
    obj.a = obj.b + 10
    B.a = 13
except: print('EE', end=' ')

try: print(obj.a, end= ' ')
except: print('E1', end=' ')

try: print(A.a, end= ' ')
except: print('E2', end=' ')

try: print(B.a, end= ' ')
except: print('E3', end=' ')
\end{verbatim}
%enditem


\newpage
11. Змінні \kw{next} та \kw{prev} вузла списку позначають наступний та попередній за ним вузол відповідно. Починаючи з голови, у список записано послідовні цілі числа від~$-57$ до~$23$. Нехай змінна \kw{p} позначає вузол, що зберігає значення~$7$.
Яке значення зберігається у вузлі \kw{p.prev.prev.prev.prev.next} ?
%enditem


\newpage
12. 
\begin{center}
	\adjustimage{max size={0.8\linewidth}{0.8\paperheight}}
	{BinTree/trees_exp/162.png}
\end{center}
Указати послідовність міток вершин дерева, зображеного на рис., що утвориться в результаті обходу дерева в прямому порядку. Мітки записувати через пробіл.\par Алгоритм починає роботу з кореня (на рис. це верхня вершина).
%enditem


\newpage
13. 
\begin{center}
	\adjustimage{max size={0.9\linewidth}{0.9\paperheight}}
	{Graphs/AdjList/317.png}
\end{center}
Для графа, зображеного на рис., вказати список суміжності вершини 1.\par
Формат оформлення відповіді:\par
список суміжності виписувати через пробіл на одному рядку, якщо пробілів десь буде більше одного - не важливо, упорядкування вершин у списку також значення не має.
%enditem


\newpage
14. Для графа на рисунку вказати послідовність, в якій алгоритм пошуку в глибину буде відмічати відвідані вершини, якщо списки суміжності вершин переглядаються алгоритмом за зростанням міток вершин та вершини відмічаються на першому вход\-женні у вершину.\par
Пошук розпочинається з вершини 2.\par

\begin{center}
	\adjustimage{max size={0.9\linewidth}{0.9\paperheight}}
	{Graphs/Traversals/349.png}
\end{center}
%enditem


\newpage
15. Для графа на рисунку з вершини 1 запущено алгоритм пошуку в ширину, що будує кістякове дерево. Списки суміжності вершин переглядаються алгоритмом за зростанням міток вершин.\par
\begin{center}
	\adjustimage{max size={0.9\linewidth}{0.9\paperheight}}
	{Graphs/Traversals/290.png}
\end{center}
\par Які з наведених ребер входять до складу отриманого кістякового дерева?\par
1) (0, 2)\par
2) (3, 4)\par
3) (1, 6)\par
4) (5, 6)\par
5) (2, 6)\par
6) (1, 2)\par
{\vskip 15pt} Номери відповідних ребер записати за зростанням через пробіл.


%enditem


\newpage
16. 

Рядки журналу через двокрапку містять ім'я користувача (ідентифікатор), час події,тип події, код події, наприклад
\begin{verbatim}
s="username : 2022-05-05 13:26 : ERROR : 404"
\end{verbatim}
у коді 
\begin{verbatim}
res = re.fullmatch('???', s) 
s = ''
code_str = ??? 
\end{verbatim}
чим треба замінити кожен з ???, щоб значенням code\_str був код події? 


\vspace{10mm}
\textbf{Завдання 17} 

Задано програму, що містить код 1-2 класів, можливе успадкування, та клієнтський код.
У наведеному коді невеликі частини закриті. 
Також наведено вихід програми.
Необхідно відновити закриті частини, щоб отриманий код відповідав наведеному виходу програми.


\textbf{Завдання 18-19} 

Кілька запитань за завданнями для самостійної роботи.
\vspace{10mm}
Скоріш за все, що завдань буде трохи менше, а деякі попередні будуть переформульовані в стилі двох останніх.

%enditem

\end {large}}

%end

\newpage
%begin

{{
\begin {large}

\newpage
1. Що буде надруковано за виконання?

\begin{verbatim}
print('R', end=' ')
f=('a','b','c','d','e',0,1,2)
print(f[-13])
\end{verbatim}

%enditem


\newpage
2. Що буде надруковано за виконання?

\begin{verbatim}
print('R', end=' ')
h=[0,1,2,3,4,5,6,7,8,9]
print(h[-14:10:-1])
\end{verbatim}

%enditem



\newpage
3. Що буде надруковано за виконання?

\begin{verbatim}
print('R', end=' ')
e=[10,11,12,13,14]
print(e.pop(4), end=' ')
print(e)
\end{verbatim}

%enditem



\newpage
4. Що буде надруковано за виконання?

\begin{verbatim}
print('R', end=' ')
f=['0','1','2','3','4']
f[:6]='12'
print(f)
\end{verbatim}

%enditem



\newpage
5. Укажіть ТИП значення виразу.\par
\begin{verbatim}
{y:y*y for y in range(-7,7)}
\end{verbatim}
%enditem


\newpage
6. Що буде надруковано за виконання?\par
\begin{verbatim}
try:
    d = {15:5, 23:0, 22:3, 22:3}
    d[15] = 6
    for x in d.keys():
        print(x, sep=' ', end=' ')
except:
    print('ERR')
\end{verbatim}
%enditem


\newpage
7. Нехай змінна \kw{s} позначає дуже довгу послідовність, по якій можна ітерувати. 

Написати вираз-генератор, що задає підпослідовність її елементів, що є рядками типу \kw{str}.

Обмеження: усі елементи шуканої підпослідовності одночасно в пам'яті розміщені бути не можуть. 
%enditem


\newpage
8. Що буде надруковано за виконання?
\begin{verbatim}
c = 0
j = 1

class B:
    j = 2
    
    def __init__(self):
        global j
        self.j = 3
        j = 4
        B.j = 5

obj = B()
try:
    print(c, end=' ')
    print(j, end=' ')
    print(B.j, end=' ')
    print(obj.j)
except:
    print('ERR')
\end{verbatim}
%enditem


\newpage
9. Що буде надруковано за виконання?
\begin{verbatim}
class B0:
    def d(self): print(1, end=' ')
    def f(self): print(2, end=' ')

class B1(B0):
    def f(self): print(3, end=' ')

try:
    B1().d()
    B1().f()
except:
    print('ERR')
\end{verbatim}
%enditem


\newpage
10. Що буде надруковано за виконання?
\begin{verbatim}
class A:
    a = 1

    def __init__(self):
        c = 2
        a = 3

class B(A):
    a = 4
    b = 5

    def __init__(self):
        super().__init__()
        self.b = 6

try:
    obj = B()
    obj.c = obj.a + 10
    B.c = 13
except: print('EE', end=' ')

try: print(obj.c, end= ' ')
except: print('E1', end=' ')

try: print(A.c, end= ' ')
except: print('E2', end=' ')

try: print(B.c, end= ' ')
except: print('E3', end=' ')
\end{verbatim}
%enditem


\newpage
11. Змінні \kw{next} та \kw{prev} вузла списку позначають наступний та попередній за ним вузол відповідно. Починаючи з голови, у список записано послідовні цілі числа від~$-27$ до~$34$. Нехай змінна \kw{p} позначає вузол, що зберігає значення~$-7$.
Яке значення зберігається у вузлі \kw{p.next.prev.next.next.next} ?
%enditem


\newpage
12. 
\begin{center}
	\adjustimage{max size={0.8\linewidth}{0.8\paperheight}}
	{BinTree/trees_exp/87.png}
\end{center}
Указати послідовність міток вершин дерева, зображеного на рис., що утвориться в результаті обходу дерева в оберненому порядку. Мітки записувати через пробіл.\par Алгоритм починає роботу з кореня (на рис. це верхня вершина).
%enditem


\newpage
13. 
\begin{center}
	\adjustimage{max size={0.9\linewidth}{0.9\paperheight}}
	{Graphs/AdjList/4.png}
\end{center}
Для графа, зображеного на рис., вказати список суміжності вершини 2.\par
Формат оформлення відповіді:\par
список суміжності виписувати через пробіл на одному рядку, якщо пробілів десь буде більше одного - не важливо, упорядкування вершин у списку також значення не має.
%enditem


\newpage
14. Для графа на рисунку вказати послідовність, в якій алгоритм пошуку в глибину буде відмічати відвідані вершини, якщо списки суміжності вершин переглядаються алгоритмом за зростанням міток вершин та вершини відмічаються на першому вход\-женні у вершину.\par
Пошук розпочинається з вершини 2.\par

\begin{center}
	\adjustimage{max size={0.9\linewidth}{0.9\paperheight}}
	{Graphs/Traversals/192.png}
\end{center}
%enditem


\newpage
15. Для графа на рисунку з вершини 3 запущено алгоритм пошуку в ширину, що будує кістякове дерево. Списки суміжності вершин переглядаються алгоритмом за зростанням міток вершин.\par
\begin{center}
	\adjustimage{max size={0.9\linewidth}{0.9\paperheight}}
	{Graphs/Traversals/180.png}
\end{center}
\par Які з наведених ребер входять до складу отриманого кістякового дерева?\par
1) (0, 2)\par
2) (2, 4)\par
3) (0, 5)\par
4) (2, 5)\par
5) (3, 5)\par
6) (0, 1)\par
{\vskip 15pt} Номери відповідних ребер записати за зростанням через пробіл.


%enditem


\newpage
16. 

Рядки журналу через двокрапку містять ім'я користувача (ідентифікатор), час події,тип події, код події, наприклад
\begin{verbatim}
s="username : 2022-05-05 13:26 : ERROR : 404"
\end{verbatim}
у коді 
\begin{verbatim}
res = re.fullmatch('???', s) 
s = ''
code_str = ??? 
\end{verbatim}
чим треба замінити кожен з ???, щоб значенням code\_str був код події? 


\vspace{10mm}
\textbf{Завдання 17} 

Задано програму, що містить код 1-2 класів, можливе успадкування, та клієнтський код.
У наведеному коді невеликі частини закриті. 
Також наведено вихід програми.
Необхідно відновити закриті частини, щоб отриманий код відповідав наведеному виходу програми.


\textbf{Завдання 18-19} 

Кілька запитань за завданнями для самостійної роботи.
\vspace{10mm}
Скоріш за все, що завдань буде трохи менше, а деякі попередні будуть переформульовані в стилі двох останніх.

%enditem

\end {large}}

%end

\newpage
%begin

{{
\begin {large}

\newpage
1. Що буде надруковано за виконання?

\begin{verbatim}
print('R', end=' ')
d=('a','b','c','d','e',0,1,2,3,4)
print(d[-6])
\end{verbatim}

%enditem


\newpage
2. Що буде надруковано за виконання?

\begin{verbatim}
print('R', end=' ')
f=['0','1','2','3','4','5','6','7','8','9']
print(f[-8:0:2])
\end{verbatim}

%enditem



\newpage
3. Що буде надруковано за виконання?

\begin{verbatim}
print('R', end=' ')
d=['0','1','2','3','4']
d.insert(-5, [4,5])
print(d)
\end{verbatim}

%enditem



\newpage
4. Що буде надруковано за виконання?

\begin{verbatim}
print('R', end=' ')
c=[10,11,12,13,14]
c[:4]=[2,3]
print(c)
\end{verbatim}

%enditem



\newpage
5. Укажіть ТИП значення виразу.\par
\begin{verbatim}
(11,2,3)
\end{verbatim}
%enditem


\newpage
6. Що буде надруковано за виконання?\par
\begin{verbatim}
try:
    t = {80:3, 74:2, 71:8, 55:6, 71:5}
    t[71] = 5
    for x in t.values():
        print(x, sep=' ', end=' ')
except:
    print('ERR')
\end{verbatim}
%enditem


\newpage
7. Нехай змінна \kw{s} позначає послідовність цілих, по якій можна ітерувати. 

Написати вираз, що перевіряє, чи правда, що послідовність  містить число 35.
%enditem


\newpage
8. Що буде надруковано за виконання?
\begin{verbatim}
d = 0
k = 1

class A:
    d = 2
    
    def __init__(self):
        self.k = 3
        d = 4
        A.d = 5

obj = A()
try:
    print(d, end=' ')
    print(k, end=' ')
    print(A.d, end=' ')
    print(obj.k)
except:
    print('ERR')
\end{verbatim}
%enditem


\newpage
9. Що буде надруковано за виконання?
\begin{verbatim}
class A0:
    def c(self): print(1, end=' ')
    def g(self): print(2, end=' ')

class A1(A0):
    def g(self): print(3, end=' ')

try:
    A1().c()
    A0().g()
except:
    print('ERR')
\end{verbatim}
%enditem


\newpage
10. Що буде надруковано за виконання?
\begin{verbatim}
class A:
    d = 1

    def h1(self): print(2, end=' ')
    def _h2(self): print(3, end=' ')

    def main(self):
        try: print(self.d, end=' ')
        except: print('E1', end=' ')

        try: self.h1()
        except: print('E2', end=' ')

        try: self._h2()
        except: print('E3', end=' ')


class B(A):
    d = 4

    def h1(self): print(5, end=' ')
    def _h2(self): print(6, end=' ')

try: B().main()
except: print('E4', end=' ')
\end{verbatim}
%enditem


\newpage
11. Змінні \kw{next} та \kw{prev} вузла списку позначають наступний та попередній за ним вузол відповідно. Починаючи з голови, у список записано послідовні цілі числа від~$-16$ до~$49$. Нехай змінна \kw{p} позначає вузол, що зберігає значення~$3$.
Яке значення зберігається у вузлі \kw{p.prev.next.prev.prev.prev} ?
%enditem


\newpage
12. 
\begin{center}
	\adjustimage{max size={0.8\linewidth}{0.8\paperheight}}
	{BinTree/trees_exp/109.png}
\end{center}
Указати послідовність міток вершин дерева, зображеного на рис., що утвориться в результаті обходу дерева в прямому порядку. Мітки записувати через пробіл.\par Алгоритм починає роботу з кореня (на рис. це верхня вершина).
%enditem


\newpage
13. 
\begin{center}
	\adjustimage{max size={0.9\linewidth}{0.9\paperheight}}
	{Graphs/AdjList/263.png}
\end{center}
Для графа, зображеного на рис., вказати список суміжності вершини 3.\par
Формат оформлення відповіді:\par
список суміжності виписувати через пробіл на одному рядку, якщо пробілів десь буде більше одного - не важливо, упорядкування вершин у списку також значення не має.
%enditem


\newpage
14. Для графа на рисунку вказати послідовність, в якій алгоритм пошуку в глибину буде відмічати відвідані вершини, якщо списки суміжності вершин переглядаються алгоритмом за зростанням міток вершин та вершини відмічаються на першому вход\-женні у вершину.\par
Пошук розпочинається з вершини 3.\par

\begin{center}
	\adjustimage{max size={0.9\linewidth}{0.9\paperheight}}
	{Graphs/Traversals/102.png}
\end{center}
%enditem


\newpage
15. Для графа на рисунку з вершини 3 запущено алгоритм пошуку в ширину, що будує кістякове дерево. Списки суміжності вершин переглядаються алгоритмом за зростанням міток вершин.\par
\begin{center}
	\adjustimage{max size={0.9\linewidth}{0.9\paperheight}}
	{Graphs/Traversals/227.png}
\end{center}
\par Які з наведених ребер входять до складу отриманого кістякового дерева?\par
1) (1, 6)\par
2) (4, 5)\par
3) (4, 6)\par
4) (1, 5)\par
5) (0, 4)\par
6) (5, 6)\par
{\vskip 15pt} Номери відповідних ребер записати за зростанням через пробіл.


%enditem


\newpage
16. 

Рядки журналу через двокрапку містять ім'я користувача (ідентифікатор), час події,тип події, код події, наприклад
\begin{verbatim}
s="username : 2022-05-05 13:26 : ERROR : 404"
\end{verbatim}
у коді 
\begin{verbatim}
res = re.fullmatch('???', s) 
s = ''
code_str = ??? 
\end{verbatim}
чим треба замінити кожен з ???, щоб значенням code\_str був код події? 


\vspace{10mm}
\textbf{Завдання 17} 

Задано програму, що містить код 1-2 класів, можливе успадкування, та клієнтський код.
У наведеному коді невеликі частини закриті. 
Також наведено вихід програми.
Необхідно відновити закриті частини, щоб отриманий код відповідав наведеному виходу програми.


\textbf{Завдання 18-19} 

Кілька запитань за завданнями для самостійної роботи.
\vspace{10mm}
Скоріш за все, що завдань буде трохи менше, а деякі попередні будуть переформульовані в стилі двох останніх.

%enditem

\end {large}}

%end

\newpage
%begin

{{
\begin {large}

\newpage
1. Що буде надруковано за виконання?

\begin{verbatim}
print('R', end=' ')
e=('a','b','c','d','e',0,1)
print(e[-13])
\end{verbatim}

%enditem


\newpage
2. Що буде надруковано за виконання?

\begin{verbatim}
print('R', end=' ')
c=('0','1','2','3','4','5','6','7','8','9')
print(c[-1:-2:-1])
\end{verbatim}

%enditem



\newpage
3. Що буде надруковано за виконання?

\begin{verbatim}
print('R', end=' ')
c=['0','1','2','3','4']
c.append('13')
print(c)
\end{verbatim}

%enditem



\newpage
4. Що буде надруковано за виконання?

\begin{verbatim}
print('R', end=' ')
g=[10,11,12,13,14]
g[8:]=[4,5]
print(g)
\end{verbatim}

%enditem



\newpage
5. Укажіть ТИП значення виразу.\par
\begin{verbatim}
{x:0 for x in range(7)}
\end{verbatim}
%enditem


\newpage
6. Що буде надруковано за виконання?\par
\begin{verbatim}
try:
    d = {18:0, 73:8, 67:2, 73:2}
    d[46] = 9
    for x in d.items():
        print(x, sep=' ', end=' ')
except:
    print('ERR')
\end{verbatim}
%enditem


\newpage
7. Нехай змінна \kw{s} позначає послідовність цілих, по якій можна ітерувати. 

Написати вираз, що перевіряє, чи правда, що послідовність  містить хоча б одне додатнє число.
%enditem


\newpage
8. Що буде надруковано за виконання?
%Оригинал был утерян. Требует отладки!
\begin{verbatim}
_a = 0

class A:
    _a = 1
    
    def __init__(self):
        self._a = 2
        _a = 3
        A._a = 4

obj = A()
A._a = 5
try:
    print(_a, end=' ')
    print(obj._a, end=' ')
    print(A._a, end=' ')
    print(obj._a, end=' ')
    print(obj._A__a)
except:
    print('ERR')
\end{verbatim}
%enditem


\newpage
9. Що буде надруковано за виконання?
\begin{verbatim}
class C0:
    def c(self): print(1, end=' ')
    def h(self): print(2, end=' ')

class C1(C0):
    def h(self): print(3, end=' ')

try:
    C1().c()
    C1().h()
except:
    print('ERR')
\end{verbatim}
%enditem


\newpage
10. Що буде надруковано за виконання?
\begin{verbatim}
class A:
    a = 1

    def _g1(self): print(2, end=' ')
    def g2(self): print(3, end=' ')

    def main(self):
        try: print(self.a, end=' ')
        except: print('E1', end=' ')

        try: self._g1()
        except: print('E2', end=' ')

        try: self.g2()
        except: print('E3', end=' ')


class B(A):
    a = 4

    def _g1(self): print(5, end=' ')
    def g2(self): print(6, end=' ')

try: B().main()
except: print('E4', end=' ')
\end{verbatim}
%enditem


\newpage
11. Змінні \kw{next} та \kw{prev} вузла списку позначають наступний та попередній за ним вузол відповідно. Починаючи з голови, у список записано послідовні цілі числа від~$-32$ до~$26$. Нехай змінна \kw{p} позначає вузол, що зберігає значення~$-3$.
Яке значення зберігається у вузлі \kw{p.prev.next.prev.next.prev} ?
%enditem


\newpage
12. 
\begin{center}
	\adjustimage{max size={0.8\linewidth}{0.8\paperheight}}
	{BinTree/trees_exp/136.png}
\end{center}
Указати послідовність міток вершин дерева, зображеного на рис., що утвориться в результаті обходу дерева в прямому порядку. Мітки записувати через пробіл.\par Алгоритм починає роботу з кореня (на рис. це верхня вершина).
%enditem


\newpage
13. 
\begin{center}
	\adjustimage{max size={0.9\linewidth}{0.9\paperheight}}
	{Graphs/AdjList/11.png}
\end{center}
Для графа, зображеного на рис., вказати список суміжності вершини 0.\par
Формат оформлення відповіді:\par
список суміжності виписувати через пробіл на одному рядку, якщо пробілів десь буде більше одного - не важливо, упорядкування вершин у списку також значення не має.
%enditem


\newpage
14. Для графа на рисунку вказати послідовність, в якій алгоритм пошуку в глибину буде відмічати відвідані вершини, якщо списки суміжності вершин переглядаються алгоритмом за зростанням міток вершин та вершини відмічаються на першому вход\-женні у вершину.\par
Пошук розпочинається з вершини 2.\par

\begin{center}
	\adjustimage{max size={0.9\linewidth}{0.9\paperheight}}
	{Graphs/Traversals/161.png}
\end{center}
%enditem


\newpage
15. Для графа на рисунку з вершини 0 запущено алгоритм пошуку в ширину, що будує кістякове дерево. Списки суміжності вершин переглядаються алгоритмом за зростанням міток вершин.\par
\begin{center}
	\adjustimage{max size={0.9\linewidth}{0.9\paperheight}}
	{Graphs/Traversals/190.png}
\end{center}
\par Які з наведених ребер входять до складу отриманого кістякового дерева?\par
1) (2, 4)\par
2) (0, 3)\par
3) (1, 6)\par
4) (1, 2)\par
5) (2, 6)\par
6) (0, 4)\par
{\vskip 15pt} Номери відповідних ребер записати за зростанням через пробіл.


%enditem


\newpage
16. 

Рядки журналу через двокрапку містять ім'я користувача (ідентифікатор), час події,тип події, код події, наприклад
\begin{verbatim}
s="username : 2022-05-05 13:26 : ERROR : 404"
\end{verbatim}
у коді 
\begin{verbatim}
res = re.fullmatch('???', s) 
s = ''
code_str = ??? 
\end{verbatim}
чим треба замінити кожен з ???, щоб значенням code\_str був код події? 


\vspace{10mm}
\textbf{Завдання 17} 

Задано програму, що містить код 1-2 класів, можливе успадкування, та клієнтський код.
У наведеному коді невеликі частини закриті. 
Також наведено вихід програми.
Необхідно відновити закриті частини, щоб отриманий код відповідав наведеному виходу програми.


\textbf{Завдання 18-19} 

Кілька запитань за завданнями для самостійної роботи.
\vspace{10mm}
Скоріш за все, що завдань буде трохи менше, а деякі попередні будуть переформульовані в стилі двох останніх.

%enditem

\end {large}}

%end

\newpage
%begin

{{
\begin {large}

\newpage
1. Що буде надруковано за виконання?

\begin{verbatim}
print('R', end=' ')
h=('a','b','c','d','e',0,1,2)
print(h[0])
\end{verbatim}

%enditem


\newpage
2. Що буде надруковано за виконання?

\begin{verbatim}
print('R', end=' ')
g=[0,1,2,3,4,5,6,7,8,9]
print(g[:-14:-3])
\end{verbatim}

%enditem



\newpage
3. Що буде надруковано за виконання?

\begin{verbatim}
print('R', end=' ')
a=[10,11,12,13,14]
print(a.index(10), end=' ')
\end{verbatim}

%enditem



\newpage
4. Що буде надруковано за виконання?

\begin{verbatim}
print('R', end=' ')
d=['0','1','2','3','4']
del d[-8:8]
print(d)
\end{verbatim}

%enditem



\newpage
5. Укажіть ТИП значення виразу.\par
\begin{verbatim}
{(1,4), 8}
\end{verbatim}
%enditem


\newpage
6. Що буде надруковано за виконання?\par
\begin{verbatim}
try:
    s = {57:4, 39:7, 30:8, 55:4, 30:5}
    s[55] = 2
    for x in s.items():
        print(*x, sep=' ', end=' ')
except:
    print('ERR')
\end{verbatim}
%enditem


\newpage
7. Нехай змінна \kw{s} позначає послідовність цілих, по якій можна ітерувати.

Написати вираз, що перевіряє, чи правда, що всі елементи послідовності більші за задане число num.
%enditem


\newpage
8. Що буде надруковано за виконання?
\begin{verbatim}
b = 0
m = 1

class D:
    m = 2
    
    def __init__(self):
        global m
        self.b = 3
        m = 4
        D.m = 5

obj = D()
try:
    print(b, end=' ')
    print(m, end=' ')
    print(D.m, end=' ')
    print(obj.b)
except:
    print('ERR')
\end{verbatim}
%enditem


\newpage
9. Що буде надруковано за виконання?
\begin{verbatim}
class D0:
    def d(self): print(1, end=' ')
    def f(self): print(2, end=' ')

class D1(D0):
    def f(self): print(3, end=' ')

try:
    D0().d()
    D1().f()
except:
    print('ERR')
\end{verbatim}
%enditem


\newpage
10. Що буде надруковано за виконання?
\begin{verbatim}
class A:
    b = 1

    def __init__(self):
        a = 2
        a = 3

class B(A):
    a = 4
    c = 5

    def __init__(self):
        super().__init__()
        self.c = 6

try:
    obj = B()
    obj.a = obj.b + 10
    B.a = 13
except: print('EE', end=' ')

try: print(obj.a, end= ' ')
except: print('E1', end=' ')

try: print(A.a, end= ' ')
except: print('E2', end=' ')

try: print(B.a, end= ' ')
except: print('E3', end=' ')
\end{verbatim}
%enditem


\newpage
11. Змінні \kw{next} та \kw{prev} вузла списку позначають наступний та попередній за ним вузол відповідно. Починаючи з голови, у список записано послідовні цілі числа від~$-22$ до~$16$. Нехай змінна \kw{p} позначає вузол, що зберігає значення~$-2$.
Яке значення зберігається у вузлі \kw{p.next.prev.next.prev.next} ?
%enditem


\newpage
12. 
\begin{center}
	\adjustimage{max size={0.8\linewidth}{0.8\paperheight}}
	{BinTree/trees_exp/3.png}
\end{center}
Указати послідовність міток вершин дерева, зображеного на рис., що утвориться в результаті обходу дерева в оберненому порядку. Мітки записувати через пробіл.\par Алгоритм починає роботу з кореня (на рис. це верхня вершина).
%enditem


\newpage
13. 
\begin{center}
	\adjustimage{max size={0.9\linewidth}{0.9\paperheight}}
	{Graphs/AdjList/162.png}
\end{center}
Для графа, зображеного на рис., вказати список суміжності вершини 6.\par
Формат оформлення відповіді:\par
список суміжності виписувати через пробіл на одному рядку, якщо пробілів десь буде більше одного - не важливо, упорядкування вершин у списку також значення не має.
%enditem


\newpage
14. Для графа на рисунку з вершини 1 запущено алгоритм пошуку в глибину, що будує кістякове дерево. Списки суміжності вершин переглядаються алгоритмом за зростанням міток вершин.\par
\begin{center}
	\adjustimage{max size={0.9\linewidth}{0.9\paperheight}}
	{Graphs/Traversals/203.png}
\end{center}
\par Які з наведених ребер входять до складу отриманого кістякового дерева?\par
1) (0, 5)\par
2) (2, 3)\par
3) (1, 2)\par
4) (0, 3)\par
5) (0, 4)\par
6) (4, 5)\par
{\vskip 15pt} Номери відповідних ребер записати за зростанням через пробіл.

%enditem


\newpage
15. Для графа на рисунку з вершини 2 запущено алгоритм пошуку в ширину, що будує кістякове дерево. Списки суміжності вершин переглядаються алгоритмом за зростанням міток вершин.\par
\begin{center}
	\adjustimage{max size={0.9\linewidth}{0.9\paperheight}}
	{Graphs/Traversals/173.png}
\end{center}
\par Які з наведених ребер входять до складу отриманого кістякового дерева?\par
1) (2, 6)\par
2) (0, 6)\par
3) (4, 6)\par
4) (3, 4)\par
5) (4, 5)\par
6) (0, 4)\par
{\vskip 15pt} Номери відповідних ребер записати за зростанням через пробіл.


%enditem


\newpage
16. 

Рядки журналу через двокрапку містять ім'я користувача (ідентифікатор), час події,тип події, код події, наприклад
\begin{verbatim}
s="username : 2022-05-05 13:26 : ERROR : 404"
\end{verbatim}
у коді 
\begin{verbatim}
res = re.fullmatch('???', s) 
s = ''
code_str = ??? 
\end{verbatim}
чим треба замінити кожен з ???, щоб значенням code\_str був код події? 


\vspace{10mm}
\textbf{Завдання 17} 

Задано програму, що містить код 1-2 класів, можливе успадкування, та клієнтський код.
У наведеному коді невеликі частини закриті. 
Також наведено вихід програми.
Необхідно відновити закриті частини, щоб отриманий код відповідав наведеному виходу програми.


\textbf{Завдання 18-19} 

Кілька запитань за завданнями для самостійної роботи.
\vspace{10mm}
Скоріш за все, що завдань буде трохи менше, а деякі попередні будуть переформульовані в стилі двох останніх.

%enditem

\end {large}}

%end

\newpage
%begin

{{
\begin {large}

\newpage
1. Що буде надруковано за виконання?

\begin{verbatim}
print('R', end=' ')
g=('a','b','c','d','e',0,1,2,3)
print(g[-4])
\end{verbatim}

%enditem


\newpage
2. Що буде надруковано за виконання?

\begin{verbatim}
print('R', end=' ')
e=(0,1,2,3,4,5,6,7,8,9)
print(e[10:8:3])
\end{verbatim}

%enditem



\newpage
3. Що буде надруковано за виконання?

\begin{verbatim}
print('R', end=' ')
f=['0','1','2','3','4']
f.insert(2, 16)
print(f)
\end{verbatim}

%enditem



\newpage
4. Що буде надруковано за виконання?

\begin{verbatim}
print('R', end=' ')
a=[10,11,12,13,14]
a[4:-2]=(15)
print(a)
\end{verbatim}

%enditem



\newpage
5. Укажіть ТИП значення виразу.\par
\begin{verbatim}
[7,2,3,4]
\end{verbatim}
%enditem


\newpage
6. Що буде надруковано за виконання?\par
\begin{verbatim}
try:
    s = {38:5, 41:5, 80:4, 46:4, 80:2}
    s[80] = 5
    for x in s.items():
        print(*x, sep=' ', end=' ')
except:
    print('ERR')
\end{verbatim}
%enditem


\newpage
7. Нехай змінна \kw{s} позначає дуже довгу послідовність, по якій можна ітерувати. 

Написати вираз-генератор, що задає підпослідовність її елементів, що не є числами типу \kw{int}.

Обмеження: усі елементи шуканої підпослідовності одночасно в пам'яті розміщені бути не можуть. 
%enditem


\newpage
8. Що буде надруковано за виконання?
%Оригинал был утерян. Требует отладки!
\begin{verbatim}
b = 0

class B:
    b = 1
    
    def __init__(self):
        self.b = 2
        b = 3
        B.b = 4

obj = B()
obj.b = 5
try:
    print(b, end=' ')
    print(obj.b, end=' ')
    print(B.b, end=' ')
    print(obj.b, end=' ')
    print(obj._B__b)
except:
    print('ERR')
\end{verbatim}
%enditem


\newpage
9. Що буде надруковано за виконання?
\begin{verbatim}
class A0:
    def b(self): print(1, end=' ')
    def h(self): print(2, end=' ')

class A1(A0):
    def b(self): print(3, end=' ')

try:
    A1().h()
    A0().b()
except:
    print('ERR')
\end{verbatim}
%enditem


\newpage
10. Що буде надруковано за виконання?
\begin{verbatim}
class A:
    c = 1

    def __init__(self):
        b = 2
        b = 3

class B(A):
    c = 4
    a = 5

    def __init__(self):
        super().__init__()
        self.a = 6

try:
    obj = B()
    obj.c = obj.a + 10
    B.c = 13
except: print('EE', end=' ')

try: print(obj.c, end= ' ')
except: print('E1', end=' ')

try: print(A.c, end= ' ')
except: print('E2', end=' ')

try: print(B.c, end= ' ')
except: print('E3', end=' ')
\end{verbatim}
%enditem


\newpage
11. Змінні \kw{next} та \kw{prev} вузла списку позначають наступний та попередній за ним вузол відповідно. Починаючи з голови, у список записано послідовні цілі числа від~$-34$ до~$30$. Нехай змінна \kw{p} позначає вузол, що зберігає значення~$2$.
Яке значення зберігається у вузлі \kw{p.prev.next.next.next.prev} ?
%enditem


\newpage
12. 
\begin{center}
	\adjustimage{max size={0.8\linewidth}{0.8\paperheight}}
	{BinTree/trees_exp/66.png}
\end{center}
Указати послідовність міток вершин дерева, зображеного на рис., що утвориться в результаті обходу дерева в оберненому порядку. Мітки записувати через пробіл.\par Алгоритм починає роботу з кореня (на рис. це верхня вершина).
%enditem


\newpage
13. 
\begin{center}
	\adjustimage{max size={0.9\linewidth}{0.9\paperheight}}
	{Graphs/AdjList/193.png}
\end{center}
Для графа, зображеного на рис., вказати список суміжності вершини 1.\par
Формат оформлення відповіді:\par
список суміжності виписувати через пробіл на одному рядку, якщо пробілів десь буде більше одного - не важливо, упорядкування вершин у списку також значення не має.
%enditem


\newpage
14. Для графа на рисунку вказати послідовність, в якій алгоритм пошуку в глибину буде відмічати відвідані вершини, якщо списки суміжності вершин переглядаються алгоритмом за зростанням міток вершин та вершини відмічаються на першому вход\-женні у вершину.\par
Пошук розпочинається з вершини 4.\par

\begin{center}
	\adjustimage{max size={0.9\linewidth}{0.9\paperheight}}
	{Graphs/Traversals/231.png}
\end{center}
%enditem


\newpage
15. Для графа на рисунку з вершини 3 запущено алгоритм пошуку в ширину, що будує кістякове дерево. Списки суміжності вершин переглядаються алгоритмом за зростанням міток вершин.\par
\begin{center}
	\adjustimage{max size={0.9\linewidth}{0.9\paperheight}}
	{Graphs/Traversals/222.png}
\end{center}
\par Які з наведених ребер входять до складу отриманого кістякового дерева?\par
1) (2, 5)\par
2) (0, 1)\par
3) (3, 5)\par
4) (1, 2)\par
5) (5, 6)\par
6) (0, 4)\par
{\vskip 15pt} Номери відповідних ребер записати за зростанням через пробіл.


%enditem


\newpage
16. 

Рядки журналу через двокрапку містять ім'я користувача (ідентифікатор), час події,тип події, код події, наприклад
\begin{verbatim}
s="username : 2022-05-05 13:26 : ERROR : 404"
\end{verbatim}
у коді 
\begin{verbatim}
res = re.fullmatch('???', s) 
s = ''
code_str = ??? 
\end{verbatim}
чим треба замінити кожен з ???, щоб значенням code\_str був код події? 


\vspace{10mm}
\textbf{Завдання 17} 

Задано програму, що містить код 1-2 класів, можливе успадкування, та клієнтський код.
У наведеному коді невеликі частини закриті. 
Також наведено вихід програми.
Необхідно відновити закриті частини, щоб отриманий код відповідав наведеному виходу програми.


\textbf{Завдання 18-19} 

Кілька запитань за завданнями для самостійної роботи.
\vspace{10mm}
Скоріш за все, що завдань буде трохи менше, а деякі попередні будуть переформульовані в стилі двох останніх.

%enditem

\end {large}}

%end

\newpage
%begin

{{
\begin {large}

\newpage
1. Що буде надруковано за виконання?

\begin{verbatim}
print('R', end=' ')
c=('a','b','c','d','e',0,1,2,4)
print(c[-11])
\end{verbatim}

%enditem


\newpage
2. Що буде надруковано за виконання?

\begin{verbatim}
print('R', end=' ')
b='0123456789'
print(b[-9::-2])
\end{verbatim}

%enditem



\newpage
3. Що буде надруковано за виконання?

\begin{verbatim}
print('R', end=' ')
b=[10,11,12,13,14]
print(b.pop(), end=' ')
print(b)
\end{verbatim}

%enditem



\newpage
4. Що буде надруковано за виконання?

\begin{verbatim}
print('R', end=' ')
b=['0','1','2','3','4']
b[:-3]=[1,2]
print(b)
\end{verbatim}

%enditem



\newpage
5. Укажіть ТИП значення виразу.\par
\begin{verbatim}
{x:x for x in range(7)}
\end{verbatim}
%enditem


\newpage
6. Що буде надруковано за виконання?\par
\begin{verbatim}
try:
    d = {62:7, 88:9, 78:9, 62:9}
    d[47] = 2
    for x, y in d.items():
        print(x, y, sep=' ', end=' ')
except:
    print('ERR')
\end{verbatim}
%enditem


\newpage
7. Задано словник \kw{d} з цілими ключами та значеннями. Записати ВИРАЗ, значенням якого є множина, що складається  з тих ключів словника, яким відповідають непарні значення.
%enditem


\newpage
8. Що буде надруковано за виконання?
%Оригинал был утерян. Требует отладки!
\begin{verbatim}
_c = 0

class C:
    _c = 1
    
    def __init__(self):
        self._c = 2
        _c = 3
        C._c = 4

obj = C()
C._c = 5
try:
    print(_c, end=' ')
    print(obj._c, end=' ')
    print(C._c, end=' ')
    print(obj._c, end=' ')
    print(obj._C__c)
except:
    print('ERR')
\end{verbatim}
%enditem


\newpage
9. Що буде надруковано за виконання?
\begin{verbatim}
class B0:
    def a(self): print(1, end=' ')
    def g(self): print(2, end=' ')

class B1(B0):
    def a(self): print(3, end=' ')

try:
    B1().g()
    B1().a()
except:
    print('ERR')
\end{verbatim}
%enditem


\newpage
10. Що буде надруковано за виконання?
\begin{verbatim}
class A:
    c = 1

    def __f1(self): print(2, end=' ')
    def _f2(self): print(3, end=' ')

    def main(self):
        try: print(self.c, end=' ')
        except: print('E1', end=' ')

        try: self.__f1()
        except: print('E2', end=' ')

        try: self._f2()
        except: print('E3', end=' ')


class B(A):
    c = 4

    def __f1(self): print(5, end=' ')
    def _f2(self): print(6, end=' ')

try: B().main()
except: print('E4', end=' ')
\end{verbatim}
%enditem


\newpage
11. Змінні \kw{next} та \kw{prev} вузла списку позначають наступний та попередній за ним вузол відповідно. Починаючи з голови, у список записано послідовні цілі числа від~$-38$ до~$48$. Нехай змінна \kw{p} позначає вузол, що зберігає значення~$6$.
Яке значення зберігається у вузлі \kw{p.next.prev.prev.prev.next} ?
%enditem


\newpage
12. 
\begin{center}
	\adjustimage{max size={0.8\linewidth}{0.8\paperheight}}
	{BinTree/trees_exp/143.png}
\end{center}
Указати послідовність міток вершин дерева, зображеного на рис., що утвориться в результаті обходу дерева в прямому порядку. Мітки записувати через пробіл.\par Алгоритм починає роботу з кореня (на рис. це верхня вершина).
%enditem


\newpage
13. 
\begin{center}
	\adjustimage{max size={0.9\linewidth}{0.9\paperheight}}
	{Graphs/AdjList/342.png}
\end{center}
Для графа, зображеного на рис., вказати список суміжності вершини 3.\par
Формат оформлення відповіді:\par
список суміжності виписувати через пробіл на одному рядку, якщо пробілів десь буде більше одного - не важливо, упорядкування вершин у списку також значення не має.
%enditem


\newpage
14. Для графа на рисунку з вершини 3 запущено алгоритм пошуку в глибину, що будує кістякове дерево. Списки суміжності вершин переглядаються алгоритмом за зростанням міток вершин.\par
\begin{center}
	\adjustimage{max size={0.9\linewidth}{0.9\paperheight}}
	{Graphs/Traversals/212.png}
\end{center}
\par Які з наведених ребер входять до складу отриманого кістякового дерева?\par
1) (0, 4)\par
2) (0, 5)\par
3) (0, 6)\par
4) (1, 5)\par
5) (1, 6)\par
6) (3, 6)\par
{\vskip 15pt} Номери відповідних ребер записати за зростанням через пробіл.

%enditem


\newpage
15. Для графа на рисунку з вершини 4 запущено алгоритм пошуку в ширину, що будує кістякове дерево. Списки суміжності вершин переглядаються алгоритмом за зростанням міток вершин.\par
\begin{center}
	\adjustimage{max size={0.9\linewidth}{0.9\paperheight}}
	{Graphs/Traversals/38.png}
\end{center}
\par Які з наведених ребер входять до складу отриманого кістякового дерева?\par
1) (1, 5)\par
2) (1, 3)\par
3) (4, 5)\par
4) (0, 2)\par
5) (2, 6)\par
6) (3, 5)\par
{\vskip 15pt} Номери відповідних ребер записати за зростанням через пробіл.


%enditem


\newpage
16. 

Рядки журналу через двокрапку містять ім'я користувача (ідентифікатор), час події,тип події, код події, наприклад
\begin{verbatim}
s="username : 2022-05-05 13:26 : ERROR : 404"
\end{verbatim}
у коді 
\begin{verbatim}
res = re.fullmatch('???', s) 
s = ''
code_str = ??? 
\end{verbatim}
чим треба замінити кожен з ???, щоб значенням code\_str був код події? 


\vspace{10mm}
\textbf{Завдання 17} 

Задано програму, що містить код 1-2 класів, можливе успадкування, та клієнтський код.
У наведеному коді невеликі частини закриті. 
Також наведено вихід програми.
Необхідно відновити закриті частини, щоб отриманий код відповідав наведеному виходу програми.


\textbf{Завдання 18-19} 

Кілька запитань за завданнями для самостійної роботи.
\vspace{10mm}
Скоріш за все, що завдань буде трохи менше, а деякі попередні будуть переформульовані в стилі двох останніх.

%enditem

\end {large}}

%end

\newpage
\end{document}
